
%% Produced by Scientific Word
%% Version 2019053000
%% Created Wed Mar 31 2021 15:11:19 GMT-0400 (Eastern Daylight Time)
%% Last revised Wed Mar 31 2021 15:52:07 GMT-0400 (Eastern Daylight Time)

\documentclass[aps, preprint]{revtex4}

%% preamble

\usepackage{amssymb,amsmath,mathtools,xcolor,graphicx,xspace,colortbl,ragged2e,rotating} % 
\usepackage{amsfonts}  %% 010
\usepackage{amsmath}  %% 010
\usepackage{amssymb}  %% 100
\usepackage{graphicx}  %% 100
\graphicspath{{../graphics/}{../tcache/}{../gcache/}}
\DeclareGraphicsExtensions{.pdf,.eps,.ps,.png,.jpg,.jpeg}
\setcounter{MaxMatrixCols}{10}
\newtheorem{theorem}{Theorem}
\newtheorem{acknowledgement}[theorem]{Acknowledgement}
\newtheorem{algorithm}[theorem]{Algorithm}
\newtheorem{axiom}[theorem]{Axiom}
\newtheorem{claim}[theorem]{Claim}
\newtheorem{conclusion}[theorem]{Conclusion}
\newtheorem{condition}[theorem]{Condition}
\newtheorem{conjecture}[theorem]{Conjecture}
\newtheorem{corollary}[theorem]{Corollary}
\newtheorem{criterion}[theorem]{Criterion}
\newtheorem{definition}[theorem]{Definition}
\newtheorem{example}[theorem]{Example}
\newtheorem{exercise}[theorem]{Exercise}
\newtheorem{lemma}[theorem]{Lemma}
\newtheorem{notation}[theorem]{Notation}
\newtheorem{problem}[theorem]{Problem}
\newtheorem{proposition}[theorem]{Proposition}
\newtheorem{remark}[theorem]{Remark}
\newtheorem{solution}[theorem]{Solution}
\newtheorem{summary}[theorem]{Summary}
\newenvironment{proof}[1][Proof]{\noindent\textbf{#1.} }{\ \rule{0.5em}{0.5em}}
\input{tcilatex}
\begin{document}
%TCIMACRO{\TeXButton{preprint}{\preprint{}}}%
%BeginExpansion
\preprint{}
%EndExpansion
\title{Clifford and Geometric Algebra} 
\author{Brian Weiner
\and } 
%TCIMACRO{\TeXButton{affiliation}{\affiliation{PSU}}}%
%BeginExpansion
\affiliation{PSU}
%EndExpansion



\begin{abstract}
\end{abstract}
\date{10/14/2012}
%TCIMACRO{\TeXButton{startpage}{\startpage{1}}}%
%BeginExpansion
\startpage{1}
%EndExpansion
\maketitle
\tableofcontents 

\section{Canonical Anti Commutation Relationship (CAR)}
The creation operators $\left \{\left .a_{j}^{ \dag } \equiv a^{ \dag } \left (\varphi _{j}\right ) \equiv a \left (\varphi _{j}\right )^{ \dag }\right \vert 1 \leq j \leq r\right \}$ belong to $\mathcal{F} \equiv \mathcal{B} \left (\mathcal{H}_{\mathcal{F}}\right )\text{,}$ where $\mathcal{H}_{\mathcal{F}}$ is a \emph{Hilbert Exterior Algebra} given by
\begin{equation}\mathcal{H}_{\mathcal{F}} \equiv  \tbigwedge \left (\mathcal{H}^{1}\right ) =\tbigoplus ???\mathcal{H}^{n}\text{.}
\end{equation}

The identity operator $I_{\mathcal{F}} \equiv I_{\mathcal{H}_{\mathcal{F}}} \in \mathcal{B} \left (\mathcal{H}_{\mathcal{F}}\right )$ can be expanded as
\begin{equation}I_{\mathcal{H}_{\mathcal{F}}} =\sum _{0 \leq n \leq r}P_{\mathcal{H}^{n}}\text{,}
\end{equation}where $\left \{\left .P_{\mathcal{H}^{n}}\right \vert 0 \leq n \leq r\right \}$ are orthogonal projectors onto the subspaces $\left \{\left . \tbigwedge ^{n}\mathcal{H}^{1} \equiv \mathcal{H}^{n}\right \vert 0 \leq n \leq r\right \}\text{.}$ The dimensions of the subspaces are given by
\begin{equation}\dim  \left \{\mathcal{H}^{n}\right \} =\ensuremath{\operatorname*{Tr}}_{\mathcal{H}_{\mathcal{F}}} \left \{P_{\mathcal{H}^{n}}\right \} =\ensuremath{\operatorname*{Tr}}_{\mathcal{H}^{n}} \left \{P_{\mathcal{H}^{n}}\right \} =\binom{r}{n} ;0 \leq n \leq r\text{,}
\end{equation}leading to
\begin{equation}\ensuremath{\operatorname*{Tr}}_{\mathcal{H}_{\mathcal{F}}} \left \{I_{\mathcal{H}_{\mathcal{F}}}\right \} =\sum _{0 \leq n \leq r}\binom{r}{n} =2^{r}\text{.}
\end{equation}The trace $\ensuremath{\operatorname*{Tr}}_{\mathcal{H}_{\mathcal{F}}}$ can be restricted to the subspaces leading to \emph{partial traces} i.e.
\begin{equation}\ensuremath{\operatorname*{Tr}}_{\mathcal{H}^{n}} =\left .\ensuremath{\operatorname*{Tr}}_{\mathcal{H}_{\mathcal{F}}}\right \vert _{\mathcal{H}^{n}}
\end{equation}giving
\begin{equation}\ensuremath{\operatorname*{Tr}}_{\mathcal{H}^{n}} \left \{P_{\mathcal{H}^{m}}\right \} =\binom{r}{n} \delta _{n m} ;0 \leq n \leq r\text{.}
\end{equation}


\begin{notation}
Unless stated otherwise we will use the identification $\ensuremath{\operatorname*{Tr}} \equiv \ensuremath{\operatorname*{Tr}}_{\mathcal{H}_{\mathcal{F}}}\text{.}$ 
\end{notation}

The creation operators $\left \{\left .a_{j}^{ \dag }\right \vert 1 \leq j \leq r\right \}$ are defined by their action on a conb $\left \{\left .\vert \varphi _{j_{1}} \cdots  \varphi _{j_{n}} \rangle \right \vert 1 \leq j_{1} <\cdots  <j_{n} \leq r ;0 \leq n \leq r\right \}$ of $\mathcal{H}_{\mathcal{F}}$ by
\begin{equation}a_{j}^{ \dag }\vert \varphi _{j_{1}} \cdots  \varphi _{j_{n}} \rangle  =\vert \varphi _{j} \varphi _{j_{1}} \cdots  \varphi _{j_{n}} \rangle  . \label{CREAT}
\end{equation}and the annihilation operators $\left \{\left .a_{j}\right \vert 1 \leq j \leq r\right \}$ are defined as their adjoints i.e.
\begin{equation}a_{j} =\left (a_{j}^{ \dag }\right )^{ \dag } ,\,1 \leq j \leq r\text{.}
\end{equation}

The creation and annihilation operators satisfy the \emph{Canonical Anticommutation Relationships
(CAR)}
\label{CAR}
\begin{subequations}
\begin{align}\left \{a_{j} ,a_{k}^{ \dag }\right \} &  =\delta _{j k} I_{\mathcal{F}} \label{CAR1} \\
\left \{a_{j}^{ \dag } ,a_{k}^{ \dag }\right \} &  =\left \{a_{j} ,a_{k}\right \} =0 \label{CAR2} \\
1 &  \leq j \leq r\text{,}\end{align}where
\end{subequations}
\begin{equation}\left \{A ,B\right \} =A B +B A\text{,}
\end{equation}


\begin{itemize}
\item The operators $\left \{\left .a_{j}^{ \dag }\right \vert 1 \leq j \leq r\right \}$ are bounded so have bounded adjoints $\left \{\left .a_{j}\right \vert 1 \leq j \leq r\right \}$ and hence
\begin{equation}a_{j}^{ \dag } =\left (a_{j}\right )^{ \dag }\text{.}
\end{equation}

\item In particular it follows from the definition eq. $\left (\ref{CREAT}\right )$ and eq. $\left (\ref{CAR1}\right )$ that
\begin{align}a_{j}^{ \dag }\vert \phi  \rangle  &  =\vert \varphi _{j} \rangle  \nonumber  \\
a_{j}\vert \phi  \rangle  &  =0 \nonumber  \\
1 &  \leq j \leq r\text{.}\end{align}

This follows from
\begin{equation} \langle \left .\phi \right \vert a_{j}^{ \dag } a_{j} \phi  \rangle  = \langle \phi \left \vert \left (I_{\mathcal{F}} -a_{j} a_{j}^{ \dag }\right )\right .\phi  \rangle  =0\text{,}
\end{equation}which implies that
\begin{equation}a_{j}\vert \phi  \rangle  =0.
\end{equation}

\item $a_{j}^{ \dag }$ and $a_{j}$ are \emph{partial isometries}, (and generate the \emph{Fermion Fock algebra} $\mathcal{F} \equiv \mathcal{B} \left ( \wedge \left (\mathcal{H}^{1}\right )\right )$), as
\begin{equation}a_{j}^{ \dag } a_{j}\text{and}a_{j} a_{j}^{ \dag }
\end{equation}are projectors and equivalently
\begin{equation}a_{j} a_{j}^{\dag } a_{j} =a_{j} \left (1 -a_{j} a_{j}^{\dag }\right ) =a_{j}\text{.}
\end{equation}In this particular case
\begin{align}a_{j}^{\dag } a_{j} +a_{j} a_{j}^{\dag } &  =I_{\mathcal{F}}\text{,} \\
a_{j}^{\dag } a_{j} a_{j} a_{j}^{\dag } &  =0 \\
a_{j} a_{j}^{\dag } a_{j}^{\dag } a_{j} &  =0\text{,}\end{align}the projectors $a_{j}^{\dag } a_{j}$ and $a_{j} a_{j}^{\dag }$ are orthogonal as $a_{j}^{\dag } a_{j} a_{j} a_{j}^{\dag } =a_{j}^{\dag } a_{j} \left (1 -a_{j}^{\dag } a_{j}\right ) =a_{j}^{\dag } a_{j} -a_{j}^{\dag } a_{j} =0$ and are also infinite dimensional if $r =\infty $. 

\item The partial isometries $\left \{\left .a_{j}^{ \dag } ,a_{j}\right \vert 1 \leq j \leq r\text{}\right \}$ associated with the projectors $\left \{\left .a_{j}^{ \dag } a_{j} ,a_{j} a_{j}^{ \dag }\right \vert 1 \leq j \leq r\text{}\right \}$ generate the CAR\ algebra $\mathcal{F} =\mathcal{B} \left (\mathcal{H}_{\mathcal{F}}\right )\text{,}$ (the space of bounded operators acting in $\mathcal{H}_{\mathcal{F}})\text{,}$ which is a \emph{noncomutative polynomial Type }$I_{2^{r}} -w^{ \ast }$\emph{-algebra}. 

\item The trace of the projectors $\left \{\left .a_{j}^{ \dag } a_{j} ,a_{j} a_{j}^{ \dag }\right \vert 1 \leq j \leq r\text{}\right \}$ are
\begin{equation}\ensuremath{\operatorname*{Tr}}_{\mathcal{H}_{\mathcal{F}}} \left \{a_{j}^{\dag } a_{j} +a_{j} a_{j}^{\dag }\right \} =\ensuremath{\operatorname*{Tr}}_{\mathcal{H}_{\mathcal{F}}} \left \{I_{\mathcal{F}}\right \} =2^{r}
\end{equation}leading to
\begin{equation}\ensuremath{\operatorname*{Tr}}_{\mathcal{H}_{\mathcal{F}}} \left \{a_{j}^{\dag } a_{j}\right \} =\ensuremath{\operatorname*{Tr}}_{\mathcal{H}_{\mathcal{F}}} \left \{a_{j} a_{j}^{\dag }\right \} =2^{r -1}\text{.}
\end{equation}and
\begin{equation}\ensuremath{\operatorname*{Tr}}_{\mathcal{H}^{n}} \left \{a_{j}^{\dag } a_{j}\right \} =\ensuremath{\operatorname*{Tr}}_{\mathcal{H}^{n}} \left \{a_{j} a_{j}^{\dag }\right \} =\binom{r -1}{n}\text{.}
\end{equation}

\item The antisymmetry of the CAR gives
\begin{equation}\ensuremath{\operatorname*{Tr}}_{\mathcal{H}_{\mathcal{F}}} \left \{a_{j}^{\dag } a_{k}\right \} =\ensuremath{\operatorname*{Tr}}_{\mathcal{H}_{\mathcal{F}}} \left \{a_{j} a_{k}^{\dag }\right \} =0 ;\,1 \leq j \neq k \leq r\text{.}
\end{equation}\end{itemize}


\paragraph{General Field Operators}
One can set the operators $\left \{\left .a_{j}^{ \dag }\right \vert 1 \leq j \leq r\right \}$ as the values of the \emph{complex linear map}
\begin{equation}\Phi ^{ \dag } :\mathcal{H}^{1} \rightarrow \mathcal{B} \left (\mathcal{H}_{\mathcal{F}}\right )
\end{equation}i.e.
\begin{equation}a_{j}^{ \dag } =\Phi ^{ \dag } \left (\vert \varphi _{j} \rangle \right )\text{,}
\end{equation}which leads to a more general (basis free) form of the CAR
\begin{subequations}
\begin{align}\left \{\Phi  \left (\vert u \rangle \right ) ,\Phi ^{ \dag } \left (\vert v \rangle \right )\right \} &  = \langle \left .u\right \vert v \rangle  I_{\mathcal{F}} \\
\left \{\Phi ^{ \dag } \left (\vert u \rangle \right ) ,\Phi ^{ \dag } \left (\vert v \rangle \right )\right \} &  =\left \{\Phi  \left (\vert u \rangle \right ) ,\Phi  \left (\vert v \rangle \right )\right \} =0 \\
 \forall \vert u \rangle  ,\vert v \rangle  &  \in \mathcal{H}^{1}\text{.}\end{align}In this case
\end{subequations}
\begin{equation}\mathcal{H}_{\mathcal{F}} = \tbigwedge \left (\mathcal{H}^{1}\right )\text{,}
\end{equation}however this can be generalized by considering different indexing Hilbert spaces $\mathfrak{H}$ for a CAR\ map into a $c^{ \ast }$\emph{-algebra} $\mathcal{A}$ followed by a GNS representation induced by a state $f \in \mathcal{S} \left (\mathcal{A}\right )$ i.e.
\begin{equation}\Phi ^{ \dag } :\mathfrak{H} \rightarrow \mathcal{A} \rightarrow \Pi _{f} \left (\mathcal{A}\right ) \subseteq \mathcal{B} \left (\mathcal{H}_{f}\right )\text{.}
\end{equation}


\begin{remark}
As $\Pi _{f} \left (\mathcal{A}\right )$ is in general included in $\mathcal{B} \left (\mathcal{H}_{f}\right )$ it may \emph{not contain} any \emph{individual} Ket-Bra operators, though it must contain linear
combinations or limits of combinations of these operators (as all operators in $\mathcal{B} \left (\mathcal{H}_{f}\right )$ can be expressed in this manner). 
\end{remark}

A map $T \in G L \left (\mathcal{B} \left (\mathcal{H}_{\mathcal{F}}\right )\right ) \subset \mathcal{B} \left (\mathcal{B} \left (\mathcal{H}_{\mathcal{F}}\right )\right )$ is a $c^{ \ast }$\emph{-algebra automorphism} if $ \forall X ,Y \in \mathcal{B} \left (\mathcal{H}_{\mathcal{F}}\right )$
\begin{align}T \left (X Y\right ) &  =T \left (X\right ) T \left (Y\right ) \\
T \left (X^{ \dag }\right ) &  =T_{\mathcal{F}} \left (X\right )^{ \dag } \\
T^{ -1} T &  =I d\text{.}\end{align}The norm of $T$ is defined by
\begin{equation}\left \Vert T\right \Vert  =\sup_{\left \Vert X\right \Vert _{\mathcal{B} \left (\mathcal{H}_{\mathcal{F}}\right )} =1}\left \Vert T \left (X\right )\right \Vert _{\mathcal{B} \left (\mathcal{H}_{\mathcal{F}}\right )}\text{,}
\end{equation}and $T$ is an isometric isomorphism if
\begin{equation}\left \Vert T \left (X\right )\right \Vert _{\mathcal{B} \left (\mathcal{H}_{\mathcal{F}}\right )} =\left \Vert X\right \Vert _{\mathcal{B} \left (\mathcal{H}_{\mathcal{F}}\right )}\text{.}
\end{equation}In particular one can see that
\label{CAR}
\begin{subequations}
\begin{align}\left \{T \left (a_{j}\right ) ,T \left (a_{k}^{ \dag }\right )\right \} &  =\left \{T \left (a_{j}\right ) ,T \left (a_{k}\right )^{ \dag }\right \} =T \left (\left \{a_{j} ,a_{k}^{ \dag }\right \}\right ) =\delta _{j k} T \left (I_{\mathcal{F}}\right ) \\
T \left (\left \{a_{j}^{ \dag } ,a_{k}^{ \dag }\right \}\right ) &  =T \left (\left \{a_{j} ,a_{k}\right \}\right ) =0 \\
1 &  \leq j ,k \leq r\text{.}\end{align}One can note that the norms
\end{subequations}
\begin{equation}\left \Vert a_{j}\right \Vert _{\mathcal{B} \left (\mathcal{H}_{\mathcal{F}}\right )} =\left \Vert a_{j}^{ \dag }\right \Vert _{\mathcal{B} \left (\mathcal{H}_{\mathcal{F}}\right )} =1.
\end{equation}Obviously
\begin{equation}T \left (I_{\mathcal{F}}\right ) \in \mathcal{Z} \left (\mathcal{F}\right )\text{,}
\end{equation}where $\mathcal{Z} \left (\mathcal{F}\right )$ is the centre of $\mathcal{F}$. In this case $\mathcal{F}$ is a simple algebra thus
\begin{equation}T \left (I_{\mathcal{F}}\right ) =z I_{\mathcal{F}} ,\text{\quad }z \in \mathbb{C}\text{.}
\end{equation}


\begin{remark}
If $T$ is a $c^{ \ast }$\emph{-algebra automorphism with the properties }
\begin{equation}T :\mathcal{F}_{1} \rightarrow \mathcal{F}_{1}
\end{equation}and
\begin{equation}T \left (I_{\mathcal{F}}\right ) =I_{\mathcal{F}}
\end{equation}then $T$ is a \emph{canonical transformation}. 
\end{remark}


\begin{remark}
One can generalize some of the preceding to non invertible elements of $\mathcal{B} \left (\mathcal{B} \left (\mathcal{H}_{\mathcal{F}}\right )\right )$ i.e. to general endomorphisms of $\mathcal{B} \left (\mathcal{H}_{\mathcal{F}}\right )\text{.}$ 
\end{remark}


\begin{notation}
The set $\left \{\mathbf{a}^{ \dag } ,\mathbf{a}\right \}$ is an example of an \emph{Adjoint Basis }of $\mathcal{F}_{1}$ 
\end{notation}


\begin{notation}
The set $\left \{\mathbf{x}_{1} ,\mathbf{x}_{2}\right \}\text{,}$ (to be defined subsequently), is an example of a \emph{Self Adjoint Basis }of $\mathcal{F}_{1}$. 
\end{notation}


\begin{notation}
A basis of $\mathcal{F}_{1}$ that is not \emph{Adjoint} or \emph{Self Adjoint }is called a \emph{General Basis}
$\mathcal{F}_{1}\text{,}$ unless otherwise qualified with another property$\text{.}$ 
\end{notation}

\section{Fat Symbol Notation}
We first note that matrices $m \times n$ matrices with entries from the \emph{Clifford Algebra / }$c^{ \ast }$\emph{-Algebra }$\mathcal{F}$ are in fact tensor products as \emph{\ }
\begin{equation}M \left (m ,n ,\mathcal{F}\right ) =M \left (m ,n ,\mathbb{C}\right ) \otimes \mathcal{F}\text{,}
\end{equation}where
\begin{align}m &  =\text{number of rows} \\
n &  =\text{number of columns.}\end{align}


\begin{notation}
Column vectors (Contravariant vectors) formed from elements of $\mathcal{F}$ are arrays
\begin{equation}\mathcal{F}^{m} =M \left (m ,1 ,\mathcal{F}\right )\text{.}
\end{equation}In finite dimensions if there is an inner product defined in $\mathcal{F}^{m}$ the dual space of $\mathcal{F}^{m}$ is isomorphic to the space of row vectors (Covariant vectors)
\begin{equation}\left (\mathcal{F}^{m}\right )^{d} =M \left (1 ,m ,\mathcal{F}\right )\text{.}
\end{equation}
\end{notation}


\begin{remark}
Any contravariant vector belonging $v \in V\text{,}$ a complex vector space of dimension $n\text{,}$ can be expanded in terms of a product of a contravariant vector $\mathbf{v} \in \mathbb{C}^{n}$ and a super covariant vector of basis vectors $\mathfrak{e}^{\#} \in V^{n}$ as
\begin{equation}v =\left (\begin{array}{ccc}\mathfrak{e}_{1} & \cdots  & \mathfrak{e}_{n}\end{array}\right ) \left (\begin{array}{c}v_{1} \\
\vdots  \\
v_{n}\end{array}\right ) =\left (\begin{array}{c}\mathfrak{e}_{1} \\
\vdots  \\
\mathfrak{e}_{n}\end{array}\right )^{\#} \left (\begin{array}{c}v_{1} \\
\vdots  \\
v_{n}\end{array}\right ) =\mathfrak{e}^{\#} \left (\begin{array}{c}v_{1} \\
\vdots  \\
v_{n}\end{array}\right ) =\mathfrak{e}^{\#} \mathbf{v} ;\,e_{j} ,1 \leq j \leq n\text{.}
\end{equation}
\end{remark}


\begin{example}
Let $\mathcal{F} =\mathcal{B} \left (\mathcal{H}_{\mathcal{F}}\right )$ and consider a Ket-Bra basis $\left \{\left .\vert \Psi _{j} \rangle  \langle \Psi _{k}\vert \right \vert 1 \leq j ,k \leq 2^{r}\right \}$ then
\begin{align}\mathcal{F}^{2^{2 r}} &  = & M \left (2^{2 r} ,1 ,\mathcal{B} \left (\mathcal{H}_{\mathcal{F}}\right )\right ) \\
\left (\mathcal{F}^{2^{2 r}}\right )^{d} &  = & M \left (1 ,2^{2 r} ,\mathcal{B} \left (\mathcal{H}_{\mathcal{F}}\right )\right )\text{.}\end{align}Any operator $X \in \mathcal{B} \left (\mathcal{H}_{\mathcal{F}}\right )$ can be expanded as
\begin{equation}X =\left (\begin{array}{ccc}\vert \Psi _{1} \rangle  \langle \Psi _{1}\vert  & \cdots  & \vert \Psi _{2^{r}} \rangle  \langle \Psi _{2^{r}}\vert \end{array}\right ) \left (\begin{array}{c}X_{1 1} \\
\vdots  \\
X_{2^{r} 2^{r}}\end{array}\right ) =\left (\vert \mathbf{\Psi } \rangle  \langle \mathbf{\Psi }\vert \right ) \mathbb{X} =\left (\begin{array}{c}\vdots  \\
\vert \Psi _{j} \rangle  \langle \Psi _{k}\vert  \\
\vdots \end{array}\right )^{\#} \mathbb{X}\text{,}
\end{equation}where $\mathbb{X} \in \mathbb{C}^{2^{2 r}}\text{.}$ 
\end{example}


\begin{example}
The action of the dual space $\mathcal{B} \left (\mathcal{H}_{\mathcal{F}}\right )$ on the space $\mathcal{B}_{1} \left (\mathcal{H}_{\mathcal{F}}\right )$ is given by
\begin{align}f_{Y} \left (X\right ) &  = & \ensuremath{\operatorname*{Tr}} \left \{Y^{ \dag } X\right \} =\left (\begin{array}{ccc}\overline{Y_{1 1}} & \cdots  & \overline{Y_{2^{r} 2^{r}}}\end{array}\right ) \ensuremath{\operatorname*{Tr}} \left \{\left (\begin{array}{c}\vert \Psi _{1} \rangle  \langle \Psi _{1}\vert ^{ \dag } \\
\vdots  \\
\vert \Psi _{2^{r}} \rangle  \langle \Psi _{2^{r}}\vert ^{ \dag }\end{array}\right ) \left (\begin{array}{ccc}\vert \Psi _{1} \rangle  \langle \Psi _{1}\vert  & \cdots  & \vert \Psi _{2^{r}} \rangle  \langle \Psi _{2^{r}}\vert \end{array}\right )\right \} \left (\begin{array}{c}X_{1 1} \\
\vdots  \\
X_{2^{r} 2^{r}}\end{array}\right ) \\
 &  = & \left (\begin{array}{ccc}\overline{Y_{1 1}} & \cdots  & \overline{Y_{2^{r} 2^{r}}}\end{array}\right ) \ensuremath{\operatorname*{Tr}} \left \{\left (\begin{array}{c}\vert \Psi _{1} \rangle  \langle \Psi _{1}\vert  \\
\vdots  \\
\vert \Psi _{2^{r}} \rangle  \langle \Psi _{2^{r}}\vert \end{array}\right )^{ \dag } \left (\begin{array}{ccc}\vert \Psi _{1} \rangle  \langle \Psi _{1}\vert  & \cdots  & \vert \Psi _{2^{r}} \rangle  \langle \Psi _{2^{r}}\vert \end{array}\right )\right \} \left (\begin{array}{c}X_{1 1} \\
\vdots  \\
X_{2^{r} 2^{r}}\end{array}\right ) \\
 &  = & \left (\begin{array}{c}Y_{1 1} \\
\vdots  \\
Y_{2^{r} 2^{r}}\end{array}\right )^{t \overline{}} \ensuremath{\operatorname*{Tr}} \left \{\left (\begin{array}{ccc}\vert \Psi _{1} \rangle  \langle \Psi _{1}\vert  & \cdots  & \vert \Psi _{2^{r}} \rangle  \langle \Psi _{2^{r}}\vert \end{array}\right )^{t \dag } \left (\begin{array}{ccc}\vert \Psi _{1} \rangle  \langle \Psi _{1}\vert  & \cdots  & \vert \Psi _{2^{r}} \rangle  \langle \Psi _{2^{r}}\vert \end{array}\right )\right \} \left (\begin{array}{c}X_{1 1} \\
\vdots  \\
X_{2^{r} 2^{r}}\end{array}\right ) \\
 &  = & \mathbb{Y}^{t \overline{}} \ensuremath{\operatorname*{Tr}} \left \{\left (\vert \mathbf{\Psi } \rangle  \langle \mathbf{\Psi }\vert \right )^{t \dag } \left (\vert \mathbf{\Psi } \rangle  \langle \mathbf{\Psi }\vert \right )\right \} \mathbb{X} =\mathbb{Y}^{ \dag } \ensuremath{\operatorname*{Tr}} \left \{\left (\vert \mathbf{\Psi } \rangle  \langle \mathbf{\Psi }\vert \right )^{t \dag } \left (\vert \mathbf{\Psi } \rangle  \langle \mathbf{\Psi }\vert \right )\right \} \mathbb{X}\text{.}\end{align}
\end{example}


\begin{example}
The tensor product $\mathcal{F}^{2^{2 r}} \otimes \left (\mathcal{F}^{2^{2 r}}\right )^{d}$ is given by
\begin{gather}\left (\begin{array}{ccc}\vert \Psi _{1} \rangle  \langle \Psi _{1}\vert  & \cdots  & \vert \Psi _{2^{r}} \rangle  \langle \Psi _{2^{r}}\vert \end{array}\right )^{t \dag } \left (\begin{array}{ccc}\vert \Psi _{1} \rangle  \langle \Psi _{1}\vert  & \cdots  & \vert \Psi _{2^{r}} \rangle  \langle \Psi _{2^{r}}\vert \end{array}\right ) \\
%TCIMACRO{\TeXButton{\shortparallel }{\shortparallel }}%
%BeginExpansion
\shortparallel 
%EndExpansion
 \\
\left (\begin{array}{c}\vert \Psi _{1} \rangle  \langle \Psi _{1}\vert  \\
\vdots  \\
\vert \Psi _{2^{r}} \rangle  \langle \Psi _{2^{r}}\vert \end{array}\right )^{ \dag } \otimes \left (\begin{array}{ccc}\vert \Psi _{1} \rangle  \langle \Psi _{1}\vert  & \cdots  & \vert \Psi _{2^{r}} \rangle  \langle \Psi _{2^{r}}\vert \end{array}\right ) =\left (\begin{array}{ccc}\vert \Psi _{1} \rangle  \langle \Psi _{1}\vert ^{ \dag }\vert \Psi _{1} \rangle  \langle \Psi _{1}\vert  & \cdots  & \vert \Psi _{1} \rangle  \langle \Psi _{1}\vert ^{ \dag }\vert \Psi _{2^{r}} \rangle  \langle \Psi _{2^{r}}\vert  \\
\vdots  & \ddots  & \vdots  \\
\vert \Psi _{1} \rangle  \langle \Psi _{2^{r}}\vert ^{ \dag }\vert \Psi _{1} \rangle  \langle \Psi _{1}\vert  & \cdots  & \vert \Psi _{1} \rangle  \langle \Psi _{2^{r}}\vert ^{ \dag }\vert \Psi _{2^{r}} \rangle  \langle \Psi _{2^{r}}\vert \end{array}\right )\text{,}\end{gather}which should be compared to the tensor product
\begin{gather}\mathbb{Y}^{t \overline{}} \otimes \mathbb{X} \\
%TCIMACRO{\TeXButton{\shortparallel }{\shortparallel }}%
%BeginExpansion
\shortparallel 
%EndExpansion
 \\
\left (\begin{array}{c}Y_{1 1} \\
\vdots  \\
Y_{2^{r} 2^{r}}\end{array}\right )^{t \overline{}} \otimes \left (\begin{array}{ccc}X_{1 1} & \cdots  & X_{2^{r} 2^{r}}\end{array}\right ) =\left (\begin{array}{ccc}\overline{Y_{1 1}} X_{1 1} & \cdots  & \overline{Y_{1 1}} X_{2^{r} 2^{r}} \\
\vdots  & \ddots  & \vdots  \\
\overline{Y_{2^{r} 2^{r}}} X_{1 1} & \cdots  & \overline{Y_{2^{r} 2^{r}}} X_{2^{r} 2^{r}}\end{array}\right )\text{.}\end{gather}
\end{example}

\subsection{Transpose}
\begin{definition}
If $\mathcal{F}$ is an\emph{\ Algebra that has a transpose operation "}$t$\emph{" defined }there are three types of transpose operations that can be defined in
\begin{equation}M \left (m ,n ,\mathcal{F}\right )\text{,}
\end{equation}
\end{definition}


\begin{enumerate}
\item The $\mathcal{F}$\emph{-Transpose}
\begin{equation}t :\mathcal{F} \rightarrow \mathcal{F}
\end{equation}

\item The \emph{Simple Transpose }$\#$
\begin{equation}\# :M \left (m ,n ,\mathcal{F}\right ) \rightarrow M \left (n ,m ,\mathcal{F}\right )\text{,}
\end{equation}is defined by
\begin{equation}\# :\overset{ \leftarrow \text{\quad \quad \quad \quad }n\text{\quad \quad \quad \quad } \rightarrow }{\left (\begin{array}{ccc}x_{1 1} & \cdots  & x_{1 m} \\
\vdots  & \, & \vdots  \\
x_{n 1} & \cdots  & x_{n m}\end{array}\right )}\, \begin{array}{c}\uparrow  \\
m \\
\downarrow \end{array} \longrightarrow \,\overset{ \leftarrow \text{\quad \quad \quad \quad }m\text{\quad \quad \quad \quad } \rightarrow }{\left (\begin{array}{ccc}x_{1 1} & \cdots  & x_{1 n} \\
\vdots  & \, & \vdots  \\
x_{m 1} & \cdots  & x_{m n}\end{array}\right )\, \begin{array}{c}\uparrow  \\
n \\
\downarrow \end{array}} ;x_{j k} \in \mathcal{F} ;1 \leq j \leq m ,1 \leq k \leq n\text{.}
\end{equation}

\item The \emph{Super Transpose }$\mathfrak{t}$
\begin{equation}\mathfrak{t} :M \left (m ,n ,\mathcal{F}\right ) \rightarrow M \left (n ,m ,\mathcal{F}\right )\text{,}
\end{equation}is defined by
\begin{equation}\mathfrak{t} :\overset{ \leftarrow \text{\quad \quad \quad \quad }n\text{\quad \quad \quad \quad } \rightarrow }{\left (\begin{array}{ccc}x_{1 1} & \cdots  & x_{1 m} \\
\vdots  & \, & \vdots  \\
x_{n 1} & \cdots  & x_{n m}\end{array}\right )}\, \begin{array}{c}\uparrow  \\
m \\
\downarrow \end{array} \longrightarrow \,\overset{ \leftarrow \text{\quad \quad \quad \quad }m\text{\quad \quad \quad \quad } \rightarrow }{\left (\begin{array}{ccc}x_{1 1}^{t} & \cdots  & x_{1 n}^{t} \\
\vdots  & \, & \vdots  \\
x_{m 1}^{t} & \cdots  & x_{m n}^{t}\end{array}\right )\, \begin{array}{c}\uparrow  \\
n \\
\downarrow \end{array}} ;x_{j k} \in \mathcal{F} ;1 \leq j \leq m ,1 \leq k \leq n\text{,}
\end{equation}where
\begin{equation}\mathfrak{t} =\# \circ t =t \circ \#\text{.}
\end{equation}\end{enumerate}



\begin{example}
Let $\left \{\mathbf{X}\right \}$ be \underline {\emph{ANY}} basis of $\mathcal{F}_{1}$ then operations on contravariant vectors
\begin{equation}\mathbf{X} =\left (\begin{array}{c}\mathfrak{X}_{1} \\
\vdots  \\
\mathfrak{X}_{2 r}\end{array}\right ) \in M \left (2 r ,1 ,\mathcal{F}_{1}\right ) =\mathcal{F}_{1}^{2 r}
\end{equation}are given by
\begin{align}\mathbf{X}^{t} &  =\left (\begin{array}{c}\mathfrak{X}_{1}^{t} \\
\vdots  \\
\mathfrak{X}_{2 r}^{t}\end{array}\right ) \in M \left (2 r ,1 ,\mathcal{F}_{1}\right ) =\mathcal{F}_{1}^{2 r}\text{,} \\
\mathbf{X}^{\#} &  =\left (\begin{array}{ccc}\mathfrak{X}_{1} & \cdots  & \mathfrak{X}_{2 r}\end{array}\right ) \in M \left (1 ,2 r ,\mathcal{F}_{1}\right ) =\left (\mathcal{F}_{1}^{2 r}\right )^{d}\text{,} \\
\mathbf{X}^{\mathfrak{t}} &  =\left (\begin{array}{ccc}\mathfrak{X}_{1}^{t} & \cdots  & \mathfrak{X}_{2 r}^{t}\end{array}\right ) \in M \left (1 ,2 r ,\mathcal{F}_{1}\right ) =\left (\mathcal{F}_{1}^{2 r}\right )^{d}\text{.}\end{align}and operations on covariant vectors
\begin{equation}\mathbf{X} =\left (\begin{array}{ccc}\mathfrak{X}_{1} & \cdots  & \mathfrak{X}_{2 r}\end{array}\right ) \in M \left (1 ,2 r ,\mathcal{F}_{1}\right ) =\left (\mathcal{F}_{1}^{2 r}\right )^{d}
\end{equation}are given by
\begin{align}\mathbf{X}^{t} &  =\left (\begin{array}{ccc}\mathfrak{X}_{1}^{t} & \cdots  & \mathfrak{X}_{2 r}^{t}\end{array}\right ) \in M \left (1 ,2 r ,\mathcal{F}_{1}\right ) =\left (\mathcal{F}_{1}^{2 r}\right )^{d} \\
\mathbf{X}^{\#} &  =\left (\begin{array}{c}\mathfrak{X}_{1} \\
\vdots  \\
\mathfrak{X}_{2 r}\end{array}\right ) \in M \left (2 r ,1 ,\mathcal{F}_{1}\right ) =\mathcal{F}_{1}^{2 r}\text{,} \\
\mathbf{X}^{\mathfrak{t}} &  =\left (\begin{array}{c}\mathfrak{X}_{1}^{t} \\
\vdots  \\
\mathfrak{X}_{2 r}^{t}\end{array}\right ) \in M \left (2 r ,1 ,\mathcal{F}_{1}\right ) =\mathcal{F}_{1}^{2 r}\text{.}\end{align}
\end{example}

One can expand $\mathcal{F}_{1}$ as the direct sum
\begin{equation}\mathcal{F}_{1} =\mathcal{F}_{1 1}  \oplus _{\mathbb{C}}\mathcal{F}_{1 2}
\end{equation}and let
\begin{align}\mathfrak{X}_{1 j} ,1 &  \leq j \leq r \in \mathcal{F}_{1 1} \\
\mathfrak{X}_{2 j} ,1 &  \leq j \leq r \in \mathcal{F}_{1 2}\end{align}and $\mathbf{X}_{1} ,\mathbf{X}_{2} \in M \left (1 ,r ,\mathcal{F}_{1}\right )$ then one obtains
\begin{equation}\left (\begin{array}{cc}\mathbf{X}_{1} & \mathbf{X}_{2}\end{array}\right )^{\mathfrak{t}} \left (\begin{array}{cc}\mathbf{X}_{1} & \mathbf{X}_{2}\end{array}\right ) =\left (\begin{array}{c}\mathbf{X}_{1}^{\mathfrak{t}} \\
\mathbf{X}_{2}^{\mathfrak{t}}\end{array}\right ) \left (\begin{array}{cc}\mathbf{X}_{1} & \mathbf{X}_{2}\end{array}\right ) =\left (\begin{array}{cc}\mathbf{X}_{1}^{\mathfrak{t}} \mathbf{X}_{1} & \mathbf{X}_{1}^{\mathfrak{t}} \mathbf{X}_{2} \\
\mathbf{X}_{2}^{\mathfrak{t}} \mathbf{X}_{1} & \mathbf{X}_{2}^{\mathfrak{t}} \mathbf{X}_{2}\end{array}\right ) \in M \left (2 r ,2 r ,\mathcal{F}_{\left (0 ,2\right )}\right )\text{,}
\end{equation}where
\begin{equation}\mathcal{F}_{\left (0 ,2\right )} =\mathcal{F}_{0}  \oplus _{\mathbb{C}}\mathcal{F}_{2}\text{.}
\end{equation}


\begin{remark}
The product $\mathbf{X}^{\mathbf{t}} \mathbf{Y}$ is a tensor product i.e.
\begin{equation}\mathbf{X}^{\mathfrak{t}} \mathbf{Y} \equiv \mathbf{X}^{\mathfrak{t}} \otimes \mathbf{Y}\text{.}
\end{equation}In particular if $\mathbf{Y} \in V^{d}$ and $\mathbf{X}^{\mathfrak{t}} \in V$ then $V \otimes V^{d} \cong \mathcal{B} \left (V\right )\text{.}$ 
\end{remark}


\begin{proposition}
Let $\left (\begin{array}{cc}\mathbf{X}_{1} & \mathbf{X}_{2}\end{array}\right ) \in M \left (2 r ,1 ,\mathcal{F}_{1}\right )\;$and consider the linear transformation
\begin{equation}T :\mathcal{F}_{1} \rightarrow \mathcal{F}_{1}
\end{equation}so that
\begin{equation}T \left (\begin{array}{cc}\mathbf{X}_{1} & \mathbf{X}_{2}\end{array}\right ) =\left (\begin{array}{cc}\mathbf{X}_{1} & \mathbf{X}_{2}\end{array}\right ) \mathbb{T}_{x}
\end{equation}then
\begin{equation}\left (T \left (\begin{array}{cc}\mathbf{X}_{1} & \mathbf{X}_{2}\end{array}\right )\right )^{\mathfrak{t}} =\left (\left (\begin{array}{cc}\mathbf{X}_{1} & \mathbf{X}_{2}\end{array}\right ) \mathbb{T}_{x}\right )^{\mathfrak{t}} =\mathbb{T}_{x}^{t} \left (\begin{array}{c}\mathbf{X}_{1}^{\mathfrak{t}} \\
\mathbf{X}_{2}^{\mathfrak{t}}\end{array}\right )\text{.}
\end{equation}


\begin{proof}
Consider the linear transformation
\begin{equation}T \left (\begin{array}{cc}\mathbf{X}_{1} & \mathbf{X}_{2}\end{array}\right ) =\left (\begin{array}{cc}T \mathbf{X}_{1} & T \mathbf{X}_{2}\end{array}\right ) =\left (\begin{array}{cc}\mathbf{X}_{1} & \mathbf{X}_{2}\end{array}\right ) \mathbb{T}_{x} =\left (\begin{array}{cc}\mathbf{X}_{1} \mathbb{T}_{x 1 1} +\mathbf{X}_{2} \mathbb{T}_{x 2 1} & \mathbf{X}_{1} \mathbb{T}_{x 1 2} +\mathbf{X}_{2} \mathbb{T}_{x 2 2}\end{array}\right ) =\left (\begin{array}{cc}\mathbf{Y}_{1} & \mathbf{Y}_{2}\end{array}\right )\text{,}
\end{equation}where
\begin{equation}\mathcal{R}_{x} \left (T\right ) \equiv \mathbb{T}_{x} =\left (\begin{array}{cc}\mathbb{T}_{x 1 1} & \mathbb{T}_{x 1 2} \\
\mathbb{T}_{x 2 1} & \mathbb{T}_{x 2 2}\end{array}\right )\text{.}
\end{equation}In more detail
\begin{align}T \left (\mathbf{X}_{1}\right ) &  =\mathbf{X}_{1} \mathbb{T}_{x 1 1} +\mathbf{X}_{2} \mathbb{T}_{x 2 1} \\
T \left (\mathbf{X}_{2}\right ) &  =\mathbf{X}_{1} \mathbb{T}_{x 1 2} +\mathbf{X}_{2} \mathbb{T}_{x 2 2}\text{,}\end{align}where
\begin{align}\sum _{j}\mathcal{X}_{1 j} T_{1 1 j k} &  =\sum _{j}T_{1 1 j k} \mathcal{X}_{1 j} =\sum _{j}\left (T_{1 1}^{t}\right )_{k j} \mathcal{X}_{1 j} \\
\sum _{j}\mathcal{X}_{2 j} T_{2 1 j k} &  =\sum _{j}T_{2 1 j k} \mathcal{X}_{2 j} =\sum _{j}\left (T_{2 1}^{t}\right )_{k j} \mathcal{X}_{2 j}\end{align}
\begin{align}\sum _{j}\mathcal{X}_{1 j} T_{1 2 j k} &  =\sum _{j}T_{1 2 j k} \mathcal{X}_{1 j} =\sum _{j}\left (T_{1 2}^{t}\right )_{k j} \mathcal{X}_{1 j} \\
\sum _{j}\mathcal{X}_{2 j} T_{2 2 j k} &  =\sum _{j}T_{2 2 j k} \mathcal{X}_{2 j} =\sum _{j}\left (T_{2 2}^{t}\right )_{k j} \mathcal{X}_{2 j}\end{align}hence
\begin{align}\left (\mathbf{X}_{1} \mathbb{T}_{x 1 1}\right )^{\mathfrak{t}} &  =\mathbb{T}_{x 1 1}^{t} \mathbf{X}_{1}^{\mathfrak{t}} \\
\left (\mathbf{X}_{2} \mathbb{T}_{x 2 1}\right )^{\mathfrak{t}} &  =\mathbb{T}_{x 2 1}^{t} \mathbf{X}_{2}^{\mathfrak{t}} \\
\left (\mathbf{X}_{1} \mathbb{T}_{x 1 2}\right )^{\mathfrak{t}} &  =\mathbb{T}_{x 1 2}^{t} \mathbf{X}_{1}^{\mathfrak{t}} \\
\left (\mathbf{X}_{2} \mathbb{T}_{x 2 2}\right )^{\mathfrak{t}} &  =\mathbb{T}_{x 2 2}^{t} \mathbf{X}_{2}^{\mathfrak{t}}\text{,}\end{align}
\begin{align}\left (T \left (\mathbf{X}_{1}\right )\right )^{\mathfrak{t}} &  =\left (\mathbf{X}_{1} \mathbb{T}_{x 1 1} +\mathbf{X}_{2} \mathbb{T}_{x 2 1}\right )^{\mathfrak{t}} =\mathbb{T}_{x 1 1}^{t} \mathbf{X}_{1}^{\mathfrak{t}} +\mathbb{T}_{x 2 1}^{t} \mathbf{X}_{2}^{\mathfrak{t}} \\
\left (T \left (\mathbf{X}_{2}\right )\right )^{\mathfrak{t}} &  =\left (\mathbf{X}_{1} \mathbb{T}_{x 1 2} +\mathbf{X}_{2} \mathbb{T}_{x 2 2}\right )^{\mathfrak{t}} =\mathbb{T}_{x 1 2}^{t} \mathbf{X}_{1}^{\mathfrak{t}} +\mathbb{T}_{x 2 2}^{t} \mathbf{X}_{2}^{\mathfrak{t}}\text{,}\end{align}and
\begin{equation}\left (T \left (\mathbf{X}\right )\right )^{\mathfrak{t}} =\left (\begin{array}{cc}\mathbb{T}_{x 1 1}^{t} & \mathbb{T}_{x 2 1}^{t} \\
\mathbb{T}_{x 1 2}^{t} & \mathbb{T}_{x 2 2}^{t}\end{array}\right ) \left (\begin{array}{c}\mathbf{X}_{1}^{\mathfrak{t}} \\
\mathbf{X}_{2}^{\mathfrak{t}}\end{array}\right ) =\left (\begin{array}{cc}\mathbb{T}_{x 1 1} & \mathbb{T}_{x 1 2} \\
\mathbb{T}_{x 2 1} & \mathbb{T}_{x 2 2}\end{array}\right )^{t} \left (\begin{array}{c}\mathbf{X}_{1}^{\mathfrak{t}} \\
\mathbf{X}_{2}^{\mathfrak{t}}\end{array}\right ) =\left (\left (\begin{array}{cc}\mathbf{X}_{1} & \mathbf{X}_{2}\end{array}\right ) \left (\begin{array}{cc}\mathbb{T}_{x 1 1} & \mathbb{T}_{x 1 2} \\
\mathbb{T}_{x 2 1} & \mathbb{T}_{x 2 2}\end{array}\right )\right )^{\mathfrak{t}}\text{.}
\end{equation}This gives
\begin{align} & \left (\begin{array}{c}\mathbf{Y}_{1}^{\mathfrak{t}} \\
\mathbf{Y}_{2}^{\mathfrak{t}}\end{array}\right ) \left (\begin{array}{cc}\mathbf{Y}_{1} & \mathbf{Y}_{2}\end{array}\right ) =\left (\left (\begin{array}{cc}\mathbf{X}_{1} & \mathbf{X}_{2}\end{array}\right ) \mathbb{T}_{x}\right )^{\mathfrak{t}} \left (\begin{array}{cc}\mathbf{X}_{1} & \mathbf{X}_{2}\end{array}\right ) \mathbb{T}_{x} =\mathbb{T}_{x}^{t} \left (\begin{array}{cc}\mathbf{X}_{1} & \mathbf{X}_{2}\end{array}\right )^{\mathfrak{t}} \left (\begin{array}{cc}\mathbf{X}_{1} & \mathbf{X}_{2}\end{array}\right ) \mathbb{T}_{x} \nonumber  \\
 &  =\mathbb{T}_{x}^{t} \left (\begin{array}{cc}\mathbf{X}_{1}^{\mathfrak{t}} \mathbf{X}_{1} & \mathbf{X}_{1}^{\mathfrak{t}} \mathbf{X}_{2} \\
\mathbf{X}_{2}^{\mathfrak{t}} \mathbf{X}_{1}^{t} & \mathbf{X}_{2} \mathbf{X}_{2}\end{array}\right ) \mathbb{T}_{x} =\left (\begin{array}{cc}\mathbb{T}_{x 1 1} & \mathbb{T}_{x 1 2} \\
\mathbb{T}_{x 2 1} & \mathbb{T}_{x 2 2}\end{array}\right )^{t} \left (\begin{array}{cc}\mathbf{X}_{1}^{t} \mathbf{X}_{1} & \mathbf{X}_{1}^{t} \mathbf{X}_{2} \\
\mathbf{X}_{2}^{t} \mathbf{X}_{1}^{t} & \mathbf{X}_{2} \mathbf{X}_{2}\end{array}\right ) \left (\begin{array}{cc}\mathbb{T}_{x 1 1} & \mathbb{T}_{x 1 2} \\
\mathbb{T}_{x 2 1} & \mathbb{T}_{x 2 2}\end{array}\right ) \nonumber  \\
 &  =\left (\begin{array}{cc}\mathbb{T}_{x 1 1}^{t} & \mathbb{T}_{x 2 1}^{t} \\
\mathbb{T}_{x 1 2}^{t} & \mathbb{T}_{x 2 2}^{t}\end{array}\right ) \left (\begin{array}{cc}\mathbf{X}_{1}^{t} \mathbf{X}_{1} & \mathbf{X}_{1}^{t} \mathbf{X}_{2} \\
\mathbf{X}_{2}^{t} \mathbf{X}_{1}^{t} & \mathbf{X}_{2} \mathbf{X}_{2}\end{array}\right ) \left (\begin{array}{cc}\mathbb{T}_{x 1 1} & \mathbb{T}_{x 1 2} \\
\mathbb{T}_{x 2 1} & \mathbb{T}_{x 2 2}\end{array}\right ) \nonumber  \\
 &  =\left (\begin{array}{cc}\mathbb{T}_{x 1 1}^{t} \mathbf{X}_{1}^{t} \mathbf{X}_{1} +\mathbb{T}_{x 2 1}^{t} \mathbf{X}_{2}^{t} \mathbf{X}_{1} & \mathbb{T}_{x 1 1}^{t} \mathbf{X}_{1}^{t} \mathbf{X}_{2} +\mathbb{T}_{x 2 1}^{t} \mathbf{X}_{2}^{t} \mathbf{X}_{2} \\
\mathbb{T}_{x 1 2}^{t} \mathbf{X}_{1}^{t} \mathbf{X}_{1} +\mathbb{T}_{x 2 2}^{t} \mathbf{X}_{2}^{t} \mathbf{X}_{1} & \mathbb{T}_{x 1 2}^{t} \mathbf{X}_{1}^{t} \mathbf{X}_{2} +\mathbb{T}_{x 2 2}^{t} \mathbf{X}_{2}^{t} \mathbf{X}_{2}\end{array}\right ) \left (\begin{array}{cc}\mathbb{T}_{x 1 1} & \mathbb{T}_{x 1 2} \\
\mathbb{T}_{x 2 1} & \mathbb{T}_{x 2 2}\end{array}\right ) \\
 &  =\left (\begin{array}{c}\mathbb{T}_{x 1 1}^{t} \mathbf{X}_{1}^{t} \mathbf{X}_{1} \mathbb{T}_{x 1 1} +\mathbb{T}_{x 2 1}^{t} \mathbf{X}_{2}^{t} \mathbf{X}_{1} \mathbb{T}_{x 1 1} +\mathbb{T}_{x 1 1}^{t} \mathbf{X}_{1}^{t} \mathbf{X}_{2} \mathbb{T}_{x 2 1} +\mathbb{T}_{x 2 1}^{t} \mathbf{X}_{2}^{t} \mathbf{X}_{2} \mathbb{T}_{x 2 1} \\
\mathbb{T}_{x 1 2}^{t} \mathbf{X}_{1}^{t} \mathbf{X}_{1} \mathbb{T}_{x 1 1} +\mathbb{T}_{x 2 2}^{t} \mathbf{X}_{2}^{t} \mathbf{X}_{1} \mathbb{T}_{x 1 1} +\mathbb{T}_{x 1 2}^{t} \mathbf{X}_{1}^{t} \mathbf{X}_{2} \mathbb{T}_{x 2 1} +\mathbb{T}_{x 2 2}^{t} \mathbf{X}_{2}^{t} \mathbf{X}_{2} \mathbb{T}_{x 2 1}\end{array}\right . \\
 & \text{\quad \quad \quad \quad \quad \quad \quad \quad \quad \quad }\left .\begin{array}{c}\mathbb{T}_{x 1 1}^{t} \mathbf{X}_{1}^{t} \mathbf{X}_{1} \mathbb{T}_{x 1 2} +\mathbb{T}_{x 2 1}^{t} \mathbf{X}_{2}^{t} \mathbf{X}_{1} \mathbb{T}_{x 1 2} +\mathbb{T}_{x 1 1}^{t} \mathbf{X}_{1}^{t} \mathbf{X}_{2} \mathbb{T}_{x 2 2} +\mathbb{T}_{x 2 1}^{t} \mathbf{X}_{2}^{t} \mathbf{X}_{2} \mathbb{T}_{x 2 2} \\
\mathbb{T}_{x 1 2}^{t} \mathbf{X}_{1}^{t} \mathbf{X}_{1} \mathbb{T}_{x 1 2} +\mathbb{T}_{x 2 2}^{t} \mathbf{X}_{2}^{t} \mathbf{X}_{1} \mathbb{T}_{x 1 2} +\mathbb{T}_{x 1 2}^{t} \mathbf{X}_{1}^{t} \mathbf{X}_{2} \mathbb{T}_{x 2 2} +\mathbb{T}_{x 2 2}^{t} \mathbf{X}_{2}^{t} \mathbf{X}_{2} \mathbb{T}_{x 2 2}\end{array}\right )\text{.}\end{align}
\end{proof}


\end{proposition}

\subsection{Adjoint}
There are two types of adjoint 


\begin{remark}
\begin{enumerate}
\item The $\mathcal{F}$\emph{-adjoint}
\begin{equation} \dag  :\mathcal{F} \rightarrow \mathcal{F}
\end{equation}and 

\item The \emph{Super Adjoint}
\begin{equation}\mathfrak{a} :M \left (\mathcal{F} ,m ,n\right ) \rightarrow M \left (\mathcal{F} ,n ,m\right )\text{,}
\end{equation}where
\begin{align}m &  =\text{number of rows} \\
n &  =\text{number of columns.}\end{align}defined by
\begin{equation}\mathfrak{a} :\overset{ \leftarrow \text{\quad \quad \quad \quad }n\text{\quad \quad \quad \quad } \rightarrow }{\left (\begin{array}{ccc}x_{1 1} & \cdots  & x_{1 m} \\
\vdots  & \, & \vdots  \\
x_{n 1} & \cdots  & x_{n m}\end{array}\right )}\, \begin{array}{c}\uparrow  \\
m \\
\downarrow \end{array} \longrightarrow \,\overset{ \leftarrow \text{\quad \quad \quad \quad }m\text{\quad \quad \quad \quad } \rightarrow }{\left (\begin{array}{ccc}x_{1 1}^{ \dag } & \cdots  & x_{1 n}^{ \dag } \\
\vdots  & \, & \vdots  \\
x_{m 1}^{ \dag } & \cdots  & x_{m n}^{ \dag }\end{array}\right )\, \begin{array}{c}\uparrow  \\
n \\
\downarrow \end{array}} ;\,x_{j k} \in \mathcal{F} ;1 \leq j \leq m ,1 \leq k \leq n\text{,}
\end{equation}i.e.
\begin{equation}\mathfrak{a} = \dag  \circ \# =\# \circ  \dag \text{.}
\end{equation}\end{enumerate}
\end{remark}

Thus if $\mathcal{X}_{j} \in $ $M \left (\mathcal{F} ,1 ,1\right ) ;$ $1 \leq j \leq 2 r$
\begin{equation}\mathfrak{a} = \dag 
\end{equation}and the super adjoint becomes
\begin{align}\mathbf{X} &  =\left (\begin{array}{ccc}\mathcal{X}_{1} & \cdots  & \mathcal{X}_{2 r}\end{array}\right ) \\
\mathbf{X}^{\mathfrak{a}} &  =\left (\begin{array}{c}\mathcal{X}_{1}^{ \dag } \\
\vdots  \\
\mathcal{X}_{2 r}^{ \dag }\end{array}\right )\text{,}\end{align}and
\begin{equation}\mathbf{X}^{\mathfrak{a}} \mathbf{X} =\left (\begin{array}{ccc}\mathcal{X}_{1}^{ \dag } \mathcal{X}_{1} & \cdots  & \mathcal{X}_{1}^{ \dag } \mathcal{X}_{2 r} \\
\vdots  & \vdots  & \vdots  \\
\mathcal{X}_{2 r}^{ \dag } \mathcal{X}_{1} & \cdots  & \mathcal{X}_{2 r}^{ \dag } \mathcal{X}_{2 r}\end{array}\right )\text{,}
\end{equation}which gives
\begin{equation}\left (\begin{array}{cc}\mathbf{X}_{1} & \mathbf{X}_{2}\end{array}\right )^{\mathfrak{a}} \left (\begin{array}{cc}\mathbf{X}_{1} & \mathbf{X}_{2}\end{array}\right ) =\left (\begin{array}{c}\mathbf{X}_{1}^{\mathfrak{a}} \\
\mathbf{X}_{2}^{\mathfrak{a}}\end{array}\right ) \left (\begin{array}{cc}\mathbf{X}_{1} & \mathbf{X}_{2}\end{array}\right ) =\left (\begin{array}{cc}\mathbf{X}_{1}^{\mathfrak{a}} \mathbf{X}_{1} & \mathbf{X}_{1}^{\mathfrak{a}} \mathbf{X}_{2} \\
\mathbf{X}_{2}^{\mathfrak{a}} \mathbf{X}_{1} & \mathbf{X}_{2}^{\mathfrak{a}} \mathbf{X}_{2}\end{array}\right )\text{.}
\end{equation}

Consider the linear transformation
\begin{align}T \left (\begin{array}{cc}\mathbf{X}_{1} & \mathbf{X}_{2}\end{array}\right ) &  =\left (\begin{array}{cc}T \mathbf{X}_{1} & T \mathbf{X}_{2}\end{array}\right ) =\left (\begin{array}{cc}\mathbf{X}_{1} & \mathbf{X}_{2}\end{array}\right ) \mathcal{R}_{x} \left (T\right ) =\left (\begin{array}{cc}\mathbf{X}_{1} & \mathbf{X}_{2}\end{array}\right ) \mathbb{T}_{x} \\
 &  =\left (\begin{array}{cc}\mathbf{X}_{1} \mathbb{T}_{x 1 1} +\mathbf{X}_{2} \mathbb{T}_{x 2 1} & \mathbf{X}_{1} \mathbb{T}_{x 1 2} +\mathbf{X}_{2} \mathbb{T}_{x 2 2}\end{array}\right ) =\left (\begin{array}{cc}\mathbf{Y}_{1} & \mathbf{Y}_{2}\end{array}\right )\text{,}\end{align}where
\begin{equation}\mathbb{T}_{x} =\left (\begin{array}{cc}\mathbb{T}_{x 1 1} & \mathbb{T}_{x 1 2} \\
\mathbb{T}_{x 2 1} & \mathbb{T}_{x 2 2}\end{array}\right )\text{.}
\end{equation}Then
\begin{align}\left (T \left (\mathbf{X}_{1}\right )\right )^{\mathfrak{a}} &  =\mathbf{X}_{1}^{\mathfrak{a}} \mathbb{T}_{x 1 1}^{ \dag } +\mathbf{X}_{2}^{\mathfrak{a}} \mathbb{T}_{x 2 1}^{ \dag } \\
\left (T \left (\mathbf{X}_{2}\right )\right )^{\mathfrak{a}} &  =\mathbf{X}_{1}^{\mathfrak{a}} \mathbb{T}_{x 1 2}^{ \dag } +\mathbf{X}_{2}^{\mathfrak{a}} \mathbb{T}_{x 2 2}^{ \dag }\text{,}\end{align}as
\begin{align}\left (\sum _{j}\mathcal{X}_{1 j} T_{1 1 j k}\right )^{ \dag } &  =\sum _{j}\mathcal{X}_{1 j}^{ \dag } T_{1 1 j k} =\sum _{j}T_{1 1 j k} \mathcal{X}_{1 j}^{ \dag } =\sum _{j}\left (T_{1 1}^{ \dag }\right )_{k j} \mathcal{X}_{1 j}^{ \dag } \\
\left (\sum _{j}\mathcal{X}_{2 j} T_{2 1 j k}\right )^{ \dag } &  =\sum _{j}\mathcal{X}_{2 j}^{ \dag } T_{2 1 j k} =\sum _{j}T_{2 1 j k} \mathcal{X}_{2 j}^{ \dag } =\sum _{j}\left (T_{2 1}^{ \dag }\right )_{k j} \mathcal{X}_{2 j}^{ \dag }\end{align}
\begin{align}\left (\sum _{j}\mathcal{X}_{1 j} T_{1 2 j k}\right )^{ \dag } &  =\sum _{j}\mathcal{X}_{1 j}^{ \dag } T_{1 2 j k} =\sum _{j}T_{1 2 j k} \mathcal{X}_{1 j}^{ \dag } =\sum _{j}\left (T_{1 2}^{ \dag }\right )_{k j} \mathcal{X}_{1 j}^{ \dag } \\
\left (\sum _{j k}\mathcal{X}_{2 j} T_{2 2 j k}\right )^{ \dag } &  =\sum _{j}\mathcal{X}_{2 j}^{ \dag } T_{2 2 j k} =\sum _{j}T_{2 2 j k} \mathcal{X}_{2 j}^{ \dag } =\sum _{j}\left (T_{2 2}^{ \dag }\right )_{k j} \mathcal{X}_{2 j}^{ \dag }\text{.}\end{align}This gives
\begin{align} & \left (\begin{array}{c}\mathbf{Y}_{1}^{\mathfrak{a}} \\
\mathbf{Y}_{2}^{\mathfrak{a}}\end{array}\right ) \left (\begin{array}{cc}\mathbf{Y}_{1} & \mathbf{Y}_{2}\end{array}\right ) =\left (\left (\begin{array}{cc}\mathbf{X}_{1} & \mathbf{X}_{2}\end{array}\right ) \mathbb{T}_{x}\right )^{\mathfrak{a}} \left (\begin{array}{cc}\mathbf{X}_{1} & \mathbf{X}_{2}\end{array}\right ) \mathbb{T}_{x} =\mathbb{T}_{x}^{ \dag } \left (\begin{array}{cc}\mathbf{X}_{1} & \mathbf{X}_{2}\end{array}\right )^{\mathfrak{a}} \left (\begin{array}{cc}\mathbf{X}_{1} & \mathbf{X}_{2}\end{array}\right ) \mathbb{T}_{x} \nonumber  \\
 &  =\mathbb{T}_{x}^{ \dag } \left (\begin{array}{cc}\mathbf{X}_{1}^{\mathfrak{a}} \mathbf{X}_{1} & \mathbf{X}_{1}^{\mathfrak{a}} \mathbf{X}_{2} \\
\mathbf{X}_{2}^{\mathfrak{a}} \mathbf{X}_{1}^{ \dag } & \mathbf{X}_{2}^{\mathfrak{a}} \mathbf{X}_{2}\end{array}\right ) \mathbb{T}_{x} =\left (\begin{array}{cc}\mathbb{T}_{x 1 1} & \mathbb{T}_{x 1 2} \\
\mathbb{T}_{x 2 1} & \mathbb{T}_{x 2 2}\end{array}\right )^{ \dag } \left (\begin{array}{cc}\mathbf{X}_{1}^{\mathfrak{a}} \mathbf{X}_{1} & \mathbf{X}_{1}^{\mathfrak{a}} \mathbf{X}_{2} \\
\mathbf{X}_{2}^{\mathfrak{a}} \mathbf{X}_{1}^{ \dag } & \mathbf{X}_{2}^{\mathfrak{a}} \mathbf{X}_{2}\end{array}\right ) \left (\begin{array}{cc}\mathbb{T}_{x 1 1} & \mathbb{T}_{x 1 2} \\
\mathbb{T}_{x 2 1} & \mathbb{T}_{x 2 2}\end{array}\right ) \nonumber  \\
 &  =\left (\begin{array}{cc}\mathbb{T}_{x 1 1}^{ \dag } & \mathbb{T}_{x 2 1}^{ \dag } \\
\mathbb{T}_{x 1 2}^{ \dag } & \mathbb{T}_{x 2 2}^{ \dag }\end{array}\right ) \left (\begin{array}{cc}\mathbf{X}_{1}^{\mathfrak{a}} \mathbf{X}_{1} & \mathbf{X}_{1}^{\mathfrak{a}} \mathbf{X}_{2} \\
\mathbf{X}_{2}^{\mathfrak{a}} \mathbf{X}_{1} & \mathbf{X}_{2}^{\mathfrak{a}} \mathbf{X}_{2}\end{array}\right ) \left (\begin{array}{cc}\mathbb{T}_{x 1 1} & \mathbb{T}_{x 1 2} \\
\mathbb{T}_{x 2 1} & \mathbb{T}_{x 2 2}\end{array}\right )\text{.}\end{align}

\subsubsection{Hilbert Schmidt Operators $\mathcal{F}_{H S} =\mathcal{B}_{2} \left (\mathcal{H}_{\mathcal{F}}\right )$}
Using the trace one can introduce 


\begin{enumerate}
\item An inner product in $\mathcal{F}$ defined by
\begin{equation} \langle \left .X\right \vert Y \rangle _{H S} =\ensuremath{\operatorname*{Tr}}_{\mathcal{H}_{\mathcal{F}}} \left \{X^{ \dag } Y\right \} \equiv \ensuremath{\operatorname*{Tr}} \left \{X^{ \dag } Y\right \} =2^{r} \ensuremath{\operatorname*{tr}}\left \{X^{ \dag } Y\right \}\text{,}
\end{equation}where $\ensuremath{\operatorname*{tr}}$ is the normalized trace, 

\item A norm
\begin{equation}\left \Vert X\right \Vert _{H S} =\left ( \langle \left .X\right \vert X \rangle _{H S}\right )^{\frac{1}{2}}\text{,}
\end{equation}

\item An inner product in $\mathcal{F}_{n}$ by restricting $X ,Y \in \mathcal{F}_{n}$ thus making $\mathcal{F}_{n H S}\text{,}$ (the closure of $\mathcal{F}_{n}$ wrt $\left \Vert \text{}\right \Vert _{H S})\text{,}$ a \emph{Hilbert subspace} of $\mathcal{F}_{H S}\text{.}$ 

\item $\mathcal{F}_{H S}$ can be expressed as a direct sum of disjoint subspaces
\begin{equation}\mathcal{F}_{H S} =\tbigoplus ???\mathcal{F}_{n H S}\text{,}
\end{equation}where $\left \{\left .\mathcal{F}_{n H S}\right \vert 0 \leq n \leq 2 r\right \}$ and $\mathcal{F}_{H S}$ are complete in the HS norm $\left \Vert X\right \Vert _{H S} \equiv \left \Vert \text{}\right \Vert _{\mathcal{B}_{2} \left (\mathcal{H}_{\mathcal{F}}\right )}\text{.}$ \end{enumerate}


Some properties of $\left \{\left .\mathcal{F}_{n H S}\right \vert 0 \leq n \leq 2 r\right \}$ and $\mathcal{F}_{H S}$ are inherited from $\left \{\left .\mathcal{F}_{n}\right \vert 0 \leq n \leq 2 r\right \}$ and $\mathcal{F}$ 


\begin{enumerate}
\item The Banach space norm of $\mathcal{F}$, which is the $c^{ \ast }$-algebra norm, is defined as
\begin{equation}\left \Vert X\right \Vert _{\mathcal{F}} \equiv \left \Vert X\right \Vert _{\mathcal{B} \left (\mathcal{H}_{\mathcal{F}}\right )} =\sup_{D \in \mathcal{S}_{n o r m a l} \left \{\mathcal{F}\right \}}\left \{\ensuremath{\operatorname*{Tr}}_{\mathcal{H}_{\mathcal{F}}} \left \{X^{ \dag } X D\right \}\right \} =\sup_{D \in \ensuremath{\operatorname*{ext}}\left \{\mathcal{S}_{n o r m a l} \left \{\mathcal{F}\right \}\right \}}\left \{\ensuremath{\operatorname*{Tr}}_{\mathcal{H}_{\mathcal{F}}} \left \{X^{ \dag } X D\right \}\right \}\text{,}
\end{equation}where $\mathcal{S}_{n o r m a l} \left \{\mathcal{F}\right \}$ is the set of normal states of $\mathcal{F}$ and $\ensuremath{\operatorname*{ext}}\left \{\mathcal{S}_{n o r m a l} \left \{\mathcal{F}\right \}\right \}$ is the set of extreme points of the convex set $\mathcal{S}_{n o r m a l} \left \{\mathcal{F}\right \}\text{.}$ 


\begin{remark}
\begin{equation}\sup_{D \in \ensuremath{\operatorname*{ext}}\left \{\mathcal{S}_{n o r m a l} \left \{\mathcal{F}\right \}\right \}}\left \{\ensuremath{\operatorname*{Tr}}_{\mathcal{H}_{\mathcal{F}}} \left \{X^{ \dag } X D\right \}\right \} =\sup_{\Psi  \in \mathcal{H}_{\mathcal{F}} ;\left \Vert \Psi \right \Vert _{\mathcal{H}_{\mathcal{F}}} =1} \langle \left .X \Psi \right \vert X \Psi  \rangle \text{.}
\end{equation}
\end{remark}

\item As a Banach space $\mathcal{F}$ can be expressed as a direct sum of disjoint subspaces
\begin{equation}\mathcal{F} =\tbigoplus ???\mathcal{F}_{n}\text{,}
\end{equation}where $\left \{\left .\mathcal{F}_{n}\right \vert 0 \leq n \leq 2 r\right \}$ and $\mathcal{F}$ are complete in the Banach space norm $\left \Vert \text{}\right \Vert _{\mathcal{F}} \equiv \left \Vert \text{}\right \Vert _{\mathcal{B} \left (\mathcal{H}_{\mathcal{F}}\right )}\text{.}$ 

\item The adjoint operation in $\mathcal{F}$ leaves the subspaces $\mathcal{F}_{n H S}$ and $\mathcal{F}_{n}$ invariant i.e.
\begin{gather} \dag  :\mathcal{F}_{n H S} \rightarrow \mathcal{F}_{n H S}\text{,} \\
 \dag  :\mathcal{F}_{n} \rightarrow \mathcal{F}_{n} ,\, \\
0 \leq n \leq 2 r\text{.}\end{gather}\end{enumerate}



\begin{lemma}
$\mathcal{F}_{H S} = \tbigwedge \left (\mathcal{F}_{1 H S}\right ) = \tbigwedge \left (\mathcal{F}_{1 1 H S} \oplus \mathcal{F}_{1 1 H S}^{d}\right ) = \tbigwedge \left (\mathcal{F}_{1 1 H S}\right ) \otimes  \tbigwedge  \left (\mathcal{F}_{1 1 H S}\right )^{d} \cong \mathcal{B} \left (\mathcal{F}_{1 H S}\right )\text{.}$ 
\end{lemma}


\begin{remark}
The unitary group $U \left (\mathcal{F}_{1 H S}\right )$ is defined as
\begin{equation}U \left (\mathcal{F}_{1 H S}\right ) =\left \{\left .T \in \mathcal{B} \left (\mathcal{F}_{1 H S}\right )\right \vert  \langle \left .T \left (X\right )\right \vert T \left (Y\right ) \rangle _{H S} = \langle \left .X\right \vert Y \rangle _{H S} , \forall X ,Y \in \mathcal{F}_{1 H S}\right \}
\end{equation}

The unitary group $U \left (\mathcal{F}_{H S}\right )$ is defined as
\begin{equation}U \left (\mathcal{F}_{H S}\right ) =\left \{\left .T \in \mathcal{B} \left (\mathcal{F}_{H S}\right )\right \vert  \langle \left .T \left (X\right )\right \vert T \left (Y\right ) \rangle _{H S} = \langle \left .X\right \vert Y \rangle _{H S} , \forall X ,Y \in \mathcal{F}_{H S}\right \}\text{.}
\end{equation}
\end{remark}


\begin{remark}
The bases $\left \{\left .a_{j}^{ \dag } ,a_{j}\right \vert 1 \leq j \leq r\text{}\right \}$ and $\left \{\left .x_{j}\right \vert 1 \leq j \leq 2 r\right \}$ are orthogonal but not normalized wrt these inner products
\begin{equation} \langle \left .a_{j}\right \vert a_{k} \rangle _{H S} = \langle \left .a_{j}^{ \dag }\right \vert a_{k}^{ \dag } \rangle _{H S} = \langle \left .x_{j}\right \vert x_{k} \rangle _{H S} =2^{r -1} \delta _{j k}\text{.}
\end{equation}
\end{remark}

Utilizing the isomorphism
\begin{equation}\mathcal{B} \left (\mathcal{F}_{H S}\right ) \cong \mathcal{F}_{H S} \otimes \mathcal{F}_{H S}^{d}
\end{equation}given by the hermitian inner product, $ \langle \left .\text{}\right \vert  \rangle _{H S}$ in $\mathcal{F}_{H S}\text{,}$ one can define the Ket-Bra operators $\vert X \rangle  \langle Y\vert  \in \mathcal{B} \left (\mathcal{F}_{H S}\right )$
\begin{equation}\vert X \rangle  \langle Y\vert  :\mathcal{F}_{H S} \rightarrow \mathcal{F}_{H S}\text{,}
\end{equation}by
\begin{equation}\vert X \rangle  \langle Y\vert A = \langle \left .Y\right \vert A \rangle _{H S} X\text{.}
\end{equation}


\begin{remark}
The space $\mathcal{F}_{H S}$ is the closure, $\mathcal{B}_{2} \left (\mathcal{H}_{\mathcal{F}}\right )\text{,}$ of finite rank operators, $\mathcal{B}_{0 0} \left (\mathcal{H}_{\mathcal{F}}\right )\text{,}$ wrt the $\left \Vert \text{}\right \Vert _{H S}$ norm and is a \emph{Hilbert Algebra}. 
\end{remark}


\begin{remark}
In \emph{infinite dimensions} $\left \{\left .a_{j}^{ \dag } ,a_{j}\right \vert 1 \leq j \leq r\text{}\right \}$ and $\left \{\left .x_{j}\right \vert 1 \leq j \leq 2 r\text{}\right \}\text{,}$ (self adjoint combnations of the creation and annihilation operators explicitly defined in the proceding), are \emph{NOT Trace
Class operators thus NOT Hilbert Schmidt operators i.e. }
\begin{align}a_{j}^{ \dag } ,a_{j} &  \notin \mathcal{B}_{2} \left (\mathcal{H}_{\mathcal{F}}\right ) \supseteq \mathcal{B}_{1} \left (\mathcal{H}_{\mathcal{F}}\right ) ;\text{\quad }1 \leq j \leq r \\
x_{j} &  \notin \mathcal{B}_{2} \left (\mathcal{H}_{\mathcal{F}}\right ) \supseteq \mathcal{B}_{1} \left (\mathcal{H}_{\mathcal{F}}\right ) ;\text{\quad }1 \leq j \leq 2 r\text{.}\end{align}\emph{\ } 
\end{remark}


\begin{remark}
The following inclusions hold and are all proper when H is infinite-dimensional: \{finite rank\} {\small\fbox{2282}}  \{trace class\} {\small\fbox{2282}}  \{Hilbert--Schmidt\}
{\small\fbox{2282}}  \{compact\} i.e.
\begin{equation}\mathcal{B}_{0 0} \left (\mathcal{H}_{\mathcal{F}}\right ) \subset \mathcal{B}_{1} \left (\mathcal{H}_{\mathcal{F}}\right ) \subset \mathcal{B}_{2} \left (\mathcal{H}_{\mathcal{F}}\right ) \subset \mathcal{B}_{0} \left (\mathcal{H}_{\mathcal{F}}\right )\text{.}
\end{equation}
\end{remark}


\begin{notation}
One can denote vectors belonging to an inner product space, where the inner products are $\left (\left .\text{}\right \vert \right )_{\ensuremath{\operatorname*{Tr}}}$ and $ \langle \left .\text{}\right \vert  \rangle _{H S}\text{,}$ \emph{produced by} \emph{elements of an algebra} by either pointed or rounded brackets as
\begin{equation}\left \vert \mathbf{X}\right ) \equiv \vert \mathbf{X} \rangle  =\left (\begin{array}{c}\mathbf{x}_{1} \\
\mathbf{x}_{2}\end{array}\right )\text{.}
\end{equation}


\begin{enumerate}
\item The $\mathcal{F}$-\emph{Transpose} of vectors belonging $\mathcal{F}_{1}^{2 r}$ are denoted by
\begin{equation}\left \vert \mathbf{X}\right )^{t} \equiv \vert \mathbf{X} \rangle ^{t} =\left \vert \mathbf{X}^{t}\right ) \equiv \vert \mathbf{X}^{t} \rangle  =\left (\begin{array}{c}\vert \mathbf{X}_{1} \rangle  \\
\vert \mathbf{X}_{2} \rangle \end{array}\right )^{t} =\left (\begin{array}{c}\left \vert \mathbf{X}_{1}\right ) \\
\left \vert \mathbf{X}_{2}\right )\end{array}\right )^{t} =\left (\begin{array}{cc}\left \vert \mathbf{X}_{1}^{t}\right ) & \left \vert \mathbf{X}_{2}^{t}\right )\end{array}\right ) =\left (\begin{array}{c}\mathbf{X}_{1} \\
\mathbf{X}_{2}\end{array}\right )^{t} =\left (\begin{array}{cc}\mathbf{X}_{1}^{t} & \mathbf{X}_{2}^{t}\end{array}\right )\text{.}
\end{equation}

\item The \emph{Simple Transpose} of vectors belonging
$\mathcal{F}_{1}^{2 r}$ are denoted by
\begin{equation}\left \vert \mathbf{X}\right )^{\#} =\left (\begin{array}{c}\left \vert \mathbf{X}_{1}\right ) \\
\left \vert \mathbf{X}_{2}\right )\end{array}\right )^{\#} =\left (\begin{array}{cc}\left \vert \mathbf{X}_{1}\right )^{\#} & \left \vert \mathbf{X}_{2}\right )^{\#}\end{array}\right ) =\left (\begin{array}{cc}\left (\mathbf{X}_{1}\right \vert  & \left (\mathbf{X}_{2}\right \vert \end{array}\right )\text{.}
\end{equation}

\item The \emph{Super Transpose} of of vectors belonging
$\mathcal{F}_{1}^{2 r}$ are denoted by
\begin{equation}\left \vert \mathbf{X}\right )^{\mathfrak{t}} =\left \vert \mathbf{X}\right )^{t \#} =\left (\begin{array}{c}\left \vert \mathbf{X}_{1}\right ) \\
\left \vert \mathbf{X}_{2}\right )\end{array}\right )^{t \#} =\left (\begin{array}{cc}\left \vert \mathbf{X}_{1}\right )^{t \#} & \left \vert \mathbf{X}_{2}\right )^{t \#}\end{array}\right ) =\left (\begin{array}{cc}\left (\mathbf{X}_{1}^{t}\right \vert  & \left (\mathbf{X}_{2}^{t}\right \vert \end{array}\right )\text{.}
\end{equation}

\item The $\mathcal{F}$-\emph{Adjoint} of of vectors belonging $\mathcal{F}_{1}^{2 r}$ are denoted by
\begin{equation}\vert \mathbf{X} \rangle ^{ \dag } =\vert \mathbf{X}^{ \dag } \rangle  =\left (\begin{array}{c}\vert \mathbf{X}_{1} \rangle  \\
\vert \mathbf{X}_{2} \rangle \end{array}\right )^{ \dag } =\left (\begin{array}{c}\vert \mathbf{X}_{1}^{ \dag } \rangle  \\
\vert \mathbf{X}_{2}^{ \dag } \rangle \end{array}\right )\text{.}
\end{equation}

\item The $\mathbf{S} \mathbf{u} \mathbf{p} \mathbf{e} \mathbf{r}$ \emph{adjoint} of such vectors if the elements belong to a star algebra denoted by
\begin{equation}\vert \mathbf{X} \rangle ^{\mathfrak{a}} =\vert \mathbf{X} \rangle ^{ \dag \#} =\left (\begin{array}{c}\vert \mathbf{X}_{1} \rangle  \\
\vert \mathbf{X}_{2} \rangle \end{array}\right )^{ \dag \#} =\left (\begin{array}{cc}\vert \mathbf{X}_{1} \rangle ^{ \dag \#} & \vert \mathbf{X}_{2} \rangle ^{ \dag \#}\end{array}\right ) =\left (\begin{array}{cc} \langle \mathbf{X}_{1}\vert  &  \langle \mathbf{X}_{2}\vert \end{array}\right ) = \langle \mathbf{X}\vert \text{.}
\end{equation}\end{enumerate}



\begin{remark}
The adjoint operation "$ \dag ""$ is the star operation in the $c^{ \ast }$-algebra $\mathcal{F}$: 


\begin{enumerate}
\item if the algebra is the complex numbers $\mathbb{C}$ then it is complex conjugation, 

\item if the algebra is formed by the bounded
operators $\mathcal{B} \left (\mathcal{H}\right )$ acting in an Hilbert space $\mathcal{H}$ it is the the adjoint operation induced by the inner product of the Hilbert space. \end{enumerate}
\end{remark}


\end{notation}

\subsubsection{Inner products expressed in terms of vector arrays:}
\begin{notation}
\begin{enumerate}
\item The \emph{simple} \emph{transpose} of bra vector arrays belonging to $\mathcal{F}_{1}^{2 r}$ $ =$ $M \left (2 r ,1 ,\mathcal{F}_{1}\right )$ are denoted by
\begin{equation}\left (\mathbf{X}\right \vert  =\left \vert \mathbf{X}\right )^{\#} =\left (\begin{array}{c}\left \vert \mathbf{X}_{1}\right ) \\
\left \vert \mathbf{X}_{2}\right )\end{array}\right )^{\#} =\left (\begin{array}{cc}\left \vert \mathbf{X}_{1}\right )^{\#} & \left \vert \mathbf{X}_{2}\right )^{\#}\end{array}\right ) =\left (\begin{array}{cc}\left (\mathbf{X}_{1}\right \vert  & \left (\mathbf{X}_{2}\right \vert \end{array}\right )
\end{equation}

giving
\begin{equation}\left (\mathbf{Y}\right \vert \left \vert \mathbf{X}\right ) =\ensuremath{\operatorname*{Tr}} \left \{\mathbf{Y} \mathbf{X}\right \} =\ensuremath{\operatorname*{Tr}} \left \{\left (\begin{array}{cc}\boldsymbol{Y}_{1} & \mathbf{Y}_{2}\end{array}\right ) \left (\begin{array}{c}\mathbf{X}_{1} \\
\mathbf{X}_{2}\end{array}\right )\right \} =\ensuremath{\operatorname*{Tr}} \left \{\left (\begin{array}{c}\mathbf{Y}_{1} \\
\mathbf{Y}_{2}\end{array}\right )^{\#} \left (\begin{array}{c}\mathbf{X}_{1} \\
\mathbf{X}_{2}\end{array}\right )\right \}\text{.}
\end{equation}

\item The \emph{super} \emph{transpose}
of bra vector arrays belonging to $\mathcal{F}_{1}^{2 r}$ $ =$ $M \left (2 r ,1 ,\mathcal{F}_{1}\right )$ are denoted by
\begin{equation}\left (\mathbf{X}\right \vert  =\left \vert \mathbf{X}\right )^{\mathfrak{t}} =\left (\begin{array}{c}\left \vert \mathbf{X}_{1}\right ) \\
\left \vert \mathbf{X}_{2}\right )\end{array}\right )^{\mathfrak{t}} =\left (\begin{array}{cc}\left \vert \mathbf{X}_{1}\right )^{\mathfrak{t}} & \left \vert \mathbf{X}_{2}\right )^{\mathfrak{t}}\end{array}\right ) =\left (\begin{array}{cc}\left (\mathbf{X}_{1}^{t}\right \vert  & \left (\mathbf{X}_{2}^{t}\right \vert \end{array}\right )
\end{equation}giving
\begin{equation}\left (\mathbf{Y}^{t}\right \vert \left \vert \mathbf{X}\right ) =\ensuremath{\operatorname*{Tr}} \left \{\mathbf{Y}^{t} \mathbf{X}\right \} =\ensuremath{\operatorname*{Tr}} \left \{\left (\begin{array}{cc}\mathbf{Y}_{1}^{t} & \mathbf{Y}_{2}^{t}\end{array}\right ) \left (\begin{array}{c}\mathbf{X}_{1} \\
\mathbf{X}_{2}\end{array}\right )\right \} =\ensuremath{\operatorname*{Tr}} \left \{\left (\begin{array}{c}\mathbf{Y}_{1} \\
\mathbf{Y}_{2}\end{array}\right )^{\# t} \left (\begin{array}{c}\mathbf{X}_{1} \\
\mathbf{X}_{2}\end{array}\right )\right \}\text{.}
\end{equation}

\item The \emph{super} \emph{adjoint}
of bra vector arrays belonging to $\mathcal{F}_{1}^{2 r}$ are denoted by
\begin{gather} \langle \mathbf{X}\vert  =\vert \mathbf{X} \rangle ^{\mathfrak{a}} =\vert \mathbf{X} \rangle ^{\# \dag } =\left (\mathbf{X}^{ \dag }\right \vert  =\left (\begin{array}{c}\vert \mathbf{X}_{1} \rangle  \\
\vert \mathbf{X}_{2} \rangle \end{array}\right )^{\# \dag } \equiv \left (\begin{array}{c}\vert \mathbf{X}_{1} \rangle  \\
\vert \mathbf{X}_{2} \rangle \end{array}\right )^{\mathfrak{a}} =\left (\begin{array}{cc}\vert \mathbf{X}_{1} \rangle ^{\mathfrak{a}} & \vert \mathbf{X}_{2} \rangle ^{\mathfrak{a}}\end{array}\right ) \\
 =\left (\begin{array}{cc} \langle \mathbf{X}_{1}\vert  &  \langle \mathbf{X}_{2}\vert \end{array}\right ) =\left (\begin{array}{cc}\left (\mathbf{X}_{1}^{ \dag }\right \vert  & \left (\mathbf{X}_{2}^{ \dag }\right \vert \end{array}\right ) \\
 \langle \mathbf{Y}\vert \vert \mathbf{X} \rangle  =\ensuremath{\operatorname*{Tr}} \left \{\mathbf{Y}^{ \dag } \mathbf{X}\right \} =\ensuremath{\operatorname*{Tr}} \left \{\left (\begin{array}{cc}\mathbf{Y}_{1}^{ \dag } & \mathbf{Y}_{2}^{ \dag }\end{array}\right ) \left (\begin{array}{c}\mathbf{X}_{1} \\
\mathbf{X}_{2}\end{array}\right )\right \} =\ensuremath{\operatorname*{Tr}} \left \{\left (\begin{array}{c}\mathbf{Y}_{1} \\
\mathbf{Y}_{2}\end{array}\right )^{\# \dag } \left (\begin{array}{c}\mathbf{X}_{1} \\
\mathbf{X}_{2}\end{array}\right )\right \} =\left (\mathbf{Y}^{ \dag }\right \vert \left \vert \mathbf{X}\right )\text{.}\end{gather}\end{enumerate}
\end{notation}


\begin{remark}
If $\mathcal{F}_{1} =\mathbb{C}^{2 r}$ then $^{ \dag }$ is complex conjugation and
\begin{equation} \langle \mathbf{Y}\vert \vert \mathbf{X} \rangle  =\left (\begin{array}{c}\mathbf{Y}_{1} \\
\mathbf{Y}_{2}\end{array}\right )^{\mathfrak{a}} \left (\begin{array}{c}\mathbf{X}_{1} \\
\mathbf{X}_{2}\end{array}\right ) =\left (\mathbf{Y}^{ \dag }\right \vert \left \vert \mathbf{X}\right ) =\left (\overline{\mathbf{Y}}\right \vert \left \vert \mathbf{X}\right ) =\sum _{1 \leq j \leq 2 r}\overline{\mathfrak{Y}_{j}} \mathfrak{X}_{j}
\end{equation}
\end{remark}

\paragraph{Expansions of Operators}
Operators acting in the Hilbert space $\mathcal{F}_{1 H S}$ can be expanded in terms of a orthonormal basis $\left \{\left .\mathfrak{X}_{1 j} ,\mathfrak{X}_{2 j}\right \vert 1 \leq j \leq r\right \}$ as finite rank operators or limits of finite rank operators as
\begin{equation}T =\left (\begin{array}{cc}\vert \mathbf{X}_{1} \rangle  & \vert \mathbf{X}_{2} \rangle \end{array}\right ) \mathbb{T}_{\mathfrak{X}} \left (\begin{array}{c} \langle \mathbf{X}_{1}\vert ^{\#} \\
 \langle \mathbf{X}_{2}\vert ^{\#}\end{array}\right )\text{,}
\end{equation}where
\begin{equation}T \left (\begin{array}{cc}\mathbf{X}_{1} & \mathbf{X}_{2}\end{array}\right ) =\left (\begin{array}{cc}\mathbf{X}_{1} & \mathbf{X}_{2}\end{array}\right ) \mathbb{T}_{\mathfrak{X}} \equiv \left (\begin{array}{cc}\mathbf{X}_{1} & \mathbf{X}_{2}\end{array}\right ) \mathcal{R}_{\mathfrak{X}} \left (T\right )
\end{equation}and
\begin{equation} \langle \left .\mathfrak{X}_{\alpha  j}\right \vert \mathfrak{X}_{\beta  k} \rangle _{H S} =\ensuremath{\operatorname*{Tr}} \left \{\mathfrak{X}_{\alpha  j}^{ \dag } \mathfrak{X}_{\beta  k}\right \} =\delta _{\alpha  \beta } \delta _{j k} ;\,1 \leq \alpha  ,\beta  \leq 2 ,1 \leq j ,k \leq r\text{.}
\end{equation}This can be further expanded as
\begin{equation*}
\end{equation*}


\begin{lemma}
$\mathcal{R}_{\mathfrak{X}}$ is a star representation of $\mathcal{B} \left (\mathcal{F}_{1 H S}\right )$ i.e.
\begin{equation}\mathcal{R}_{\mathfrak{X}} \left (T\right )^{ \dag } =\mathcal{R}_{\mathfrak{X}} \left (T^{ \dag }\right )\text{,}
\end{equation}where $"" \dag ""$ is wrt the HS\ inner product in $\mathcal{F}_{1 H S}\text{.}$ 
\end{lemma}


\begin{proof}
One can expand $T$ and $T^{ \dag }$
\begin{equation*}
\end{equation*}showing that
\begin{equation*}
\end{equation*}Thus $\mathcal{R}_{\mathfrak{X}}$ is a star representation. 
\end{proof}


\begin{remark}
$T$ in general is not a star transformation wrt $"" \dag ""$in $\mathcal{F}_{1 H S}$ i.e.
\begin{equation}T \left (X^{ \dag }\right ) \neq T \left (X\right )^{ \dag } ,\text{\quad }X \in \mathcal{F}_{1 H S}\text{.}
\end{equation}
\end{remark}


\begin{lemma}
In the space $\mathcal{B}_{a n t i} \left (\mathcal{F}_{H S}\right )$ of antilinear operators acting in the Hilbert Algebra $\mathcal{F}_{H S}$ the adjoint operation $ \ddag  \in \mathcal{B}_{a n t i} \left (\mathcal{F}_{H S}\right )$ generalizes the adjoint operation $ \dag  \in \mathcal{B}_{a n t i} \left (\mathcal{F}_{H S}\right )\text{,}$where $ \ddag $ is induced by the inner product $ \langle \left .\text{}\right \vert  \rangle _{H S}$ and $ \dag $ is induced by the inner product $ \langle \left .\text{}\right \vert  \rangle _{\mathcal{H}_{\mathcal{F}}}$ 
\end{lemma}


\begin{proof}
The adjoint mapping defined by $ \langle \left .\text{}\right \vert  \rangle _{\mathcal{H}_{\mathcal{F}}}$ is defined by
\begin{equation} \langle \left .X^{ \dag } \left (\Phi \right )\right \vert \Psi  \rangle _{\mathcal{H}_{\mathcal{F}}} = \langle \left .\Phi \right \vert X \left (\Psi \right ) \rangle _{\mathcal{H}_{\mathcal{F}}}\, \forall \Psi  \in \mathcal{H}_{\mathcal{F}} ,\Phi  \in \mathcal{H}_{\mathcal{F}}\text{.}
\end{equation}The adjoint mapping $ \ddag $ induced by $ \langle \left .\text{}\right \vert  \rangle _{H S}$ is defined by
\begin{equation} \langle \left .T^{ \ddag } \left (X\right )\right \vert Y \rangle _{H S} = \langle \left .X\right \vert T \left (Y\right ) \rangle _{H S}\, \forall Y \in \mathcal{F}_{H S} ,X \in \mathcal{F}_{H S}\text{.}
\end{equation}In particular consider the left multiplication map
\begin{equation}L \left (Z\right ) \left (Y\right ) =Z Y\text{\thinspace \quad } \forall Y \in \mathcal{F}_{H S}
\end{equation}and its adjoint defined by
\begin{equation} \langle \left .L \left (Z\right )^{ \ddag } X\right \vert Y \rangle _{H S} = \langle \left .X\right \vert L \left (Z\right ) \left (Y\right ) \rangle _{H S}\, \forall Y \in \mathcal{F}_{H S} ,Z ,X \in \mathcal{F}_{H S}\text{.}
\end{equation}Expanding the inner product in terms of the trace one obtains
\begin{equation} \langle \left .L \left (Z\right )^{ \ddag } \left (X\right )\right \vert Y \rangle _{H S} =\ensuremath{\operatorname*{Tr}} \left \{\left (L \left (Z\right )^{ \ddag } \left (X\right )\right )^{ \dag } Y\right \}\text{,}
\end{equation}which by the definition of the adjoint gives
\begin{equation} \langle \left .X\right \vert L \left (Z\right ) \left (Y\right ) \rangle _{H S} =\ensuremath{\operatorname*{Tr}} \left \{X^{ \dag } L \left (Z\right ) \left (Y\right )\right \}\text{,}
\end{equation}hence
\begin{gather}\ensuremath{\operatorname*{Tr}} \left \{\left (L \left (Z\right )^{ \ddag } \left (X\right )\right )^{ \dag } Y\right \} =\ensuremath{\operatorname*{Tr}} \left \{X^{ \dag } L \left (Z\right ) \left (Y\right )\right \} =\ensuremath{\operatorname*{Tr}} \left \{X^{ \dag } Z Y\right \} \\
\left (L \left (Z\right )^{ \ddag } \left (X\right )\right )^{ \dag } =X^{ \dag } Z \\
L \left (Z\right )^{ \ddag } \left (X\right ) =Z^{ \dag } X =L \left (Z^{ \dag }\right ) \left (X\right )\text{.}\end{gather}This shows that
\begin{equation}L \left (Z\right )^{ \ddag } =L \left (Z^{ \dag }\right )
\end{equation}and $ \ddag $ is a generalization of $ \dag \text{.}$ 
\end{proof}

\paragraph{General Basis}
In general one can expand any operator $A \in \mathcal{B} \left (\mathcal{F}_{1}\right )$ wrt any basis $\left \{\left .\mathfrak{Z}_{j}\right \vert 1 \leq j \leq 2 r\text{}\right \}$ of $\mathcal{F}_{1}$ as
\begin{equation}A\vert \mathbf{Z} \rangle  =\vert \mathbf{Z} \rangle \mathcal{R}_{\mathfrak{Z}} \left (A\right ) \equiv \vert \mathbf{Z} \rangle \mathbb{A}_{\mathbf{Z}}\text{.}
\end{equation}


\begin{remark}
\begin{align}\vert \mathbf{Z} \rangle ^{\mathfrak{a}} &  = &  \langle \mathbf{Z}\vert  =\left (\mathbf{Z}^{ \dag }\right \vert  \\
\vert \mathbf{Z}_{1} ,\mathbf{Z}_{2} \rangle ^{\mathfrak{a}} &  = &  \langle \mathbf{Z}_{1}^{\#} ,\mathbf{Z}_{2}^{\#}\vert  =\left (\mathbf{Z}_{1}^{ \dag \#} ,\mathbf{Z}_{2}^{ \dag \#}\right \vert \text{.}\end{align}
\end{remark}

Using the HS\ inner product one obtains
\begin{equation}\mathcal{M}_{\mathbf{Z}} \left (A\right ) = \langle \mathbf{Z}\vert A\vert \mathbf{Z} \rangle  = \langle \mathbf{Z}\vert \vert \mathbf{Z} \rangle \mathbb{A}_{\mathbf{Z}} =\mathbb{S}_{\mathbf{Z}} \mathbb{A}_{\mathbf{Z}} \equiv \mathbb{S}_{\mathbf{Z}} \mathcal{R}_{\mathfrak{Z}} \left (A\right )\text{,}
\end{equation}where
\begin{equation}\mathbb{S}_{\mathbf{Z}} = \langle \mathbf{Z}\vert \vert \mathbf{Z} \rangle \text{,}
\end{equation}leading to the following relationship between the representation $\mathcal{R}_{\mathfrak{Z}} \left (A\right )$ and the matrix $\mathcal{M}_{\mathbf{Z}} \left (A\right )$ of $A\;$
\begin{equation}\mathcal{R}_{\mathfrak{Z}} \left (A\right ) =\mathbb{S}_{\mathbf{Z}}^{ -1} \mathcal{M}_{\mathbf{Z}} \left (A\right )
\end{equation}wrt the $\left \{\mathbf{Z}\right \}$ basis. 

$\lceil $The representation of the product is the product of the representations i.e.
\begin{align}\mathcal{R}_{\mathfrak{Z}} \left (A B\right ) &  = & \mathcal{R}_{\mathfrak{Z}} \left (A\right ) \mathcal{R}_{\mathfrak{Z}} \left (B\right ) =\mathbb{S}_{\mathbf{Z}}^{ -1} \mathcal{M}_{\mathbf{Z}} \left (A\right ) \mathbb{S}_{\mathbf{Z}}^{ -1} \mathcal{M}_{\mathbf{Z}} \left (B\right ) \\
\mathcal{R}_{\mathfrak{Z}} \left (A B\right ) &  = & \mathbb{S}_{\mathbf{Z}}^{ -1} \mathcal{M}_{\mathbf{Z}} \left (A B\right )\end{align}thus
\begin{align}\mathcal{M}_{\mathbf{Z}} \left (A B\right ) &  = & \mathcal{M}_{\mathbf{Z}} \left (A\right ) \mathbb{S}_{\mathbf{Z}}^{ -1} \mathcal{M}_{\mathbf{Z}} \left (B\right ) \\
 &  = &  \langle \mathbf{Z}\vert A \left [\vert \mathbf{Z} \rangle  \langle \mathbf{Z}^{t}\vert \vert \mathbf{Z} \rangle ^{ -1} \langle \mathbf{Z}^{t}\vert \right ] B\vert \mathbf{Z} \rangle  = \langle \mathbf{Z}\vert A I_{\mathcal{F}_{H S 1}} B\vert \mathbf{Z} \rangle \text{.}\end{align}$\rfloor $ 

In order to generate expansions of elements of $\mathcal{B} \left (\mathcal{F}_{H S 1}\right )$ one can utilize the \emph{resolution of the identity }$I_{\mathcal{F}_{H S 1}} \in \mathcal{B} \left (\mathcal{F}_{H S 1}\right )$
\begin{equation}I_{\mathcal{F}_{H S 1}} =\vert \mathbf{Z} \rangle  \langle \mathbf{Z}\vert \vert \mathbf{Z} \rangle ^{ -1} \langle \mathbf{Z}\vert  \equiv \vert \mathbf{Z} \rangle \mathbb{S}_{\mathbf{Z}}^{ -1} \langle \mathbf{Z}\vert 
\end{equation}to obtain
\begin{gather}A =I_{\mathcal{F}_{H S 1}} A I_{\mathcal{F}_{H S 1}} =\vert \mathbf{Z} \rangle  \langle \mathbf{Z}\vert \vert \mathbf{Z} \rangle ^{ -1} \langle \mathbf{Z}\vert A\vert \mathbf{Z} \rangle  \langle \mathbf{Z}\vert \vert \mathbf{Z} \rangle ^{ -1}  \langle \mathbf{Z}\vert ^{\mathfrak{t}} =\vert \mathbf{Z} \rangle  \langle \mathbf{Z}\vert \vert \mathbf{Z} \rangle ^{ -1} \mathcal{M}_{\mathbf{Z}} \left (A\right ) \langle \mathbf{Z}\vert \vert \mathbf{Z} \rangle ^{ -1} \langle \mathbf{Z}\vert  \\
 =\fbox{=}\end{gather}i.e.
\begin{gather}A =\vert \mathbf{Z} \rangle \mathbb{S}_{\mathbf{Z}}^{ -1} \mathcal{M}_{\mathbf{Z}} \left (A\right ) \mathbb{S}_{\mathbf{Z}}^{ -1} \langle \mathbf{Z}\vert  =\vert \mathbf{Z} \rangle \mathcal{R}_{\mathfrak{Z}} \left (A\right ) \mathbb{S}_{\mathbf{Z}}^{ -1} \langle \mathbf{Z}\vert  \\
\fbox{=}\text{.}\end{gather}The product of two operators belonging to $\mathcal{B} \left (\mathcal{F}_{1 H S}\right )$ is given by
\begin{equation}A B\vert \mathbf{Z} \rangle  =\vert \mathbf{Z} \rangle \mathbb{A}_{\mathbf{Z}} \mathbb{B}_{\mathbf{Z}} \equiv \vert \mathbf{Z} \rangle \mathcal{R}_{\mathfrak{Z}} \left (A\right ) \mathcal{R}_{\mathfrak{Z}} \left (B\right ) =\vert \mathbf{Z} \rangle \mathcal{R}_{\mathfrak{Z}} \left (A B\right )\text{,}
\end{equation}
\begin{equation}A B =\left [\vert \mathbf{Z} \rangle \mathbb{S}_{\mathbf{Z}}^{ -1} \mathcal{M}_{\mathbf{Z}} \left (A\right ) \mathbb{S}_{\mathbf{Z}}^{ -1} \langle \mathbf{Z}\vert \right ] \left [\vert \mathbf{Z} \rangle \mathbb{S}_{\mathbf{Z}}^{ -1} \mathcal{M}_{\mathbf{Z}} \left (B\right ) \mathbb{S}_{\mathbf{Z}}^{ -1} \langle \mathbf{Z}\vert \right ] =\vert \mathbf{Z} \rangle \mathbb{S}_{\mathbf{Z}}^{ -1} \mathcal{M}_{\mathbf{Z}} \left (A\right ) \mathbb{S}_{\mathbf{Z}}^{ -1} \mathcal{M}_{\mathbf{Z}} \left (B\right ) \mathbb{S}_{\mathbf{Z}}^{ -1} \langle \mathbf{Z}\vert 
\end{equation}and
\begin{equation}A B =\vert \mathbf{Z} \rangle \mathcal{R}_{\mathfrak{Z}} \left (A\right ) \mathbb{S}_{\mathbf{Z}}^{ -1} \langle \mathbf{Z}\vert \vert \mathbf{Z} \rangle \mathcal{R}_{\mathfrak{Z}} \left (B\right ) \mathbb{S}_{\mathbf{Z}}^{ -1} \langle \mathbf{Z}\vert  =\vert \mathbf{Z} \rangle \mathcal{R}_{\mathfrak{Z}} \left (A\right ) \mathcal{R}_{\mathfrak{Z}} \left (B\right ) \mathbb{S}_{\mathbf{Z}}^{ -1} \langle \mathbf{Z}\vert  =\vert \mathbf{Z} \rangle \mathcal{R}_{\mathfrak{Z}} \left (A B\right ) \mathbb{S}_{\mathbf{Z}}^{ -1} \langle \mathbf{Z}\vert \text{.}
\end{equation}


\begin{lemma}
Expansion of Inverse
\begin{equation}A^{ -1} =\vert \mathbf{Z} \rangle \mathcal{R}_{\mathfrak{Z}} \left (A\right )^{ -1} \mathbb{S}_{\mathbf{Z}}^{ -1} \langle \mathbf{Z}\vert  =\vert \mathbf{Z} \rangle \mathcal{M}_{\mathbf{Z}} \left (A\right )^{ -1} \langle \mathbf{Z}\vert \text{.}
\end{equation}
\end{lemma}


\begin{proof}
\begin{align}A^{ -1} &  = & \vert \mathbf{Z} \rangle \mathcal{R}_{\mathfrak{Z}} \left (A^{ -1}\right ) \mathbb{S}_{\mathbf{Z}}^{ -1} \langle \mathbf{Z}\vert  =\vert \mathbf{Z} \rangle \mathcal{R}_{\mathfrak{Z}} \left (A\right )^{ -1} \mathbb{S}_{\mathbf{Z}}^{ -1} \langle \mathbf{Z}\vert  =\vert \mathbf{Z} \rangle \left (\mathbb{S}_{\mathbf{Z}}^{ -1} \mathcal{M}_{\mathbf{Z}} \left (A\right )\right )^{ -1} \mathbb{S}_{\mathbf{Z}}^{ -1} \langle \mathbf{Z}\vert  \\
 &  = & \vert \mathbf{Z} \rangle \mathcal{M}_{\mathbf{Z}} \left (A\right )^{ -1} \mathbb{S}_{\mathbf{Z}} \mathbb{S}_{\mathbf{Z}}^{ -1} \langle \mathbf{Z}\vert  =\vert \mathbf{Z} \rangle \mathcal{M}_{\mathbf{Z}} \left (A\right )^{ -1} \langle \mathbf{Z}\vert \text{.}\end{align}
\end{proof}


\begin{lemma}
Expansion of Adjoint
\begin{equation}A^{ \dag } =\vert \mathbf{Z} \rangle \mathcal{R}_{\mathfrak{Z}} \left (A^{ \dag }\right ) \mathbb{S}_{\mathbf{Z}}^{ -1} \langle \mathbf{Z}\vert  =\left (\vert \mathbf{Z} \rangle \mathcal{R}_{\mathfrak{Z}} \left (A\right ) \mathbb{S}_{\mathbf{Z}}^{ -1} \langle \mathbf{Z}\vert \right )^{ \dag } =\vert \mathbf{Z}^{ \dag } \rangle \mathbb{S}_{\mathbf{Z}}^{ -1} \mathcal{R}_{\mathfrak{Z}} \left (A\right )^{ \dag } \langle \mathbf{Z}^{ \dag }\vert  =\vert \mathbf{Z}^{ \dag } \rangle \mathbb{S}_{\mathbf{Z}}^{ -1} \mathcal{M}_{\mathbf{Z}} \left (A\right )^{ \dag } \mathbb{S}_{\mathbf{Z}}^{ -1}\, \langle \mathbf{Z}^{ \dag }\vert \text{.}
\end{equation}
\end{lemma}


\begin{proof}
\begin{align} & \left (\left (\begin{array}{cc}\vert z_{1} \rangle  & \vert z_{2} \rangle \end{array}\right ) \mathcal{R}_{\mathfrak{Z}} \left (A\right ) \mathbb{S}_{\mathbf{Z}}^{ -1} \left (\begin{array}{c} \langle z_{1}\vert  \\
 \langle z_{2}\vert \end{array}\right )\right )^{ \dag } =\left (\vert z_{1} \rangle X_{1 1}  \langle z_{1}\vert  +\vert z_{1} \rangle  X_{1 2}  \langle z_{2}\vert  +\vert z_{2} \rangle  X_{2 1}  \langle z_{1}\vert  +\vert z_{2} \rangle  X_{2 2} \langle z_{2}\vert \right )^{ \dag } \\
 &  =\left (\vert z_{1}^{ \dag } \rangle \overline{X_{1 1}}  \langle z_{1}^{ \dag }\vert  +\vert z_{2}^{ \dag } \rangle  \overline{X_{1 2}}  \langle z_{1}^{ \dag }\vert  +\vert z_{1}^{ \dag } \rangle  \overline{X_{2 1}}  \langle z_{2}^{ \dag }\vert  +\vert z_{2}^{ \dag } \rangle  \overline{X_{2 2}} \langle z_{2}^{ \dag }\vert \right ) \\
 &  =\left (\vert z_{1}^{ \dag } \rangle X_{1 1}^{ \dag }  \langle z_{1}^{ \dag }\vert  +\vert z_{1}^{ \dag } \rangle  X_{1 2}^{ \dag }  \langle z_{2}^{ \dag }\vert  +\vert z_{2}^{ \dag } \rangle  X_{2 1}^{ \dag }  \langle z_{1}^{ \dag }\vert  +\vert z_{2}^{ \dag } \rangle  X_{2 2}^{ \dag } \langle z_{2}^{ \dag }\vert \right ) \\
 &  =\left (\begin{array}{cc}\vert z_{1}^{ \dag } \rangle  & \vert z_{2}^{ \dag } \rangle \end{array}\right ) \left (\mathcal{R}_{\mathfrak{Z}} \left (A\right ) \mathbb{S}_{\mathbf{Z}}^{ -1}\right )^{ \dag } \left (\begin{array}{c} \langle z_{1}^{ \dag }\vert  \\
 \langle z_{2}\vert ^{ \dag }\end{array}\right ) =\left (\begin{array}{cc}\vert z_{1}^{ \dag } \rangle  & \vert z_{2}^{ \dag } \rangle \end{array}\right ) \mathbb{S}_{\mathbf{Z}}^{ -1} \mathcal{R}_{\mathfrak{Z}} \left (A\right )^{ \dag } \left (\begin{array}{c} \langle z_{1}^{ \dag }\vert  \\
 \langle z_{2}^{ \dag }\vert \end{array}\right )\text{.}\end{align}
\end{proof}


\begin{example}
If the basis is self adjoint then
\begin{gather}A^{ \dag } =\left (\left (\begin{array}{cc}\vert \mathbf{x}_{1} \rangle  & \vert \mathbf{x}_{2} \rangle \end{array}\right ) \mathcal{R}_{\mathfrak{Z}} \left (A^{ \dag }\right ) \mathbb{S} \left (\begin{array}{c} \langle \mathbf{x}_{1}\vert  \\
 \langle \mathbf{x}_{2}\vert \end{array}\right )\right )^{ \dag } =\left (\begin{array}{cc}\vert \mathbf{x}_{1} \rangle  & \vert \mathbf{x}_{2} \rangle \end{array}\right ) \left (\mathcal{R}_{\mathfrak{Z}} \left (A^{ \dag }\right ) \mathbb{S}_{\mathbf{Z}}^{ -1}\right )^{ \dag } \left (\begin{array}{c} \langle \mathbf{x}_{1}\vert  \\
 \langle \mathbf{x}_{2}\vert \end{array}\right ) \\
 =\left (\begin{array}{cc}\vert \mathbf{x}_{1} \rangle  & \vert \mathbf{x}_{2} \rangle \end{array}\right ) \mathbb{S}_{\mathbf{Z}}^{ -1} \mathcal{R}_{\mathfrak{Z}} \left (A\right )^{ \dag } \left (\begin{array}{c} \langle \mathbf{x}_{1}\vert  \\
 \langle \mathbf{x}_{2}\vert \end{array}\right ) =\left (\begin{array}{cc}\vert \mathbf{x}_{1} \rangle  & \vert \mathbf{x}_{2} \rangle \end{array}\right ) \mathbb{S}_{\mathbf{Z}}^{ -1} \mathcal{R}_{\mathfrak{Z}} \left (A\right )^{ \dag } \left (\begin{array}{c} \langle \mathbf{x}_{1}\vert  \\
 \langle \mathbf{x}_{2}\vert \end{array}\right )\text{.}\end{gather}
\end{example}


\begin{example}
If the basis is adjoint then
\begin{gather}A^{ \dag } =\left (\left (\begin{array}{cc}\vert \mathbf{a}^{ \dag } \rangle  & \vert \mathbf{a} \rangle \end{array}\right ) \mathcal{R}_{\mathfrak{Z}} \left (A^{ \dag }\right ) \mathbb{S} \left (\begin{array}{c} \langle \mathbf{a}^{ \dag }\vert  \\
 \langle \mathbf{a}\vert \end{array}\right )\right )^{ \dag } =\left (\begin{array}{cc}\vert \mathbf{a}^{ \dag } \rangle  & \vert \mathbf{a} \rangle \end{array}\right ) \left (\begin{array}{cc}\mathbb{O} & \mathbb{I}_{r} \\
\mathbb{I}_{r} & \mathbb{O}\end{array}\right ) \left (\mathcal{R}_{\mathfrak{Z}} \left (A^{ \dag }\right ) \mathbb{S}_{\mathbf{Z}}^{ -1}\right )^{ \dag } \left (\begin{array}{cc}\mathbb{O} & \mathbb{I}_{r} \\
\mathbb{I}_{r} & \mathbb{O}\end{array}\right ) \left (\begin{array}{c} \langle \mathbf{a}^{ \dag }\vert  \\
 \langle \mathbf{a}\vert \end{array}\right ) \\
 =\left (\begin{array}{cc}\vert \mathbf{a}^{ \dag } \rangle  & \vert \mathbf{a} \rangle \end{array}\right ) \left (\begin{array}{cc}\mathbb{O} & \mathbb{I}_{r} \\
\mathbb{I}_{r} & \mathbb{O}\end{array}\right ) \mathbb{S}_{\mathbf{Z}}^{ -1} \mathcal{R}_{\mathfrak{Z}} \left (A\right )^{ \dag } \left (\begin{array}{cc}\mathbb{O} & \mathbb{I}_{r} \\
\mathbb{I}_{r} & \mathbb{O}\end{array}\right ) \left (\begin{array}{c} \langle \mathbf{a}^{ \dag }\vert  \\
 \langle \mathbf{a}\vert \end{array}\right )\text{.}\end{gather}
\end{example}

\subsubsection{Complex Symmetric Inner Product}
One can form another set of transformations based on the complex symmetric inner product $\left (\left .X\right \vert Y\right )_{\ensuremath{\operatorname*{Tr}}}\text{,}$ defined on Hilbert Schmidt operators by
\begin{equation}\left (\left .X\right \vert Y\right )_{\ensuremath{\operatorname*{Tr}}} =\ensuremath{\operatorname*{Tr}} \left \{X Y\right \}\text{,}
\end{equation}given by
\begin{equation}T_{c} =\left (\begin{array}{cc}\left \vert \mathbf{X}_{\mathbf{c}}_{1}\right ) & \left \vert \mathbf{X}_{\mathbf{c}}_{2}\right )\end{array}\right ) \mathbb{T}_{c \mathfrak{X}} \left (\begin{array}{c}\left (\mathbf{X}_{c 1}\right \vert  \\
\left (\mathbf{X}_{c 2}\right \vert \end{array}\right )\text{.}
\end{equation}


\begin{lemma}
If $X \in \mathcal{F}_{H S}$ then $X^{2} \in \mathcal{F}_{H S}$ and $\left \vert \ensuremath{\operatorname*{Tr}} \left \{X^{2}\right \}\right \vert  =\left \vert \left (\left .X\right \vert X\right )_{\ensuremath{\operatorname*{Tr}}}\right \vert  <\infty $ 
\end{lemma}


\begin{proof}
$\mathcal{F}_{H S}$ is a 2-sided self adjoint ideal contained in $\mathcal{F}\text{.}$ 
\end{proof}


\begin{remark}
The two inner products are related by
\begin{equation} \langle \left .X\right \vert Y \rangle _{H S} =\left (\left .X^{ \dag }\right \vert Y\right )_{\ensuremath{\operatorname*{Tr}}}\text{.}
\end{equation}
\end{remark}


\begin{theorem}
$\left (\left .X\right \vert Y\right )_{\ensuremath{\operatorname*{Tr}}}$ is nondegenerate. 


\begin{proof}
degeneracy is defined as $ \exists X \neq 0$ such that
\begin{equation}\left (\left .X\right \vert Y\right )_{\ensuremath{\operatorname*{Tr}}} =0\text{\quad } \forall Y \in \mathcal{F}\text{.}
\end{equation}As the trace is faithful nondegeneracy is proven. 
\end{proof}


\end{theorem}


\begin{theorem}
One can always find a conb wrt $\left (\left .\text{}\right \vert \right )_{\ensuremath{\operatorname*{Tr}}}$ of $\mathcal{F}_{1}$ 


\begin{proof}
Any non singular complex symmetric matrix $\mathbb{S}_{c} \in G L \left (2 r ,\mathbb{C}\right )$ can be transformed into the identity i.e.
\begin{equation}\mathbb{P} \mathbb{S} \mathbb{P}^{t} =\mathbb{I}_{\mathcal{F}_{1}}\text{,}
\end{equation}where $\mathbb{P} \in G L \left (2 r ,\mathbb{C}\right )$ and $\mathbb{S}$ is the metric matrix of any basis of $\mathcal{F}_{1}\text{.}$ 
\end{proof}


\end{theorem}

The operator $T_{c}$ can be further expanded as
\begin{align}T_{c} &  =\left \vert \mathbf{X}_{\mathbf{c}}_{1}\right )\mathbb{T}_{c \mathfrak{X}_{c} 11} \left (\mathbf{X}_{\mathbf{c}}_{1}\right \vert  +\left \vert \mathbf{X}_{\mathbf{c}}_{1}\right ) \mathbb{T}_{c \mathfrak{X}_{c} 12} \left (\mathbf{X}_{\mathbf{c}}_{2}\right \vert  +\left \vert \mathbf{X}_{\mathbf{c}}_{2}\right ) \mathbb{T}_{c \mathfrak{X}_{c} 21} \left (\mathbf{X}_{\mathbf{c}}_{1}\right \vert  +\left \vert \mathbf{X}_{\mathbf{c}}_{2}\right ) \mathbb{T}_{c \mathfrak{X}_{c} 22}\left (\mathbf{X}_{\mathbf{c}}_{2}\right \vert  \\
 &  =\sum _{1 \leq j ,k \leq r}\left \{\left \vert \mathfrak{X}_{c 1 j}\right )\mathbb{T}_{c \mathfrak{X}_{c} 11 j k} \left (\mathfrak{X}_{c 1 k}\right \vert  +\left \vert \mathfrak{X}_{c 1 j}\right ) \mathbb{T}_{c \mathfrak{X}_{c} 12 j k} \left (\mathfrak{X}_{c 2 k}\right \vert  +\left \vert \mathfrak{X}_{c 2 j}\right ) \mathbb{T}_{c \mathfrak{X}_{c} 21 j k} \left (\mathfrak{X}_{c 1 k}\right \vert  +\left \vert \mathfrak{X}_{c 2 j}\right ) \mathbb{T}_{c \mathfrak{X}_{c} 22 j k}\left (\mathfrak{X}_{c 2 k}\right \vert \right \} \\
 &  =\sum _{1 \leq j ,k \leq r}\left \{\mathbb{T}_{c \mathfrak{X}_{c} 11 j k}\left \vert \mathfrak{X}_{c 1 j}\right )\left (\mathfrak{X}_{c 1 k}\right \vert  +\mathbb{T}_{c \mathfrak{X}_{c} 12 j k}\left \vert \mathfrak{X}_{c 1 j}\right )\left (\mathfrak{X}_{c 2 k}\right \vert  +\mathbb{T}_{c \mathfrak{X}_{c} 21 j k}\left \vert \mathfrak{X}_{c 2 j}\right )\left (\mathfrak{X}_{c 1 k}\right \vert  +\mathbb{T}_{c \mathfrak{X}_{c} 22 j k}\left \vert \mathfrak{X}_{c 2 j}\right )\left (\mathfrak{X}_{c 2 k}\right \vert \right \}\text{,}\end{align}where $\left \{\left .\mathbf{X}_{c 1} ,\mathbf{X}_{c 2}\right \vert 1 \leq j \leq 2 r\right \}$ is a conb wrt $\left (\left .\text{}\right \vert \right )_{\ensuremath{\operatorname*{Tr}}}\text{.}$ The transformations $\left \vert X\right )\left (Y\right \vert $ are analogous to Ket-Bra operators and are defined by
\begin{equation}\left \vert X\right )\left (Y\right \vert A =\left (\left .Y\right \vert A\right )_{\ensuremath{\operatorname*{Tr}}} X\text{.}
\end{equation}


\begin{remark}
Being nondegenerate does not mean that $\left (\left .X\right \vert X\right )_{\ensuremath{\operatorname*{Tr}}} \neq 0$ $ \forall X \in \mathcal{F}\text{.}$ Consider the case $r =2$
\begin{equation}X =\left (\begin{array}{cccc}z_{1} & 0 & 0 & 0 \\
0 & z_{2} & 0 & 0 \\
0 & 0 & 0 & 0 \\
0 & 0 & 0 & 0\end{array}\right )
\end{equation}then
\begin{align}\ensuremath{\operatorname*{Tr}} \left \{X^{2}\right \} &  =z_{1}^{2} +z_{2}^{2} =\left (x_{1} +i y_{1}\right ) \left (x_{1} +i y_{1}\right ) +\left (x_{2} +i y_{2}\right ) \left (x_{2} +i y_{2}\right ) \\
 &  =x_{1}^{2} -y_{1}^{2} +2 i \left (x_{1} y_{1}\right ) +x_{2}^{2} -y_{2}^{2} +2 i \left (x_{2} y_{2}\right ) \\
 &  =x_{1}^{2} +x_{2}^{2} -\left (y_{1}^{2} +y_{2}^{2}\right ) +2 i \left (x_{1} y_{1} +x_{2} y_{2}\right )\text{.}\end{align}Let
\begin{align}z_{1} &  =1 +i \\
z_{1}^{2} &  =2 i \\
z_{2} &  =1 -i \\
z_{2}^{2} &  = -2 i\end{align}then
\begin{equation}z_{1}^{2} +z_{2}^{2} =0.
\end{equation}
\end{remark}


\begin{definition}
The transpose operation, "$t ""\text{,}$ in $\mathcal{B} \left (\mathcal{F}_{H S}\right )$ is defined by
\begin{equation}\left (\left .T^{t} X\right \vert Y\right )_{\ensuremath{\operatorname*{Tr}}} =\left (\left .X\right \vert T Y\right )_{\ensuremath{\operatorname*{Tr}}}\text{\quad } \forall Y \in \mathcal{F}_{H S} ,\,X \in \mathcal{F}_{H S}
\end{equation}and
\begin{equation}T \in \mathcal{B} \left (\mathcal{F}_{H S}\right )\text{.}
\end{equation}
\end{definition}


\begin{remark}
There is no unique operation in $\mathcal{B} \left (\mathcal{H}_{\mathcal{F}}\right )$ that generalizes to this transpose operation. 
\end{remark}


\begin{lemma}
The transpose operation defined by $\left (\left .\text{}\right \vert \right )_{\ensuremath{\operatorname*{Tr}}}$ corresponds to the transpose defined in the Clifford Algebra $\mathcal{F}\text{.}$ 
\end{lemma}


\begin{remark}
The inner product $\left (\left .\text{}\right \vert \right )_{\ensuremath{\operatorname*{Tr}}}$ induces inner products in $\left \{\left .\mathcal{F}_{n H S}\right \vert 0 \leq n \leq 2 r\right \}$. 
\end{remark}


\begin{lemma}
The representation $\mathcal{R}_{\mathfrak{X}}$ of $\mathcal{F}$ has the properties
\begin{align}\mathcal{R}_{\mathfrak{X}} \left (T\right )^{t} &  = & \mathcal{R}_{\mathfrak{X}} \left (T^{t}\right )\text{,} \\
\mathcal{R}_{\mathfrak{X}} \left (\overline{T}\right ) &  = & \overline{\mathcal{R}_{\mathfrak{X}} \left (T\right )}\text{.}\end{align}
\end{lemma}


\begin{remark}
As $\mathcal{F}_{H S}$ is a \emph{Hilbert algebra} \emph{complex conjugation has a basis free definition} and is defined
by the adjoint operation $"" \dag ""$and the transpose "$t ""$ through the trace operation by
\begin{equation}\overline{X} =\left (X^{ \dag }\right )^{t}\text{.}
\end{equation}
\end{remark}


\begin{remark}
The orthoganol group $O \left (\mathcal{F}_{1 H S} ,\mathbb{C}\right )$ is defined as
\begin{equation}O \left (\mathcal{F}_{1 H S} ,\mathbb{C}\right ) =\left \{\left .T \in \mathcal{B} \left (\mathcal{F}_{1 H S}\right )\right \vert \left (\left .T \left (X\right )\right \vert T \left (Y\right )\right )_{\ensuremath{\operatorname*{Tr}}} =\left (\left .X\right \vert Y\right )_{\ensuremath{\operatorname*{Tr}}} , \forall X ,Y \in \mathcal{F}_{1 H S}\right \}\text{.}
\end{equation}
\end{remark}


\begin{remark}
The unitary group $U \left (\mathcal{F}_{1 H S} ,\mathbb{C}\right )$ is defined as
\begin{equation}U \left (\mathcal{F}_{1 H S}\right ) =\left \{\left .T \in \mathcal{B} \left (\mathcal{F}_{1 H S}\right )\right \vert  \langle \left .T \left (X\right )\right \vert T \left (Y\right ) \rangle _{H S} = \langle \left .X\right \vert Y \rangle _{H S} , \forall X ,Y \in \mathcal{F}_{1 H S}\right \}\text{.}
\end{equation}
\end{remark}

\paragraph{Generalization of the operation on $\mathcal{F} -$arrays}
One can consider any mapping (not necessarily linear)
\begin{equation}\mathfrak{m} :\mathcal{F} \rightarrow \mathcal{F}
\end{equation}and compose it with
\begin{equation}\# :M \left (m ,n ,\mathcal{F}\right ) \rightarrow M \left (n ,m ,\mathcal{F}\right )\text{,}
\end{equation}to form
\begin{equation}\# \circ \mathfrak{m} :M \left (m ,n ,\mathcal{F}\right ) \rightarrow M \left (n ,m ,\mathcal{F}\right )\text{.}
\end{equation}

\subsubsection{Matrix Representations of the Inner Products in $\mathcal{F}_{1 H S}$}
\begin{enumerate}
\item Complex Symmetric Inner Product \end{enumerate}



\begin{notation}
\begin{equation}A =\left (\begin{array}{ccc}\mathfrak{X}_{1} & \cdots  & \mathfrak{X}_{2 r}\end{array}\right ) \left (\begin{array}{c}a_{1} \\
\vdots  \\
a_{2 r}\end{array}\right ) \equiv \left (\begin{array}{ccc}\left \vert \mathfrak{X}_{1}\right ) & \cdots  & \left \vert \mathfrak{X}_{2 r}\right )\end{array}\right ) \left (\begin{array}{c}a_{1} \\
\vdots  \\
a_{2 r}\end{array}\right ) =\left \vert \mathbf{X}\right )\mathbf{a}\text{.}
\end{equation}
\end{notation}

\paragraph{Invariance Group of the Complex Symmetric $\left (\left .A\right \vert B\right )_{\ensuremath{\operatorname*{Tr}}} =\ensuremath{\operatorname*{Tr}} \left \{A B\right \}$}
\begin{notation}
\begin{align}\left (\left (\begin{array}{ccc}\mathfrak{X}_{1} & \cdots  & \mathfrak{X}_{2 r}\end{array}\right ) \left (\begin{array}{c}a_{1} \\
\vdots  \\
a_{2 r}\end{array}\right )\right \vert  &  = & \left (\left (\begin{array}{ccc}\left (\mathfrak{X}_{1}\right \vert  & \cdots  & \left (\mathfrak{X}_{2 r}\right \vert \end{array}\right ) \left (\begin{array}{c}a_{1} \\
\vdots  \\
a_{2 r}\end{array}\right )\right )^{\#} =\left (\begin{array}{c}a_{1} \\
\vdots  \\
a_{2 r}\end{array}\right )^{t} \left (\begin{array}{ccc}\left (\mathfrak{X}_{1}\right \vert  & \cdots  & \left (\mathfrak{X}_{2 r}\right \vert \end{array}\right )^{\#} \\
 &  = & \left (\begin{array}{ccc}a_{1} & \cdots  & a_{2 r}\end{array}\right ) \left (\begin{array}{c}\left (\mathfrak{X}_{1}\right \vert  \\
\vdots  \\
\left (\mathfrak{X}_{2 r}\right \vert \end{array}\right )\end{align}The complex symmetric inner product in $\mathcal{F}$ can be expanded as
\begin{align}\left (\left .A\right \vert B\right ) &  = & \ensuremath{\operatorname*{Tr}} \left \{A B\right \} =\left (\left .\left (\begin{array}{ccc}\mathfrak{X}_{1} & \cdots  & \mathfrak{X}_{2 r}\end{array}\right ) \left (\begin{array}{c}a_{1} \\
\vdots  \\
a_{2 r}\end{array}\right )\right \vert \left (\begin{array}{ccc}\mathfrak{X}_{1} & \cdots  & \mathfrak{X}_{2 r}\end{array}\right ) \left (\begin{array}{c}b_{1} \\
\vdots  \\
b_{2 r}\end{array}\right )\right ) \\
 &  = & \left (\begin{array}{ccc}a_{1} & \cdots  & a_{2 r}\end{array}\right ) \left (\begin{array}{c}\left (\mathfrak{X}_{1}\right \vert  \\
\vdots  \\
\left (\mathfrak{X}_{2 r}\right \vert \end{array}\right ) \left (\begin{array}{ccc}\left \vert \mathfrak{X}_{1}\right ) & \cdots  & \left \vert \mathfrak{X}_{2 r}\right )\end{array}\right ) \left (\begin{array}{c}b_{1} \\
\vdots  \\
b_{2 r}\end{array}\right ) \\
 &  = & \left (\begin{array}{ccc}a_{1} & \cdots  & a_{2 r}\end{array}\right ) \left (\begin{array}{ccc}\left (\mathfrak{X}_{1}\right \vert \left \vert \mathfrak{X}_{1}\right ) & \cdots  & \left (\mathfrak{X}_{1}\right \vert \left \vert \mathfrak{X}_{2 r}\right ) \\
\vdots  & \ddots  & \vdots  \\
\left (\mathfrak{X}_{2 r}\right \vert \left \vert \mathfrak{X}_{1}\right ) & \cdots  & \left (\mathfrak{X}_{2 r}\right \vert \left \vert \mathfrak{X}_{2 r}\right )\end{array}\right ) \left (\begin{array}{c}b_{1} \\
\vdots  \\
b_{2 r}\end{array}\right ) \\
 &  = & \left (\begin{array}{ccc}a_{1} & \cdots  & a_{2 r}\end{array}\right ) \left (\begin{array}{ccc}\ensuremath{\operatorname*{Tr}} \left \{\mathfrak{X}_{1} \mathfrak{X}_{1}\right \} & \cdots  & \ensuremath{\operatorname*{Tr}} \left \{\mathfrak{X}_{1} \mathfrak{X}_{2 r}\right \} \\
\vdots  & \ddots  & \vdots  \\
\ensuremath{\operatorname*{Tr}} \left \{\mathfrak{X}_{2 r} \mathfrak{X}_{1}\right \} & \cdots  & \ensuremath{\operatorname*{Tr}} \left \{\mathfrak{X}_{2 r} \mathfrak{X}_{2 r}\right \}\end{array}\right ) \left (\begin{array}{c}b_{1} \\
\vdots  \\
b_{2 r}\end{array}\right ) \\
 &  = & \left (\begin{array}{ccc}a_{1} & \cdots  & a_{2 r}\end{array}\right ) \mathbb{S}_{c \mathfrak{X}} \left (\begin{array}{c}b_{1} \\
\vdots  \\
b_{2 r}\end{array}\right ) =\mathbf{a}^{t} \mathbb{S}_{c \mathfrak{X}} \mathbf{b}\end{align}The invariance group of this inner product is $O \left (2 r ,\mathbb{C}\right )$
\begin{align}\left (\left .T \left (A\right )\right \vert T \left (B\right )\right )_{\ensuremath{\operatorname*{Tr}}} &  = & \ensuremath{\operatorname*{Tr}} \left \{T \left (A\right ) T \left (B\right )\right \} =\left (\left .A\right \vert B\right )_{\ensuremath{\operatorname*{Tr}}} \\
 &  = & \left (\left .T \left (\begin{array}{ccc}\mathfrak{X}_{1} & \cdots  & \mathfrak{X}_{2 r}\end{array}\right ) \left (\begin{array}{c}a_{1} \\
\vdots  \\
a_{2 r}\end{array}\right )\right \vert T \left (\begin{array}{ccc}\mathfrak{X}_{1} & \cdots  & \mathfrak{X}_{2 r}\end{array}\right ) \left (\begin{array}{c}b_{1} \\
\vdots  \\
b_{2 r}\end{array}\right )\right ) \\
 &  = & \left (\left .\left (\begin{array}{ccc}\mathfrak{X}_{1} & \cdots  & \mathfrak{X}_{2 r}\end{array}\right ) \mathbb{T}_{\mathfrak{X}} \left (\begin{array}{c}a_{1} \\
\vdots  \\
a_{2 r}\end{array}\right )\right \vert \left (\begin{array}{ccc}\mathfrak{X}_{1} & \cdots  & \mathfrak{X}_{2 r}\end{array}\right ) \mathbb{T}_{\mathfrak{X}} \left (\begin{array}{c}b_{1} \\
\vdots  \\
b_{2 r}\end{array}\right )\right ) \\
 &  = & \left (\begin{array}{ccc}a_{1} & \cdots  & a_{2 r}\end{array}\right ) \mathbb{T}_{\mathfrak{X}}^{t} \left (\begin{array}{c}\left (\mathfrak{X}_{1}\right \vert  \\
\vdots  \\
\left (\mathfrak{X}_{2 r}\right \vert \end{array}\right ) \left (\begin{array}{ccc}\left \vert \mathfrak{X}_{1}\right ) & \cdots  & \left \vert \mathfrak{X}_{2 r}\right )\end{array}\right ) \mathbb{T}_{\mathfrak{X}} \left (\begin{array}{c}b_{1} \\
\vdots  \\
b_{2 r}\end{array}\right ) \\
 &  = & \left (\begin{array}{ccc}a_{1} & \cdots  & a_{2 r}\end{array}\right ) \mathbb{T}_{\mathfrak{X}}^{t} \mathbb{S}_{c \mathfrak{X}} \mathbb{T}_{\mathfrak{X}} \left (\begin{array}{c}b_{1} \\
\vdots  \\
b_{2 r}\end{array}\right ) =\mathbf{a}^{t} \mathbb{T}_{\mathfrak{X}}^{t} \mathbb{S}_{c \mathfrak{X}} \mathbb{T}_{\mathfrak{X}} \mathbf{b} =\mathbf{a}^{t} \mathbb{S}_{c \mathfrak{X}} \mathbf{b}\text{,}\end{align}showing that
\begin{equation}\mathbb{T}_{\mathfrak{X}}^{t} \mathbb{S}_{c \mathfrak{X}} \mathbb{T}_{\mathfrak{X}} =\mathbb{S}_{c \mathfrak{X}}\text{.}
\end{equation}
\end{notation}

The invariance transformations $T \in O \left (\mathcal{F}_{1 H S} ,\mathbb{C}\right ) \cong O \left (2 r ,\mathbb{C}\right )$ can be expressed as inner transformations of $\mathcal{F}$ generated by the Lie algebra $\mathcal{F}_{2}$ i.e.
\begin{equation}T \left (X\right ) =\widehat{o} \left (T\right ) =\exp  \left \{g\right \} X \exp  \left \{ -g\right \}\text{,}
\end{equation}where
\begin{equation}g =\sum _{1 \leq j <k \leq 2 r}g_{j k} x_{j} x_{k} \in \mathcal{F}_{2} ;\,g_{j k} \in \mathbb{C} ,1 \leq j <k \leq 2 r\text{.}
\end{equation}


\begin{enumerate}
\item [2] Hermitian Inner Product $ \langle \left .X\right \vert Y \rangle  =\ensuremath{\operatorname*{Tr}} \left \{X^{ \dag } Y\right \}$ \end{enumerate}



\begin{notation}
\begin{equation}A =\left (\begin{array}{ccc}\mathfrak{X}_{1} & \cdots  & \mathfrak{X}_{2 r}\end{array}\right ) \left (\begin{array}{c}a_{1} \\
\vdots  \\
a_{2 r}\end{array}\right ) \equiv \left (\begin{array}{ccc}\vert \mathfrak{X}_{1} \rangle  & \cdots  & \vert \mathfrak{X}_{2 r} \rangle \end{array}\right ) \left (\begin{array}{c}a_{1} \\
\vdots  \\
a_{2 r}\end{array}\right ) =\vert \mathbf{X} \rangle \mathbf{a}
\end{equation}
\end{notation}


\begin{notation}
\begin{align} \langle \left (\begin{array}{ccc}\vert \mathfrak{X}_{1} \rangle  & \cdots  & \vert \mathfrak{X}_{2 r} \rangle \end{array}\right ) \left (\begin{array}{c}a_{1} \\
\vdots  \\
a_{2 r}\end{array}\right )\vert  &  = & \left (\left (\begin{array}{ccc}\vert \mathfrak{X}_{1} \rangle  & \cdots  & \vert \mathfrak{X}_{2 r} \rangle \end{array}\right ) \left (\begin{array}{c}a_{1} \\
\vdots  \\
a_{2 r}\end{array}\right )\right )^{\mathfrak{a}} =\left (\begin{array}{c}a_{1} \\
\vdots  \\
a_{2 r}\end{array}\right )^{\mathfrak{a}} \left (\begin{array}{ccc}\vert \mathfrak{X}_{1} \rangle  & \cdots  & \vert \mathfrak{X}_{2 r} \rangle \end{array}\right )^{\mathfrak{a}} \\
 &  = & \left (\begin{array}{ccc}\overline{a_{1}} & \cdots  & \overline{a_{2 r}}\end{array}\right ) \left (\begin{array}{c} \langle \mathfrak{X}_{1}\vert  \\
\vdots  \\
 \langle \mathfrak{X}_{2 r}\vert \end{array}\right )\end{align}The Hermitian inner product in $\mathcal{F}_{1 H S}$ can be expanded as
\begin{align} \langle \left .A\right \vert B \rangle  &  = & \ensuremath{\operatorname*{Tr}} \left \{A^{ \dag } B\right \} = \langle \left .\left (\begin{array}{ccc}\mathfrak{X}_{1} & \cdots  & \mathfrak{X}_{2 r}\end{array}\right ) \left (\begin{array}{c}a_{1} \\
\vdots  \\
a_{2 r}\end{array}\right )\right \vert \left (\begin{array}{ccc}\mathfrak{X}_{1} & \cdots  & \mathfrak{X}_{2 r}\end{array}\right ) \left (\begin{array}{c}b_{1} \\
\vdots  \\
b_{2 r}\end{array}\right ) \rangle  \\
 &  = & \left (\begin{array}{ccc}\overline{a_{1}} & \cdots  & \overline{a_{2 r}}\end{array}\right ) \left (\begin{array}{c} \langle \mathfrak{X}_{1}\vert  \\
\vdots  \\
 \langle \mathfrak{X}_{2 r}\vert \end{array}\right ) \left (\begin{array}{ccc}\vert \mathfrak{X}_{1} \rangle  & \cdots  & \vert \mathfrak{X}_{2 r} \rangle \end{array}\right ) \left (\begin{array}{c}b_{1} \\
\vdots  \\
b_{2 r}\end{array}\right ) \\
 &  = & \left (\begin{array}{ccc}\overline{a_{1}} & \cdots  & \overline{a_{2 r}}\end{array}\right ) \left (\begin{array}{ccc} \langle \mathfrak{X}_{1}\vert \vert \mathfrak{X}_{1} \rangle  & \cdots  &  \langle \mathfrak{X}_{1}\vert \left \vert \mathfrak{X}_{2 r}\right ) \\
\vdots  & \ddots  & \vdots  \\
 \langle \mathfrak{X}_{1}\vert \vert \mathfrak{X}_{2 r} \rangle  & \cdots  &  \langle \mathfrak{X}_{2 r}\vert \vert \mathfrak{X}_{2 r} \rangle \end{array}\right ) \left (\begin{array}{c}b_{1} \\
\vdots  \\
b_{2 r}\end{array}\right ) \\
 &  = & \left (\begin{array}{ccc}\overline{a_{1}} & \cdots  & \overline{a_{2 r}}\end{array}\right ) \left (\begin{array}{ccc}\ensuremath{\operatorname*{Tr}} \left \{\mathfrak{X}_{1}^{ \dag } \mathfrak{X}_{1}\right \} & \cdots  & \ensuremath{\operatorname*{Tr}} \left \{\mathfrak{X}_{1}^{ \dag } \mathfrak{X}_{2 r}\right \} \\
\vdots  & \ddots  & \vdots  \\
\ensuremath{\operatorname*{Tr}} \left \{\mathfrak{X}_{2 r}^{ \dag } \mathfrak{X}_{1}\right \} & \cdots  & \ensuremath{\operatorname*{Tr}} \left \{\mathfrak{X}_{2 r}^{ \dag } \mathfrak{X}_{2 r}\right \}\end{array}\right ) \left (\begin{array}{c}b_{1} \\
\vdots  \\
b_{2 r}\end{array}\right ) \\
 &  = & \left (\begin{array}{ccc}\overline{a_{1}} & \cdots  & \overline{a_{2 r}}\end{array}\right ) \mathbb{S}_{\mathfrak{X}} \left (\begin{array}{c}b_{1} \\
\vdots  \\
b_{2 r}\end{array}\right ) =\mathbf{a}^{ \dag } \mathbb{S}_{\mathfrak{X}} \mathbf{b}\end{align}The invariance group of this inner product is $U \left (2 r\right )$
\begin{align} \langle \left .T \left (A\right )\right \vert T \left (B\right ) \rangle  &  = & \ensuremath{\operatorname*{Tr}} \left \{T \left (A\right )^{ \dag } T \left (B\right )\right \} = \langle \left .A\right \vert B \rangle  \\
 &  = &  \langle \left .T \left (\begin{array}{ccc}\mathfrak{X}_{1} & \cdots  & \mathfrak{X}_{2 r}\end{array}\right ) \left (\begin{array}{c}a_{1} \\
\vdots  \\
a_{2 r}\end{array}\right )\right \vert T \left (\begin{array}{ccc}\mathfrak{X}_{1} & \cdots  & \mathfrak{X}_{2 r}\end{array}\right ) \left (\begin{array}{c}b_{1} \\
\vdots  \\
b_{2 r}\end{array}\right ) \rangle  \\
 &  = &  \langle \left .\left (\begin{array}{ccc}\mathfrak{X}_{1} & \cdots  & \mathfrak{X}_{2 r}\end{array}\right ) \mathbb{T}_{\mathfrak{X}} \left (\begin{array}{c}a_{1} \\
\vdots  \\
a_{2 r}\end{array}\right )\right \vert \left (\begin{array}{ccc}\mathfrak{X}_{1} & \cdots  & \mathfrak{X}_{2 r}\end{array}\right ) \mathbb{T}_{\mathfrak{X}} \left (\begin{array}{c}b_{1} \\
\vdots  \\
b_{2 r}\end{array}\right ) \rangle  \\
 &  = & \left (\begin{array}{ccc}a_{1} & \cdots  & a_{2 r}\end{array}\right ) \mathbb{T}_{\mathfrak{X}}^{ \dag } \mathbb{S}_{\mathfrak{X}} \mathbb{T}_{\mathfrak{X}} \left (\begin{array}{c}b_{1} \\
\vdots  \\
b_{2 r}\end{array}\right ) =\mathbf{a}^{ \dag } \mathbb{T}_{\mathfrak{X}}^{ \dag } \mathbb{S}_{\mathfrak{X}} \mathbb{T}_{\mathfrak{X}} \mathbf{b} =\mathbf{a}^{ \dag } \mathbb{S}_{\mathfrak{X}} \mathbf{b}\text{,}\end{align}showing that
\begin{equation}\mathbb{T}_{\mathfrak{X}}^{ \dag } \mathbb{S}_{\mathfrak{X}} \mathbb{T}_{\mathfrak{X}} =\mathbb{S}_{\mathfrak{X}}\text{.}
\end{equation}
\end{notation}

The invariance transformations $T \in U \left (\mathcal{F}_{1 H S}\right ) \cong U \left (2 r\right )$ \emph{cannot in general } be expressed as inner transformations of $\mathcal{F}$ generated by the Lie algebra $\mathcal{F}_{2}$ unless they belong to $O \left (\mathcal{F}_{1 H S} ,\mathbb{R}\right ) \cong U \left (\mathcal{F}_{1 H S}\right ) \cap $ $O \left (\mathcal{F}_{1 H S} ,\mathbb{C}\right ) \equiv $ $O \left (2 r ,\mathbb{R}\right ) \cong U \left (2 r\right ) \cap $ $O \left (2 r ,\mathbb{C}\right )$ i.e.
\begin{equation}T \left (X\right ) =\widehat{u} \left (T\right ) =\exp  \left \{g\right \} X \exp  \left \{ -g\right \}\text{,}
\end{equation}where
\begin{equation}g =\sum _{1 \leq j <k \leq 2 r}g_{j k} x_{j} x_{k} \in \mathcal{F}_{2} ;\,g_{j k} \in \mathbb{R} ,1 \leq j <k \leq 2 r\text{.}
\end{equation}

\section{Cannonical Transformations}
From the preceding one has the following theorem. 


\begin{theorem}
Cannonical Transformations
\begin{equation}T :\mathcal{F}_{1} \rightarrow \mathcal{F}_{1}
\end{equation}belong to the group $O \left (\mathcal{F}_{1 H S} ,\mathbb{R}\right ) =O \left (\mathcal{F}_{1 H S} ,\mathbb{C}\right ) \cap U \left (\mathcal{F}_{1 H S}\right )\text{.}$ 
\end{theorem}


\begin{proof}
As $\mathbb{T} \in U \left (2 r\right )$ and $\mathbb{T} \in O \left (2 r ,\mathbb{C}\right )$
\begin{gather}\mathbb{T}^{t} \mathbb{T} =\mathbb{T}^{ \dag } \mathbb{T} \\
\mathbb{T}^{t} =\mathbb{T}^{ \dag } =\overline{\mathbb{T}}^{t} \\
\Downarrow  \\
\mathbb{T} =\ensuremath{\operatorname*{Re}} \left \{\mathbb{T}\right \}\end{gather}hence
\begin{equation}\mathbb{T} \in O \left (2 r ,\mathbb{R}\right )\text{.}
\end{equation}
\end{proof}

\subsection{Standard Definition}
The standard definition of a \emph{cannonical transformation}
\begin{equation}T :\mathcal{F}_{1} \rightarrow \mathcal{F}_{1}
\end{equation}is that it is a unitary transformation $T \in U \left (\mathcal{F}_{H S 1}\right ) \cong $ $U \left (2 r\right )$ that preserves the CAR's and is generated by
\begin{equation}T \left (\begin{array}{cc}\mathbf{a}^{ \dag } & \mathbf{a}\end{array}\right ) =\left (\begin{array}{cc}\mathbf{a}^{ \dag } & \mathbf{a}\end{array}\right ) \mathbb{T}_{a} =\left (\begin{array}{cc}\mathbf{a}^{ \dag } & \mathbf{a}\end{array}\right ) \mathcal{R}_{a} \left (T\right ) =\left (\begin{array}{cc}\mathbf{a}^{ \dag } & \mathbf{a}\end{array}\right ) \left (\begin{array}{cc}\mathbb{U} & \overline{\mathbb{V}} \\
\mathbb{V} & \overline{\mathbb{U}}\end{array}\right ) =\left (\begin{array}{cc}\mathbf{b}^{ \dag } & \mathbf{b}\end{array}\right )\text{.}
\end{equation}Transformed creation and annihilation operators are given by
\begin{align}b_{k}^{ \dag } &  =\sum _{1 \leq j \leq r}\left (a_{j}^{ \dag } U_{j k} +a_{j} V_{j k}\right ) \\
b_{k} &  =\sum _{1 \leq j \leq r}\left (a_{j}^{ \dag } \bar{V}_{j k} +a_{j} \bar{U}_{j k}\right )\end{align}and the CAR\ relationships become
\begin{align}\delta _{k l} &  =\left \{b_{k}^{ \dag } ,b_{l}\right \} =\left \{\sum _{1 \leq j \leq r}\left (a_{j}^{ \dag } U_{j k} +a_{j} V_{j k}\right ) ,\sum _{1 \leq j \leq r}\left (a_{m}^{ \dag } \bar{V}_{m l} +a_{m} \bar{U}_{m l}\right )\right \} \nonumber  \\
 &  =\sum _{1 \leq j ,m \leq r}\left (\left \{a_{j}^{ \dag } U_{j k} ,a_{m}^{ \dag } \bar{V}_{m l}\right \} +\left \{a_{j}^{ \dag } U_{j k} ,a_{m} \bar{U}_{m l}\right \} +\left \{a_{j} V_{j k} ,a_{m}^{ \dag } \bar{V}_{m l}\right \} +\left \{a_{j} V_{j k} ,a_{m} \bar{U}_{m k}\right \}\right ) \nonumber  \\
 &  =\sum _{1 \leq j ,m \leq r}\left (U_{j k} \left \{a_{j}^{ \dag } ,a_{m}^{ \dag }\right \}\bar{V}_{m l} +U_{j k} \left \{a_{j}^{ \dag } ,a_{m}\right \}\bar{U}_{m l} +V_{j k} \left \{a_{j} ,a_{m}^{ \dag }\right \}\bar{V}_{m l} +V_{j k} \left \{a_{j} ,a_{m}\right \}\bar{U}_{m k}\right ) \nonumber  \\
 &  =\sum _{1 \leq j ,m \leq r}\left (U_{j k} \delta _{j m} \bar{U}_{m l} +V_{j k} \delta _{j m} \bar{V}_{m l}\right ) =\sum _{1 \leq j \leq r}\left (U_{j k} \bar{U}_{j l} +V_{j k} \bar{V}_{j l}\right ) =\sum _{1 \leq j \leq r}\left (U_{l j}^{ \dag } U_{j k} +V_{l j}^{ \dag } V_{j k}\right ) \nonumber  \\
 &  =\left (\mathbb{U}^{ \dag } \mathbb{U}\right )_{l k} +\left (\mathbb{V}^{ \dag } \mathbb{V}\right )_{l k}\text{,}\end{align}i.e.
\begin{equation}\delta _{k l} =\left (\mathbb{U}^{ \dag } \mathbb{U}\right )_{l k} +\left (\mathbb{V}^{ \dag } \mathbb{V}\right )_{l k}\text{.}
\end{equation}In the same way
\begin{align}0 &  =\left \{b_{k}^{ \dag } ,b_{l}^{ \dag }\right \} =\left \{\sum _{1 \leq j \leq r}\left (a_{j}^{ \dag } U_{j k} +a_{j} V_{j k}\right ) ,\sum _{1 \leq j \leq r}\left (a_{m}^{ \dag } U_{m l} +a_{m} V_{m l}\right )\right \} \nonumber  \\
 &  =\sum _{1 \leq j ,m \leq r}\left (U_{j k} \delta _{j m} V_{m l} +V_{j k} \delta _{j m} U_{m l}\right ) =\sum _{1 \leq j \leq r}\left (V_{l j}^{t} U_{j k} +U_{l j}^{t} V_{j k}\right ) \nonumber  \\
 &  =\left (\mathbb{V}^{t} \mathbb{U}\right )_{l k} +\left (\mathbb{U}^{t} \mathbb{V}\right )_{l k}\end{align}and
\begin{equation}0 =\left \{b_{k} ,b_{l}\right \} =\left (\mathbb{V}^{ \dag } \bar{\mathbb{U}}\right )_{l k} +\left (\mathbb{U}^{ \dag } \bar{\mathbb{V}}\right )_{l k}\text{,}
\end{equation}i.e.
\begin{equation}\left (\mathbb{V}^{t} \mathbb{U}\right )_{l k} = -\left (\mathbb{U}^{t} \mathbb{V}\right )_{l k}\text{.}
\end{equation}Thus
\begin{align} & \left (\begin{array}{cc}\mathbb{U} & \overline{\mathbb{V}} \\
\mathbb{V} & \overline{\mathbb{U}}\end{array}\right )^{ \dag } \left (\begin{array}{cc}\mathbb{U} & \overline{\mathbb{V}} \\
\mathbb{V} & \overline{\mathbb{U}}\end{array}\right ) =\left (\begin{array}{cc}\mathbb{U}^{ \dag } & \mathbb{V}^{ \dag } \\
\mathbb{V}^{t} & \mathbb{U}^{t}\end{array}\right ) \left (\begin{array}{cc}\mathbb{U} & \overline{\mathbb{V}} \\
\mathbb{V} & \overline{\mathbb{U}}\end{array}\right ) \\
 &  =\left (\begin{array}{cc}\mathbb{U}^{ \dag } \mathbb{U} +\mathbb{V}^{ \dag } \mathbb{V} & \mathbb{U}^{ \dag } \overline{\mathbb{V}} +\mathbb{V}^{ \dag } \overline{\mathbb{U}} \\
\mathbb{V}^{t} \mathbb{U} +\mathbb{U}^{t} \mathbb{V} & \mathbb{V}^{t} \overline{\mathbb{V}} +\mathbb{U}^{t} \overline{\mathbb{U}}\end{array}\right ) =\left (\begin{array}{cc}\mathbb{I}_{r} & \mathbb{O}_{r} \\
\mathbb{O}_{r} & \mathbb{I}_{r}\end{array}\right )\text{.}\end{align}Hence $T$ and $\mathbb{T}_{a} =\left (\begin{array}{cc}\mathbb{U} & \overline{\mathbb{V}} \\
\mathbb{V} & \overline{\mathbb{U}}\end{array}\right )$ $ \in U \left (2 r\right )$ and
\begin{equation}\mathbb{U}^{ \dag } \mathbb{U} +\mathbb{V}^{ \dag } \mathbb{V}\, =\mathbb{I}_{r}
\end{equation}
\begin{equation}\mathbb{U}^{ \dag } \overline{\mathbb{V}} +\mathbb{V}^{ \dag } \overline{\mathbb{U}} =\mathbb{O}_{r}\text{,}
\end{equation}\emph{which we show below are the conditions that }$\mathbb{T}_{a}$\emph{\ }$ \in O \left (2 r ,\mathbb{R}\right ) \equiv T \in O \left (\mathcal{F}_{1} ,\mathbb{R}\right )\text{.}$ 

\subsection{Representations of Canonical Transformations}
\begin{remark}
\begin{gather}T\vert \mathbf{Z} \rangle  =\vert \mathbf{Z} \rangle \mathcal{R}_{\mathfrak{Z}} \left (T\right ) \\
\mathcal{M}_{\mathbf{Z}} \left (T\right )\, = \langle \mathbf{Z}\vert T\vert \mathbf{Z} \rangle  = \langle \mathbf{Z}\vert \vert \mathbf{Z} \rangle \mathcal{R}_{\mathfrak{Z}} \left (T\right ) =\mathbb{S}_{\mathbf{Z}} \mathcal{R}_{\mathfrak{Z}} \left (T\right )\text{.}\end{gather}
\end{remark}


\begin{remark}
\begin{gather}T =\vert \mathbf{Z} \rangle \mathbb{S}_{\mathbf{Z}}^{ -1} \mathcal{M}_{\mathbf{Z}} \left (T\right ) \mathbb{S}_{\mathbf{Z}}^{ -1} \langle \mathbf{Z}\vert  =\vert \mathbf{Z} \rangle \mathcal{R}_{\mathfrak{Z}} \left (T\right ) \mathbb{S}_{\mathbf{Z}}^{ -1} \langle \mathbf{Z}\vert  \\
\mathbb{S}_{\mathbf{Z}}^{ -1} \mathcal{M}_{\mathbf{Z}} \left (T\right ) \mathbb{S}_{\mathbf{Z}}^{ -1}\, =\,\mathcal{R}_{\mathfrak{Z}} \left (T\right ) \mathbb{S}_{\mathbf{Z}}^{ -1} \\
\mathcal{M}_{\mathbf{Z}} \left (T\right )\, =\frac{1}{2^{r -1}} \left (\begin{array}{cc}\mathbb{U} & \overline{\mathbb{V}} \\
\mathbb{V} & \overline{\mathbb{U}}\end{array}\right ) \\
\mathcal{R}_{\mathfrak{Z}} \left (T\right ) =\left (\begin{array}{cc}\mathbb{U} & \overline{\mathbb{V}} \\
\mathbb{V} & \overline{\mathbb{U}}\end{array}\right )\text{.}\end{gather}
\end{remark}

Canonical transformations can be expanded in an adjoint and a self adjoint \emph{orthogonal
bases wrt the Hilbert Schmidt trace inner product }as
\begin{align}T &  =\left [\left (\begin{array}{cc}\vert \mathbf{a}^{ \dag } \rangle  & \vert \mathbf{a} \rangle \end{array}\right ) \frac{1}{2^{\frac{r -1}{2}}} \left (\begin{array}{cc}\mathbb{U} & \overline{\mathbb{V}} \\
\mathbb{V} & \overline{\mathbb{U}}\end{array}\right ) \left (\begin{array}{c} \langle \mathbf{a}^{ \dag }\vert ^{\#} \\
 \langle \mathbf{a}\vert ^{\#}\end{array}\right )\right ] \left (\begin{array}{cc}\vert \mathbf{a}^{ \dag } \rangle  & \vert \mathbf{a} \rangle \end{array}\right ) \\
 &  =\left [\frac{1}{2^{r -1}} \left (\begin{array}{cc}\vert \mathbf{a}^{ \dag } \rangle  & \vert \mathbf{a} \rangle \end{array}\right ) \left (\begin{array}{cc}\mathbb{U} & \overline{\mathbb{V}} \\
\mathbb{V} & \overline{\mathbb{U}}\end{array}\right ) \left (\begin{array}{c} \langle \mathbf{a}^{ \dag }\vert ^{\#} \\
 \langle \mathbf{a}\vert ^{\#}\end{array}\right )\right ] \left (\begin{array}{cc}\vert \mathbf{a}^{ \dag } \rangle  & \vert \mathbf{a} \rangle \end{array}\right ) =\left (\begin{array}{cc}\vert \mathbf{a}^{ \dag } \rangle  & \vert \mathbf{a} \rangle \end{array}\right ) \left (\begin{array}{cc}\mathbb{U} & \overline{\mathbb{V}} \\
\mathbb{V} & \overline{\mathbb{U}}\end{array}\right )\text{,}\end{align}which corresponds to
\begin{equation}\left (\begin{array}{c} \langle \mathbf{a}^{ \dag }\vert ^{\#} \\
 \langle \mathbf{a}\vert ^{\#}\end{array}\right ) T \left (\begin{array}{cc}\vert \mathbf{a}^{ \dag } \rangle  & \vert \mathbf{a} \rangle \end{array}\right ) =\frac{1}{2^{r -1}} \left (\begin{array}{c} \langle \mathbf{a}^{ \dag }\vert ^{\#} \\
 \langle \mathbf{a}\vert ^{\#}\end{array}\right ) \left [\left (\begin{array}{cc}\vert \mathbf{a}^{ \dag } \rangle  & \vert \mathbf{a} \rangle \end{array}\right ) \left (\begin{array}{cc}\mathbb{U} & \overline{\mathbb{V}} \\
\mathbb{V} & \overline{\mathbb{U}}\end{array}\right ) \left (\begin{array}{c} \langle \mathbf{a}^{ \dag }\vert ^{\#} \\
 \langle \mathbf{a}\vert ^{\#}\end{array}\right )\right ] \left (\begin{array}{cc}\vert \mathbf{a}^{ \dag } \rangle  & \vert \mathbf{a} \rangle \end{array}\right )\text{.}
\end{equation}


\begin{remark}
The $K$ map transforms the $\left \{\mathbf{a}^{ \dag } ,\mathbf{a}\right \}$ basis into the $\left \{\mathbf{x}_{1} ,\mathbf{x}_{2}\right \}$ basis and has the effect of transforming $\mathcal{F}$ from a CAR\ algebra into a Clifford algebra and is defined by
\begin{gather}\left (\begin{array}{cc}\mathbf{x}_{1} & \mathbf{x}_{2}\end{array}\right ) =K \left (\begin{array}{cc}\mathbf{a}^{ \dag } & \mathbf{a}\end{array}\right ) =\left (\begin{array}{cc}\mathbf{a}^{ \dag } & \mathbf{a}\end{array}\right ) \mathbb{K}_{a} =\left (\begin{array}{cc}\mathbf{a}^{ \dag } & \mathbf{a}\end{array}\right ) \left (\begin{array}{cc}\frac{1}{\sqrt{2}} \mathbb{I}_{r} & \frac{i}{\sqrt{2}} \mathbb{I}_{r} \\
\frac{1}{\sqrt{2}} \mathbb{I}_{r} &  -\frac{i}{\sqrt{2}} \mathbb{I}_{r}\end{array}\right ) \\
\mathcal{R}_{a} \left (K\right ) =\mathbb{K}_{a} =\left (\begin{array}{cc}\frac{1}{\sqrt{2}} \mathbb{I}_{r} & \frac{i}{\sqrt{2}} \mathbb{I}_{r} \\
\frac{1}{\sqrt{2}} \mathbb{I}_{r} &  -\frac{i}{\sqrt{2}} \mathbb{I}_{r}\end{array}\right )\text{.}\end{gather}
\end{remark}


\begin{lemma}
$(1)$ $K \in U \left (2 r\right ) ;\text{\quad }(2) K \notin O \left (2 r ,\mathbb{C}\right )\text{.}$ 
\end{lemma}


\begin{proof}
\begin{enumerate}
\item 

\item
\begin{equation}\left (\begin{array}{cc}\frac{1}{\sqrt{2}} \mathbb{I}_{r} & \frac{i}{\sqrt{2}} \mathbb{I}_{r} \\
\frac{1}{\sqrt{2}} \mathbb{I}_{r} &  -\frac{i}{\sqrt{2}} \mathbb{I}_{r}\end{array}\right )^{ \dag } \left (\begin{array}{cc}\mathbb{I}_{r} & \mathbb{O}_{r} \\
\mathbb{O}_{r} & \mathbb{I}_{r}\end{array}\right ) \left (\begin{array}{cc}\frac{1}{\sqrt{2}} \mathbb{I}_{r} & \frac{i}{\sqrt{2}} \mathbb{I}_{r} \\
\frac{1}{\sqrt{2}} \mathbb{I}_{r} &  -\frac{i}{\sqrt{2}} \mathbb{I}_{r}\end{array}\right ) =\frac{1}{2} \left (\begin{array}{cc}\mathbb{I}_{r} & \mathbb{I}_{r} \\
 -i \mathbb{I}_{r} & i \mathbb{I}_{r}\end{array}\right ) \left (\begin{array}{cc}\mathbb{I}_{r} & i \mathbb{I}_{r} \\
\mathbb{I}_{r} &  -i \mathbb{I}_{r}\end{array}\right ) =\left (\begin{array}{cc}\mathbb{I}_{r} & \mathbb{O}_{r} \\
\mathbb{O}_{r} & \mathbb{I}_{r}\end{array}\right )
\end{equation}

Hence
\begin{equation}\mathbb{K}_{a}^{ \dag } \mathbb{S}_{a} \mathbb{K}_{a} =\mathbb{S}_{a}\text{.}
\end{equation}

\item
\begin{align} &  & \left (\begin{array}{cc}\frac{1}{\sqrt{2}} \mathbb{I}_{r} & \frac{i}{\sqrt{2}} \mathbb{I}_{r} \\
\frac{1}{\sqrt{2}} \mathbb{I}_{r} &  -\frac{i}{\sqrt{2}} \mathbb{I}_{r}\end{array}\right )^{t} \left (\begin{array}{cc}\mathbb{O}_{r} & \mathbb{I}_{r} \\
\mathbb{I}_{r} & \mathbb{O}_{r}\end{array}\right ) \left (\begin{array}{cc}\frac{1}{\sqrt{2}} \mathbb{I}_{r} & \frac{i}{\sqrt{2}} \mathbb{I}_{r} \\
\frac{1}{\sqrt{2}} \mathbb{I}_{r} &  -\frac{i}{\sqrt{2}} \mathbb{I}_{r}\end{array}\right ) =\frac{1}{2} \left (\begin{array}{cc}\mathbb{I}_{r} & \mathbb{I}_{r} \\
i \mathbb{I}_{r} &  -i \mathbb{I}_{r}\end{array}\right ) \left (\begin{array}{cc}\mathbb{O}_{r} & \mathbb{I}_{r} \\
\mathbb{I}_{r} & \mathbb{O}_{r}\end{array}\right ) \left (\begin{array}{cc}\mathbb{I}_{r} & i \mathbb{I}_{r} \\
\mathbb{I}_{r} &  -i \mathbb{I}_{r}\end{array}\right ) \\
 &  = & \frac{1}{2} \left (\begin{array}{cc}\mathbb{I}_{r} & \mathbb{I}_{r} \\
 -i \mathbb{I}_{r} & i \mathbb{I}_{r}\end{array}\right ) \left (\begin{array}{cc}\mathbb{I}_{r} & i \mathbb{I}_{r} \\
\mathbb{I}_{r} &  -i \mathbb{I}_{r}\end{array}\right ) =\left (\begin{array}{cc}\mathbb{I}_{r} & \mathbb{O}_{r} \\
\mathbb{O}_{r} & \mathbb{I}_{r}\end{array}\right )\text{.}\end{align}Hence
\begin{equation}\mathbb{K}_{a}^{t} \mathbb{S}_{c a} \mathbb{K}_{a} \neq \mathbb{S}_{c a}
\end{equation}thus
\begin{equation}K \notin O \left (2 r ,\mathbb{C}\right )\text{.}
\end{equation}\end{enumerate}
\end{proof}


\begin{remark}
In general one can consider $\mathcal{F}_{1}$ as sub Banach space of $\mathcal{F}$ wrt to the operator norm generated by $\mathcal{H}_{\mathcal{F}}$ and $K$ as a Banach space isometry. Both bases $\left \{\left .a_{j}^{ \dag } ,a_{j}\right \vert 1 \leq j \leq r\text{}\right \}$ and $\left \{\left .x_{j}\right \vert 1 \leq j \leq 2 r\text{}\right \}$ are formed by bounded operators, but not in general Hilbert-Schimdt (thus not Trace Class). The preceding expansions of $K$ can be generalized by using $\mathcal{F}_{1} \otimes \mathcal{F}_{1}^{d}$ instead of $\mathcal{F}_{1 H S} \otimes \mathcal{F}_{1 H S}\text{.}$ 
\end{remark}

\subsection{Invariance Groups of Clifford and CAR Algebras}
In \emph{finite dimensions} the CAR and Clifford generated by algebras $\left \{\left .a_{j}^{ \dag } ,a_{j}\right \vert 1 \leq j \leq r\text{}\right \}$ and $\left \{\left .x_{j}\right \vert 1 \leq j \leq 2 r\text{}\right \}$ have the following identical properties: 


\begin{enumerate}
\item both are complex vector spaces of the same dimension, 

\item their
algebraic multiplication is identical, 

\item they have identical normalized trace and they both have
the same adjoint operation defined. \end{enumerate}


However they have different invariance groups, the
CAR algebra has the invariance group $O \left (2 r ,\mathbb{R}\right )\text{,}$ while the Clifford algebra has the invariance group $O \left (2 r ,\mathbb{C}\right )\text{,}$ that obviously contains $O \left (2 r ,\mathbb{R}\right )$. 

The properties of their invariance groups are: 


\begin{enumerate}
\item The invariance group of the CAR algebra has the properties 


\begin{enumerate}
\item It maps $\mathcal{F}_{1}$ to $\mathcal{F}_{1}$ 

\item It conserves the CAR's 

\item It
is an algebraic homomorphism 

\item It is a star mapping 

\item It
is isometric \end{enumerate}


\item The invariance group of the Clifford
algebra has the properties 


\begin{enumerate}
\item It maps $\mathcal{F}_{1}$ to $\mathcal{F}_{1}$ 

\item It conserves the Clifford Commutation Relationships 

\item It
is an algebraic homomorphism \end{enumerate}
\end{enumerate}



\begin{remark}
If $T$ is Clifford Algebra isomorphism i.e belongs to $O \left (\mathcal{F}_{1} ,\mathbb{C}\right ) \cong O \left (2 r ,\mathbb{C}\right )$ and its matrix representation wrt an orthonormal complex symmetric inner product basis has the form
\begin{equation}\mathbb{T} =\left (\begin{array}{cc}\mathbb{T}_{1 1} & \mathbb{T}_{1 2} \\
\mathbb{T}_{2 1} & \mathbb{T}_{2 2}\end{array}\right ) \in O \left (2 r ,\mathbb{C}\right )\text{,}
\end{equation}hence
\begin{equation}\mathbb{T} =\left (\begin{array}{cc}\cos  \left \{\mathbf{\eta }\right \} &  -\sin  \left \{\mathbf{\eta }\right \} \\
\sin  \left \{\mathbf{\eta }\right \} & \cos  \left \{\mathbf{\eta }\right \}\end{array}\right ) ,\,\mathbf{\eta } \in M \left (r ,r ,\mathbb{C}\right )\text{.}
\end{equation}
\end{remark}

$\lceil $
\begin{equation}\left (\begin{array}{cc}\cos  z & \sin  z \\
 -\sin  z & \cos  z\end{array}\right ) \left (\begin{array}{cc}\cos  z &  -\sin  z \\
\sin  z & \cos  z\end{array}\right ) =\left (\begin{array}{cc}\cos ^{2} z +\sin ^{2} z & \sin  z \cos  z -\cos  z \sin  z \\
\sin  z \cos  z -\cos  z \sin  z & \cos ^{2} z +\sin ^{2} z\end{array}\right ) =\left (\begin{array}{cc}1 & 0 \\
0 & 1\end{array}\right )
\end{equation}$\rfloor $ 


\begin{remark}
If $T$ is canonical i.e belongs to $O \left (\mathcal{F}_{1} ,\mathbb{R}\right ) \cong O \left (2 r ,\mathbb{R}\right )$ and it is a CAR algebra isomorphism. Its matrix representation wrt an adjoint basis has the form
\begin{equation}\left (\begin{array}{cc}\mathbb{U} & \overline{\mathbb{V}} \\
\mathbb{V} & \overline{\mathbb{U}}\end{array}\right ) \in M \left (2 r ,\mathbb{C}\right )\text{,}
\end{equation}while its matrix representation wrt a self adjoint basis has the form
\begin{equation}\left (\begin{array}{cc}\ensuremath{\operatorname*{Re}} \left (\mathbb{U} +\mathbb{V}\right ) &  -\ensuremath{\operatorname*{Im}} \left (\mathbb{U} +\mathbb{V}\right ) \\
\ensuremath{\operatorname*{Im}} \left (\mathbb{U} -\mathbb{V}\right ) & \ensuremath{\operatorname*{Re}} \left (\mathbb{U} -\mathbb{V}\right )\end{array}\right ) \in M \left (2 r ,\mathbb{R}\right )\text{.}
\end{equation}
\end{remark}


\begin{remark}
$\text{.}$ 


\begin{enumerate}
\item All \emph{adjoint basis} are on \emph{one }orbit, $\mathfrak{O}_{a}\text{,}$ of $O \left (2 r ,\mathbb{R}\right )$ in the set of bases of $\mathcal{F}_{1}\text{.}$ 

\item All \emph{self adjoint basis} are on \emph{one}
orbit, $\mathfrak{O}_{s a}\text{,}$of $O \left (2 r ,\mathbb{R}\right )$ in the set of bases of $\mathcal{F}_{1}\text{.}$ \end{enumerate}
\end{remark}

\subsection{Induced Map $ \boxtimes T \equiv T_{\mathcal{F}}$}
In finite dimensions $\mathcal{F}$ is a polynomial algebra
\begin{equation}\mathcal{F} =\mathcal{P} \left (\left \{\left .a_{j}^{ \dag } ,a_{j}\right \vert 1 \leq j \leq r\right \}\right ) =\mathcal{P} \left (\left \{\left .x_{j}\right \vert 1 \leq j \leq 2 r\right \}\right )\text{,}
\end{equation}where $\left \{\left .x_{j}\right \vert 1 \leq j \leq 2 r\right \}$ is a self adjoint basis and $\left \{\left .a_{j}^{ \dag } ,a_{j}\right \vert 1 \leq j \leq r\right \}$ an adjoint one. In infinite dimensions topologies and convergence must be considered and play a significant role. 

As
a \emph{vector space} $\mathcal{F}$ one has two \emph{filterings }of $\mathcal{F}$ based on the \emph{graded} decompositions \emph{\ }
\begin{equation}\mathcal{F} = \tbigwedge \left (\mathcal{F}_{1}\right ) =\tbigoplus ??? \tbigwedge ^{n}\mathcal{F}_{1} =\tbigoplus ???\mathcal{F}_{n}
\end{equation}and
\begin{equation}\mathcal{F} =\tbigoplus ???\mathfrak{F}_{n m}\text{,}
\end{equation}where 

\begin{equation}\mathcal{F}_{n} =\mathbb{C}\text{-}\ensuremath{\operatorname*{linspan}}\left \{\left .x_{j_{1}} \cdots  x_{j_{n}}\right \vert 1 \leq j_{1} <\cdots  <j_{n} \leq 2 r ;0 \leq n \leq 2 r\right \}
\end{equation}and
\begin{equation}\mathfrak{F}_{n m} =\mathbb{C}\text{-}\ensuremath{\operatorname*{linspan}}\left \{\left .a_{j_{1}}^{ \dag } \cdots  a_{j_{n}}^{ \dag } a_{k_{1}} \cdots  a_{k_{m}}\right \vert 1 \leq j_{1} <\cdots  <j_{n} \leq r ,1 \leq k_{1} <\cdots  <k_{m} \leq r ;0 \leq n \leq r ,0 \leq m \leq r\right \}\text{.}
\end{equation}


\begin{proposition}
One only needs \emph{normal} ordered products as the identity operator $I_{\mathcal{F}}$ is part of the basis and the CAR's allows one to express all products in \emph{normal} ordered form. 
\end{proposition}

A self adjoint basis $\left \{\mathbf{x}_{1} ,\mathbf{x}_{2}\right \}$ of $\mathcal{F}_{1}$ can be produced from an adjoint basis $\left \{\mathbf{a}^{ \dag } ,\mathbf{a}\right \}$ of $\mathcal{F}_{1}$ by the unitary transformation $K \in U \left (\mathcal{F}_{1}\right ) \cong U \left (2 r\right )$
\begin{equation}\left (\begin{array}{cc}\mathbf{x}_{1} & \mathbf{x}_{2}\end{array}\right ) =K \left (\begin{array}{cc}\mathbf{a}^{ \dag } & \mathbf{a}\end{array}\right ) =\left (\begin{array}{cc}\mathbf{a}^{ \dag } & \mathbf{a}\end{array}\right ) \mathbb{K}_{a} =\left (\begin{array}{cc}\mathbf{a}^{ \dag } & \mathbf{a}\end{array}\right ) \left (\begin{array}{cc}\frac{1}{\sqrt{2}} \mathbb{I}_{r} & \frac{1}{\sqrt{2}} \mathbb{I}_{r} \\
 -\frac{i}{\sqrt{2}} \mathbb{I}_{r} & \frac{i}{\sqrt{2}} \mathbb{I}_{r}\end{array}\right )\text{.}
\end{equation}

These two gradings are related by
\begin{align}\mathcal{F}_{0} &  =\mathfrak{F}_{0 0} \\
\mathcal{F}_{1} &  =\mathfrak{F}_{1 0} \oplus \mathfrak{F}_{0 1} \\
\mathcal{F}_{2} &  =\mathfrak{F}_{0 0} \oplus \mathfrak{F}_{2 0} \oplus \mathfrak{F}_{1 1} \oplus \mathfrak{F}_{0 2} \\
\mathcal{F}_{3} &  =\mathfrak{F}_{0 0} \oplus \mathfrak{F}_{3 0} \oplus \mathfrak{F}_{1 2} \oplus \mathfrak{F}_{2 1} \oplus \mathfrak{F}_{0 3} \\
\mathcal{F}_{4} &  =\mathfrak{F}_{0 0} \oplus \mathfrak{F}_{4 0} \oplus \mathfrak{F}_{1 3} \oplus \mathfrak{F}_{3 1} \oplus \mathfrak{F}_{2 2} \oplus \mathfrak{F}_{0 4}\end{align}leading to
\begin{equation}\mathcal{F}_{n} =\mathfrak{F}_{0 0} \oplus \tbigoplus ???\mathfrak{F}_{p n -p} ;\text{\quad }2 \leq n \leq r\text{.}
\end{equation}


\begin{remark}
The vector space
\begin{equation}\mathcal{F}_{p a r a} =\mathcal{F}_{\left (1\right )} =\mathcal{F}_{0} \oplus \mathcal{F}_{1}
\end{equation}is the space of \emph{paravectors}, which has the invariance group $G L \left (\mathcal{F}_{0} \oplus \mathcal{F}_{1} ,\mathbb{C}\right ) \cong G L \left (2 r +1 ,\mathbb{C}\right )\text{.}$ 
\end{remark}


\begin{remark}
The vector space
\begin{equation}\mathcal{F}_{p a r a} =\mathcal{F}_{\left (1\right )} =\mathcal{F}_{0} \oplus \mathcal{F}_{1} =\mathfrak{F}_{0 0} \oplus \mathfrak{F}_{1 0} \oplus \mathfrak{F}_{0 1}
\end{equation}is left invaraint by the $c^{ \ast }$-algebra inner automorphism group $O \left (\mathcal{F}_{0} \oplus \mathcal{F}_{1} ,\mathbb{R}\right ) \equiv O \left (2 r +1 ,\mathbb{R}\right )\text{.}$ 
\end{remark}


\begin{remark}
The vector space $\mathcal{F}_{1}$ is left invariant by the $c^{ \ast }$-algebra inner automorphism group $O \left (\mathcal{F}_{1} ,\mathbb{R}\right ) \equiv O \left (2 r ,\mathbb{R}\right )$. 
\end{remark}


\begin{proposition}
If $T$ is bounded linear map
\begin{equation}T :\mathcal{F}_{1} \rightarrow \mathcal{F}_{1}
\end{equation}then it induces a bounded linear map
\begin{equation}T_{\mathcal{F}}  \equiv  \boxtimes T :\mathcal{F} \rightarrow \mathcal{F}\text{,}
\end{equation}that leaves the subspaces $\mathcal{F}_{n} ,0 \leq n \leq 2 r$ invariant i.e.
\begin{equation} \boxtimes T :\mathcal{F}_{n} \rightarrow \mathcal{F}_{n}
\end{equation}and
\begin{equation}\left . \boxtimes T\right \vert _{\mathcal{F}_{n}} = \boxtimes ^{n}T\text{.}
\end{equation}


\begin{proof}
In finite dimensions if $T \in G L \left (\mathcal{F}_{1}\right )$ then $ \boxtimes T \in G L \left (\mathcal{F}\right )\text{.}$ 


\begin{lemma}
$ \boxtimes T$ is an exterior algebra endomorphism i.e
\begin{equation} \boxtimes T \left (X \wedge Y\right ) = \boxtimes T \left (X\right )  \wedge  \boxtimes T \left (Y\right )\text{\quad } \forall X ,Y \in \mathcal{F}\text{.}
\end{equation}


\begin{proof}
It is sufficient to consider $X ,Y \in \mathcal{F}_{1}$ as $\mathcal{F}$ is formed by exterior products of elements of $\mathcal{F}_{1}\text{.}$ In this case we have
\begin{equation} \boxtimes T \left (X \wedge Y\right ) =\left ( \boxtimes ^{2}T\right ) \left (X \wedge Y\right )
\end{equation}and
\begin{equation} \boxtimes T \left (X\right )  \wedge  \boxtimes T \left (Y\right ) = \boxtimes ^{1}T \left (X\right ) \wedge  \boxtimes ^{1}T \left (Y\right )\text{.}
\end{equation}As
\begin{equation}\left ( \boxtimes ^{2}T\right ) \left (X \wedge Y\right )\overset{\triangle }{ =} \boxtimes ^{1}T \left (X\right ) \wedge  \boxtimes ^{1}T \left (Y\right )
\end{equation}one has
\begin{equation} \boxtimes T \left (X \wedge Y\right ) =\left ( \boxtimes ^{2}T\right ) \left (X \wedge Y\right )\text{.}
\end{equation}
\end{proof}


\end{lemma}


\begin{lemma}
$ \boxtimes T$ preserves the two gradings and filterings described previously. 


\begin{proof}
The gradings are exterior algebra gradings induced by the vector $\mathcal{F}_{1}\text{.}$ 
\end{proof}


\end{lemma}


\begin{lemma}
If $T$ is canonical i.e. $T \in O \left (2 r ,\mathbb{R}\right )$ then
\begin{equation} \boxtimes  :O \left (2 r ,\mathbb{R}\right ) \rightarrow \mathcal{B} \left (\mathcal{F}\right )
\end{equation}is a reducible representation of $O \left (2 r ,\mathbb{R}\right )\text{.}$ 


\begin{proof}
$ \boxtimes $ is the direct sum of the representations $\left \{\left . \boxtimes ^{n}\right \vert 0 \leq n \leq 2 r\right \}$ i.e.
\begin{equation} \boxtimes \, =\tbigoplus ??? \boxtimes ^{n}\text{.}
\end{equation}
\end{proof}


\end{lemma}


\end{proof}


\end{proposition}


\begin{theorem}
Canonical transformations, $T \in O \left (\mathcal{F}_{1} ,\mathbb{R}\right ) \cong O \left (2 r ,\mathbb{R}\right )$, induce $c^{ \ast }$-algebra inner automorphisms $ \boxtimes \left (T\right ) \equiv T_{\mathcal{F}}$
\begin{equation} \boxtimes \left (T\right ) \equiv T_{\mathcal{F}} :\mathcal{F} \rightarrow \mathcal{F}
\end{equation}i.e
\begin{align}T_{\mathcal{F}} \left (X^{ \dag }\right ) &  =T_{\mathcal{F}} \left (X\right )^{ \dag } \\
T_{\mathcal{F}} \left (X Y\right ) &  =T_{\mathcal{F}} \left (X\right ) T_{\mathcal{F}} \left (Y\right ) \\
T_{\mathcal{F}} &  \in O \left (\mathcal{F} ,\mathbb{R}\right ) \subset \mathcal{B} \left (\mathcal{F}\right ) ,\text{\quad } \forall X ,Y \in \mathcal{F}\end{align}and
\begin{equation}T_{\mathcal{F}} \left (X\right ) =u X u^{ \dag } ;\text{\quad } \forall X ,u \in \mathcal{F}\text{.}
\end{equation}If $T$ is canonical then 
\end{theorem}


\begin{enumerate}
\item
\begin{equation}\mathbb{T}_{a} =\left (\begin{array}{cc}\mathbb{U} & \overline{\mathbb{V}} \\
\mathbb{V} & \overline{\mathbb{U}}\end{array}\right )\text{,}
\end{equation}

\item
\begin{equation}\mathbb{T}_{x} =\left (\begin{array}{cc}\ensuremath{\operatorname*{Re}} \left (\mathbb{U} +\mathbb{V}\right ) &  -\ensuremath{\operatorname*{Im}} \left (\mathbb{U} +\mathbb{V}\right ) \\
\ensuremath{\operatorname*{Im}} \left (\mathbb{U} -\mathbb{V}\right ) & \ensuremath{\operatorname*{Re}} \left (\mathbb{U} -\mathbb{V}\right )\end{array}\right )\text{.}
\end{equation}

\item
\begin{equation}\mathbb{T}_{x} =\mathbb{K}^{ \dag } \mathbb{T}_{a} \mathbb{K}\text{.}
\end{equation}

One can show explicitly this by computing the action of $\widehat{u}$ the bases $\left \{\left .a_{j}^{ \dag } ,a_{j}\right \vert 1 \leq j \leq r\text{}\right \}$ and $\left \{\left .x_{j}\right \vert 1 \leq j \leq 2 r\text{}\right \}$ of $\mathcal{F}_{1}$ which gives
\begin{align}T \left (\begin{array}{cc}\mathbf{x}_{1} & \mathbf{x}_{2}\end{array}\right ) &  =\widehat{u} \left (\begin{array}{cc}\mathbf{x}_{1} & \mathbf{x}_{2}\end{array}\right ) =\left (\begin{array}{cc}u \mathbf{x}_{1} u^{ \dag } & u \mathbf{x}_{2} u^{ \dag }\end{array}\right ) =\left (\begin{array}{cc}\mathbf{x}_{1} & \mathbf{x}_{2}\end{array}\right ) \left (\begin{array}{cc}\ensuremath{\operatorname*{Re}} \left (\mathbb{U} +\mathbb{V}\right ) &  -\ensuremath{\operatorname*{Im}} \left (\mathbb{U} +\mathbb{V}\right ) \\
\ensuremath{\operatorname*{Im}} \left (\mathbb{U} -\mathbb{V}\right ) & \ensuremath{\operatorname*{Re}} \left (\mathbb{U} -\mathbb{V}\right )\end{array}\right ) \\
T \left (\begin{array}{cc}\mathbf{a}^{ \dag } & \mathbf{a}\end{array}\right ) &  =\widehat{u} \left (\begin{array}{cc}\mathbf{a}^{ \dag } & \mathbf{a}\end{array}\right ) =\left (\begin{array}{cc}u \mathbf{a}^{ \dag } u^{ \dag } & u \mathbf{a} u^{ \dag }\end{array}\right ) \left (\begin{array}{cc}\mathbb{U} & \overline{\mathbb{V}} \\
\mathbb{V} & \overline{\mathbb{U}}\end{array}\right )\text{,}\end{align}where $u \in O \left (2 r ,\mathbb{R}\right )\text{.}$ \end{enumerate}



\begin{corollary}
The transformation
\begin{equation}\widehat{u} :\mathcal{F} \rightarrow \mathcal{F}
\end{equation}is defined by
\begin{equation}\widehat{u} \left (X\right ) =u X u^{ \dag } , \forall \,X \in \mathcal{F}\text{,}
\end{equation}where
\begin{equation}u =e^{{}} \exp  \left \{ \sum _{1 \leq j <k \leq 2 r}g_{j k} x_{j} x_{k}\right \} ;g_{j k} \in \mathbb{R} ,1 \leq j <k \leq 2 r\text{.}
\end{equation}


\begin{proof}
As $ \sum _{1 \leq j <k \leq 2 r}g_{j k} x_{j} x_{k}$ belongs to the Lie Algebra $s o \left (2 r ,\mathbb{R}\right )$ of $S O \left (2 r ,\mathbb{R}\right )$ then $u \in S O \left (2 r ,\mathbb{R}\right )\text{.}$The generated canonical transformation
\begin{equation}\widehat{u} :\mathcal{F} \rightarrow \mathcal{F}
\end{equation}is a representation of $S O \left (2 r ,\mathbb{R}\right )\text{.}$ 
\end{proof}


\end{corollary}


\begin{corollary}
The transformations $\widehat{u}$ leave the subspaces $\mathcal{F}_{n} ,\,0 \leq n \leq 2 r$ invariant and is $c^{ \ast }$-algebra automorphism. 


\begin{proof}
The action of $\widehat{u}$ on self adjoint and adjoint bases of $\mathcal{F}$ is given by
\begin{equation}\widehat{u} \left (x_{j_{1}} \cdots  x_{j_{n}}\right ) =\widehat{u} \left (x_{j_{1}}\right ) \cdots  \widehat{u} \left (x_{j_{n}}\right ) ;1 \leq j_{1} <\cdots  <j_{n} \leq 2 r ,0 \leq n \leq 2 r
\end{equation}and
\begin{equation}\widehat{u} \left (a_{j_{1}}^{ \dag } \cdots  a_{j_{n}}^{ \dag } a_{k_{1}} \cdots  a_{k_{m}}\right ) =\widehat{u} \left (a_{j_{1}}^{ \dag }\right ) \cdots  \widehat{u} \left (a_{j_{n}}^{ \dag }\right ) \widehat{u} \left (a_{k_{1}}\right ) \cdots  \widehat{u} \left (a_{k_{m}}\right ) ;1 \leq j_{1} <\cdots  <j_{n} \leq r ,1 \leq k_{1} <\cdots  <k_{m} \leq r ,0 \leq n ,m \leq 2 r
\end{equation}leading to
\begin{gather}\widehat{u} \left (X Y\right ) =\widehat{u} \left (X\right ) \widehat{u} \left (Y\right )\, \\
\widehat{u} \left (X^{ \dag }\right ) =\widehat{u} \left (X\right )^{ \dag } \\
 \forall X ,Y \in \mathcal{F}\text{.}\end{gather}
\end{proof}


\end{corollary}


\begin{corollary}
\begin{equation} \boxtimes \left (\left .\widehat{u}\right \vert _{\mathcal{F}_{1}}\right ) =\widehat{u}
\end{equation}


\begin{proof}
Consider the action of $ \boxtimes \left (\left .\widehat{u}\right \vert _{\mathcal{F}_{1}}\right )$ and $\widehat{u}$ on a basis of $\mathcal{F}$ i.e.
\begin{align} \boxtimes \left (\left .\widehat{u}\right \vert _{\mathcal{F}_{1}}\right ) \left (x_{j_{1}} \cdots  x_{j_{n}}\right ) &  = \boxtimes ^{n}\left (\left .\widehat{u}\right \vert _{\mathcal{F}_{1}}\right ) \left (x_{j_{1}} \cdots  x_{j_{n}}\right ) =\left (\left (\left .\widehat{u}\right \vert _{\mathcal{F}_{1}}\right ) x_{j_{1}} \wedge \cdots  \wedge \left (\left .\widehat{u}\right \vert _{\mathcal{F}_{1}}\right ) x_{j_{n}}\right ) \\
\widehat{u} \left (x_{j_{1}} \cdots  x_{j_{n}}\right ) &  =\widehat{u} \left (x_{j_{1}}\right ) \cdots  \widehat{u} \left (x_{j_{n}}\right ) \equiv \widehat{u} \left (x_{j_{1}} \wedge \cdots  \wedge x_{j_{n}}\right ) =\widehat{u} \left (x_{j_{1}}\right ) \wedge \cdots  \wedge \widehat{u} \left (x_{j_{n}}\right ) \\
1 &  \leq j_{1} <\cdots  <j_{n} \leq 2 r ,1 \leq n \leq 2 r\end{align}The action of $\widehat{u}$ on the identity is given by zeroth order term
\begin{equation}\widehat{u} \left (I_{\mathcal{F}}\right ) =I_{\mathcal{F}}
\end{equation}and as $ \boxtimes \left (\left .\widehat{u}\right \vert _{\mathcal{F}_{1}}\right )$ is unitary
\begin{equation} \boxtimes ^{0}\left (\left .\widehat{u}\right \vert _{\mathcal{F}_{1}}\right ) \left (I_{\mathcal{F}}\right ) =I_{\mathcal{F}}\text{.}
\end{equation}
\end{proof}


\end{corollary}

$\lceil $In order to study whether a canonical transformation belongs to a special or general linear group we consider
\begin{equation}\det \left \{T T^{ \dag }\right \} =1 =\left \vert \det \left \{T\right \}\right \vert ^{2} \Rightarrow \det \left \{T\right \} =e^{i \theta }\text{.}
\end{equation}In particular if 


\begin{enumerate}
\item
\begin{equation}\mathbb{T} =\left (\begin{array}{cc}\mathbb{U} & \overline{\mathbb{V}} \\
\mathbb{V} & \overline{\mathbb{U}}\end{array}\right ) =\left (\begin{array}{cc}\mathbb{O}_{r} & \mathbb{I}_{r} \\
\mathbb{I}_{r} & \mathbb{O}_{r}\end{array}\right )\text{,}
\end{equation}i.e. a particle-hole transfromation, then
\begin{equation}\det \left \{\left (\begin{array}{cc}\mathbb{O}_{r} & \mathbb{I}_{r} \\
\mathbb{I}_{r} & \mathbb{O}_{r}\end{array}\right )\right \} =\left ( -1\right )^{r}\text{.}
\end{equation}$\lceil $
\begin{align}\det \left \{\left (\begin{array}{cc}0 & 1 \\
1 & 0\end{array}\right )\right \} &  =\left ( -1\right ) \left ( +1\right ) = -1 \\
\det \left \{\left (\begin{array}{cccc}0 & 0 & 0 & 1 \\
0 & 0 & 1 & 0 \\
0 & 1 & 0 & 0 \\
1 & 0 & 0 & 0\end{array}\right )\right \} &  =\left ( -1\right ) \left ( +1\right ) \left ( -1\right ) \left ( +1\right ) =1 \\
\det \left \{\left (\begin{array}{cccccc}0 & 0 & 0 & 0 & 0 & 1 \\
0 & 0 & 0 & 0 & 1 & 0 \\
0 & 0 & 0 & 1 & 0 & 0 \\
0 & 0 & 1 & 0 & 0 & 0 \\
0 & 1 & 0 & 0 & 0 & 0 \\
1 & 0 & 0 & 0 & 0 & 0\end{array}\right )\right \} &  =\left ( -1\right ) \left ( +1\right ) \left ( -1\right ) \left ( +1\right ) \left ( -1\right ) \left ( +1\right ) = -1\end{align}$\rfloor $ 

\item If
\begin{equation}\mathbb{T} =\left (\begin{array}{cc}\mathbb{U} & \overline{\mathbb{V}} \\
\mathbb{V} & \overline{\mathbb{U}}\end{array}\right ) =\left (\begin{array}{cc}\mathbb{O}_{r} & i \mathbb{I}_{r} \\
 -i \mathbb{I}_{r} & \mathbb{O}_{r}\end{array}\right )
\end{equation}then
\begin{equation}\det \left \{\left (\begin{array}{cc}\mathbb{O}_{r} & i \mathbb{I}_{r} \\
 -i \mathbb{I}_{r} & \mathbb{O}_{r}\end{array}\right )\right \} =i^{2 r} \left ( -1\right )^{r} \left ( -1\right )^{r} = -\left ( -1\right )^{2 r} = -1\text{,}
\end{equation}

\item if
\begin{equation}\mathbb{T} =\left (\begin{array}{cc}\mathbb{U} & \overline{\mathbb{V}} \\
\mathbb{V} & \overline{\mathbb{U}}\end{array}\right ) =\left (\begin{array}{cc}\mathbb{I}_{r} & \mathbb{O}_{r} \\
\mathbb{O}_{r} & \mathbb{I}_{r}\end{array}\right )
\end{equation}then
\begin{equation}\det \left \{\left (\begin{array}{cc}\mathbb{I}_{r} & \mathbb{O}_{r} \\
\mathbb{O}_{r} & \mathbb{I}_{r}\end{array}\right )\right \} =1\text{,}
\end{equation}

\item if
\begin{equation}\mathbb{T} =\left (\begin{array}{cc}\mathbb{U} & \overline{\mathbb{V}} \\
\mathbb{V} & \overline{\mathbb{U}}\end{array}\right ) =\left (\begin{array}{cc}i \mathbb{I}_{r} & \mathbb{O}_{r} \\
\mathbb{O}_{r} &  -i \mathbb{I}_{r}\end{array}\right )
\end{equation}then
\begin{equation}\det \left \{\left (\begin{array}{cc}i \mathbb{I}_{r} & \mathbb{O}_{r} \\
\mathbb{O}_{r} &  -i \mathbb{I}_{r}\end{array}\right )\right \} =i^{2 r} \left ( -1\right )^{r} = -\left ( -1\right )^{r} =\left ( -1\right )^{r +1}\text{.}
\end{equation}$\rfloor $ \end{enumerate}



\begin{remark}
$\widehat{u}$ is a spinor representation of $O \left (2 r ,\mathbb{R}\right )$ and produces the \emph{PIN} group, which is the universal covering group of $O \left (2 r ,\mathbb{R}\right )\text{.}$ 
\end{remark}

\paragraph{$O \left (2 r ,\mathbb{C}\right )$}
We can extend the preceding to complex transformations by: 


\begin{theorem}
The elements $u \left (\mathbf{g}\right ) \in \mathcal{F}$ defined by
\begin{equation}u \left (\mathbf{g}\right ) =\exp  \left \{ \sum _{1 \leq j <k \leq 2 r}g_{j k} x_{j} x_{k}\right \} ;g_{j k} \in \mathbb{C} ,1 \leq j <k \leq 2 r\text{,}
\end{equation}are a representaion of the group $O \left (2 r ,\mathbb{C}\right )\text{.}$Inner transformations, $\widehat{u \left (\mathbf{g}\right )}$ of $\mathcal{F}$ generated by $u \left (\mathbf{g}\right )$ and defined by
\begin{equation}\widehat{u \left (\mathbf{g}\right )} \left (X\right ) =u \left (\mathbf{g}\right ) X u \left (\mathbf{g}\right )^{ -1} =u \left (\mathbf{g}\right ) X u \left (\mathbf{g}\right )^{t}
\end{equation}are \underline {\emph{Clifford Algebra}} automorphisms \emph{but NOT} $c^{ \ast }$\emph{-algebra} automorphisms unless $\mathbf{g} \in M \left (2 r ,\mathbb{R}\right )\text{,}$ i.e $u \left (\mathbf{g}\right ) \in O \left (2 r ,\mathbb{R}\right )\text{.}$ 
\end{theorem}


\begin{corollary}
If $\mathbf{g}$ is not real then $\widehat{u \left (\mathbf{g}\right )}$ is not a star mapping i.e in general
\begin{equation}\widehat{u \left (\mathbf{g}\right )} \left (X^{ \dag }\right ) \neq \left (\widehat{u \left (\mathbf{g}\right )} \left (X\right )\right )^{ \dag }\text{.}
\end{equation}. 
\end{corollary}


\begin{corollary}
\emph{Clifford Algebra} automorphisms
\begin{equation}T :\mathcal{F} \rightarrow \mathcal{F}
\end{equation}are Exterior algebra automorphisms
\begin{equation}T : \tbigwedge \left (\mathcal{F}_{1}\right ) \rightarrow  \tbigwedge \left (\mathcal{F}_{1}\right )\text{,}
\end{equation}\emph{BUT }Exterior algebra automorphisms are not necessarily \emph{Clifford Algebra}
automorphisms. 
\end{corollary}


\begin{proposition}
The following are Lie algebras contained in $\mathcal{F}$ 


\begin{enumerate}
\item $\mathcal{F}$ as
\begin{equation}\left [X ,Y\right ] \in \mathcal{F}\text{\thinspace \quad } \forall X ,Y \in \mathcal{F}
\end{equation}generates the Lie algebra $g l \left (\mathcal{F} ,\mathbb{C}\right ) \cong g l \left (2^{2 r} ,\mathbb{C}\right )\text{.}$ 

\item $\mathcal{F}_{\left (2\right )} =\mathcal{F}_{0} \oplus \mathcal{F}_{1} \oplus \mathcal{F}_{2}$
\begin{equation}\left [X ,Y\right ] \in \mathcal{F}_{\left (2\right )}\text{\thinspace \quad } \forall X ,Y \in \mathcal{F}_{\left (2\right )}
\end{equation}generates the Lie algebra $s l \left (\mathcal{F}_{\left (2\right )} ,\mathbb{C}\right ) \cong s l \left (2 r +1 ,\mathbb{C}\right )\text{.}$ 

\item $\mathcal{F}_{\left (0 ,2\right )} =\mathcal{F}_{0} \oplus \mathcal{F}_{2}$
\begin{equation}\left [X ,Y\right ] \in \mathcal{F}_{\left (0 ,2\right )}\text{\thinspace \quad } \forall X ,Y \in \mathcal{F}_{\left (0 ,2\right )}
\end{equation}generates the Lie algebra $s l \left (\mathcal{F}_{\left (0 ,2\right )} ,\mathbb{C}\right ) \cong s l \left (2 r ,\mathbb{C}\right )\text{.}$ 

\item $\mathcal{F}_{0}$
\begin{equation}\left [X ,Y\right ] \in \mathcal{F}_{0}\text{\thinspace \quad } \forall X ,Y \in \mathcal{F}_{0}
\end{equation}generates the Lie algebra $\mathbb{C}\text{.}$ \end{enumerate}
\end{proposition}


\begin{lemma}
The induced map $ \boxtimes \left (K\right )$
\begin{equation} \boxtimes \left (K\right ) :\mathcal{F} \rightarrow \mathcal{F}
\end{equation}is an \emph{Exterior Algebra} automorphism, where the $K$ transformation
\begin{equation}K :\mathcal{F}_{1} \rightarrow \mathcal{F}_{1}
\end{equation}belongs to $U \left (\mathcal{F}_{1}\right ) \cong U \left (2 r\right )$ and is defined in the preceding$\text{.}$ 
\end{lemma}


\begin{lemma}
The induced map $ \boxtimes \left (K\right )$
\begin{equation} \boxtimes \left (K\right ) :\mathcal{F} \rightarrow \mathcal{F}
\end{equation}is \underline {\emph{NOT}} a \emph{star} automorphism, where the $K$ transformation
\begin{equation}K :\mathcal{F}_{1} \rightarrow \mathcal{F}_{1}
\end{equation}belongs to $U \left (\mathcal{F}_{1}\right ) \cong U \left (2 r\right )$ and is defined in the preceding$\text{.}$ 


\begin{proof}
We need to show that
\begin{equation}\left ( \boxtimes \left (K\right ) \left (X^{ \dag }\right )\right ) \neq \left ( \boxtimes \left (K\right ) \left (X\right )\right )^{ \dag }
\end{equation}
\end{proof}


\end{lemma}


\begin{theorem}
The induced map $ \boxtimes \left (K\right )$
\begin{equation} \boxtimes \left (K\right ) :\mathcal{F} \rightarrow \mathcal{F}
\end{equation}is \underline {\emph{NOT}} a \emph{Clifford Algebra} automorphism, where
the $K$ transformation
\begin{equation}K :\mathcal{F}_{1} \rightarrow \mathcal{F}_{1}
\end{equation}belongs to $U \left (\mathcal{F}_{1}\right ) \cong U \left (2 r\right )$ and is defined in the preceding$\text{.}$ 
\end{theorem}


\begin{proof}
We need to show that
\begin{equation}\left ( \boxtimes \left (K\right ) \left (X_{1}\right )\right ) \circ \left ( \boxtimes \left (K\right ) \left (X_{2}\right )\right )  \neq  \boxtimes \left (K\right ) \left (X_{1} X_{2}\right )\text{.}
\end{equation}It is suffient to consider $X_{1} ,X_{2} \in \mathcal{F}_{1}$ as $\mathcal{F}$ is formed by products of elements of $\mathcal{F}_{1}\text{.}$ In this case we have
\begin{equation}\left ( \boxtimes \left (K\right ) \left (X_{1}\right )\right ) =K \left (X_{1}\right ) ;\text{\quad }\left ( \boxtimes \left (K\right ) \left (X_{2}\right )\right ) =K \left (X_{2}\right )\text{.}
\end{equation}
\end{proof}

The Clifford product is given by
\begin{equation}X_{1} X_{2} =\sum _{0 \leq m ,n \leq 1}X_{1 m} X_{2 n}\text{,}
\end{equation}where
\begin{equation}X_{1 m} X_{2 n} =\sum _{j =0}^{\min \left (m ,n\right )} \langle X_{1 m} X_{2 n} \rangle _{\left \vert m -n\right \vert  +2 j}\text{.}
\end{equation}In this case this gives
\begin{equation}X_{1} X_{2} =X_{1 1} X_{2 1} =\sum _{j =0}^{1} \langle X_{1 1} X_{2 1} \rangle _{2 j} = \langle X_{1 1} X_{2 1} \rangle _{0} + \langle X_{1 1} X_{2 1} \rangle _{2}\text{.}
\end{equation}Hence
\begin{align} \boxtimes \left (K\right ) \left (X_{1} X_{2}\right ) &   = \boxtimes \left (K\right ) \left ( \langle X_{1 1} X_{2 1} \rangle _{0} + \langle X_{1 1} X_{2 1} \rangle _{2}\right ) = \boxtimes ^{0}\left (K\right ) \left ( \langle X_{1 1} X_{2 1} \rangle _{0} + \boxtimes ^{2}\left (K\right )  \langle X_{1 1} X_{2 1} \rangle _{2}\right ) \\
 &  = \langle X_{1 1} X_{2 1} \rangle _{0} +K \left (X_{1 1}\right ) \wedge K \left (X_{2 1}\right ) =\left (\left .X_{1}\right \vert X_{2}\right ) +K \left (X_{1}\right ) \wedge K \left (X_{2}\right )\end{align}and
\begin{equation}\left ( \boxtimes \left (K\right ) \left (X_{1}\right )\right ) \circ \left ( \boxtimes \left (K\right ) \left (X_{2}\right )\right ) =K \left (X_{1}\right ) \circ K \left (X_{2}\right ) =\left (\left .K \left (X_{1}\right )\right \vert K \left (X_{2}\right )\right ) +K \left (X_{1}\right ) \wedge K \left (X_{2}\right )\text{.}
\end{equation}Thus
\begin{equation} \boxtimes \left (K\right ) \left (X_{1} X_{2}\right ) =\left ( \boxtimes \left (K\right ) \left (X_{1}\right )\right ) \circ \left ( \boxtimes \left (K\right ) \left (X_{2}\right )\right )
\end{equation}\emph{if and only if }
\begin{equation}K \in O \left (\mathcal{F}_{1} ,\mathbb{C}\right ) \cong O \left (2 r ,\mathbb{C}\right )\text{.}
\end{equation}One \emph{concludes this is untrue} by considering the representation $\mathcal{R}_{a} \left (K\right )$ of $K$
\begin{equation}\mathcal{R}_{a} \left (K\right ) =\mathbb{K}_{a} =\left (\begin{array}{cc}\frac{1}{\sqrt{2}} \mathbb{I}_{r} & \frac{i}{\sqrt{2}} \mathbb{I}_{r} \\
\frac{1}{\sqrt{2}} \mathbb{I}_{r} &  -\frac{i}{\sqrt{2}} \mathbb{I}_{r}\end{array}\right )\text{.}
\end{equation}The product $\mathcal{R}_{a} \left (K\right )^{t} \mathcal{R}_{a} \left (K\right )$ is given by
\begin{equation}\left (\begin{array}{cc}\frac{1}{\sqrt{2}} \mathbb{I}_{r} & \frac{i}{\sqrt{2}} \mathbb{I}_{r} \\
\frac{1}{\sqrt{2}} \mathbb{I}_{r} &  -\frac{i}{\sqrt{2}} \mathbb{I}_{r}\end{array}\right )^{t} \left (\begin{array}{cc}\frac{1}{\sqrt{2}} \mathbb{I}_{r} & \frac{i}{\sqrt{2}} \mathbb{I}_{r} \\
\frac{1}{\sqrt{2}} \mathbb{I}_{r} &  -\frac{i}{\sqrt{2}} \mathbb{I}_{r}\end{array}\right ) =\frac{1}{2} \left (\begin{array}{cc}\mathbb{I}_{r} & \mathbb{I}_{r} \\
i \mathbb{I}_{r} &  -i \mathbb{I}_{r}\end{array}\right ) \left (\begin{array}{cc}\mathbb{I}_{r} & i \mathbb{I}_{r} \\
\mathbb{I}_{r} &  -i \mathbb{I}_{r}\end{array}\right ) =\left (\begin{array}{cc}\mathbb{I}_{r} & \mathbb{O} \\
\mathbb{O} &  -\mathbb{I}_{r}\end{array}\right )\text{,}
\end{equation}which clearly shows that $\mathcal{R}_{a} \left (K\right ) \notin O \left (2 r ,\mathbb{C}\right )\text{.}$ 


\begin{remark}
The metric matrix for the basis $\left \{\mathbf{a}^{ \dag } ,\mathbf{a}\right \}$ wrt the inner product $\left (\left .\text{}\right \vert \right )_{\ensuremath{\operatorname*{Tr}}}$ restricted to $\mathcal{F}_{1 H S}$ is
\begin{equation}\left (\begin{array}{cccccc}\left (\left .a_{1}^{ \dag }\right \vert a_{1}^{ \dag }\right )_{\ensuremath{\operatorname*{Tr}}} & \cdots  & \left (\left .a_{1}^{ \dag }\right \vert a_{r}^{ \dag }\right )_{\ensuremath{\operatorname*{Tr}}} & \left (\left .a_{1}^{ \dag }\right \vert a_{1}\right )_{\ensuremath{\operatorname*{Tr}}} & \cdots  & \left (\left .a_{1}^{ \dag }\right \vert a_{r}\right )_{\ensuremath{\operatorname*{Tr}}} \\
\vdots  & \ddots  & \vdots  & \vdots  & \ddots  & \vdots  \\
\left (\left .a_{r}^{ \dag }\right \vert a_{1}^{ \dag }\right )_{\ensuremath{\operatorname*{Tr}}} & \cdots  & \left (\left .a_{r}^{ \dag }\right \vert a_{r}^{ \dag }\right )_{\ensuremath{\operatorname*{Tr}}} & \left (\left .a_{r}^{ \dag }\right \vert a_{1}\right )_{\ensuremath{\operatorname*{Tr}}} & \cdots  & \left (\left .a_{r}^{ \dag }\right \vert a_{r}\right )_{\ensuremath{\operatorname*{Tr}}} \\
\left (\left .a_{1}\right \vert a_{1}^{ \dag }\right )_{\ensuremath{\operatorname*{Tr}}} & \cdots  & \left (\left .a_{1}\right \vert a_{r}^{ \dag }\right )_{\ensuremath{\operatorname*{Tr}}} & \left (\left .a_{1}\right \vert a_{1}\right )_{\ensuremath{\operatorname*{Tr}}} & \cdots  & \left (\left .a_{1}\right \vert a_{r}\right )_{\ensuremath{\operatorname*{Tr}}} \\
\vdots  & \ddots  & \vdots  & \vdots  & \ddots  & \vdots  \\
\left (\left .a_{r}\right \vert a_{1}^{ \dag }\right )_{\ensuremath{\operatorname*{Tr}}} & \cdots  & \left (\left .a_{r}\right \vert a_{r}^{ \dag }\right )_{\ensuremath{\operatorname*{Tr}}} & \left (\left .a_{r}\right \vert a_{1}\right )_{\ensuremath{\operatorname*{Tr}}} & \cdots  & \left (\left .a_{r}\right \vert a_{r}\right )_{\ensuremath{\operatorname*{Tr}}}\end{array}\right ) =\left (\begin{array}{cc}\left (\left .\mathbf{a}^{ \dag }\right \vert \mathbf{a}^{ \dag }\right )_{\ensuremath{\operatorname*{Tr}}} & \left (\left .\mathbf{a}^{ \dag }\right \vert \mathbf{a}\right )_{\ensuremath{\operatorname*{Tr}}} \\
\left (\left .\mathbf{a}\right \vert \mathbf{a}^{ \dag }\right )_{\ensuremath{\operatorname*{Tr}}} & \left (\left .\mathbf{a}\right \vert \mathbf{a}\right )_{\ensuremath{\operatorname*{Tr}}}\end{array}\right ) =\left (r -1\right ) \left (\begin{array}{cc}\mathbb{O}_{r} & \mathbb{I}_{r} \\
\mathbb{I}_{r} & \mathbb{O}_{r}\end{array}\right )
\end{equation}
\end{remark}


\begin{remark}
The metric matrix for the basis $\left \{\mathbf{x}\right \}$ wrt the inner product $\left (\left .\text{}\right \vert \right )_{\ensuremath{\operatorname*{Tr}}}$ restricted to $\mathcal{F}_{1 H S}$ is
\begin{equation}\left (\begin{array}{ccc}\left (\left .x_{1}\right \vert x_{1}\right )_{\ensuremath{\operatorname*{Tr}}} & \cdots  & \left (\left .x_{1}\right \vert x_{2 r}\right )_{\ensuremath{\operatorname*{Tr}}} \\
\vdots  & \ddots  & \vdots  \\
\left (\left .x_{2 r}\right \vert x_{1}\right )_{\ensuremath{\operatorname*{Tr}}} & \cdots  & \left (\left .x_{2 r}\right \vert x_{2 r}\right )_{\ensuremath{\operatorname*{Tr}}}\end{array}\right ) =\left (r -1\right ) \mathbb{I}_{2 r}
\end{equation}
\end{remark}


\begin{remark}
The transformation
\begin{equation}M :\mathcal{F}_{1} \rightarrow \mathcal{F}_{1}
\end{equation}with representation
\begin{gather}\left (\begin{array}{cc}\mathbf{x}_{1} & \mathbf{x}_{2}\end{array}\right ) =M \left (\begin{array}{cc}\mathbf{a}^{ \dag } & \mathbf{a}\end{array}\right ) =\left (\begin{array}{cc}\mathbf{a}^{ \dag } & \mathbf{a}\end{array}\right ) \mathbb{M}_{a} =\left (\begin{array}{cc}\mathbf{a}^{ \dag } & \mathbf{a}\end{array}\right ) \frac{1}{\sqrt{2}} \left (\begin{array}{cc}\mathbb{I}_{r} &  -\mathbb{I}_{r} \\
\mathbb{I}_{r} & \mathbb{I}_{r}\end{array}\right ) =\frac{1}{\sqrt{2}} \left (\begin{array}{cc}\mathbf{a}^{ \dag } +\mathbf{a} & \mathbf{a} -\mathbf{a}^{ \dag }\end{array}\right ) \\
\mathbf{x}_{1}^{ \dag } =\mathbf{x}_{1} ;\text{\quad }\mathbf{x}_{2}^{ \dag } = -\mathbf{x}_{2}\end{gather}\emph{does not belong to} $O \left (\mathcal{F}_{1} ,\mathbb{R}\right ) \cong O \left (2 r ,\mathbb{R}\right ) \subset O \left (\mathcal{F}_{1} ,\mathbb{C}\right ) \cong O \left (2 r ,\mathbb{C}\right )$ as
\begin{equation}\frac{1}{\sqrt{2}} \left (\begin{array}{cc}\mathbb{I}_{r} &  -\mathbb{I}_{r} \\
\mathbb{I}_{r} & \mathbb{I}_{r}\end{array}\right )^{t} \left (r -1\right ) \left (\begin{array}{cc}\mathbb{O}_{r} & \mathbb{I}_{r} \\
\mathbb{I}_{r} & \mathbb{O}_{r}\end{array}\right ) \frac{1}{\sqrt{2}} \left (\begin{array}{cc}\mathbb{I}_{r} &  -\mathbb{I}_{r} \\
\mathbb{I}_{r} & \mathbb{I}_{r}\end{array}\right ) \neq \left (r -1\right ) \left (\begin{array}{cc}\mathbb{O}_{r} & \mathbb{I}_{r} \\
\mathbb{I}_{r} & \mathbb{O}_{r}\end{array}\right )\text{,}
\end{equation}$\lceil $
\begin{equation}\frac{1}{2} \left (\begin{array}{cc}\mathbb{I}_{r} & \mathbb{I}_{r} \\
 -\mathbb{I}_{r} & \mathbb{I}_{r}\end{array}\right ) \left (\begin{array}{cc}\mathbb{O}_{r} & \mathbb{I}_{r} \\
\mathbb{I}_{r} & \mathbb{O}_{r}\end{array}\right ) \left (\begin{array}{cc}\mathbb{I}_{r} &  -\mathbb{I}_{r} \\
\mathbb{I}_{r} & \mathbb{I}_{r}\end{array}\right ) =\frac{1}{2} \left (\begin{array}{cc}\mathbb{I}_{r} & \mathbb{I}_{r} \\
 -\mathbb{I}_{r} & \mathbb{I}_{r}\end{array}\right ) \left (\begin{array}{cc}\mathbb{I}_{r} & \mathbb{I}_{r} \\
\mathbb{I}_{r} &  -\mathbb{I}_{r}\end{array}\right ) =\left (\begin{array}{cc}\mathbb{I}_{r} & \mathbb{O}_{r} \\
\mathbb{O}_{r} &  -\mathbb{I}_{r}\end{array}\right )
\end{equation}$\rfloor $ 
\end{remark}

even though
\begin{equation}\mathcal{R}_{a} \left (M\right )^{t} \mathcal{R}_{a} \left (M\right ) =\frac{1}{\sqrt{2}} \left (\begin{array}{cc}\mathbb{I}_{r} &  -\mathbb{I}_{r} \\
\mathbb{I}_{r} & \mathbb{I}_{r}\end{array}\right )^{t} \frac{1}{\sqrt{2}} \left (\begin{array}{cc}\mathbb{I}_{r} &  -\mathbb{I}_{r} \\
\mathbb{I}_{r} & \mathbb{I}_{r}\end{array}\right ) =\left (\begin{array}{cc}\mathbb{I}_{r} & \mathbb{O}_{r} \\
\mathbb{O}_{r} & \mathbb{I}_{r}\end{array}\right )\text{.}
\end{equation}


\begin{lemma}
The induced map $ \boxtimes \left (K\right )$ is NOT a $ \ast $ automorphism 


\begin{proof}
We need to show that
\begin{equation}K \left (X^{ \dag }\right ) \neq K \left (X\right )^{ \dag }\text{.}
\end{equation}First we expand $X$ and $X^{ \dag }$
\begin{align}X &  =\sum _{1 \leq j \leq r}\left \{X_{1 j} a_{j}^{ \dag } +X_{2 j} a_{j}\right \} \\
X^{ \dag } &  =\sum _{1 \leq j \leq r}\left \{\overline{X_{1 j}} a_{j} +\overline{X_{2 j}} a_{j}^{ \dag }\right \}\text{,}\end{align}then apply the mapping $K$ and compare the results of taking the adjoint before and after the application obtaining
\begin{align}K \left (X\right ) &  =K \left (\sum _{1 \leq j \leq r}\left \{X_{1 j} a_{j}^{ \dag } +X_{2 j} a_{j}\right \}\right ) =\sum _{1 \leq j \leq r}\left \{X_{1 j} K \left (a_{j}^{ \dag }\right ) +X_{2 j} K \left (a_{j}\right )\right \} =\frac{1}{\sqrt{2}} \sum _{1 \leq j \leq r}\left \{X_{1 j} \left (a_{j}^{ \dag } +a_{j}\right ) +X_{2 j} i \left (a_{j}^{ \dag } -a_{j}\right )\right \} \\
K \left (X^{ \dag }\right ) &  =K \left (\sum _{1 \leq j \leq r}\left \{\overline{X_{1 j}} a_{j} +\overline{X_{2 j}} a_{j}^{ \dag }\right \}\right ) =\sum _{1 \leq j \leq r}\left \{\overline{X_{1 j}} K \left (a_{j}\right ) +\overline{X_{2 j}} K \left (a_{j}^{ \dag }\right )\right \} =\frac{1}{\sqrt{2}} \sum _{1 \leq j \leq r}\left \{\overline{X_{1 j}} i \left (a_{j}^{ \dag } -a_{j}\right ) +\overline{X_{2 j}} \left (a_{j}^{ \dag } +a_{j}\right )\right \} \\
 &  =\frac{1}{\sqrt{2}} \sum _{1 \leq j \leq r}\left \{\overline{X_{1 j}} i a_{j}^{ \dag } -\overline{X_{1 j}} i a_{j} +\overline{X_{2 j}} a_{j}^{ \dag } +\overline{X_{2 j}} a_{j}\right \} =\frac{1}{\sqrt{2}} \sum _{1 \leq j \leq r}\left \{\left (\overline{X_{1 j}} i +\overline{X_{2 j}}\right ) a_{j}^{ \dag } +\left (\overline{X_{2 j}} -\overline{X_{1 j}} i\right ) a_{j}\right \} \\
K \left (X\right )^{ \dag } &  =\frac{1}{\sqrt{2}} \sum _{1 \leq j \leq r}\left \{\overline{X_{1 j}} \left (a_{j} +a_{j}^{ \dag }\right ) -\overline{X_{2 j}} i \left (a_{j} -a_{j}^{ \dag }\right )\right \} =\frac{1}{\sqrt{2}} \sum _{1 \leq j \leq r}\left \{\overline{X_{1 j}} a_{j} +\overline{X_{1 j}} a_{j}^{ \dag } -\overline{X_{2 j}} i a_{j} +\overline{X_{2 j}} i a_{j}^{ \dag }\right \} \\
 &  =\frac{1}{\sqrt{2}} \sum _{1 \leq j \leq r}\left \{\left (\overline{X_{1 j}} +\overline{X_{2 j}} i\right ) a_{j}^{ \dag } +\left (\overline{X_{1 j}} -\overline{X_{2 j}} i\right ) a_{j}\right \}\text{.}\end{align}In fat notation this can be expressed as
\begin{align}X &  =\left (\begin{array}{cc}\mathbf{a}^{ \dag } & \mathbf{a}\end{array}\right ) \left (\begin{array}{c}\mathbf{X}_{1} \\
\mathbf{X}_{2}\end{array}\right ) \\
X^{ \dag } &  =\left (\begin{array}{cc}\mathbf{a} & \mathbf{a}^{ \dag }\end{array}\right ) \left (\begin{array}{c}\overline{\mathbf{X}_{1}} \\
\overline{\mathbf{X}_{2}}\end{array}\right )\end{align}and
\begin{align}K \left (X\right ) &  =\left (\begin{array}{cc}K \left (\mathbf{a}^{ \dag }\right ) & K \left (\mathbf{a}\right )\end{array}\right ) \left (\begin{array}{c}\mathbf{X}_{1} \\
\mathbf{X}_{2}\end{array}\right ) =\left (\begin{array}{cc}\mathbf{a}^{ \dag } & \mathbf{a}\end{array}\right ) \left (\begin{array}{cc}\frac{1}{\sqrt{2}} \mathbb{I}_{r} & \frac{i}{\sqrt{2}} \mathbb{I}_{r} \\
\frac{1}{\sqrt{2}} \mathbb{I}_{r} &  -\frac{i}{\sqrt{2}} \mathbb{I}_{r}\end{array}\right ) \left (\begin{array}{c}\mathbf{X}_{1} \\
\mathbf{X}_{2}\end{array}\right ) \\
K \left (\mathbf{a}^{ \dag }\right ) &  =\frac{1}{\sqrt{2}} \left (\mathbf{a}^{ \dag } +\mathbf{a}\right ) \\
K \left (\mathbf{a}\right ) &  =\frac{i}{\sqrt{2}} \left (\mathbf{a}^{ \dag } -\mathbf{a}\right ) \\
K \left (X\right ) &  =\frac{1}{\sqrt{2}} \left (\begin{array}{cc}\mathbf{a}^{ \dag } +\mathbf{a} & i \left (\mathbf{a}^{ \dag } -\mathbf{a}\right )\end{array}\right ) \left (\begin{array}{c}\mathbf{X}_{1} \\
\mathbf{X}_{2}\end{array}\right )\end{align}
\begin{align*}K \left (X^{ \dag }\right ) &  =\left (\begin{array}{cc}K \left (\mathbf{a}\right ) & K \left (\mathbf{a}^{ \dag }\right )\end{array}\right ) \left (\begin{array}{c}\overline{\mathbf{X}_{1}} \\
\overline{\mathbf{X}_{2}}\end{array}\right ) =\left (\begin{array}{cc}\mathbf{a} & \mathbf{a}^{ \dag }\end{array}\right ) \left (\begin{array}{cc}\frac{1}{\sqrt{2}} \mathbb{I}_{r} & \frac{i}{\sqrt{2}} \mathbb{I}_{r} \\
\frac{1}{\sqrt{2}} \mathbb{I}_{r} &  -\frac{i}{\sqrt{2}} \mathbb{I}_{r}\end{array}\right ) \left (\begin{array}{c}\overline{\mathbf{X}_{1}} \\
\overline{\mathbf{X}_{2}}\end{array}\right ) \\
 &  =\frac{1}{\sqrt{2}} \left (\begin{array}{cc}i \left (\mathbf{a}^{ \dag } -\mathbf{a}\right ) & \mathbf{a}^{ \dag } +\mathbf{a}\end{array}\right ) \left (\begin{array}{c}\mathbf{X}_{1} \\
\mathbf{X}_{2}\end{array}\right )\text{.}\end{align*}Hence
\begin{equation}K \left (X^{ \dag }\right ) \neq K \left (X\right )^{ \dag }
\end{equation}
\end{proof}


\end{lemma}


\begin{corollary}
$K$ is not a canonical transformation 
\end{corollary}


\begin{lemma}
NOT\ TRUE\ By \underline {the \emph{Skolem-Noether
}} theorem one can conclude that in finite dimensions
\begin{equation}K \left (X\right ) =\exp  \left \{i  \sum _{0 \leq n \leq 2 r} \sum _{1 \leq j_{1} <\cdots  <j_{n} \leq 2 r}\lambda _{j_{1} \cdots  j_{n}} x_{j_{1}} \cdots  x_{j_{n}}\right \} X \exp  \left \{ -i  \sum _{0 \leq n \leq 2 r} \sum _{1 \leq j_{1} <\cdots  <j_{n} \leq 2 r}\lambda _{j_{1} \cdots  j_{n}} x_{j_{1}} \cdots  x_{j_{n}}\right \}\text{,}
\end{equation}where
\begin{equation} \sum _{0 \leq n \leq 2 r} \sum _{1 \leq j_{1} <\cdots  <j_{n} \leq 2 r}\lambda _{j_{1} \cdots  j_{n}} x_{j_{1}} \cdots  x_{j_{n}} =\left ( \sum _{0 \leq n \leq 2 r} \sum _{1 \leq j_{1} <\cdots  <j_{n} \leq 2 r}\lambda _{j_{1} \cdots  j_{n}} x_{j_{1}} \cdots  x_{j_{n}}\right )^{ \dag }\text{.}
\end{equation}
\end{lemma}

\section{Endomorphisms of $\mathcal{F}$}


\subsection{General Transformations of $\mathcal{F}_{1}$}
Consider a general isomorphism
\begin{equation}T :\mathcal{F}_{1} \rightarrow \mathcal{F}_{1}
\end{equation}given by
\begin{equation}T \left (\begin{array}{cc}\mathbf{a}^{ \dag } & \mathbf{a}\end{array}\right ) =\left (\begin{array}{cc}T \left (\mathbf{a}^{ \dag }\right ) & T \left (\mathbf{a}\right )\end{array}\right ) =\left (\begin{array}{cc}\mathbf{a}^{ \dag } & \mathbf{a}\end{array}\right ) \left (\begin{array}{cc}\mathbb{T}_{a 1 1} & \mathbb{T}_{a 1 2} \\
\mathbb{T}_{a 2 1} & \mathbb{T}_{a 2 2}\end{array}\right ) =\left (\begin{array}{cc}\mathbf{\beta } & \mathbf{\alpha }\end{array}\right )\text{.}
\end{equation}


\begin{enumerate}
\item If $T$ is a $^{ \ast }$- transformation then
\begin{equation}T (x^{ \dag }) =T (x)^{ \dag }\text{.}
\end{equation}In general
\begin{align}T \left (a_{j}^{ \dag }\right ) &  =\sum _{1 \leq k \leq r}\left \{a_{k}^{ \dag } T_{1 1 k j} +a_{k} T_{2 1 k j}\right \} \\
T \left (a_{j}\right ) &  =\sum _{1 \leq k \leq r}\left \{a_{k}^{ \dag } T_{1 2 k j} +a_{k} T_{2 2 k j}\right \} \\
T \left (a_{j}\right )^{ \dag } &  =\sum _{1 \leq k \leq r}\left \{a_{k}^{ \dag } \overset{\underline{}}{T}_{2 2 k j} +a_{k} \overset{\underline{}}{T}_{1 2 k j}\right \}\end{align}showing that $T$ is a $^{ \ast }$- transformation if
\begin{align}\overset{\underline{}}{T}_{2 2 k j} &  =T_{1 1 k j} \\
T_{2 1 k j} &  =\overset{\underline{}}{T}_{1 2 k j}\end{align}i.e.
\begin{equation}\mathbb{T}_{a} =\left (\begin{array}{cc}\mathbb{T}_{a 1 1} & \overset{\underline{}}{T}_{a 1 2} \\
\mathbb{T}_{a 1 2} & \overset{\underline{}}{T}_{a 1 1}\end{array}\right )\text{.}
\end{equation}

\item In finite dimensions one can always transform to a self adjoint
basis from a non selfadjoint basis using the transformation $K$ (which is a $ \ast $ transformation)
\begin{align*}T &  =\left (\begin{array}{cc}\mathbf{x}_{1} & \mathbf{x}_{2}\end{array}\right ) \mathbb{T}_{x} =\left (\begin{array}{cc}\mathbf{a}^{ \dag } & \mathbf{a}\end{array}\right ) \mathbb{K}_{a} \mathbb{T}_{a} \mathbb{K}_{a}^{ \dag } =\left (\begin{array}{cc}\mathbf{a}^{ \dag } & \mathbf{a}\end{array}\right ) \mathbb{T}_{a} \\
\mathbb{T}_{x} &  =\mathbb{K}_{a} \mathbb{T}_{a} \mathbb{K}_{a}^{ \dag } \\
\mathbb{T}_{x} &  =\left (\begin{array}{cc}\frac{1}{\sqrt{2}} \mathbb{I}_{r} & \frac{1}{\sqrt{2}} \mathbb{I}_{r} \\
 -\frac{i}{\sqrt{2}} \mathbb{I}_{r} & \frac{i}{\sqrt{2}} \mathbb{I}_{r}\end{array}\right ) \left (\begin{array}{cc}\mathbb{T}_{a 1 1} & \overset{\underline{}}{T}_{a 1 2} \\
\mathbb{T}_{a 1 2} & \overset{\underline{}}{T}_{a 1 1}\end{array}\right ) \left (\begin{array}{cc}\frac{1}{\sqrt{2}} \mathbb{I}_{r} & \frac{i}{\sqrt{2}} \mathbb{I}_{r} \\
\frac{1}{\sqrt{2}} \mathbb{I}_{r} & \frac{ -i}{\sqrt{2}} \mathbb{I}_{r}\end{array}\right ) =\left (\begin{array}{cc}\ensuremath{\operatorname*{Re}} \left (\mathbb{T}_{a 1 1} +\mathbb{T}_{a 2 1}\right ) &  -\ensuremath{\operatorname*{Im}} \left (\mathbb{T}_{a 1 1} +\mathbb{T}_{a 2 1}\right ) \\
\ensuremath{\operatorname*{Im}} \left (\mathbb{T}_{a 1 1} -\mathbb{T}_{a 2 1}\right ) & \ensuremath{\operatorname*{Re}} \left (\mathbb{T}_{a 1 1} -\mathbb{T}_{a 2 1}\right )\end{array}\right )\text{.}\end{align*}Hence one has 


\begin{proposition}
\emph{If }$T$\emph{\ is invertible and a }$^{ \ast }$\emph{- transformation }$T \in G L \left (2 r ,\mathbb{R}\right )\text{.}$ 
\end{proposition}

\item $T \in U \left (2 r\right )$ if
\begin{equation*}\left (\begin{array}{cc}\mathbb{T}_{a 1 1} & \mathbb{T}_{a 1 2} \\
\mathbb{T}_{a 2 1} & \mathbb{T}_{a 2 2}\end{array}\right )^{ \dag } \left (\begin{array}{cc}\mathbb{I}_{r} & \mathbb{O}_{r} \\
\mathbb{O}_{r} & \mathbb{I}_{r}\end{array}\right ) \left (\begin{array}{cc}\mathbb{T}_{a 1 1} & \mathbb{T}_{a 1 2} \\
\mathbb{T}_{a 2 1} & \mathbb{T}_{a 2 2}\end{array}\right ) =\left (\begin{array}{cc}\mathbb{I}_{r} & \mathbb{O}_{r} \\
\mathbb{O}_{r} & \mathbb{I}_{r}\end{array}\right )
\end{equation*}

\item $T$ is a $^{ \ast }$- transformation and $T \in U \left (2 r\right )$ if
\begin{equation}\left (\begin{array}{cc}\ensuremath{\operatorname*{Re}} \left (\mathbb{T}_{a 1 1} +\mathbb{T}_{a 2 1}\right ) &  -\ensuremath{\operatorname*{Im}} \left (\mathbb{T}_{a 1 1} +\mathbb{T}_{a 2 1}\right ) \\
\ensuremath{\operatorname*{Im}} \left (\mathbb{T}_{a 1 1} -\mathbb{T}_{a 2 1}\right ) & \ensuremath{\operatorname*{Re}} \left (\mathbb{T}_{a 1 1} -\mathbb{T}_{a 2 1}\right )\end{array}\right )^{t} \left (\begin{array}{cc}\ensuremath{\operatorname*{Re}} \left (\mathbb{T}_{a 1 1} +\mathbb{T}_{a 2 1}\right ) &  -\ensuremath{\operatorname*{Im}} \left (\mathbb{T}_{a 1 1} +\mathbb{T}_{a 2 1}\right ) \\
\ensuremath{\operatorname*{Im}} \left (\mathbb{T}_{a 1 1} -\mathbb{T}_{a 2 1}\right ) & \ensuremath{\operatorname*{Re}} \left (\mathbb{T}_{a 1 1} -\mathbb{T}_{a 2 1}\right )\end{array}\right ) =\left (\begin{array}{cc}\mathbb{I}_{r} & \mathbb{O}_{r} \\
\mathbb{O}_{r} & \mathbb{I}_{r}\end{array}\right )
\end{equation}i.e. $T \in O \left (2 r ,\mathbb{R}\right )\text{.}$ \end{enumerate}


\subsection{Endomorphisms of $\mathcal{F}$ Induced by Endomorphisms of $\mathcal{F}_{1}$}
We mainly concentrate on Automorphisms, which can be classified into inner and outer automorphisms. 


\begin{proposition}
If $T$ is a $^{ \ast }$- transformation and $T \in U \left (2 r\right )$ then it induces a $c^{ \ast }$-algebra $\mathbf{i} \mathbf{n} \mathbf{n} \mathbf{e} \mathbf{r}$ automorphism, $T_{\mathcal{F}}\text{,}$ i.e.
\begin{equation}T_{\mathcal{F}} \left (X Y) =T_{\mathcal{F}} \left (X\right ) T_{\mathcal{F}} \left (Y\right )\right )\text{\quad } \forall X ,Y \in \mathcal{F} ;\text{\quad }T_{\mathcal{F}} \in \mathcal{B} \left (\mathcal{F}\right )
\end{equation}and
\begin{equation}T_{\mathcal{F}} \left (X\right ) =u X u^{ \dag }\text{,}
\end{equation}where
\begin{equation}u \in U \left (2 r\right )\text{and}u \in \mathcal{B} \left (\mathcal{F}\right ) \cong \mathcal{B} \left (\mathbb{C}^{2^{r}}\right )\text{.}
\end{equation}
\end{proposition}

\subsection{CAR Preserving Transformations}


\subsubsection{Adjoint Form of the CAR}
\begin{notation}
In matrix form the Fermion comutators metric matrix can be expressd as
\begin{equation}\left (\begin{array}{cc}\left \{\mathbf{a} ,\mathbf{a}^{ \dag }\right \} & \left \{\mathbf{a} ,\mathbf{a}\right \} \\
\left \{\mathbf{a}^{ \dag } ,\mathbf{a}^{ \dag }\right \} & \left \{\mathbf{a}^{ \dag } ,\mathbf{a}\right \}\end{array}\right )\overset{\triangle }{ =}\left (\begin{array}{cc}\mathbf{a}^{\mathfrak{t}} \mathbf{a}^{ \dag } & \mathbf{a}^{\mathfrak{t}} \mathbf{a} \\
\mathbf{a}^{\mathfrak{a}} \mathbf{a}^{ \dag } & \mathbf{a}^{\mathfrak{a}} \mathbf{a}\end{array}\right ) +\left (\begin{array}{cc}\left (\mathbf{a}^{\mathfrak{a}} \mathbf{a}\right )^{\mathfrak{t}} & \left (\mathbf{a}^{\mathfrak{t}} \mathbf{a}\right )^{\mathfrak{t}} \\
\left (\mathbf{a}^{\mathfrak{a}} \mathbf{a}^{ \dag }\right )^{\mathfrak{t}} & \left (\mathbf{a}^{\mathfrak{t}} \mathbf{a}^{ \dag }\right )^{\mathfrak{t}}\end{array}\right )\text{.}
\end{equation}
\end{notation}

The CAR can be expanded as
\begin{gather}\left (\begin{array}{cc}\mathbf{a}^{\mathfrak{t}} \mathbf{a}^{ \dag } & \mathbf{a}^{\mathfrak{t}} \mathbf{a} \\
\mathbf{a}^{\mathfrak{a}} \mathbf{a}^{ \dag } & \mathbf{a}^{\mathfrak{a}} \mathbf{a}\end{array}\right ) +\left (\begin{array}{cc}\left (\mathbf{a}^{\mathfrak{a}} \mathbf{a}\right )^{\mathfrak{t}} & \left (\mathbf{a}^{\mathfrak{t}} \mathbf{a}\right )^{\mathfrak{t}} \\
\left (\mathbf{a}^{\mathfrak{a}} \mathbf{a}^{ \dag }\right )^{\mathfrak{t}} & \left (\mathbf{a}^{\mathfrak{t}} \mathbf{a}^{ \dag }\right )^{\mathfrak{t}}\end{array}\right ) \nonumber  \\
 =\left (\begin{array}{c}\mathbf{a}^{\mathfrak{t}} \\
\mathbf{a}^{a}\end{array}\right ) \left (\begin{array}{cc}\mathbf{a}^{ \dag } & \mathbf{a}\end{array}\right ) +\left (\begin{array}{cc}\mathbb{O}_{r} & \mathbb{I}_{r} \\
\mathbb{I}_{r} & \mathbb{O}_{r}\end{array}\right ) \left (\left (\begin{array}{c}\mathbf{a}^{\mathfrak{t}} \\
\mathbf{a}^{a}\end{array}\right ) \left (\begin{array}{cc}\mathbf{a}^{ \dag } & \mathbf{a}\end{array}\right )\right )^{\mathfrak{t}} \left (\begin{array}{cc}\mathbb{O}_{r} & \mathbb{I}_{r} \\
\mathbb{I}_{r} & \mathbb{O}_{r}\end{array}\right ) \nonumber  \\
 =\left (\begin{array}{cc}\mathbf{a}^{\mathfrak{t}} \mathbf{a}^{ \dag } & \mathbf{a}^{\mathfrak{t}} \mathbf{a} \\
\mathbf{a}^{\mathfrak{a}} \mathbf{a}^{ \dag } & \mathbf{a}^{\mathfrak{a}} \mathbf{a}\end{array}\right ) +\left (\begin{array}{cc}\mathbb{O}_{r} & \mathbb{I}_{r} \\
\mathbb{I}_{r} & \mathbb{O}_{r}\end{array}\right ) \left (\begin{array}{cc}\left (\mathbf{a}^{\mathfrak{a}} \mathbf{a}\right )^{\mathfrak{t}} & \left (\mathbf{a}^{\mathfrak{t}} \mathbf{a}\right )^{\mathfrak{t}} \\
\left (\mathbf{a}^{\mathfrak{a}} \mathbf{a}^{ \dag }\right )^{\mathfrak{t}} & \left (\mathbf{a}^{\mathfrak{t}} \mathbf{a}^{ \dag }\right )^{\mathfrak{t}}\end{array}\right ) \left (\begin{array}{cc}\mathbb{O}_{r} & \mathbb{I}_{r} \\
\mathbb{I}_{r} & \mathbb{O}_{r}\end{array}\right ) \nonumber  \\
 =\left (\begin{array}{cc}\mathbb{I}_{r} & \mathbb{O}_{r} \\
\mathbb{O}_{r} & \mathbb{I}_{r}\end{array}\right )\text{,}\end{gather}where
\begin{equation}\left (\begin{array}{cc}\mathbf{a}^{\mathfrak{t}} \mathbf{a}^{ \dag } & \mathbf{a}^{\mathfrak{t}} \mathbf{a} \\
\mathbf{a}^{\mathfrak{a}} \mathbf{a}^{ \dag } & \mathbf{a}^{\mathfrak{a}} \mathbf{a}\end{array}\right ) =\left (\begin{array}{cc}\mathbf{a}^{ \dag } & \mathbf{a}\end{array}\right )^{\mathfrak{a}} \left (\begin{array}{cc}\mathbf{a}^{ \dag } & \mathbf{a}\end{array}\right ) =\left (\begin{array}{c}\mathbf{a}^{\mathfrak{t}} \\
\mathbf{a}^{\mathfrak{a}}\end{array}\right ) \left (\begin{array}{cc}\mathbf{a}^{ \dag } & \mathbf{a}\end{array}\right )\text{.}
\end{equation}


\begin{notation}
\begin{example}
Let $r =1$
\begin{align}\mathbf{a}^{\mathfrak{t}} \mathbf{a}^{ \dag } &  =\left (\begin{array}{c}a_{1} \\
a_{2}\end{array}\right ) \left (\begin{array}{cc}a_{1}^{ \dag } & a_{2}^{ \dag }\end{array}\right ) =\left (\begin{array}{cc}a_{1} a_{1}^{ \dag } & a_{1} a_{2}^{ \dag } \\
a_{2} a_{1}^{ \dag } & a_{2} a_{2}^{ \dag }\end{array}\right ) \\
\left (\mathbf{a}^{\mathfrak{a}} \mathbf{a}\right )^{\mathfrak{t}} &  =\left (\left (\begin{array}{c}a_{1}^{ \dag } \\
a_{2}^{ \dag }\end{array}\right ) \left (\begin{array}{cc}a_{1} & a_{2}\end{array}\right )\right )^{\mathfrak{t}} =\left (\begin{array}{cc}a_{1}^{ \dag } a_{1} & a_{2}^{ \dag } a_{1} \\
a_{1}^{ \dag } a_{2} & a_{2}^{ \dag } a_{2}\end{array}\right ) \\
\mathbf{a}^{\mathfrak{a}} \mathbf{a} &  =\left (\begin{array}{c}a_{1}^{ \dag } \\
a_{2}^{ \dag }\end{array}\right ) \left (\begin{array}{cc}a_{1} & a_{2}\end{array}\right ) =\left (\begin{array}{cc}a_{1}^{ \dag } a_{1} & a_{2}^{ \dag } a_{1} \\
a_{1}^{ \dag } a_{2} & a_{2}^{ \dag } a_{2}\end{array}\right ) \\
\left (\mathbf{a}^{\mathfrak{t}} \mathbf{a}^{ \dag }\right )^{\mathfrak{t}} &  =\left (\left (\begin{array}{c}a_{1} \\
a_{2}\end{array}\right ) \left (\begin{array}{cc}a_{1}^{ \dag } & a_{2}^{ \dag }\end{array}\right )\right )^{\mathfrak{t}} =\left (\begin{array}{cc}a_{1} a_{1}^{ \dag } & a_{2} a_{1}^{ \dag } \\
a_{1} a_{2}^{ \dag } & a_{2} a_{2}^{ \dag }\end{array}\right )\end{align}
\begin{align}\mathbf{a}^{\mathfrak{t}} \mathbf{a} &  =\left (\begin{array}{c}a_{1} \\
a_{2}\end{array}\right ) \left (\begin{array}{cc}a_{1} & a_{2}\end{array}\right ) =\left (\begin{array}{cc}a_{1} a_{1} & a_{1} a_{2} \\
a_{2} a_{1} & a_{2} a_{2}\end{array}\right ) \\
\left (\mathbf{a}^{\mathfrak{t}} \mathbf{a}\right )^{\mathfrak{t}} &  =\left (\left (\begin{array}{c}a_{1} \\
a_{2}\end{array}\right ) \left (\begin{array}{cc}a_{1} & a_{2}\end{array}\right )\right )^{\mathfrak{t}} =\left (\begin{array}{cc}a_{1} a_{1} & a_{2} a_{1} \\
a_{1} a_{2} & a_{2} a_{2}\end{array}\right ) \\
\mathbf{a}^{\mathfrak{a}} \mathbf{a}^{ \dag } &  =\left (\begin{array}{c}a_{1}^{ \dag } \\
a_{2}^{ \dag }\end{array}\right ) \left (\begin{array}{cc}a_{1}^{ \dag } & a_{2}^{ \dag }\end{array}\right ) =\left (\begin{array}{cc}a_{1}^{ \dag } a_{1}^{ \dag } & a_{2}^{ \dag } a_{1}^{ \dag } \\
a_{1}^{ \dag } a_{2}^{ \dag } & a_{2}^{ \dag } a_{2}^{ \dag }\end{array}\right ) \\
\left (\mathbf{a}^{\mathfrak{a}} \mathbf{a}^{ \dag }\right )^{\mathfrak{t}} &  =\left (\left (\begin{array}{c}a_{1}^{ \dag } \\
a_{2}^{ \dag }\end{array}\right ) \left (\begin{array}{cc}a_{1}^{ \dag } & a_{2}^{ \dag }\end{array}\right )\right )^{\mathfrak{t}} =\left (\begin{array}{cc}a_{1}^{ \dag } a_{1}^{ \dag } & a_{1}^{ \dag } a_{2}^{ \dag } \\
a_{2}^{ \dag } a_{1}^{ \dag } & a_{2}^{ \dag } a_{2}^{ \dag }\end{array}\right )\end{align}
\end{example}


\end{notation}


\begin{example}
For $r =2$ the CAR is
\begin{equation}\left (\begin{array}{cccc}a_{1} a_{1}^{ \dag } & a_{1} a_{2}^{ \dag } & a_{1} a_{1} & a_{1} a_{2} \\
a_{2} a_{1}^{ \dag } & a_{2} a_{2}^{ \dag } & a_{2} a_{1} & a_{2} a_{2} \\
a_{1}^{ \dag } a_{1}^{ \dag } & a_{1}^{ \dag } a_{2}^{ \dag } & a_{1}^{ \dag } a_{1} & a_{1}^{ \dag } a_{2} \\
a_{2}^{ \dag } a_{1}^{ \dag } & a_{2}^{ \dag } a_{2}^{ \dag } & a_{2}^{ \dag } a_{1} & a_{2}^{ \dag } a_{2}\end{array}\right ) +\left (\begin{array}{cccc}a_{1}^{ \dag } a_{1} & a_{2}^{ \dag } a_{1} & a_{1} a_{1} & a_{2} a_{1} \\
a_{1}^{ \dag } a_{2} & a_{2}^{ \dag } a_{2} & a_{1} a_{2} & a_{2} a_{2} \\
a_{1}^{ \dag } a_{1}^{ \dag } & a_{2}^{ \dag } a_{1}^{ \dag } & a_{1} a_{1}^{ \dag } & a_{2} a_{1}^{ \dag } \\
a_{1}^{ \dag } a_{2}^{ \dag } & a_{2}^{ \dag } a_{2}^{ \dag } & a_{1} a_{2}^{ \dag } & a_{2} a_{2}^{ \dag }\end{array}\right ) =\left (\begin{array}{cccc}1 & 0 & 0 & 0 \\
0 & 1 & 0 & 0 \\
0 & 0 & 1 & 0 \\
0 & 0 & 0 & 1\end{array}\right )\text{,}
\end{equation}giving
\begin{gather}\left (\begin{array}{cc}\left (\begin{array}{c}a_{1} \\
a_{2}\end{array}\right ) \left (\begin{array}{cc}a_{1}^{ \dag } & a_{2}^{ \dag }\end{array}\right ) & \left (\begin{array}{c}a_{1} \\
a_{2}\end{array}\right ) \left (\begin{array}{cc}a_{1} & a_{2}\end{array}\right ) \\
\left (\begin{array}{c}a_{1}^{ \dag } \\
a_{2}^{ \dag }\end{array}\right ) \left (\begin{array}{cc}a_{1}^{ \dag } & a_{2}^{ \dag }\end{array}\right ) & \left (\begin{array}{c}a_{1}^{ \dag } \\
a_{2}^{ \dag }\end{array}\right ) \left (\begin{array}{cc}a_{1} & a_{2}\end{array}\right )\end{array}\right ) +\left (\begin{array}{cc}\left (\left (\begin{array}{c}a_{1}^{ \dag } \\
a_{2}^{ \dag }\end{array}\right ) \left (\begin{array}{cc}a_{1} & a_{2}\end{array}\right )\right )^{\mathfrak{t}} & \left (\left (\begin{array}{c}a_{1}^{ \dag } \\
a_{2}^{ \dag }\end{array}\right ) \left (\begin{array}{cc}a_{1}^{ \dag } & a_{2}^{ \dag }\end{array}\right )\right )^{\mathfrak{a}} \\
\left (\left (\begin{array}{c}a_{1} \\
a_{2}\end{array}\right ) \left (\begin{array}{cc}a_{1} & a_{2}\end{array}\right )\right )^{\mathfrak{a}} & \left (\left (\begin{array}{c}a_{1} \\
a_{2}\end{array}\right ) \left (\begin{array}{cc}a_{1}^{ \dag } & a_{2}^{ \dag }\end{array}\right )\right )^{\mathfrak{t}}\end{array}\right ) \\
 =\left (\begin{array}{cc}\left (\begin{array}{c}a_{1} \\
a_{2}\end{array}\right ) \left (\begin{array}{cc}a_{1}^{ \dag } & a_{2}^{ \dag }\end{array}\right ) & \left (\begin{array}{c}a_{1} \\
a_{2}\end{array}\right ) \left (\begin{array}{cc}a_{1} & a_{2}\end{array}\right ) \\
\left (\begin{array}{c}a_{1}^{ \dag } \\
a_{2}^{ \dag }\end{array}\right ) \left (\begin{array}{cc}a_{1}^{ \dag } & a_{2}^{ \dag }\end{array}\right ) & \left (\begin{array}{c}a_{1}^{ \dag } \\
a_{2}^{ \dag }\end{array}\right ) \left (\begin{array}{cc}a_{1} & a_{2}\end{array}\right )\end{array}\right ) +\left (\begin{array}{cc}\left (\begin{array}{c}a_{1} \\
a_{2}\end{array}\right ) \left (\begin{array}{cc}a_{1}^{ \dag } & a_{2}^{ \dag }\end{array}\right ) & \left (\begin{array}{c}a_{1} \\
a_{2}\end{array}\right ) \left (\begin{array}{cc}a_{1} & a_{2}\end{array}\right ) \\
\left (\begin{array}{c}a_{1}^{ \dag } \\
a_{2}^{ \dag }\end{array}\right ) \left (\begin{array}{cc}a_{1}^{ \dag } & a_{2}^{ \dag }\end{array}\right ) & \left (\begin{array}{c}a_{1}^{ \dag } \\
a_{2}^{ \dag }\end{array}\right ) \left (\begin{array}{cc}a_{1} & a_{2}\end{array}\right )\end{array}\right )^{\mathfrak{t}} \\
 =\left (\begin{array}{cccc}1 & 0 & 0 & 0 \\
0 & 1 & 0 & 0 \\
0 & 0 & 1 & 0 \\
0 & 0 & 0 & 1\end{array}\right )\text{,}\end{gather}where
\begin{align}\left (\begin{array}{cccc}a_{1} a_{1}^{ \dag } & a_{1} a_{2}^{ \dag } & a_{1} a_{1} & a_{1} a_{2} \\
a_{2} a_{1}^{ \dag } & a_{2} a_{2}^{ \dag } & a_{2} a_{1} & a_{2} a_{2} \\
a_{1}^{ \dag } a_{1}^{ \dag } & a_{1}^{ \dag } a_{2}^{ \dag } & a_{1}^{ \dag } a_{1} & a_{1}^{ \dag } a_{2} \\
a_{2}^{ \dag } a_{1}^{ \dag } & a_{2}^{ \dag } a_{2}^{ \dag } & a_{2}^{ \dag } a_{1} & a_{2}^{ \dag } a_{2}\end{array}\right ) &  =\left (\begin{array}{c}a_{1} \\
a_{2} \\
a_{1}^{ \dag } \\
a_{2}^{ \dag }\end{array}\right ) \left (\begin{array}{cccc}a_{1}^{ \dag } & a_{2}^{ \dag } & a_{1} & a_{2}\end{array}\right ) \\
 &  =\left (\begin{array}{cccc}a_{1}^{ \dag } & a_{2}^{ \dag } & a_{1} & a_{2}\end{array}\right )^{\mathfrak{a}} \left (\begin{array}{cccc}a_{1}^{ \dag } & a_{2}^{ \dag } & a_{1} & a_{2}\end{array}\right )\end{align}$\lceil $
\begin{equation}\left (\begin{array}{c}a_{1}^{ \dag } \\
a_{2}^{ \dag } \\
a_{1} \\
a_{2}\end{array}\right ) \left (\begin{array}{cccc}a_{1} & a_{2} & a_{1}^{ \dag } & a_{2}^{ \dag }\end{array}\right ) =\left (\begin{array}{cccc}a_{1}^{ \dag } a_{1} & a_{1}^{ \dag } a_{2} & a_{1}^{ \dag } a_{1}^{ \dag } & a_{1}^{ \dag } a_{2}^{ \dag } \\
a_{2}^{ \dag } a_{1} & a_{2}^{ \dag } a_{2} & a_{2}^{ \dag } a_{1}^{ \dag } & a_{2}^{ \dag } a_{2}^{ \dag } \\
a_{1} a_{1} & a_{1} a_{2} & a_{1} a_{1}^{ \dag } & a_{1} a_{2}^{ \dag } \\
a_{2} a_{1} & a_{2} a_{2} & a_{2} a_{1}^{ \dag } & a_{2} a_{2}^{ \dag }\end{array}\right )
\end{equation}$\rfloor $ and
\begin{align} & \left (\begin{array}{cccc}a_{1}^{ \dag } a_{1} & a_{2}^{ \dag } a_{1} & a_{1} a_{1} & a_{2} a_{1} \\
a_{1}^{ \dag } a_{2} & a_{2}^{ \dag } a_{2} & a_{1} a_{2} & a_{2} a_{2} \\
a_{1}^{ \dag } a_{1}^{ \dag } & a_{2}^{ \dag } a_{1}^{ \dag } & a_{1} a_{1}^{ \dag } & a_{2} a_{1}^{ \dag } \\
a_{1}^{ \dag } a_{2}^{ \dag } & a_{2}^{ \dag } a_{2}^{ \dag } & a_{1} a_{2}^{ \dag } & a_{2} a_{2}^{ \dag }\end{array}\right ) =\left (\left (\begin{array}{c}a_{1}^{ \dag } \\
a_{2}^{ \dag } \\
a_{1} \\
a_{2}\end{array}\right ) \left (\begin{array}{cccc}a_{1} & a_{2} & a_{1}^{ \dag } & a_{2}^{ \dag }\end{array}\right )\right )^{\mathfrak{t}} =\left (\left (\begin{array}{cccc}a_{1} & a_{2} & a_{1}^{ \dag } & a_{2}^{ \dag }\end{array}\right )^{\mathfrak{a}} \left (\begin{array}{cccc}a_{1} & a_{2} & a_{1}^{ \dag } & a_{2}^{ \dag }\end{array}\right )\right )^{\mathfrak{t}} \nonumber  \\
 &  =\left (\left (\begin{array}{cc}\mathbb{O}_{2} & \mathbb{I}_{2} \\
\mathbb{I}_{2} & \mathbb{O}_{2}\end{array}\right ) \left (\begin{array}{c}a_{1} \\
a_{2} \\
a_{1}^{ \dag } \\
a_{2}^{ \dag }\end{array}\right ) \left (\begin{array}{cccc}a_{1}^{ \dag } & a_{2}^{ \dag } & a_{1} & a_{2}\end{array}\right ) \left (\begin{array}{cc}\mathbb{O}_{2} & \mathbb{I}_{2} \\
\mathbb{I}_{2} & \mathbb{O}_{2}\end{array}\right )\right )^{\mathfrak{t}} \\
 &  =\left (\left (\begin{array}{cc}\mathbb{O}_{2} & \mathbb{I}_{2} \\
\mathbb{I}_{2} & \mathbb{O}_{2}\end{array}\right ) \left (\begin{array}{cccc}a_{1} a_{1}^{ \dag } & a_{1} a_{2}^{ \dag } & a_{1} a_{1} & a_{1} a_{2} \\
a_{2} a_{1}^{ \dag } & a_{2} a_{2}^{ \dag } & a_{2} a_{1} & a_{2} a_{2} \\
a_{1}^{ \dag } a_{1}^{ \dag } & a_{1}^{ \dag } a_{2}^{ \dag } & a_{1}^{ \dag } a_{1} & a_{1}^{ \dag } a_{2} \\
a_{2}^{ \dag } a_{1}^{ \dag } & a_{2}^{ \dag } a_{2}^{ \dag } & a_{2}^{ \dag } a_{1} & a_{2}^{ \dag } a_{2}\end{array}\right ) \left (\begin{array}{cc}\mathbb{O}_{2} & \mathbb{I}_{2} \\
\mathbb{I}_{2} & \mathbb{O}_{2}\end{array}\right )\right )^{\mathfrak{t}} \nonumber  \\
 &  =\left (\left (\begin{array}{cc}\mathbb{O}_{2} & \mathbb{I}_{2} \\
\mathbb{I}_{2} & \mathbb{O}_{2}\end{array}\right ) \left (\begin{array}{cccc}a_{1} a_{1} & a_{1} a_{2} & a_{1} a_{1}^{ \dag } & a_{1} a_{2}^{ \dag } \\
a_{2} a_{1} & a_{2} a_{2} & a_{2} a_{1}^{ \dag } & a_{2} a_{2}^{ \dag } \\
a_{1}^{ \dag } a_{1} & a_{1}^{ \dag } a_{2} & a_{1}^{ \dag } a_{1}^{ \dag } & a_{1}^{ \dag } a_{2}^{ \dag } \\
a_{2}^{ \dag } a_{1} & a_{2}^{ \dag } a_{2} & a_{2}^{ \dag } a_{1}^{ \dag } & a_{2}^{ \dag } a_{2}^{ \dag }\end{array}\right )\right )^{\mathfrak{t}} \nonumber  \\
 &  =\left (\begin{array}{cccc}a_{1}^{ \dag } a_{1} & a_{1}^{ \dag } a_{2} & a_{1}^{ \dag } a_{1}^{ \dag } & a_{1}^{ \dag } a_{2}^{ \dag } \\
a_{2}^{ \dag } a_{1} & a_{2}^{ \dag } a_{2} & a_{2}^{ \dag } a_{1}^{ \dag } & a_{2}^{ \dag } a_{2}^{ \dag } \\
a_{1} a_{1} & a_{1} a_{2} & a_{1} a_{1}^{ \dag } & a_{1} a_{2}^{ \dag } \\
a_{2} a_{1} & a_{2} a_{2} & a_{2} a_{1}^{ \dag } & a_{2} a_{2}^{ \dag }\end{array}\right )^{\mathfrak{t}} =\left (\begin{array}{cccc}a_{1}^{ \dag } a_{1} & a_{2}^{ \dag } a_{1} & a_{1} a_{1} & a_{2} a_{1} \\
a_{1}^{ \dag } a_{2} & a_{2}^{ \dag } a_{2} & a_{1} a_{2} & a_{2} a_{2} \\
a_{1}^{ \dag } a_{1}^{ \dag } & a_{2}^{ \dag } a_{1}^{ \dag } & a_{1} a_{1}^{ \dag } & a_{2} a_{1}^{ \dag } \\
a_{1}^{ \dag } a_{2}^{ \dag } & a_{2}^{ \dag } a_{2}^{ \dag } & a_{1} a_{2}^{ \dag } & a_{2} a_{2}^{ \dag }\end{array}\right )\text{.}\end{align}This leads to the CAR being expressed as
\begin{equation}\left (\begin{array}{c}a_{1} \\
a_{2} \\
a_{1}^{ \dag } \\
a_{2}^{ \dag }\end{array}\right ) \left (\begin{array}{cccc}a_{1}^{ \dag } & a_{2}^{ \dag } & a_{1} & a_{2}\end{array}\right ) +\left (\left (\begin{array}{cc}\mathbb{O}_{2} & \mathbb{I}_{2} \\
\mathbb{I}_{2} & \mathbb{O}_{2}\end{array}\right ) \left (\begin{array}{c}a_{1} \\
a_{2} \\
a_{1}^{ \dag } \\
a_{2}^{ \dag }\end{array}\right ) \left (\begin{array}{cccc}a_{1}^{ \dag } & a_{2}^{ \dag } & a_{1} & a_{2}\end{array}\right ) \left (\begin{array}{cc}\mathbb{O}_{2} & \mathbb{I}_{2} \\
\mathbb{I}_{2} & \mathbb{O}_{2}\end{array}\right )\right )^{\mathfrak{t}} =\left (\begin{array}{cccc}1 & 0 & 0 & 0 \\
0 & 1 & 0 & 0 \\
0 & 0 & 1 & 0 \\
0 & 0 & 0 & 1\end{array}\right )
\end{equation}$\text{.}$ 
\end{example}

\subsubsection{Canonical Transfromations}
\begin{example}
For a canonical transformation $T$
\begin{align} & \left (\begin{array}{cccc}\bar{u}_{1 1} & \bar{u}_{2 1} & \bar{v}_{1 1} & \bar{v}_{2 1} \\
\bar{u}_{1 2} & \bar{u}_{2 2} & \bar{v}_{1 2} & \bar{v}_{2 2} \\
v_{1 1} & v_{2 1} & u_{1 1} & u_{2 1} \\
v_{1 2} & v_{2 2} & u_{1 2} & u_{2 2}\end{array}\right ) \left (\begin{array}{c}a_{1} \\
a_{2} \\
a_{1}^{ \dag } \\
a_{2}^{ \dag }\end{array}\right ) \left (\begin{array}{cccc}a_{1}^{ \dag } & a_{2}^{ \dag } & a_{1} & a_{2}\end{array}\right ) \left (\begin{array}{cccc}u_{1 1} & u_{1 2} & \bar{v}_{1 1} & \bar{v}_{1 2} \\
u_{2 1} & u_{2 2} & \bar{v}_{2 1} & \bar{v}_{2 2} \\
v_{1 1} & v_{1 2} & \bar{u}_{1 1} & \bar{u}_{1 2} \\
v_{2 1} & v_{2 2} & \bar{u}_{2 1} & \bar{u}_{2 2}\end{array}\right ) +\left (\begin{array}{cc}\mathbb{O}_{2} & \mathbb{I}_{2} \\
\mathbb{I}_{2} & \mathbb{O}_{2}\end{array}\right ) \times  \nonumber  \\
 & \left (\left (\begin{array}{cccc}\bar{u}_{1 1} & \bar{u}_{2 1} & \bar{v}_{1 1} & \bar{v}_{2 1} \\
\bar{u}_{1 2} & \bar{u}_{2 2} & \bar{v}_{1 2} & \bar{v}_{2 2} \\
v_{1 1} & v_{2 1} & u_{1 1} & u_{2 1} \\
v_{1 2} & v_{2 2} & u_{1 2} & u_{2 2}\end{array}\right ) \left (\begin{array}{c}a_{1} \\
a_{2} \\
a_{1}^{ \dag } \\
a_{2}^{ \dag }\end{array}\right ) \left (\begin{array}{cccc}a_{1}^{ \dag } & a_{2}^{ \dag } & a_{1} & a_{2}\end{array}\right ) \left (\begin{array}{cccc}u_{1 1} & u_{1 2} & \bar{v}_{1 1} & \bar{v}_{1 2} \\
u_{2 1} & u_{2 2} & \bar{v}_{2 1} & \bar{v}_{2 2} \\
v_{1 1} & v_{1 2} & \bar{u}_{1 1} & \bar{u}_{1 2} \\
v_{2 1} & v_{2 2} & \bar{u}_{2 1} & \bar{u}_{2 2}\end{array}\right )\right )^{\mathfrak{t}} \left (\begin{array}{cc}\mathbb{O}_{2} & \mathbb{I}_{2} \\
\mathbb{I}_{2} & \mathbb{O}_{2}\end{array}\right ) \nonumber  \\
 &  =\left (\begin{array}{cccc}1 & 0 & 0 & 0 \\
0 & 1 & 0 & 0 \\
0 & 0 & 1 & 0 \\
0 & 0 & 0 & 1\end{array}\right )\text{.}\end{align}This gives
\begin{align} & \left (\begin{array}{cccc}\bar{u}_{1 1} & \bar{u}_{2 1} & \bar{v}_{1 1} & \bar{v}_{2 1} \\
\bar{u}_{1 2} & \bar{u}_{2 2} & \bar{v}_{1 2} & \bar{v}_{2 2} \\
v_{1 1} & v_{2 1} & u_{1 1} & u_{2 1} \\
v_{1 2} & v_{2 2} & u_{1 2} & u_{2 2}\end{array}\right ) \left (\begin{array}{c}a_{1} \\
a_{2} \\
a_{1}^{ \dag } \\
a_{2}^{ \dag }\end{array}\right ) \left (\begin{array}{cccc}a_{1}^{ \dag } & a_{2}^{ \dag } & a_{1} & a_{2}\end{array}\right ) \left (\begin{array}{cccc}u_{1 1} & u_{1 2} & \bar{v}_{1 1} & \bar{v}_{1 2} \\
u_{2 1} & u_{2 2} & \bar{v}_{2 1} & \bar{v}_{2 2} \\
v_{1 1} & v_{1 2} & \bar{u}_{1 1} & \bar{u}_{1 2} \\
v_{2 1} & v_{2 2} & \bar{u}_{2 1} & \bar{u}_{2 2}\end{array}\right ) +\left (\begin{array}{cc}\mathbb{O}_{2} & \mathbb{I}_{2} \\
\mathbb{I}_{2} & \mathbb{O}_{2}\end{array}\right ) \times  \nonumber  \\
 & \left (\begin{array}{cccc}u_{1 1} & u_{1 2} & \bar{v}_{1 1} & \bar{v}_{1 2} \\
u_{2 1} & u_{2 2} & \bar{v}_{2 1} & \bar{v}_{2 2} \\
v_{1 1} & v_{1 2} & \bar{u}_{1 1} & \bar{u}_{1 2} \\
v_{2 1} & v_{2 2} & \bar{u}_{2 1} & \bar{u}_{2 2}\end{array}\right )^{t} \left (\left (\begin{array}{c}a_{1} \\
a_{2} \\
a_{1}^{ \dag } \\
a_{2}^{ \dag }\end{array}\right ) \left (\begin{array}{cccc}a_{1}^{ \dag } & a_{2}^{ \dag } & a_{1} & a_{2}\end{array}\right )\right )^{\mathfrak{t}} \left (\begin{array}{cccc}\bar{u}_{1 1} & \bar{u}_{2 1} & \bar{v}_{1 1} & \bar{v}_{2 1} \\
\bar{u}_{1 2} & \bar{u}_{2 2} & \bar{v}_{1 2} & \bar{v}_{2 2} \\
v_{1 1} & v_{2 1} & u_{1 1} & u_{2 1} \\
v_{1 2} & v_{2 2} & u_{1 2} & u_{2 2}\end{array}\right )^{t} \left (\begin{array}{cc}\mathbb{O}_{2} & \mathbb{I}_{2} \\
\mathbb{I}_{2} & \mathbb{O}_{2}\end{array}\right ) \nonumber  \\
 &  =\left (\begin{array}{cccc}1 & 0 & 0 & 0 \\
0 & 1 & 0 & 0 \\
0 & 0 & 1 & 0 \\
0 & 0 & 0 & 1\end{array}\right )\end{align}and
\begin{align} & \left (\begin{array}{cccc}\bar{u}_{1 1} & \bar{u}_{2 1} & \bar{v}_{1 1} & \bar{v}_{2 1} \\
\bar{u}_{1 2} & \bar{u}_{2 2} & \bar{v}_{1 2} & \bar{v}_{2 2} \\
v_{1 1} & v_{2 1} & u_{1 1} & u_{2 1} \\
v_{1 2} & v_{2 2} & u_{1 2} & u_{2 2}\end{array}\right ) \left (\begin{array}{c}a_{1} \\
a_{2} \\
a_{1}^{ \dag } \\
a_{2}^{ \dag }\end{array}\right ) \left (\begin{array}{cccc}a_{1}^{ \dag } & a_{2}^{ \dag } & a_{1} & a_{2}\end{array}\right ) \left (\begin{array}{cccc}u_{1 1} & u_{1 2} & \bar{v}_{1 1} & \bar{v}_{1 2} \\
u_{2 1} & u_{2 2} & \bar{v}_{2 1} & \bar{v}_{2 2} \\
v_{1 1} & v_{1 2} & \bar{u}_{1 1} & \bar{u}_{1 2} \\
v_{2 1} & v_{2 2} & \bar{u}_{2 1} & \bar{u}_{2 2}\end{array}\right ) +\left (\begin{array}{cc}\mathbb{O}_{2} & \mathbb{I}_{2} \\
\mathbb{I}_{2} & \mathbb{O}_{2}\end{array}\right ) \times  \nonumber  \\
 & \left (\begin{array}{cccc}u_{1 1} & u_{2 1} & v_{1 1} & v_{1 2} \\
u_{1 2} & u_{2 2} & v_{2 1} & v_{2 2} \\
\bar{v}_{1 1} & \bar{v}_{2 1} & \bar{u}_{1 1} & \bar{u}_{2 1} \\
\bar{v}_{1 2} & \bar{v}_{2 2} & \bar{u}_{1 2} & \bar{u}_{2 2}\end{array}\right ) \left (\left (\begin{array}{c}a_{1} \\
a_{2} \\
a_{1}^{ \dag } \\
a_{2}^{ \dag }\end{array}\right ) \left (\begin{array}{cccc}a_{1}^{ \dag } & a_{2}^{ \dag } & a_{1} & a_{2}\end{array}\right )\right )^{\mathfrak{t}} \left (\begin{array}{cccc}\bar{u}_{1 1} & \bar{u}_{2 1} & v_{1 1} & v_{1 2} \\
\bar{u}_{1 2} & \bar{u}_{2 2} & v_{2 1} & v_{2 2} \\
\bar{v}_{1 1} & \bar{v}_{2 1} & u_{1 1} & u_{2 1} \\
\bar{v}_{1 2} & \bar{v}_{2 2} & u_{1 2} & u_{2 2}\end{array}\right ) \left (\begin{array}{cc}\mathbb{O}_{2} & \mathbb{I}_{2} \\
\mathbb{I}_{2} & \mathbb{O}_{2}\end{array}\right ) \nonumber  \\
 &  =\left (\begin{array}{cccc}1 & 0 & 0 & 0 \\
0 & 1 & 0 & 0 \\
0 & 0 & 1 & 0 \\
0 & 0 & 0 & 1\end{array}\right )\text{.}\end{align}Now
\begin{align} & \left (\begin{array}{cccc}\bar{u}_{1 1} & \bar{u}_{2 1} & v_{1 1} & v_{1 2} \\
\bar{u}_{1 2} & \bar{u}_{2 2} & v_{2 1} & v_{2 2} \\
\bar{v}_{1 1} & \bar{v}_{2 1} & u_{1 1} & u_{2 1} \\
\bar{v}_{1 2} & \bar{v}_{2 2} & u_{1 2} & u_{2 2}\end{array}\right ) \left (\begin{array}{cc}\mathbb{O}_{2} & \mathbb{I}_{2} \\
\mathbb{I}_{2} & \mathbb{O}_{2}\end{array}\right ) \nonumber  \\
 &  =\left (\begin{array}{cc}\mathbb{O}_{2} & \mathbb{I}_{2} \\
\mathbb{I}_{2} & \mathbb{O}_{2}\end{array}\right ) \left (\begin{array}{cc}\mathbb{O}_{2} & \mathbb{I}_{2} \\
\mathbb{I}_{2} & \mathbb{O}_{2}\end{array}\right ) \left (\begin{array}{cccc}\bar{u}_{1 1} & \bar{u}_{2 1} & v_{1 1} & v_{1 2} \\
\bar{u}_{1 2} & \bar{u}_{2 2} & v_{2 1} & v_{2 2} \\
\bar{v}_{1 1} & \bar{v}_{2 1} & u_{1 1} & u_{2 1} \\
\bar{v}_{1 2} & \bar{v}_{2 2} & u_{1 2} & u_{2 2}\end{array}\right ) \left (\begin{array}{cc}\mathbb{O}_{2} & \mathbb{I}_{2} \\
\mathbb{I}_{2} & \mathbb{O}_{2}\end{array}\right ) \\
 &  =\left (\begin{array}{cccc}\bar{v}_{1 1} & \bar{v}_{2 1} & u_{1 1} & u_{2 1} \\
\bar{v}_{1 2} & \bar{v}_{2 2} & u_{1 2} & u_{2 2} \\
\bar{u}_{1 1} & \bar{u}_{2 1} & v_{1 1} & v_{1 2} \\
\bar{u}_{1 2} & \bar{u}_{2 2} & v_{2 1} & v_{2 2}\end{array}\right ) \left (\begin{array}{cc}\mathbb{O}_{2} & \mathbb{I}_{2} \\
\mathbb{I}_{2} & \mathbb{O}_{2}\end{array}\right ) \left (\begin{array}{cc}\mathbb{O}_{2} & \mathbb{I}_{2} \\
\mathbb{I}_{2} & \mathbb{O}_{2}\end{array}\right ) =\left (\begin{array}{cc}\mathbb{O}_{2} & \mathbb{I}_{2} \\
\mathbb{I}_{2} & \mathbb{O}_{2}\end{array}\right ) \left (\begin{array}{cccc}u_{1 1} & u_{2 1} & \bar{v}_{1 1} & \bar{v}_{2 1} \\
u_{1 2} & u_{2 2} & \bar{v}_{1 2} & \bar{v}_{2 2} \\
v_{1 1} & v_{1 2} & \bar{u}_{1 1} & \bar{u}_{2 1} \\
v_{2 1} & v_{2 2} & \bar{u}_{1 2} & \bar{u}_{2 2}\end{array}\right )\end{align}i.e.
\begin{equation}\frame{$\left (\begin{array}{cccc}\bar{u}_{1 1} & \bar{u}_{2 1} & v_{1 1} & v_{1 2} \\
\bar{u}_{1 2} & \bar{u}_{2 2} & v_{2 1} & v_{2 2} \\
\bar{v}_{1 1} & \bar{v}_{2 1} & u_{1 1} & u_{2 1} \\
\bar{v}_{1 2} & \bar{v}_{2 2} & u_{1 2} & u_{2 2}\end{array}\right ) \left (\begin{array}{cc}\mathbb{O}_{2} & \mathbb{I}_{2} \\
\mathbb{I}_{2} & \mathbb{O}_{2}\end{array}\right ) =\left (\begin{array}{cc}\mathbb{O}_{2} & \mathbb{I}_{2} \\
\mathbb{I}_{2} & \mathbb{O}_{2}\end{array}\right ) \left (\begin{array}{cccc}u_{1 1} & u_{2 1} & \bar{v}_{1 1} & \bar{v}_{2 1} \\
u_{1 2} & u_{2 2} & \bar{v}_{1 2} & \bar{v}_{2 2} \\
v_{1 1} & v_{1 2} & \bar{u}_{1 1} & \bar{u}_{2 1} \\
v_{2 1} & v_{2 2} & \bar{u}_{1 2} & \bar{u}_{2 2}\end{array}\right )$}
\end{equation}and
\begin{align} & \left (\begin{array}{cc}\mathbb{O}_{2} & \mathbb{I}_{2} \\
\mathbb{I}_{2} & \mathbb{O}_{2}\end{array}\right ) \left (\begin{array}{cccc}u_{1 1} & u_{2 1} & v_{1 1} & v_{1 2} \\
u_{1 2} & u_{2 2} & v_{2 1} & v_{2 2} \\
\bar{v}_{1 1} & \bar{v}_{2 1} & \bar{u}_{1 1} & \bar{u}_{2 1} \\
\bar{v}_{1 2} & \bar{v}_{2 2} & \bar{u}_{1 2} & \bar{u}_{2 2}\end{array}\right ) \\
 &  =\left (\begin{array}{cc}\mathbb{O}_{2} & \mathbb{I}_{2} \\
\mathbb{I}_{2} & \mathbb{O}_{2}\end{array}\right ) \left (\begin{array}{cccc}u_{1 1} & u_{2 1} & v_{1 1} & v_{1 2} \\
u_{1 2} & u_{2 2} & v_{2 1} & v_{2 2} \\
\bar{v}_{1 1} & \bar{v}_{2 1} & \bar{u}_{1 1} & \bar{u}_{2 1} \\
\bar{v}_{1 2} & \bar{v}_{2 2} & \bar{u}_{1 2} & \bar{u}_{2 2}\end{array}\right ) \left (\begin{array}{cc}\mathbb{O}_{2} & \mathbb{I}_{2} \\
\mathbb{I}_{2} & \mathbb{O}_{2}\end{array}\right ) \left (\begin{array}{cc}\mathbb{O}_{2} & \mathbb{I}_{2} \\
\mathbb{I}_{2} & \mathbb{O}_{2}\end{array}\right ) \\
 &  =\left (\begin{array}{cccc}\bar{v}_{1 1} & \bar{v}_{2 1} & \bar{u}_{1 1} & \bar{u}_{2 1} \\
\bar{v}_{1 2} & \bar{v}_{2 2} & \bar{u}_{1 2} & \bar{u}_{2 2} \\
u_{1 1} & u_{2 1} & v_{1 1} & v_{1 2} \\
u_{1 2} & u_{2 2} & v_{2 1} & v_{2 2}\end{array}\right ) \left (\begin{array}{cc}\mathbb{O}_{2} & \mathbb{I}_{2} \\
\mathbb{I}_{2} & \mathbb{O}_{2}\end{array}\right ) \left (\begin{array}{cc}\mathbb{O}_{2} & \mathbb{I}_{2} \\
\mathbb{I}_{2} & \mathbb{O}_{2}\end{array}\right ) \\
 &  =\left (\begin{array}{cccc}\bar{u}_{1 1} & \bar{u}_{2 1} & \bar{v}_{1 1} & \bar{v}_{2 1} \\
\bar{u}_{1 2} & \bar{u}_{2 2} & \bar{v}_{1 2} & \bar{v}_{2 2} \\
v_{1 1} & v_{2 1} & u_{1 1} & u_{2 1} \\
v_{1 2} & v_{2 2} & u_{1 2} & u_{2 2}\end{array}\right ) \left (\begin{array}{cc}\mathbb{O}_{2} & \mathbb{I}_{2} \\
\mathbb{I}_{2} & \mathbb{O}_{2}\end{array}\right )\end{align}i.e.
\begin{equation}\frame{$\left (\begin{array}{cc}\mathbb{O}_{2} & \mathbb{I}_{2} \\
\mathbb{I}_{2} & \mathbb{O}_{2}\end{array}\right ) \left (\begin{array}{cccc}u_{1 1} & u_{2 1} & v_{1 1} & v_{2 1} \\
u_{1 2} & u_{2 2} & v_{1 2} & v_{2 2} \\
\bar{v}_{1 1} & \bar{v}_{2 1} & \bar{u}_{1 1} & \bar{u}_{2 1} \\
\bar{v}_{1 2} & \bar{v}_{2 2} & \bar{u}_{1 2} & \bar{u}_{2 2}\end{array}\right ) =\left (\begin{array}{cccc}\bar{u}_{1 1} & \bar{u}_{2 1} & \bar{v}_{1 1} & \bar{v}_{2 1} \\
\bar{u}_{1 2} & \bar{u}_{2 2} & \bar{v}_{1 2} & \bar{v}_{2 2} \\
v_{1 1} & v_{2 1} & u_{1 1} & u_{2 1} \\
v_{1 2} & v_{2 2} & u_{1 2} & u_{2 2}\end{array}\right ) \left (\begin{array}{cc}\mathbb{O}_{2} & \mathbb{I}_{2} \\
\mathbb{I}_{2} & \mathbb{O}_{2}\end{array}\right )\text{.}$}
\end{equation}Hence
\begin{align} & \left (\begin{array}{cccc}\bar{u}_{1 1} & \bar{u}_{2 1} & \bar{v}_{1 1} & \bar{v}_{2 1} \\
\bar{u}_{1 2} & \bar{u}_{2 2} & \bar{v}_{1 2} & \bar{v}_{2 2} \\
v_{1 1} & v_{2 1} & u_{1 1} & u_{2 1} \\
v_{1 2} & v_{2 2} & u_{1 2} & u_{2 2}\end{array}\right ) \left (\begin{array}{c}a_{1} \\
a_{2} \\
a_{1}^{ \dag } \\
a_{2}^{ \dag }\end{array}\right ) \left (\begin{array}{cccc}a_{1}^{ \dag } & a_{2}^{ \dag } & a_{1} & a_{2}\end{array}\right ) \left (\begin{array}{cccc}u_{1 1} & u_{1 2} & \bar{v}_{1 1} & \bar{v}_{1 2} \\
u_{2 1} & u_{2 2} & \bar{v}_{2 1} & \bar{v}_{2 2} \\
v_{1 1} & v_{1 2} & \bar{u}_{1 1} & \bar{u}_{1 2} \\
v_{2 1} & v_{2 2} & \bar{u}_{2 1} & \bar{u}_{2 2}\end{array}\right ) + \nonumber  \\
 & \left (\begin{array}{cccc}\bar{u}_{1 1} & \bar{u}_{2 1} & \bar{v}_{1 1} & \bar{v}_{2 1} \\
\bar{u}_{1 2} & \bar{u}_{2 2} & \bar{v}_{1 2} & \bar{v}_{2 2} \\
v_{1 1} & v_{2 1} & u_{1 1} & u_{2 1} \\
v_{1 2} & v_{2 2} & u_{1 2} & u_{2 2}\end{array}\right ) \left (\begin{array}{cc}\mathbb{O}_{2} & \mathbb{I}_{2} \\
\mathbb{I}_{2} & \mathbb{O}_{2}\end{array}\right ) \left (\left (\begin{array}{c}a_{1} \\
a_{2} \\
a_{1}^{ \dag } \\
a_{2}^{ \dag }\end{array}\right ) \left (\begin{array}{cccc}a_{1}^{ \dag } & a_{2}^{ \dag } & a_{1} & a_{2}\end{array}\right )\right )^{\mathfrak{t}} \times  \nonumber  \\
 & \left (\begin{array}{cc}\mathbb{O}_{2} & \mathbb{I}_{2} \\
\mathbb{I}_{2} & \mathbb{O}_{2}\end{array}\right ) \left (\begin{array}{cccc}u_{1 1} & u_{2 1} & v_{1 1} & v_{1 2} \\
u_{1 2} & u_{2 2} & v_{2 1} & v_{2 2} \\
\bar{v}_{1 1} & \bar{v}_{2 1} & \bar{u}_{1 1} & \bar{u}_{2 1} \\
\bar{v}_{1 2} & \bar{v}_{2 2} & \bar{u}_{1 2} & \bar{u}_{2 2}\end{array}\right ) =\left (\begin{array}{cccc}1 & 0 & 0 & 0 \\
0 & 1 & 0 & 0 \\
0 & 0 & 1 & 0 \\
0 & 0 & 0 & 1\end{array}\right )\end{align}giving that
\begin{equation}\left (\begin{array}{cccc}\bar{u}_{1 1} & \bar{u}_{2 1} & \bar{v}_{1 1} & \bar{v}_{2 1} \\
\bar{u}_{1 2} & \bar{u}_{2 2} & \bar{v}_{1 2} & \bar{v}_{2 2} \\
v_{1 1} & v_{2 1} & u_{1 1} & u_{2 1} \\
v_{1 2} & v_{2 2} & u_{1 2} & u_{2 2}\end{array}\right ) \left (\begin{array}{cccc}u_{1 1} & u_{1 2} & \bar{v}_{1 1} & \bar{v}_{1 2} \\
u_{2 1} & u_{2 2} & \bar{v}_{2 1} & \bar{v}_{2 2} \\
v_{1 1} & v_{1 2} & \bar{u}_{1 1} & \bar{u}_{1 2} \\
v_{2 1} & v_{2 2} & \bar{u}_{2 1} & \bar{u}_{2 2}\end{array}\right ) =\left (\begin{array}{cccc}1 & 0 & 0 & 0 \\
0 & 1 & 0 & 0 \\
0 & 0 & 1 & 0 \\
0 & 0 & 0 & 1\end{array}\right )
\end{equation}
\end{example}


\begin{remark}
Recall
\begin{align}\mathbf{a} &  =\left (a_{1} ,\cdots  ,a_{r}\right ) \\
\mathbf{a}^{ \dag } &  =\left (a_{1}^{ \dag } ,\cdots  ,a_{r}^{ \dag }\right ) \\
\mathbf{a}^{\mathfrak{a}} &  =\left (\begin{array}{c}a_{1}^{ \dag } \\
\vdots  \\
a_{r}^{ \dag }\end{array}\right ) =\left (\mathbf{a}^{ \dag }\right )^{\mathfrak{t}} \\
\mathbf{a}^{\mathfrak{t}} &  =\left (\begin{array}{c}a_{1} \\
\vdots  \\
a_{r}\end{array}\right )\text{.}\end{align}
\end{remark}

Canonical transformation have the following effect
\begin{align} & \left (\begin{array}{cc}\mathbb{U} & \overline{\mathbb{V}} \\
\mathbb{V} & \overline{\mathbb{U}}\end{array}\right )^{ \dag } \left (\begin{array}{c}\mathbf{a}^{\mathfrak{t}} \\
\mathbf{a}^{\mathfrak{a}}\end{array}\right ) \left (\begin{array}{cc}\mathbf{a}^{ \dag } & \mathbf{a}\end{array}\right ) \left (\begin{array}{cc}\mathbb{U} & \overline{\mathbb{V}} \\
\mathbb{V} & \overline{\mathbb{U}}\end{array}\right ) \nonumber  \\
 &  +\left (\begin{array}{cc}\mathbb{O}_{r} & \mathbb{I}_{r} \\
\mathbb{I}_{r} & \mathbb{O}_{r}\end{array}\right ) \left (\left (\begin{array}{cc}\mathbb{U} & \overline{\mathbb{V}} \\
\mathbb{V} & \overline{\mathbb{U}}\end{array}\right )^{ \dag } \left (\begin{array}{c}\mathbf{a}^{\mathfrak{t}} \\
\mathbf{a}^{\mathfrak{a}}\end{array}\right ) \left (\begin{array}{cc}\mathbf{a}^{ \dag } & \mathbf{a}\end{array}\right ) \left (\begin{array}{cc}\mathbb{U} & \overline{\mathbb{V}} \\
\mathbb{V} & \overline{\mathbb{U}}\end{array}\right )\right )^{t} \left (\begin{array}{cc}\mathbb{O}_{r} & \mathbb{I}_{r} \\
\mathbb{I}_{r} & \mathbb{O}_{r}\end{array}\right ) \nonumber  \\
 &  =\left (\begin{array}{cc}\mathbb{U}^{ \dag } & \mathbb{V}^{ \dag } \\
\mathbb{V}^{t} & \mathbb{U}^{t}\end{array}\right ) \left (\begin{array}{c}\mathbf{a}^{\mathfrak{t}} \\
\mathbf{a}^{\mathfrak{a}}\end{array}\right ) \left (\begin{array}{cc}\mathbf{a}^{ \dag } & \mathbf{a}\end{array}\right ) \left (\begin{array}{cc}\mathbb{U} & \overline{\mathbb{V}} \\
\mathbb{V} & \overline{\mathbb{U}}\end{array}\right ) \nonumber  \\
 &  +\left (\begin{array}{cc}\mathbb{O}_{r} & \mathbb{I}_{r} \\
\mathbb{I}_{r} & \mathbb{O}_{r}\end{array}\right ) \left (\begin{array}{cc}\mathbb{U}^{t} & \mathbb{V}^{t} \\
\mathbb{V}^{ \dag } & \mathbb{U}^{ \dag }\end{array}\right ) \left (\left (\begin{array}{c}\mathbf{a} \\
\mathbf{a}^{\mathfrak{a}}\end{array}\right ) \left (\begin{array}{cc}\mathbf{a}^{ \dag } & \mathbf{a}\end{array}\right )\right )^{t} \left (\begin{array}{cc}\overline{\mathbb{U}} & \mathbb{V} \\
\overline{\mathbb{V}} & \mathbb{U}\end{array}\right ) \left (\begin{array}{cc}\mathbb{O}_{r} & \mathbb{I}_{r} \\
\mathbb{I}_{r} & \mathbb{O}_{r}\end{array}\right ) \nonumber  \\
 &  =\left (\begin{array}{cc}\mathbb{U}^{ \dag } & \mathbb{V}^{ \dag } \\
\mathbb{V}^{t} & \mathbb{U}^{t}\end{array}\right ) \left (\begin{array}{c}\mathbf{a}^{\mathfrak{t}} \\
\mathbf{a}^{\mathfrak{a}}\end{array}\right ) \left (\begin{array}{cc}\mathbf{a}^{ \dag } & \mathbf{a}\end{array}\right ) \left (\begin{array}{cc}\mathbb{U} & \overline{\mathbb{V}} \\
\mathbb{V} & \overline{\mathbb{U}}\end{array}\right ) \nonumber  \\
 &  +\left (\begin{array}{cc}\mathbb{V}^{ \dag } & \mathbb{U}^{ \dag } \\
\mathbb{U}^{t} & \mathbb{V}^{t}\end{array}\right ) \left (\begin{array}{cc}\mathbb{O}_{r} & \mathbb{I}_{r} \\
\mathbb{I}_{r} & \mathbb{O}_{r}\end{array}\right ) \left (\left (\begin{array}{c}\mathbf{a}^{\mathfrak{t}} \\
\mathbf{a}^{\mathfrak{a}}\end{array}\right ) \left (\begin{array}{cc}\mathbf{a}^{ \dag } & \mathbf{a}\end{array}\right )\right )^{t} \left (\begin{array}{cc}\mathbb{O}_{r} & \mathbb{I}_{r} \\
\mathbb{I}_{r} & \mathbb{O}_{r}\end{array}\right ) \left (\begin{array}{cc}\mathbb{V} & \overline{\mathbb{U}} \\
\mathbb{U} & \overline{\mathbb{V}}\end{array}\right ) \nonumber  \\
 &  =\left (\begin{array}{cc}\mathbb{U}^{ \dag } & \mathbb{V}^{ \dag } \\
\mathbb{V}^{t} & \mathbb{U}^{t}\end{array}\right ) \left \{\left (\begin{array}{c}\mathbf{a}^{\mathfrak{t}} \\
\mathbf{a}^{\mathfrak{a}}\end{array}\right ) \left (\begin{array}{cc}\mathbf{a}^{ \dag } & \mathbf{a}\end{array}\right ) +\left (\begin{array}{cc}\mathbb{O}_{r} & \mathbb{I}_{r} \\
\mathbb{I}_{r} & \mathbb{O}_{r}\end{array}\right ) \left (\left (\begin{array}{c}\mathbf{a}^{\mathfrak{t}} \\
\mathbf{a}^{\mathfrak{a}}\end{array}\right ) \left (\begin{array}{cc}\mathbf{a}^{ \dag } & \mathbf{a}\end{array}\right )\right )^{t} \left (\begin{array}{cc}\mathbb{O}_{r} & \mathbb{I}_{r} \\
\mathbb{I}_{r} & \mathbb{O}_{r}\end{array}\right )\right \} \left (\begin{array}{cc}\mathbb{V} & \overline{\mathbb{U}} \\
\mathbb{U} & \overline{\mathbb{V}}\end{array}\right ) \nonumber  \\
 &  =\left (\begin{array}{cc}\mathbb{U}^{ \dag } & \mathbb{V}^{ \dag } \\
\mathbb{V}^{t} & \mathbb{U}^{t}\end{array}\right ) \left (\begin{array}{cc}\mathbb{I}_{r} & \mathbb{O}_{r} \\
\mathbb{O}_{r} & \mathbb{I}_{r}\end{array}\right ) \left (\begin{array}{cc}\mathbb{V} & \overline{\mathbb{U}} \\
\mathbb{U} & \overline{\mathbb{V}}\end{array}\right ) =\left (\begin{array}{cc}\mathbb{I}_{r} & \mathbb{O}_{r} \\
\mathbb{O}_{r} & \mathbb{I}_{r}\end{array}\right )\text{.}\end{align}\emph{That confirms that }$T \in U \left (2 r ,\mathbb{R}\right )\text{.}$ 


\begin{enumerate}
\item
\begin{align}\left (\begin{array}{cc}\mathbb{T}_{1 1} & \mathbb{T}_{1 2} \\
\overline{\mathbb{T}}_{1 2} & \overline{\mathbb{T}}_{1 1}\end{array}\right ) \left (\begin{array}{cc}\mathbb{O}_{r} & \mathbb{I}_{r} \\
\mathbb{I}_{r} & \mathbb{O}_{r}\end{array}\right ) &  =\left (\begin{array}{cc}\mathbb{T}_{1 2} & \mathbb{T}_{1 1} \\
\overline{\mathbb{T}}_{1 1} & \overline{\mathbb{T}}_{1 2}\end{array}\right ) \\
\left (\begin{array}{cc}\mathbb{O}_{r} & \mathbb{I}_{r} \\
\mathbb{I}_{r} & \mathbb{O}_{r}\end{array}\right ) \left (\begin{array}{cc}\mathbb{T}_{1 1} & \mathbb{T}_{1 2} \\
\overline{\mathbb{T}}_{1 2} & \overline{\mathbb{T}}_{1 1}\end{array}\right ) &  =\left (\begin{array}{cc}\overline{\mathbb{T}}_{1 2} & \overline{\mathbb{T}}_{1 1} \\
\mathbb{T}_{1 1} & \mathbb{T}_{1 2}\end{array}\right ) \\
\left (\begin{array}{cc}\mathbb{O}_{r} & \mathbb{I}_{r} \\
\mathbb{I}_{r} & \mathbb{O}_{r}\end{array}\right ) \left (\begin{array}{cc}\mathbb{T}_{1 1} & \mathbb{T}_{1 2} \\
\overline{\mathbb{T}}_{1 2} & \overline{\mathbb{T}}_{1 1}\end{array}\right ) \left (\begin{array}{cc}\mathbb{O}_{r} & \mathbb{I}_{r} \\
\mathbb{I}_{r} & \mathbb{O}_{r}\end{array}\right ) &  =\left (\begin{array}{cc}\mathbb{O}_{r} & \mathbb{I}_{r} \\
\mathbb{I}_{r} & \mathbb{O}_{r}\end{array}\right ) \left (\begin{array}{cc}\mathbb{T}_{1 2} & \mathbb{T}_{1 1} \\
\overline{\mathbb{T}}_{1 1} & \overline{\mathbb{T}}_{1 2}\end{array}\right ) =\left (\begin{array}{cc}\mathbb{O}_{r} & \mathbb{I}_{r} \\
\mathbb{I}_{r} & \mathbb{O}_{r}\end{array}\right )\text{.}\end{align}

\item If $T$ preserves the CAR and is $^{ \ast }$- transformation then
\begin{equation}\mathbb{T} =\left (\begin{array}{cc}\mathbb{T}_{1 1} & \mathbb{T}_{1 2} \\
\overline{\mathbb{T}}_{1 2} & \overline{\mathbb{T}}_{1 1}\end{array}\right )
\end{equation}and
\begin{equation}\mathbb{T} \in U \left (2 r\right )\text{.}
\end{equation}\end{enumerate}


\subsubsection{Transpose Form of the CAR}
\begin{example}
\begin{equation}\left (\begin{array}{cccc}a_{1}^{ \dag } a_{1}^{ \dag } & a_{1}^{ \dag } a_{2}^{ \dag } & a_{1}^{ \dag } a_{1} & a_{1}^{ \dag } a_{2} \\
a_{2}^{ \dag } a_{1}^{ \dag } & a_{2}^{ \dag } a_{2}^{ \dag } & a_{2}^{ \dag } a_{1} & a_{2}^{ \dag } a_{2} \\
a_{1} a_{1}^{ \dag } & a_{1} a_{2}^{ \dag } & a_{1} a_{1} & a_{1} a_{2} \\
a_{2} a_{1}^{ \dag } & a_{2} a_{2}^{ \dag } & a_{2} a_{1} & a_{2} a_{2}\end{array}\right ) +\left (\begin{array}{cccc}a_{1}^{ \dag } a_{1}^{ \dag } & a_{2}^{ \dag } a_{1}^{ \dag } & a_{1} a_{1}^{ \dag } & a_{2} a_{1}^{ \dag } \\
a_{1}^{ \dag } a_{2}^{ \dag } & a_{2}^{ \dag } a_{2}^{ \dag } & a_{1} a_{2}^{ \dag } & a_{2} a_{2}^{ \dag } \\
a_{1}^{ \dag } a_{1} & a_{2}^{ \dag } a_{1} & a_{1} a_{1} & a_{2} a_{1} \\
a_{1}^{ \dag } a_{2} & a_{2}^{ \dag } a_{2} & a_{1} a_{2} & a_{2} a_{2}\end{array}\right ) =\left (\begin{array}{cccc}0 & 0 & 1 & 0 \\
0 & 0 & 0 & 1 \\
1 & 0 & 0 & 0 \\
0 & 1 & 0 & 0\end{array}\right )\text{,}
\end{equation}giving
\begin{gather}\left (\begin{array}{cc}\left (\begin{array}{c}a_{1}^{ \dag } \\
a_{2}^{ \dag }\end{array}\right ) \left (\begin{array}{cc}a_{1}^{ \dag } & a_{2}^{ \dag }\end{array}\right ) & \left (\begin{array}{c}a_{1}^{ \dag } \\
a_{2}^{ \dag }\end{array}\right ) \left (\begin{array}{cc}a_{1} & a_{2}\end{array}\right ) \\
\left (\begin{array}{c}a_{1} \\
a_{2}\end{array}\right ) \left (\begin{array}{cc}a_{1}^{ \dag } & a_{2}^{ \dag }\end{array}\right ) & \left (\begin{array}{c}a_{1} \\
a_{2}\end{array}\right ) \left (\begin{array}{cc}a_{1} & a_{2}\end{array}\right )\end{array}\right ) +\left (\begin{array}{cc}\left (\left (\begin{array}{c}a_{1}^{ \dag } \\
a_{2}^{ \dag }\end{array}\right ) \left (\begin{array}{cc}a_{1}^{ \dag } & a_{2}^{ \dag }\end{array}\right )\right )^{\mathfrak{t}} & \left (\left (\begin{array}{c}a_{1} \\
a_{2}\end{array}\right ) \left (\begin{array}{cc}a_{1}^{ \dag } & a_{2}^{ \dag }\end{array}\right )\right )^{\mathfrak{t}} \\
\left (\left (\begin{array}{c}a_{1}^{ \dag } \\
a_{2}^{ \dag }\end{array}\right ) \left (\begin{array}{cc}a_{1} & a_{2}\end{array}\right )\right )^{\mathfrak{t}} & \left (\left (\begin{array}{c}a_{1} \\
a_{2}\end{array}\right ) \left (\begin{array}{cc}a_{1} & a_{2}\end{array}\right )\right )^{\mathfrak{t}}\end{array}\right ) \\
 =\left (\begin{array}{cc}\left (\begin{array}{c}a_{1}^{ \dag } \\
a_{2}^{ \dag }\end{array}\right ) \left (\begin{array}{cc}a_{1}^{ \dag } & a_{2}^{ \dag }\end{array}\right ) & \left (\begin{array}{c}a_{1}^{ \dag } \\
a_{2}^{ \dag }\end{array}\right ) \left (\begin{array}{cc}a_{1} & a_{2}\end{array}\right ) \\
\left (\begin{array}{c}a_{1} \\
a_{2}\end{array}\right ) \left (\begin{array}{cc}a_{1}^{ \dag } & a_{2}^{ \dag }\end{array}\right ) & \left (\begin{array}{c}a_{1} \\
a_{2}\end{array}\right ) \left (\begin{array}{cc}a_{1} & a_{2}\end{array}\right )\end{array}\right ) +\left (\begin{array}{cc}\left (\begin{array}{c}a_{1}^{ \dag } \\
a_{2}^{ \dag }\end{array}\right ) \left (\begin{array}{cc}a_{1}^{ \dag } & a_{2}^{ \dag }\end{array}\right ) & \left (\begin{array}{c}a_{1}^{ \dag } \\
a_{2}^{ \dag }\end{array}\right ) \left (\begin{array}{cc}a_{1} & a_{2}\end{array}\right ) \\
\left (\begin{array}{c}a_{1} \\
a_{2}\end{array}\right ) \left (\begin{array}{cc}a_{1}^{ \dag } & a_{2}^{ \dag }\end{array}\right ) & \left (\begin{array}{c}a_{1} \\
a_{2}\end{array}\right ) \left (\begin{array}{cc}a_{1} & a_{2}\end{array}\right )\end{array}\right )^{\mathfrak{t}} = \\
 =\left (\begin{array}{cccc}0 & 0 & 1 & 0 \\
0 & 0 & 0 & 1 \\
1 & 0 & 0 & 0 \\
0 & 1 & 0 & 0\end{array}\right )\end{gather}
\end{example}

In general the transpose form of the CAR is
\begin{align} & \left (\begin{array}{cc}\left \{\mathbf{a}^{ \dag } ,\mathbf{a}^{ \dag }\right \} & \left \{\mathbf{a}^{ \dag } ,\mathbf{a}\right \} \\
\left \{\mathbf{a} ,\mathbf{a}^{ \dag }\right \} & \left \{\mathbf{a} ,\mathbf{a}\right \}\end{array}\right )\overset{\triangle }{ =}\left (\begin{array}{cc}\mathbf{a}^{\mathfrak{a}} \mathbf{a}^{ \dag } & \mathbf{a}^{\mathfrak{a}} \mathbf{a} \\
\mathbf{a}^{\mathfrak{t}} \mathbf{a}^{ \dag } & \mathbf{a}^{\mathfrak{t}} \mathbf{a}\end{array}\right ) +\left (\begin{array}{cc}\mathbf{a}^{\mathfrak{a}} \mathbf{a}^{ \dag } & \mathbf{a}^{\mathfrak{a}} \mathbf{a} \\
\mathbf{a}^{\mathfrak{t}} \mathbf{a}^{ \dag } & \mathbf{a}^{\mathfrak{t}} \mathbf{a}\end{array}\right )^{\mathfrak{t}} \nonumber  \\
 &  =\left (\begin{array}{cc}\mathbf{a}^{\mathfrak{a}} \mathbf{a}^{ \dag } & \mathbf{a}^{\mathfrak{a}} \mathbf{a} \\
\mathbf{a}^{\mathfrak{t}} \mathbf{a}^{ \dag } & \mathbf{a}^{\mathfrak{t}} \mathbf{a}\end{array}\right ) +\left (\begin{array}{cc}\left (\mathbf{a}^{\mathfrak{a}} \mathbf{a}^{ \dag }\right )^{\mathfrak{t}} & \left (\mathbf{a}^{\mathfrak{a}} \mathbf{a}^{ \dag }\right )^{\mathfrak{t}} \\
\left (\mathbf{a}^{\mathfrak{a}} \mathbf{a}\right )^{\mathfrak{t}} & \left (\mathbf{a}^{\mathfrak{t}} \mathbf{a}\right )^{\mathfrak{t}}\end{array}\right ) \nonumber  \\
 &  =\left (\begin{array}{c}\mathbf{a}^{\mathfrak{a}} \\
\mathbf{a}^{\mathfrak{t}}\end{array}\right ) \left (\begin{array}{cc}\mathbf{a}^{ \dag } & \mathbf{a}\end{array}\right ) +\left (\left (\begin{array}{c}\mathbf{a}^{\mathfrak{a}} \\
\mathbf{a}^{\mathfrak{t}}\end{array}\right ) \left (\begin{array}{cc}\mathbf{a}^{ \dag } & \mathbf{a}\end{array}\right )\right )^{\mathfrak{t}} \nonumber  \\
 &  =\left (\begin{array}{cc}\mathbf{a}^{ \dag } & \mathbf{a}\end{array}\right )^{\mathfrak{t}} \left (\begin{array}{cc}\mathbf{a}^{ \dag } & \mathbf{a}\end{array}\right ) +\left (\left (\begin{array}{cc}\mathbf{a}^{ \dag } & \mathbf{a}\end{array}\right )^{\mathfrak{t}} \left (\begin{array}{cc}\mathbf{a}^{ \dag } & \mathbf{a}\end{array}\right )\right )^{\mathfrak{t}} \nonumber  \\
 &  =\left (\begin{array}{cc}\mathbb{O}_{r} & \mathbb{I}_{r} \\
\mathbb{I}_{r} & \mathbb{O}_{r}\end{array}\right )\text{.}\end{align}

\paragraph{Canonical Transformations for transpose form}
Canonical transformation gives
\begin{align} & \left (\begin{array}{cc}\left \{\mathbf{b}^{ \dag } ,\mathbf{b}^{ \dag }\right \} & \left \{\mathbf{b}^{ \dag } ,\mathbf{b}\right \} \\
\left \{\mathbf{b} ,\mathbf{b}\right \} & \left \{\mathbf{b}^{ \dag } ,\mathbf{b}\right \}\end{array}\right ) =\left (\begin{array}{cc}\mathbb{O}_{r} & \mathbb{I}_{r} \\
\mathbb{I}_{r} & \mathbb{O}_{r}\end{array}\right ) \nonumber  \\
 &  =\left (\begin{array}{cc}\mathbb{U} & \overline{\mathbb{V}} \\
\mathbb{V} & \overline{\mathbb{U}}\end{array}\right )^{t} \left (\begin{array}{cc}\mathbf{a}^{ \dag } & \mathbf{a}\end{array}\right )^{\mathfrak{t}} \left (\begin{array}{cc}\mathbf{a}^{ \dag } & \mathbf{a}\end{array}\right ) \left (\begin{array}{cc}\mathbb{U} & \overline{\mathbb{V}} \\
\mathbb{V} & \overline{\mathbb{U}}\end{array}\right ) +\left (\left (\begin{array}{cc}\mathbb{U} & \overline{\mathbb{V}} \\
\mathbb{V} & \overline{\mathbb{U}}\end{array}\right )^{t} \left (\begin{array}{cc}\mathbf{a}^{ \dag } & \mathbf{a}\end{array}\right )^{\mathfrak{t}} \left (\begin{array}{cc}\mathbf{a}^{ \dag } & \mathbf{a}\end{array}\right ) \left (\begin{array}{cc}\mathbb{U} & \overline{\mathbb{V}} \\
\mathbb{V} & \overline{\mathbb{U}}\end{array}\right )\right )^{t} \\
 &  =\left (\begin{array}{cc}\mathbb{U} & \overline{\mathbb{V}} \\
\mathbb{V} & \overline{\mathbb{U}}\end{array}\right )^{t} \left (\begin{array}{cc}\mathbf{a}^{ \dag } & \mathbf{a}\end{array}\right )^{\mathfrak{t}} \left (\begin{array}{cc}\mathbf{a}^{ \dag } & \mathbf{a}\end{array}\right ) \left (\begin{array}{cc}\mathbb{U} & \overline{\mathbb{V}} \\
\mathbb{V} & \overline{\mathbb{U}}\end{array}\right ) +\left (\begin{array}{cc}\mathbb{U} & \overline{\mathbb{V}} \\
\mathbb{V} & \overline{\mathbb{U}}\end{array}\right )^{t} \left (\left (\begin{array}{cc}\mathbf{a}^{ \dag } & \mathbf{a}\end{array}\right )^{\mathfrak{t}} \left (\begin{array}{cc}\mathbf{a}^{ \dag } & \mathbf{a}\end{array}\right )\right )^{\mathfrak{t}} \left (\begin{array}{cc}\mathbb{U} & \overline{\mathbb{V}} \\
\mathbb{V} & \overline{\mathbb{U}}\end{array}\right ) \nonumber  \\
 &  =\left (\begin{array}{cc}\mathbb{U} & \overline{\mathbb{V}} \\
\mathbb{V} & \overline{\mathbb{U}}\end{array}\right )^{t} \left (\begin{array}{cc}\left \{\mathbf{a}^{ \dag } ,\mathbf{a}^{ \dag }\right \} & \left \{\mathbf{a}^{ \dag } ,\mathbf{a}\right \} \\
\left \{\mathbf{a} ,\mathbf{a}\right \} & \left \{\mathbf{a}^{ \dag } ,\mathbf{a}\right \}\end{array}\right ) \left (\begin{array}{cc}\mathbb{U} & \overline{\mathbb{V}} \\
\mathbb{V} & \overline{\mathbb{U}}\end{array}\right ) \nonumber  \\
 &  =\left (\begin{array}{cc}\mathbb{U} & \overline{\mathbb{V}} \\
\mathbb{V} & \overline{\mathbb{U}}\end{array}\right )^{t} \left (\begin{array}{cc}\mathbb{O}_{r} & \mathbb{I}_{r} \\
\mathbb{I}_{r} & \mathbb{O}_{r}\end{array}\right ) \left (\begin{array}{cc}\mathbb{U} & \overline{\mathbb{V}} \\
\mathbb{V} & \overline{\mathbb{U}}\end{array}\right ) =\left (\begin{array}{cc}\mathbb{U}^{t} & \mathbb{V}^{t} \\
\overline{\mathbb{V}}^{t} & \overline{\mathbb{U}}^{t}\end{array}\right ) \left (\begin{array}{cc}\mathbb{O}_{r} & \mathbb{I}_{r} \\
\mathbb{I}_{r} & \mathbb{O}_{r}\end{array}\right ) \left (\begin{array}{cc}\mathbb{U} & \overline{\mathbb{V}} \\
\mathbb{V} & \overline{\mathbb{U}}\end{array}\right ) \nonumber  \\
 &  =\left (\begin{array}{cc}\mathbb{V}^{t} & \mathbb{U}^{t} \\
\overline{\mathbb{U}}^{t} & \overline{\mathbb{V}}^{t}\end{array}\right ) \left (\begin{array}{cc}\mathbb{U} & \overline{\mathbb{V}} \\
\mathbb{V} & \overline{\mathbb{U}}\end{array}\right ) =\left (\begin{array}{cc}\mathbb{V}^{t} \mathbb{U} +\mathbb{U}^{t} \mathbb{V} & \mathbb{V}^{t} \overline{\mathbb{V}} +\mathbb{U}^{t} \overline{\mathbb{U}} \\
\overline{\mathbb{U}}^{t} \mathbb{U} +\overline{\mathbb{V}}^{t} \mathbb{V} & \overline{\mathbb{U}}^{t} \overline{\mathbb{V}} +\overline{\mathbb{V}}^{t} \overline{\mathbb{U}}\end{array}\right ) =\left (\begin{array}{cc}\mathbb{O}_{r} & \mathbb{I}_{r} \\
\mathbb{I}_{r} & \mathbb{O}_{r}\end{array}\right )\text{.}\end{align}i.e.
\begin{equation}\frame{$\left (\begin{array}{cc}\mathbb{U} & \overline{\mathbb{V}} \\
\mathbb{V} & \overline{\mathbb{U}}\end{array}\right )^{t} \left (\begin{array}{cc}\mathbb{O}_{r} & \mathbb{I}_{r} \\
\mathbb{I}_{r} & \mathbb{O}_{r}\end{array}\right ) \left (\begin{array}{cc}\mathbb{U} & \overline{\mathbb{V}} \\
\mathbb{V} & \overline{\mathbb{U}}\end{array}\right ) =\left (\begin{array}{cc}\mathbb{O}_{r} & \mathbb{I}_{r} \\
\mathbb{I}_{r} & \mathbb{O}_{r}\end{array}\right )$}
\end{equation}implies that
\begin{align} & \left (\left (\begin{array}{cc}\mathbb{O}_{r} & \mathbb{I}_{r} \\
\mathbb{I}_{r} & \mathbb{O}_{r}\end{array}\right ) \left (\begin{array}{cc}\mathbb{U} & \overline{\mathbb{V}} \\
\mathbb{V} & \overline{\mathbb{U}}\end{array}\right ) \left (\begin{array}{cc}\mathbb{O}_{r} & \mathbb{I}_{r} \\
\mathbb{I}_{r} & \mathbb{O}_{r}\end{array}\right )\right )^{t} =\left (\left (\begin{array}{cc}\mathbb{V} & \overline{\mathbb{U}} \\
\mathbb{U} & \overline{\mathbb{V}}\end{array}\right ) \left (\begin{array}{cc}\mathbb{O}_{r} & \mathbb{I}_{r} \\
\mathbb{I}_{r} & \mathbb{O}_{r}\end{array}\right )\right )^{t} \nonumber  \\
 &  =\left (\begin{array}{cc}\overline{\mathbb{U}} & \mathbb{V} \\
\overline{\mathbb{V}} & \mathbb{U}\end{array}\right )^{t} =\left (\begin{array}{cc}\overline{\mathbb{U}}^{t} & \overline{\mathbb{V}}^{t} \\
\mathbb{V}^{t} & \mathbb{U}^{t}\end{array}\right ) =\left (\begin{array}{cc}\mathbb{U}^{ \dag } & \mathbb{V}^{ \dag } \\
\mathbb{V}^{t} & \mathbb{U}^{t}\end{array}\right ) =\left (\begin{array}{cc}\mathbb{U} & \overline{\mathbb{V}} \\
\mathbb{V} & \overline{\mathbb{U}}\end{array}\right )^{ \dag }\text{.}\end{align}


\begin{lemma}
The conditions
\begin{equation}\left (\begin{array}{cc}\mathbb{T}_{1 1} & \mathbb{T}_{1 2} \\
\mathbb{T}_{2 1} & \mathbb{T}_{2 2}\end{array}\right )^{t} \left (\begin{array}{cc}\mathbb{O}_{r} & \mathbb{I}_{r} \\
\mathbb{I}_{r} & \mathbb{O}_{r}\end{array}\right ) \left (\begin{array}{cc}\mathbb{T}_{1 1} & \mathbb{T}_{1 2} \\
\mathbb{T}_{2 1} & \mathbb{T}_{2 2}\end{array}\right ) =\left (\begin{array}{cc}\mathbb{O}_{r} & \mathbb{I}_{r} \\
\mathbb{I}_{r} & \mathbb{O}_{r}\end{array}\right )
\end{equation}and $T \in U \left (2 r\right )$
\begin{equation}\left (\begin{array}{cc}\mathbb{T}_{1 1} & \mathbb{T}_{1 2} \\
\mathbb{T}_{2 1} & \mathbb{T}_{2 2}\end{array}\right )^{ \dag } \left (\begin{array}{cc}\mathbb{T}_{1 1} & \mathbb{T}_{1 2} \\
\mathbb{T}_{2 1} & \mathbb{T}_{2 2}\end{array}\right ) =\left (\begin{array}{cc}\mathbb{I}_{r} & \mathbb{O}_{r} \\
\mathbb{O}_{r} & \mathbb{I}_{r}\end{array}\right )
\end{equation}imply $T$ is cannonical i.e. $T \in O \left (2 r ,\mathbb{R}\right )\text{.}$ 


\begin{proof}
These two conditions imply that
\begin{align}\left (\begin{array}{cc}\mathbb{O}_{r} & \mathbb{I}_{r} \\
\mathbb{I}_{r} & \mathbb{O}_{r}\end{array}\right ) \left (\begin{array}{cc}\mathbb{T}_{1 1} & \mathbb{T}_{1 2} \\
\mathbb{T}_{2 1} & \mathbb{T}_{2 2}\end{array}\right )^{t} \left (\begin{array}{cc}\mathbb{O}_{r} & \mathbb{I}_{r} \\
\mathbb{I}_{r} & \mathbb{O}_{r}\end{array}\right ) &  =\left (\begin{array}{cc}\mathbb{T}_{1 1} & \mathbb{T}_{1 2} \\
\mathbb{T}_{2 1} & \mathbb{T}_{2 2}\end{array}\right )^{ \dag } \\
\left (\begin{array}{cc}\mathbb{O}_{r} & \mathbb{I}_{r} \\
\mathbb{I}_{r} & \mathbb{O}_{r}\end{array}\right ) \left (\begin{array}{cc}\mathbb{T}_{1 1}^{t} & \mathbb{T}_{2 1}^{t} \\
\mathbb{T}_{1 2}^{t} & \mathbb{T}_{2 2}^{t}\end{array}\right ) \left (\begin{array}{cc}\mathbb{O}_{r} & \mathbb{I}_{r} \\
\mathbb{I}_{r} & \mathbb{O}_{r}\end{array}\right ) &  =\left (\begin{array}{cc}\mathbb{T}_{1 1}^{ \dag } & \mathbb{T}_{2 1}^{ \dag } \\
\mathbb{T}_{1 2}^{ \dag } & \mathbb{T}_{2 2}^{ \dag }\end{array}\right ) \\
\left (\begin{array}{cc}\mathbb{O}_{r} & \mathbb{I}_{r} \\
\mathbb{I}_{r} & \mathbb{O}_{r}\end{array}\right ) \left (\begin{array}{cc}\mathbb{T}_{2 1}^{t} & \mathbb{T}_{1 1}^{t} \\
\mathbb{T}_{2 2}^{t} & \mathbb{T}_{1 2}^{t}\end{array}\right ) &  =\left (\begin{array}{cc}\mathbb{T}_{1 1}^{ \dag } & \mathbb{T}_{2 1}^{ \dag } \\
\mathbb{T}_{1 2}^{ \dag } & \mathbb{T}_{2 2}^{ \dag }\end{array}\right ) \\
\left (\begin{array}{cc}\mathbb{T}_{2 2}^{t} & \mathbb{T}_{1 2}^{t} \\
\mathbb{T}_{2 1}^{t} & \mathbb{T}_{1 1}^{t}\end{array}\right ) &  =\left (\begin{array}{cc}\mathbb{T}_{1 1}^{ \dag } & \mathbb{T}_{2 1}^{ \dag } \\
\mathbb{T}_{1 2}^{ \dag } & \mathbb{T}_{2 2}^{ \dag }\end{array}\right )\end{align}
\begin{equation} \Rightarrow \left (\begin{array}{cc}\mathbb{T}_{2 2} & \mathbb{T}_{2 1} \\
\mathbb{T}_{1 2} & \mathbb{T}_{1 1}\end{array}\right ) =\left (\begin{array}{cc}\bar{\mathbb{T}}_{1 1} & \bar{\mathbb{T}}_{1 2} \\
\bar{\mathbb{T}}_{2 1} & \bar{\mathbb{T}}_{2 2}\end{array}\right )
\end{equation}giving
\begin{equation}\left (\begin{array}{cc}\mathbb{T}_{1 1} & \mathbb{T}_{1 2} \\
\mathbb{T}_{2 1} & \mathbb{T}_{2 2}\end{array}\right ) =\left (\begin{array}{cc}\mathbb{T}_{1 1} & \bar{\mathbb{T}}_{2 1} \\
\mathbb{T}_{1 2} & \bar{\mathbb{T}}_{1 1}\end{array}\right )\text{,}
\end{equation}thus $T$ is cannonical. 
\end{proof}


\end{lemma}

\subsection{General Transformations $ \in G L \left (2 r ,\mathbb{C}\right )$}
In order to study more general transformation it is convenient to considers the self adjoint basis
\begin{equation}\left (\begin{array}{cc}\mathbf{x}_{1} & \mathbf{x}_{2}\end{array}\right ) =K \left (\begin{array}{cc}\mathbf{a}^{ \dag } & \mathbf{a}\end{array}\right ) =\left (\begin{array}{cc}\mathbf{a}^{ \dag } & \mathbf{a}\end{array}\right ) \mathbb{K}
\end{equation}and the self adjoint form of the CAR i.e. the Clifford Commutation Relationships
\begin{equation*}\left (\begin{array}{cc}\left \{\mathbf{x}_{1} ,\mathbf{x}_{1}\right \} & \left \{\mathbf{x}_{1} ,\mathbf{x}_{2}\right \} \\
\left \{\mathbf{x}_{2} ,\mathbf{x}_{1}\right \} & \left \{\mathbf{x}_{2} ,\mathbf{x}_{2}\right \}\end{array}\right ) =\left (\begin{array}{cc}\mathbf{x}_{1}^{\mathfrak{t}} \mathbf{x}_{1} & \mathbf{x}_{1}^{\mathfrak{t}} \mathbf{x}_{2} \\
\mathbf{x}_{2}^{\mathfrak{t}} \mathbf{x}_{1} & \mathbf{x}_{2}^{\mathfrak{t}} \mathbf{x}_{2}\end{array}\right ) +\left (\begin{array}{cc}\mathbf{x}_{1}^{\mathfrak{t}} \mathbf{x}_{1} & \mathbf{x}_{1}^{\mathfrak{t}} \mathbf{x}_{2} \\
\mathbf{x}_{2}^{\mathfrak{t}} \mathbf{x}_{1} & \mathbf{x}_{2}^{\mathfrak{t}} \mathbf{x}_{2}\end{array}\right )^{\mathfrak{t}} =\left (\begin{array}{cc}\mathbb{I}_{r} & \mathbb{O} \\
\mathbb{O} & \mathbb{I}_{r}\end{array}\right )\text{,}
\end{equation*}where
\begin{equation}\left (\begin{array}{cc}\mathbf{x}_{1}^{\mathfrak{t}} \mathbf{x}_{1} & \mathbf{x}_{1}^{\mathfrak{t}} \mathbf{x}_{2} \\
\mathbf{x}_{2}^{\mathfrak{t}} \mathbf{x}_{1} & \mathbf{x}_{2}^{\mathfrak{t}} \mathbf{x}_{2}\end{array}\right ) =\left (\mathbf{x}_{1} \mathbf{x}_{2}\right )^{\mathfrak{t}} \left (\mathbf{x}_{1} \mathbf{x}_{2}\right ) =\left (\begin{array}{c}\mathbf{x}_{1}^{\mathfrak{t}} \\
\mathbf{x}_{2}^{\mathfrak{t}}\end{array}\right ) \left (\mathbf{x}_{1} \mathbf{x}_{2}\right )\text{.}
\end{equation}$\lceil $e.g. for $r =2$
\begin{equation}\left (\begin{array}{cccc}x_{1 1} x_{1 1} & x_{1 1} x_{1 2} & x_{1 1} x_{2 1} & x_{1 1} x_{2 2} \\
x_{1 2} x_{1 1} & x_{1 2} x_{1 2} & x_{1 2} x_{2 1} & x_{1 2} x_{2 2} \\
x_{2 1} x_{1 1} & x_{2 1} x_{1 2} & x_{2 1} x_{2 1} & x_{2 1} x_{2 2} \\
x_{2 2} x_{1 1} & x_{2 2} x_{1 2} & x_{2 2} x_{2 1} & x_{2 2} x_{2 2}\end{array}\right ) =\left (\begin{array}{c}x_{1 1} \\
x_{1 2} \\
x_{2 1} \\
x_{2 2}\end{array}\right ) \left (\begin{array}{cccc}x_{1 1} & x_{1 2} & x_{2 1} & x_{2 2}\end{array}\right )\text{.}
\end{equation}The CAR\ can then be expressed as
\begin{align} & \left (\begin{array}{cccc}x_{1 1} x_{1 1} & x_{1 1} x_{1 2} & x_{1 1} x_{2 1} & x_{1 1} x_{2 2} \\
x_{1 2} x_{1 1} & x_{1 2} x_{1 2} & x_{1 2} x_{2 1} & x_{1 2} x_{2 2} \\
x_{2 1} x_{1 1} & x_{2 1} x_{1 2} & x_{2 1} x_{2 1} & x_{2 1} x_{2 2} \\
x_{2 2} x_{1 1} & x_{2 2} x_{1 2} & x_{2 2} x_{2 1} & x_{2 2} x_{2 2}\end{array}\right ) +\left (\begin{array}{cccc}x_{1 1} x_{1 1} & x_{1 1} x_{1 2} & x_{1 1} x_{2 1} & x_{1 1} x_{2 2} \\
x_{1 2} x_{1 1} & x_{1 2} x_{1 2} & x_{1 2} x_{2 1} & x_{1 2} x_{2 2} \\
x_{2 1} x_{1 1} & x_{2 1} x_{1 2} & x_{2 1} x_{2 1} & x_{2 1} x_{2 2} \\
x_{2 2} x_{1 1} & x_{2 2} x_{1 2} & x_{2 2} x_{2 1} & x_{2 2} x_{2 2}\end{array}\right )^{\mathfrak{t}} \nonumber  \\
 &  =\left (\begin{array}{cccc}x_{1 1} x_{1 1} +x_{1 1} x_{1 1} & x_{1 1} x_{1 2} +x_{1 2} x_{1 1} & x_{1 1} x_{2 1} +x_{2 1} x_{1 1} & x_{1 1} x_{2 2} +x_{2 2} x_{1 1} \\
x_{1 2} x_{1 1} +x_{1 1} x_{1 2} & x_{1 2} x_{1 2} +x_{1 2} x_{1 2} & x_{1 2} x_{2 1} +x_{2 1} x_{1 2} & x_{1 2} x_{2 2} +x_{2 2} x_{1 2} \\
x_{2 1} x_{1 1} +x_{1 1} x_{2 1} & x_{2 1} x_{1 2} +x_{1 2} x_{2 1} & x_{2 1} x_{2 1} +x_{2 1} x_{2 1} & x_{2 1} x_{2 2} +x_{2 2} x_{2 1} \\
x_{2 2} x_{1 1} +x_{1 1} x_{2 2} & x_{2 2} x_{1 2} +x_{1 2} x_{2 2} & x_{2 2} x_{2 1} +x_{2 1} x_{2 2} & x_{2 2} x_{2 2} +x_{2 2} x_{2 2}\end{array}\right ) =\mathbb{I}_{4}\text{.}\end{align}$\rfloor $ 

Consider the transformation of the basis by $T \in G L \left (2 r ,\mathbb{C}\right )\text{,}$ \emph{which in general is not a }$ \ast $\emph{\ transformation},
\begin{equation}T \left (\begin{array}{cc}\mathbf{x}_{1} & \mathbf{x}_{2}\end{array}\right ) =\left (\mathbf{x}_{1} \mathbf{x}_{2}\right ) \mathbb{T} =\left (\mathbf{x}_{1} \mathbf{x}_{2}\right ) \left (\begin{array}{cc}\mathbb{T}_{1 1} & \mathbb{T}_{1 2} \\
\mathbb{T}_{2 1} & \mathbb{T}_{2 2}\end{array}\right ) =\left (\begin{array}{cc}\mathbf{y}_{1} & \mathbf{y}_{2}\end{array}\right )
\end{equation}giving
\begin{align} & \left (\begin{array}{cc}\left \{\mathbf{y}_{1} ,\mathbf{y}_{1}\right \} & \left \{\mathbf{y}_{1} ,\mathbf{y}_{2}\right \} \\
\left \{\mathbf{y}_{2} ,\mathbf{y}_{1}\right \} & \left \{\mathbf{y}_{2} ,\mathbf{y}_{2}\right \}\end{array}\right )\overset{\triangle }{ =}\left (\begin{array}{cc}\mathbf{y}_{1}^{\mathfrak{t}} \mathbf{y}_{1} & \mathbf{y}_{1}^{\mathfrak{t}} \mathbf{y}_{2} \\
\mathbf{y}_{2}^{\mathfrak{t}} \mathbf{y}_{1} & \mathbf{y}_{2}^{\mathfrak{t}} \mathbf{y}_{2}\end{array}\right ) +\left (\begin{array}{cc}\mathbf{y}_{1}^{\mathfrak{t}} \mathbf{y}_{1} & \mathbf{y}_{1}^{\mathfrak{t}} \mathbf{y}_{2} \\
\mathbf{y}_{2}^{\mathfrak{t}} \mathbf{y}_{1} & \mathbf{y}_{2}^{\mathfrak{t}} \mathbf{y}_{2}\end{array}\right )^{\mathfrak{t}} \nonumber  \\
 &  =\mathbb{T}^{t} \left (\begin{array}{cc}\mathbf{x}_{1}^{\mathfrak{t}} \mathbf{x}_{1} & \mathbf{x}_{1}^{\mathfrak{t}} \mathbf{x}_{2} \\
\mathbf{x}_{2}^{\mathfrak{t}} \mathbf{x}_{1} & \mathbf{x}_{2}^{\mathfrak{t}} \mathbf{x}_{2}\end{array}\right ) \mathbb{T} +\left (\mathbb{T}^{t} \left (\begin{array}{cc}\mathbf{x}_{1}^{\mathfrak{t}} \mathbf{x}_{1} & \mathbf{x}_{1}^{\mathfrak{t}} \mathbf{x}_{2} \\
\mathbf{x}_{2}^{\mathfrak{t}} \mathbf{x}_{1} & \mathbf{x}_{2}^{\mathfrak{t}} \mathbf{x}_{2}\end{array}\right ) \mathbb{T}\right )^{\mathfrak{t}} \\
 &  =\mathbb{T}^{t} \left (\begin{array}{cc}\mathbf{x}_{1}^{\mathfrak{t}} \mathbf{x}_{1} & \mathbf{x}_{1}^{\mathfrak{t}} \mathbf{x}_{2} \\
\mathbf{x}_{2}^{\mathfrak{t}} \mathbf{x}_{1} & \mathbf{x}_{2}^{\mathfrak{t}} \mathbf{x}_{2}\end{array}\right ) \mathbb{T} +\mathbb{T}^{t} \left (\begin{array}{cc}\mathbf{x}_{1}^{\mathfrak{t}} \mathbf{x}_{1} & \mathbf{x}_{1}^{\mathfrak{t}} \mathbf{x}_{2} \\
\mathbf{x}_{2}^{\mathfrak{t}} \mathbf{x}_{1} & \mathbf{x}_{2}^{\mathfrak{t}} \mathbf{x}_{2}\end{array}\right )^{\mathfrak{t}} \mathbb{T} \\
 &  =\mathbb{T}^{t} \left (\left (\begin{array}{cc}\left \{\mathbf{x}_{1} ,\mathbf{x}_{1}\right \} & \left \{\mathbf{x}_{1} ,\mathbf{x}_{2}\right \} \\
\left \{\mathbf{x}_{2} ,\mathbf{x}_{1}\right \} & \left \{\mathbf{x}_{2} ,\mathbf{x}_{2}\right \}\end{array}\right )\right ) \mathbb{T}\end{align}Thus $T$ preserves the Clifford Commutation Relationship form of the CAR $\ensuremath{\operatorname*{Iff}}$
\begin{equation}\mathbb{T}^{t} \mathbb{T} =\mathbb{I}_{2 r}
\end{equation}i.e. iff $T \in O (2 r ,\mathbb{C})\text{.}$ \emph{However such }$T^{ \prime } s$\emph{\ are not in general }$ \ast $\emph{-transformations} i.e.
\begin{equation}T \in O (2 r ,\mathbb{C}) \nRightarrow T \left (X^{ \dag }\right ) =T \left (X\right )^{ \dag }\text{,}
\end{equation}thus in general do not transform self adjoint bases to self adjoint bases or non self adjoint bases to non self adjoint
bases. 

\paragraph{K Transformation}
\begin{remark}
$K$ is a star transformation i.e.
\begin{equation}K \left (X^{ \dag }\right ) =K \left (X\right )^{ \dag } ;\text{\quad }X \in \mathcal{F}_{1}
\end{equation}
\end{remark}

If $\left \{\mathbf{x}_{1} ,\mathbf{x}_{2}\right \}$ are self adjoint then
\begin{align}\left (\begin{array}{cc}\mathbf{x}_{1} & \mathbf{x}_{2}\end{array}\right ) &  =K \left (\begin{array}{cc}\mathbf{a}^{ \dag } & \mathbf{a}\end{array}\right ) =\left (\begin{array}{cc}\mathbf{a}^{ \dag } & \mathbf{a}\end{array}\right ) \mathbb{K} \\
 &  =\left (\begin{array}{cc}\mathbf{a}^{ \dag } & \mathbf{a}\end{array}\right ) \left (\begin{array}{cc}\frac{1}{\sqrt{2}} \mathbb{I}_{r} & \frac{i}{\sqrt{2}} \mathbb{I}_{r} \\
\frac{1}{\sqrt{2}} \mathbb{I}_{r} &  -\frac{i}{\sqrt{2}} \mathbb{I}_{r}\end{array}\right )\text{,}\end{align}and the inverse transformation by
\begin{align}K^{ -1} \left (\begin{array}{cc}\mathbf{x}_{1} & \mathbf{x}_{2}\end{array}\right ) &  =\left (\begin{array}{cc}\mathbf{x}_{1} & \mathbf{x}_{2}\end{array}\right ) \mathbb{K}^{ \dag } =\left (\begin{array}{cc}\mathbf{x}_{1} & \mathbf{x}_{2}\end{array}\right ) \left (\begin{array}{cc}\frac{1}{\sqrt{2}} \mathbb{I}_{r} & \frac{1}{\sqrt{2}} \mathbb{I}_{r} \\
 -\frac{i}{\sqrt{2}} \mathbb{I}_{r} & \frac{i}{\sqrt{2}} \mathbb{I}_{r}\end{array}\right ) \\
 &  =\left (\begin{array}{cc}\mathbf{a}^{ \dag } & \mathbf{a}\end{array}\right ) \mathbb{K} \mathbb{K}^{ \dag } =\left (\begin{array}{cc}\mathbf{a}^{ \dag } & \mathbf{a}\end{array}\right )\text{,}\end{align}where 


\begin{enumerate}
\item
\begin{equation}\mathbb{K} =\left (\begin{array}{cc}\frac{1}{\sqrt{2}} \mathbb{I}_{r} & \frac{i}{\sqrt{2}} \mathbb{I}_{r} \\
\frac{1}{\sqrt{2}} \mathbb{I}_{r} &  -\frac{i}{\sqrt{2}} \mathbb{I}_{r}\end{array}\right ) \in U \left (2 r\right )
\end{equation}

\item
\begin{equation}\mathbb{K}^{t} =\left (\begin{array}{cc}\frac{1}{\sqrt{2}} \mathbb{I}_{r} & \frac{1}{\sqrt{2}} \mathbb{I}_{r} \\
\frac{i}{\sqrt{2}} \mathbb{I}_{r} &  -\frac{i}{\sqrt{2}} \mathbb{I}_{r}\end{array}\right ) \in U \left (2 r\right )
\end{equation}

\item
\begin{align}\mathbb{K}^{t} \mathbb{K} &  =\left (\begin{array}{cc}\frac{1}{\sqrt{2}} \mathbb{I}_{r} & \frac{1}{\sqrt{2}} \mathbb{I}_{r} \\
\frac{i}{\sqrt{2}} \mathbb{I}_{r} &  -\frac{i}{\sqrt{2}} \mathbb{I}_{r}\end{array}\right ) \left (\begin{array}{cc}\frac{1}{\sqrt{2}} \mathbb{I}_{r} & \frac{i}{\sqrt{2}} \mathbb{I}_{r} \\
\frac{1}{\sqrt{2}} \mathbb{I}_{r} &  -\frac{i}{\sqrt{2}} \mathbb{I}_{r}\end{array}\right ) =\left (\begin{array}{cc}\mathbb{I}_{r} & \mathbb{O} \\
\mathbb{O} &  -\mathbb{I}_{r}\end{array}\right ) \\
\mathbb{K} \mathbb{K}^{t} &  =\left (\begin{array}{cc}\frac{1}{\sqrt{2}} \mathbb{I}_{r} & \frac{i}{\sqrt{2}} \mathbb{I}_{r} \\
\frac{1}{\sqrt{2}} \mathbb{I}_{r} &  -\frac{i}{\sqrt{2}} \mathbb{I}_{r}\end{array}\right ) \left (\begin{array}{cc}\frac{1}{\sqrt{2}} \mathbb{I}_{r} & \frac{1}{\sqrt{2}} \mathbb{I}_{r} \\
\frac{i}{\sqrt{2}} \mathbb{I}_{r} &  -\frac{i}{\sqrt{2}} \mathbb{I}_{r}\end{array}\right ) =\left (\begin{array}{cc}\mathbb{O} & \mathbb{I}_{r} \\
\mathbb{I}_{r} & \mathbb{O}\end{array}\right )\end{align}

\item
\begin{equation}\mathbb{K} \mathbb{K}^{t} \left (\begin{array}{cc}\mathbb{O} & \mathbb{I}_{r} \\
\mathbb{I}_{r} & \mathbb{O}\end{array}\right ) =\mathbb{K} \left (\left (\begin{array}{cc}\mathbb{O} & \mathbb{I}_{r} \\
\mathbb{I}_{r} & \mathbb{O}\end{array}\right ) \mathbb{K}\right )^{t} =\left (\begin{array}{cc}\mathbb{I}_{r} & \mathbb{O} \\
\mathbb{O} & \mathbb{I}_{r}\end{array}\right )
\end{equation}

\item
\begin{equation}\left (\mathbb{K}^{t}\right )^{ \dag } \mathbb{K}^{t} =\left (\begin{array}{cc}\frac{1}{\sqrt{2}} \mathbb{I}_{r} &  -\frac{i}{\sqrt{2}} \mathbb{I}_{r} \\
\frac{1}{\sqrt{2}} \mathbb{I}_{r} & \frac{i}{\sqrt{2}} \mathbb{I}_{r}\end{array}\right ) \left (\begin{array}{cc}\frac{1}{\sqrt{2}} \mathbb{I}_{r} & \frac{1}{\sqrt{2}} \mathbb{I}_{r} \\
\frac{i}{\sqrt{2}} \mathbb{I}_{r} &  -\frac{i}{\sqrt{2}} \mathbb{I}_{r}\end{array}\right ) =\left (\begin{array}{cc}\mathbb{I}_{r} & \mathbb{O} \\
\mathbb{O} & \mathbb{I}_{r}\end{array}\right )
\end{equation}\end{enumerate}


$K$ is not cannonical as
\begin{align} & \left (\begin{array}{cc}\frac{1}{\sqrt{2}} \mathbb{I}_{r} & \frac{i}{\sqrt{2}} \mathbb{I}_{r} \\
\frac{1}{\sqrt{2}} \mathbb{I}_{r} &  -\frac{i}{\sqrt{2}} \mathbb{I}_{r}\end{array}\right )^{t} \left (\begin{array}{cc}\mathbb{O}_{r} & \mathbb{I}_{r} \\
\mathbb{I}_{r} & \mathbb{O}_{r}\end{array}\right ) \left (\begin{array}{cc}\frac{1}{\sqrt{2}} \mathbb{I}_{r} & \frac{i}{\sqrt{2}} \mathbb{I}_{r} \\
\frac{1}{\sqrt{2}} \mathbb{I}_{r} &  -\frac{i}{\sqrt{2}} \mathbb{I}_{r}\end{array}\right ) \nonumber  \\
 &  =\left (\begin{array}{cc}\frac{1}{\sqrt{2}} \mathbb{I}_{r} & \frac{1}{\sqrt{2}} \mathbb{I}_{r} \\
\frac{i}{\sqrt{2}} \mathbb{I}_{r} &  -\frac{i}{\sqrt{2}} \mathbb{I}_{r}\end{array}\right ) \left (\begin{array}{cc}\frac{1}{\sqrt{2}} \mathbb{I}_{r} &  -\frac{i}{\sqrt{2}} \mathbb{I}_{r} \\
\frac{1}{\sqrt{2}} \mathbb{I}_{r} & \frac{i}{\sqrt{2}} \mathbb{I}_{r}\end{array}\right ) =\left (\begin{array}{cc}\mathbb{I}_{r} & \mathbb{O}_{r} \\
\mathbb{O}_{r} & \mathbb{I}_{r}\end{array}\right ) \neq \left (\begin{array}{cc}\mathbb{O}_{r} & \mathbb{I}_{r} \\
\mathbb{I}_{r} & \mathbb{O}_{r}\end{array}\right )\text{,}\end{align}as
\begin{equation}\mathbb{K}^{t} \left (\begin{array}{cc}\mathbb{O}_{r} & \mathbb{I}_{r} \\
\mathbb{I}_{r} & \mathbb{O}_{r}\end{array}\right ) \mathbb{K} =\left (\begin{array}{cc}\mathbb{I}_{r} & \mathbb{O}_{r} \\
\mathbb{O}_{r} & \mathbb{I}_{r}\end{array}\right )\text{.}
\end{equation}This also shows that
\begin{equation}\left (\mathbb{K}^{ \dag }\right )^{t} \mathbb{K}^{t} \left (\begin{array}{cc}\mathbb{O}_{r} & \mathbb{I}_{r} \\
\mathbb{I}_{r} & \mathbb{O}_{r}\end{array}\right ) \mathbb{K} \mathbb{K}^{ \dag } =\left (\begin{array}{cc}\mathbb{O}_{r} & \mathbb{I}_{r} \\
\mathbb{I}_{r} & \mathbb{O}_{r}\end{array}\right ) =\left (\mathbb{K}^{ \dag }\right )^{t} \left (\begin{array}{cc}\mathbb{I}_{r} & \mathbb{O}_{r} \\
\mathbb{O}_{r} & \mathbb{I}_{r}\end{array}\right ) \mathbb{K}^{ \dag }\text{.}
\end{equation}

Let
\begin{equation}T :\mathcal{F}_{1} \rightarrow \mathcal{F}_{1}
\end{equation}then
\begin{equation}T \left (\begin{array}{cc}\mathbf{x}_{1} & \mathbf{x}_{2}\end{array}\right ) =\left (\begin{array}{cc}\mathbf{x}_{1} & \mathbf{x}_{2}\end{array}\right ) \mathbb{T}_{x} =\left (\begin{array}{cc}\mathbf{x}_{1} & \mathbf{x}_{2}\end{array}\right ) \mathbb{K}^{ \dag } \mathbb{K} \mathbb{T}_{x} \mathbb{K}^{ \dag } =\left (\begin{array}{cc}\mathbf{a}^{ \dag } & \mathbf{a}\end{array}\right ) \mathbb{K} \mathbb{T}_{x} \mathbb{K}^{ \dag } =\left (\begin{array}{cc}\mathbf{a}^{ \dag } & \mathbf{a}\end{array}\right ) \mathbb{T}_{a}\text{.}
\end{equation}

Let $\left \{\mathbf{x}_{1} ,\mathbf{x}_{2}\right \}$ be unrestricted and
\begin{equation}K^{ -1} =K^{ \dag }
\end{equation}then
\begin{align}K^{ -1} \left (\begin{array}{cc}\mathbf{x}_{1} & \mathbf{x}_{2}\end{array}\right ) &  =\left (\begin{array}{cc}\mathbf{x}_{1} & \mathbf{x}_{2}\end{array}\right ) \mathbb{K}^{ \dag } =\left (\begin{array}{cc}\mathbf{x}_{1} & \mathbf{x}_{2}\end{array}\right ) \left (\begin{array}{cc}\frac{1}{\sqrt{2}} \mathbb{I}_{r} & \frac{1}{\sqrt{2}} \mathbb{I}_{r} \\
 -\frac{i}{\sqrt{2}} \mathbb{I}_{r} & \frac{i}{\sqrt{2}} \mathbb{I}_{r}\end{array}\right ) \\
 &  =\left (\begin{array}{cc}\mathbf{a}^{ \dag } & \mathbf{a}\end{array}\right ) =\frac{1}{\sqrt{2}} \left (\begin{array}{cc}\mathbf{x}_{1} -i \mathbf{x}_{2} & \mathbf{x}_{1} +i \mathbf{x}_{2}\end{array}\right )\end{align}and $\left \{\mathbf{a}_{1} ,\mathbf{a}_{2}\right \}$ satisfy the CAR in block anti diagonal form
\begin{align} & \left (\begin{array}{cc}\left \{\mathbf{a}_{1} ,\mathbf{a}_{1}\right \} & \left \{\mathbf{a}_{1} ,\mathbf{a}_{2}\right \} \\
\left \{\mathbf{a}_{2} ,\mathbf{a}_{1}\right \} & \left \{\mathbf{a}_{2} ,\mathbf{a}_{2}\right \}\end{array}\right ) =\left (\begin{array}{cc}\mathbf{a}_{1} \mathbf{a}_{1} & \mathbf{a}_{1} \mathbf{a}_{2} \\
\mathbf{a}_{2} \mathbf{a}_{1} & \mathbf{a}_{2} \mathbf{a}_{2}\end{array}\right ) +\left (\begin{array}{cc}\mathbf{a}_{1} \mathbf{a}_{1} & \mathbf{a}_{1} \mathbf{a}_{2} \\
\mathbf{a}_{2} \mathbf{a}_{1} & \mathbf{a}_{2} \mathbf{a}_{2}\end{array}\right )^{t} \nonumber  \\
 &  =\left (\begin{array}{cc}\frac{1}{\sqrt{2}} \mathbb{I}_{r} & \frac{1}{\sqrt{2}} \mathbb{I}_{r} \\
 -\frac{i}{\sqrt{2}} \mathbb{I}_{r} & \frac{i}{\sqrt{2}} \mathbb{I}_{r}\end{array}\right )^{t} \left (\left (\begin{array}{cc}\mathbf{x}_{1} \mathbf{x}_{1} & \mathbf{x}_{1} \mathbf{x}_{2} \\
\mathbf{x}_{2} \mathbf{x}_{1} & \mathbf{x}_{2} \mathbf{x}_{2}\end{array}\right ) +\left (\begin{array}{cc}\mathbf{x}_{1} \mathbf{x}_{1} & \mathbf{x}_{1} \mathbf{x}_{2} \\
\mathbf{x}_{2} \mathbf{x}_{1} & \mathbf{x}_{2} \mathbf{x}_{2}\end{array}\right )^{t}\right ) \left (\begin{array}{cc}\frac{1}{\sqrt{2}} \mathbb{I}_{r} & \frac{1}{\sqrt{2}} \mathbb{I}_{r} \\
 -\frac{i}{\sqrt{2}} \mathbb{I}_{r} & \frac{i}{\sqrt{2}} \mathbb{I}_{r}\end{array}\right ) \nonumber  \\
 &  =\left (\begin{array}{cc}\frac{1}{\sqrt{2}} \mathbb{I}_{r} &  -\frac{i}{\sqrt{2}} \mathbb{I}_{r} \\
\frac{1}{\sqrt{2}} \mathbb{I}_{r} & \frac{i}{\sqrt{2}} \mathbb{I}_{r}\end{array}\right ) \left (\begin{array}{cc}\frac{1}{\sqrt{2}} \mathbb{I}_{r} & \frac{1}{\sqrt{2}} \mathbb{I}_{r} \\
 -\frac{i}{\sqrt{2}} \mathbb{I}_{r} & \frac{i}{\sqrt{2}} \mathbb{I}_{r}\end{array}\right ) =\left (\begin{array}{cc}\mathbb{O} & \mathbb{I}_{r} \\
\mathbb{I}_{r} & \mathbb{O}\end{array}\right )\text{.}\end{align}As
\begin{equation}\left (\begin{array}{cc}\left \{\mathbf{a}_{1} ,\mathbf{a}_{1}\right \} & \left \{\mathbf{a}_{1} ,\mathbf{a}_{2}\right \} \\
\left \{\mathbf{a}_{2} ,\mathbf{a}_{1}\right \} & \left \{\mathbf{a}_{2} ,\mathbf{a}_{2}\right \}\end{array}\right ) =\left (\mathbb{K}^{ \dag }\right )^{t} \left (\left (\begin{array}{cc}\mathbf{x}_{1} \mathbf{x}_{1} & \mathbf{x}_{1} \mathbf{x}_{2} \\
\mathbf{x}_{2} \mathbf{x}_{1} & \mathbf{x}_{2} \mathbf{x}_{2}\end{array}\right ) +\left (\begin{array}{cc}\mathbf{x}_{1} \mathbf{x}_{1} & \mathbf{x}_{1} \mathbf{x}_{2} \\
\mathbf{x}_{2} \mathbf{x}_{1} & \mathbf{x}_{2} \mathbf{x}_{2}\end{array}\right )^{t}\right ) \mathbb{K}^{ \dag }
\end{equation}a transformed $\left \{\mathbf{x}_{1} ,\mathbf{x}_{2}\right \}$ basis gives
\begin{align} & \left (\mathbb{K}^{ \dag }\right )^{t} \mathbb{T}^{t} \left (\left (\begin{array}{cc}\mathbf{x}_{1} \mathbf{x}_{1} & \mathbf{x}_{1} \mathbf{x}_{2} \\
\mathbf{x}_{2} \mathbf{x}_{1} & \mathbf{x}_{2} \mathbf{x}_{2}\end{array}\right ) +\left (\begin{array}{cc}\mathbf{x}_{1} \mathbf{x}_{1} & \mathbf{x}_{1} \mathbf{x}_{2} \\
\mathbf{x}_{2} \mathbf{x}_{1} & \mathbf{x}_{2} \mathbf{x}_{2}\end{array}\right )^{t}\right ) \mathbb{T} \mathbb{K}^{ \dag } \\
 &  =\left (\mathbb{K}^{ \dag }\right )^{t} \mathbb{T}^{t} \left (\mathbb{K}^{ \dag } \mathbb{K}\right )^{t} \left (\left (\begin{array}{cc}\mathbf{x}_{1} \mathbf{x}_{1} & \mathbf{x}_{1} \mathbf{x}_{2} \\
\mathbf{x}_{2} \mathbf{x}_{1} & \mathbf{x}_{2} \mathbf{x}_{2}\end{array}\right ) +\left (\begin{array}{cc}\mathbf{x}_{1} \mathbf{x}_{1} & \mathbf{x}_{1} \mathbf{x}_{2} \\
\mathbf{x}_{2} \mathbf{x}_{1} & \mathbf{x}_{2} \mathbf{x}_{2}\end{array}\right )^{t}\right ) \mathbb{K}^{ \dag } \mathbb{K} \mathbb{T} \mathbb{K}^{ \dag } \\
 &  =\left (\mathbb{K}^{ \dag }\right )^{t} \mathbb{T}^{t} \mathbb{K}^{t} \left (\mathbb{K}^{ \dag }\right )^{t} \left (\left (\begin{array}{cc}\mathbf{x}_{1} \mathbf{x}_{1} & \mathbf{x}_{1} \mathbf{x}_{2} \\
\mathbf{x}_{2} \mathbf{x}_{1} & \mathbf{x}_{2} \mathbf{x}_{2}\end{array}\right ) +\left (\begin{array}{cc}\mathbf{x}_{1} \mathbf{x}_{1} & \mathbf{x}_{1} \mathbf{x}_{2} \\
\mathbf{x}_{2} \mathbf{x}_{1} & \mathbf{x}_{2} \mathbf{x}_{2}\end{array}\right )^{t}\right ) \mathbb{K}^{ \dag } \mathbb{K} \mathbb{T} \mathbb{K}^{ \dag } \\
 &  =\left (\mathbb{K} \mathbb{T} \mathbb{K}^{ \dag }\right )^{t} \left (\mathbb{K}^{ \dag }\right )^{t} \left (\left (\begin{array}{cc}\mathbf{x}_{1} \mathbf{x}_{1} & \mathbf{x}_{1} \mathbf{x}_{2} \\
\mathbf{x}_{2} \mathbf{x}_{1} & \mathbf{x}_{2} \mathbf{x}_{2}\end{array}\right ) +\left (\begin{array}{cc}\mathbf{x}_{1} \mathbf{x}_{1} & \mathbf{x}_{1} \mathbf{x}_{2} \\
\mathbf{x}_{2} \mathbf{x}_{1} & \mathbf{x}_{2} \mathbf{x}_{2}\end{array}\right )^{t}\right ) \mathbb{K}^{ \dag } \mathbb{K} \mathbb{T} \mathbb{K}^{ \dag } \\
 &  =\mathbb{S}^{t} \left (\mathbb{K}^{ \dag }\right )^{t} \left (\left (\begin{array}{cc}\mathbf{x}_{1} \mathbf{x}_{1} & \mathbf{x}_{1} \mathbf{x}_{2} \\
\mathbf{x}_{2} \mathbf{x}_{1} & \mathbf{x}_{2} \mathbf{x}_{2}\end{array}\right ) +\left (\begin{array}{cc}\mathbf{x}_{1} \mathbf{x}_{1} & \mathbf{x}_{1} \mathbf{x}_{2} \\
\mathbf{x}_{2} \mathbf{x}_{1} & \mathbf{x}_{2} \mathbf{x}_{2}\end{array}\right )^{t}\right ) \mathbb{K}^{ \dag } \mathbb{S} =\mathbb{S}^{t} \left (\begin{array}{cc}\left \{\mathbf{a}_{1} ,\mathbf{a}_{1}\right \} & \left \{\mathbf{a}_{1} ,\mathbf{a}_{2}\right \} \\
\left \{\mathbf{a}_{2} ,\mathbf{a}_{1}\right \} & \left \{\mathbf{a}_{2} ,\mathbf{a}_{2}\right \}\end{array}\right ) \mathbb{S}\text{,}\end{align}where
\begin{equation}\mathbb{S} =\mathbb{K} \mathbb{T} \mathbb{K}^{ \dag }\text{,}
\end{equation}i.e.
\begin{align}\left (\begin{array}{cc}\left \{\mathbf{y}_{1} ,\mathbf{y}_{1}\right \} & \left \{\mathbf{y}_{1} ,\mathbf{y}_{2}\right \} \\
\left \{\mathbf{y}_{2} ,\mathbf{y}_{1}\right \} & \left \{\mathbf{y}_{2} ,\mathbf{y}_{2}\right \}\end{array}\right ) &  =\mathbb{T}^{t} \left (\left (\begin{array}{cc}\mathbf{x}_{1} \mathbf{x}_{1} & \mathbf{x}_{1} \mathbf{x}_{2} \\
\mathbf{x}_{2} \mathbf{x}_{1} & \mathbf{x}_{2} \mathbf{x}_{2}\end{array}\right ) +\left (\begin{array}{cc}\mathbf{x}_{1} \mathbf{x}_{1} & \mathbf{x}_{1} \mathbf{x}_{2} \\
\mathbf{x}_{2} \mathbf{x}_{1} & \mathbf{x}_{2} \mathbf{x}_{2}\end{array}\right )^{t}\right ) \mathbb{T} \nonumber  \\
 &  =\mathbb{T}^{t} \left (\left (\begin{array}{cc}\mathbb{I}_{r} & \mathbb{O}_{r} \\
\mathbb{O}_{r} & \mathbb{I}_{r}\end{array}\right )\right ) \mathbb{T} =\left (\begin{array}{cc}\mathbb{I}_{r} & \mathbb{O}_{r} \\
\mathbb{O}_{r} & \mathbb{I}_{r}\end{array}\right )\end{align}and
\begin{align}\left (\begin{array}{cc}\left \{\mathbf{b}_{1} ,\mathbf{b}_{1}\right \} & \left \{\mathbf{b}_{1} ,\mathbf{b}_{2}\right \} \\
\left \{\mathbf{b}_{2} ,\mathbf{b}_{1}\right \} & \left \{\mathbf{b}_{2} ,\mathbf{b}_{2}\right \}\end{array}\right ) &  =\mathbb{S}^{t} \left (\left (\begin{array}{cc}\mathbf{a}_{1} \mathbf{a}_{1} & \mathbf{a}_{1} \mathbf{a}_{2} \\
\mathbf{a}_{2} \mathbf{a}_{1} & \mathbf{a}_{2} \mathbf{a}_{2}\end{array}\right ) +\left (\begin{array}{cc}\mathbf{a}_{1} \mathbf{a}_{1} & \mathbf{a}_{1} \mathbf{a}_{2} \\
\mathbf{a}_{2} \mathbf{a}_{1} & \mathbf{a}_{2} \mathbf{a}_{2}\end{array}\right )^{t}\right ) \mathbb{S} \nonumber  \\
 &  =\mathbb{S}^{t} \left (\left (\begin{array}{cc}\mathbb{O}_{r} & \mathbb{I}_{r} \\
\mathbb{I}_{r} & \mathbb{O}_{r}\end{array}\right )\right ) \mathbb{S} =\left (\begin{array}{cc}\mathbb{O}_{r} & \mathbb{I}_{r} \\
\mathbb{I}_{r} & \mathbb{O}_{r}\end{array}\right )\text{.}\end{align}If $S$ is \emph{also} canonical then
\begin{align}\mathbb{S}^{ \dag } \mathbb{S} &  =\left (\mathbb{K} \mathbb{T} \mathbb{K}^{ \dag }\right )^{ \dag } \mathbb{K} \mathbb{T} \mathbb{K}^{ \dag } =\mathbb{K} \mathbb{T}^{ \dag } \mathbb{K}^{ \dag } \mathbb{K} \mathbb{T} \mathbb{K}^{ \dag } =\mathbb{K} \mathbb{T}^{ \dag } \mathbb{T} \mathbb{K}^{ \dag } =\mathbb{I}_{2 r} \nonumber  \\
 &  \Leftrightarrow \mathbb{T}^{ \dag } \mathbb{T} =\mathbb{I}_{2 r}\end{align}and $T \in $ $O \left (2 r ,\mathbb{R}\right ) =O \left (2 r ,\mathbb{C}\right ) \cap U \left (2 r\right )$ is canonical and vice-versa. 

One can further check to see that
\begin{align}\mathbb{S}^{t} \mathbb{S} &  =\left (\mathbb{K} \mathbb{T} \mathbb{K}^{ \dag }\right )^{t} \mathbb{K} \mathbb{T} \mathbb{K}^{ \dag } =\overline{\mathbb{K}} \mathbb{T}^{t} \mathbb{K}^{t} \mathbb{K} \mathbb{T} \mathbb{K}^{ \dag } =\left (\mathbb{K}^{ \dag }\right )^{t} \mathbb{T}^{t} \left (\begin{array}{cc}\mathbb{O} & \mathbb{I}_{r} \\
\mathbb{I}_{r} & \mathbb{O}\end{array}\right ) \mathbb{T} \mathbb{K}^{ \dag } \nonumber  \\
 &  =\left (\mathbb{K}^{ \dag }\right )^{t} \left (\begin{array}{cc}\mathbb{I}_{r} & \mathbb{O}_{r} \\
\mathbb{O}_{r} & \mathbb{I}_{r}\end{array}\right ) \mathbb{K}^{ \dag } =\left (\mathbb{K}^{ \dag }\right )^{t} \mathbb{K}^{ \dag } \Longleftrightarrow \mathbb{K} \mathbb{K}^{t} =\mathbb{I}_{2 r}\end{align}which confirms $S \in $ $O \left (2 r ,\mathbb{C}\right )\text{.}$ 

\emph{NOTE} that the definition of $K$ involves the complex unit "$i ""$ in an essential way. 

\section{Dressed (Quasi-Free) Vacuum Vectors ($O \left (2 r ,\mathbb{R}\right )/U \left (r\right )$-Coherent States)}
The action of the orthogonal group $O \left (2 r ,\mathbb{R}\right )$ on $\mathcal{F}$ is given by cannonical transformation, which are given by the adjoint representation of $O \left (2 r ,\mathbb{R}\right )$ on $\mathcal{F}_{1}$ leading to
\begin{equation}\widehat{u} \left (g\right ) \left (A\right ) =u \left (g\right ) A u \left (g\right )^{ \dag }\text{,}
\end{equation}where
\begin{equation}\widehat{u} \left (g\right ) :\mathcal{F}_{n} \rightarrow \mathcal{F}_{n} ;0 \leq n \leq 2 r ,g \in O \left (2 r ,\mathbb{R}\right )\text{.}
\end{equation}The subgroup $U \left (r\right )$ is a maximal compact subgroup of $O \left (2 r ,\mathbb{R}\right )$ given by
\begin{equation}u \left (\mathbf{\lambda }\right ) =\exp  \left \{i \sum _{1 \leq j ,k \leq r}\lambda _{j k} a_{j}^{ \dag } a_{k}\right \} ;\,\lambda _{j k} =\,\overline{\lambda }_{k j} ,1 \leq j ,k \leq r
\end{equation}and one can consider the coset $O \left (2 r ,\mathbb{R}\right )/U \left (r\right )$ and thus the $O \left (2 r ,\mathbb{R}\right )/U \left (r\right )$-Coherent States
\begin{equation}\vert \phi  \left (\gamma \right ) \rangle  =\mathcal{N} \left (\gamma \right ) \exp  \left \{\sum _{1 \leq j <k \leq r}\left (\gamma _{j k} a_{j}^{ \dag } a_{k}^{ \dag } -\overline{\gamma }_{j k} a_{k} a_{j}\right )\right \}\vert \phi _{a} \rangle \text{,}
\end{equation}where
\begin{equation}u \left (g\right )\vert \phi _{a} \rangle  =e^{i \theta  \left (g\right )}\vert \phi _{a} \rangle  , \forall g \in O \left (2 r ,\mathbb{R}\right ) ,\theta  \left (g\right ) \in \mathbb{R}\text{.}
\end{equation}$\lceil $The exponent of the coset is anti self adjoint
\begin{equation}\left (\gamma _{j k} a_{j}^{ \dag } a_{k}^{ \dag } -\overline{\gamma }_{j k} a_{k} a_{j}\right )^{ \dag } =\overline{\gamma }_{j k} a_{k} a_{j} -\gamma _{j k} a_{j}^{ \dag } a_{k}^{ \dag } = -\left (\gamma _{j k} a_{j}^{ \dag } a_{k}^{ \dag } -\overline{\gamma }_{j k} a_{k} a_{j}\right )
\end{equation}as it should be in order to produce a unitary transformation belonging to $U \left (2 r\right )$ and can be expressed as
\begin{align} & \sum _{1 \leq j <k \leq r}\left (\gamma _{j k} a_{j}^{ \dag } a_{k}^{ \dag } -\overline{\gamma }_{j k} a_{k} a_{j}\right ) =\sum _{1 \leq j <k \leq r}\left (\gamma _{j k} a_{j}^{ \dag } a_{k}^{ \dag } +\overline{\gamma }_{j k} a_{j} a_{k}\right ) =\frac{1}{2} \left (\begin{array}{cc}\mathbf{a}^{ \dag } & \mathbf{a}\end{array}\right ) \left (\begin{array}{cc}\mathbb{O} & \mathbf{\gamma } \\
\overline{\mathbf{\gamma }} & \mathbb{O}\end{array}\right ) \left (\begin{array}{c}\mathbf{a} \\
\mathbf{a}^{ \dag }\end{array}\right ) \\
 &  =\frac{1}{2} \left (\begin{array}{cc}\mathbf{a}^{ \dag } & \mathbf{a}\end{array}\right ) \left (\begin{array}{c}\mathbf{\gamma } \mathbf{a}^{ \dag } \\
\overline{\mathbf{\gamma }} \mathbf{a}\end{array}\right ) =\frac{1}{2} \left (\mathbf{a}^{ \dag } \mathbf{\gamma } \mathbf{a}^{ \dag } +\mathbf{a} \overline{\mathbf{\gamma }} \mathbf{a}\right )\text{,}\end{align}where $\gamma $ is a $r \times r$ complex antisymmetric matrix i.e.
\begin{equation}\mathbf{\gamma }  = -\mathbf{\gamma }^{t}
\end{equation}so that
\begin{align}\mathbf{\gamma } &  =\mathbf{\gamma }_{r} +i \mathbf{\gamma }_{i} \\
\mathbf{\gamma }^{t} &  =\mathbf{\gamma }_{r}^{t} +i \mathbf{\gamma }_{i}^{t} = -\left (\mathbf{\gamma }_{r} +i \mathbf{\gamma }_{i}\right ) \\
\mathbf{\gamma }^{ \dag } &  =\mathbf{\gamma }_{r}^{t} -i \mathbf{\gamma }_{i}^{t} = -\mathbf{\gamma }_{r} +i \mathbf{\gamma }_{i} = -\overline{\mathbf{\gamma }}\text{.}\end{align}$\lceil $An $r =2$ example is given by
\begin{align} & \frac{1}{2} \left (\begin{array}{cccc}a_{j}^{ \dag } & a_{k}^{ \dag } & a_{j} & a_{k}\end{array}\right ) \left (\begin{array}{cccc}0 & 0 & 0 & \gamma _{j k} \\
0 & 0 &  -\gamma _{j k} & 0 \\
0 & \overline{\gamma }_{j k} & 0 & 0 \\
 -\overline{\gamma }_{j k} & 0 & 0 & 0\end{array}\right ) \left (\begin{array}{c}a_{j} \\
a_{k} \\
a_{j}^{ \dag } \\
a_{k}^{ \dag }\end{array}\right ) =\frac{1}{2} \left (\begin{array}{cccc}a_{j}^{ \dag } & a_{k}^{ \dag } & a_{j} & a_{k}\end{array}\right ) \left (\begin{array}{c}\gamma _{j k} a_{k}^{ \dag } \\
 -\gamma _{j k} a_{j}^{ \dag } \\
\overline{\gamma }_{j k} a_{k} \\
 -\overline{\gamma }_{j k} a_{j}\end{array}\right ) \\
 &  =\frac{1}{2} \left (\gamma _{j k} a_{j}^{ \dag } a_{k}^{ \dag } -\gamma _{j k} a_{k}^{ \dag } a_{j}^{ \dag } +\overline{\gamma }_{j k} a_{j} a_{k} -\overline{\gamma }_{j k} a_{k} a_{j}\right ) =\gamma _{j k} a_{j}^{ \dag } a_{k}^{ \dag } -\overline{\gamma }_{j k} a_{k} a_{j} =\gamma _{j k} a_{j}^{ \dag } a_{k}^{ \dag } +\overline{\gamma }_{j k} a_{j} a_{k}\end{align}$\rfloor $ 

The real form of $\left (\begin{array}{cc}\mathbb{O} & \mathbf{\gamma } \\
\overline{\mathbf{\gamma }} & \mathbb{O}\end{array}\right )$ is given by
\begin{align}\mathbb{K} \left (\begin{array}{cc}\mathbb{O} & \mathbf{\gamma } \\
\overline{\mathbf{\gamma }} & \mathbb{O}\end{array}\right ) \mathbb{K}^{ \dag } &  =\left (\begin{array}{cc}\frac{1}{\sqrt{2}} \mathbb{I}_{r} & \frac{1}{\sqrt{2}} \mathbb{I}_{r} \\
 -\frac{i}{\sqrt{2}} \mathbb{I}_{r} & \frac{i}{\sqrt{2}} \mathbb{I}_{r}\end{array}\right ) \left (\begin{array}{cc}\mathbb{O} & \mathbf{\gamma } \\
\overline{\mathbf{\gamma }} & \mathbb{O}\end{array}\right ) \left (\begin{array}{cc}\frac{1}{\sqrt{2}} \mathbb{I}_{r} & \frac{i}{\sqrt{2}} \mathbb{I}_{r} \\
\frac{1}{\sqrt{2}} \mathbb{I}_{r} & \frac{ -i}{\sqrt{2}} \mathbb{I}_{r}\end{array}\right ) \\
 &  =\left (\begin{array}{cc}\ensuremath{\operatorname*{Re}} \left (\overline{\mathbf{\gamma }}\right ) &  -\ensuremath{\operatorname*{Im}} \left (\overline{\mathbf{\gamma }}\right ) \\
 -\ensuremath{\operatorname*{Im}} \left (\overline{\mathbf{\gamma }}\right ) &  -\ensuremath{\operatorname*{Re}} \left (\overline{\mathbf{\gamma }}\right )\end{array}\right ) =\left (\begin{array}{cc}\ensuremath{\operatorname*{Re}} \left (\mathbf{\gamma }\right ) & \ensuremath{\operatorname*{Im}} \left (\mathbf{\gamma }\right ) \\
\ensuremath{\operatorname*{Im}} \left (\mathbf{\gamma }\right ) &  -\ensuremath{\operatorname*{Re}} \left (\mathbf{\gamma }\right )\end{array}\right )\text{,}\end{align}so that
\begin{align} & \frac{1}{2} \left (\begin{array}{cc}\mathbf{a}^{ \dag } & \mathbf{a}\end{array}\right ) \mathcal{R}_{a} \left (\left (\begin{array}{cc}\mathbb{O} & \mathbf{\gamma } \\
\overline{\mathbf{\gamma }} & \mathbb{O}\end{array}\right )\right ) \left (\begin{array}{c}\mathbf{a} \\
\mathbf{a}^{ \dag }\end{array}\right ) =\frac{1}{2} \left (\begin{array}{cc}\mathbf{x}_{1} & \mathbf{x}_{2}\end{array}\right ) \mathcal{R}_{x} \left (\left (\begin{array}{cc}\mathbb{O} & \mathbf{\gamma } \\
\overline{\mathbf{\gamma }} & \mathbb{O}\end{array}\right )\right ) \left (\begin{array}{c}\mathbf{x}_{1} \\
\mathbf{x}_{2}\end{array}\right ) \\
 &  =\frac{1}{2} \left (\begin{array}{cc}\mathbf{x}_{1} & \mathbf{x}_{2}\end{array}\right ) \mathbb{K} \left (\begin{array}{cc}\mathbb{O} & \mathbf{\gamma } \\
\overline{\mathbf{\gamma }} & \mathbb{O}\end{array}\right ) \mathbb{K}^{ \dag } \left (\begin{array}{c}\mathbf{x}_{1} \\
\mathbf{x}_{2}\end{array}\right ) =\frac{1}{2} \left (\begin{array}{cc}\mathbf{x}_{1} & \mathbf{x}_{2}\end{array}\right ) \left (\begin{array}{cc}\ensuremath{\operatorname*{Re}} \left (\overline{\mathbf{\gamma }}\right ) &  -\ensuremath{\operatorname*{Im}} \left (\overline{\mathbf{\gamma }}\right ) \\
 -\ensuremath{\operatorname*{Im}} \left (\overline{\mathbf{\gamma }}\right ) &  -\ensuremath{\operatorname*{Re}} \left (\overline{\mathbf{\gamma }}\right )\end{array}\right ) \left (\begin{array}{c}\mathbf{x}_{1} \\
\mathbf{x}_{2}\end{array}\right ) \\
 &  =\frac{1}{2} \left (\begin{array}{cc}\mathbf{x}_{1} & \mathbf{x}_{2}\end{array}\right ) \left (\begin{array}{c}\ensuremath{\operatorname*{Re}} \left (\overline{\mathbf{\gamma }}\right ) \mathbf{x}_{1} -\ensuremath{\operatorname*{Im}} \left (\overline{\mathbf{\gamma }}\right ) \mathbf{x}_{2} \\
 -\ensuremath{\operatorname*{Im}} \left (\overline{\mathbf{\gamma }}\right ) \mathbf{x}_{1} -\ensuremath{\operatorname*{Re}} \left (\overline{\mathbf{\gamma }}\right ) \mathbf{x}_{2}\end{array}\right ) =\frac{1}{2} \left (\mathbf{x}_{1} \ensuremath{\operatorname*{Re}} \left (\overline{\mathbf{\gamma }}\right ) \mathbf{x}_{1} -\mathbf{x}_{1} \ensuremath{\operatorname*{Im}} \left (\overline{\mathbf{\gamma }}\right ) \mathbf{x}_{2} -\mathbf{x}_{2} \ensuremath{\operatorname*{Im}} \left (\overline{\mathbf{\gamma }}\right ) \mathbf{x}_{1} -\mathbf{x}_{2} \ensuremath{\operatorname*{Re}} \left (\overline{\mathbf{\gamma }}\right ) \mathbf{x}_{2}\right )\end{align}$\lceil $A $r =2$ example is given by
\begin{align} & \left (\gamma _{j k} a_{j}^{ \dag } a_{k}^{ \dag } -\overline{\gamma }_{j k} a_{k} a_{j}\right )^{ \dag } =\overline{\gamma }_{j k} a_{k} a_{j} -\gamma _{j k} a_{j}^{ \dag } a_{k}^{ \dag } = -\left (\gamma _{j k} a_{j}^{ \dag } a_{k}^{ \dag } -\overline{\gamma }_{j k} a_{k} a_{j}\right ) \\
 & \left (\gamma _{j k} a_{j}^{ \dag } a_{k}^{ \dag } -\overline{\gamma }_{j k} a_{k} a_{j}\right ) =\frac{1}{2} \left (\gamma _{j k} \left (x_{j} -i x_{j +r}\right ) \left (x_{k} -i x_{k +r}\right ) -\overline{\gamma }_{j k} \left (x_{k} +i x_{k +r}\right ) \left (x_{j} +i x_{j +r}\right )\right ) \\
 &  =\frac{1}{2} \left (\gamma _{j k} x_{j} x_{k} -\gamma _{j k} x_{j +r} x_{k +r} -i \gamma _{j k} x_{j +r} x_{k} -i \gamma _{j k} x_{j} x_{k +r} -\overline{\gamma }_{j k} x_{k} x_{j} +\overline{\gamma }_{j k} x_{k +r} x_{j +r} -i \overline{\gamma }_{j k} x_{k} x_{j +r} -i \overline{\gamma }_{j k} x_{k +r} x_{j}\right ) \\
 &  =\frac{1}{2} \left (\gamma _{j k} +\overline{\gamma }_{j k}\right ) x_{j} x_{k} -\left (\overline{\gamma }_{j k} +\gamma _{j k}\right ) x_{j +r} x_{k +r} -i \left (\gamma _{j k} -\overline{\gamma }_{j k}\right ) x_{j +r} x_{k} -i \left (\gamma _{j k} -\overline{\gamma }_{j k}\right ) x_{j} x_{k +r} \\
 &  =\left (\ensuremath{\operatorname*{Re}} \gamma _{j k}\right ) x_{j} x_{k} -\left (\ensuremath{\operatorname*{Re}} \gamma _{j k}\right ) x_{j +r} x_{k +r} +\left (\ensuremath{\operatorname*{Im}} \gamma _{j k}\right ) x_{j +r} x_{k} +\left (\ensuremath{\operatorname*{Im}} \gamma _{j k}\right ) x_{j} x_{k +r}\end{align}$\rfloor $ 

and
\begin{align}\exp  \left \{\frac{1}{2} \left (\mathbf{a}^{ \dag } \mathbf{\gamma } \mathbf{a}^{ \dag } +\mathbf{a} \overline{\mathbf{\gamma }} \mathbf{a}\right )\right \} &  =\exp  \left \{\frac{1}{2} \left (\begin{array}{cc}\mathbf{a}^{ \dag } & \mathbf{a}\end{array}\right ) \left (\begin{array}{cc}\mathbb{O} & \mathbf{\gamma } \\
\overline{\mathbf{\gamma }} & \mathbb{O}\end{array}\right ) \left (\begin{array}{c}\mathbf{a} \\
\mathbf{a}^{ \dag }\end{array}\right )\right \} \\
 &  =\exp  \left \{\frac{1}{2} \left (\mathbf{x}_{1} \ensuremath{\operatorname*{Re}} \left (\mathbf{\gamma }\right ) \mathbf{x}_{1} +\mathbf{x}_{1} \ensuremath{\operatorname*{Im}} \left (\mathbf{\gamma }\right ) \mathbf{x}_{2} +\mathbf{x}_{2} \ensuremath{\operatorname*{Im}} \left (\mathbf{\gamma }\right ) \mathbf{x}_{1} -\mathbf{x}_{2} \ensuremath{\operatorname*{Re}} \left (\mathbf{\gamma }\right ) \mathbf{x}_{2}\right )\right \} \\
 &  =\exp  \left \{\frac{1}{2} \left (\begin{array}{cc}\mathbf{x}_{1} & \mathbf{x}_{2}\end{array}\right ) \left (\begin{array}{cc}\ensuremath{\operatorname*{Re}} \left (\mathbf{\gamma }\right ) & \ensuremath{\operatorname*{Im}} \left (\mathbf{\gamma }\right ) \\
\ensuremath{\operatorname*{Im}} \left (\mathbf{\gamma }\right ) &  -\ensuremath{\operatorname*{Re}} \left (\mathbf{\gamma }\right )\end{array}\right ) \left (\begin{array}{c}\mathbf{x}_{1} \\
\mathbf{x}_{2}\end{array}\right )\right \}\text{,}\end{align}after one has observed that
\begin{align}\ensuremath{\operatorname*{Re}} \left (\mathbf{\gamma }\right ) &  =\ensuremath{\operatorname*{Re}} \left (\overline{\mathbf{\gamma }}\right ) \\
\ensuremath{\operatorname*{Im}} \left (\mathbf{\gamma }\right ) &  = -\ensuremath{\operatorname*{Im}} \left (\overline{\mathbf{\gamma }}\right )\text{.}\end{align}Note that the matrix of coefficients of the self adjoint operators are \emph{real antisymmetric}
\begin{equation}\left (\begin{array}{cc}\ensuremath{\operatorname*{Re}} \left (\mathbf{\gamma }\right ) & \ensuremath{\operatorname*{Im}} \left (\mathbf{\gamma }\right ) \\
\ensuremath{\operatorname*{Im}} \left (\mathbf{\gamma }\right ) &  -\ensuremath{\operatorname*{Re}} \left (\mathbf{\gamma }\right )\end{array}\right ) = -\left (\begin{array}{cc}\ensuremath{\operatorname*{Re}} \left (\mathbf{\gamma }\right ) & \ensuremath{\operatorname*{Im}} \left (\mathbf{\gamma }\right ) \\
\ensuremath{\operatorname*{Im}} \left (\mathbf{\gamma }\right ) &  -\ensuremath{\operatorname*{Re}} \left (\mathbf{\gamma }\right )\end{array}\right )^{t}\text{.}
\end{equation}Complex (and thus real) antisymmetric matrices have the following properties:
\begin{align}\mathbb{T} \mathbf{\gamma } \mathbb{T}^{t} &  =\left (\begin{array}{cc}\mathbf{0} & \mathbf{\zeta } \\
 -\mathbf{\zeta } & \mathbf{0}\end{array}\right ) \\
\mathbf{\gamma } &  =\mathbb{T}^{ \dag } \left (\begin{array}{cc}\mathbf{0} & \mathbf{\zeta } \\
 -\mathbf{\zeta } & \mathbf{0}\end{array}\right ) \overline{\mathbb{T}} \\
\mathbf{\gamma } \mathbf{\gamma }^{ \dag } &  =\mathbb{T}^{ \dag } \left (\begin{array}{cc}\mathbf{0} & \mathbf{\zeta } \\
 -\mathbf{\zeta } & \mathbf{0}\end{array}\right ) \overline{\mathbb{T}} \left (\mathbb{T}^{ \dag } \left (\begin{array}{cc}\mathbf{0} & \mathbf{\zeta } \\
 -\mathbf{\zeta } & \mathbf{0}\end{array}\right ) \overline{\mathbb{T}}\right )^{ \dag } =\mathbb{T}^{ \dag } \left (\begin{array}{cc}\mathbf{0} & \mathbf{\zeta } \\
 -\mathbf{\zeta } & \mathbf{0}\end{array}\right ) \overline{\mathbb{T}} \mathbb{T}^{t} \left (\begin{array}{cc}\mathbf{0} & \mathbf{\zeta } \\
 -\mathbf{\zeta } & \mathbf{0}\end{array}\right )^{ \dag } \mathbb{T}^{ \dag } \\
 &  =\mathbb{T}^{ \dag } \left (\begin{array}{cc}\mathbf{0} & \mathbf{\zeta } \\
 -\mathbf{\zeta } & \mathbf{0}\end{array}\right ) \left (\begin{array}{cc}\mathbf{0} & \mathbf{\zeta } \\
 -\mathbf{\zeta } & \mathbf{0}\end{array}\right )^{ \dag } \mathbb{T} =\mathbb{T}^{ \dag } \left (\begin{array}{cc}\mathbf{0} & \mathbf{\zeta } \\
 -\mathbf{\zeta } & \mathbf{0}\end{array}\right ) \left (\begin{array}{cc}\mathbf{0} &  -\overline{\mathbf{\zeta }} \\
\overline{\mathbf{\zeta }} & \mathbf{0}\end{array}\right )^{ \dag } \mathbb{T} =\mathbb{T}^{ \dag } \left (\begin{array}{cc}\left \vert \mathbf{\zeta }\right \vert ^{2} & \mathbf{0} \\
\mathbf{0} & \left \vert \mathbf{\zeta }\right \vert ^{2}\end{array}\right ) \mathbb{T} \\
\mathbf{\gamma }^{ \dag } \mathbf{\gamma } &  =\left (\mathbb{T}^{ \dag } \left (\begin{array}{cc}\mathbf{0} & \mathbf{\zeta } \\
 -\mathbf{\zeta } & \mathbf{0}\end{array}\right ) \overline{\mathbb{T}}\right )^{ \dag } \mathbb{T}^{ \dag } \left (\begin{array}{cc}\mathbf{0} & \mathbf{\zeta } \\
 -\mathbf{\zeta } & \mathbf{0}\end{array}\right ) \overline{\mathbb{T}} =\mathbb{T}^{t} \left (\begin{array}{cc}\mathbf{0} &  -\overline{\mathbf{\zeta }} \\
\overline{\mathbf{\zeta }} & \mathbf{0}\end{array}\right ) \mathbb{T} \mathbb{T}^{ \dag } \left (\begin{array}{cc}\mathbf{0} & \mathbf{\zeta } \\
 -\mathbf{\zeta } & \mathbf{0}\end{array}\right ) \overline{\mathbb{T}} \\
 &  =\mathbb{T}^{t} \left (\begin{array}{cc}\left \vert \mathbf{\zeta }\right \vert ^{2} & \mathbf{0} \\
\mathbf{0} & \left \vert \mathbf{\zeta }\right \vert ^{2}\end{array}\right ) \overline{\mathbb{T}} \Rightarrow \mathbf{\gamma }^{ \dag } \mathbf{\gamma } =\overline{\mathbf{\gamma } \mathbf{\gamma }^{ \dag }}\text{,}\end{align}where $\mathbf{\zeta }$ is a real diagonal $s \times s$ real matrix i.e.
\begin{equation}\mathbf{\zeta } =\left (\begin{array}{cccc}\zeta _{1} & 0 & \cdots  & 0 \\
0 & \ddots  & \adots  & \vdots  \\
\vdots  & \adots  & \ddots  & 0 \\
0 & \cdots  & 0 & \zeta _{s}\end{array}\right ) ;\zeta _{j} \geq 0 \in \mathbb{R} ,1 \leq j \leq s\text{.}
\end{equation}One can always transform to the canonical basis of $\mathbf{\gamma }$ that produces the skew block diagonal form of $\mathbf{\gamma }$, and obtain
\begin{equation*}\gamma  = \sum _{1 \leq j \leq s}\zeta _{j} \left (c_{j}^{ \dag } c_{j +s}^{ \dag } +c_{j} c_{j +s}\right ) = \sum _{1 \leq j \leq s}\zeta _{j} \left (y_{j} y_{j +s} -y_{j +r} y_{j +s +r}\right )
\end{equation*}$\lceil $
\begin{align}\gamma  &  = \sum _{1 \leq j \leq s}\zeta _{j} \left (c_{j}^{ \dag } c_{j +s}^{ \dag } +c_{j} c_{j +s}\right ) =\frac{1}{2}  \sum _{1 \leq j \leq s}\zeta _{j} \left (\left (y_{j} -i y_{j +r}\right ) \left (y_{j +s} -i y_{j +s +r}\right ) +\left (y_{j} +i y_{j +r}\right ) \left (y_{j +s} +i y_{j +s +r}\right )\right ) \\
 &  =\frac{1}{2}  \sum _{1 \leq j \leq s}\zeta _{j} \left (\left (y_{j} \left (y_{j +s} -i y_{j +s +r}\right ) -i y_{j +r} \left (y_{j +s} -i y_{j +s +r}\right )\right ) +\left (y_{j} \left (y_{j +s} +i y_{j +s +r}\right ) +i y_{j +r} \left (y_{j +s} +i y_{j +s +r}\right )\right )\right ) \\
 &  =\frac{1}{2}  \sum _{1 \leq j \leq s}\zeta _{j} \left (\left (\left (y_{j} y_{j +s} -i y_{j} y_{j +s +r}\right ) -i \left (y_{j +r} y_{j +s} -i y_{j +r} y_{j +s +r}\right )\right ) +\left (\left (y_{j} y_{j +s} +i y_{j} y_{j +s +r}\right ) +i \left (y_{j +r} y_{j +s} +i y_{j +r} y_{j +s +r}\right )\right )\right ) \\
 &  =\frac{1}{2}  \sum _{1 \leq j \leq s}\zeta _{j} \left (y_{j} y_{j +s} -i y_{j} y_{j +s +r} -i y_{j +r} y_{j +s} -y_{j +r} y_{j +s +r} +y_{j} y_{j +s} +i y_{j} y_{j +s +r} +i y_{j +r} y_{j +s} -y_{j +r} y_{j +s +r}\right ) \\
 &  = \sum _{1 \leq j \leq s}\zeta _{j} \left (y_{j} y_{j +s} -y_{j +r} y_{j +s +r}\right )\end{align}$\rfloor $ 

and
\begin{align}\left (\mathbf{c}^{ \dag } ,\mathbf{c}\right ) &  =\left (\mathbf{a}^{ \dag } ,\mathbf{a}\right )\left (\begin{array}{cc}\mathbb{T} & \mathbb{O} \\
\mathbb{O} & \overline{\mathbb{T}}\end{array}\right ) \\
\left (\mathbf{y}_{1} ,\mathbf{y}_{2}\right ) &  =\left (\mathbf{x}_{1} ,\mathbf{x}_{2}\right )\left (\begin{array}{cc}\ensuremath{\operatorname*{Re}} \left (\mathbb{T}\right ) &  -\ensuremath{\operatorname*{Im}} \left (\mathbb{T}\right ) \\
\ensuremath{\operatorname*{Im}} \left (\mathbb{T}\right ) & \ensuremath{\operatorname*{Re}} \left (\mathbb{T}\right )\end{array}\right )\text{,}\end{align}where $\mathbb{T} \in U \left (r\right )$ and produces the $r \times r$ matrix
\begin{equation}\mathbf{\Xi } =\left (\begin{array}{cc}\mathbf{0} & \mathbf{\zeta } \\
 -\mathbf{\zeta } & \mathbf{0}\end{array}\right )
\end{equation}from $\mathbf{\gamma }\text{.}$ 

It can be shown that the adjoint action gives the spinor representation of $O \left (2 r ,\mathbb{R}\right )$ as
\begin{align} & \widehat{\exp  \left \{\sum _{1 \leq j <k \leq r}\left (\gamma _{j k} a_{j}^{ \dag } a_{k}^{ \dag } +\overline{\gamma }_{j k} a_{k} a_{j}\right )\right \} \exp  \left \{\sum _{1 \leq j ,k \leq r}\lambda _{j k} a_{j}^{ \dag } a_{k}\right \}} \left (\begin{array}{cc}\mathbf{a}^{ \dag } & \mathbf{a}\end{array}\right ) \\
 &  =\left (\begin{array}{cc}\mathbf{a}^{ \dag } & \mathbf{a}\end{array}\right ) \left (\begin{array}{cc}\sqrt{\left (\mathbb{I}_{r} -\mathbb{Z} \mathbb{Z}^{ \dag }\right )} & \mathbb{Z} \\
 -\mathbb{Z}^{ \dag } & \sqrt{\left (\mathbb{I}_{r} -\mathbb{Z}^{ \dag } \mathbb{Z}\right )}\end{array}\right ) \left (\begin{array}{cc}\mathbb{W} & \mathbb{O} \\
\mathbb{O} & \overline{\mathbb{W}}\end{array}\right ) =\left (\begin{array}{cc}\mathbf{a}^{ \dag } & \mathbf{a}\end{array}\right ) \left (\begin{array}{cc}\mathbb{U} & \overline{\mathbb{V}} \\
\mathbb{V} & \overline{\mathbb{U}}\end{array}\right )\text{,}\end{align}$\lceil $
\begin{equation}\left (\begin{array}{cc}\sqrt{\left (\mathbb{I}_{r} -\mathbb{Z} \mathbb{Z}^{ \dag }\right )} & \mathbb{Z} \\
 -\mathbb{Z}^{ \dag } & \sqrt{\left (\mathbb{I}_{r} -\mathbb{Z}^{ \dag } \mathbb{Z}\right )}\end{array}\right ) \left (\begin{array}{cc}\mathbb{W} & \mathbb{O} \\
\mathbb{O} & \overline{\mathbb{W}}\end{array}\right ) =\left (\begin{array}{cc}\sqrt{\left (\mathbb{I}_{r} -\mathbb{Z} \mathbb{Z}^{ \dag }\right )} \mathbb{W} & \mathbb{Z} \overline{\mathbb{W}} \\
 -\mathbb{Z}^{ \dag } \mathbb{W} & \sqrt{\left (\mathbb{I}_{r} -\mathbb{Z}^{ \dag } \mathbb{Z}\right )} \overline{\mathbb{W}}\end{array}\right )
\end{equation}and
\begin{equation} -\overline{\mathbb{Z}^{ \dag } \mathbb{W}} = -\mathbb{Z}^{t} \overline{\mathbb{W}} =\mathbb{Z} \overline{\mathbb{W}}
\end{equation}
\begin{equation}\left (\mathbb{Z} \overline{\mathbb{W}}\right )^{t} =\mathbb{W}^{ \dag } \mathbb{Z}^{t} = -\mathbb{W}^{ \dag } \mathbb{Z} \neq \mathbb{Z} \overline{\mathbb{W}}
\end{equation}
\begin{align}\mathbb{U} &  =\sqrt{\left (\mathbb{I}_{r} -\mathbb{Z} \mathbb{Z}^{ \dag }\right )} \mathbb{W} \\
\mathbb{V} &  =\mathbb{Z} \overline{\mathbb{W}}\end{align}$\rfloor \;$ 

where
\begin{align*} & \mathbb{Z} =\overline{\mathbb{V}} =\mathbf{\gamma } \frac{\sin  \left (\sqrt{\mathbf{\gamma }^{ \dag } \mathbf{\gamma }}\right )}{\sqrt{\mathbf{\gamma }^{ \dag } \mathbf{\gamma }}} =\mathbf{\gamma } \ensuremath{\operatorname*{sinc}} \left (\sqrt{\mathbf{\gamma }^{ \dag } \mathbf{\gamma }}\right ) = -\mathbb{Z}^{t} \\
 & \mathbb{Z}^{ \dag } =\overline{\mathbb{Z}}^{t} = -\overline{\mathbb{Z}}\text{,} \\
 & \mathbf{\gamma }  = -\mathbf{\gamma }^{t}\text{,}\end{align*}so that
\begin{equation}\left (\begin{array}{cc}\mathbb{U} & \overline{\mathbb{V}} \\
\mathbb{V} & \overline{\mathbb{U}}\end{array}\right ) =\left (\begin{array}{cc}\sqrt{\left (\mathbb{I}_{r} -\mathbb{Z} \mathbb{Z}^{ \dag }\right )} & \mathbb{Z} \\
 -\mathbb{Z}^{ \dag } & \sqrt{\left (\mathbb{I}_{r} -\mathbb{Z}^{ \dag } \mathbb{Z}\right )}\end{array}\right ) \left (\begin{array}{cc}\mathbb{W} & \mathbb{O} \\
\mathbb{O} & \overline{\mathbb{W}}\end{array}\right ) =\left (\begin{array}{cc}\cos  \left (\sqrt{\mathbf{\gamma }^{ \dag } \mathbf{\gamma }}\right ) & \mathbf{\gamma } \frac{\sin  \left (\sqrt{\mathbf{\gamma }^{ \dag } \mathbf{\gamma }}\right )}{\sqrt{\mathbf{\gamma }^{ \dag } \mathbf{\gamma }}} \\
 -\frac{\sin  \left (\sqrt{\mathbf{\gamma }^{ \dag } \mathbf{\gamma }}\right )}{\sqrt{\mathbf{\gamma }^{ \dag } \mathbf{\gamma }}} \mathbf{\gamma }^{ \dag } & \cos  \left (\sqrt{\mathbf{\gamma }^{ \dag } \mathbf{\gamma }}\right )\end{array}\right ) \left (\begin{array}{cc}\mathbb{W} & \mathbb{O} \\
\mathbb{O} & \overline{\mathbb{W}}\end{array}\right )\text{.}
\end{equation}Observe
\begin{align}\mathbb{Z} &  =\overline{\mathbb{V}} \\
 -\overline{\mathbb{Z}}^{t} &  =\mathbb{V} \Rightarrow \overline{\mathbb{Z}} = -\mathbb{Z}^{ \dag }\text{.}\end{align}

The canonical form of this action is
\begin{align} & \widehat{\exp  \left \{ \sum _{1 \leq j \leq s}\zeta _{j} \left (c_{j}^{ \dag } c_{j +s}^{ \dag } +c_{j} c_{j +s}\right )\right \}} \exp  \left \{i \sum _{1 \leq j ,k \leq r}\lambda _{j k} a_{j}^{ \dag } a_{k}\right \} \left (\begin{array}{cc}\mathbf{c}^{ \dag } & \mathbf{c}\end{array}\right ) \\
 &  =\left (\begin{array}{cc}\mathbf{c}^{ \dag } & \mathbf{c}\end{array}\right ) \left (\begin{array}{cc}\mathbb{W}_{c} & \mathbb{O} \\
\mathbb{O} & \overline{\mathbb{W}}_{c}\end{array}\right ) \left (\begin{array}{cc}\left (\sqrt{\left (\mathbb{I}_{r} -\mathbb{Z} \mathbb{Z}^{ \dag }\right )}\right )_{c} & \left (\mathbb{Z}\right )_{c} \\
\left ( -\mathbb{Z}^{ \dag }\right )_{c} & \left (\sqrt{\left (\mathbb{I}_{r} -\mathbb{Z}^{ \dag } \mathbb{Z}\right )}\right )_{c}\end{array}\right ) =\left (\begin{array}{cc}\mathbf{c}^{ \dag } & \mathbf{c}\end{array}\right ) \left (\begin{array}{cc}\left (\mathbb{U}\right )_{c} & \left (\overline{\mathbb{V}}\right )_{c} \\
\left (\mathbb{V}\right )_{c} & \left (\overline{\mathbb{U}}\right )_{c}\end{array}\right )\text{,}\end{align}

\subparagraph{Unitarity Conditions}
Canonical constraints give
\begin{align}\left (\begin{array}{cc}\mathbb{U} & \overline{\mathbb{V}} \\
\mathbb{V} & \overline{\mathbb{U}}\end{array}\right )^{ \dag } \left (\begin{array}{cc}\mathbb{U} & \overline{\mathbb{V}} \\
\mathbb{V} & \overline{\mathbb{U}}\end{array}\right ) &  =\left (\begin{array}{cc}\mathbb{U}^{ \dag } & \mathbb{V}^{ \dag } \\
\mathbb{V}^{t} & \mathbb{U}^{t}\end{array}\right ) \left (\begin{array}{cc}\mathbb{U} & \overline{\mathbb{V}} \\
\mathbb{V} & \overline{\mathbb{U}}\end{array}\right ) =\left (\begin{array}{cc}\mathbb{U}^{ \dag } \mathbb{U} +\mathbb{V}^{ \dag } \mathbb{V} & \mathbb{U}^{ \dag } \overline{\mathbb{V}} +\mathbb{V}^{ \dag } \overline{\mathbb{U}} \\
\mathbb{V}^{t} \mathbb{U} +\mathbb{U}^{t} \mathbb{V} & \mathbb{V}^{t} \overline{\mathbb{V}} +\mathbb{U}^{t} \overline{\mathbb{U}}\end{array}\right ) =\left (\begin{array}{cc}\mathbb{I}_{r} & \mathbb{O} \\
\mathbb{O} & \mathbb{I}_{r}\end{array}\right ) \\
\mathbb{U}^{ \dag } \mathbb{U} +\mathbb{V}^{ \dag } \mathbb{V} &  =\mathbb{I}_{r} \\
\mathbb{V}^{t} \mathbb{U} &  = -\mathbb{U}^{t} \mathbb{V} \Leftrightarrow \left (\overline{\mathbb{Z}}\right ) \left (\sqrt{\left (\mathbb{I}_{r} -\mathbb{Z} \mathbb{Z}^{ \dag }\right )}\right ) =\left (\sqrt{\left (\mathbb{I}_{r} -\mathbb{Z}^{ \dag } \mathbb{Z}\right )}\right ) \left (\overline{\mathbb{Z}}\right )\text{,}\end{align}where
\begin{equation}\left (\begin{array}{cc}\mathbf{a}^{ \dag } & \mathbf{a}\end{array}\right ) \left (\begin{array}{cc}\sqrt{\left (\mathbb{I}_{r} -\mathbb{Z} \mathbb{Z}^{ \dag }\right )} & \mathbb{Z} \\
 -\mathbb{Z}^{ \dag } & \sqrt{\left (\mathbb{I}_{r} -\mathbb{Z}^{ \dag } \mathbb{Z}\right )}\end{array}\right ) =\left (\begin{array}{cc}\mathbf{a}^{ \dag } & \mathbf{a}\end{array}\right ) \left (\begin{array}{cc}\mathbb{U} & \overline{\mathbb{V}} \\
\mathbb{V} & \overline{\mathbb{U}}\end{array}\right )\text{.}
\end{equation}$\lceil $Let $\mathbb{Z}$ be unitary then
\begin{equation}\left (\begin{array}{cc}\mathbb{O} & \mathbb{Z} \\
 -\mathbb{Z}^{ \dag } & \mathbb{O}\end{array}\right ) \left (\begin{array}{cc}\mathbb{O} &  -\mathbb{Z} \\
\mathbb{Z}^{ \dag } & \mathbb{O}\end{array}\right ) =\left (\begin{array}{cc}\mathbb{Z} \mathbb{Z}^{ \dag } & \mathbb{O} \\
\mathbb{O} & \mathbb{Z}^{ \dag } \mathbb{Z}\end{array}\right ) =\left (\begin{array}{cc}\mathbb{I}_{r} & \mathbb{O} \\
\mathbb{O} & \mathbb{I}_{r}\end{array}\right )
\end{equation}If
\begin{equation}\mathbb{Z} = -\mathbb{Z}^{ \dag } = -\overline{\mathbb{Z}}^{t} = -\mathbb{Z}^{t}
\end{equation}then the transformation becomes a particle hole transformation
\begin{equation}\left (\begin{array}{cc}\mathbb{O} & \mathbb{Z} \\
\mathbb{Z} & \mathbb{O}\end{array}\right )\text{.}
\end{equation}$\rfloor $ 

The unitary constraint is shown to be fulfilled by observing
\begin{align} & \left (\begin{array}{cc}\sqrt{\left (\mathbb{I}_{r} -\mathbb{Z} \mathbb{Z}^{ \dag }\right )} & \mathbb{Z} \\
 -\mathbb{Z}^{ \dag } & \sqrt{\left (\mathbb{I}_{r} -\mathbb{Z}^{ \dag } \mathbb{Z}\right )}\end{array}\right )^{ \dag } \left (\begin{array}{cc}\sqrt{\left (\mathbb{I}_{r} -\mathbb{Z} \mathbb{Z}^{ \dag }\right )} & \mathbb{Z} \\
 -\mathbb{Z}^{ \dag } & \sqrt{\left (\mathbb{I}_{r} -\mathbb{Z}^{ \dag } \mathbb{Z}\right )}\end{array}\right ) \\
 &  =\left (\begin{array}{cc}\sqrt{\left (\mathbb{I}_{r} -\mathbb{Z} \mathbb{Z}^{ \dag }\right )} &  -\mathbb{Z} \\
\mathbb{Z}^{ \dag } & \sqrt{\left (\mathbb{I}_{r} -\mathbb{Z}^{ \dag } \mathbb{Z}\right )}\end{array}\right ) \left (\begin{array}{cc}\sqrt{\left (\mathbb{I}_{r} -\mathbb{Z} \mathbb{Z}^{ \dag }\right )} & \mathbb{Z} \\
 -\mathbb{Z}^{ \dag } & \sqrt{\left (\mathbb{I}_{r} -\mathbb{Z}^{ \dag } \mathbb{Z}\right )}\end{array}\right ) \\
 &  =\left (\begin{array}{cc}\left (\mathbb{I}_{r} -\mathbb{Z} \mathbb{Z}^{ \dag }\right ) +\mathbb{Z} \mathbb{Z}^{ \dag } & \sqrt{\left (\mathbb{I}_{r} -\mathbb{Z} \mathbb{Z}^{ \dag }\right )} \mathbb{Z} -\mathbb{Z} \sqrt{\left (\mathbb{I}_{r} -\mathbb{Z}^{ \dag } \mathbb{Z}\right )} \\
\mathbb{Z}^{ \dag } \sqrt{\left (\mathbb{I}_{r} -\mathbb{Z} \mathbb{Z}^{ \dag }\right )} -\sqrt{\left (\mathbb{I}_{r} -\mathbb{Z}^{ \dag } \mathbb{Z}\right )} \mathbb{Z}^{ \dag } & \mathbb{Z}^{ \dag } \mathbb{Z} +\left (\mathbb{I}_{r} -\mathbb{Z}^{ \dag } \mathbb{Z}\right )\end{array}\right ) \\
 &  =\left (\begin{array}{cc}\mathbb{I}_{r} & \mathbb{O} \\
\mathbb{O} & \mathbb{I}_{r}\end{array}\right )\end{align}as
\begin{equation}\mathbb{Z}^{ \dag } \sqrt{\left (\mathbb{I}_{r} -\mathbb{Z} \mathbb{Z}^{ \dag }\right )} =\sqrt{\left (\mathbb{I}_{r} -\mathbb{Z}^{ \dag } \mathbb{Z}\right )} \mathbb{Z}^{ \dag }\text{,}
\end{equation}which follows from
\begin{equation}\overline{\mathbb{Z}} \sqrt{\left (\mathbb{I}_{r} -\mathbb{Z} \mathbb{Z}^{ \dag }\right )} =\sqrt{\left (\mathbb{I}_{r} -\mathbb{Z}^{ \dag } \mathbb{Z}\right )} \overline{\mathbb{Z}}
\end{equation}and
\begin{equation}\mathbb{Z} = -\mathbb{Z}^{t}\text{.}
\end{equation}

The unitary constraint in the $\gamma $ parameterization involves the following relationships
\begin{align} & \mathbf{\gamma } \frac{\sin  \left (\sqrt{\mathbf{\gamma }^{ \dag } \mathbf{\gamma }}\right )}{\sqrt{\mathbf{\gamma }^{ \dag } \mathbf{\gamma }}} \frac{\sin  \left (\sqrt{\mathbf{\gamma }^{ \dag } \mathbf{\gamma }}\right )}{\sqrt{\mathbf{\gamma }^{ \dag } \mathbf{\gamma }}} \mathbf{\gamma }^{ \dag } =\mathbf{\gamma }\sin \left (\sqrt{\mathbf{\gamma }^{ \dag } \mathbf{\gamma }}\right )\left (\mathbf{\gamma }^{ \dag } \mathbf{\gamma }\right )^{ -1}\sin \left (\sqrt{\mathbf{\gamma }^{ \dag } \mathbf{\gamma }}\right )\mathbf{\gamma }^{ \dag } \\
 &  =\sin \left (\sqrt{\mathbf{\gamma } \mathbf{\gamma }^{ \dag }}\right )\mathbf{\gamma } \left (\mathbf{\gamma }^{ \dag } \mathbf{\gamma }\right )^{ -1} \mathbf{\gamma }^{ \dag }\sin \left (\sqrt{\mathbf{\gamma } \mathbf{\gamma }^{ \dag }}\right ) =\sin \left (\sqrt{\mathbf{\gamma } \mathbf{\gamma }^{ \dag }}\right )\mathbf{\gamma } \mathbf{\gamma }^{ -\mathbf{1}} \mathbf{\gamma }^{ - \dag } \mathbf{\gamma }^{ \dag }\sin \left (\sqrt{\mathbf{\gamma } \mathbf{\gamma }^{ \dag }}\right ) \\
 &  =\sin \left (\sqrt{\mathbf{\gamma } \mathbf{\gamma }^{ \dag }}\right )\sin \left (\sqrt{\mathbf{\gamma } \mathbf{\gamma }^{ \dag }}\right ) =\sin \left (\sqrt{\overline{\mathbf{\gamma }^{ \dag } \mathbf{\gamma }}}\right )\sin \left (\sqrt{\overline{\mathbf{\gamma }^{ \dag } \mathbf{\gamma }}}\right ) =\sin ^{\mathbf{2}}\left (\sqrt{\mathbf{\gamma }^{ \dag } \mathbf{\gamma }}\right )\end{align}and
\begin{equation}\mathbf{\gamma } \frac{\sin  \left (\sqrt{\left (\mathbf{\gamma }^{ \dag } \mathbf{\gamma }\right )^{t}}\right )}{\sqrt{\left (\mathbf{\gamma }^{ \dag } \mathbf{\gamma }\right )^{t}}} \cos  \left (\sqrt{\mathbf{\gamma }^{ \dag } \mathbf{\gamma }}\right ) =\left (\mathbf{\gamma } \frac{\sin  \left (\sqrt{\left (\mathbf{\gamma }^{ \dag } \mathbf{\gamma }\right )^{t}}\right )}{\sqrt{\left (\mathbf{\gamma }^{ \dag } \mathbf{\gamma }\right )^{t}}} \cos  \left (\sqrt{\mathbf{\gamma }^{ \dag } \mathbf{\gamma }}\right )\right )^{t} = -\cos  \left (\sqrt{\left (\mathbf{\gamma }^{ \dag } \mathbf{\gamma }\right )^{t}}\right ) \frac{\sin  \left (\sqrt{\mathbf{\gamma }^{ \dag } \mathbf{\gamma }}\right )}{\sqrt{\mathbf{\gamma }^{ \dag } \mathbf{\gamma }}} \mathbf{\gamma }
\end{equation}$\lceil $Note that
\begin{equation}\left (\mathbf{\gamma }^{ \dag } \mathbf{\gamma }\right )^{t} =\left (\overline{\mathbf{\gamma }}^{t} \mathbf{\gamma }\right )^{t} =\left (\mathbf{\gamma }^{t} \overline{\mathbf{\gamma }}\right )^{t} =\mathbf{\gamma }^{ \dag } \mathbf{\gamma }
\end{equation}In scalar form
\begin{equation}z =x +i y =x^{t} -i y^{t} = -x +i y = -\overline{z}
\end{equation}
\begin{align}\left (x +i y\right )^{ \dag } \left (x +i y\right ) &  =\left (x^{t} -i y^{t}\right ) \left (x +i y\right ) =\left ( -x +i y\right ) \left (x +i y\right ) = -x^{2} -i x y +i y x -y^{2} \\
\left (\left (x +i y\right )^{ \dag } \left (x +i y\right )\right )^{t} &  =\left (x +i y\right )^{t} \overline{\left (x +i y\right )} = -\left (x +i y\right ) \left (x -i y\right ) = -x^{2} +i x y -i y x -y^{2} \\
\left (x +i y\right )^{ \dag } \left (x +i y\right ) &  =\overline{\left (\left (x +i y\right )^{ \dag } \left (x +i y\right )\right )} =\overline{\left (\left (x +i y\right )^{ \dag } \left (x +i y\right )\right )}^{t} \\
 &  \Rightarrow \left (x +i y\right )^{ \dag } \left (x +i y\right ) =\left (\left (x +i y\right )^{ \dag } \left (x +i y\right )\right )^{t}\end{align}$\rfloor $ 

The unitary constraint in the $\gamma $ parameterization is given by
\begin{align} & \left (\begin{array}{cc}\cos  \left (\sqrt{\mathbf{\gamma }^{ \dag } \mathbf{\gamma }}\right ) & \mathbf{\gamma } \frac{\sin  \left (\sqrt{\mathbf{\gamma }^{ \dag } \mathbf{\gamma }}\right )}{\sqrt{\mathbf{\gamma }^{ \dag } \mathbf{\gamma }}} \\
 -\frac{\sin  \left (\sqrt{\mathbf{\gamma }^{ \dag } \mathbf{\gamma }}\right )}{\sqrt{\mathbf{\gamma }^{ \dag } \mathbf{\gamma }}} \mathbf{\gamma }^{ \dag } & \cos  \left (\sqrt{\mathbf{\gamma }^{ \dag } \mathbf{\gamma }}\right )\end{array}\right )^{ \dag } \left (\begin{array}{cc}\cos  \left (\sqrt{\mathbf{\gamma }^{ \dag } \mathbf{\gamma }}\right ) & \mathbf{\gamma } \frac{\sin  \left (\sqrt{\mathbf{\gamma }^{ \dag } \mathbf{\gamma }}\right )}{\sqrt{\mathbf{\gamma }^{ \dag } \mathbf{\gamma }}} \\
 -\frac{\sin  \left (\sqrt{\mathbf{\gamma }^{ \dag } \mathbf{\gamma }}\right )}{\sqrt{\mathbf{\gamma }^{ \dag } \mathbf{\gamma }}} \mathbf{\gamma }^{ \dag } & \cos  \left (\sqrt{\mathbf{\gamma }^{ \dag } \mathbf{\gamma }}\right )\end{array}\right ) \\
 &  =\left (\begin{array}{cc}\cos  \left (\sqrt{\mathbf{\gamma }^{ \dag } \mathbf{\gamma }}\right ) &  -\mathbf{\gamma } \frac{\sin  \left (\sqrt{\mathbf{\gamma }^{ \dag } \mathbf{\gamma }}\right )}{\sqrt{\mathbf{\gamma }^{ \dag } \mathbf{\gamma }}} \\
\frac{\sin  \left (\sqrt{\mathbf{\gamma }^{ \dag } \mathbf{\gamma }}\right )}{\sqrt{\mathbf{\gamma }^{ \dag } \mathbf{\gamma }}} \mathbf{\gamma }^{ \dag } & \cos  \left (\sqrt{\mathbf{\gamma }^{ \dag } \mathbf{\gamma }}\right )\end{array}\right ) \left (\begin{array}{cc}\cos  \left (\sqrt{\mathbf{\gamma }^{ \dag } \mathbf{\gamma }}\right ) & \mathbf{\gamma } \frac{\sin  \left (\sqrt{\mathbf{\gamma }^{ \dag } \mathbf{\gamma }}\right )}{\sqrt{\mathbf{\gamma }^{ \dag } \mathbf{\gamma }}} \\
 -\frac{\sin  \left (\sqrt{\mathbf{\gamma }^{ \dag } \mathbf{\gamma }}\right )}{\sqrt{\mathbf{\gamma }^{ \dag } \mathbf{\gamma }}} \mathbf{\gamma }^{ \dag } & \cos  \left (\sqrt{\mathbf{\gamma }^{ \dag } \mathbf{\gamma }}\right )\end{array}\right ) \\
 &  =\left (\begin{array}{cc}\cos ^{2} \left (\sqrt{\mathbf{\gamma }^{ \dag } \mathbf{\gamma }}\right ) +\mathbf{\gamma } \frac{\sin  \left (\sqrt{\mathbf{\gamma }^{ \dag } \mathbf{\gamma }}\right )}{\sqrt{\mathbf{\gamma }^{ \dag } \mathbf{\gamma }}} \frac{\sin  \left (\sqrt{\mathbf{\gamma }^{ \dag } \mathbf{\gamma }}\right )}{\sqrt{\mathbf{\gamma }^{ \dag } \mathbf{\gamma }}} \mathbf{\gamma }^{ \dag } & \cos  \left (\sqrt{\mathbf{\gamma }^{ \dag } \mathbf{\gamma }}\right ) \mathbf{\gamma } \frac{\sin  \left (\sqrt{\mathbf{\gamma }^{ \dag } \mathbf{\gamma }}\right )}{\sqrt{\mathbf{\gamma }^{ \dag } \mathbf{\gamma }}} -\mathbf{\gamma } \frac{\sin  \left (\sqrt{\mathbf{\gamma }^{ \dag } \mathbf{\gamma }}\right )}{\sqrt{\mathbf{\gamma }^{ \dag } \mathbf{\gamma }}} \cos  \left (\sqrt{\mathbf{\gamma }^{ \dag } \mathbf{\gamma }}\right ) \\
\frac{\sin  \left (\sqrt{\mathbf{\gamma }^{ \dag } \mathbf{\gamma }}\right )}{\sqrt{\mathbf{\gamma }^{ \dag } \mathbf{\gamma }}} \mathbf{\gamma }^{ \dag } \cos  \left (\sqrt{\mathbf{\gamma }^{ \dag } \mathbf{\gamma }}\right ) -\cos  \left (\sqrt{\mathbf{\gamma }^{ \dag } \mathbf{\gamma }}\right ) \frac{\sin  \left (\sqrt{\mathbf{\gamma }^{ \dag } \mathbf{\gamma }}\right )}{\sqrt{\mathbf{\gamma }^{ \dag } \mathbf{\gamma }}} \mathbf{\gamma }^{ \dag } & \frac{\sin  \left (\sqrt{\mathbf{\gamma }^{ \dag } \mathbf{\gamma }}\right )}{\sqrt{\mathbf{\gamma }^{ \dag } \mathbf{\gamma }}} \mathbf{\gamma }^{ \dag } \mathbf{\gamma } \frac{\sin  \left (\sqrt{\mathbf{\gamma }^{ \dag } \mathbf{\gamma }}\right )}{\sqrt{\mathbf{\gamma }^{ \dag } \mathbf{\gamma }}} +\cos ^{2} \left (\sqrt{\mathbf{\gamma }^{ \dag } \mathbf{\gamma }}\right )\end{array}\right ) \\
 &  =\left (\begin{array}{cc}\sin ^{2} \left (\sqrt{\mathbf{\gamma }^{ \dag } \mathbf{\gamma }}\right ) +\cos ^{2} \left (\sqrt{\mathbf{\gamma }^{ \dag } \mathbf{\gamma }}\right ) & \cos  \left (\sqrt{\mathbf{\gamma }^{ \dag } \mathbf{\gamma }}\right ) \mathbf{\gamma } \frac{\sin  \left (\sqrt{\mathbf{\gamma }^{ \dag } \mathbf{\gamma }}\right )}{\sqrt{\mathbf{\gamma }^{ \dag } \mathbf{\gamma }}} -\mathbf{\gamma } \frac{\sin  \left (\sqrt{\mathbf{\gamma }^{ \dag } \mathbf{\gamma }}\right )}{\sqrt{\mathbf{\gamma }^{ \dag } \mathbf{\gamma }}} \cos  \left (\sqrt{\mathbf{\gamma }^{ \dag } \mathbf{\gamma }}\right ) \\
\frac{\sin  \left (\sqrt{\mathbf{\gamma }^{ \dag } \mathbf{\gamma }}\right )}{\sqrt{\mathbf{\gamma }^{ \dag } \mathbf{\gamma }}} \mathbf{\gamma }^{ \dag } \cos  \left (\sqrt{\mathbf{\gamma }^{ \dag } \mathbf{\gamma }}\right ) -\cos  \left (\sqrt{\mathbf{\gamma }^{ \dag } \mathbf{\gamma }}\right ) \frac{\sin  \left (\sqrt{\mathbf{\gamma }^{ \dag } \mathbf{\gamma }}\right )}{\sqrt{\mathbf{\gamma }^{ \dag } \mathbf{\gamma }}} \mathbf{\gamma }^{ \dag } & \sin ^{2} \left (\sqrt{\mathbf{\gamma }^{ \dag } \mathbf{\gamma }}\right ) +\cos ^{2} \left (\sqrt{\mathbf{\gamma }^{ \dag } \mathbf{\gamma }}\right )\end{array}\right ) \\
 &  =\left (\begin{array}{cc}\mathbb{I}_{r} & \mathbb{O} \\
\mathbb{O} & \mathbb{I}_{r}\end{array}\right )\end{align}$\rfloor $ 

\paragraph{Particle-Hole Transformations}
In particular a maximal particle-hole transformation is given by letting
\begin{equation}\left (\begin{array}{cc}\mathbb{U} & \overline{\mathbb{V}} \\
\mathbb{V} & \overline{\mathbb{U}}\end{array}\right ) =\left (\begin{array}{cc}\mathbb{O} & \overline{\mathbb{V}} \\
\mathbb{V} & \mathbb{O}\end{array}\right )\text{,}
\end{equation}so that
\begin{equation}\mathbb{U} =\mathbb{O}\text{.}
\end{equation}$\lceil $
\begin{equation}\left (\begin{array}{cc}\mathbb{O} & \overline{\mathbb{V}} \\
\mathbb{V} & \mathbb{O}\end{array}\right ) \left (\begin{array}{cc}\mathbb{O} & \overline{\mathbb{V}} \\
\mathbb{V} & \mathbb{O}\end{array}\right )^{ \dag } =\left (\begin{array}{cc}\mathbb{O} & \overline{\mathbb{V}} \\
\mathbb{V} & \mathbb{O}\end{array}\right ) \left (\begin{array}{cc}\mathbb{O} & \mathbb{V}^{ \dag } \\
\mathbb{V}^{t} & \mathbb{O}\end{array}\right ) =\left (\begin{array}{cc}\overline{\mathbb{V}} \mathbb{V}^{t} & \mathbb{O} \\
\mathbb{O} & \mathbb{V} \mathbb{V}^{ \dag }\end{array}\right ) =\left (\begin{array}{cc}\mathbb{I}_{r} & \mathbb{O} \\
\mathbb{O} & \mathbb{I}_{r}\end{array}\right )
\end{equation}
\begin{align}\mathbb{U} &  =\sqrt{\left (\mathbb{I}_{r} -\mathbb{Z} \mathbb{Z}^{ \dag }\right )} \mathbb{W} \\
\mathbb{V} &  =\mathbb{Z} \overline{\mathbb{W}}\end{align}The $U \left (r\right )$ subgroup acts as
\begin{equation}\left (\begin{array}{cc}\mathbf{a}^{ \dag } & \mathbf{a}\end{array}\right ) \left (\begin{array}{cc}\mathbb{W} & \mathbb{O} \\
\mathbb{O} & \overline{\mathbb{W}}\end{array}\right ) =\left (\begin{array}{cc}\mathbf{a}^{ \dag } \mathbb{W} & \mathbf{a} \overline{\mathbb{W}}\end{array}\right )\text{.}
\end{equation}Consider in particular
\begin{equation}\left (\begin{array}{cc}\mathbb{W} & \mathbb{O} \\
\mathbb{O} & \overline{\mathbb{W}}\end{array}\right ) =\left (\begin{array}{cc}e^{i \theta } \mathbb{I} & \mathbb{O} \\
\mathbb{O} & e^{ -i \theta } \mathbb{I}\end{array}\right )\text{,}
\end{equation}where
\begin{align}e^{i \theta } &  =\cos  \theta  +i \sin  \theta  \\
e^{i \frac{\pi }{2}} &  =i ,e^{i \pi } = -1 ,e^{i \frac{3 \pi }{2}} = -i ,e^{i 2 \pi } =1.\end{align}This gives
\begin{align}\left (\begin{array}{cc}e^{0} \mathbb{I} & \mathbb{O} \\
\mathbb{O} & e^{0} \mathbb{I}\end{array}\right ) &  =\left (\begin{array}{cc}\mathbb{I} & \mathbb{O} \\
\mathbb{O} & \mathbb{I}\end{array}\right ) ;\left (\begin{array}{cc}e^{i \frac{\pi }{2}} \mathbb{I} & \mathbb{O} \\
\mathbb{O} & e^{ -i \frac{\pi }{2}} \mathbb{I}\end{array}\right ) =\left (\begin{array}{cc}i \mathbb{I} & \mathbb{O} \\
\mathbb{O} &  -i \mathbb{I}\end{array}\right ) ;\left (\begin{array}{cc}e^{i \pi } \mathbb{I} & \mathbb{O} \\
\mathbb{O} & e^{ -i \pi } \mathbb{I}\end{array}\right ) =\left (\begin{array}{cc} -\mathbb{I} & \mathbb{O} \\
\mathbb{O} &  -\mathbb{I}\end{array}\right ) ; \\
\left (\begin{array}{cc}e^{i \frac{3 \pi }{2}} \mathbb{I} & \mathbb{O} \\
\mathbb{O} & e^{ -i \frac{3 \pi }{2}} \mathbb{I}\end{array}\right ) &  =\left (\begin{array}{cc} -i \mathbb{I} & \mathbb{O} \\
\mathbb{O} & i \mathbb{I}\end{array}\right ) ;\left (\begin{array}{cc}e^{i 2 \pi } \mathbb{I} & \mathbb{O} \\
\mathbb{O} & e^{ -i 2 \pi } \mathbb{I}\end{array}\right ) =\left (\begin{array}{cc}\mathbb{I} & \mathbb{O} \\
\mathbb{O} & \mathbb{I}\end{array}\right )\text{,}\end{align}thus produces the following bases
\begin{equation}\left (\begin{array}{cc}\mathbf{a}^{ \dag } & \mathbf{a}\end{array}\right ) ,\left (\begin{array}{cc}i \mathbf{a}^{ \dag } &  -i \mathbf{a}\end{array}\right ) ,\left (\begin{array}{cc} -\mathbf{a}^{ \dag } &  -\mathbf{a}\end{array}\right ) ,\left (\begin{array}{cc} -i \mathbf{a}^{ \dag } & i \mathbf{a}\end{array}\right )\text{.}
\end{equation}$\rfloor $ 

The condition $\mathbb{U} =\mathbb{O}$ gives ($\mathbb{W} \neq \mathbb{O})$
\begin{equation*}\mathbb{U} =\cos  \left (\sqrt{\mathbf{\gamma }^{ \dag } \mathbf{\gamma }}\right ) \mathbb{W} \Rightarrow \cos  \left (\mathbf{\zeta }\right ) =\mathbb{O} \Rightarrow \mathbf{\zeta } =\left (\begin{array}{ccc}\zeta _{1} & \cdots  & 0 \\
\vdots  & \ddots  & \vdots  \\
0 & \cdots  & \zeta _{r}\end{array}\right ) =\left (2 \mathbf{n} +1\right ) \frac{\pi }{2}\text{,}
\end{equation*}where $\mathbf{n}$ is a $r \times r$ matrix
\begin{equation}\mathbf{n} =\left (\begin{array}{ccc} \pm 1 ,0 & \cdots  & 0 \\
\vdots  & \ddots  & \vdots  \\
0 & \cdots  &  \pm 1 ,0\end{array}\right )\text{,}
\end{equation}which implies that
\begin{align}\mathbb{V} &  =\mathbb{T}^{ \dag } \left (\begin{array}{cc}\mathbb{O} & \mathbf{\zeta } \\
 -\overline{\mathbf{\zeta }} & \mathbb{O}\end{array}\right ) \overline{\mathbb{T}} \left (\mathbb{T}^{ \dag } \left (\begin{array}{cc}\mathbf{\zeta } & \mathbb{O} \\
\mathbb{O} & \mathbf{\zeta }\end{array}\right ) \mathbb{T}\right )^{ -1} \sin  \left (\mathbb{T}^{ \dag } \left (\begin{array}{cc}\mathbf{\zeta } & \mathbb{O} \\
\mathbb{O} & \mathbf{\zeta }\end{array}\right ) \mathbb{T}\right ) \\
 &  =\left (2 n_{j} +1\right ) \frac{\pi }{2} \mathbb{T}^{ \dag } \left (\begin{array}{cc}\mathbb{O} & \mathbb{I}_{s} \\
 -\mathbb{I}_{s} & \mathbb{O}\end{array}\right ) \overline{\mathbb{T}} \left (\mathbb{T}^{ \dag } \left (2 n_{j} +1\right ) \frac{\pi }{2} \left (\begin{array}{cc}\mathbb{I}_{s} & \mathbb{O} \\
\mathbb{O} & \mathbb{I}_{s}\end{array}\right ) \mathbb{T}\right )^{ -1} \sin  \left (\left (2 \mathbf{n} +1\right ) \frac{\pi }{2}\right ) \\
 &  =\mathbb{T}^{ \dag } \left (\begin{array}{cc}\mathbb{O} & \mathbb{I}_{s} \\
 -\mathbb{I}_{s} & \mathbb{O}\end{array}\right ) \overline{\mathbb{T}} \mathbb{T}^{ \dag } \left (\begin{array}{cc}\mathbb{I}_{s} & \mathbb{O} \\
\mathbb{O} & \mathbb{I}_{s}\end{array}\right ) \mathbb{T} \sin  \left (\left (2 \mathbf{n} +1\right ) \frac{\pi }{2}\right ) \\
 &  =\mathbb{T}^{ \dag } \left (\begin{array}{cc}\mathbb{O} & \mathbb{I}_{s} \\
 -\mathbb{I}_{s} & \mathbb{O}\end{array}\right ) \overline{\mathbb{T}} \sin  \left (\left (2 \mathbf{n} +1\right ) \frac{\pi }{2}\right ) = \pm \mathbb{T}^{ \dag } \left (\begin{array}{cc}\mathbb{O} & \mathbb{I}_{s} \\
 -\mathbb{I}_{s} & \mathbb{O}\end{array}\right ) \overline{\mathbb{T}} ;n_{j} =0 ,1 \leq j \leq r\end{align}where $s =\frac{r}{2}$ and we have utilized
\begin{align}\mathbb{T} \sqrt{\mathbf{\gamma }^{ \dag } \mathbf{\gamma }} \mathbb{T}^{ \dag } &  =\left (\begin{array}{cc}\mathbf{\zeta } & \mathbb{O} \\
\mathbb{O} & \mathbf{\zeta }\end{array}\right ) \\
\mathbb{T} \mathbf{\gamma } \mathbb{T}^{t} &  =\left (\begin{array}{cc}\mathbb{O} & \mathbf{\zeta } \\
 -\mathbf{\zeta } & \mathbb{O}\end{array}\right )\text{,}\end{align}and
\begin{align}\sqrt{\mathbf{\gamma }^{ \dag } \mathbf{\gamma }} &  =\mathbb{T}^{ \dag } \left (\begin{array}{cc}\mathbf{\zeta } & \mathbb{O} \\
\mathbb{O} & \mathbf{\zeta }\end{array}\right ) \mathbb{T} \\
\mathbf{\gamma } &  =\mathbb{T}^{ \dag } \left (\begin{array}{cc}\mathbb{O} & \mathbf{\zeta } \\
 -\overline{\mathbf{\zeta }} & \mathbb{O}\end{array}\right ) \overline{\mathbb{T}}\text{.}\end{align}$\lceil $ E.g. for $r =2$ 


\begin{align}\mathbb{T}_{k} &  =\left (\begin{array}{cccc}a_{k}^{ \dag } & a_{k +s}^{ \dag } & a_{k} & a_{k +s}\end{array}\right ) \left (\begin{array}{cc}\mathbb{O}_{2} & \left (\begin{array}{cc}0 & 1 \\
 -1 & 0\end{array}\right ) \\
\left (\begin{array}{cc}0 & 1 \\
 -1 & 0\end{array}\right ) & \mathbb{O}_{2}\end{array}\right ) \\
 &  =\left (\begin{array}{cccc} -a_{k +s} & a_{k} &  -a_{k +s}^{ \dag } & a_{k}^{ \dag }\end{array}\right ) =\left (\begin{array}{cccc}b_{k}^{ \dag } & b_{k +s}^{ \dag } & b_{k} & b_{k +s}\end{array}\right )\end{align}(Both the matrix $\mathbf{\gamma }$ and the non zero blocks have to be antisymmetric.) This particle hole transformation is unitary as
\begin{equation}\left (\begin{array}{cccc}0 & 0 & 0 & 1 \\
0 & 0 &  -1 & 0 \\
0 & 1 & 0 & 0 \\
 -1 & 0 & 0 & 0\end{array}\right ) \left (\begin{array}{cccc}0 & 0 & 0 &  -1 \\
0 & 0 & 1 & 0 \\
0 &  -1 & 0 & 0 \\
1 & 0 & 0 & 0\end{array}\right ) =\left (\begin{array}{cccc}1 & 0 & 0 & 0 \\
0 & 1 & 0 & 0 \\
0 & 0 & 1 & 0 \\
0 & 0 & 0 & 1\end{array}\right )\text{.}
\end{equation}Consider the $r =2$ example, one can observe that
\begin{gather}\exp  \left \{\frac{1}{2} \left (2 n +1\right )^{\frac{1}{2}} \genfrac{(}{)}{}{}{\pi }{2}^{\frac{1}{2}} \left (\begin{array}{cc}\mathbf{a}^{ \dag } & \mathbf{a}\end{array}\right ) \left (\begin{array}{cccc}0 & 0 & 0 & 1 \\
0 & 0 &  -1 & 0 \\
0 & 1 & 0 & 0 \\
 -1 & 0 & 0 & 0\end{array}\right ) \left (\begin{array}{c}\mathbf{a} \\
\mathbf{a}^{ \dag }\end{array}\right )\right \} \\
 =\left (\begin{array}{cc}\mathbf{a}^{ \dag } & \mathbf{a}\end{array}\right ) \left (\begin{array}{cccc}0 & 0 & 0 & 1 \\
0 & 0 &  -1 & 0 \\
0 & 1 & 0 & 0 \\
 -1 & 0 & 0 & 0\end{array}\right ) \left (\begin{array}{c}\mathbf{a} \\
\mathbf{a}^{ \dag }\end{array}\right ) \\
 \parallel  \\
\exp  \left \{\frac{1}{2} \left (2 n +1\right )^{\frac{1}{2}} \genfrac{(}{)}{}{}{\pi }{2}^{\frac{1}{2}} \left (\begin{array}{cc}\mathbf{x}_{1} & \mathbf{x}_{2}\end{array}\right ) \left (\begin{array}{cccc}0 & 1 & 0 & 0 \\
 -1 & 0 & 0 & 0 \\
0 & 0 & 0 &  -1 \\
0 & 0 & 1 & 0\end{array}\right ) \left (\begin{array}{c}\mathbf{x}_{1} \\
\mathbf{x}_{2}\end{array}\right )\right \} \\
 =\left (\begin{array}{cc}\mathbf{x}_{1} & \mathbf{x}_{2}\end{array}\right ) \left (\begin{array}{cccc}0 & 1 & 0 & \;0 \\
 -1 & 0 & 0 & 0 \\
0 & 0 & 0 &  -1 \\
0 & 0 & 1 & 0\end{array}\right ) \left (\begin{array}{c}\mathbf{x}_{1} \\
\mathbf{x}_{2}\end{array}\right )\text{.}\end{gather}
\begin{align} & \frac{1}{2} \left (\begin{array}{cccc}x_{1} & x_{2} & x_{3} & x_{4}\end{array}\right ) \left (\begin{array}{cccc}0 & 1 & 0 & \;0 \\
 -1 & 0 & 0 & 0 \\
0 & 0 & 0 &  -1 \\
0 & 0 & 1 & 0\end{array}\right ) \left (\begin{array}{c}x_{1} \\
x_{2} \\
x_{3} \\
x_{4}\end{array}\right ) =\frac{1}{2} \left (\begin{array}{cccc}x_{1} & x_{2} & x_{3} & x_{4}\end{array}\right ) \left (\begin{array}{c}x_{2} \\
 -x_{1} \\
 -x_{4} \\
x_{3}\end{array}\right ) \\
 &  =\frac{1}{2} \left (x_{1} x_{2} -x_{2} x_{1} -x_{3} x_{4} +x_{4} x_{3}\right ) =x_{1} x_{2} -x_{3} x_{4} =\frac{1}{2} \left (a_{1}^{ \dag } +a_{1}\right ) \left (a_{2}^{ \dag } +a_{2}\right ) +\frac{1}{2} \left (a_{1}^{ \dag } -a_{1}\right ) \left (a_{2}^{ \dag } -a_{2}\right ) \\
 &  =\frac{1}{2} \left (a_{1}^{ \dag } a_{2}^{ \dag } +a_{1} a_{2}^{ \dag } +a_{1}^{ \dag } a_{2} +a_{1} a_{2}\right ) +\frac{1}{2} \left (a_{1}^{ \dag } a_{2}^{ \dag } -a_{1} a_{2}^{ \dag } -a_{1}^{ \dag } a_{2} +a_{1} a_{2}\right ) =a_{1}^{ \dag } a_{2}^{ \dag } +a_{1} a_{2}\end{align}In general
\begin{equation}\exp  \left \{\frac{1}{2} \left (\mathbf{x}_{1} \ensuremath{\operatorname*{Re}} \left (\mathbf{\gamma }\right ) \mathbf{x}_{1} +\mathbf{x}_{1} \ensuremath{\operatorname*{Im}} \left (\mathbf{\gamma }\right ) \mathbf{x}_{2} +\mathbf{x}_{2} \ensuremath{\operatorname*{Im}} \left (\mathbf{\gamma }\right ) \mathbf{x}_{1} -\mathbf{x}_{2} \ensuremath{\operatorname*{Re}} \left (\mathbf{\gamma }\right ) \mathbf{x}_{2}\right )\right \}\text{.}
\end{equation}$\lceil $
\begin{align} & \left (\begin{array}{cccc}a_{k}^{ \dag } & a_{k +s}^{ \dag } & a_{k} & a_{k +s}\end{array}\right ) \left (\begin{array}{cccc}0 & 0 & 0 & 1 \\
0 & 0 &  -1 & 0 \\
0 & 1 & 0 & 0 \\
 -1 & 0 & 0 & 0\end{array}\right ) \left (\begin{array}{c}a_{k} \\
a_{k +s} \\
a_{k}^{ \dag } \\
a_{k +s}^{ \dag }\end{array}\right ) =\left (\begin{array}{cccc} -a_{k +s} & a_{k} &  -a_{k +s}^{ \dag } & a_{k}^{ \dag }\end{array}\right ) \left (\begin{array}{c}a_{k} \\
a_{k +s} \\
a_{k}^{ \dag } \\
a_{k +s}^{ \dag }\end{array}\right ) \\
 &  =\left (\begin{array}{cccc}a_{k}^{ \dag } & a_{k +s}^{ \dag } & a_{k} & a_{k +s}\end{array}\right ) \left (\begin{array}{cc}\mathbb{O} & \mathbb{J}_{2} \\
\mathbb{J}_{2} & \mathbb{O}\end{array}\right ) \left (\begin{array}{c}a_{k} \\
a_{k +s} \\
a_{k}^{ \dag } \\
a_{k +s}^{ \dag }\end{array}\right ) \\
 &  = -a_{k +s} a_{k} +a_{k} a_{k +s} -a_{k +s}^{ \dag } a_{k}^{ \dag } +a_{k}^{ \dag } a_{k +s}^{ \dag } =2 \left (a_{k} a_{k +s} +a_{k}^{ \dag } a_{k +s}^{ \dag }\right ) =2 \left (a_{k}^{ \dag } a_{k +s}^{ \dag } -a_{k +s} a_{k}\right ) =2 \left (a_{k}^{ \dag } a_{k +s}^{ \dag } +a_{k} a_{k +s}\right )\end{align}
\begin{align}\left (\begin{array}{cc}\mathbb{O} & \mathbb{J}_{2} \\
\mathbb{J}_{2} & \mathbb{O}\end{array}\right )^{t} &  = -\left (\begin{array}{cc}\mathbb{O} & \mathbb{J}_{2} \\
\mathbb{J}_{2} & \mathbb{O}\end{array}\right ) \\
\left (\begin{array}{cc}\mathbb{O} & \mathbb{J}_{2} \\
\mathbb{J}_{2} & \mathbb{O}\end{array}\right )^{ \dag } \left (\begin{array}{cc}\mathbb{O} & \mathbb{J}_{2} \\
\mathbb{J}_{2} & \mathbb{O}\end{array}\right ) &  =\left (\begin{array}{cc}\mathbb{O} &  -\mathbb{J}_{2} \\
 -\mathbb{J}_{2} & \mathbb{O}\end{array}\right ) \left (\begin{array}{cc}\mathbb{O} & \mathbb{J}_{2} \\
\mathbb{J}_{2} & \mathbb{O}\end{array}\right ) = -\left (\begin{array}{cc}\mathbb{I}_{2} & \mathbb{O} \\
\mathbb{O} & \mathbb{I}_{2}\end{array}\right ) \\
\left (\begin{array}{cccc}0 & 0 & 0 & 1 \\
0 & 0 &  -1 & 0 \\
0 & 1 & 0 & 0 \\
 -1 & 0 & 0 & 0\end{array}\right ) \left (\begin{array}{cccc}0 & 0 & 0 & 1 \\
0 & 0 &  -1 & 0 \\
0 & 1 & 0 & 0 \\
 -1 & 0 & 0 & 0\end{array}\right ) &  =\left (\begin{array}{cccc} -1 & 0 & 0 & 0 \\
0 &  -1 & 0 & 0 \\
0 & 0 &  -1 & 0 \\
0 & 0 & 0 &  -1\end{array}\right )\end{align}
\begin{equation}\mathbb{O} \mathbb{J}_{4} \mathbb{O}^{ \dag } \mathbb{O} \mathbb{J}_{4} \mathbb{O}^{ \dag } =\mathbb{O} \mathbb{J}_{4}^{2} \mathbb{O}^{ \dag } = -\mathbb{I}_{4}
\end{equation}
\begin{equation}\mathbb{T}_{k} = \pm \left (\begin{array}{cccc}a_{k}^{ \dag } & a_{k +s}^{ \dag } & a_{k} & a_{k +s}\end{array}\right ) \left (\begin{array}{cc}\mathbb{O} & \mathbb{I} \\
\mathbb{I} & \mathbb{O}\end{array}\right ) \left (\begin{array}{c}a_{k} \\
a_{k +s} \\
a_{k}^{ \dag } \\
a_{k +s}^{ \dag }\end{array}\right ) =\left (\begin{array}{cccc}a_{k}^{ \dag } & a_{k +s}^{ \dag } & a_{k} & a_{k +s}\end{array}\right ) \exp  \left \{\left (2 n +1\right ) \pi  \mathbb{O} \mathbb{J}_{4} \mathbb{O}^{ \dag }\right \} \left (\begin{array}{c}a_{k} \\
a_{k +s} \\
a_{k}^{ \dag } \\
a_{k +s}^{ \dag }\end{array}\right )
\end{equation}$\rfloor $ 

$\rfloor $
\begin{align}T &  =\left (\begin{array}{cc}\mathbf{a}^{ \dag } & \mathbf{a}\end{array}\right ) \left (\begin{array}{cc}\mathbb{U} & \overline{\mathbb{V}} \\
\mathbb{V} & \overline{\mathbb{U}}\end{array}\right ) \left (\begin{array}{c}\mathbf{a} \\
\mathbf{a}^{ \dag }\end{array}\right ) \\
 &  =\left (\begin{array}{cc}\mathbf{x}_{1} & \mathbf{x}_{2}\end{array}\right ) \left (\begin{array}{cc}\ensuremath{\operatorname*{Re}} \left (\mathbb{U} +\mathbb{V}\right ) &  -\ensuremath{\operatorname*{Im}} \left (\mathbb{U} +\mathbb{V}\right ) \\
\ensuremath{\operatorname*{Im}} \left (\mathbb{U} -\mathbb{V}\right ) & \ensuremath{\operatorname*{Re}} \left (\mathbb{U} -\mathbb{V}\right )\end{array}\right ) \left (\begin{array}{c}\mathbf{x}_{1} \\
\mathbf{x}_{2}\end{array}\right ) \\
 &  =\exp  \left \{\frac{1}{2} \left (\begin{array}{cc}\mathbf{a}^{ \dag } & \mathbf{a}\end{array}\right ) \left (\begin{array}{cc}\mathbb{O} & \mathbf{\gamma } \\
\overline{\mathbf{\gamma }} & \mathbb{O}\end{array}\right ) \left (\begin{array}{c}\mathbf{a} \\
\mathbf{a}^{ \dag }\end{array}\right )\right \} \\
 &  =\exp  \left \{\frac{1}{2} \left (\begin{array}{cc}\mathbf{x}_{1} & \mathbf{x}_{2}\end{array}\right ) \left (\begin{array}{cc}\ensuremath{\operatorname*{Re}} \left (\mathbf{\gamma }\right ) & \ensuremath{\operatorname*{Im}} \left (\mathbf{\gamma }\right ) \\
\ensuremath{\operatorname*{Im}} \left (\mathbf{\gamma }\right ) &  -\ensuremath{\operatorname*{Re}} \left (\mathbf{\gamma }\right )\end{array}\right ) \left (\begin{array}{c}\mathbf{x}_{1} \\
\mathbf{x}_{2}\end{array}\right )\right \}\end{align}

\subparagraph{Canonical Transformations in Different Bases}
First we recall the relationships between the different bases of $\mathcal{F}_{1}$: 


\begin{enumerate}
\item Transforming between real and non self adjoint bases
\begin{equation}K \left (\begin{array}{cc}\mathbf{a}^{ \dag } & \mathbf{a}\end{array}\right ) =\left (\begin{array}{cc}\mathbf{a}^{ \dag } & \mathbf{a}\end{array}\right ) \mathbb{K} =\left (\begin{array}{cc}\mathbf{a}^{ \dag } & \mathbf{a}\end{array}\right ) \left (\begin{array}{cc}\frac{1}{\sqrt{2}} \mathbb{I}_{r} & \frac{i}{\sqrt{2}} \mathbb{I}_{r} \\
\frac{1}{\sqrt{2}} \mathbb{I}_{r} & \frac{ -i}{\sqrt{2}} \mathbb{I}_{r}\end{array}\right ) =\left (\begin{array}{cc}\mathbf{x}_{1} & \mathbf{x}_{2}\end{array}\right )\text{.}
\end{equation}

\item Transforming canonicaly transformed bases
\begin{align*} & K \left (\begin{array}{cc}\mathbf{b}^{ \dag } & \mathbf{b}\end{array}\right ) =K \left (\begin{array}{cc}\mathbf{a}^{ \dag } & \mathbf{a}\end{array}\right ) \left (\begin{array}{cc}\mathbb{U} & \overline{\mathbb{V}} \\
\mathbb{V} & \overline{\mathbb{U}}\end{array}\right ) =\left (\begin{array}{cc}\mathbf{y}_{1} & \mathbf{y}_{2}\end{array}\right ) \\
 &  =\left (\begin{array}{cc}\mathbf{b}^{ \dag } & \mathbf{b}\end{array}\right ) \mathbb{K} =\left (\begin{array}{cc}\mathbf{a}^{ \dag } & \mathbf{a}\end{array}\right ) \mathbb{K} \left (\begin{array}{cc}\mathbb{U} & \overline{\mathbb{V}} \\
\mathbb{V} & \overline{\mathbb{U}}\end{array}\right ) =\left (\begin{array}{cc}\mathbf{y}_{1} & \mathbf{y}_{2}\end{array}\right )\text{.}\end{align*}

\item Canonical transformation between real basis
\begin{align*}T_{r} \left (\begin{array}{cc}\mathbf{x}_{1} & \mathbf{x}_{2}\end{array}\right ) &  =\left (\begin{array}{cc}\mathbf{y}_{1} & \mathbf{y}_{2}\end{array}\right ) \\
K^{ \dag } T_{r} K \left (\begin{array}{cc}\mathbf{a}^{ \dag } & \mathbf{a}\end{array}\right ) &  =\left (\begin{array}{cc}\mathbf{a}^{ \dag } & \mathbf{a}\end{array}\right ) \mathbb{K}^{ \dag } \mathbb{T}_{r} \mathbb{K} =K^{ \dag } \left (\begin{array}{cc}\mathbf{y}_{1} & \mathbf{y}_{2}\end{array}\right ) =\left (\begin{array}{cc}\mathbf{b}^{ \dag } & \mathbf{b}\end{array}\right ) =T \left (\begin{array}{cc}\mathbf{a}^{ \dag } & \mathbf{a}\end{array}\right ) =\left (\begin{array}{cc}\mathbf{a}^{ \dag } & \mathbf{a}\end{array}\right ) \mathbb{T} \\
\mathbb{T} &  =\mathbb{K} ^{ \dag } \mathbb{T}_{r} \mathbb{K} \equiv \mathbb{T}_{r} =\mathbb{K} \mathbb{T} \mathbb{K}^{ \dag }\text{.}\end{align*}

\item Canonical transformation between non self adjoint basis
\begin{align*}T \left (\begin{array}{cc}\mathbf{a}^{ \dag } & \mathbf{a}\end{array}\right ) &  =\left (\begin{array}{cc}\mathbf{b}^{ \dag } & \mathbf{b}\end{array}\right ) =\left (\begin{array}{cc}\mathbf{a}^{ \dag } & \mathbf{a}\end{array}\right ) \mathbb{T} =\left (\begin{array}{cc}\mathbf{a}^{ \dag } & \mathbf{a}\end{array}\right ) \left (\begin{array}{cc}\mathbb{U} & \overline{\mathbb{V}} \\
\mathbb{V} & \overline{\mathbb{U}}\end{array}\right ) \\
K T K^{ \dag } \left (\begin{array}{cc}\mathbf{x}_{1} & \mathbf{x}_{2}\end{array}\right ) &  =\left (\begin{array}{cc}\mathbf{x}_{1} & \mathbf{x}_{2}\end{array}\right ) \mathbb{K} \mathbb{T} \mathbb{K}^{ \dag } =K \left (\begin{array}{cc}\mathbf{b}^{ \dag } & \mathbf{b}\end{array}\right ) =\left (\begin{array}{cc}\mathbf{y}_{1} & \mathbf{y}_{2}\end{array}\right ) =T_{r} \left (\begin{array}{cc}\mathbf{x}_{1} & \mathbf{x}_{2}\end{array}\right ) =\left (\begin{array}{cc}\mathbf{x}_{1} & \mathbf{x}_{2}\end{array}\right ) \mathbb{T}_{r} \\
\mathbb{T} &  =\mathbb{K}^{ \dag } \mathbb{T}_{r} \mathbb{K} \equiv \mathbb{T}_{r} =\mathbb{K} \mathbb{T} \mathbb{K}^{ \dag }\text{.}\end{align*}\end{enumerate}


Canonical transformation between non self adjoint basis
are of the form
\begin{equation}\mathbb{T} =\left (\begin{array}{cc}\mathbb{U} & \overline{\mathbb{V}} \\
\mathbb{V} & \overline{\mathbb{U}}\end{array}\right )\text{,}
\end{equation}while transformation between self adjoint basis are of the form
\begin{align}\mathbb{T}_{r} &  =\left (\begin{array}{cc}\frac{1}{\sqrt{2}} \mathbb{I}_{r} & \frac{1}{\sqrt{2}} \mathbb{I}_{r} \\
 -\frac{i}{\sqrt{2}} \mathbb{I}_{r} & \frac{i}{\sqrt{2}} \mathbb{I}_{r}\end{array}\right ) \left (\begin{array}{cc}\mathbb{U} & \overline{\mathbb{V}} \\
\mathbb{V} & \overline{\mathbb{U}}\end{array}\right ) \left (\begin{array}{cc}\frac{1}{\sqrt{2}} \mathbb{I}_{r} & \frac{i}{\sqrt{2}} \mathbb{I}_{r} \\
\frac{1}{\sqrt{2}} \mathbb{I}_{r} & \frac{ -i}{\sqrt{2}} \mathbb{I}_{r}\end{array}\right ) \\
 &  =\left (\begin{array}{cc}\frac{1}{\sqrt{2}} \mathbb{U} +\frac{1}{\sqrt{2}} \mathbb{V} & \frac{1}{\sqrt{2}} \overline{\mathbb{V}} +\frac{1}{\sqrt{2}} \overline{\mathbb{U}} \\
 -\frac{i}{\sqrt{2}} \mathbb{U} +\frac{i}{\sqrt{2}} \mathbb{V} &  -\frac{i}{\sqrt{2}} \overline{\mathbb{V}} +\frac{i}{\sqrt{2}} \overline{\mathbb{U}}\end{array}\right ) \left (\begin{array}{cc}\frac{1}{\sqrt{2}} \mathbb{I}_{r} & \frac{i}{\sqrt{2}} \mathbb{I}_{r} \\
\frac{1}{\sqrt{2}} \mathbb{I}_{r} & \frac{ -i}{\sqrt{2}} \mathbb{I}_{r}\end{array}\right ) \\
 &  =\left (\begin{array}{cc}\frac{1}{2} \mathbb{U} +\frac{1}{2} \overline{\mathbb{U}} +\frac{1}{2} \mathbb{V} +\frac{1}{2} \overline{\mathbb{V}} & \frac{i}{2} \mathbb{U} -\frac{i}{2} \overline{\mathbb{U}} +\frac{i}{2} \mathbb{V} -\frac{i}{2} \overline{\mathbb{V}} \\
\frac{i}{2} \overline{\mathbb{U}} -\frac{i}{2} \mathbb{U} +\frac{i}{2} \mathbb{V} -\frac{i}{2} \overline{\mathbb{V}} & \frac{1}{2} \mathbb{U} +\frac{1}{2} \overline{\mathbb{U}} -\frac{1}{2} \mathbb{V} -\frac{1}{2} \overline{\mathbb{V}}\end{array}\right ) \\
 &  =\frac{1}{2} \left (\begin{array}{cc}2 \ensuremath{\operatorname*{Re}} \left (\mathbb{U} +\mathbb{V}\right ) &  -2 \ensuremath{\operatorname*{Im}} \left (\mathbb{U} +\mathbb{V}\right ) \\
2 \ensuremath{\operatorname*{Im}} \left (\mathbb{U} -\mathbb{V}\right ) & 2 \ensuremath{\operatorname*{Re}} \left (\mathbb{U} -\mathbb{V}\right )\end{array}\right ) =\left (\begin{array}{cc}\ensuremath{\operatorname*{Re}} \left (\mathbb{U} +\mathbb{V}\right ) &  -\ensuremath{\operatorname*{Im}} \left (\mathbb{U} +\mathbb{V}\right ) \\
\ensuremath{\operatorname*{Im}} \left (\mathbb{U} -\mathbb{V}\right ) & \ensuremath{\operatorname*{Re}} \left (\mathbb{U} -\mathbb{V}\right )\end{array}\right )\text{.}\end{align}$\lceil $ We check this using scalers
\begin{align} & \left (\begin{array}{cc}\frac{1}{\sqrt{2}} & \frac{1}{\sqrt{2}} \\
 -\frac{i}{\sqrt{2}} & \frac{i}{\sqrt{2}}\end{array}\right ) \left (\begin{array}{cc}u & \overline{v} \\
v & \overline{u}\end{array}\right ) \left (\begin{array}{cc}\frac{1}{\sqrt{2}} & \frac{i}{\sqrt{2}} \\
\frac{1}{\sqrt{2}} & \frac{ -i}{\sqrt{2}}\end{array}\right ) =\frac{1}{2} \left (\begin{array}{cc}1 & 1 \\
 -i & i\end{array}\right ) \left (\begin{array}{cc}u & \overline{v} \\
v & \overline{u}\end{array}\right ) \left (\begin{array}{cc}1 & i \\
1 &  -i\end{array}\right ) \\
 &  =\frac{1}{2} \left (\begin{array}{cc}u +v & \overline{v} +\overline{u} \\
i v -i u & i \overline{u} -i \overline{v}\end{array}\right ) \left (\begin{array}{cc}1 & i \\
1 &  -i\end{array}\right ) =\frac{1}{2} \left (\begin{array}{cc}u +v +\overline{v} +\overline{u} & i \left (u +v -\overline{v} -\overline{u}\right ) \\
i \left (v -u +\overline{u} -\overline{v}\right ) &  -v +u +\overline{u} -\overline{v}\end{array}\right ) \\
 &  =\left (\begin{array}{cc}\ensuremath{\operatorname*{Re}} \left (u +v\right ) &  -\ensuremath{\operatorname*{Im}} \left (u +v\right ) \\
\ensuremath{\operatorname*{Im}} \left (u -v\right ) & \ensuremath{\operatorname*{Re}} \left (u -v\right )\end{array}\right )\text{,}\end{align}where we have used
\begin{align}i \left (x +i y\right ) -i \left (x -i y\right ) &  =i x -y -i x -y = -2 y \\
i \left (x -i y\right ) -i \left (x +i y\right ) &  =i x +y -i x +y =2 y \\
\left (x +i y\right ) +\left (x -i y\right ) &  =2 x\text{.}\end{align}$\rfloor $ 

Noting that
\begin{equation}\left (\begin{array}{cc}\mathbf{x}_{1} & \mathbf{x}_{2}\end{array}\right ) =\left (\begin{array}{cc}\mathbf{a}^{ \dag } & \mathbf{a}\end{array}\right ) \left (\begin{array}{cc}\frac{1}{\sqrt{2}} \mathbb{I}_{r} & \frac{i}{\sqrt{2}} \mathbb{I}_{r} \\
\frac{1}{\sqrt{2}} \mathbb{I}_{r} &  -\frac{i}{\sqrt{2}} \mathbb{I}_{r}\end{array}\right )\text{,}
\end{equation}and
\begin{equation}\left (\begin{array}{cc}\mathbf{a}^{ \dag } & \mathbf{a}\end{array}\right ) =\left (\begin{array}{cc}\mathbf{x}_{1} & \mathbf{x}_{2}\end{array}\right ) \left (\begin{array}{cc}\frac{1}{\sqrt{2}} \mathbb{I}_{r} & \frac{1}{\sqrt{2}} \mathbb{I}_{r} \\
 -\frac{i}{\sqrt{2}} \mathbb{I}_{r} & \frac{i}{\sqrt{2}} \mathbb{I}_{r}\end{array}\right )
\end{equation}one obtains 

\begin{equation}\left (\begin{array}{cc}\mathbf{b}^{ \dag } & \mathbf{b}\end{array}\right ) =\widehat{\exp  \left \{\sum _{1 \leq j <k \leq r}\left (\gamma _{j k} a_{j}^{ \dag } a_{k}^{ \dag } +\overline{\gamma }_{j k} a_{k} a_{j}\right )\right \} \exp  \left \{\sum _{1 \leq j ,k \leq r}\lambda _{j k} a_{j}^{ \dag } a_{k}\right \}} \left (\begin{array}{cc}\mathbf{a}^{ \dag } & \mathbf{a}\end{array}\right ) =\left (\begin{array}{cc}\mathbf{a}^{ \dag } & \mathbf{a}\end{array}\right ) \left (\begin{array}{cc}\mathbb{U} & \overline{\mathbb{V}} \\
\mathbb{V} & \overline{\mathbb{U}}\end{array}\right )\text{.}
\end{equation}This leads to
\begin{align}K \left (\begin{array}{cc}\mathbf{b}^{ \dag } & \mathbf{b}\end{array}\right ) &  =\left (\begin{array}{cc}\mathbf{y}_{1} & \mathbf{y}_{2}\end{array}\right ) =K \widehat{\exp  \left \{\sum _{1 \leq j <k \leq r}\left (\gamma _{j k} a_{j}^{ \dag } a_{k}^{ \dag } +\overline{\gamma }_{j k} a_{k} a_{j}\right )\right \} \exp  \left \{\sum _{1 \leq j ,k \leq r}\lambda _{j k} a_{j}^{ \dag } a_{k}\right \}} K^{ \dag } K \left (\begin{array}{cc}\mathbf{a}^{ \dag } & \mathbf{a}\end{array}\right ) \\
 &  =K \widehat{\exp  \left \{\sum _{1 \leq j <k \leq r}\left (\gamma _{j k} a_{j}^{ \dag } a_{k}^{ \dag } +\overline{\gamma }_{j k} a_{k} a_{j}\right )\right \} \exp  \left \{\sum _{1 \leq j ,k \leq r}\lambda _{j k} a_{j}^{ \dag } a_{k}\right \}} K^{ \dag } \left (\begin{array}{cc}\mathbf{x}_{1} & \mathbf{x}_{2}\end{array}\right ) =K \left (\begin{array}{cc}\mathbf{a}^{ \dag } & \mathbf{a}\end{array}\right ) \left (\begin{array}{cc}\mathbb{U} & \overline{\mathbb{V}} \\
\mathbb{V} & \overline{\mathbb{U}}\end{array}\right ) \\
 &  =\left (\begin{array}{cc}\mathbf{x}_{1} & \mathbf{x}_{2}\end{array}\right ) \left (\begin{array}{cc}\ensuremath{\operatorname*{Re}} \left (\mathbb{U} +\mathbb{V}\right ) &  -\ensuremath{\operatorname*{Im}} \left (\mathbb{U} +\mathbb{V}\right ) \\
\ensuremath{\operatorname*{Im}} \left (\mathbb{U} -\mathbb{V}\right ) & \ensuremath{\operatorname*{Re}} \left (\mathbb{U} -\mathbb{V}\right )\end{array}\right )\end{align}and
\begin{align} & \left (\begin{array}{cc}\mathbf{y}_{1} & \mathbf{y}_{2}\end{array}\right ) =K \widehat{\exp  \left \{\sum _{1 \leq j <k \leq r}\left (\gamma _{j k} a_{j}^{ \dag } a_{k}^{ \dag } +\overline{\gamma }_{j k} a_{k} a_{j}\right )\right \} \exp  \left \{\sum _{1 \leq j ,k \leq r}\lambda _{j k} a_{j}^{ \dag } a_{k}\right \}} K^{ \dag } \left (\begin{array}{cc}\mathbf{x}_{1} & \mathbf{x}_{2}\end{array}\right ) \\
 &  =\left (\begin{array}{cc}\mathbf{x}_{1} & \mathbf{x}_{2}\end{array}\right ) \left (\begin{array}{cc}\ensuremath{\operatorname*{Re}} \left (\mathbb{U} +\mathbb{V}\right ) &  -\ensuremath{\operatorname*{Im}} \left (\mathbb{U} +\mathbb{V}\right ) \\
\ensuremath{\operatorname*{Im}} \left (\mathbb{U} -\mathbb{V}\right ) & \ensuremath{\operatorname*{Re}} \left (\mathbb{U} -\mathbb{V}\right )\end{array}\right ) \\
 &  =K \widehat{\exp  \left \{\sum _{1 \leq j <k \leq r}\left (\gamma _{j k} a_{j}^{ \dag } a_{k}^{ \dag } +\overline{\gamma }_{j k} a_{k} a_{j}\right )\right \}} K^{ \dag } \left (\begin{array}{cc}\mathbf{x}_{1} & \mathbf{x}_{2}\end{array}\right ) =K \left (\begin{array}{cc}\mathbf{a}^{ \dag } & \mathbf{a}\end{array}\right ) \left (\begin{array}{cc}\mathbb{U} & \overline{\mathbb{V}} \\
\mathbb{V} & \overline{\mathbb{U}}\end{array}\right ) \\
 &  =\left (\begin{array}{cc}\mathbf{x}_{1} & \mathbf{x}_{2}\end{array}\right ) \mathbb{K} \left (\begin{array}{cc}\mathbb{U} & \overline{\mathbb{V}} \\
\mathbb{V} & \overline{\mathbb{U}}\end{array}\right ) \mathbb{K}^{ \dag }\text{.}\end{align}

\emph{Reiterating the preceeding}
\begin{align} & \exp  \left \{\frac{1}{2} \left (\mathbf{a}^{ \dag } \mathbf{\gamma } \mathbf{a}^{ \dag } +\mathbf{a} \overline{\mathbf{\gamma }} \mathbf{a}\right )\right \} =\exp  \left \{\frac{1}{2} \left (\begin{array}{cc}\mathbf{a}^{ \dag } & \mathbf{a}\end{array}\right ) \left (\begin{array}{cc}\mathbb{O} & \mathbf{\gamma } \\
\overline{\mathbf{\gamma }} & \mathbb{O}\end{array}\right ) \left (\begin{array}{c}\mathbf{a} \\
\mathbf{a}^{ \dag }\end{array}\right )\right \} =\left (\begin{array}{cc}\mathbf{a}^{ \dag } & \mathbf{a}\end{array}\right ) \left (\begin{array}{cc}\mathbb{U}_{\gamma } & \overline{\mathbb{V}}_{\gamma } \\
\mathbb{V}_{\gamma } & \overline{\mathbb{U}}_{\gamma }\end{array}\right ) \left (\begin{array}{c}\mathbf{a} \\
\mathbf{a}^{ \dag }\end{array}\right ) \\
 &  =\exp  \left \{\frac{1}{2} \left (\mathbf{x}_{1} \ensuremath{\operatorname*{Re}} \left (\mathbf{\gamma }\right ) \mathbf{x}_{1} +\mathbf{x}_{1} \ensuremath{\operatorname*{Im}} \left (\mathbf{\gamma }\right ) \mathbf{x}_{2} +\mathbf{x}_{2} \ensuremath{\operatorname*{Im}} \left (\mathbf{\gamma }\right ) \mathbf{x}_{1} -\mathbf{x}_{2} \ensuremath{\operatorname*{Re}} \left (\mathbf{\gamma }\right ) \mathbf{x}_{2}\right )\right \} \\
 &  =\exp  \left \{\frac{1}{2} \left (\begin{array}{cc}\mathbf{x}_{1} & \mathbf{x}_{2}\end{array}\right ) \left (\begin{array}{cc}\ensuremath{\operatorname*{Re}} \left (\mathbf{\gamma }\right ) & \ensuremath{\operatorname*{Im}} \left (\mathbf{\gamma }\right ) \\
\ensuremath{\operatorname*{Im}} \left (\mathbf{\gamma }\right ) &  -\ensuremath{\operatorname*{Re}} \left (\mathbf{\gamma }\right )\end{array}\right ) \left (\begin{array}{c}\mathbf{x}_{1} \\
\mathbf{x}_{2}\end{array}\right )\right \} \\
 &  =\left (\begin{array}{cc}\mathbf{x}_{1} & \mathbf{x}_{2}\end{array}\right ) \left (\begin{array}{cc}\ensuremath{\operatorname*{Re}} \left (\mathbb{U}_{\gamma } +\mathbb{V}_{\gamma }\right ) &  -\ensuremath{\operatorname*{Im}} \left (\mathbb{U}_{\gamma } +\mathbb{V}_{\gamma }\right ) \\
\ensuremath{\operatorname*{Im}} \left (\mathbb{U}_{\gamma } -\mathbb{V}_{\gamma }\right ) & \ensuremath{\operatorname*{Re}} \left (\mathbb{U}_{\gamma } -\mathbb{V}_{\gamma }\right )\end{array}\right ) \left (\begin{array}{c}\mathbf{x}_{1} \\
\mathbf{x}_{2}\end{array}\right )\text{.}\end{align}

The bases $\left \{\mathbf{x}_{1} ,\mathbf{x}_{2}\right \} ,\left \{\mathbf{y}_{1} ,\mathbf{y}_{2}\right \} ,\left \{\mathbf{a}^{ \dag } ,\mathbf{a}\right \}$ and $\left \{\mathbf{b}^{ \dag } ,\mathbf{b}\right \}$ are orthogonal but not normalized as
\begin{equation}\ensuremath{\operatorname*{tr}}\left \{x^{ \dag } x\right \} =\ensuremath{\operatorname*{tr}}\left \{y^{ \dag } y\right \} =\ensuremath{\operatorname*{tr}}\left \{a^{ \dag } a\right \} =\ensuremath{\operatorname*{tr}}\left \{b^{ \dag } b\right \} =\frac{1}{2}
\end{equation}but that does not alter the matrix representation relationships. 


\begin{remark}
\begin{align}\left (\begin{array}{cc}\mathbf{e}_{1} & \mathbf{e}_{2}\end{array}\right ) \left (\begin{array}{cc}\mathbb{A} & \mathbb{B} \\
\mathbb{C} & \mathbb{D}\end{array}\right ) \left (\begin{array}{c}\mathbf{e}_{1} \\
\mathbf{e}_{2}\end{array}\right ) &  \equiv \left (\begin{array}{cc}\left \vert \mathbf{e}_{1}\right ) & \left \vert \mathbf{e}_{2}\right )\end{array}\right ) \left (\begin{array}{cc}\mathbb{A} & \mathbb{B} \\
\mathbb{C} & \mathbb{D}\end{array}\right ) \left (\begin{array}{c}\left (\mathbf{e}_{1}\right \vert  \\
\left (\mathbf{e}_{2}\right \vert \end{array}\right ) =\left (\begin{array}{c}\left (\mathbf{e}_{1}\right \vert  \\
\left (\mathbf{e}_{2}\right \vert \end{array}\right )^{ \dag } \left (\begin{array}{cc}\mathbb{A} & \mathbb{B} \\
\mathbb{C} & \mathbb{D}\end{array}\right ) \left (\begin{array}{c}\left (\mathbf{e}_{1}\right \vert  \\
\left (\mathbf{e}_{2}\right \vert \end{array}\right ) \\
 &  =\left \vert \mathbf{e}_{1}\right )\mathbb{A} \left (\mathbf{e}_{1}\right \vert  +\left \vert \mathbf{e}_{1}\right ) \mathbb{B} \left (\mathbf{e}_{2}\right \vert  +\left \vert \mathbf{e}_{2}\right ) \mathbb{C} \left (\mathbf{e}_{1}\right \vert  +\left \vert \mathbf{e}_{2}\right ) \mathbb{D}\left (\mathbf{e}_{2}\right \vert \text{.}\end{align}
\end{remark}

Using the canonical form of $\gamma $
\begin{equation}\gamma  = \sum _{1 \leq j \leq s}\zeta _{j} \left (c_{j}^{ \dag } c_{j +s}^{ \dag } +c_{j} c_{j +s}\right )
\end{equation}one obtains
\begin{align} & \exp  \left \{\frac{1}{2} \left (\mathbf{c}^{ \dag } \mathbf{\Xi } \mathbf{c}^{ \dag } +\mathbf{c} \mathbf{\Xi } \mathbf{c}\right )\right \} =\exp  \left \{\frac{1}{2} \left (\begin{array}{cc}\mathbf{c}^{ \dag } & \mathbf{c}\end{array}\right ) \left (\begin{array}{cc}\mathbb{O} & \mathbf{\Xi } \\
\mathbf{\Xi } & \mathbb{O}\end{array}\right ) \left (\begin{array}{c}\mathbf{c} \\
\mathbf{c}^{ \dag }\end{array}\right )\right \} \\
 &  =\left (\begin{array}{cc}\mathbf{c}^{ \dag } & \mathbf{c}\end{array}\right ) \widetilde{\mathbb{T}}^{ \dag } \left (\begin{array}{cc}\mathbb{U}_{\gamma } & \overline{\mathbb{V}}_{\gamma } \\
\mathbb{V}_{\gamma } & \overline{\mathbb{U}}_{\gamma }\end{array}\right ) \widetilde{\mathbb{T}} \left (\begin{array}{c}\mathbf{c} \\
\mathbf{c}^{ \dag }\end{array}\right ) \\
 &  =\exp  \left \{\frac{1}{2} \left (\mathbf{y}_{1} \mathbf{\Xi } \mathbf{y}_{1} -\mathbf{y}_{2} \mathbf{\Xi } \mathbf{y}_{2}\right )\right \} =\exp  \left \{\frac{1}{2} \left (\begin{array}{cc}\mathbf{y}_{1} & \mathbf{y}_{2}\end{array}\right ) \left (\begin{array}{cc}\mathbf{\Xi } & \mathbb{O} \\
\mathbb{O} &  -\mathbf{\Xi }\end{array}\right ) \left (\begin{array}{c}\mathbf{y}_{1} \\
\mathbf{y}_{2}\end{array}\right )\right \} \\
 &  =\left (\begin{array}{cc}\mathbf{y}_{1} & \mathbf{y}_{2}\end{array}\right ) \widetilde{\mathbb{T}}^{ \dag } \left (\begin{array}{cc}\ensuremath{\operatorname*{Re}} \left (\mathbb{U}_{\gamma } +\mathbb{V}_{\gamma }\right ) &  -\ensuremath{\operatorname*{Im}} \left (\mathbb{U}_{\gamma } +\mathbb{V}_{\gamma }\right ) \\
\ensuremath{\operatorname*{Im}} \left (\mathbb{U}_{\gamma } -\mathbb{V}_{\gamma }\right ) & \ensuremath{\operatorname*{Re}} \left (\mathbb{U}_{\gamma } -\mathbb{V}_{\gamma }\right )\end{array}\right ) \widetilde{\mathbb{T}} \left (\begin{array}{c}\mathbf{y}_{1} \\
\mathbf{y}_{2}\end{array}\right ) \\
 &  =\left (\begin{array}{cc}\mathbf{y}_{1} & \mathbf{y}_{2}\end{array}\right ) \left (\begin{array}{cc}\mathbb{T}^{ \dag } \ensuremath{\operatorname*{Re}} \left (\mathbb{U}_{\gamma } +\mathbb{V}_{\gamma }\right ) \mathbb{T} &  -\mathbb{T}^{ \dag } \ensuremath{\operatorname*{Im}} \left (\mathbb{U}_{\gamma } +\mathbb{V}_{\gamma }\right ) \mathbb{T} \\
\mathbb{T}^{ \dag } \ensuremath{\operatorname*{Im}} \left (\mathbb{U}_{\gamma } -\mathbb{V}_{\gamma }\right ) \mathbb{T} & \mathbb{T}^{ \dag } \ensuremath{\operatorname*{Re}} \left (\mathbb{U}_{\gamma } -\mathbb{V}_{\gamma }\right ) \mathbb{T}\end{array}\right ) \left (\begin{array}{c}\mathbf{y}_{1} \\
\mathbf{y}_{2}\end{array}\right )\text{,}\end{align}where
\begin{equation}\exp  \left \{\frac{1}{2} \left (\mathbf{c}^{ \dag } \mathbf{\Xi } \mathbf{c}^{ \dag } +\mathbf{c} \mathbf{\Xi } \mathbf{c}\right )\right \} =\exp  \left \{\frac{1}{2} \left (\begin{array}{cc}\mathbf{c}^{ \dag } & \mathbf{c}\end{array}\right ) \left (\begin{array}{cc}\mathbb{O} & \left (\begin{array}{cc}\mathbf{0} & \mathbf{\zeta } \\
 -\mathbf{\zeta } & \mathbf{0}\end{array}\right ) \\
\left (\begin{array}{cc}\mathbf{0} & \mathbf{\zeta } \\
 -\mathbf{\zeta } & \mathbf{0}\end{array}\right ) & \mathbb{O}\end{array}\right ) \left (\begin{array}{c}\mathbf{c} \\
\mathbf{c}^{ \dag }\end{array}\right )\right \}
\end{equation}and
\begin{equation}\widetilde{\mathbb{T}} =\left (\begin{array}{cc}\mathbb{T} & \mathbb{O} \\
\mathbb{O} & \overline{\mathbb{T}}\end{array}\right )\text{.}
\end{equation}

In particular
\begin{equation}\exp  \left \{\frac{1}{2} \left (\mathbf{c}^{ \dag } \mathbf{\Xi } \mathbf{c}^{ \dag } +\mathbf{c} \mathbf{\Xi } \mathbf{c}\right )\right \} =\exp  \left \{ \sum _{1 \leq j \leq s}\zeta _{j} \left (c_{j}^{ \dag } c_{j +s}^{ \dag } +c_{j} c_{j +s}\right )\right \} = \tprod _{1 \leq j \leq s}\exp  \left \{\zeta _{j} \left (c_{j}^{ \dag } c_{j +s}^{ \dag } +c_{j} c_{j +s}\right )\right \}\text{.}
\end{equation}Expanding each exponential gives
\begin{align}\exp  \left \{\zeta _{j} \left (c_{j}^{ \dag } c_{j +s}^{ \dag } +c_{j} c_{j +s}\right )\right \} &  =I_{\mathcal{F}} +\zeta _{j} \left (c_{j}^{ \dag } c_{j +s}^{ \dag } +c_{j} c_{j +s}\right ) +\frac{1}{2} \zeta _{j}^{2} \left (c_{j}^{ \dag } c_{j +s}^{ \dag } +c_{j} c_{j +s}\right )^{2} \\
 &  +\frac{1}{3 !} \zeta _{j}^{3} \left (c_{j}^{ \dag } c_{j +s}^{ \dag } +c_{j} c_{j +s}\right )^{3} +\cdots \end{align}$\lceil $ Anti - self adjointness is displayed by
\begin{equation}\left (c_{j}^{ \dag } c_{j +s}^{ \dag } +c_{j} c_{j +s}\right )^{ \dag } =c_{j +s} c_{j} +c_{j +s}^{ \dag } c_{j}^{ \dag } = -\left (c_{j}^{ \dag } c_{j +s}^{ \dag } +c_{j} c_{j +s}\right )
\end{equation}$\rfloor $ 

$\lceil $ The squared term gives
\begin{align} & \frac{1}{2} \zeta _{j}^{2} \left (c_{j}^{ \dag } c_{j +s}^{ \dag } +c_{j} c_{j +s}\right )^{2} =\frac{1}{2} \zeta _{j}^{2} \left (c_{j}^{ \dag } c_{j +s}^{ \dag } +c_{j} c_{j +s}\right ) \left (c_{j}^{ \dag } c_{j +s}^{ \dag } +c_{j} c_{j +s}\right ) \\
 &  =\frac{1}{2} \zeta _{j}^{2} \left (c_{j}^{ \dag } c_{j +s}^{ \dag } c_{j} c_{j +s} +c_{j} c_{j +s} c_{j}^{ \dag } c_{j +s}^{ \dag }\right ) = -\frac{1}{2} \zeta _{j}^{2} \left (c_{j}^{ \dag } c_{j} c_{j +s}^{ \dag } c_{j +s} +c_{j} c_{j}^{ \dag } c_{j +s} c_{j +s}^{ \dag }\right )\text{,}\end{align}while the cubic term gives
\begin{align} & \frac{1}{3 !} \zeta _{j}^{3} \left (c_{j}^{ \dag } c_{j +s}^{ \dag } +c_{j} c_{j +s}\right )^{3} = -\frac{1}{3 !} \zeta _{j}^{3} \left (c_{j}^{ \dag } c_{j} c_{j +s}^{ \dag } c_{j +s} +c_{j} c_{j}^{ \dag } c_{j +s} c_{j +s}^{ \dag }\right ) \left (c_{j}^{ \dag } c_{j +s}^{ \dag } +c_{j} c_{j +s}\right ) \\
 &  = -\frac{1}{3 !} \zeta _{j}^{3} \left (c_{j}^{ \dag } c_{j} c_{j +s}^{ \dag } c_{j +s} \left (c_{j}^{ \dag } c_{j +s}^{ \dag } +c_{j} c_{j +s}\right ) +c_{j} c_{j}^{ \dag } c_{j +s} c_{j +s}^{ \dag } \left (c_{j}^{ \dag } c_{j +s}^{ \dag } +c_{j} c_{j +s}\right )\right ) \\
 &  = -\frac{1}{3 !} \zeta _{j}^{3} \left (c_{j}^{ \dag } c_{j} c_{j +s}^{ \dag } c_{j +s} c_{j}^{ \dag } c_{j +s}^{ \dag } +c_{j}^{ \dag } c_{j} c_{j +s}^{ \dag } c_{j +s} c_{j} c_{j +s} +c_{j} c_{j}^{ \dag } c_{j +s} c_{j +s}^{ \dag } c_{j}^{ \dag } c_{j +s}^{ \dag } +c_{j} c_{j}^{ \dag } c_{j +s} c_{j +s}^{ \dag } c_{j} c_{j +s}\right ) \\
 &  = -\frac{1}{3 !} \zeta _{j}^{3} \left (c_{j}^{ \dag } c_{j} c_{j}^{ \dag } c_{j +s}^{ \dag } c_{j +s} c_{j +s}^{ \dag } +c_{j} c_{j}^{ \dag } c_{j} c_{j +s} c_{j +s}^{ \dag } c_{j +s}\right ) \\
 &  = -\frac{1}{3 !} \zeta _{j}^{3} \left (c_{j}^{ \dag } \left (1 -c_{j}^{ \dag } c_{j}\right ) c_{j +s}^{ \dag } c_{j +s} c_{j +s}^{ \dag } +c_{j} \left (1 -c_{j} c_{j}^{ \dag }\right ) c_{j +s} c_{j +s}^{ \dag } c_{j +s}\right ) \\
 &  = -\frac{1}{3 !} \zeta _{j}^{3} \left (c_{j}^{ \dag } \left (1 -c_{j +s} c_{j +s}^{ \dag }\right ) c_{j +s}^{ \dag } +c_{j} c_{j +s} \left (1 -c_{j +s} c_{j +s}^{ \dag }\right )\right ) = -\frac{1}{3 !} \zeta _{j}^{3} \left (c_{j}^{ \dag } c_{j +s}^{ \dag } +c_{j} c_{j +s}\right )\text{.}\end{align}$\rfloor $ 

Thus
\begin{align} & \exp  \left \{\zeta _{j} \left (c_{j}^{ \dag } c_{j +s}^{ \dag } +c_{j} c_{j +s}\right )\right \} =\zeta _{j} \left (c_{j}^{ \dag } c_{j +s}^{ \dag } +c_{j} c_{j +s}\right ) -\frac{1}{3 !} \zeta _{j}^{3} \left (c_{j}^{ \dag } c_{j +s}^{ \dag } -c_{j} c_{j +s}\right ) +\frac{1}{5 !} \zeta _{j}^{5} \left (c_{j}^{ \dag } c_{j +s}^{ \dag } -c_{j} c_{j +s}\right ) +\cdots  \\
 &  +I_{\mathcal{F}} -\frac{1}{2} \zeta _{j}^{2} \left (c_{j}^{ \dag } c_{j} c_{j +s}^{ \dag } c_{j +s} +c_{j} c_{j}^{ \dag } c_{j +s} c_{j +s}^{ \dag }\right ) +\frac{1}{4 !} \zeta _{j}^{4} \left (c_{j}^{ \dag } c_{j} c_{j +s}^{ \dag } c_{j +s} +c_{j} c_{j}^{ \dag } c_{j +s} c_{j +s}^{ \dag }\right ) +\cdots  \\
 &  =\left (c_{j}^{ \dag } c_{j} c_{j +s}^{ \dag } c_{j +s} +c_{j} c_{j}^{ \dag } c_{j +s} c_{j +s}^{ \dag }\right ) \cos  \left (\zeta _{j}\right ) +\left (c_{j}^{ \dag } c_{j +s}^{ \dag } +c_{j} c_{j +s}\right ) \sin  \left (\zeta _{j}\right ) \\
 &  =\left (n_{j} n_{j +s} +h_{j} h_{j +s}\right ) \cos  \left (\zeta _{j}\right ) +\left (c_{j}^{ \dag } c_{j +s}^{ \dag } +c_{j} c_{j +s}\right ) \sin  \left (\zeta _{j}\right )\text{.}\end{align}This can be expressed generally as
\begin{align} & \left (\begin{array}{cccc}\mathbf{n}_{1} & \mathbf{n}_{2} & \mathbf{h}_{1} & \mathbf{h}_{2}\end{array}\right ) \left (\begin{array}{cccc}0 & \cos  \left (\mathbf{\zeta }\right ) & \mathbf{0} & \mathbf{0} \\
\cos  \left (\mathbf{\zeta }\right ) & 0 & \mathbf{0} & \mathbf{0} \\
\mathbf{0} & \mathbf{0} & 0 & \cos  \left (\mathbf{\zeta }\right ) \\
\mathbf{0} & \mathbf{0} & \cos  \left (\mathbf{\zeta }\right ) & 0\end{array}\right ) \left (\begin{array}{c}\mathbf{n}_{1} \\
\mathbf{n}_{2} \\
\mathbf{h}_{1} \\
\mathbf{h}_{2}\end{array}\right ) = \\
 & \left (\begin{array}{cccc}\mathbf{n}_{2} \cos  \left (\mathbf{\zeta }\right ) & \mathbf{n}_{1} \cos  \left (\mathbf{\zeta }\right ) & \mathbf{h}_{2} \cos  \left (\mathbf{\zeta }\right ) & \mathbf{h}_{1} \cos  \left (\mathbf{\zeta }\right )\end{array}\right ) \left (\begin{array}{c}\mathbf{n}_{1} \\
\mathbf{n}_{2} \\
\mathbf{h}_{1} \\
\mathbf{h}_{2}\end{array}\right ) \\
 &  = \langle \mathbf{n}_{1}\left \vert \mathbf{n}_{2}\right . \rangle  \cos  \left (\mathbf{\zeta }\right ) + \langle \mathbf{n}_{1}\left \vert \mathbf{n}_{2}\right . \rangle  \cos  \left (\mathbf{\zeta }\right ) + \langle \mathbf{h}_{1}\left \vert \mathbf{h}_{2}\right . \rangle  \cos  \left (\mathbf{\zeta }\right ) + \langle \mathbf{h}_{1}\left \vert \mathbf{h}_{2}\right . \rangle  \cos  \left (\mathbf{\zeta }\right ) \\
 &  =2 \sum _{1 \leq j \leq s}\left (n_{j} n_{j +s} +h_{j} h_{j +s}\right ) \cos  \left (\zeta _{j}\right )\text{,}\end{align}where
\begin{align}\mathbf{n}_{1} &  =\left (n_{1} ,\cdots  ,n_{s}\right ) \\
\mathbf{n}_{2} &  =\left (n_{s +1} ,\cdots  ,n_{2 s}\right )\end{align}
\begin{align}\mathbf{h}_{1} &  =\left (h_{1} ,\cdots  ,h_{s}\right ) \\
\mathbf{h}_{2} &  =\left (h_{s +1} ,\cdots  ,h_{2 s}\right )\end{align}and
\begin{align} & \left (\begin{array}{cccc}\mathbf{c}_{1}^{ \dag } & \mathbf{c}_{2}^{ \dag } & \mathbf{c}_{1} & \mathbf{c}_{2}\end{array}\right ) \left (\begin{array}{cccc}0 & 0 & \mathbf{0} & \sin  \left (\mathbf{\zeta }\right ) \\
0 & 0 &  -\sin  \left (\mathbf{\zeta }\right ) & \mathbf{0} \\
\mathbf{0} &  -\sin  \left (\mathbf{\zeta }\right ) & 0 & 0 \\
\sin  \left (\mathbf{\zeta }\right ) & \mathbf{0} & 0 & 0\end{array}\right ) \left (\begin{array}{c}\mathbf{c}_{1} \\
\mathbf{c}_{2} \\
\mathbf{c}_{1}^{ \dag } \\
\mathbf{c}_{2}^{ \dag }\end{array}\right ) = \\
 & \left (\begin{array}{cccc}\mathbf{c}_{2} \sin  \left (\mathbf{\zeta }\right ) &  -\mathbf{c}_{1} \sin  \left (\mathbf{\zeta }\right ) &  -\mathbf{c}_{2}^{ \dag } \sin  \left (\mathbf{\zeta }\right ) & \mathbf{c}_{1}^{ \dag } \sin  \left (\mathbf{\zeta }\right )\end{array}\right ) \left (\begin{array}{c}\mathbf{c}_{1} \\
\mathbf{c}_{2} \\
\mathbf{c}_{1}^{ \dag } \\
\mathbf{c}_{2}^{ \dag }\end{array}\right ) \\
 &  =\sum _{1 \leq j \leq s}\left \{c_{j +s} c_{j} \sin  \left (\zeta _{j}\right ) -c_{j} c_{j +s} \sin  \left (\zeta _{j}\right ) -c_{j +s}^{ \dag } c_{j}^{ \dag } \sin  \left (\zeta _{j}\right ) +c_{j}^{ \dag } c_{j +s}^{ \dag } \sin  \left (\zeta _{j}\right )\right \} \\
 &  =2 \sum _{1 \leq j \leq s}\left (c_{j}^{ \dag } c_{j +s}^{ \dag } -c_{j} c_{j +s}\right ) \sin  \left (\zeta _{j}\right )\text{.}\end{align}$\rfloor $ 

Thus if
\begin{equation*}\zeta _{j} =\left (2 n_{j} +1\right ) \frac{\pi }{2}
\end{equation*}then
\begin{align}\sin  \left (\zeta _{j}\right ) &  = \pm 1 \\
\cos  \left (\zeta _{j}\right ) &  =0\end{align}and
\begin{equation}\exp  \left \{\left (2 n_{j} +1\right ) \frac{\pi }{2} \left (c_{j}^{ \dag } c_{j +s}^{ \dag } +c_{j} c_{j +s}\right )\right \} = \pm \left (c_{j}^{ \dag } c_{j +s}^{ \dag } +c_{j} c_{j +s}\right )\text{.}
\end{equation}Letting $n_{j} =0$ one obtains
\begin{align} & \exp  \left \{\frac{\pi }{2} \frac{1}{2} \left (\begin{array}{cccc}c_{j}^{ \dag } & c_{j +s}^{ \dag } & c_{j} & c_{j +s}\end{array}\right ) \left (\begin{array}{cc}\begin{array}{cc}0 & 0 \\
0 & 0\end{array} & \begin{array}{cc}0 & 1 \\
 -1 & 0\end{array} \\
\begin{array}{cc}0 & 1 \\
 -1 & 0\end{array} & \begin{array}{cc}0 & 0 \\
0 & 0\end{array}\end{array}\right ) \left (\begin{array}{c}c_{j} \\
c_{j +s} \\
c_{j}^{ \dag } \\
c_{j +s}^{ \dag }\end{array}\right )\right \} \\
 &  =\frac{1}{2} \left (\begin{array}{cccc}c_{j}^{ \dag } & c_{j +s}^{ \dag } & c_{j} & c_{j +s}\end{array}\right ) \left (\begin{array}{cc}\begin{array}{cc}0 & 0 \\
0 & 0\end{array} & \begin{array}{cc}0 & 1 \\
 -1 & 0\end{array} \\
\begin{array}{cc}0 & 1 \\
 -1 & 0\end{array} & \begin{array}{cc}0 & 0 \\
0 & 0\end{array}\end{array}\right ) \left (\begin{array}{c}c_{j} \\
c_{j +s} \\
c_{j}^{ \dag } \\
c_{j +s}^{ \dag }\end{array}\right )\text{.}\end{align}
\begin{equation} =\frac{1}{2} \left (\begin{array}{cccc} -c_{j +s} & c_{j} &  -c_{j +s}^{ \dag } & c_{j}^{ \dag }\end{array}\right ) \left (\begin{array}{c}c_{j} \\
c_{j +s} \\
c_{j}^{ \dag } \\
c_{j +s}^{ \dag }\end{array}\right ) =\frac{1}{2} \left ( -c_{j +s} c_{j} +c_{j} c_{j +s} -c_{j +s}^{ \dag } c_{j}^{ \dag } +c_{j}^{ \dag } c_{j +s}^{ \dag }\right ) =c_{j}^{ \dag } c_{j +s}^{ \dag } +c_{j} c_{j +s}
\end{equation}$\rfloor $ 

In the maximal particle hole case this gives
\begin{align} & \exp  \left \{\frac{\pi }{2} \frac{1}{2} \left (\begin{array}{cc}\mathbf{c}^{ \dag } & \mathbf{c}\end{array}\right ) \left (\begin{array}{cc}\mathbb{O}_{r} & \left (\begin{array}{cc}\mathbb{O}_{s} & \mathbb{I}_{s} \\
 -\mathbb{I}_{s} & \mathbb{O}_{s}\end{array}\right ) \\
\left (\begin{array}{cc}\mathbf{0} & \mathbb{I}_{s} \\
 -\mathbb{I}_{s} & \mathbf{0}\end{array}\right ) & \mathbb{O}_{r}\end{array}\right ) \left (\begin{array}{c}\mathbf{c} \\
\mathbf{c}^{ \dag }\end{array}\right )\right \} \\
 &  =\frac{1}{2} \left (\begin{array}{cc}\mathbf{c}^{ \dag } & \mathbf{c}\end{array}\right ) \left (\begin{array}{cc}\mathbb{O}_{r} & \left (\begin{array}{cc}\mathbb{O}_{s} & \mathbb{I}_{s} \\
 -\mathbb{I}_{s} & \mathbb{O}_{s}\end{array}\right ) \\
\left (\begin{array}{cc}\mathbf{0} & \mathbb{I}_{s} \\
 -\mathbb{I}_{s} & \mathbf{0}\end{array}\right ) & \mathbb{O}_{r}\end{array}\right ) \left (\begin{array}{c}\mathbf{c} \\
\mathbf{c}^{ \dag }\end{array}\right )\text{.}\end{align}

If $\overline{\mathbb{U}}^{ -1}$ exists then the coherent state can be expressed as
\begin{equation}\mathcal{N} \left (\phi  \left (c\right )\right ) {\displaystyle\prod _{1 \leq j \leq s}}\exp  \left \{c_{j} a_{j}^{ \dag } a_{j +s}^{ \dag }\right \}\vert \phi  \rangle \text{,}
\end{equation}where
\begin{equation}\mathbf{C} =\left (\begin{array}{ccc}c_{1} & 0 & 0 \\
0 & \ddots  & 0 \\
0 & 0 & c_{s}\end{array}\right ) =\frac{\mathbb{Z}_{c}}{\sqrt{\left (\mathbb{I}_{r} -\mathbb{Z}_{c}^{ \dag } \mathbb{Z}_{c}\right )}} =\overline{\mathbb{V}}_{c} \overline{\mathbb{U}}_{c}^{ -1}\text{.}
\end{equation}

If $\overline{\mathbb{U}}_{c}^{ -1}$ \emph{does not exist} then the unnormalized coherent state can be expressed more generally as
\begin{equation} \tprod _{1 \leq j \leq p}\exp  \left \{\frac{\pi }{2} \left (c_{j}^{ \dag } c_{j +s}^{ \dag } +c_{j} c_{j +s}\right )\right \} \exp  \left \{\sum _{p +1 \leq j \leq s}c_{j} \left (g\right ) c_{j}^{ \dag } c_{j +s}^{ \dag }\right \}\vert \phi  \rangle  =\exp  \left \{\sum _{p +1 \leq j \leq s}c_{j} \left (g\right ) c_{j}^{ \dag } c_{j +s}^{ \dag }\right \}  \tprod _{1 \leq j \leq p}c_{j}^{ \dag } c_{j +s}^{ \dag }\vert \phi  \rangle \text{.}
\end{equation}

If $\overline{\mathbb{U}}_{c}^{ -1}$ \emph{does exist} one can expand dressed vacuum states as
\begin{align}\vert \phi  \left (g\right ) \rangle  &  =\mathcal{N} \left (\phi  \left (g\right )\right ) \exp  \left \{\sum _{1 \leq j \leq s}c_{j} \left (g\right ) a_{j}^{ \dag } a_{j +s}^{ \dag }\right \}\vert \phi  \rangle  =\mathcal{N} \left (\phi  \left (g\right )\right ) {\displaystyle\prod _{1 \leq j \leq s}}\exp  \left \{c_{j} \left (g\right ) a_{j}^{ \dag } a_{j +s}^{ \dag }\right \}\vert \phi  \rangle  \\
 &  =\mathcal{N} \left (\phi  \left (g\right )\right ) {\displaystyle\prod _{1 \leq j \leq s}}\left (I_{\mathcal{F}} +c_{j} \left (g\right ) a_{j}^{ \dag } a_{j +s}^{ \dag }\right )\vert \phi  \rangle  ={\displaystyle\prod _{1 \leq j \leq s}}\left (\mathfrak{d}_{j} \left (g\right ) I_{\mathcal{F}} +\mathfrak{c}_{j} \left (g\right ) a_{j}^{ \dag } a_{j +s}^{ \dag }\right )\vert \phi  \rangle  \\
 &  ={\displaystyle\prod _{1 \leq j \leq s}}\mathfrak{d}_{j} \left (g\right ) \left (I_{\mathcal{F}} +\frac{\mathfrak{c}_{j} \left (g\right )}{\mathfrak{d}_{j} \left (g\right )} a_{j}^{ \dag } a_{j +s}^{ \dag }\right )\vert \phi  \rangle  =\left ({\displaystyle\prod _{1 \leq j \leq s}}\mathfrak{d}_{j} \left (g\right )\right ) {\displaystyle\prod _{1 \leq j \leq s}}\left (I_{\mathcal{F}} +\frac{\mathfrak{c}_{j} \left (g\right )}{\mathfrak{d}_{j} \left (g\right )} a_{j}^{ \dag } a_{j +s}^{ \dag }\right )\vert \phi  \rangle \text{.}\end{align}This shows that
\begin{equation}\frac{\mathfrak{c}_{k} \left (g\right )}{\mathfrak{d}_{k} \left (g\right )} =c_{k} \left (g\right )\text{and}\mathcal{N} \left (\phi  \left (g\right )\right ) ={\displaystyle\prod _{1 \leq j \leq s}}\mathfrak{d}_{j} \left (g\right )\text{.}
\end{equation}The normalized wedge product square of dressed vacuum states is given by in terms of naked vacuum states as
\begin{align}\vert \phi  \left (g\right ) \wedge \phi  \left (g\right ) \rangle  &  =\mathcal{N}^{2}\vert \exp  \left \{\sum _{1 \leq j \leq s}c_{j} a_{j}^{ \dag } a_{j +s}^{ \dag }\right \} \phi  \wedge \exp  \left \{\sum _{1 \leq j \leq s}c_{j} a_{j}^{ \dag } a_{j +s}^{ \dag }\right \} \phi  \rangle  \\
 &  =\mathcal{N}^{2} \exp  \left \{\sum _{1 \leq j \leq s}c_{j} a_{j}^{ \dag } a_{j +s}^{ \dag }\right \} \boxtimes \exp  \left \{\sum _{1 \leq j \leq s}c_{j} a_{j}^{ \dag } a_{j +s}^{ \dag }\right \}\vert \phi  \wedge \phi  \rangle  \\
 & \overset{\text{?}}{ =}\mathcal{N}^{2} \exp  \left \{\sum _{1 \leq j \leq s}c_{j} a_{j}^{ \dag } a_{j +s}^{ \dag }\right \} \exp  \left \{\sum _{1 \leq j \leq s}c_{j} a_{j}^{ \dag } a_{j +s}^{ \dag }\right \}\vert \phi  \wedge \phi  \rangle  \\
 &  =\mathcal{N}^{2} \exp  \left \{2 \sum _{1 \leq j \leq s}c_{j} a_{j}^{ \dag } a_{j +s}^{ \dag }\right \}\vert \phi  \wedge \phi  \rangle  \\
 &  =\mathcal{N}^{2} \exp  \left \{2 \sum _{1 \leq j \leq s}c_{j} a_{j}^{ \dag } a_{j +s}^{ \dag }\right \}\vert \phi  \rangle \end{align}and is not equal to the dressed vacuum state. 

$\lceil $
\begin{align}A\vert \varphi _{1} \rangle  &  =\vert \psi _{1} \rangle  \\
B\vert \varphi _{1} \rangle  &  =\vert \psi _{2} \rangle  \\
 \langle \left .\psi _{1}\right \vert \psi _{2} \rangle  &  =0 \\
\vert \psi _{1} \wedge \psi _{2} \rangle  &  =\vert \vert A \varphi _{1} \rangle  \wedge \vert B \varphi _{1} \rangle  \rangle  \neq \left (A \boxtimes B\right )\vert \varphi _{1} \wedge \varphi _{1} \rangle  =0\end{align}
\begin{align}A \boxtimes B &  =\mathcal{A}_{2} \left (A \otimes B\right ) \mathcal{A}_{2} \\
\psi _{1} \wedge \psi _{2} &  =\mathcal{A}_{2} \left (\psi _{1} \otimes \psi _{2}\right ) \\
\left (A \otimes B\right ) \left (\psi _{1} \otimes \psi _{2}\right ) &  =A \psi _{1} \otimes B \psi _{2} \\
\left (A \otimes B\right ) \left (\varphi _{1} \otimes \varphi _{1}\right ) &  =A \varphi _{1} \otimes B \varphi _{1} \\
\left (A \boxtimes B\right )\vert \varphi _{1} \wedge \varphi _{1} \rangle  &  =\mathcal{A}_{2} \left (A \otimes B\right ) \mathcal{A}_{2} \mathcal{A}_{2} \left (\varphi _{1} \otimes \varphi _{1}\right ) =\mathcal{A}_{2} \left (A \otimes B\right ) \mathcal{A}_{2} \left (\varphi _{1} \otimes \varphi _{1}\right )\end{align}
\begin{align}\vert \vert A \psi _{1} \rangle  \wedge \vert B \psi _{2} \rangle  \rangle  &  =\mathcal{A}_{2} \left (\left (A \psi _{1}\right ) \otimes \left (B \psi _{2}\right )\right ) =\mathcal{A}_{2} \left (\left (A \otimes B\right ) \left (\psi _{1} \otimes \psi _{2}\right )\right ) \\
\left (A \boxtimes B\right )\vert \psi _{1} \wedge \psi _{2} \rangle  &  =\mathcal{A}_{2} \left (A \otimes B\right ) \mathcal{A}_{2} \mathcal{A}_{2} \left (\psi _{1} \otimes \psi _{2}\right ) =\mathcal{A}_{2} \left (A \otimes B\right ) \mathcal{A}_{2} \left (\psi _{1} \otimes \psi _{2}\right )\end{align}\emph{NOTE!!!!}
\begin{equation}\mathcal{A}_{2} \left (\left (A \otimes B\right ) \left (\psi _{1} \otimes \psi _{2}\right )\right ) =\mathcal{A}_{2} \left (A \psi _{1} \otimes B \psi _{2}\right ) \neq \mathcal{A}_{2} \left (A \otimes B\right ) \mathcal{A}_{2} \left (\psi _{1} \otimes \psi _{2}\right )\text{,}
\end{equation}and is not equal in general even when
\begin{equation}A =B\text{.}
\end{equation}This can be observed by considering
\begin{align} & \left (A \otimes B\right ) \left (\psi _{1} \otimes \psi _{2}\right ) =A \psi _{1} \otimes B \psi _{2} \\
 &  =\left (\mathcal{A}_{2} \left (A \otimes B\right ) \mathcal{A}_{2} +\mathcal{A}_{2} \left (A \otimes B\right ) \mathcal{A}_{2}^{\bot } +\mathcal{A}_{2}^{\bot } \left (A \otimes B\right ) \mathcal{A}_{2} +\mathcal{A}_{2}^{\bot } \left (A \otimes B\right ) \mathcal{A}_{2}^{\bot }\right ) \left (\mathcal{A}_{2} +\mathcal{A}_{2}^{\bot }\right ) \left (\psi _{1} \otimes \psi _{2}\right ) \\
 &  =\left (\mathcal{A}_{2} \left (A \otimes B\right ) \mathcal{A}_{2} +\mathcal{A}_{2} \left (A \otimes B\right ) \mathcal{S}_{2} +\mathcal{S}_{2} \left (A \otimes B\right ) \mathcal{A}_{2} +\mathcal{S}_{2} \left (A \otimes B\right ) \mathcal{S}_{2}\right ) \left (\mathcal{A}_{2} +\mathcal{S}_{2}\right ) \left (\psi _{1} \otimes \psi _{2}\right ) \\
 &  =\mathcal{A}_{2} \left (A \otimes B\right ) \mathcal{A}_{2} \mathcal{A}_{2} \left (\psi _{1} \otimes \psi _{2}\right ) +\mathcal{A}_{2} \left (A \otimes B\right ) \mathcal{A}_{2} \mathcal{S}_{2} \left (\psi _{1} \otimes \psi _{2}\right ) \\
 &  +\mathcal{A}_{2} \left (A \otimes B\right ) \mathcal{S}_{2} \mathcal{A}_{2} \left (\psi _{1} \otimes \psi _{2}\right ) +\mathcal{A}_{2} \left (A \otimes B\right ) \mathcal{S}_{2} \mathcal{S}_{2} \left (\psi _{1} \otimes \psi _{2}\right ) \\
 &  +\mathcal{S}_{2} \left (A \otimes B\right ) \mathcal{A}_{2} \mathcal{A}_{2} \left (\psi _{1} \otimes \psi _{2}\right ) +\mathcal{S}_{2} \left (A \otimes B\right ) \mathcal{A}_{2} \mathcal{S}_{2} \left (\psi _{1} \otimes \psi _{2}\right ) \\
 &  +\mathcal{S}_{2} \left (A \otimes B\right ) \mathcal{S}_{2} \mathcal{A}_{2} \left (\psi _{1} \otimes \psi _{2}\right ) +\mathcal{S}_{2} \left (A \otimes B\right ) \mathcal{S}_{2} \mathcal{S}_{2} \left (\psi _{1} \otimes \psi _{2}\right ) \\
 &  =\mathcal{A}_{2} \left (A \otimes B\right ) \mathcal{A}_{2} \left (\psi _{1} \otimes \psi _{2}\right ) +\mathcal{A}_{2} \left (A \otimes B\right ) \mathcal{S}_{2} \left (\psi _{1} \otimes \psi _{2}\right ) \\
 &  +\mathcal{S}_{2} \left (A \otimes B\right ) \mathcal{A}_{2} \left (\psi _{1} \otimes \psi _{2}\right ) +\mathcal{S}_{2} \left (A \otimes B\right ) \mathcal{S}_{2} \left (\psi _{1} \otimes \psi _{2}\right )\end{align}
\begin{align*} &  =\left (A \boxtimes B\right ) \vert \psi _{1} \wedge \psi _{2} \rangle  +\mathcal{A}_{2} \left (A \otimes B\right ) \mathcal{S}_{2} \left (\psi _{1} \otimes \psi _{2}\right ) +\mathcal{S}_{2} \left (A \otimes B\right ) \mathcal{A}_{2} \left (\psi _{1} \otimes \psi _{2}\right ) +\mathcal{S}_{2} \left (A \otimes B\right ) \mathcal{S}_{2} \left (\psi _{1} \otimes \psi _{2}\right ) \\
 &  =\left (A \boxtimes B\right ) \vert \psi _{1} \wedge \psi _{2} \rangle  +\left (A \vee B\right ) \left (\psi _{1} \vee \psi _{2}\right ) +\mathcal{A}_{2} \left (A \otimes B\right ) \mathcal{S}_{2} \left (\psi _{1} \otimes \psi _{2}\right ) +\mathcal{S}_{2} \left (A \otimes B\right ) \mathcal{A}_{2} \left (\psi _{1} \otimes \psi _{2}\right ) \\
 &  =\left (A \boxtimes B\right ) \vert \psi _{1} \wedge \psi _{2} \rangle  +\left (A \vee B\right ) \left (\psi _{1} \vee \psi _{2}\right ) +\mathcal{A}_{2} \left (A \otimes B\right ) \left (\psi _{1} \vee \psi _{2}\right ) +\mathcal{S}_{2} \left (A \otimes B\right ) \left (\psi _{1} \wedge \psi _{2}\right )\text{.}\end{align*}
\begin{align}\left (A \otimes A\right ) \left (\psi _{1} \otimes \psi _{2}\right ) &  =A \psi _{1} \otimes A \psi _{2} \\
\mathcal{S}_{2} \left (A \psi _{1} \otimes A \psi _{2}\right ) &  =\mathcal{S}_{2} \left (A \psi _{2} \otimes A \psi _{1}\right )\end{align}$\rfloor $ 

$\lceil $
\begin{align} & \mathcal{N} \left (\phi  \left (g\right )\right ) {\displaystyle\prod _{1 \leq j \leq s}}\exp  \left \{c_{j} \left (g\right ) a_{j}^{ \dag } a_{j +s}^{ \dag }\right \}\vert \phi  \rangle  =\mathcal{N} \left (\phi  \left (g\right )\right ) \exp  \left \{c_{1} \left (g\right ) a_{1}^{ \dag } a_{1 +s}^{ \dag }\right \} \cdots  \exp  \left \{c_{s} \left (g\right ) a_{s}^{ \dag } a_{s +s}^{ \dag }\right \}\vert \phi  \rangle  \\
 &  =\mathcal{N} \left (\phi  \left (g\right )\right ) \left (I_{\mathcal{F}} +c_{1} \left (g\right ) a_{1}^{ \dag } a_{1 +s}^{ \dag }\right ) \cdots  \left (I_{\mathcal{F}} +c_{s} \left (g\right ) a_{s}^{ \dag } a_{s +s}^{ \dag }\right )\vert \phi  \rangle  \\
 &  =\mathcal{N} \left (\phi  \left (g\right )\right ) \times  \\
 & \left (I_{\mathcal{F}} +\sum _{1 \leq j \leq s}c_{j} \left (g\right ) a_{j}^{ \dag } a_{j +s}^{ \dag } +\sum _{1 \leq j_{1} <j_{2} \leq s}c_{j_{1}} \left (g\right ) c_{j_{2}} \left (g\right ) a_{j_{1}}^{ \dag } a_{j_{1} +s}^{ \dag } a_{j_{2}}^{ \dag } a_{j_{2} +s}^{ \dag } +\cdots  +c_{j_{1}} \left (g\right ) \cdots  c_{j_{s}} \left (g\right ) a_{1}^{ \dag } a_{1 +s}^{ \dag } \cdots  a_{s}^{ \dag } a_{s +s}^{ \dag }\right )\vert \phi  \rangle \end{align}
\begin{equation}\mathcal{N} \left (\phi  \left (g\right )\right ) =\left (1 +S_{1} \left (g\right ) +S_{2} \left (g\right ) +\cdots  +S_{s} \left (g\right )\right )^{ -\frac{1}{2}} =\left (\det \left \{1 +G^{ \dag } G\right \}\right )^{ -\frac{1}{2}} =\left (\det \left \{1 +\mathbf{c}^{ \dag } \mathbf{c}\right \}\right )^{ -\frac{1}{2}}
\end{equation}For
\begin{equation}r =4 ,s =2
\end{equation}compare
\begin{align} & \mathcal{N} \left (\phi  \left (g\right )\right ) \left (I_{\mathcal{F}} +c_{1} \left (g\right ) a_{1}^{ \dag } a_{3}^{ \dag }\right ) \left (I_{\mathcal{F}} +c_{2} \left (g\right ) a_{2}^{ \dag } a_{4}^{ \dag }\right )\vert \phi  \rangle  \\
 &  =\left (\mathcal{N} \left (\phi  \left (g\right )\right ) I_{\mathcal{F}} +2 \mathcal{N} \left (\phi  \left (g\right )\right ) c_{1} \left (g\right ) a_{1}^{ \dag } a_{3}^{ \dag } +2 \mathcal{N} \left (\phi  \left (g\right )\right ) c_{2} \left (g\right ) a_{2}^{ \dag } a_{4}^{ \dag } +\mathcal{N} \left (\phi  \left (g\right )\right ) c_{1} \left (g\right ) c_{2} \left (g\right ) a_{1}^{ \dag } a_{3}^{ \dag } a_{2}^{ \dag } a_{4}^{ \dag }\right )\vert \phi  \rangle \end{align}to
\begin{align} & \left (\mathfrak{d}_{1} \left (g\right ) I_{\mathcal{F}} +\mathfrak{c}_{1} \left (g\right ) a_{1}^{ \dag } a_{3}^{ \dag }\right ) \left (\mathfrak{d}_{2} \left (g\right ) I_{\mathcal{F}} +\mathfrak{c}_{2} \left (g\right ) a_{2}^{ \dag } a_{4}^{ \dag }\right )\vert \phi  \rangle  \\
 &  =\left (\mathfrak{d}_{1} \left (g\right ) \mathfrak{d}_{2} \left (g\right ) I_{\mathcal{F}} +\mathfrak{d}_{1} \left (g\right ) \mathfrak{c}_{2} \left (g\right ) a_{2}^{ \dag } a_{4}^{ \dag } +\mathfrak{c}_{1} \left (g\right ) \mathfrak{d}_{2} \left (g\right ) a_{1}^{ \dag } a_{3}^{ \dag } +\mathfrak{c}_{1} \left (g\right ) \mathfrak{c}_{2} \left (g\right ) a_{1}^{ \dag } a_{3}^{ \dag } a_{2}^{ \dag } a_{4}^{ \dag }\right )\vert \phi  \rangle \end{align}to obtain
\begin{align}\mathfrak{d}_{1} \left (g\right ) \mathfrak{d}_{2} \left (g\right ) &  =\left (\det \left \{1 +\mathbf{c}^{ \dag } \mathbf{c}\right \}\right )^{ -\frac{1}{2}} \\
\mathfrak{c}_{1} \left (g\right ) \mathfrak{d}_{2} \left (g\right ) &  =2 \left (\det \left \{1 +\mathbf{c}^{ \dag } \mathbf{c}\right \}\right )^{ -\frac{1}{2}} c_{1} \left (g\right ) \\
\mathfrak{c}_{2} \left (g\right ) \mathfrak{d}_{1} \left (g\right ) &  =2 \left (\det \left \{1 +\mathbf{c}^{ \dag } \mathbf{c}\right \}\right )^{ -\frac{1}{2}} c_{2} \left (g\right ) \\
\mathfrak{c}_{1} \left (g\right ) \mathfrak{c}_{2} \left (g\right ) &  =\left (\det \left \{1 +\mathbf{c}^{ \dag } \mathbf{c}\right \}\right )^{ -\frac{1}{2}} c_{1} \left (g\right ) c_{2} \left (g\right )\text{.}\end{align}For
\begin{equation}r =2 ,s =1
\end{equation}compare
\begin{equation*}\mathcal{N} \left (\phi  \left (g\right )\right ) \left (I_{\mathcal{F}} +c_{1} \left (g\right ) a_{1}^{ \dag } a_{2}^{ \dag }\right )\vert \phi  \rangle  =\mathcal{N} \left (\phi  \left (g\right )\right ) \left \{\vert \phi  \rangle  +c_{1} \left (g\right )\vert \varphi _{1} \varphi _{2} \rangle \right \} =\left (1 +c_{1} \left (g\right )^{2}\right )^{ -\frac{1}{2}} \vert \phi  \rangle  +\left (1 +c_{1} \left (g\right )^{2}\right )^{ -\frac{1}{2}} c_{1} \left (g\right )\vert \varphi _{1} \varphi _{2} \rangle 
\end{equation*}to
\begin{equation}\left (\mathfrak{d}_{1} \left (g\right ) I_{\mathcal{F}} +\mathfrak{c}_{1} \left (g\right ) a_{1}^{ \dag } a_{2}^{ \dag }\right )\vert \phi  \rangle  =\mathfrak{d}_{1} \left (g\right ) \vert \phi  \rangle  +\mathfrak{c}_{1} \left (g\right ) a_{1}^{ \dag } a_{2}^{ \dag }\vert \phi  \rangle 
\end{equation}to obtain
\begin{align}\mathfrak{d}_{1} \left (g\right ) &  =\left (\det \left \{1 +c_{1}^{ \dag } c_{1}\right \}\right )^{ -\frac{1}{2}} \\
\mathfrak{c}_{1} \left (g\right ) &  =\left (\det \left \{1 +c_{1}^{ \dag } c_{1}\right \}\right )^{ -\frac{1}{2}} c_{1} \left (g\right )\text{.}\end{align}$\rfloor $ 

These states have the properties: 

1.
\begin{align} & a_{k}^{ \dag } {\displaystyle\prod _{1 \leq j \leq s}}\left (\mathfrak{d}_{j} \left (g\right ) I_{\mathcal{F}} +\mathfrak{c}_{j} \left (g\right ) a_{j}^{ \dag } a_{j +s}^{ \dag }\right )\vert \phi  \rangle  ={\displaystyle\prod _{1 \leq j \neq k \leq s}}\left (\mathfrak{d}_{j} \left (g\right ) I_{\mathcal{F}} +\mathfrak{c}_{j} \left (g\right ) a_{j}^{ \dag } a_{j +s}^{ \dag }\right ) a_{k}^{ \dag } \left (d_{k} I_{\mathcal{F}} +\mathfrak{c}_{k} \left (g\right ) a_{k}^{ \dag } a_{k +s}^{ \dag }\right )\vert \phi  \rangle  \nonumber  \\
 &  ={\displaystyle\prod _{1 \leq j \neq k \leq s}}\left (\mathfrak{d}_{j} \left (g\right ) I_{\mathcal{F}} +\mathfrak{c}_{j} \left (g\right ) a_{j}^{ \dag } a_{j +s}^{ \dag }\right ) a_{k}^{ \dag }\vert \phi  \rangle  =\mathfrak{d}_{k} \left (g\right ) a_{k}^{ \dag } {\displaystyle\prod _{1 \leq j \neq k \leq s}}\left (\mathfrak{d}_{j} \left (g\right ) I_{\mathcal{F}} +\mathfrak{c}_{j} \left (g\right ) a_{j}^{ \dag } a_{j +s}^{ \dag }\right )\vert \phi  \rangle \text{,}\end{align}i.e.
\begin{equation*}a_{k}^{ \dag }\vert \phi  \left (g\right ) \rangle  =\mathfrak{d}_{k} \left (g\right ) a_{k}^{ \dag }\vert \phi _{g} \left (\hat{k}\right ) \rangle \text{.}
\end{equation*}2.
\begin{align} & a_{k} {\displaystyle\prod _{1 \leq j \leq s}}\left (\mathfrak{d}_{j} \left (g\right ) I_{\mathcal{F}} +\mathfrak{c}_{j} \left (g\right ) a_{j}^{ \dag } a_{j +s}^{ \dag }\right )\vert \phi  \rangle  ={\displaystyle\prod _{1 \leq j \neq k \leq s}}\left (\mathfrak{d}_{j} \left (g\right ) I_{\mathcal{F}} +\mathfrak{c}_{j} \left (g\right ) a_{j}^{ \dag } a_{j +s}^{ \dag }\right ) a_{k} \left (\mathfrak{d}_{k} \left (g\right ) I_{\mathcal{F}} +\mathfrak{c}_{k} \left (g\right ) a_{k}^{ \dag } a_{k +s}^{ \dag }\right )\vert \phi  \rangle  \nonumber  \\
 &  ={\displaystyle\prod _{1 \leq j \neq k \leq s}}\left (\mathfrak{d}_{j} \left (g\right ) I_{\mathcal{F}} +\mathfrak{c}_{j} \left (g\right ) a_{j}^{ \dag } a_{j +s}^{ \dag }\right ) \left (\mathfrak{d}_{k} \left (g\right ) a_{k} +\mathfrak{c}_{k} \left (g\right ) a_{k} a_{k}^{ \dag } a_{k +s}^{ \dag }\right )\vert \phi  \rangle  \nonumber  \\
 &  ={\displaystyle\prod _{1 \leq j \neq k \leq s}}\left (\mathfrak{d}_{j} \left (g\right ) I_{\mathcal{F}} +\mathfrak{c}_{j} \left (g\right ) a_{j}^{ \dag } a_{j +s}^{ \dag }\right ) \left (\mathfrak{d}_{k} \left (g\right ) a_{k} +\mathfrak{c}_{k} \left (g\right ) \left (1 -a_{k}^{ \dag } a_{k}\right ) a_{k +s}^{ \dag }\right )\vert \phi  \rangle  \nonumber  \\
 &  ={\displaystyle\prod _{1 \leq j \neq k \leq s}}\left (\mathfrak{d}_{j} \left (g\right ) I_{\mathcal{F}} +\mathfrak{c}_{j} \left (g\right ) a_{j}^{ \dag } a_{j +s}^{ \dag }\right ) \left (\mathfrak{d}_{k} \left (g\right ) a_{k +s} -\mathfrak{c}_{k} \left (g\right ) a_{k +s}^{ \dag }\right )\vert \phi  \rangle  \\
 &  =\mathfrak{c}_{k} \left (g\right ) a_{k +s}^{ \dag } {\displaystyle\prod _{1 \leq j \neq k \leq s}}\left (\mathfrak{d}_{j} \left (g\right ) I_{\mathcal{F}} +\mathfrak{c}_{j} \left (g\right ) a_{j}^{ \dag } a_{j +s}^{ \dag }\right )\vert \phi  \rangle \end{align}i.e.
\begin{equation*}a_{k}\vert \phi _{g} \rangle  =\mathfrak{c}_{k} \left (g\right ) a_{k +s}^{ \dag }\vert \phi _{g} \left (\hat{k}\right ) \rangle \text{,}
\end{equation*}3.
\begin{align} & a_{k +s}^{ \dag } {\displaystyle\prod _{1 \leq j \leq s}}\left (\mathfrak{d}_{j} I_{\mathcal{F}} +\mathfrak{c}_{j} \left (g\right ) a_{j}^{ \dag } a_{j +s}^{ \dag }\right )\vert \phi  \rangle  ={\displaystyle\prod _{1 \leq j \neq k \leq s}}\left (\mathfrak{d}_{j} I_{\mathcal{F}} +\mathfrak{c}_{j} \left (g\right ) a_{j}^{ \dag } a_{j +s}^{ \dag }\right ) a_{k +s}^{ \dag } \left (\mathfrak{d}_{j} I_{\mathcal{F}} +\mathfrak{c}_{k} \left (g\right ) a_{k}^{ \dag } a_{k +s}^{ \dag }\right )\vert \phi  \rangle  \nonumber  \\
 &  ={\displaystyle\prod _{1 \leq j \neq k \leq s}}\left (\mathfrak{d}_{j} I_{\mathcal{F}} +\mathfrak{c}_{j} \left (g\right ) a_{j}^{ \dag } a_{j +s}^{ \dag }\right ) a_{k +s}^{ \dag }\vert \phi  \rangle  =d_{k} a_{k +s}^{ \dag } {\displaystyle\prod _{1 \leq j \neq k \leq s}}\left (\mathfrak{d}_{j} I_{\mathcal{F}} +\mathfrak{c}_{j} \left (g\right ) a_{j}^{ \dag } a_{j +s}^{ \dag }\right )\vert \phi  \rangle \text{,}\end{align}i.e.
\begin{equation*}a_{k +s}^{ \dag }\vert \phi _{g} \rangle  =\mathfrak{d}_{k} a_{k +s}^{ \dag }\vert \phi _{g} \left (\hat{k}\right ) \rangle 
\end{equation*}4.
\begin{align} & a_{k +s} {\displaystyle\prod _{1 \leq j \leq s}}\left (\mathfrak{d}_{j} I_{\mathcal{F}} +\mathfrak{c}_{j} a_{j}^{ \dag } a_{j +s}^{ \dag }\right )\vert \phi  \rangle  ={\displaystyle\prod _{1 \leq j \neq k \leq s}}\left (\mathfrak{d}_{j} I_{\mathcal{F}} +\mathfrak{c}_{j} a_{j}^{ \dag } a_{j +s}^{ \dag }\right ) a_{k +s} \left (\mathfrak{d}_{k} I_{\mathcal{F}} +\mathfrak{c}_{k} a_{k}^{ \dag } a_{k +s}^{ \dag }\right )\vert \phi  \rangle  \nonumber  \\
 &  ={\displaystyle\prod _{1 \leq j \neq k \leq s}}\left (\mathfrak{d}_{j} I_{\mathcal{F}} +\mathfrak{c}_{j} a_{j}^{ \dag } a_{j +s}^{ \dag }\right ) \left (\mathfrak{d}_{k} a_{k +s} +\mathfrak{c}_{k} a_{k +s} a_{k}^{ \dag } a_{k +s}^{ \dag }\right )\vert \phi  \rangle  ={\displaystyle\prod _{1 \leq j \neq k \leq s}}\left (\mathfrak{d}_{j} I_{\mathcal{F}} +\mathfrak{c}_{j} a_{j}^{ \dag } a_{j +s}^{ \dag }\right ) \left (\mathfrak{d}_{k} a_{k +s} -\mathfrak{c}_{k} a_{k}^{ \dag } \left (1 -a_{k +s}^{ \dag } a_{k +s}\right )\right )\vert \phi  \rangle  \nonumber  \\
 &  ={\displaystyle\prod _{1 \leq j \neq k \leq s}}\left (\mathfrak{d}_{j} I_{\mathcal{F}} +\mathfrak{c}_{j} a_{j}^{ \dag } a_{j +s}^{ \dag }\right ) \left (\mathfrak{d}_{k} a_{k +s} -\mathfrak{c}_{k} a_{k +s}^{ \dag }\right )\vert \phi  \rangle  = -\mathfrak{c}_{k} a_{k}^{ \dag } {\displaystyle\prod _{1 \leq j \neq k \leq s}}\left (\mathfrak{d}_{j} I_{\mathcal{F}} +\mathfrak{c}_{j} a_{j}^{ \dag } a_{j +s}^{ \dag }\right )\vert \phi  \rangle \end{align}i.e.
\begin{equation*}a_{k +s}\vert \phi  \left (g\right ) \rangle  = -\mathfrak{c}_{k} \left (g\right ) a_{k}^{ \dag }\vert \phi _{g} \left (\hat{k}\right ) \rangle \text{.}
\end{equation*}Summarizing
\begin{align}a_{k}^{ \dag }\vert \phi  \left (g\right ) \rangle  &  =\mathfrak{d}_{k} a_{k}^{ \dag }\vert \phi _{g} \left (\hat{k}\right ) \rangle  \nonumber  \\
a_{k}\vert \phi  \left (g\right ) \rangle  &  =\mathfrak{c}_{k} a_{k +s}^{ \dag }\vert \phi _{g} \left (\hat{k}\right ) \rangle  \nonumber  \\
a_{k +s}^{ \dag }\vert \phi  \left (g\right ) \rangle  &  =\mathfrak{d}_{k} a_{k +s}^{ \dag }\vert \phi _{g} \left (\hat{k}\right ) \rangle  \nonumber  \\
a_{k +s}\vert \phi  \left (g\right ) \rangle  &  = -\mathfrak{c}_{k} a_{k}^{ \dag }\vert \phi _{g} \left (\hat{k}\right ) \rangle \text{,}\end{align}which leads to dressed creation and annihilation operators
\begin{align}b_{k} &  =\mathfrak{c}_{k} \left (g\right ) a_{k}^{ \dag } +\mathfrak{d}_{k} \left (g\right ) a_{k +s} \nonumber  \\
b_{k +s} &  =\mathfrak{c}_{k} \left (g\right ) a_{k +s}^{ \dag } -\mathfrak{d}_{k} \left (g\right ) a_{k} \nonumber  \\
b_{k}^{ \dag } &  =\mathfrak{d}_{k} \left (g\right ) a_{k +s}^{ \dag } +\mathfrak{c}_{k} \left (g\right ) a_{k} \nonumber  \\
b_{k +s}^{ \dag } &  = -\mathfrak{d}_{k} \left (g\right ) a_{k}^{ \dag } +\mathfrak{c}_{k} \left (g\right ) a_{k +s}\end{align}$\lceil $If
\begin{align}b_{k} &  =a_{k}^{ \dag }\text{, i.e}\mathfrak{c}_{k} \left (g\right ) =1 ,\mathfrak{d}_{k} \left (g\right ) =0 \\
b_{k +s} &  =a_{k +s}^{ \dag }\text{, i.e}\mathfrak{c}_{k} \left (g\right ) =1 ,\mathfrak{d}_{k} \left (g\right ) =0 \\
b_{k}^{ \dag } &  =a_{k} \\
b_{k +s}^{ \dag } &  =a_{k +s}\end{align}one has vacuums formed by GHFB states
\begin{equation}\vert \phi  \left (g\right ) \rangle  =\mathcal{N} \left (\phi  \left (g\right )\right ) \exp  \left \{\sum _{1 \leq j \leq s}c_{j} \left (g\right ) a_{j}^{ \dag } a_{j +s}^{ \dag }\right \}\vert \phi  \rangle  =\mathcal{N} \left (\phi  \left (g\right )\right ) {\displaystyle\prod _{1 \leq j \leq s}}\exp  \left \{c_{j} \left (g\right ) a_{j}^{ \dag } a_{j +s}^{ \dag }\right \}\vert \phi  \rangle  ={\displaystyle\prod _{1 \leq j \leq s}}\left (\mathfrak{d}_{j} \left (g\right ) I_{\mathcal{F}} +\mathfrak{c}_{j} \left (g\right ) a_{j}^{ \dag } a_{j +s}^{ \dag }\right )\vert \phi  \rangle \text{.}
\end{equation}Consider the particle - hole transformation for $r =2$
\begin{align} & \left (\begin{array}{cccc}a_{k}^{ \dag } & a_{k +s}^{ \dag } & a_{k} & a_{k +s}\end{array}\right ) \mathbb{T}_{k} =\left (\begin{array}{cccc}a_{k}^{ \dag } & a_{k +s}^{ \dag } & a_{k} & a_{k +s}\end{array}\right ) \left (\begin{array}{cc}\mathbb{U}_{k} & \overline{\mathbb{V}}_{k} \\
\mathbb{V}_{k} & \overline{\mathbb{U}}_{k}\end{array}\right ) \nonumber  \\
 &  =\left (\begin{array}{cccc}a_{k}^{ \dag } & a_{k +s}^{ \dag } & a_{k} & a_{k +s}\end{array}\right ) \left (\begin{array}{cccc}0 & 0 & 1 & 0 \\
0 & 0 & 0 & 1 \\
1 & 0 & 0 & 0 \\
0 & 1 & 0 & 0\end{array}\right ) =\left (\begin{array}{cccc}b_{k}^{ \dag } & b_{k +s}^{ \dag } & b_{k} & b_{k +s}\end{array}\right ) =\left (\begin{array}{cccc}a_{k} & a_{k +s} & a_{k}^{ \dag } & a_{k +s}^{ \dag }\end{array}\right )\text{.}\end{align}This is a \emph{proper Bogoliubov} transformation i.e. $ \in O \left (4 ,\mathbb{R}\right )$\emph{\ and} $ \in S O \left (4 ,\mathbb{R}\right )$ as
\begin{equation}\det \left \{\mathbb{T}_{k}\right \} =\left \vert \begin{array}{cccc}0 & 0 & 1 & 0 \\
0 & 0 & 0 & 1 \\
1 & 0 & 0 & 0 \\
0 & 1 & 0 & 0\end{array}\right \vert  =1 \left \vert \begin{array}{ccc}0 & 0 & 1 \\
1 & 0 & 0 \\
0 & 1 & 0\end{array}\right \vert  =1 \left \vert \begin{array}{cc}1 & 0 \\
0 & 1\end{array}\right \vert  =1.
\end{equation}The real form of this transformation is
\begin{equation}\mathbb{K}_{k} \mathbb{T}_{k} \mathbb{K}_{k}^{ \dag } =\left (\begin{array}{cc}\ensuremath{\operatorname*{Re}} \left (\mathbb{U}_{k} +\mathbb{V}_{k}\right ) & \ensuremath{\operatorname*{Im}} \left (\mathbb{V}_{k} +\mathbb{U}_{k}\right ) \\
\ensuremath{\operatorname*{Im}} \left (\mathbb{V}_{k} -\mathbb{U}_{k}\right ) & \ensuremath{\operatorname*{Re}} \left (\mathbb{U}_{k} -\mathbb{V}_{k}\right )\end{array}\right ) =\left (\begin{array}{cccc}1 & 0 & 0 & 0 \\
0 & 1 & 0 & 0 \\
0 & 0 &  -1 & 0 \\
0 & 0 & 0 &  -1\end{array}\right )\text{,}
\end{equation}which gives
\begin{align} & \left (\begin{array}{cccc}a_{k}^{ \dag } & a_{k +s}^{ \dag } & a_{k} & a_{k +s}\end{array}\right ) \mathbb{T}_{k} =\left (\begin{array}{cccc}a_{k}^{ \dag } & a_{k +s}^{ \dag } & a_{k} & a_{k +s}\end{array}\right ) \left (\begin{array}{cc}\mathbb{U}_{k} & \overline{\mathbb{V}}_{k} \\
\mathbb{V}_{k} & \overline{\mathbb{U}}_{k}\end{array}\right ) \nonumber  \\
 &  =\left (\begin{array}{cccc}a_{k}^{ \dag } & a_{k +s}^{ \dag } & a_{k} & a_{k +s}\end{array}\right ) \left (\begin{array}{cccc}1 & 0 & 0 & 0 \\
0 & 0 & 0 & 1 \\
1 & 0 & 0 & 0 \\
0 & 0 & 0 & 1\end{array}\right ) =\left (\begin{array}{cccc}b_{k}^{ \dag } & b_{k +s}^{ \dag } & b_{k} & b_{k +s}\end{array}\right ) =\left (\begin{array}{cccc}a_{k} & a_{k +s}^{ \dag } & a_{k}^{ \dag } & a_{k +s}\end{array}\right )\text{.}\end{align}This is again a proper Bogoliubov transformation i.e. $ \in S O \left (4 ,\mathbb{R}\right )$ as
\begin{equation}\det \left \{\mathbb{T}_{k}\right \} =\left \vert \begin{array}{cccc}0 & 0 & 1 & 0 \\
0 & 0 & 0 & 1 \\
1 & 0 & 0 & 0 \\
0 & 1 & 0 & 0\end{array}\right \vert  =1 \left \vert \begin{array}{ccc}0 & 0 & 1 \\
1 & 0 & 0 \\
0 & 1 & 0\end{array}\right \vert  =1 \left \vert \begin{array}{cc}1 & 0 \\
0 & 1\end{array}\right \vert  =1.
\end{equation}

Let $r =1$ then
\begin{equation}\mathbb{T}_{k} =\left (\begin{array}{cc}\mathbb{U}_{k} & \overline{\mathbb{V}}_{k} \\
\mathbb{V}_{k} & \overline{\mathbb{U}}_{k}\end{array}\right ) =\left (\begin{array}{cc}0 & 1 \\
1 & 0\end{array}\right )
\end{equation}This is an \emph{improper} Bogoliubov transformation i.e. $ \in O \left (2 r ,\mathbb{R}\right )$ and $ \notin S O \left (2 r ,\mathbb{R}\right )$ as
\begin{equation}\det \left \{\mathbb{T}_{k}\right \} =\left \vert \begin{array}{cc}0 & 1 \\
1 & 0\end{array}\right \vert  = -1.
\end{equation}The real form of this transformation is
\begin{equation}\mathbb{K}_{k} \mathbb{T}_{k} \mathbb{K}_{k}^{ \dag } =\left (\begin{array}{cc}\ensuremath{\operatorname*{Re}} \left (\mathbb{U}_{k} +\mathbb{V}_{k}\right ) & \ensuremath{\operatorname*{Im}} \left (\mathbb{V}_{k} +\mathbb{U}_{k}\right ) \\
\ensuremath{\operatorname*{Im}} \left (\mathbb{V}_{k} -\mathbb{U}_{k}\right ) & \ensuremath{\operatorname*{Re}} \left (\mathbb{U}_{k} -\mathbb{V}_{k}\right )\end{array}\right ) =\left (\begin{array}{cc}1 & 0 \\
0 &  -1\end{array}\right )\text{.}
\end{equation}$\rfloor $
\begin{align} & \left (\begin{array}{cccc}a_{k}^{ \dag } & a_{k +s}^{ \dag } & a_{k} & a_{k +s}\end{array}\right ) \mathbb{T}_{k} =\left (\begin{array}{cccc}a_{k}^{ \dag } & a_{k +s}^{ \dag } & a_{k} & a_{k +s}\end{array}\right ) \left (\begin{array}{cc}\mathbb{U}_{k} & \overline{\mathbb{V}}_{k} \\
\mathbb{V}_{k} & \overline{\mathbb{U}}_{k}\end{array}\right ) \nonumber  \\
 &  =\left (\begin{array}{cccc}a_{k}^{ \dag } & a_{k +s}^{ \dag } & a_{k} & a_{k +s}\end{array}\right ) \left (\begin{array}{cccc}0 &  -\mathfrak{d}_{k} & \mathfrak{c}_{k} & 0 \\
\mathfrak{d}_{k} & 0 & 0 & \mathfrak{c}_{k} \\
\mathfrak{c}_{k} & 0 & 0 &  -\mathfrak{d}_{k} \\
0 & \mathfrak{c}_{k} & \mathfrak{d}_{k} & 0\end{array}\right ) =\left (\begin{array}{cccc}b_{k}^{ \dag } & b_{k +s}^{ \dag } & b_{k} & b_{k +s}\end{array}\right )\end{align}i.e.
\begin{equation}\left (\begin{array}{cc}\mathbb{U}_{k} & \overline{\mathbb{V}}_{k} \\
\mathbb{V}_{k} & \overline{\mathbb{U}}_{k}\end{array}\right ) =\left (\begin{array}{cccc}0 &  -\mathfrak{d}_{k} & \mathfrak{c}_{k} & 0 \\
\mathfrak{d}_{k} & 0 & 0 & \mathfrak{c}_{k} \\
\mathfrak{c}_{k} & 0 & 0 &  -\mathfrak{d}_{k} \\
0 & \mathfrak{c}_{k} & \mathfrak{d}_{k} & 0\end{array}\right )\text{.}
\end{equation}This gives
\begin{equation}T \left (\begin{array}{cc}\mathbf{a}^{ \dag } & \mathbf{a}\end{array}\right ) =\left (\begin{array}{cc}\mathbf{a}^{ \dag } & \mathbf{a}\end{array}\right ) \mathbb{T} =\left (\begin{array}{cc}\mathbf{a}^{ \dag } & \mathbf{a}\end{array}\right ) \left (\begin{array}{cc}\mathbb{U} & \overline{\mathbb{V}} \\
\mathbb{V} & \overline{\mathbb{U}}\end{array}\right ) =\left (\begin{array}{cc}\mathbf{b}^{ \dag } & \mathbf{b}\end{array}\right )\text{,}
\end{equation}where
\begin{equation*}\mathbb{T} =\left (\begin{array}{ccc}\mathbb{T}_{1} & \mathbb{O} & \mathbb{O} \\
\mathbb{O} & \ddots  & \mathbb{O} \\
\mathbb{O} & \mathbb{O} & \mathbb{T}_{r}\end{array}\right )\text{.}
\end{equation*}The conditions
\begin{equation}\mathbb{U}^{ \dag } \mathbb{U} +\mathbb{V}^{ \dag } \mathbb{V}\, =\mathbb{I}_{r}
\end{equation}
\begin{equation}\mathbb{U}^{ \dag } \overline{\mathbb{V}} +\mathbb{V}^{ \dag } \overline{\mathbb{U}} =\mathbb{O}_{r}
\end{equation}give
\begin{equation*}\left (\begin{array}{cc}0 & \mathfrak{d}_{k} \\
 -\mathfrak{d}_{k} & 0\end{array}\right ) \left (\begin{array}{cc}0 &  -\mathfrak{d}_{k} \\
\mathfrak{d}_{k} & 0\end{array}\right ) +\left (\begin{array}{cc}\mathfrak{c}_{k} & 0 \\
0 & \mathfrak{c}_{k}\end{array}\right ) \left (\begin{array}{cc}\mathfrak{c}_{k} & 0 \\
0 & \mathfrak{c}_{k}\end{array}\right ) =\left (\begin{array}{cc}1 & 0 \\
0 & 1\end{array}\right )
\end{equation*}i.e.
\begin{equation*}\mathfrak{d}_{k}^{2} +\mathfrak{c}_{k}^{2} =1.
\end{equation*}By construction $\vert \phi  \left (g\right ) \rangle $ is a vacuum for $\left \{b_{k}\right \}$ i.e.
\begin{equation*}b_{k}\vert \phi  \left (g\right ) \rangle  =\left (\mathfrak{c}_{k} a_{k}^{ \dag } +\mathfrak{d}_{k} a_{k +s}\right )\vert \phi  \left (g\right ) \rangle  =0\text{,}
\end{equation*}and as the transformation $\mathbb{T}$ is canonical and $\vert \phi  \left (g\right ) \rangle $ is normalized then
\begin{align*} \langle \phi  \left (g\right )\left \vert \left \{b_{k}^{ \dag } ,b_{k}\right \}\phi  \left (g\right )\right . \rangle  &  = \langle \phi  \left (g\right )\left \vert \phi  \left (g\right )\right . \rangle  = \langle \phi  \left (g\right )\left \vert \left (b_{k}^{ \dag } b_{k} +b_{k} b_{k}^{ \dag }\right ) \phi  \left (g\right )\right . \rangle  \\
 &  = \langle \phi  \left (g\right )\left \vert b_{k} b_{k}^{ \dag } \phi  \left (g\right )\right . \rangle  =1\text{,}\end{align*}which implies
\begin{equation*}\vert b_{k} \rangle  =b_{k}^{ \dag }\vert \phi  \left (g\right ) \rangle 
\end{equation*}is a normalized state. 

One can show that
\begin{align}\left (\overline{\mathbb{V}}_{k} \overline{\mathbb{U}}_{k}^{ -1}\right ) &  =\mathbb{G}_{k} =\left (\begin{array}{cc}\mathfrak{c}_{k} & 0 \\
0 & \mathfrak{c}_{k}\end{array}\right ) \left (\begin{array}{cc}0 &  -\mathfrak{d}_{k} \\
\mathfrak{d}_{k} & 0\end{array}\right )^{ -1} \\
 &  =\left (\begin{array}{cc}\mathfrak{c}_{k} & 0 \\
0 & \mathfrak{c}_{k}\end{array}\right ) \left (\begin{array}{cc}0 & \mathfrak{d}_{k}^{ -1} \\
 -\mathfrak{d}_{k}^{ -1} & 0\end{array}\right ) =\left (\begin{array}{cc}0 & \mathfrak{c}_{k} \mathfrak{d}_{k}^{ -1} \\
 -\mathfrak{c}_{k} \mathfrak{d}_{k}^{ -1} & 0\end{array}\right ) \\
 &  =\left (\begin{array}{cc}0 & \frac{\mathfrak{c}_{k}}{\mathfrak{d}_{k}} \\
 -\frac{\mathfrak{c}_{k}}{\mathfrak{d}_{k}} & 0\end{array}\right ) =\left (\begin{array}{cc}0 & c_{k} \\
 -c_{k +s} & 0\end{array}\right )\text{.}\end{align}i.e.
\begin{equation}c_{k} =\frac{\mathfrak{c}_{k}}{\mathfrak{d}_{k}}
\end{equation}$\lceil $
\begin{equation}\left (\begin{array}{cc}0 & \mathfrak{d}_{k}^{ -1} \\
 -\mathfrak{d}_{k}^{ -1} & 0\end{array}\right ) \left (\begin{array}{cc}0 &  -\mathfrak{d}_{k} \\
\mathfrak{d}_{k} & 0\end{array}\right ) =\left (\begin{array}{cc}1 & 0 \\
0 & 1\end{array}\right )
\end{equation}$\rfloor $
\begin{equation}\left (\begin{array}{cc}\mathbb{U}_{k} & \overline{\mathbb{V}}_{k} \\
\mathbb{V}_{k} & \overline{\mathbb{U}}_{k}\end{array}\right ) =\left (\begin{array}{cccc}0 &  -\mathfrak{d}_{k} & \mathfrak{c}_{k} & 0 \\
\mathfrak{d}_{k} & 0 & 0 & \mathfrak{c}_{k} \\
\mathfrak{c}_{k} & 0 & 0 &  -\mathfrak{d}_{k} \\
0 & \mathfrak{c}_{k} & \mathfrak{d}_{k} & 0\end{array}\right )
\end{equation}
\begin{align}\vert \phi  \left (g\right ) \rangle  &  =\mathcal{N} \left (\phi  \left (g\right )\right ) \exp  \left \{\sum _{1 \leq j \leq s}\frac{\mathfrak{c}_{j}}{\mathfrak{d}_{j}} a_{j}^{ \dag } a_{j +s}^{ \dag }\right \}\vert \phi  \rangle  =\mathcal{N} \left (\phi  \left (g\right )\right ) {\displaystyle\prod _{1 \leq j \leq s}}\exp  \left \{\frac{\mathfrak{c}_{j}}{\mathfrak{d}_{j}} a_{j}^{ \dag } a_{j +s}^{ \dag }\right \}\vert \phi  \rangle  \\
 &  ={\displaystyle\prod _{1 \leq j \leq s}}\left (\mathfrak{d}_{j} \left (g\right ) I_{\mathcal{F}} +\mathfrak{c}_{j} \left (g\right ) a_{j}^{ \dag } a_{j +s}^{ \dag }\right )\vert \phi  \rangle \text{.}\end{align}

$\lceil $If
\begin{equation*}\mathfrak{c}_{k} =\mathfrak{d}_{k}
\end{equation*}then
\begin{align}b_{k} &  =\mathfrak{c}_{k} \left (a_{k}^{ \dag } +a_{k +s}\right ) \nonumber  \\
b_{k +s} &  =\mathfrak{c}_{k} \left (a_{k +s}^{ \dag } -a_{k}\right ) \nonumber  \\
b_{k}^{ \dag } &  =\mathfrak{c}_{k} \left (a_{k +s}^{ \dag } +a_{k}\right ) \nonumber  \\
b_{k +s}^{ \dag } &  =\mathfrak{c}_{k} \left ( -a_{k}^{ \dag } +a_{k +s}\right )\end{align}and
\begin{equation*}b_{k} +b_{k +s} =\mathfrak{c}_{k} \left (a_{k}^{ \dag } +a_{k +s} +a_{k +s}^{ \dag } -a_{k}\right )
\end{equation*}$\rfloor $ 

The states $\phi  \left (g\right )$ are $O (2 r ,\mathbb{R})/U (r)$ coherent states that are parameterized by $g$ and thus form a K{\"a}hler manifold. 

\section{General Self Adjoint Field Operators}
In general one can define self adjoint operators by the \emph{real} linear maps
\begin{equation}X ,Y :\mathfrak{h} \rightarrow \mathcal{F}\text{,}
\end{equation}from the complex Hilbert space $\mathfrak{h}$ given by
\begin{align}X \left (\vert u \rangle \right ) &  =\frac{1}{\sqrt{2}} \left (\Phi ^{ \dag } \left (\vert u \rangle \right ) +\Phi  \left (\vert u \rangle \right )\right ) =X \left (\vert u \rangle \right )^{ \dag } \\
Y \left (\vert u \rangle \right ) &  =\frac{i}{\sqrt{2}} \left (\Phi ^{ \dag } \left (\vert u \rangle \right ) -\Phi  \left (\vert u \rangle \right )\right ) =Y \left (\vert u \rangle \right )^{ \dag }\text{,}\end{align}which give
\begin{align}\Phi ^{ \dag } \left (\vert u \rangle \right ) &  =\frac{1}{\sqrt{2}} \left (X \left (\vert u \rangle \right ) -i Y \left (\vert u \rangle \right )\right ) \\
\Phi  \left (\vert u \rangle \right ) &  =\frac{1}{\sqrt{2}} \left (X \left (\vert u \rangle \right ) +i Y \left (\vert u \rangle \right )\right )\text{.}\end{align}One can observe that $X$ and $Y$ are related by
\begin{equation}X \left (i\vert u \rangle \right ) =\frac{1}{\sqrt{2}} \left (\Phi ^{ \dag } \left (i\vert u \rangle \right ) +\Phi  \left (\vert i u \rangle \right )\right ) =\frac{1}{\sqrt{2}} \left (i \Phi ^{ \dag } \left (\vert u \rangle \right ) -i \Phi  \left (\vert u \rangle \right )\right ) =\frac{i}{\sqrt{2}} \left (\Phi ^{ \dag } \left (\vert u \rangle \right ) -\Phi  \left (\vert u \rangle \right )\right ) =Y \left (\vert u \rangle \right )\text{.}
\end{equation}The action of the \emph{complex linear map} $\Phi ^{ \dag }$ on the complex Hilbert space $\mathfrak{h}$ in terms of the \emph{real map} $X$ is then given by
\begin{align}\Phi ^{ \dag } \left (\vert u \rangle \right ) &  =\frac{1}{\sqrt{2}} \left (X \left (\vert u \rangle \right ) -i X \left (i\vert u \rangle \right )\right ) \\
\Phi ^{ \dag } \left (i\vert v \rangle \right ) &  =\frac{1}{\sqrt{2}} \left (X \left (i\vert v \rangle \right ) +i X \left (\vert v \rangle \right )\right ) \\
\Phi ^{ \dag } \left (\vert u \rangle  +i\vert v \rangle \right ) &  =\frac{1}{\sqrt{2}} \left (X \left (\vert u \rangle \right ) -i X \left (i\vert u \rangle \right )\right ) +\frac{1}{\sqrt{2}} \left (X \left (i\vert v \rangle \right ) +i X \left (\vert v \rangle \right )\right ) \\
 &  =\frac{1}{\sqrt{2}} \left (X \left (\vert u \rangle  +X \left (i\vert v \rangle \right )\right )\right ) +\frac{1}{\sqrt{2}} \left (i X \left (\vert v \rangle  -i X \left (i\vert u \rangle \right )\right )\right ) \\
 &  =\frac{1}{\sqrt{2}} X \left (\vert u \rangle  +i\vert v \rangle \right ) +i X \left (\vert v \rangle  -i\vert u \rangle \right )\text{.}\end{align}If one makes the restriction $\vert u \rangle  ,\vert v \rangle  \in \mathfrak{h}_{\mathbb{R}}\text{,}$ where $\mathfrak{h}$ is a complexification of $\mathfrak{h}_{\mathbb{R}}$ then the $\mathbb{C} -$linear map $\Phi ^{ \dag }$
\begin{equation}\Phi ^{ \dag } :\mathfrak{h} \rightarrow \mathcal{F}
\end{equation}corresponds to the $\mathbb{R} -$linear map
\begin{equation}X :\mathfrak{h} \oplus i \overline{\mathfrak{h}} \rightarrow \mathcal{F}
\end{equation}from the real spaces $\mathfrak{h} ,\overline{\mathfrak{h}}\text{.}$ 

$\lceil $\emph{Note that} $X$ and $Y$ are $\mathbb{R}$-linear i.e.
\begin{align}X \left (\sum _{1 \leq j \leq p}\lambda _{j}\vert u_{j} \rangle \right ) &  =\sum _{1 \leq j \leq p}\lambda _{j} X \left (\vert u_{j} \rangle \right ) \\
Y \left (\sum _{1 \leq j \leq p}\lambda _{j}\vert u_{j} \rangle \right ) &  =\sum _{1 \leq j \leq p}\lambda _{j} Y \left (\vert u_{j} \rangle \right )\text{,}\end{align}when
\begin{equation}\lambda _{j} \in \mathbb{R} ,1 \leq j \leq p\text{,}
\end{equation}while
\begin{align}X \left (\sum _{1 \leq j \leq p}\lambda _{j}\vert u_{j} \rangle \right ) &  \neq \sum _{1 \leq j \leq p}\lambda _{j} X \left (\vert u_{j} \rangle \right ) \\
Y \left (\sum _{1 \leq j \leq p}\lambda _{j}\vert u_{j} \rangle \right ) &  \neq \sum _{1 \leq j \leq p}\lambda _{j} Y \left (\vert u_{j} \rangle \right )\text{,}\end{align}when
\begin{equation}\lambda _{j} \in \mathbb{C} , \notin  ,\mathbb{R} ;1 \leq j \leq p\text{\quad }\text{i.e} .\lambda _{j} \neq \overline{\lambda _{j}}\text{.}
\end{equation}$\rfloor $. 

\subsubsection{The real mapping $\mathfrak{X}$}
The maps $X$ and $Y$ can be consolidated into a single \emph{real} map, $\mathfrak{X}\text{,}$ that depends on the choice of basis (if one starts with a complex Hilbert space) or more basically on the \emph{complexification
of a real space} $\mathfrak{h}_{\mathbb{R}}$
\begin{equation}\mathfrak{X} :\mathfrak{h}_{\mathbb{R}} \oplus i \mathfrak{h}_{\mathbb{R}} \rightarrow \mathcal{F}\text{,}
\end{equation}where
\begin{align}\mathfrak{X} \left (\vert \varphi _{j} \rangle \right ) &  =X \left (\vert \varphi _{j} \rangle \right ) \\
\mathfrak{X} \left (i\vert \varphi _{j} \rangle \right ) &  =Y \left (\vert \varphi _{j} \rangle \right )\end{align}and
\begin{equation}\mathfrak{h}_{\mathbb{R}} =\mathbb{R}\text{-}\ensuremath{\operatorname*{linspan}}\left \{\left \{\left .\vert \varphi _{j} \rangle \right \vert 1 \leq j \leq r\right \}\right \}\text{.}
\end{equation}$\lceil $
\begin{align}X \left (\vert \varphi _{j} \rangle \right ) &  =\frac{1}{\sqrt{2}} \left (\Phi ^{ \dag } \left (\vert \varphi _{j} \rangle \right ) +\Phi  \left (\vert \varphi _{j} \rangle \right )\right ) \\
Y \left (\vert \varphi _{j} \rangle \right ) &  =\frac{i}{\sqrt{2}} \left (\Phi ^{ \dag } \left (\vert \varphi _{j} \rangle \right ) -\Phi  \left (\vert \varphi _{j} \rangle \right )\right ) \\
 &  =\frac{1}{\sqrt{2}} \left (\Phi ^{ \dag } \left (\vert i \varphi _{j} \rangle \right ) +\Phi  \left (i\vert \varphi _{j} \rangle \right )\right ) =X \left (i\vert \varphi _{j} \rangle \right )\text{,}\end{align}i.e.
\begin{equation}Y \left (\vert u \rangle \right ) =X \left (i\vert u \rangle \right )\text{.}
\end{equation}$\rfloor $ 

Non self adjoint elements can be expressed in terms of non self adjoint operators as
\begin{align}\Phi ^{ \dag } \left (\vert u \rangle \right ) &  =\frac{1}{\sqrt{2}} \left (\mathfrak{X} \left (\vert u \rangle \right ) -i \mathfrak{X} \left (i\vert u \rangle \right )\right ) \\
\Phi  \left (\vert u \rangle \right ) &  =\frac{1}{\sqrt{2}} \left (\mathfrak{X} \left (\vert u \rangle \right ) +i \mathfrak{X} \left (i\vert u \rangle \right )\right )\text{,}\end{align}where
\begin{equation}\vert u \rangle  \in \mathfrak{h}_{\mathbb{R}}\text{.}
\end{equation}

The maps $\left \{\Phi ^{ \dag } ,\Phi \right \}$ could also be named $\left \{a^{ \dag } ,a\right \}$ as
\begin{equation}a_{k}^{ \dag } =a^{ \dag } \left (\vert \varphi _{k} \rangle \right ) =\Phi ^{ \dag } \left (\vert \varphi _{k} \rangle \right ) =\Phi  \left (\vert \varphi _{k} \rangle \right )^{ \dag }\text{,}
\end{equation}but this has the dissadvantage of treating $\left \{a_{k}^{ \dag }\right \}$ differently to $\left \{b_{k}^{ \dag }\right \}\text{.}$ Hence
\begin{align} & \left \{X \left (\vert u \rangle \right ) ,X \left (\vert v \rangle \right )\right \} =\frac{1}{2} \left \{\left (\Phi ^{ \dag } \left (\vert u \rangle \right ) +\Phi  \left (\vert u \rangle \right )\right ) ,\left (\Phi ^{ \dag } \left (\vert v \rangle \right ) +\Phi  \left (\vert v \rangle \right )\right )\right \} \nonumber  \\
 &  =\left \{\Phi ^{ \dag } \left (\vert u \rangle \right ) ,\Phi ^{ \dag } \left (\vert v \rangle \right )\right \} +\left \{\Phi ^{ \dag } \left (\vert u \rangle \right ) ,\Phi  \left (\vert v \rangle \right )\right \} +\left \{\Phi  \left (\vert u \rangle \right ) ,\Phi ^{ \dag } \left (\vert v \rangle \right )\right \} +\left \{\Phi  \left (\vert u \rangle \right ) ,\Phi  \left (\vert v \rangle \right )\right \} \nonumber  \\
 &  =0 +\overline{ \langle \left .u\right \vert v \rangle } I_{\mathcal{F}} + \langle \left .u\right \vert v \rangle  I_{\mathcal{F}} +0 =\ensuremath{\operatorname*{Re}}  \langle \left .u\right \vert v \rangle  I_{\mathcal{F}}\text{,}\end{align}
\begin{align} & \left \{X \left (\vert u \rangle \right ) ,Y \left (\vert v \rangle \right )\right \} =\left \{\left (\Phi ^{ \dag } \left (\vert u \rangle \right ) +\Phi  \left (\vert u \rangle \right )\right ) ,i \left (\Phi ^{ \dag } \left (\vert v \rangle \right ) -\Phi  \left (\vert v \rangle \right )\right )\right \} \nonumber  \\
 &  =i \left \{\Phi ^{ \dag } \left (\vert u \rangle \right ) ,\Phi ^{ \dag } \left (\vert v \rangle \right )\right \} -i \left \{\Phi ^{ \dag } \left (\vert u \rangle \right ) ,\Phi  \left (\vert v \rangle \right )\right \} +i \left \{\Phi  \left (\vert u \rangle \right ) ,\Phi ^{ \dag } \left (\vert v \rangle \right )\right \} -i \left \{\Phi  \left (\vert u \rangle \right ) ,\Phi  \left (\vert v \rangle \right )\right \} \nonumber  \\
 &  =0 -i \overline{ \langle \left .u\right \vert v \rangle } I_{\mathcal{F}} +i  \langle \left .u\right \vert v \rangle  I_{\mathcal{F}} =\ensuremath{\operatorname*{Im}}  \langle \left .u\right \vert v \rangle  I_{\mathcal{F}}\end{align}and
\begin{align} & \left \{Y \left (\vert u \rangle \right ) ,Y \left (\vert v \rangle \right )\right \} =\left \{i \left (\Phi ^{ \dag } \left (\vert u \rangle \right ) -\Phi  \left (\vert u \rangle \right )\right ) ,i \left (\Phi ^{ \dag } \left (\vert v \rangle \right ) -\Phi  \left (\vert v \rangle \right )\right )\right \} \nonumber  \\
 &  = -\left \{\Phi ^{ \dag } \left (\vert u \rangle \right ) ,\Phi ^{ \dag } \left (\vert v \rangle \right )\right \} +\left \{\Phi ^{ \dag } \left (\vert u \rangle \right ) ,\Phi  \left (\vert v \rangle \right )\right \} +\left \{\Phi  \left (\vert u \rangle \right ) ,\Phi ^{ \dag } \left (\vert v \rangle \right )\right \} -\left \{\Phi  \left (\vert u \rangle \right ) ,\Phi  \left (\vert v \rangle \right )\right \} \nonumber  \\
 &  =0 +\overline{ \langle \left .u\right \vert v \rangle } I_{\mathcal{F}} + \langle \left .u\right \vert v \rangle  I_{\mathcal{F}} +0 =\ensuremath{\operatorname*{Re}}  \langle \left .u\right \vert v \rangle  I_{\mathcal{F}}\text{.}\end{align}This leads to
\begin{align}\left \{X \left (\vert u \rangle \right ) ,X \left (\vert u \rangle \right )\right \} &  = \langle \left .u\right \vert u \rangle  I_{\mathcal{F}} \\
\left \{X \left (\vert u \rangle \right ) ,Y \left (\vert u \rangle \right )\right \} &  =0 \\
\left \{Y \left (\vert u \rangle \right ) ,Y \left (\vert u \rangle \right )\right \} &  = \langle \left .u\right \vert u \rangle  I_{\mathcal{F}}\text{,}\end{align}and in particular for a conb $\left \{\left .\vert \varphi _{i} \rangle \right \vert 1 \leq j \leq r\right \}$
\begin{align}\left \{X \left (\vert \varphi _{i} \rangle \right ) ,X \left (\vert \varphi _{j} \rangle \right )\right \} &  =\delta _{i j} I_{\mathcal{F}} \\
\left \{X \left (\vert \varphi _{i} \rangle \right ) ,Y \left (\vert \varphi _{j} \rangle \right )\right \} &  =0 \\
\left \{Y \left (\vert \varphi _{i} \rangle \right ) ,Y \left (\vert \varphi _{j} \rangle \right )\right \} &  =\delta _{i j} I_{\mathcal{F}}\text{.}\end{align}As the discrete field operators $\left \{a_{j}^{ \dag } ,a_{j}\right \}$ are defined by
\begin{equation}a_{j}^{ \dag } =\Phi ^{ \dag } \left (\vert \varphi _{j} \rangle \right )\text{,}
\end{equation}one has
\begin{align}x_{j} &  =\frac{1}{\sqrt{2}} \left (a_{j}^{ \dag } +a_{j}\right ) =x_{j}^{ \dag } =X \left (\vert \varphi _{i} \rangle \right ) \\
x_{j +r} &  =\frac{i}{\sqrt{2}} \left (a_{j}^{ \dag } -a_{j}\right ) =x_{j +r}^{ \dag } =Y \left (\vert \varphi _{i} \rangle \right ) =X \left (i\vert \varphi _{i} \rangle \right )\text{,}\end{align}$\lceil $
\begin{align}x_{j}^{2} &  =\frac{1}{2} \left (a_{j}^{ \dag } +a_{j}\right ) \left (a_{j}^{ \dag } +a_{j}\right ) =\frac{1}{2} \left (a_{j} a_{j}^{ \dag } +a_{j}^{ \dag } a_{j}\right ) =\frac{1}{2} I_{\mathcal{F}} \\
x_{j +r}^{2} &  = -\frac{1}{2} \left (a_{j}^{ \dag } -a_{j}\right ) \left (a_{j}^{ \dag } -a_{j}\right ) =\frac{1}{2} \left (a_{j} a_{j}^{ \dag } +a_{j}^{ \dag } a_{j}\right ) =\frac{1}{2} I_{\mathcal{F}}\text{.}\end{align}$\rfloor $ 

Summarizing one has
\begin{align}x_{j} &  =\frac{1}{\sqrt{2}} \left (a_{j}^{ \dag } +a_{j}\right ) \\
x_{j +r} &  =\frac{i}{\sqrt{2}} \left (a_{j}^{ \dag } -a_{j}\right )\end{align}and the inverse relationships
\begin{align}a_{j}^{ \dag } &  =\frac{1}{\sqrt{2}} \left (x_{j} -i x_{j +r}\right ) \\
a_{j} &  =\frac{1}{\sqrt{2}} \left (x_{j} +i x_{j +r}\right )\text{.}\end{align}$\lceil $
\begin{align}a_{j}^{ \dag } a_{k}^{ \dag } &  =\frac{1}{\sqrt{2}} \left (x_{j} -i x_{j +r}\right ) \frac{1}{\sqrt{2}} \left (x_{k} -i x_{k +r}\right ) =\frac{1}{2} \left (x_{j} x_{k} -i x_{j +r} x_{k} -i x_{j} x_{k +r} +x_{j +r} x_{k +r}\right ) \\
 &  =\frac{1}{2} \left (x_{j} x_{k} +x_{j +r} x_{k +r} -i \left (x_{j +r} x_{k} +x_{j} x_{k +r}\right )\right )\end{align}$\rfloor $ 

This can be checked by
\begin{align}\frac{1}{\sqrt{2}} \left (a_{j}^{ \dag } +a_{j}\right ) &  =\frac{1}{2} \left (\left (x_{j} -i x_{j +r}\right ) +\left (x_{j} +i x_{j +r}\right )\right ) =x_{j} \nonumber  \\
\frac{i}{\sqrt{2}} \left (a_{j}^{ \dag } -a_{j}\right ) &  =\frac{i}{2} \left (\left (x_{j} -i x_{j +r}\right ) -\left (x_{j} +i x_{j +r}\right )\right ) =x_{j +s}\text{.}\end{align}

In matrix notation the transformation between the $\left \{\left .a_{j} ,a_{j}^{ \dag }\right \vert 1 \leq j \leq r\right \}$ and $\left \{\left .x_{j}\right \vert 1 \leq j \leq 2 r\right \}$ bases is given by
\begin{align}\left (\mathbf{x}\right ) &  =\left (\begin{array}{cc}\mathbf{x}_{1} & \mathbf{x}_{2}\end{array}\right ) =K \left (\begin{array}{cc}\mathbf{a}^{ \dag } & \mathbf{a}\end{array}\right ) =\left (\begin{array}{cc}K \left (\mathbf{a}^{ \dag }\right ) & K \left (\mathbf{a}\right )\end{array}\right ) \\
 &  =\left (\begin{array}{cc}\mathbf{a}^{ \dag } & \mathbf{a}\end{array}\right ) \mathbb{K} =\left (\begin{array}{cc}\mathbf{a}^{ \dag } & \mathbf{a}\end{array}\right ) \left (\begin{array}{cc}\frac{1}{\sqrt{2}} \mathbb{I}_{r} & \frac{i}{\sqrt{2}} \mathbb{I}_{r} \\
\frac{1}{\sqrt{2}} \mathbb{I}_{r} &  -\frac{i}{\sqrt{2}} \mathbb{I}_{r}\end{array}\right )\text{,}\end{align}where
\begin{equation}\mathbb{K} =\left (\begin{array}{cc}\frac{1}{\sqrt{2}} \mathbb{I}_{r} & \frac{i}{\sqrt{2}} \mathbb{I}_{r} \\
\frac{1}{\sqrt{2}} \mathbb{I}_{r} &  -\frac{i}{\sqrt{2}} \mathbb{I}_{r}\end{array}\right ) \in U \left (2 r ,\mathbb{R}\right )\text{,}
\end{equation}i.e.
\begin{align}\left (\begin{array}{cc}\frac{1}{\sqrt{2}} \mathbb{I}_{r} & \frac{i}{\sqrt{2}} \mathbb{I}_{r} \\
\frac{1}{\sqrt{2}} \mathbb{I}_{r} &  -\frac{i}{\sqrt{2}} \mathbb{I}_{r}\end{array}\right )^{ \dag } \left (\begin{array}{cc}\frac{1}{\sqrt{2}} \mathbb{I}_{r} & \frac{i}{\sqrt{2}} \mathbb{I}_{r} \\
\frac{1}{\sqrt{2}} \mathbb{I}_{r} &  -\frac{i}{\sqrt{2}} \mathbb{I}_{r}\end{array}\right ) &  = \\
\left (\begin{array}{cc}\frac{1}{\sqrt{2}} \mathbb{I}_{r} & \frac{1}{\sqrt{2}} \mathbb{I}_{r} \\
\frac{ -i}{\sqrt{2}} \mathbb{I}_{r} & \frac{i}{\sqrt{2}} \mathbb{I}_{r}\end{array}\right ) \left (\begin{array}{cc}\frac{1}{\sqrt{2}} \mathbb{I}_{r} & \frac{i}{\sqrt{2}} \mathbb{I}_{r} \\
\frac{1}{\sqrt{2}} \mathbb{I}_{r} &  -\frac{i}{\sqrt{2}} \mathbb{I}_{r}\end{array}\right ) &  =\left (\begin{array}{cc}\mathbb{I}_{r} & \mathbb{O}_{r} \\
\mathbb{O}_{r} & \mathbb{I}_{r}\end{array}\right )\text{,}\end{align}\emph{but is NOT a canonical transformation }as it is not a\emph{\ }$ \ast $- map\emph{. } 

\subsection{$K$ is a Linear Map but Not an Algebra Map}
The $K$ transformation is
\begin{equation}\left (\mathbf{x}\right ) =\left (\begin{array}{cc}\mathbf{x}_{1} & \mathbf{x}_{2}\end{array}\right ) =K \left (\begin{array}{cc}\mathbf{a}^{ \dag } & \mathbf{a}\end{array}\right ) =\left (\begin{array}{cc}K \left (\mathbf{a}^{ \dag }\right ) & K \left (\mathbf{a}\right )\end{array}\right ) =\left (\begin{array}{cc}\mathbf{a}^{ \dag } & \mathbf{a}\end{array}\right ) \mathbb{K}\text{,}
\end{equation}so in particular
\begin{align}x_{k} &  =K \left (a_{k}^{ \dag }\right ) ;1 \leq k \leq r \\
x_{k +r} &  =K \left (a_{k}\right ) ;1 \leq k \leq r\end{align}thus
\begin{equation}x_{1} x_{1} =K \left (a_{1}^{ \dag }\right ) K \left (a_{1}^{ \dag }\right ) =\frac{1}{2} \mathbb{I}_{r} \neq K \left (a_{1}^{ \dag } a_{1}^{ \dag }\right ) =0.
\end{equation}However it is algebraic on products of different basis elements as
\begin{align}a_{1}^{ \dag } a_{2}^{ \dag } &  =a_{1}^{ \dag } \wedge a_{2}^{ \dag } \\
x_{1} x_{2} &  =K \left (a_{1}^{ \dag }\right ) K \left (a_{2}^{ \dag }\right ) =K \left (a_{1}^{ \dag }\right ) \wedge K \left (a_{2}^{ \dag }\right ) =\frac{1}{\sqrt{2}} \left (a_{1}^{ \dag } +a_{1}\right ) \frac{1}{\sqrt{2}} \left (a_{2}^{ \dag } +a_{2}\right ) =\frac{1}{2} \left (a_{1}^{ \dag } a_{2}^{ \dag } +a_{1}^{ \dag } a_{2} +a_{1} a_{2}^{ \dag } +a_{1} a_{2}\right )\end{align}and
\begin{align}\left (K \boxtimes K\right ) \left (a_{1}^{ \dag } a_{2}^{ \dag }\right ) &  =K \left (a_{1}^{ \dag }\right ) \wedge K \left (a_{2}^{ \dag }\right ) =\frac{1}{\sqrt{2}} \left (a_{1}^{ \dag } +a_{1}\right ) \wedge \frac{1}{\sqrt{2}} \left (a_{2}^{ \dag } +a_{2}\right ) \\
 &  =\frac{1}{2} \left (a_{1}^{ \dag } a_{2}^{ \dag } +a_{1}^{ \dag } a_{2} +a_{1} a_{2}^{ \dag } +a_{1} a_{2}\right ) =K \left (a_{1}^{ \dag }\right ) K \left (a_{2}^{ \dag }\right )\text{.}\end{align}One should note though that even on products of different basis elements it is not a $ \ast $-algebra map as
\begin{align}K \left (a_{1}\right )^{ \dag } &  \neq K \left (a_{1}^{ \dag }\right ) \\
K \left (a_{1}\right )^{ \dag } &  =x_{3} \\
K \left (a_{1}^{ \dag }\right ) &  =x_{1}\text{.}\end{align}

\subsection{Non Definite Bases}
One should also note that one can define bases formed by a mixture of self and anti self adjoint operators
\begin{equation}\left \{\left .x_{j} ,x_{j +r}\right \vert 1 \leq j \leq p\right \} ,\left \{\left .i x_{j +p} ,i x_{j +p +r}\right \vert 1 \leq j \leq q\right \}\text{,}
\end{equation}where
\begin{align}\left (x_{j}\right )^{2} &  =\left (x_{j +r}\right )^{2} =\frac{1}{2} \mathbb{I}_{r} ;1 \leq j \leq p \\
\left (i x_{j +p}\right )^{2} &  =\left (i x_{j +p +r}\right )^{2} = -\frac{1}{2} \mathbb{I}_{r} ;1 \leq j \leq q\end{align}that correspond to the unitary transformations
\begin{equation}\mathbb{K}_{p ,q} =\left (\begin{array}{cc}\frac{1}{\sqrt{2}} \mathbb{I}_{p ,q} & \frac{i}{\sqrt{2}} \mathbb{I}_{p ,q} \\
\frac{1}{\sqrt{2}} \mathbb{I}_{p ,q} &  -\frac{i}{\sqrt{2}} \mathbb{I}_{p ,q}\end{array}\right )\text{,}
\end{equation}where
\begin{align}\left (\begin{array}{cc}\frac{1}{\sqrt{2}} \mathbb{I}_{p ,q} & \frac{i}{\sqrt{2}} \mathbb{I}_{p ,q} \\
\frac{1}{\sqrt{2}} \mathbb{I}_{p ,q} &  -\frac{i}{\sqrt{2}} \mathbb{I}_{p ,q}\end{array}\right )^{ \dag } \left (\begin{array}{cc}\frac{1}{\sqrt{2}} \mathbb{I}_{p ,q} & \frac{i}{\sqrt{2}} \mathbb{I}_{p ,q} \\
\frac{1}{\sqrt{2}} \mathbb{I}_{p ,q} &  -\frac{i}{\sqrt{2}} \mathbb{I}_{p ,q}\end{array}\right ) &  = \\
\left (\begin{array}{cc}\frac{1}{\sqrt{2}} \mathbb{I}_{p ,q} & \frac{1}{\sqrt{2}} \mathbb{I}_{p ,q} \\
\frac{ -i}{\sqrt{2}} \mathbb{I}_{p ,q} & \frac{i}{\sqrt{2}} \mathbb{I}_{p ,q}\end{array}\right ) \left (\begin{array}{cc}\frac{1}{\sqrt{2}} \mathbb{I}_{p ,q} & \frac{i}{\sqrt{2}} \mathbb{I}_{p ,q} \\
\frac{1}{\sqrt{2}} \mathbb{I}_{p ,q} &  -\frac{i}{\sqrt{2}} \mathbb{I}_{p ,q}\end{array}\right ) &  =\left (\begin{array}{cc}\mathbb{I}_{r} & \mathbb{O}_{r} \\
\mathbb{O}_{r} & \mathbb{I}_{r}\end{array}\right )\text{,}\end{align}and
\begin{equation}\mathbb{I}_{p ,q} =\left (\begin{array}{cc}\mathbb{I}_{p} & \mathbb{O} \\
\mathbb{O} &  -\mathbb{I}_{q}\end{array}\right )\text{and}p +q =r\text{.}
\end{equation}

\subsubsection{Null Bases}
One can construct \emph{Null bases }elements\emph{\ }$\left \{\left .\nu _{j}\right \vert 1 \leq j \leq 2 \max \left (p ,q\right )\right \}$ wrt $ \langle \left .\text{}\right \vert  \rangle _{p ,q}$ and $\left (\left .\text{}\right \vert \right )_{p ,q}$ that are selfadjoint and nilpotent
\begin{align}\nu _{j} &  =\frac{1}{\sqrt{2}} \left (x_{j} +x_{j +\max \left (p ,q\right )/2}\right ) \\
\left (\nu _{j}\right )^{2} &  =\frac{1}{2} \left (x_{j} +x_{j +\max \left (p ,q\right )/2}\right ) \left (x_{j} +x_{j +\max \left (p ,q\right )/2}\right ) =0 \\
\nu _{j +p} &  =\frac{1}{\sqrt{2}} \left (x_{j} -x_{j +\max \left (p ,q\right )/2}\right ) \\
\left (\nu _{j +p}\right )^{2} &  =\frac{1}{2} \left (x_{j} -x_{j +\max \left (p ,q\right )/2}\right ) \left (x_{j} -x_{j +\max \left (p ,q\right )/2}\right ) =0 \\
\nu _{2 \max \left (p ,q\right ) +j} &  =x_{2 \max \left (p ,q\right ) +j}\end{align}Nilpotenticy
\begin{equation}\alpha ^{2} =0
\end{equation}implies that
\begin{equation}\left (\left .\alpha \right \vert \alpha \right )_{p ,q} =0
\end{equation}but not necessarily
\begin{equation} \langle \left .\alpha \right \vert \alpha  \rangle _{p ,q} =0\text{,}
\end{equation}unless
\begin{equation}\alpha  =\alpha ^{ \dag }\text{.}
\end{equation}

One can construct \emph{Null bases }wrt\emph{\ }$\left (\left .\text{}\right \vert \right )_{p ,q}$ that \emph{are not }selfadjoint but are \emph{nilpotent} e.g. in the case $p =2 r$ the creation and annihiltion operators
\begin{align}\left (a_{j}^{ \dag }\right )^{2} &  =\frac{1}{\sqrt{2}} \left (x_{j} +i x_{j +s}\right ) \frac{1}{\sqrt{2}} \left (x_{j} +i x_{j +s}\right ) =\frac{1}{2} \left (x_{j}^{2} +i x_{j} x_{j +s} +i x_{j +s} x_{j} -x_{j +s}^{2}\right ) =0 \\
\left (a_{j}\right )^{2} &  =\frac{1}{\sqrt{2}} \left (x_{j} -i x_{j +s}\right ) \frac{1}{\sqrt{2}} \left (x_{j} -i x_{j +s}\right ) =\frac{1}{2} \left (x_{j}^{2} -i x_{j} x_{j +s} -i x_{j +s} x_{j} -x_{j +s}^{2}\right ) =0.\end{align}

\subsection{Transformations Between Self Adjoint Bases}
\bigskip The self adjoint basis produced by the canonically transformed $\left \{\left .a_{j}^{ \dag } ,a_{j}\right \vert 1 \leq j \leq r\right \}$ is
\begin{align*}\left (\begin{array}{cc}\mathbf{y}_{1} & \mathbf{y}_{2}\end{array}\right ) &  =K \left (\begin{array}{cc}\mathbf{b}^{ \dag } & \mathbf{b}\end{array}\right ) =\left (\begin{array}{cc}\mathbf{b}^{ \dag } & \mathbf{b}\end{array}\right ) \left (\begin{array}{cc}\frac{1}{\sqrt{2}} \mathbb{I}_{r} & \frac{i}{\sqrt{2}} \mathbb{I}_{r} \\
\frac{1}{\sqrt{2}} \mathbb{I}_{r} &  -\frac{i}{\sqrt{2}} \mathbb{I}_{r}\end{array}\right ) \\
 &  =\left (\begin{array}{cc}\mathbf{a}^{ \dag } & \mathbf{a}\end{array}\right ) \left (\begin{array}{cc}\frac{1}{\sqrt{2}} \mathbb{I}_{r} & \frac{i}{\sqrt{2}} \mathbb{I}_{r} \\
\frac{1}{\sqrt{2}} \mathbb{I}_{r} &  -\frac{i}{\sqrt{2}} \mathbb{I}_{r}\end{array}\right ) \left (\begin{array}{cc}\frac{1}{\sqrt{2}} \mathbb{I}_{r} & \frac{i}{\sqrt{2}} \mathbb{I}_{r} \\
\frac{1}{\sqrt{2}} \mathbb{I}_{r} &  -\frac{i}{\sqrt{2}} \mathbb{I}_{r}\end{array}\right )^{ \dag } \left (\begin{array}{cc}\mathbb{U} & \overline{\mathbb{V}} \\
\mathbb{V} & \overline{\mathbb{U}}\end{array}\right ) \left (\begin{array}{cc}\frac{1}{\sqrt{2}} \mathbb{I}_{r} & \frac{i}{\sqrt{2}} \mathbb{I}_{r} \\
\frac{1}{\sqrt{2}} \mathbb{I}_{r} &  -\frac{i}{\sqrt{2}} \mathbb{I}_{r}\end{array}\right ) \\
 &  =\left (\begin{array}{cc}\mathbf{x}_{1} & \mathbf{x}_{2}\end{array}\right ) \left (\begin{array}{cc}\frac{1}{\sqrt{2}} \mathbb{I}_{r} & \frac{i}{\sqrt{2}} \mathbb{I}_{r} \\
\frac{1}{\sqrt{2}} \mathbb{I}_{r} &  -\frac{i}{\sqrt{2}} \mathbb{I}_{r}\end{array}\right )^{ \dag } \left (\begin{array}{cc}\mathbb{U} & \overline{\mathbb{V}} \\
\mathbb{V} & \overline{\mathbb{U}}\end{array}\right ) \left (\begin{array}{cc}\frac{1}{\sqrt{2}} \mathbb{I}_{r} & \frac{i}{\sqrt{2}} \mathbb{I}_{r} \\
\frac{1}{\sqrt{2}} \mathbb{I}_{r} &  -\frac{i}{\sqrt{2}} \mathbb{I}_{r}\end{array}\right )\end{align*}


\begin{align} &  =\left (\begin{array}{cc}\mathbf{x}_{1} & \mathbf{x}_{2}\end{array}\right ) \left (\begin{array}{cc}\frac{1}{\sqrt{2}} \mathbb{U} +\frac{1}{\sqrt{2}} \mathbb{V} & \frac{1}{\sqrt{2}} \overline{\mathbb{V}} +\frac{1}{\sqrt{2}} \overline{\mathbb{U}} \\
\frac{ -i}{\sqrt{2}} \mathbb{U} +\frac{i}{\sqrt{2}} \mathbb{V} & \frac{ -i}{\sqrt{2}} \overline{\mathbb{V}} +\frac{i}{\sqrt{2}} \overline{\mathbb{U}}\end{array}\right ) \left (\begin{array}{cc}\frac{1}{\sqrt{2}} \mathbb{I}_{r} & \frac{i}{\sqrt{2}} \mathbb{I}_{r} \\
\frac{1}{\sqrt{2}} \mathbb{I}_{r} &  -\frac{i}{\sqrt{2}} \mathbb{I}_{r}\end{array}\right ) \nonumber  \\
 &  =\left (\begin{array}{cc}\mathbf{x}_{1} & \mathbf{x}_{2}\end{array}\right ) \left (\begin{array}{cc}\frac{1}{2} \mathbb{U} +\frac{1}{2} \mathbb{V} +\frac{1}{2} \overline{\mathbb{V}} +\frac{1}{2} \overline{\mathbb{U}} & \frac{i}{2} \mathbb{U} -\frac{i}{2} \mathbb{V} +\frac{i}{2} \overline{\mathbb{V}} -\frac{i}{2} \overline{\mathbb{U}} \\
\frac{ -i}{2} \mathbb{U} +\frac{i}{2} \mathbb{V} -\frac{i}{2} \overline{\mathbb{V}} +\frac{i}{2} \overline{\mathbb{U}} & \frac{1}{2} \mathbb{U} +\frac{1}{2} \mathbb{V} -\frac{1}{2} \overline{\mathbb{V}} +\frac{1}{2} \overline{\mathbb{U}}\end{array}\right ) \nonumber  \\
 &  =\left (\begin{array}{cc}\mathbf{x}_{1} & \mathbf{x}_{2}\end{array}\right ) \left (\begin{array}{cc}\ensuremath{\operatorname*{Re}} \left (\mathbb{U} +\mathbb{V}\right ) & \ensuremath{\operatorname*{Im}} \left (\mathbb{V} +\mathbb{U}\right ) \\
\ensuremath{\operatorname*{Im}} \left (\mathbb{V} -\mathbb{U}\right ) & \ensuremath{\operatorname*{Re}} \left (\mathbb{U} -\mathbb{V}\right )\end{array}\right )\text{.}\end{align}

i.e.
\begin{equation}\left (\begin{array}{cc}\mathbf{y}_{1} & \mathbf{y}_{2}\end{array}\right ) =\left (\begin{array}{cc}\mathbf{x}_{1} & \mathbf{x}_{2}\end{array}\right ) \left (\begin{array}{cc}\ensuremath{\operatorname*{Re}} \left (\mathbb{U} +\mathbb{V}\right ) & \ensuremath{\operatorname*{Im}} \left (\mathbb{V} +\mathbb{U}\right ) \\
\ensuremath{\operatorname*{Im}} \left (\mathbb{V} -\mathbb{U}\right ) & \ensuremath{\operatorname*{Re}} \left (\mathbb{U} -\mathbb{V}\right )\end{array}\right )\text{.}
\end{equation}Thus as
\begin{align}\left (\begin{array}{cc}\mathbf{x}_{1} & \mathbf{x}_{2}\end{array}\right ) &  =K \left (\begin{array}{cc}\mathbf{a}^{ \dag } & \mathbf{a}\end{array}\right ) =\left (\begin{array}{cc}\mathbf{a}^{ \dag } & \mathbf{a}\end{array}\right ) \mathbb{K} \\
\left (\begin{array}{cc}\mathbf{y}_{1} & \mathbf{y}_{2}\end{array}\right ) &  =K \left (\begin{array}{cc}\mathbf{b}^{ \dag } & \mathbf{b}\end{array}\right ) =\left (\begin{array}{cc}\mathbf{b}^{ \dag } & \mathbf{b}\end{array}\right ) \mathbb{K}\end{align}one obtains
\begin{align}\left (\begin{array}{cc}\mathbf{y}_{1} & \mathbf{y}_{2}\end{array}\right ) &  =\left (\begin{array}{cc}\mathbf{x}_{1} & \mathbf{x}_{2}\end{array}\right ) \left (\begin{array}{cc}\ensuremath{\operatorname*{Re}} \left (\mathbb{U} +\mathbb{V}\right ) & \ensuremath{\operatorname*{Im}} \left (\mathbb{V} +\mathbb{U}\right ) \\
\ensuremath{\operatorname*{Im}} \left (\mathbb{V} -\mathbb{U}\right ) & \ensuremath{\operatorname*{Re}} \left (\mathbb{U} -\mathbb{V}\right )\end{array}\right ) \\
 &  =\left (\begin{array}{cc}\mathbf{a}^{ \dag } & \mathbf{a}\end{array}\right ) \mathbb{K} \mathbb{K}^{ \dag } \left (\begin{array}{cc}\mathbb{U} & \overline{\mathbb{V}} \\
\mathbb{V} & \overline{\mathbb{U}}\end{array}\right ) \mathbb{K}\end{align}and hence
\begin{equation}\left (\begin{array}{cc}\ensuremath{\operatorname*{Re}} \left (\mathbb{U} +\mathbb{V}\right ) & \ensuremath{\operatorname*{Im}} \left (\mathbb{V} +\mathbb{U}\right ) \\
\ensuremath{\operatorname*{Im}} \left (\mathbb{V} -\mathbb{U}\right ) & \ensuremath{\operatorname*{Re}} \left (\mathbb{U} -\mathbb{V}\right )\end{array}\right ) =\mathbb{K}^{ \dag } \left (\begin{array}{cc}\mathbb{U} & \overline{\mathbb{V}} \\
\mathbb{V} & \overline{\mathbb{U}}\end{array}\right ) \mathbb{K}\text{.}
\end{equation}As both $\left (\begin{array}{cc}\mathbb{U} & \overline{\mathbb{V}} \\
\mathbb{V} & \overline{\mathbb{U}}\end{array}\right )$ and $\left (\begin{array}{cc}\frac{1}{\sqrt{2}} \mathbb{I}_{r} & \frac{i}{\sqrt{2}} \mathbb{I}_{r} \\
\frac{1}{\sqrt{2}} \mathbb{I}_{r} &  -\frac{i}{\sqrt{2}} \mathbb{I}_{r}\end{array}\right )$ are unitary one has that
\begin{align} & \left (\begin{array}{cc}\ensuremath{\operatorname*{Re}} \left (\mathbb{U} +\mathbb{V}\right ) & \ensuremath{\operatorname*{Im}} \left (\mathbb{V} +\mathbb{U}\right ) \\
\ensuremath{\operatorname*{Im}} \left (\mathbb{V} -\mathbb{U}\right ) & \ensuremath{\operatorname*{Re}} \left (\mathbb{U} -\mathbb{V}\right )\end{array}\right ) \left (\begin{array}{cc}\ensuremath{\operatorname*{Re}} \left (\mathbb{U} +\mathbb{V}\right ) & \ensuremath{\operatorname*{Im}} \left (\mathbb{V} +\mathbb{U}\right ) \\
\ensuremath{\operatorname*{Im}} \left (\mathbb{V} -\mathbb{U}\right ) & \ensuremath{\operatorname*{Re}} \left (\mathbb{U} -\mathbb{V}\right )\end{array}\right )^{ \dag } \nonumber  \\
 &  =\left (\begin{array}{cc}\ensuremath{\operatorname*{Re}} \left (\mathbb{U} +\mathbb{V}\right ) & \ensuremath{\operatorname*{Im}} \left (\mathbb{V} +\mathbb{U}\right ) \\
\ensuremath{\operatorname*{Im}} \left (\mathbb{V} -\mathbb{U}\right ) & \ensuremath{\operatorname*{Re}} \left (\mathbb{U} -\mathbb{V}\right )\end{array}\right )^{ \dag } \left (\begin{array}{cc}\ensuremath{\operatorname*{Re}} \left (\mathbb{U} +\mathbb{V}\right ) & \ensuremath{\operatorname*{Im}} \left (\mathbb{V} +\mathbb{U}\right ) \\
\ensuremath{\operatorname*{Im}} \left (\mathbb{V} -\mathbb{U}\right ) & \ensuremath{\operatorname*{Re}} \left (\mathbb{U} -\mathbb{V}\right )\end{array}\right ) \nonumber  \\
 &  =\mathbb{I}_{2 r}\end{align}and as $\left (\begin{array}{cc}\ensuremath{\operatorname*{Re}} \left (\mathbb{U} +\mathbb{V}\right ) & \ensuremath{\operatorname*{Im}} \left (\mathbb{V} +\mathbb{U}\right ) \\
\ensuremath{\operatorname*{Im}} \left (\mathbb{V} -\mathbb{U}\right ) & \ensuremath{\operatorname*{Re}} \left (\mathbb{U} -\mathbb{V}\right )\end{array}\right )$ is a real matrix
\begin{equation}\left (\begin{array}{cc}\ensuremath{\operatorname*{Re}} \left (\mathbb{U} +\mathbb{V}\right ) & \ensuremath{\operatorname*{Im}} \left (\mathbb{V} +\mathbb{U}\right ) \\
\ensuremath{\operatorname*{Im}} \left (\mathbb{V} -\mathbb{U}\right ) & \ensuremath{\operatorname*{Re}} \left (\mathbb{U} -\mathbb{V}\right )\end{array}\right )^{ \dag } =\left (\begin{array}{cc}\ensuremath{\operatorname*{Re}} \left (\mathbb{U} +\mathbb{V}\right ) & \ensuremath{\operatorname*{Im}} \left (\mathbb{V} +\mathbb{U}\right ) \\
\ensuremath{\operatorname*{Im}} \left (\mathbb{V} -\mathbb{U}\right ) & \ensuremath{\operatorname*{Re}} \left (\mathbb{U} -\mathbb{V}\right )\end{array}\right )^{t}
\end{equation}one has
\begin{equation}\left (\begin{array}{cc}\ensuremath{\operatorname*{Re}} \left (\mathbb{U} +\mathbb{V}\right ) & \ensuremath{\operatorname*{Im}} \left (\mathbb{V} +\mathbb{U}\right ) \\
\ensuremath{\operatorname*{Im}} \left (\mathbb{V} -\mathbb{U}\right ) & \ensuremath{\operatorname*{Re}} \left (\mathbb{U} -\mathbb{V}\right )\end{array}\right ) \in O \left (2 r ,\mathbb{R}\right )\text{,}
\end{equation}i.e. \emph{the group of canonical transformations is }$O \left (2 r ,\mathbb{R}\right ) \subset U \left (2 r\right )\text{.}$ It should be observed that $O \left (2 r ,\mathbb{R}\right ) \subset $ $O \left (2 r ,\mathbb{C}\right )$ and in fact $O \left (2 r ,\mathbb{R}\right ) =U \left (2 r\right ) \cap $ $O \left (2 r ,\mathbb{C}\right )\text{.}$ There are other transformation of $\mathcal{F}_{1}$ that preserve the CAR's that are not canonical, in particular $\mathbb{K}_{r} =\left (\begin{array}{cc}\frac{1}{\sqrt{2}} \mathbb{I}_{r} & \frac{i}{\sqrt{2}} \mathbb{I}_{r} \\
\frac{1}{\sqrt{2}} \mathbb{I}_{r} &  -\frac{i}{\sqrt{2}} \mathbb{I}_{r}\end{array}\right ) \in U \left (2 r\right ) , \notin O \left (2 r ,\mathbb{R}\right )\text{,}$ and any $T \in O \left (2 r ,\mathbb{C}\right )$ but $T \notin O \left (2 r ,\mathbb{R}\right )$ that are not $ \ast  -$ maps i.e.
\begin{equation}T \left (x_{j}\right )^{ \dag } \neq T \left (x_{j}^{ \dag }\right )\text{.}
\end{equation}

\emph{Note that the CAR's with indefinite signiture are \textbf{also conserved by}}$O \left (2 r ,\mathbb{C}\right )\text{,}$ I\ conjecture that they are canonical if they $ \in O \left (p ,q ,\mathbb{R}\right )$ $\text{.}$ 

The action of $O \left (2 r ,\mathbb{R}\right )$ on $\mathcal{F}$ can be expressed in terms of inner automorphisms i.e. 

\begin{equation}y_{j} =\hat{O} \left (x_{j}\right ) =O \left (\mathbf{\lambda }\right ) x_{j} O \left (\mathbf{\lambda }\right )^{t} =\exp  \left \{\sum _{1 \leq k <l \leq 2 r}\lambda _{k l} x_{k} x_{l}\right \} x_{j} \exp  \left \{ -\sum _{1 \leq k <l \leq 2 r}\lambda _{k l} x_{k} x_{l}\right \}\text{,}
\end{equation}where
\begin{equation}\mathbf{\lambda }  = -\mathbf{\lambda }^{t} ,\,\lambda _{j k} \in \mathbb{R} ;\mathbf{\lambda } \in M \left (2 r ,\mathbb{R}\right )\text{,}
\end{equation}thus
\begin{align}\left (\begin{array}{cc}\mathbf{y}_{1} & \mathbf{y}_{2}\end{array}\right ) &  =\widehat{O \left (\mathbf{\lambda }\right )} \left (\begin{array}{cc}\mathbf{x}_{1} & \mathbf{x}_{2}\end{array}\right ) =\left (\begin{array}{cc}\mathbf{x}_{1} & \mathbf{x}_{2}\end{array}\right ) \mathbb{T} =\left (\begin{array}{cc}\mathbf{x}_{1} & \mathbf{x}_{2}\end{array}\right ) \left (\begin{array}{cc}\ensuremath{\operatorname*{Re}} \left (\mathbb{U} +\mathbb{V}\right ) & \ensuremath{\operatorname*{Im}} \left (\mathbb{V} +\mathbb{U}\right ) \\
\ensuremath{\operatorname*{Im}} \left (\mathbb{V} -\mathbb{U}\right ) & \ensuremath{\operatorname*{Re}} \left (\mathbb{U} -\mathbb{V}\right )\end{array}\right ) \\
\left (\begin{array}{cc}\mathbf{y}_{1} & \mathbf{y}_{2}\end{array}\right ) &  =\widehat{O \left (\mathbf{\lambda }\right )} \left (\begin{array}{cc}\mathbf{x}_{1} & \mathbf{x}_{2}\end{array}\right ) =\left (\begin{array}{cc}O \left (\mathbf{\lambda }\right ) \mathbf{x}_{1} O \left (\mathbf{\lambda }\right )^{t} & O \left (\mathbf{\lambda }\right ) \mathbf{x}_{2} O \left (\mathbf{\lambda }\right )^{t}\end{array}\right ) \\
 &  =\left (\begin{array}{cc}\mathbf{x}_{1} & \mathbf{x}_{2}\end{array}\right ) \widehat{\mathbb{O}} =\left (\begin{array}{cc}\mathbf{x}_{1} & \mathbf{x}_{2}\end{array}\right ) \exp  \left \{\ensuremath{\operatorname*{adj}} \left (\mathbf{\lambda }\right )\right \}\text{.}\end{align}

\section{Clifford Algebra}
The squares of the self adjoint operators $\left \{\left .x_{j}\right \vert 1 \leq j \leq 2 r\right \}$ are multiples of the identity i.e.
\begin{equation}x_{j}^{2} =\frac{1}{\sqrt{2}} \left (a_{j}^{ \dag } +a_{j}\right ) \frac{1}{\sqrt{2}} \left (a_{j}^{ \dag } +a_{j}\right ) =\frac{1}{2} I_{\mathcal{F}}
\end{equation}
\begin{equation}x_{j +r}^{2} =\frac{i}{\sqrt{2}} \left (a_{j}^{ \dag } -a_{j}\right ) \frac{i}{\sqrt{2}} \left (a_{j}^{ \dag } -a_{j}\right ) =\frac{1}{2} I_{\mathcal{F}}
\end{equation}Thus
\begin{equation}\left \Vert x_{j}\right \Vert  \equiv \left \Vert x_{j}\right \Vert _{\mathcal{B} \left (\mathcal{H}_{\mathcal{F}}\right )} \equiv \left \Vert x_{j}\right \Vert _{\mathcal{F}} =\frac{1}{\sqrt{2}} ,1 \leq j \leq 2 r\text{.}
\end{equation}

The products for $1 \leq j \neq k \leq r$ are
\begin{align}x_{j} x_{k} &  =\frac{1}{\sqrt{2}} \left (a_{j}^{ \dag } +a_{j}\right ) \frac{1}{\sqrt{2}} \left (a_{k}^{ \dag } +a_{k}\right ) =\frac{1}{2} \left (a_{j}^{ \dag } a_{k} +a_{j} a_{k} +a_{j}^{ \dag } a_{k}^{ \dag } +a_{j} a_{k}^{ \dag }\right ) \nonumber  \\
x_{k} x_{j} &  =\frac{1}{\sqrt{2}} \left (a_{k}^{ \dag } +a_{k}\right ) \frac{1}{\sqrt{2}} \left (a_{j}^{ \dag } +a_{j}\right ) =\frac{1}{2} \left (a_{k}^{ \dag } a_{j} +a_{k} a_{j} +a_{k}^{ \dag } a_{j}^{ \dag } +a_{k} a_{j}^{ \dag }\right )\end{align}hence
\begin{equation}x_{j} x_{k} = -x_{k} x_{j}\text{.}
\end{equation}

The products for $r +1 \leq j +r \neq k +r \leq 2 r$ are
\begin{align}x_{j +r} x_{k +r} &  =\frac{i}{\sqrt{2}} \left (a_{j}^{ \dag } -a_{j}\right ) \frac{i}{\sqrt{2}} \left (a_{k}^{ \dag } -a_{k}\right ) =\frac{1}{2} \left ( -a_{j}^{ \dag } a_{k}^{ \dag } +a_{j}^{ \dag } a_{k} +a_{j} a_{k}^{ \dag } -a_{j} a_{k}\right ) \nonumber  \\
x_{k +r} x_{j +r} &  =\frac{i}{\sqrt{2}} \left (a_{k}^{ \dag } -a_{k}\right ) \frac{i}{\sqrt{2}} \left (a_{j}^{ \dag } -a_{j}\right ) =\frac{1}{2} \left ( -a_{k}^{ \dag } a_{j}^{ \dag } +a_{k}^{ \dag } a_{j} +a_{k} a_{j}^{ \dag } -a_{k} a_{j}\right )\end{align}hence
\begin{equation}x_{j +r} x_{k +r} = -x_{k +r} x_{j +r}\text{.}
\end{equation}

The products for $1 \leq j \leq r ;$ $r +1 \leq k +r \leq 2 r$ are
\begin{align}x_{j} x_{k +r} &  =\frac{1}{\sqrt{2}} \left (a_{j}^{ \dag } +a_{j}\right ) \frac{i}{\sqrt{2}} \left (a_{k}^{ \dag } -a_{k}\right ) =\frac{1}{2} \left (i a_{j}^{ \dag } a_{k}^{ \dag } -i a_{j}^{ \dag } a_{k} +i a_{j} a_{k}^{ \dag } -i a_{j} a_{k}\right ) \nonumber  \\
x_{k +r} x_{j} &  =\frac{i}{\sqrt{2}} \left (a_{k}^{ \dag } -a_{k}\right ) \frac{1}{\sqrt{2}} \left (a_{j}^{ \dag } +a_{j}\right ) =\left (x_{j} x_{k +s}\right )^{ \dag } =\frac{1}{2} \left (i a_{k}^{ \dag } a_{j}^{ \dag } +i a_{k}^{ \dag } a_{j} -i a_{k} a_{j}^{ \dag } -i a_{k} a_{j}\right )\end{align}hence
\begin{equation}x_{j} x_{k +r} = -x_{k +r} x_{j}\text{.}
\end{equation}

The squares for $1 \leq j \leq r$ are 


\begin{align}x_{j}^{2} &  =x_{j} x_{j} =\frac{1}{\sqrt{2}} \left (a_{j}^{ \dag } +a_{j}\right ) \frac{1}{\sqrt{2}} \left (a_{j}^{ \dag } +a_{j}\right ) =\frac{1}{2} \left (a_{j} a_{j}^{ \dag } +a_{j} a_{j} +a_{j}^{ \dag } a_{j}^{ \dag } +a_{j}^{ \dag } a_{j}\right ) =\frac{1}{2} I_{\mathcal{F}} \nonumber  \\
x_{j +r}^{2} &  =x_{j +r} x_{j +r} =\frac{i}{\sqrt{2}} \left (a_{j}^{ \dag } -a_{j}\right ) \frac{i}{\sqrt{2}} \left (a_{j}^{ \dag } -a_{j}\right ) \nonumber  \\
 &  = -\frac{1}{2} \left (a_{j}^{ \dag } a_{j}^{ \dag } -a_{j}^{ \dag } a_{j} -a_{j} a_{j}^{ \dag } +a_{j} a_{j}\right ) =\frac{1}{2} I_{\mathcal{F}}\text{.}\end{align}

Thus the self adjoint operators $\left \{\left .x_{j}\right \vert 1 \leq j \leq 2 r\right \}$ satisfy the CAR's
\begin{equation}\left \{x_{j} ,x_{k}\right \} =\delta _{j k} I_{\mathcal{F}} ;\,1 \leq i ,j \leq 2 r\text{.}
\end{equation}

The Fermi Fock algebra $\mathcal{F}\text{,}$ which is representation of a CAR algebra, is generated by \emph{complex} polynomials of the generators $\left \{\left .x_{j}\right \vert 1 \leq j \leq 2 r\right \}$ that satisfy the conditions
\begin{align}x_{j} x_{k} &  = -x_{k} x_{j} ,\,1 \leq j <k \leq 2 r \\
x_{j}^{2} &  =\frac{1}{2} I_{\mathcal{F}} ,\,1 \leq j \leq 2 r\end{align}or equivalently
\begin{equation}\left \{x_{j} ,x_{k}\right \} =\delta _{j k} I_{\mathcal{F}} ;\,1 \leq i ,j \leq 2 r\text{,}
\end{equation}which is a \emph{Clifford Algebra}. This is a $\mathbb{N} -$graded (special case of a \emph{Filtered}) algebra and can be expressed as the exterior algebra over $\mathcal{F}_{1}$ i.e. as the vector space direct sum
\begin{equation}\mathcal{F} =\bigwedge \left (\mathcal{F}_{1}\right ) =\underset{0 \leq n \leq 2 r}{\bigoplus }\mathcal{F}_{n}\text{,}
\end{equation}where $\mathcal{F}_{n}$ are multilinear homogeneous polynomials of degree $n$ in $\left \{x_{j}\right \}$ and the scalar part is defined as
\begin{equation}\mathcal{F}_{0} =\left \{\left .z I_{\mathcal{F}}\right \vert z \in \mathbb{C}\right \}\text{.}
\end{equation}

$\lceil $From Greu$\beta $ one has
\begin{equation}\mathcal{F} =\bigwedge \left (\mathcal{F}_{1}\right ) =\bigwedge \left (\mathfrak{f}_{1} \oplus \mathfrak{f}_{1}^{d}\right ) =\bigwedge \left (\mathfrak{f}_{1}\right ) \otimes \bigwedge  \left (\mathfrak{f}_{1}^{d}\right ) =\mathcal{B} \left (\mathcal{H}_{\mathcal{F}}\right ) =\mathcal{B} \left (\bigwedge \left (\mathfrak{f}_{1}\right )\right )\text{.}
\end{equation}

If $\mathcal{F}_{1}$ and $\mathfrak{f}_{1}$ have inner products i.e. $\mathcal{F}_{1} \equiv \mathcal{F}_{1 H S}$ then $\mathfrak{f}_{1} \cong \mathfrak{f}_{1}^{d}$ and
\begin{equation}\bigwedge \left (\mathfrak{f}_{1}\right ) \otimes \bigwedge  \left (\mathfrak{f}_{1}\right ) \cong \bigwedge \left (\mathfrak{f}_{1}\right ) \otimes \bigwedge  \left (\mathfrak{f}_{1}^{d}\right ) \cong \mathcal{B} \left (\bigwedge \left (\mathfrak{f}_{1}\right )\right )\text{.}
\end{equation}$\rfloor $ 

The Clifford algebra has two fundamental algebraic products, under which it is closed, the antisymmetric wedge
product "$A \wedge B ""$ and the Clifford/Geometric product "$A \circ B \equiv A B ""$. The wedge product can be defined in terms of the Clifford product. 

The subspaces $\left \{\left .\mathcal{F}_{n}\right \vert 0 \leq n \leq 2 r\right \}$ are disjoint and
\begin{equation*}I_{\mathcal{F}} \neq \sum _{1 \leq n \leq 2 r}\sum _{1 \leq j_{1} <\cdots  <j_{n} \leq 2 r}A_{j_{1} \cdots  j_{2 r}} x_{j_{1} \cdots } x_{j_{n}}\text{.}
\end{equation*}One can show this do by considering the operator valued trace and noting that it is zero on elements $n >0.\;$ 

This is a different decomposition to the ones introduced before i.e. in terms of degrees of creation and annihilation
operators
\begin{equation}\mathcal{F} =\underset{0 \leq n ,m \leq r}{\bigoplus }\mathfrak{F}_{n m}
\end{equation}and mappings between the component Hilbert spaces
\begin{equation}\mathcal{F} =\underset{0 \leq n ,m \leq r}{\bigoplus }\mathcal{B}_{n m} \left (\mathcal{H}_{\mathcal{F}}\right ) \equiv \underset{0 \leq n ,m \leq r}{\bigoplus }\mathcal{B} \left (\mathcal{H}^{n} ,\mathcal{H}^{m}\right )\text{.}
\end{equation}One should observe that
\begin{equation}\mathcal{F}_{n} =\underset{\begin{array}{c}p +q =n \\
0 \leq p ,q \leq r\end{array}}{\bigoplus }\mathfrak{F}_{p q}\text{.}
\end{equation}
\begin{equation}\mathcal{F}_{n} =\mathfrak{F}_{0 0} \oplus \tbigoplus ???\mathfrak{F}_{p n -p} ;\text{\quad }2 \leq n \leq r\text{.}
\end{equation}The expansion
\begin{equation}\mathcal{F} =\underset{0 \leq n \leq 2 r}{\bigoplus }\mathcal{F}_{n}
\end{equation}shows that $\mathcal{F}$ is $\mathbb{N} -$graded. One also has the \emph{Super Algebra }$\mathbb{N}_{2} -$grading
\begin{equation}\mathcal{F} =\mathcal{F}_{e v e n} \oplus \mathcal{F}_{o d d}\text{,}
\end{equation}where
\begin{align}\mathcal{F}_{e v e n} &  =\underset{1 \leq n \leq r}{\bigoplus }\mathcal{F}_{2 n} \\
\mathcal{F}_{o d d} &  =\underset{1 \leq n \leq r}{\bigoplus }\mathcal{F}_{2 n -1}\text{.}\end{align}

\bigskip If $A$ is a general element of $\mathcal{F}$ then
\begin{equation}A =\sum _{0 \leq n \leq 2 r} \langle A \rangle _{n} =\sum _{0 \leq n \leq 2 r}A_{n}\text{,}
\end{equation}where
\begin{equation} \langle A \rangle _{n} =A_{n} \in \mathcal{F}_{n}\text{.}
\end{equation}Thus
\begin{equation} \langle A \rangle _{n} =\sum _{1 \leq j_{1} <\cdots  <j_{n} \leq 2 r}A_{n j_{1} \cdots  j_{n}} x_{j_{1}} \cdots  x_{j_{n}} =\sum _{1 \leq j_{1} <\cdots  <j_{n} \leq 2 r}A_{n j_{1} \cdots  j_{n}} x_{j_{1}} \wedge \cdots  \wedge x_{j_{n}}
\end{equation}and
\begin{equation} \langle A \rangle _{0} =\alpha  I_{\mathcal{F}} =A_{0} ;\alpha  \in \mathbb{C}\text{.}
\end{equation}Thus
\begin{equation} \langle A_{n} \rangle _{m} =\delta _{n m} A_{n}\text{.}
\end{equation}

\subsection{Reversion}
\qquad Reversion is defined by
\begin{equation}\left (x_{j_{1}} \cdots  x_{j_{n}}\right )^{t} =x_{j_{n}} \cdots  x_{j_{1}} =\left ( -1\right )^{\binom{n}{2}} x_{j_{1}} \cdots  x_{j_{n}}
\end{equation}and reversion plus complex conjugation by
\begin{equation}\left (x_{j_{1}} \cdots  x_{j_{n}}\right )^{ \dag } =\overline{x_{j_{n}} \cdots  x_{j_{1}}} =\left ( -1\right )^{\binom{n}{2}} \overline{x_{j_{1}} \cdots  x_{j_{n}}}\text{.}
\end{equation}This definition gives
\begin{equation}\left (A B\right )^{ \dag } =B^{ \dag } A^{ \dag }
\end{equation}for arbitary $A$ and $B\text{.}$ 

\textbf{Reversion composed with complex conjugation is actually the operator adjoint operation} as can be
seen by
\begin{align}\left (x_{1} x_{2}\right )^{ \dag } &  =\left (\frac{1}{\sqrt{2}} \left (a_{1}^{ \dag } +a_{1}\right ) \frac{1}{\sqrt{2}} \left (a_{2}^{ \dag } +a_{2}\right )\right )^{ \dag } =\frac{1}{\sqrt{2}} \left (a_{2}^{ \dag } +a_{2}\right ) \frac{1}{\sqrt{2}} \left (a_{1}^{ \dag } +a_{1}\right ) \nonumber  \\
 &  =\frac{1}{2} \left (a_{2}^{ \dag } a_{1}^{ \dag } +a_{2}^{ \dag } a_{1} +a_{2} a_{1}^{ \dag } +a_{2} a_{1}\right ) = -\frac{1}{2} \left (a_{1}^{ \dag } a_{2}^{ \dag } +a_{1}^{ \dag } a_{2} +a_{1} a_{2}^{ \dag } +a_{1} a_{2}\right ) \nonumber  \\
 &  = -x_{2} x_{1}\end{align}and
\begin{align}\left (x_{1} x_{1 +r}\right )^{ \dag } &  =\left (\frac{1}{\sqrt{2}} \left (a_{1}^{ \dag } +a_{1}\right ) \frac{i}{\sqrt{2}} \left (a_{1}^{ \dag } -a_{1}\right )\right )^{ \dag } =\frac{i}{\sqrt{2}} \left (a_{1}^{ \dag } -a_{1}\right ) \frac{1}{\sqrt{2}} \left (a_{1}^{ \dag } +a_{1}\right ) \nonumber  \\
 &  =\frac{i}{2} \left (a_{1}^{ \dag } a_{1}^{ \dag } +a_{1}^{ \dag } a_{1} -a_{1} a_{1}^{ \dag } -a_{1} a_{1}\right ) = -\frac{i}{2} \left (a_{1}^{ \dag } a_{1}^{ \dag } +a_{1}^{ \dag } a_{1} -a_{1} a_{1}^{ \dag } +a_{1} a_{1}\right ) \nonumber  \\
 &  = -x_{1 +r} x_{1}\text{.}\end{align}

\section{Geometric Algebra I}


\subsection{Scaler Product}
The product $C_{1} D_{1}$ of two grade 1 operators (vectors) can be expressed as the sum of a symmetric and antisymmetric part of the operator product in $\mathcal{F}$ i.e.
\begin{equation}C_{1} D_{1} =\frac{1}{2} \left \{C_{1} D_{1} +D_{1} C_{1}\right \} +\frac{1}{2} \left \{C_{1} D_{1} -D_{1} C_{1}\right \} =\frac{1}{2} \left \{C_{1} ,D_{1}\right \} +\frac{1}{2} \left [C_{1} ,D_{1}\right ]
\end{equation}which allows one to define a "scalar product" $( \cdot \vert  \cdot )$
\begin{equation}( \cdot \vert  \cdot ) :\mathcal{F}_{1} \longrightarrow \mathbb{C} I_{\mathcal{F}}
\end{equation}between vectors as
\begin{equation}(C_{1}\vert D_{1}) =\frac{1}{2} \left \{C_{1} ,D_{1}\right \}\text{.}
\end{equation}This scalar product can be extended to all of $\mathcal{F}$ by using the exterior algebra property of $\mathcal{F}$ (i.e. $\mathcal{F} = \wedge \left (\mathcal{F}_{1}\right ))\text{.}$ 

The basis $\left \{\left .x_{j}\right \vert \,1 \leq j \leq 2 r\right \}$ is "orthogonal" with respect to this scalar product i.e.
\begin{equation}(x_{j}\vert x_{k}) =x_{j} \vee x_{k} =\frac{1}{2} \left \{x_{j} x_{k} +x_{k} x_{j}\right \} =\left \{x_{j} ,x_{k}\right \} =\frac{1}{2} \delta _{j k} I_{\mathcal{F}} ;1 \leq j \leq 2 r\text{,}
\end{equation}but not normalized. One can define normalized vectors by
\begin{equation}\tilde{x}_{j} =\sqrt{2} x_{j} ;1 \leq j \leq 2 r\text{,}
\end{equation}but these vectors do not satisfy the CAR as
\begin{equation*}\left \{\tilde{x}_{j} ,\tilde{x}_{k}\right \} =2 \delta _{j k} I_{\mathcal{F}}\text{.}
\end{equation*}

One can expand any two vectors $C_{1}$ and $D_{1}$ in terms of this orthonormal basis as
\begin{align}C_{1} &  =\sum _{1 \leq j \leq 2 r}c_{j} \tilde{x}_{j} \\
D_{1} &  =\sum _{1 \leq j \leq 2 r}d_{j} \tilde{x}_{j}\end{align}and obtain
\begin{equation}(C_{1}\vert D_{1}) =(D_{1}\vert C_{1}) =\sum _{1 \leq j \leq 2 r}c_{j} d_{j} I_{\mathcal{F}} =\left (\left (C_{1}\vert D_{1}\right )\right ) I_{\mathcal{F}}\text{,}
\end{equation}where
\begin{equation}\left (\left (C_{1}\vert D_{1}\right )\right ) =\sum _{1 \leq j \leq 2 r}c_{j} d_{j}\text{.}
\end{equation}

The CAR's can then be expressed in terms of this scalar product as
\begin{equation}\left \{C_{1} ,D_{1}\right \} =\sum _{1 \leq j ,k \leq 2 r}c_{j} d_{k} \left \{\tilde{x}_{j} ,\tilde{x}_{k}\right \}I_{\mathcal{F}} =\sum _{1 \leq j \leq 2 r}c_{j} d_{j} I_{\mathcal{F}} =2 (C_{1}\vert D_{1})\text{.}
\end{equation}

The invariance group of the scalar product $( \cdot \vert  \cdot )$ is $O \left (2 r ,\mathbb{C}\right )\text{.}$ 

One can define a quadratic function $Q$ in terms of this scalar product by
\begin{equation}Q \left (C_{1}\right ) =\left (C_{1}\left \vert C_{1}\right .\right )
\end{equation}and then one can recover the scalar product by using the polarization identity
\begin{equation}\left (C_{1}\left \vert D_{1}\right .\right ) =Q \left (C_{1} +D_{1}\right ) -Q \left (C_{1}\right ) -Q \left (D_{1}\right )\text{.}
\end{equation}

\subsection{Contraction}
Left and right contraction are defined by
\begin{align**}v\rfloor D =\frac{1}{2} \left \{v D -\hat{D} v\right \} \\
v\lfloor D =\frac{1}{2} \left \{D v -v \hat{D}\right \}\text{,}\end{align**}where the \emph{grade involution operation }
\begin{equation*}\widehat{} :D \rightarrow \hat{D}\text{,}
\end{equation*}is defined by
\begin{align*}D &  =\sum _{0 \leq n \leq 2 r}D_{n} \\
\hat{D}_{n} &  =\left ( -1\right )^{n} D_{n} \\
\widehat{D} &  =\sum _{0 \leq n \leq 2 r}\left ( -1\right )^{n} D_{n} =\sum _{0 \leq m \leq r}D_{2 m} -\sum _{1 \leq m \leq r}D_{2 m -1}\text{.}\end{align*}The scalar products
\begin{align}\left (A\left \vert B\right .\right ) &  = \langle A^{t} B \rangle _{0} \\
 \langle A\left \vert B\right . \rangle  &  = \langle A^{ \dag } B \rangle _{0}\end{align}can be expressed in terms of the contraction operation. 

Contraction can be generalized to act between
any elements of the GA by
\begin{align}A\rfloor B =\sum _{0 \leq p ,q \leq 2 r} \langle A_{p} B_{q} \rangle _{p -q} \\
A\lfloor B =\sum _{0 \leq p ,q \leq 2 r} \langle A_{p} B_{q} \rangle _{q -p}\text{.}\end{align}

\subsection{Wedge Product}
The wedge product \emph{between vectors} is defined as the antisymmetric part of the operator product i.e.
\begin{equation*}C_{1} \wedge D_{1} = \langle C_{1} D_{1} \rangle _{2} =\frac{1}{2} \left [C_{1} ,D_{1}\right ]\text{(The last identity works only for vectors, definately does not work for grade 0 elements).}
\end{equation*}In general the wedge product in a Clifford algebra is defined by
\begin{equation}C \wedge D =\sum _{0 \leq m ,n \leq 2 r} \langle C_{m} D_{n} \rangle _{m +n}\text{.}
\end{equation}

In order to investigate the connection between the products $ \circ  , \boxtimes $ and $ \wedge $ we consider the first quantized expansions of the creation and annihilation operators for \underline {\emph{Type I}}\emph{\ }\underline {\emph{vacuum}\emph{representations}} of a CAR algebra, e.g. the Fock algebra $\mathcal{F}\text{,}$
\begin{align}\Pi _{\vert \phi  \rangle  \langle \phi \vert } \left (a_{1}^{ \dag }\right ) &  \equiv a_{1}^{ \dag } =\vert \varphi _{1} \rangle  \langle \phi \vert  +\sum _{2 \leq n \leq r}\sum _{2 \leq j_{1} <\cdots  <j_{n -1} <r}\vert \varphi _{1} \varphi _{j_{1}} \cdots  \varphi _{j_{n -1}} \rangle  \langle \varphi _{j_{1}} \cdots  \varphi _{j_{n -1}}\vert \text{,} \\
a_{1} &  =\vert \phi  \rangle  \langle \varphi _{1}\vert  +\sum _{2 \leq n \leq r}\sum _{2 \leq j_{1} <\cdots  <j_{n -1} <r}\vert \varphi _{j_{1}} \cdots  \varphi _{j_{n -1}} \rangle  \langle \varphi _{1} \varphi _{j_{1}} \cdots  \varphi _{j_{n -1}}\vert \text{.}\end{align}Then the GA\ wedge product between a creator and an annihilator gives 

\begin{equation}a_{1}^{ \dag } \wedge a_{1} =\frac{1}{2} \left [a_{1}^{ \dag } ,a_{1}\right ] =\frac{1}{2} \left (a_{1}^{ \dag } a_{1} -a_{1} a_{1}^{ \dag }\right ) =\frac{1}{2} \left (a_{1}^{ \dag } a_{1} -I_{\mathcal{F}} +a_{1}^{ \dag } a_{1}\right ) =a_{1}^{ \dag } a_{1} -\frac{1}{2} I_{\mathcal{F}}
\end{equation}$\lceil $ which is antisymmetric as
\begin{equation}a_{1} \wedge a_{1}^{ \dag } =\frac{1}{2} \left [a_{1} ,a_{1}^{ \dag }\right ] =\frac{1}{2} \left (a_{1} a_{1}^{ \dag } -a_{1}^{ \dag } a_{1}\right ) =\frac{1}{2} \left (I_{\mathcal{F}} -a_{1}^{ \dag } a_{1} -a_{1}^{ \dag } a_{1}\right ) =\frac{1}{2} I_{\mathcal{F}} -a_{1}^{ \dag } a_{1}
\end{equation}$\rfloor $ 

while the operator product $ \boxtimes $ gives
\begin{align} & \left (\vert \varphi _{1} \rangle  \langle \phi \vert  +\sum _{2 \leq n \leq r}\sum _{2 \leq j_{1} <\cdots  <j_{n -1} <r}\vert \varphi _{1} \varphi _{j_{1}} \cdots  \varphi _{j_{n -1}} \rangle  \langle \varphi _{j_{1}} \cdots  \varphi _{j_{n -1}}\vert \right ) \\
 &  \boxtimes \left (\vert \phi  \rangle  \langle \varphi _{1}\vert  +\sum _{2 \leq n \leq r}\sum _{2 \leq j_{1} <\cdots  <j_{n -1} <r}\vert \varphi _{j_{1}} \cdots  \varphi _{j_{n -1}} \rangle  \langle \varphi _{1} \varphi _{j_{1}} \cdots  \varphi _{j_{n -1}}\vert \right ) \\
 &  =\vert \varphi _{1} \rangle  \langle \varphi _{1}\vert  +\sum _{2 \leq n \leq r}\sum _{2 \leq j_{1} <\cdots  <j_{n -1} <r}\vert \varphi _{1} \varphi _{j_{1}} \cdots  \varphi _{j_{n -1}} \rangle   \langle \;\varphi _{j_{1}} \cdots  \varphi _{j_{n -1}}\vert  +\sum _{2 \leq n \leq r}\sum _{2 \leq j_{1} <\cdots  <j_{n -1} <r}\vert \varphi _{j_{1}} \cdots  \varphi _{j_{n -1}} \rangle  \langle \varphi _{1} \varphi _{j_{1}} \cdots  \varphi _{j_{n -1}}\vert \; \\
 &  +\sum _{2 \leq n \leq r}\sum _{2 \leq j_{1} <\cdots  <j_{n -1} <r}\vert \varphi _{1} \varphi _{j_{1}} \cdots  \varphi _{j_{n -1}} \rangle   \langle \varphi _{j_{1}} \cdots  \varphi _{j_{n -1}}\vert  +\sum _{2 \leq n \leq r}\sum _{2 \leq j_{1} <\cdots  <j_{n -1} <r}\vert \varphi _{j_{1}} \cdots  \varphi _{j_{n -1}} \rangle  \langle \varphi _{1} \varphi _{j_{1}} \cdots  \varphi _{j_{n -1}}\vert \text{,}\end{align}\emph{which are not equal i.e }
\begin{equation}a_{1}^{ \dag } \wedge a_{1} \neq a_{1}^{ \dag } \boxtimes a_{1}
\end{equation}or equivalently
\begin{equation}\Pi _{\vert \phi  \rangle  \langle \phi \vert } \left (a_{1}^{ \dag } \wedge a_{1}\right ) \neq \Pi _{\vert \phi  \rangle  \langle \phi \vert } \left (a_{1}^{ \dag }\right ) \boxtimes \Pi _{\vert \phi  \rangle  \langle \phi \vert } \left (a_{1}\right )\text{.}
\end{equation}Incidently
\begin{equation}\Pi _{\vert \phi  \rangle  \langle \phi \vert } \left (a_{1}^{ \dag } \wedge a_{1}\right ) =\Pi _{\vert \phi  \rangle  \langle \phi \vert } \left (a_{1}^{ \dag }\right ) \wedge \Pi _{\vert \phi  \rangle  \langle \phi \vert } \left (a_{1}\right )
\end{equation}as $ \wedge $ is defined in terms of the Clifford$ \equiv $Operator product, while
\begin{equation}a_{1}^{ \dag } \boxtimes a_{1}
\end{equation}only makes sense in terms of representations i.e.
\begin{equation}\Pi _{\vert \phi  \rangle  \langle \phi \vert } \left (a_{1}^{ \dag }\right ) \boxtimes \Pi _{\vert \phi  \rangle  \langle \phi \vert } \left (a_{1}\right )
\end{equation}as the $ \boxtimes $ product is only defined for representations. 

The GA\ wedge product
between two creators gives
\begin{equation}a_{1}^{ \dag } \wedge a_{2}^{ \dag } =a_{1}^{ \dag } a_{2}^{ \dag } =\vert \varphi _{1} \varphi _{2} \rangle  \langle \phi \vert  +\sum _{3 \leq n \leq r}\sum _{3 \leq j_{1} <\cdots  <j_{n -2} <r}\vert \varphi _{1} \varphi _{2} \varphi _{j_{1}} \cdots  \varphi _{j_{n -2}} \rangle  \langle \varphi _{j_{1}} \cdots  \varphi _{j_{n -2}}\vert \text{,}
\end{equation}and we need to see if $a_{1}^{ \dag } \boxtimes a_{2}^{ \dag } =a_{1}^{ \dag } a_{2}^{ \dag }\text{.}$ In first quantization
\begin{align}a_{1}^{ \dag } \boxtimes a_{2}^{ \dag } &  =\left (\vert \varphi _{1} \rangle  \langle \phi \vert  +\sum _{2 \leq n \leq r}\sum _{2 \leq j_{1} <\cdots  <j_{n -2} <r}\vert \varphi _{1} \varphi _{j_{1}} \cdots  \varphi _{j_{n -1}} \rangle  \langle \varphi _{j_{1}} \cdots  \varphi _{j_{n -1}}\vert \right ) \\
 &  \boxtimes \left (\vert \varphi _{2} \rangle  \langle \phi \vert  +\sum _{2 \leq n \leq r}\sum _{1 \leq j_{1} <\cdots  <j_{n -2} \neq 2 <r}\vert \varphi _{2} \varphi _{j_{1}} \cdots  \varphi _{j_{n -1}} \rangle  \langle \varphi _{j_{1}} \cdots  \varphi _{j_{n -1}}\vert \right ) \\
 &  =\vert \varphi _{1} \rangle  \langle \phi \vert  \boxtimes \vert \varphi _{2} \rangle   \langle \phi \vert  +\vert \varphi _{1} \rangle   \langle \phi \vert  \boxtimes \left (\sum _{2 \leq n \leq r}\sum _{1 \leq j_{1} <\cdots  <j_{n -2} \neq 2 <r}\vert \varphi _{2} \varphi _{j_{1}} \cdots  \varphi _{j_{n -1}} \rangle  \langle \varphi _{j_{1}} \cdots  \varphi _{j_{n -1}}\vert \right ) \\
 &  +\left (\sum _{2 \leq n \leq r}\sum _{2 \leq j_{1} <\cdots  <j_{n -2} <r}\vert \varphi _{1} \varphi _{j_{1}} \cdots  \varphi _{j_{n -1}} \rangle  \langle \varphi _{j_{1}} \cdots  \varphi _{j_{n -1}}\vert \right ) \boxtimes \vert \varphi _{2} \rangle  \langle \phi \vert  \\
 &  +\left (\sum _{2 \leq n \leq r}\sum _{2 \leq j_{1} <\cdots  <j_{n -2} <r}\vert \varphi _{1} \varphi _{j_{1}} \cdots  \varphi _{j_{n -1}} \rangle  \langle \varphi _{j_{1}} \cdots  \varphi _{j_{n -1}}\vert \right ) \boxtimes \left (\sum _{2 \leq n \leq r}\sum _{1 \leq j_{1} <\cdots  <j_{n -2} \neq 2 <r}\vert \varphi _{2} \varphi _{j_{1}} \cdots  \varphi _{j_{n -1}} \rangle  \langle \varphi _{j_{1}} \cdots  \varphi _{j_{n -1}}\vert \right )\end{align}while
\begin{align}a_{1}^{ \dag } a_{2}^{ \dag } &  =\left (\vert \varphi _{1} \rangle  \langle \phi \vert  +\sum _{2 \leq n \leq r}\sum _{2 \leq j_{1} <\cdots  <j_{n -2} <r}\vert \varphi _{1} \varphi _{j_{1}} \cdots  \varphi _{j_{n -1}} \rangle  \langle \varphi _{j_{1}} \cdots  \varphi _{j_{n -1}}\vert \right ) \times  \\
 & \left (\vert \varphi _{2} \rangle  \langle \phi \vert  +\sum _{2 \leq n \leq r}\sum _{1 \leq j_{1} <\cdots  <j_{n -2} \neq 2 <r}\vert \varphi _{2} \varphi _{j_{1}} \cdots  \varphi _{j_{n -1}} \rangle  \langle \varphi _{j_{1}} \cdots  \varphi _{j_{n -1}}\vert \right ) \\
 &  =\left (\sum _{2 \leq n \leq r}\sum _{2 \leq j_{1} <\cdots  <j_{n -2} <r}\vert \varphi _{1} \varphi _{j_{1}} \cdots  \varphi _{j_{n -1}} \rangle  \langle \varphi _{j_{1}} \cdots  \varphi _{j_{n -1}}\vert \right )\vert \varphi _{2} \rangle  \langle \phi \vert  \\
 &  +\left (\sum _{2 \leq n \leq r}\sum _{2 \leq j_{1} <\cdots  <j_{n -2} <r}\vert \varphi _{1} \varphi _{j_{1}} \cdots  \varphi _{j_{n -1}} \rangle  \langle \varphi _{j_{1}} \cdots  \varphi _{j_{n -1}}\vert \right ) \left (\sum _{2 \leq n \leq r}\sum _{1 \leq j_{1} <\cdots  <j_{n -2} \neq 2 <r}\vert \varphi _{2} \varphi _{j_{1}} \cdots  \varphi _{j_{n -1}} \rangle  \langle \varphi _{j_{1}} \cdots  \varphi _{j_{n -1}}\vert \right )\text{.}\end{align}\emph{so indeed} $a_{1}^{ \dag } \boxtimes a_{2}^{ \dag } =a_{1}^{ \dag } a_{2}^{ \dag }\text{.}$ 

The Clifford wedge product between vectors, is defined as
\begin{equation}v_{1} \wedge \cdots  \wedge v_{n} =\frac{1}{n !} \sum _{\sigma  \in S_{n}}v_{\sigma  \left (1\right )} \cdots  v_{\sigma  \left (n\right )}
\end{equation}$\lceil $
\begin{align}v_{1} \wedge v_{2} &  =\frac{1}{2} \left (v_{1} v_{2} -v_{2} v_{1}\right ) =v_{1} v_{2} -\frac{1}{2} \left (v_{1} v_{2} +v_{2} v_{1}\right ) =v_{1} v_{2} -\left (\left .v_{1}\right \vert v_{2}\right ) \\
 &  =v_{1} v_{2} - \langle v_{1} v_{2} \rangle _{0} = \langle v_{1} v_{2} \rangle _{2} - \langle v_{1} v_{2} \rangle _{0} \\
 &  = \langle  \langle v_{1} \rangle _{1}  \langle v_{2} \rangle _{1} \rangle _{2} - \langle  \langle v_{1} \rangle _{1}  \langle v_{2} \rangle _{1} \rangle _{0} = \langle  \langle v_{1} \rangle _{1}  \langle v_{2} \rangle _{1} \rangle _{1 +1} - \langle  \langle v_{1} \rangle _{1}  \langle v_{2} \rangle _{1} \rangle _{1 -1} =\sum _{0 \leq p ,q \leq 2 r}\sum _{0 \leq j \leq p +q} \langle  \langle v_{1} \rangle _{p}  \langle v_{2} \rangle _{q} \rangle _{j}\end{align}
\begin{equation*}v_{1} \wedge v_{2} \wedge v_{3} =\frac{1}{6} \left (v_{1} v_{2} v_{3} -v_{2} v_{1} v_{3} +v_{3} v_{2} v_{1} -v_{2} v_{3} v_{1} -v_{1} v_{3} v_{2} +v_{3} v_{1} v_{2}\right )
\end{equation*}
\begin{equation*}v_{1} \wedge \left (v_{2} \wedge v_{3}\right ) =\frac{1}{6} \left (v_{1} v_{2} v_{3} +v_{1} v_{2} v_{3} -2 v_{2} v_{3} v_{1} +v_{1} v_{2} v_{3} +v_{1} v_{2} v_{3}\right )
\end{equation*}
\begin{equation}v_{1} \wedge v_{2} \wedge v_{3} =\left ( \langle v_{1} v_{2} \rangle _{2} - \langle v_{1} v_{2} \rangle _{0}\right ) \wedge v_{3} =
\end{equation}
\begin{equation}v_{1} v_{2} v_{3} =v_{1} \wedge v_{2} \wedge v_{3} + \langle v_{2} v_{3} \rangle _{0} v_{1} - \langle v_{1} v_{3} \rangle _{0} v_{2} + \langle v_{1} v_{2} \rangle _{0} v_{3}
\end{equation}$\rfloor $ 

and can be generalized to give 

\begin{equation}v \wedge B =\frac{1}{2} \left \{v B +\hat{B} v\right \}
\end{equation}so that
\begin{equation}v B =v\rfloor B +v \wedge B\text{.}
\end{equation}This can be extended to general elements to give
\begin{align}A\rfloor B =\sum _{0 \leq p ,q \leq 2 r} \langle A_{p} B_{q} \rangle _{p -q} \\
A\lfloor B =\sum _{0 \leq p ,q \leq 2 r} \langle A_{p} B_{q} \rangle _{q -p}\text{.}\end{align}This leads to
\begin{equation}A B =A\rfloor B +A \wedge B
\end{equation}giving
\begin{equation}A B =\sum _{0 \leq p ,q \leq 2 r} \langle A_{p} B_{q} \rangle _{p -q} +A \wedge B\text{,}
\end{equation}so that
\begin{align}A \wedge B &  =A B -\sum _{0 \leq p ,q \leq 2 r} \langle A_{p} B_{q} \rangle _{p -q} \\
 &  =\sum _{0 \leq p ,q \leq 2 r}A_{p} B_{q} -\sum _{0 \leq p ,q \leq 2 r} \langle A_{p} B_{q} \rangle _{p -q} \\
 &  =\sum _{0 \leq p ,q \leq 2 r}\sum _{j =0}^{\min \left (p ,q\right )} \langle A_{p} B_{q} \rangle _{\left \vert p -q\right \vert  +2 j} -\sum _{0 \leq p ,q \leq 2 r} \langle A_{p} B_{q} \rangle _{p -q} =\sum _{0 \leq p ,q \leq 2 r} \langle A_{p} B_{q} \rangle _{p +q}\text{?}\end{align}$\lceil $
\begin{align}v_{1} \wedge v_{2} &  =\sum _{0 \leq p ,q \leq 4}\sum _{j =0}^{\min \left (p ,q\right )} \langle v_{1 p} v_{2 q} \rangle _{\left \vert p -q\right \vert  +2 j} -\sum _{0 \leq p ,q \leq 4} \langle v_{1 p} v_{2 q} \rangle _{p -q} \\
 &  = \langle v_{1 1} v_{2 1} \rangle _{0} + \langle v_{1 1} v_{2 1} \rangle _{2} - \langle v_{1 1} v_{2 1} \rangle _{0}\end{align}
\begin{align}A_{2} \wedge B_{2} &  =\sum _{j =0}^{2} \langle A_{2} B_{2} \rangle _{\left \vert p -q\right \vert  +2 j} - \langle A_{2} B_{2} \rangle _{0} \\
 &  = \langle A_{2} B_{2} \rangle _{0} + \langle A_{2} B_{2} \rangle _{2} + \langle A_{2} B_{2} \rangle _{4} - \langle A_{2} B_{2} \rangle _{0} = \langle A_{2} B_{2} \rangle _{2} + \langle A_{2} B_{2} \rangle _{4}\text{?}\end{align}$\rfloor $ 

The wedge product can be fully generalized by
\begin{equation}A \wedge B =\sum _{0 \leq p ,q \leq 2 r} \langle A_{p} B_{q} \rangle _{q +p}\text{.}
\end{equation}

\subsection{Products of Basis Vectors}
The general geometric product is given by
\begin{equation}A B =\sum _{0 \leq m ,n \leq 2 r}A_{m} B_{n}\text{,}
\end{equation}where
\begin{equation}A_{m} B_{n} =\sum _{j =0}^{\min \left (m ,n\right )} \langle A_{m} B_{n} \rangle _{\left \vert m -n\right \vert  +2 j}\text{.}
\end{equation}

The product of the vectors $C_{1} ,D_{1}$ can be expressed as
\begin{equation}C_{1} D_{1} = \langle C_{1} D_{1} \rangle _{0} + \langle C_{1} D_{1} \rangle _{2} =(C_{1}\vert D_{1}) +C_{1} \wedge D_{1}
\end{equation}and in particular the product of two basis vectors as
\begin{align}x_{j_{1}} x_{j_{2}} &  =\left (x_{j_{1}}\left \vert x_{j_{2}}\right .\right ) +x_{j_{1}} \wedge x_{j_{2}} = \langle x_{j_{1}} x_{j_{2}} \rangle _{0} + \langle x_{j_{1}} x_{j_{2}} \rangle _{2} \\
 &  = \Delta _{j_{1} j_{2}} +x_{j_{1}} \wedge x_{j_{2}} =\ensuremath{\operatorname*{Pf}} \left ( \Delta _{j_{1} j_{2}}\right ) +x_{j_{1}} \wedge x_{j_{2}}\text{,}\end{align}which is equivalent to
\begin{equation}x_{j_{1}} \wedge x_{j_{2}} =x_{j_{1}} x_{j_{2}} -\left (x_{j_{1}}\left \vert x_{j_{2}}\right .\right )\text{.}
\end{equation}$\lceil $
\begin{equation}A B =\sum _{0 \leq m ,n \leq 2 r}\sum _{j =0}^{\min \left (m ,n\right )} \langle A_{m} B_{n} \rangle _{\left \vert m -n\right \vert  +2 j}
\end{equation}Let
\begin{align}A &  =A_{1} \\
B &  =B_{1} \\
A B &  =A_{1} B_{1} =\sum _{j =0}^{1} \langle A_{1} B_{1} \rangle _{2 j} = \langle A_{1} B_{1} \rangle _{0} + \langle A_{1} B_{1} \rangle _{2} \\
A_{1} \wedge B_{1} &  =A B - \langle A_{1} B_{1} \rangle _{0}\end{align}
\begin{equation}A \wedge B =\sum _{0 \leq p ,q \leq 1} \langle A_{p} B_{q} \rangle _{q +p} = \langle A_{0} B_{0} \rangle _{0} + \langle A_{1} B_{0} \rangle _{1} + \langle A_{0} B_{1} \rangle _{1} + \langle A_{1} B_{1} \rangle _{2}\text{.}
\end{equation}The exterior product is defined by
\begin{align}A_{p} \wedge B_{q} &  =A_{p} B_{q} -\sum _{j =0}^{\min \left (p ,q\right ) -1} \langle A_{p} B_{q} \rangle _{\left \vert p -q\right \vert  +2 j} \\
 &  =\sum _{j =0}^{\min \left (p ,q\right )} \langle A_{p} B_{q} \rangle _{\left \vert p -q\right \vert  +2 j} -\sum _{j =0}^{\min \left (p ,q\right ) -1} \langle A_{p} B_{q} \rangle _{\left \vert p -q\right \vert  +2 j} \\
 &  = \langle A_{p} B_{q} \rangle _{\left \vert p -q\right \vert  +2 \min \left (p ,q\right )} = \langle A_{p} B_{q} \rangle _{p +q}\end{align}which gives in this case
\begin{equation}A_{0} \wedge B_{0} =A_{0} B_{0} -\sum _{j =0}^{\min \left (0 ,0\right ) -1} \langle A_{0} B_{0} \rangle _{\left \vert 0\right \vert  +2 j} =A_{0} B_{0} - \langle A_{0} B_{0} \rangle _{ -2} =A_{0} B_{0}
\end{equation}
\begin{equation}A_{1} \wedge B_{0} =A_{1} B_{0} -\sum _{j =0}^{\min \left (1 ,0\right ) -1} \langle A_{1} B_{0} \rangle _{\left \vert 1 -0\right \vert  +2 j} =A_{1} B_{0} - \langle A_{1} B_{0} \rangle _{0} =A_{1} B_{0}
\end{equation}
\begin{equation}A_{1} \wedge B_{1} =A_{1} B_{1} -\sum _{j =0}^{\min \left (1 ,1\right ) -1} \langle A_{1} B_{1} \rangle _{\left \vert 1 -1\right \vert  +2 j} =A_{1} B_{1} - \langle A_{1} B_{1} \rangle _{0}
\end{equation}
\begin{align*}A_{1} B_{1} &  =\sum _{j =0}^{1} \langle A_{1} B_{1} \rangle _{\left \vert 1 -1\right \vert  +2 j} = \langle A_{1} B_{1} \rangle _{0} + \langle A_{1} B_{1} \rangle _{2} = \langle A_{1} B_{1} \rangle _{0} +A_{1} \wedge B_{1} \\
A_{1} \wedge B_{1} &  = \langle A_{1} B_{1} \rangle _{2}\end{align*}
\begin{align}A_{2} \wedge B_{1} &  =A_{2} B_{1} -\sum _{j =0}^{\min \left (2 ,1\right ) -1} \langle A_{2} B_{1} \rangle _{\left \vert 2 -1\right \vert  +2 j} =A_{2} B_{1} - \langle A_{2} B_{1} \rangle _{1 +0} =A_{2} B_{1} - \langle A_{2} B_{1} \rangle _{1} \\
A_{2} B_{1} &  =\sum _{j =0}^{1} \langle A_{2} B_{1} \rangle _{1 +2 j} = \langle A_{2} B_{1} \rangle _{1} + \langle A_{2} B_{1} \rangle _{3} \\
A_{2} \wedge B_{1} &  = \langle A_{2} B_{1} \rangle _{3}\end{align}
\begin{equation}A_{2} \wedge B_{2} =A_{2} B_{2} -\sum _{j =0}^{\min \left (2 ,2\right ) -1} \langle A_{2} B_{2} \rangle _{\left \vert 2 -2\right \vert  +2 j} =A_{2} B_{2} -\sum _{j =0}^{1} \langle A_{2} B_{2} \rangle _{2 j} =A_{2} B_{2} - \langle A_{2} B_{2} \rangle _{2} - \langle A_{2} B_{2} \rangle _{0}
\end{equation}
\begin{align}A_{2} B_{2} &  =\sum _{j =0}^{2} \langle A_{2} B_{2} \rangle _{\left \vert 2 -2\right \vert  +2 j} = \langle A_{2} B_{2} \rangle _{0} + \langle A_{2} B_{2} \rangle _{2} + \langle A_{2} B_{2} \rangle _{4} = \langle A_{2} B_{2} \rangle _{0} + \langle A_{2} B_{2} \rangle _{2} +A_{2} \wedge B_{2} \\
A_{2} \wedge B_{2} &  = \langle A_{2} B_{2} \rangle _{4}\end{align}
\begin{equation}A_{4} \wedge B_{2} =A_{4} B_{2} -\sum _{j =0}^{\min \left (4 ,2\right ) -1} \langle A_{4} B_{2} \rangle _{\left \vert 4 -2\right \vert  +2 j} =A_{4} B_{2} -\sum _{j =0}^{1} \langle A_{4} B_{2} \rangle _{2 +2 j} =A_{4} B_{2} - \langle A_{4} B_{2} \rangle _{2} - \langle A_{4} B_{2} \rangle _{4}
\end{equation}
\begin{align}A_{4} B_{2} &  =\sum _{j =0}^{2} \langle A_{4} B_{2} \rangle _{2 +2 j} = \langle A_{4} B_{2} \rangle _{2} + \langle A_{4} B_{2} \rangle _{4} + \langle A_{4} B_{2} \rangle _{6} \\
A_{4} \wedge B_{2} &  = \langle A_{4} B_{2} \rangle _{6}\end{align}
\begin{equation}A_{4} \wedge B_{4} =A_{4} B_{4} -\sum _{j =0}^{\min \left (4 ,4\right ) -1} \langle A_{4} B_{4} \rangle _{\left \vert 4 -4\right \vert  +2 j} =A_{4} B_{4} -\sum _{j =0}^{3} \langle A_{4} B_{4} \rangle _{2 j} =A_{4} B_{4} - \langle A_{4} B_{4} \rangle _{0} - \langle A_{4} B_{4} \rangle _{2} - \langle A_{4} B_{4} \rangle _{6}
\end{equation}
\begin{align}A_{4} B_{4} &  =\sum _{j =0}^{\min \left (4 ,4\right )} \langle A_{4} B_{4} \rangle _{0 +2 j} =\sum _{j =0}^{4} \langle A_{4} B_{4} \rangle _{2 j} = \langle A_{4} B_{4} \rangle _{0} + \langle A_{4} B_{4} \rangle _{2} + \langle A_{4} B_{4} \rangle _{6} + \langle A_{4} B_{4} \rangle _{8} \\
A_{4} \wedge B_{4} &  = \langle A_{4} B_{4} \rangle _{8}\end{align}
\begin{align}A &  =A_{0} +A_{2} +A_{4} \\
B &  =B_{0} +B_{2} +B_{4} \\
A B &  =\sum _{0 \leq m ,n \leq 4}\sum _{j =0}^{\min \left (m ,n\right )} \langle A_{m} B_{n} \rangle _{\left \vert m -n\right \vert  +2 j} \\
 &  = \langle A_{0} B_{0} \rangle _{0} + \langle A_{2} B_{0} \rangle _{2} + \langle A_{0} B_{2} \rangle _{2} + \langle A_{4} B_{0} \rangle _{4} \\
 &  + \langle A_{0} B_{4} \rangle _{4} + \langle A_{2} B_{2} \rangle _{0} + \langle A_{2} B_{2} \rangle _{2} + \langle A_{2} B_{2} \rangle _{4} \\
 &  + \langle A_{4} B_{2} \rangle _{2} + \langle A_{4} B_{2} \rangle _{4} + \langle A_{4} B_{2} \rangle _{6} + \langle A_{2} B_{4} \rangle _{2} + \langle A_{2} B_{4} \rangle _{4} + \langle A_{2} B_{4} \rangle _{6} \\
 &  + \langle A_{4} B_{4} \rangle _{0} + \langle A_{4} B_{4} \rangle _{2} + \langle A_{4} B_{4} \rangle _{4} + \langle A_{4} B_{4} \rangle _{6} + \langle A_{4} B_{4} \rangle _{8}\end{align}
\begin{align}A \wedge B &  =A B -\sum _{0 \leq m ,n \leq 2}\sum _{j =0}^{\min \left (m ,n ,1\right ) -1} \langle A_{m} B_{n} \rangle _{\left \vert m -n\right \vert  +2 j} \\
 &  = \langle A_{0} B_{0} \rangle _{0} + \langle A_{0} B_{2} \rangle _{2} + \langle A_{2} B_{0} \rangle _{2} + \langle A_{2} B_{2} \rangle _{4} + \langle A_{2} B_{4} \rangle _{6} + \langle A_{4} B_{2} \rangle _{6} + \langle A_{4} B_{4} \rangle _{8}\end{align}
\begin{equation}A \wedge B =\sum _{0 \leq p ,q \leq 4} \langle A_{p} B_{q} \rangle _{q +p}\text{.}
\end{equation}$\rfloor $ 

$\lceil $
\begin{align}x_{j_{1}} \wedge x_{j_{2}} &  =\sum _{0 \leq p ,q \leq 2 r} \langle  \langle x_{j_{1}} \rangle _{q}  \langle x_{j_{2}} \rangle _{p} \rangle _{q +p} = \langle  \langle x_{j_{1}} \rangle _{0}  \langle x_{j_{2}} \rangle _{0} \rangle _{0} + \langle  \langle x_{j_{1}} \rangle _{1}  \langle x_{j_{2}} \rangle _{0} \rangle _{1} + \langle  \langle x_{j_{1}} \rangle _{0}  \langle x_{j_{2}} \rangle _{1} \rangle _{1} + \langle  \langle x_{j_{1}} \rangle _{1}  \langle x_{j_{2}} \rangle _{1} \rangle _{2} \\
 &  = \langle  \langle x_{j_{1}} \rangle _{1}  \langle x_{j_{2}} \rangle _{1} \rangle _{2} = \langle x_{j_{1}} x_{j_{2}} \rangle _{2} =\frac{1}{2} \left [x_{j_{1}} ,x_{j_{2}}\right ]\text{.}\end{align}which is correct. 

$\rfloor $ 

\subparagraph{Pfaffian}
The Pfaffian, is a multilinear, complex valued function
\begin{equation}\ensuremath{\operatorname*{Pf}} :\overset{\left .n \left (n -1\right )\right /2}{\overbrace{\mathbb{C} \times \cdots  \times \mathbb{C}}}\, \cong \mathbb{C}^{\left .n \left (n -1\right )\right /2} \rightarrow \mathbb{C}
\end{equation}defined in general by
\begin{align}\ensuremath{\operatorname*{Pf}} \left (\mathbf{\Delta }_{j_{1} \cdots  j_{2 n}}\right ) &  =\ensuremath{\operatorname*{Pf}} \left (\mathbf{\Delta }_{j_{1} j_{2}} ,\cdots  ,\mathbf{\Delta }_{j_{1} j_{n}} ,\mathbf{\Delta }_{j_{2} j_{3}} ,\cdots  ,\mathbf{\Delta }_{j_{2} j_{n}} ,\cdots  ,\mathbf{\Delta }_{j_{2 n -1} j_{2 n}}\right ) \\
 &  ={\displaystyle\sum _{\rho  \in C_{2}^{2 n} \left (\mathbb{N}_{1}^{2 n}\right )}}\pi  \left (\rho \right ) {\displaystyle\prod _{1 \leq  \leq k \leq n}} \Delta _{j_{\rho  \left (2 k -1\right )} j_{\rho  \left (2 k\right )}} = \frac{1}{2^{n} n !} {\displaystyle\sum _{\sigma  \in S_{2 n}}}\pi  \left (\sigma \right ) {\displaystyle\prod _{1 \leq k \leq n}} \Delta _{j_{\sigma  \left (2 k -1\right )} j_{\sigma  \left (2 k\right )}}\text{,} \\
\mathbf{\Delta }_{j_{1} \cdots  j_{2 n}} &  \equiv \left (\mathbf{\Delta }_{j_{1} j_{2}} ,\cdots  ,\mathbf{\Delta }_{j_{1} j_{n}} ,\mathbf{\Delta }_{j_{2} j_{3}} ,\cdots  ,\mathbf{\Delta }_{j_{2} j_{n}} ,\cdots  ,\mathbf{\Delta }_{j_{2 n -1} j_{2 n}}\right )\text{,}\end{align}where $\rho $ and $\sigma $ are permutations
\begin{equation}\sigma  ,\rho  :\left \{1 ,\cdots  ,2 n\right \} \longrightarrow \left \{\begin{array}{c}\left \{\rho  \left (1\right ) ,\cdots  ,\rho  \left (2 n\right )\right \} \\
\left \{\sigma  \left (1\right ) ,\cdots  ,\sigma  \left (2 n\right )\right \}\end{array}\right .\text{,}
\end{equation}constrained by
\begin{align}\rho  \left (2 k -1\right ) &  <\rho  \left (2 k\right ) ;1 \leq k \leq n\text{,} \\
\sigma  \left (2 k -1\right ) &  <\sigma  \left (2 k\right ) ;1 \leq k \leq n\text{,}\end{align}leading to the induced permutation
\begin{equation}\sigma  ,\rho  :\left \{j_{1} ,\cdots  ,j_{2 n}\right \} \longrightarrow \left \{\begin{array}{c}\left \{j_{\rho  \left (1\right )} ,\cdots  ,j_{\rho  \left (2 n\right )}\right \} \\
\left \{j_{\sigma  \left (1\right )} ,\cdots  ,j_{\sigma  \left (2 n\right )}\right \}\end{array}\right .
\end{equation}and $\pi  ( \cdot )$ is its parity. \emph{Note that }$\left \{j_{1} ,\cdots  ,j_{2 n}\right \}$ \emph{do not have to have distinct values i.e. it is possible, for example, that }$j_{1} =j_{2} =\cdots  =j_{2 n}$. 

The Pfaffian can also be defined by as
\begin{equation}\frac{1}{n !} g^{n \wedge } =\,\overset{n}{\overbrace{g \wedge \cdots  \wedge g}}\, =P f \left (\mathbf{g}\right ) \left (x_{j_{1}} \wedge \cdots  \wedge x_{j_{2 n}}\right )\text{,}
\end{equation}when $\left \{j_{1} ,\cdots  ,j_{2 n}\right \}$ are all unequal, where
\begin{equation}g ={\displaystyle\sum _{1 \leq k \leq n}}g_{j_{2 k -1 ,2 k}} x_{j_{2 k -1}} \wedge x_{j_{2 k}}\text{.}
\end{equation}$\lceil \text{?}$This can also be extended to the case of equal subscripts by
\begin{equation}\frac{1}{n !} g^{n \wedge } =\,\overset{n}{\overbrace{g \wedge \cdots  \wedge g}}\, =P f \left (\mathbf{g}\right ) \left (x_{1} \wedge \cdots  \wedge x_{2 n}\right )\text{,}
\end{equation}where
\begin{equation}g ={\displaystyle\sum _{1 \leq k \leq n}}g_{2 k -1 ,2 k} x_{2 k -1} \wedge x_{2 k}
\end{equation}and setting
\begin{equation}g_{2 k -1 ,2 k} =\mathbf{\Delta }_{j_{2 k -1} j_{2 k}} ;1 \leq k \leq n\text{.}
\end{equation}

$\rfloor $ 

$\lceil $All odd order Pfaffians are zero. The second order pfaffian is
\begin{equation}\ensuremath{\operatorname*{Pf}} \left (\mathbf{\Delta }_{j_{1} j_{2}}\right ) =\mathbf{\Delta }_{j_{1} j_{2}}\text{,}
\end{equation}while the fourth order pfaffian is
\begin{equation}\ensuremath{\operatorname*{Pf}} \left (\mathbf{\Delta }_{j_{1} j_{2} j_{3} j_{4}}\right ) =\ensuremath{\operatorname*{Pf}} \left (\mathbf{\Delta }_{j_{1} j_{2}} ,\mathbf{\Delta }_{j_{1} j_{3}} ,\mathbf{\Delta }_{j_{1} j_{4}} ,\mathbf{\Delta }_{j_{2} j_{3}} ,\mathbf{\Delta }_{j_{2} j_{4}} ,\mathbf{\Delta }_{j_{3} j_{4}}\right ) = \Delta _{j_{1} j_{2}} \Delta _{j_{3} j_{4}} -  \Delta _{j_{1} j_{3}} \Delta _{j_{2} j_{4}} +  \Delta _{j_{1} j_{4}} \Delta _{j_{2} j_{3}}\text{.}
\end{equation}

The pfaffian $\ensuremath{\operatorname*{Pf}} \left (\mathbf{\Delta }_{1 2 3 4}\right )$ of the $4 \times 4$ antisymmetric matrix
\begin{equation}\left (\begin{array}{cccc}0 & \mathbf{\Delta }_{1 2} & \mathbf{\Delta }_{1 3} & \mathbf{\Delta }_{1 4} \\
 -\mathbf{\Delta }_{1 2} & 0 & \mathbf{\Delta }_{2 3} & \mathbf{\Delta }_{2 4} \\
 -\mathbf{\Delta }_{1 3} &  -\mathbf{\Delta }_{2 3} & 0 & \mathbf{\Delta }_{3 4} \\
 -\mathbf{\Delta }_{1 4} &  -\mathbf{\Delta }_{2 4} &  -\mathbf{\Delta }_{3 4} & 0\end{array}\right )
\end{equation}is
\begin{equation}\ensuremath{\operatorname*{Pf}} \left (\mathbf{\Delta }_{1 2 3 4}\right ) =\ensuremath{\operatorname*{Pf}} \left ( \Delta _{1 2} , \Delta _{1 3} , \Delta _{1 4} , \Delta _{2 3} , \Delta _{2 4} , \Delta _{3 4}\right ) = \Delta _{1 2} \Delta _{3 4} -  \Delta _{1 3} \Delta _{2 4} +  \Delta _{1 4} \Delta _{2 3}\text{,}
\end{equation}which is the same as the pfaffian $\ensuremath{\operatorname*{Pf}} \left (\mathbf{\Delta }_{1 2 3 4}\right )$ of the general $4 \times 4$ matrix
\begin{equation}\mathbf{\Delta }_{4} =\left (\begin{array}{cccc}\mathbf{\Delta }_{1 1} & \mathbf{\Delta }_{1 2} & \mathbf{\Delta }_{1 3} & \mathbf{\Delta }_{1 4} \\
\mathbf{\Delta }_{2 1} & \mathbf{\Delta }_{2 2} & \mathbf{\Delta }_{2 3} & \mathbf{\Delta }_{2 4} \\
\mathbf{\Delta }_{3 1} & \mathbf{\Delta }_{3 2} & \mathbf{\Delta }_{3 3} & \mathbf{\Delta }_{3 4} \\
\mathbf{\Delta }_{4 1} & \mathbf{\Delta }_{4 2} & \mathbf{\Delta }_{4 3} & \mathbf{\Delta }_{4 4}\end{array}\right )\text{,}
\end{equation}when we only consider the top triangle. 

The pfaffian $\ensuremath{\operatorname*{Pf}} \left (\mathbf{\Delta }_{1 1 2 2}\right )$ based on the matrix $\mathbf{\Delta }_{4}$ is
\begin{equation}\ensuremath{\operatorname*{Pf}} \left (\mathbf{\Delta }_{1 1 2 2}\right ) =\ensuremath{\operatorname*{Pf}} \left ( \Delta _{1 1} , \Delta _{1 2} , \Delta _{1 2} , \Delta _{1 2} , \Delta _{1 2} , \Delta _{2 2}\right ) = \Delta _{1 1} \Delta _{2 2} -  \Delta _{1 2} \Delta _{1 2} +  \Delta _{1 2} \Delta _{1 2} =  \Delta _{1 1} \Delta _{2 2}
\end{equation}and is the pfaffian of
\begin{equation}\left (\begin{array}{cccc}0 & \mathbf{\Delta }_{1 1} & \mathbf{\Delta }_{1 2} & \mathbf{\Delta }_{1 2} \\
 -\mathbf{\Delta }_{1 1} & 0 & \mathbf{\Delta }_{1 2} & \mathbf{\Delta }_{1 2} \\
 -\mathbf{\Delta }_{1 3} &  -\mathbf{\Delta }_{1 2} & 0 & \mathbf{\Delta }_{2 2} \\
 -\mathbf{\Delta }_{1 2} &  -\mathbf{\Delta }_{1 2} &  -\mathbf{\Delta }_{2 2} & 0\end{array}\right )\text{.}
\end{equation}

The sixth order pfaffian can be expanded in terms of sub Pfaffians as
\begin{align} & \ensuremath{\operatorname*{Pf}} \left (\mathbf{\Delta }_{j_{1} j_{2} j_{3} j_{4} j_{5} j_{6}}\right ) =\ensuremath{\operatorname*{Pf}} \left (\mathbf{\Delta }_{j_{1} j_{2} j_{3} j_{4}}\right )  \Delta _{j_{5} j_{6}} -\ensuremath{\operatorname*{Pf}} \left (\mathbf{\Delta }_{j_{1} j_{2} j_{3} j_{5}}\right ) \Delta _{j_{4} j_{6}} \nonumber  \\
 &  +\ensuremath{\operatorname*{Pf}} \left (\mathbf{\Delta }_{j_{1} j_{2} j_{4} j_{5}}\right )  \Delta _{j_{3} j_{6}} -\ensuremath{\operatorname*{Pf}} \left (\mathbf{\Delta }_{j_{1} j_{3} j_{4} j_{5}}\right )  \Delta _{j_{2} j_{6}} +\ensuremath{\operatorname*{Pf}} \left (\mathbf{\Delta }_{j_{2} j_{3} j_{4} j_{5}}\right ) \Delta _{j_{1} j_{6}}\text{.}\end{align}$\rfloor $ 

\subparagraph{Continue main theme}
The third order products are given by (noting that
\begin{equation} \langle  \langle x_{j_{1}} x_{j_{2}} \rangle _{2} x_{j_{3}} \rangle _{1} = \langle x_{j_{1}} x_{j_{2}} \rangle _{2}\rfloor x_{j_{3}} = \langle x_{j_{2}} x_{j_{3}} \rangle _{0} x_{j_{1}} - \langle x_{j_{1}} x_{j_{3}} \rangle _{0} x_{j_{2}}
\end{equation})
\begin{align}x_{j_{1}} x_{j_{2}} x_{j_{3}} &  =\left (x_{j_{1}} x_{j_{2}}\right ) x_{j_{3}} =\left (x_{j_{1}} x_{j_{2}}\right )\rfloor x_{j_{3}} +\left (x_{j_{1}} x_{j_{2}}\right ) \wedge x_{j_{3}} \nonumber  \\
 &  = \langle  \langle x_{j_{1}} x_{j_{2}} \rangle _{2} x_{j_{3}} \rangle _{1} +\left ( \langle x_{j_{1}} x_{j_{2}} \rangle _{0} + \langle x_{j_{1}} x_{j_{2}} \rangle _{2}\right ) \wedge x_{j_{3}} \nonumber  \\
 &  = \langle x_{j_{2}} x_{j_{3}} \rangle _{0} x_{j_{1}} - \langle x_{j_{1}} x_{j_{3}} \rangle _{0} x_{j_{2}} + \langle x_{j_{1}} x_{j_{2}} \rangle _{0} x_{j_{3}} + \langle x_{j_{1}} x_{j_{2}} x_{j_{3}} \rangle _{3} \nonumber  \\
 &  = \langle x_{j_{1}} x_{j_{2}} x_{j_{3}} \rangle _{3} + \langle x_{j_{2}} x_{j_{3}} \rangle _{0} x_{j_{1}} - \langle x_{j_{1}} x_{j_{3}} \rangle _{0} x_{j_{2}} + \langle x_{j_{1}} x_{j_{2}} \rangle _{0} x_{j_{3}} \nonumber  \\
 &  =x_{j_{1}} \wedge x_{j_{2}} \wedge x_{j_{3}} +\left (\left .x_{j_{2}}\right \vert x_{j_{3}}\right ) x_{j_{1}} -\left (\left .x_{j_{1}}\right \vert x_{j_{3}}\right ) x_{j_{2}} +\left (\left .x_{j_{1}}\right \vert x_{j_{2}}\right ) x_{j_{3}}\text{.}\end{align}giving
\begin{equation}x_{j_{1}} x_{j_{2}} x_{j_{3}} =x_{j_{1}} \wedge x_{j_{2}} \wedge x_{j_{3}} +\ensuremath{\operatorname*{Pf}} \left ( \Delta _{j_{2} j_{3}}\right ) x_{j_{1}} -\ensuremath{\operatorname*{Pf}} \left ( \Delta _{j_{1} j_{3}}\right ) x_{j_{2}} +\ensuremath{\operatorname*{Pf}} \left ( \Delta _{j_{1} j_{2}}\right ) x_{j_{3}}
\end{equation}and
\begin{equation}x_{j_{1}} \wedge x_{j_{2}} \wedge x_{j_{3}} =x_{j_{1}} x_{j_{2}} x_{j_{3}} -\ensuremath{\operatorname*{Pf}} \left ( \Delta _{j_{2} j_{3}}\right ) x_{j_{1}} +\ensuremath{\operatorname*{Pf}} \left ( \Delta _{j_{1} j_{3}}\right ) x_{j_{2}} -\ensuremath{\operatorname*{Pf}} \left ( \Delta _{j_{1} j_{2}}\right ) x_{j_{3}}\text{.}
\end{equation}

The fourth order products are given by
\begin{align} & x_{j_{1}} x_{j_{2}} x_{j_{3}} x_{j_{4}} =\left (x_{j_{1}} x_{j_{2}} x_{j_{3}}\right ) x_{j_{4}} =\left (x_{j_{1}} x_{j_{2}} x_{j_{3}}\right )\rfloor x_{j_{4}} +\left (x_{j_{1}} x_{j_{2}} x_{j_{3}}\right ) \wedge x_{j_{4}} \nonumber  \\
 &  =\left ( \langle x_{j_{1}} x_{j_{2}} x_{j_{3}} \rangle _{3} + \langle x_{j_{2}} x_{j_{3}} \rangle _{0} x_{j_{1}} - \langle x_{j_{1}} x_{j_{3}} \rangle _{0} x_{j_{2}} + \langle x_{j_{1}} x_{j_{2}} \rangle _{0} x_{j_{3}}\right )\rfloor x_{j_{4}} \nonumber  \\
 &  +\left ( \langle x_{j_{1}} x_{j_{2}} x_{j_{3}} \rangle _{3} + \langle x_{j_{2}} x_{j_{3}} \rangle _{0} x_{j_{1}} - \langle x_{j_{1}} x_{j_{3}} \rangle _{0} x_{j_{2}} + \langle x_{j_{1}} x_{j_{2}} \rangle _{0} x_{j_{3}}\right ) \wedge x_{j_{4}} \\
 &  = \langle x_{j_{1}} x_{j_{2}} x_{j_{3}} \rangle _{3}\rfloor x_{j_{4}} + \langle x_{j_{2}} x_{j_{3}} \rangle _{0} x_{j_{1}}\rfloor x_{j_{4}} - \langle x_{j_{1}} x_{j_{3}} \rangle _{0} x_{j_{2}}\rfloor x_{j_{4}} + \langle x_{j_{1}} x_{j_{2}} \rangle _{0} x_{j_{3}}\rfloor x_{j_{4}} \nonumber  \\
 &  + \langle x_{j_{1}} x_{j_{2}} x_{j_{3}} \rangle _{3} \wedge x_{j_{4}} + \langle x_{j_{2}} x_{j_{3}} \rangle _{0} x_{j_{1}} \wedge x_{j_{4}} - \langle x_{j_{1}} x_{j_{3}} \rangle _{0} x_{j_{2}} \wedge x_{j_{4}} + \langle x_{j_{1}} x_{j_{2}} \rangle _{0} x_{j_{3}} \wedge x_{j_{4}} \nonumber  \\
 &  = \langle x_{j_{1}} x_{j_{2}} \rangle _{2} x_{j_{3}}\rfloor x_{j_{4}} - \langle x_{j_{1}} x_{j_{3}} \rangle _{2} x_{j_{2}}\rfloor x_{j_{4}} + \langle x_{j_{2}} x_{j_{3}} \rangle _{2} x_{j_{1}}\rfloor x_{j_{4}} + \langle x_{j_{2}} x_{j_{3}} \rangle _{0} x_{j_{1}}\rfloor x_{j_{4}} - \langle x_{j_{1}} x_{j_{3}} \rangle _{0} x_{j_{2}}\rfloor x_{j_{4}} + \langle x_{j_{1}} x_{j_{2}} \rangle _{0} x_{j_{3}}\rfloor x_{j_{4}} \nonumber  \\
 &  + \langle x_{j_{1}} x_{j_{2}} x_{j_{3}} x_{j_{4}} \rangle _{3} + \langle x_{j_{2}} x_{j_{3}} \rangle _{0} x_{j_{1}} \wedge x_{j_{4}} - \langle x_{j_{1}} x_{j_{3}} \rangle _{0} x_{j_{2}} \wedge x_{j_{4}} + \langle x_{j_{1}} x_{j_{2}} \rangle _{0} x_{j_{3}} \wedge x_{j_{4}} \nonumber  \\
 &  = \langle x_{j_{3}} x_{j_{4}} \rangle _{0} x_{j_{1}} \wedge x_{j_{2}} - \langle x_{j_{2}} x_{j_{4}} \rangle _{0} x_{j_{1}} \wedge x_{j_{3}} + \langle x_{j_{1}} x_{j_{4}} \rangle _{0} x_{j_{2}} \wedge x_{j_{3}} \nonumber  \\
 &  + \langle x_{j_{2}} x_{j_{3}} \rangle _{0}  \langle x_{j_{1}} x_{j_{4}} \rangle _{0} - \langle x_{j_{1}} x_{j_{3}} \rangle _{0}  \langle x_{j_{2}} x_{j_{4}} \rangle _{0} + \langle x_{j_{1}} x_{j_{2}} \rangle _{0}  \langle x_{j_{3}} x_{j_{4}} \rangle _{0} \\
 &  + \langle x_{j_{1}} x_{j_{2}} x_{j_{3}} x_{j_{4}} \rangle _{3} + \langle x_{j_{2}} x_{j_{3}} \rangle _{0} x_{j_{1}} \wedge x_{j_{4}} - \langle x_{j_{1}} x_{j_{3}} \rangle _{0} x_{j_{2}} \wedge x_{j_{4}} + \langle x_{j_{1}} x_{j_{2}} \rangle _{0} x_{j_{3}} \wedge x_{j_{4}} \nonumber  \\
 &  =x_{j_{1}} \wedge x_{j_{2}} \wedge x_{j_{3}} \wedge x_{j_{4}} \nonumber  \\
 &  +\left (\left .x_{j_{1}}\right \vert x_{j_{2}}\right ) x_{j_{3}} \wedge x_{j_{4}} -\left (\left .x_{j_{1}}\right \vert x_{j_{3}}\right ) x_{j_{2}} \wedge x_{j_{4}} +\left (\left .x_{j_{1}}\right \vert x_{j_{4}}\right ) x_{j_{2}} \wedge x_{j_{3}} \nonumber  \\
 &  +\left (\left .x_{j_{2}}\right \vert x_{j_{3}}\right ) x_{j_{1}} \wedge x_{j_{4}} -\left (\left .x_{j_{2}}\right \vert x_{j_{4}}\right ) x_{j_{1}} \wedge x_{j_{3}} +\left (\left .x_{j_{3}}\right \vert x_{j_{4}}\right ) x_{j_{1}} \wedge x_{j_{2}} \nonumber  \\
 &  +\left (\left .x_{j_{1}}\right \vert x_{j_{2}}\right ) \left (\left .x_{j_{3}}\right \vert x_{j_{4}}\right ) -\left (\left .x_{j_{1}}\right \vert x_{j_{3}}\right ) \left (\left .x_{j_{2}}\right \vert x_{j_{4}}\right ) +\left (\left .x_{j_{1}}\right \vert x_{j_{4}}\right ) \left (\left .x_{j_{2}}\right \vert x_{j_{3}}\right )\text{.}\end{align}Giving
\begin{align}x_{j_{1}} x_{j_{2}} x_{j_{3}} x_{j_{4}} &  =x_{j_{1}} \wedge x_{j_{2}} \wedge x_{j_{3}} \wedge x_{j_{4}} \nonumber  \\
 &  +\left (\left .x_{j_{1}}\right \vert x_{j_{2}}\right ) x_{j_{3}} \wedge x_{j_{4}} -\left (\left .x_{j_{1}}\right \vert x_{j_{3}}\right ) x_{j_{2}} \wedge x_{j_{4}} +\left (\left .x_{j_{1}}\right \vert x_{j_{4}}\right ) x_{j_{2}} \wedge x_{j_{3}} \nonumber  \\
 &  +\left (\left .x_{j_{2}}\right \vert x_{j_{3}}\right ) x_{j_{1}} \wedge x_{j_{4}} -\left (\left .x_{j_{2}}\right \vert x_{j_{4}}\right ) x_{j_{1}} \wedge x_{j_{3}} +\left (\left .x_{j_{3}}\right \vert x_{j_{4}}\right ) x_{j_{1}} \wedge x_{j_{2}} \nonumber  \\
 &  +\left (\left .x_{j_{1}}\right \vert x_{j_{2}}\right ) \left (\left .x_{j_{3}}\right \vert x_{j_{4}}\right ) -\left (\left .x_{j_{1}}\right \vert x_{j_{3}}\right ) \left (\left .x_{j_{2}}\right \vert x_{j_{4}}\right ) +\left (\left .x_{j_{1}}\right \vert x_{j_{4}}\right ) \left (\left .x_{j_{2}}\right \vert x_{j_{3}}\right )\end{align}and thus
\begin{align}x_{j_{1}} x_{j_{2}} x_{j_{3}} x_{j_{4}} &  =x_{j_{1}} \wedge x_{j_{2}} \wedge x_{j_{3}} \wedge x_{j_{4}} \nonumber  \\
 &  + \Delta _{j_{1} j_{2}}x_{j_{3}} \wedge x_{j_{4}} - \Delta _{j_{1} j_{3}}x_{j_{2}} \wedge x_{j_{4}} + \Delta _{j_{1} j_{4}}x_{j_{2}} \wedge x_{j_{3}} \nonumber  \\
 &  + \Delta _{j_{2} j_{2 3}}x_{j_{1}} \wedge x_{j_{4}} - \Delta _{j_{2} j_{4}}x_{j_{1}} \wedge x_{j_{3}} + \Delta _{j_{3} j_{4}}x_{j_{1}} \wedge x_{j_{2}} \nonumber  \\
 &  + \Delta _{j_{1} j_{2}} \Delta _{j_{3} j_{4}} -  \Delta _{j_{1} j_{3}} \Delta _{j_{2} j_{4}} +  \Delta _{j_{1} j_{4}} \Delta _{j_{2} j_{2 3}}\text{.}\end{align}which can also be expressed as
\begin{align}x_{j_{1}} x_{j_{2}} x_{j_{3}} x_{j_{4}} &  =x_{j_{1}} \wedge x_{j_{2}} \wedge x_{j_{3}} \wedge x_{j_{4}} \nonumber  \\
 &  +\ensuremath{\operatorname*{Pf}} \left ( \Delta _{j_{1} j_{2}}\right ) x_{j_{3}} \wedge x_{j_{4}} -\ensuremath{\operatorname*{Pf}} \left ( \Delta _{j_{1} j_{3}}\right ) x_{j_{2}} \wedge x_{j_{4}} +\ensuremath{\operatorname*{Pf}} \left ( \Delta _{j_{1} j_{4}}\right ) x_{j_{2}} \wedge x_{j_{3}} \nonumber  \\
 &  +\ensuremath{\operatorname*{Pf}} \left ( \Delta _{j_{2} j_{3}}\right ) x_{j_{1}} \wedge x_{j_{4}} -\ensuremath{\operatorname*{Pf}} \left ( \Delta _{j_{2} j_{4}}\right ) x_{j_{1}} \wedge x_{j_{3}} +\ensuremath{\operatorname*{Pf}} \left ( \Delta _{j_{3} j_{4}}\right ) x_{j_{1}} \wedge x_{j_{2}} \nonumber  \\
 &  +\ensuremath{\operatorname*{Pf}} \left ( \Delta _{j_{1} j_{2} j_{3} j_{4}}\right )\text{.}\end{align}$\lceil $To check we derive this again by
\begin{align} & x_{i} x_{j} x_{k} x_{l} =\left \{\left (x_{i}\left \vert x_{j}\right .\right ) +x_{i} \wedge x_{j}\right \}\left \{\left (x_{k}\left \vert x_{l}\right .\right ) +x_{k} \wedge x_{l}\right \} \nonumber  \\
 &  =\left (x_{i}\left \vert x_{j}\right .\right ) \left (x_{k}\left \vert x_{l}\right .\right ) +\left (x_{i}\left \vert x_{j}\right .\right ) x_{k} \wedge x_{l} +\left (x_{k}\left \vert x_{l}\right .\right ) x_{i} \wedge x_{j} +\left (x_{i} \wedge x_{j}\right ) \left (x_{k} \wedge x_{l}\right ) \nonumber  \\
 &  =\left (x_{i}\left \vert x_{j}\right .\right ) \left (x_{k}\left \vert x_{l}\right .\right ) +\left (x_{i}\left \vert x_{j}\right .\right ) x_{k} \wedge x_{l} +\left (x_{k}\left \vert x_{l}\right .\right ) x_{i} \wedge x_{j} \nonumber  \\
 &  -\left (\left .x_{i}\right \vert x_{k}\right ) x_{j} \wedge x_{l} +\left (\left .x_{i}\right \vert x_{l}\right ) x_{j} \wedge x_{k} +\left (\left .x_{j}\right \vert x_{k}\right ) x_{i} \wedge x_{l} -\left (\left .x_{j}\right \vert x_{l}\right ) x_{i} \wedge x_{k} \nonumber  \\
 &  +x_{i} \wedge x_{j} \wedge x_{k} \wedge x_{l}\end{align}$\rfloor $ 

The fifth order by
\begin{align}x_{i} x_{j} x_{k} x_{l} x_{m} &  =\left \{x_{i} \wedge x_{j} \wedge x_{k} \wedge x_{l} +\left (\left .x_{i}\right \vert x_{j}\right ) x_{k} \wedge x_{l} -\left (\left .x_{i}\right \vert x_{k}\right ) x_{j} \wedge x_{l} +\left (\left .x_{i}\right \vert x_{l}\right ) x_{j} \wedge x_{k}\right . \nonumber  \\
 &  +\left (\left .x_{j}\right \vert x_{k}\right ) x_{i} \wedge x_{l} -\left (\left .x_{j}\right \vert x_{l}\right ) x_{i} \wedge x_{k} +\left (\left .x_{k}\right \vert x_{l}\right ) x_{i} \wedge x_{j} \nonumber  \\
 &  +\left .\left (\left .x_{i}\right \vert x_{j}\right ) \left (\left .x_{k}\right \vert x_{l}\right ) -\left (\left .x_{i}\right \vert x_{k}\right ) \left (\left .x_{j}\right \vert x_{l}\right ) +\left (\left .x_{i}\right \vert x_{l}\right ) \left (\left .x_{j}\right \vert x_{k}\right )\right \} x_{m}\text{.}\end{align}$\lceil $
\begin{align}\left \{x_{i} \wedge x_{j} \wedge x_{k} \wedge x_{l}\right \}x_{m} &  =x_{i} \wedge x_{j} \wedge x_{k} \wedge x_{l} \wedge x_{m} +x_{i} \wedge x_{j} \wedge x_{k} \left (\left .x_{l}\right \vert x_{m}\right ) -x_{i} \wedge x_{j} \wedge x_{l} \left (\left .x_{k}\right \vert x_{m}\right ) \nonumber  \\
 &  +x_{i} \wedge x_{k} \wedge x_{l} \left (\left .x_{j}\right \vert x_{m}\right ) -x_{j} \wedge x_{k} \wedge x_{l} \left (\left .x_{i}\right \vert x_{m}\right ) \nonumber  \\
\left \{\left (\left .x_{i}\right \vert x_{j}\right ) x_{k} \wedge x_{l}\right \}x_{m} &  =\left (\left .x_{i}\right \vert x_{j}\right ) \left \{x_{k} \wedge x_{l} \wedge x_{m} +\left (\left .x_{l}\right \vert x_{m}\right ) x_{k} -\left (\left .x_{k}\right \vert x_{m}\right ) x_{l}\right \} \nonumber  \\
 &  =\left (\left .x_{i}\right \vert x_{j}\right ) x_{k} \wedge x_{l} \wedge x_{m} +\left (\left .x_{i}\right \vert x_{j}\right ) \left (\left .x_{l}\right \vert x_{m}\right ) x_{k} -\left (\left .x_{i}\right \vert x_{j}\right ) \left (\left .x_{k}\right \vert x_{m}\right ) x_{l} \nonumber  \\
\left (\left .x_{i}\right \vert x_{j}\right ) \left (\left .x_{k}\right \vert x_{l}\right ) x_{m} &  =\left (\left .x_{i}\right \vert x_{j}\right ) \left (\left .x_{k}\right \vert x_{l}\right ) x_{m}\end{align}$\rfloor $ giving
\begin{align} & x_{i} x_{j} x_{k} x_{l} x_{m} =x_{i} \wedge x_{j} \wedge x_{k} \wedge x_{l} \wedge x_{m} \nonumber  \\
 &  +\left (\left .x_{i}\right \vert x_{j}\right ) x_{k} \wedge x_{l} \wedge x_{m} -\left (\left .x_{i}\right \vert x_{k}\right ) x_{j} \wedge x_{l} \wedge x_{m} +\left (\left .x_{i}\right \vert x_{l}\right ) x_{j} \wedge x_{k} \wedge x_{m} \nonumber  \\
 &  +\left (\left .x_{j}\right \vert x_{k}\right ) x_{i} \wedge x_{l} \wedge x_{m} -\left (\left .x_{j}\right \vert x_{l}\right ) x_{i} \wedge x_{k} \wedge x_{m} +\left (\left .x_{k}\right \vert x_{l}\right ) x_{i} \wedge x_{j} \wedge x_{m} \nonumber  \\
 &  +\left (\left .x_{l}\right \vert x_{m}\right ) x_{i} \wedge x_{j} \wedge x_{k} -\left (\left .x_{k}\right \vert x_{m}\right ) x_{i} \wedge x_{j} \wedge x_{l} +\left (\left .x_{j}\right \vert x_{m}\right ) x_{i} \wedge x_{k} \wedge x_{l} -\left (\left .x_{i}\right \vert x_{m}\right ) x_{j} \wedge x_{k} \wedge x_{l} \nonumber  \\
 &  +\left (\left .x_{i}\right \vert x_{j}\right ) \left (\left .x_{l}\right \vert x_{m}\right ) x_{k} -\left (\left .x_{i}\right \vert x_{j}\right ) \left (\left .x_{k}\right \vert x_{m}\right ) x_{l} \nonumber  \\
 &  -\left (\left .x_{i}\right \vert x_{k}\right ) \left (\left .x_{l}\right \vert x_{m}\right ) x_{j} +\left (\left .x_{i}\right \vert x_{k}\right ) \left (\left .x_{j}\right \vert x_{m}\right ) x_{l} \nonumber  \\
 &  -\left (\left .x_{i}\right \vert x_{l}\right ) \left (\left .x_{j}\right \vert x_{m}\right ) x_{k} +\left (\left .x_{i}\right \vert x_{l}\right ) \left (\left .x_{k}\right \vert x_{m}\right ) x_{j} \nonumber  \\
 &  +\left (\left .x_{j}\right \vert x_{k}\right ) \left (\left .x_{l}\right \vert x_{m}\right ) x_{i} -\left (\left .x_{j}\right \vert x_{k}\right ) \left (\left .x_{i}\right \vert x_{m}\right ) x_{l} \nonumber  \\
 &  -\left (\left .x_{j}\right \vert x_{l}\right ) \left (\left .x_{k}\right \vert x_{m}\right ) x_{i} +\left (\left .x_{j}\right \vert x_{l}\right ) \left (\left .x_{k}\right \vert x_{m}\right ) x_{k} \nonumber  \\
 &  +\left (\left .x_{k}\right \vert x_{l}\right ) \left (\left .x_{j}\right \vert x_{m}\right ) x_{i} -\left (\left .x_{k}\right \vert x_{l}\right ) \left (\left .x_{i}\right \vert x_{m}\right ) x_{j} \nonumber  \\
 &  +\left \{\left (\left .x_{i}\right \vert x_{j}\right ) \left (\left .x_{k}\right \vert x_{l}\right ) -\left (\left .x_{i}\right \vert x_{k}\right ) \left (\left .x_{j}\right \vert x_{l}\right ) +\left (\left .x_{i}\right \vert x_{l}\right ) \left (\left .x_{j}\right \vert x_{k}\right )\right \} x_{m}\text{.}\end{align}then
\begin{align} & x_{i} x_{j} x_{k} x_{l} x_{m} =x_{i} \wedge x_{j} \wedge x_{k} \wedge x_{l} \wedge x_{m} \nonumber  \\
 &  +\left (\left .x_{i}\right \vert x_{j}\right ) x_{k} \wedge x_{l} \wedge x_{m} -\left (\left .x_{i}\right \vert x_{k}\right ) x_{j} \wedge x_{l} \wedge x_{m} +\left (\left .x_{i}\right \vert x_{l}\right ) x_{j} \wedge x_{k} \wedge x_{m} \nonumber  \\
 &  +\left (\left .x_{j}\right \vert x_{k}\right ) x_{i} \wedge x_{l} \wedge x_{m} -\left (\left .x_{j}\right \vert x_{l}\right ) x_{i} \wedge x_{k} \wedge x_{m} +\left (\left .x_{k}\right \vert x_{l}\right ) x_{i} \wedge x_{j} \wedge x_{m} \nonumber  \\
 &  +\left (\left .x_{l}\right \vert x_{m}\right ) x_{i} \wedge x_{j} \wedge x_{k} -\left (\left .x_{k}\right \vert x_{m}\right ) x_{i} \wedge x_{j} \wedge x_{l} +\left (\left .x_{j}\right \vert x_{m}\right ) x_{i} \wedge x_{k} \wedge x_{l} -\left (\left .x_{i}\right \vert x_{m}\right ) x_{j} \wedge x_{k} \wedge x_{l} \nonumber  \\
 &  +\left \{\left (\left .x_{i}\right \vert x_{j}\right ) \left (\left .x_{l}\right \vert x_{m}\right ) -\left (\left .x_{i}\right \vert x_{l}\right ) \left (\left .x_{j}\right \vert x_{m}\right ) +\left (\left .x_{j}\right \vert x_{l}\right ) \left (\left .x_{k}\right \vert x_{m}\right )\right \} x_{k} \nonumber  \\
 &  +\left \{\left (\left .x_{i}\right \vert x_{k}\right ) \left (\left .x_{j}\right \vert x_{m}\right ) -\left (\left .x_{j}\right \vert x_{k}\right ) \left (\left .x_{i}\right \vert x_{m}\right ) -\left (\left .x_{i}\right \vert x_{j}\right ) \left (\left .x_{k}\right \vert x_{m}\right )\right \} x_{l} \nonumber  \\
 &  +\left \{\left (\left .x_{i}\right \vert x_{l}\right ) \left (\left .x_{k}\right \vert x_{m}\right ) -\left (\left .x_{k}\right \vert x_{l}\right ) \left (\left .x_{i}\right \vert x_{m}\right ) -\left (\left .x_{i}\right \vert x_{k}\right ) \left (\left .x_{l}\right \vert x_{m}\right )\right \} x_{j} \nonumber  \\
 &  +\left \{\left (\left .x_{j}\right \vert x_{k}\right ) \left (\left .x_{l}\right \vert x_{m}\right ) -\left (\left .x_{j}\right \vert x_{l}\right ) \left (\left .x_{k}\right \vert x_{m}\right ) +\left (\left .x_{k}\right \vert x_{l}\right ) \left (\left .x_{j}\right \vert x_{m}\right )\right \} x_{i} \nonumber  \\
 &  +\left \{\left (\left .x_{i}\right \vert x_{j}\right ) \left (\left .x_{k}\right \vert x_{l}\right ) -\left (\left .x_{i}\right \vert x_{k}\right ) \left (\left .x_{j}\right \vert x_{l}\right ) +\left (\left .x_{i}\right \vert x_{l}\right ) \left (\left .x_{j}\right \vert x_{k}\right )\right \} x_{m}\text{.}\end{align}Thus
\begin{align} & x_{i} x_{j} x_{k} x_{l} x_{m} =x_{i} \wedge x_{j} \wedge x_{k} \wedge x_{l} \wedge x_{m} \nonumber  \\
 &  + \Delta _{i j}x_{k} \wedge x_{l} \wedge x_{m} - \Delta _{i k}x_{j} \wedge x_{l} \wedge x_{m} + \Delta _{i l}x_{j} \wedge x_{k} \wedge x_{m} + \Delta _{j k}x_{i} \wedge x_{l} \wedge x_{m} - \Delta _{j l}x_{i} \wedge x_{k} \wedge x_{m} \nonumber  \\
 &  + \Delta _{k l}x_{i} \wedge x_{j} \wedge x_{m} + \Delta _{l m}x_{i} \wedge x_{j} \wedge x_{k} - \Delta _{k m}x_{i} \wedge x_{j} \wedge x_{l} + \Delta _{j m}x_{i} \wedge x_{k} \wedge x_{l} - \Delta _{i m}x_{j} \wedge x_{k} \wedge x_{l} \nonumber  \\
 &  +\left \{ \Delta _{i j} \Delta _{l m} -  \Delta _{i l} \Delta _{j m} +  \Delta _{j l} \Delta _{k m}\right \} x_{k} \nonumber  \\
 &  +\left \{ \Delta _{i k} \Delta _{j m} -  \Delta _{j k} \Delta _{i m} -  \Delta _{i j} \Delta _{k m}\right \} x_{l} \nonumber  \\
 &  +\left \{ \Delta _{i l} \Delta _{k m} -  \Delta _{k l} \Delta _{i m} -  \Delta _{i k} \Delta _{l m}\right \} x_{j} \nonumber  \\
 &  +\left \{ \Delta _{l m} \Delta _{j k} -  \Delta _{j l} \Delta _{k m} +  \Delta _{k l} \Delta _{j m}\right \} x_{i} \nonumber  \\
 &  +\left \{ \Delta _{i j} \Delta _{k l} -  \Delta _{i k} \Delta _{j l} +  \Delta _{i l} \Delta _{j k}\right \} x_{m}\text{.}\end{align}Hence the fifth order products are
\begin{align} & x_{j_{1}} x_{j_{2}} x_{j_{3}} x_{j_{4}} x_{j_{5}} =x_{j_{1}} \wedge x_{j_{2}} \wedge x_{j_{3}} \wedge x_{j_{4}} \wedge x_{j_{5}} \nonumber  \\
 &  +\ensuremath{\operatorname*{Pf}} \left ( \Delta _{j_{1} j_{2}}\right ) x_{j_{3}} \wedge x_{j_{4}} \wedge x_{j_{5}} -\ensuremath{\operatorname*{Pf}} \left ( \Delta _{j_{1} j_{3}}\right ) x_{j_{2}} \wedge x_{j_{4}} \wedge x_{j_{5}} +\ensuremath{\operatorname*{Pf}} \left ( \Delta _{j_{1} j_{4}}\right ) x_{j_{2}} \wedge x_{j_{3}} \wedge x_{j_{5}} \nonumber  \\
 &  +\ensuremath{\operatorname*{Pf}} \left ( \Delta _{j_{2} j_{3}}\right ) x_{j_{1}} \wedge x_{j_{4}} \wedge x_{j_{5}} -\ensuremath{\operatorname*{Pf}} \left ( \Delta _{j_{2} j_{4}}\right ) x_{j_{1}} \wedge x_{j_{3}} \wedge x_{j_{5}} +\ensuremath{\operatorname*{Pf}} \left ( \Delta _{j_{3} j_{4}}\right ) x_{j_{1}} \wedge x_{j_{2}} \wedge x_{j_{5}} \nonumber  \\
 &  +\ensuremath{\operatorname*{Pf}} \left ( \Delta _{j_{4} j_{5}}\right ) x_{j_{1}} \wedge x_{j_{2}} \wedge x_{j_{3}} -\ensuremath{\operatorname*{Pf}} \left ( \Delta _{j_{3} j_{5}}\right ) x_{j_{1}} \wedge x_{j_{2}} \wedge x_{j_{4}} +\ensuremath{\operatorname*{Pf}} \left ( \Delta _{j_{2} j_{5}}\right ) x_{j_{1}} \wedge x_{j_{3}} \wedge x_{j_{4}} -\ensuremath{\operatorname*{Pf}} \left ( \Delta _{j_{1} j_{5}}\right ) x_{j_{2}} \wedge x_{j_{3}} \wedge x_{j_{4}} \nonumber  \\
 &  +\ensuremath{\operatorname*{Pf}} \left ( \Delta _{j_{2} j_{3} j_{4} j_{5}}\right ) x_{j_{1}} -\ensuremath{\operatorname*{Pf}} \left ( \Delta _{j_{1} j_{3} j_{4} j_{5}}\right ) x_{j_{2}} +\ensuremath{\operatorname*{Pf}} \left ( \Delta _{j_{1} j_{2} j_{4} j_{5}}\right ) x_{j_{3}} -\ensuremath{\operatorname*{Pf}} \left ( \Delta _{j_{1} j_{2} j_{3} j_{5}}\right ) x_{j_{4}} +\ensuremath{\operatorname*{Pf}} \left ( \Delta _{j_{1} j_{2} j_{3} j_{4}}\right ) x_{j_{5}}\text{,}\end{align}

The sixth order products decompose as
\begin{align} & x_{i} x_{j} x_{k} x_{l} x_{m} x_{n} =\left \{x_{i} \wedge x_{j} \wedge x_{k} \wedge x_{l} \wedge x_{m} \wedge x_{n}\right . \nonumber  \\
 &  + \Delta _{i j}x_{k} \wedge x_{l} \wedge x_{m} - \Delta _{i k}x_{j} \wedge x_{l} \wedge x_{m} + \Delta _{i l}x_{j} \wedge x_{k} \wedge x_{m} + \Delta _{j k}x_{i} \wedge x_{l} \wedge x_{m} - \Delta _{j l}x_{i} \wedge x_{k} \wedge x_{m} \nonumber  \\
 &  + \Delta _{k l}x_{i} \wedge x_{j} \wedge x_{m} + \Delta _{l m}x_{i} \wedge x_{j} \wedge x_{k} - \Delta _{k m}x_{i} \wedge x_{j} \wedge x_{l} + \Delta _{j m}x_{i} \wedge x_{k} \wedge x_{l} - \Delta _{i m}x_{j} \wedge x_{k} \wedge x_{l} \nonumber  \\
 &  +\left .\ensuremath{\operatorname*{Pf}} \left ( \Delta _{j k l m}\right ) x_{i} -\ensuremath{\operatorname*{Pf}} \left ( \Delta _{i k l m}\right ) x_{j} +\ensuremath{\operatorname*{Pf}} \left ( \Delta _{i j l m}\right ) x_{k} -\ensuremath{\operatorname*{Pf}} \left ( \Delta _{i j k m}\right ) x_{l} +\ensuremath{\operatorname*{Pf}} \left ( \Delta _{i j k l}\right ) x_{m}\right \} x_{n}\end{align}leading to
\begin{align} & x_{i} x_{j} x_{k} x_{l} x_{m} x_{n} =x_{i} \wedge x_{j} \wedge x_{k} \wedge x_{l} \wedge x_{m} \wedge x_{n} \nonumber  \\
 &  + \Delta _{m n}x_{i} \wedge x_{j} \wedge x_{k} \wedge x_{l} - \Delta _{\ln }x_{i} \wedge x_{j} \wedge x_{k} \wedge x_{m} + \Delta _{k n}x_{i} \wedge x_{j} \wedge x_{l} \wedge x_{m} \nonumber  \\
 &  - \Delta _{j n}x_{i} \wedge x_{k} \wedge x_{l} \wedge x_{m} + \Delta _{i n}x_{j} \wedge x_{k} \wedge x_{l} \wedge x_{m} \nonumber  \\
 &  + \Delta _{i j}x_{k} \wedge x_{l} \wedge x_{m} \wedge x_{n} + \Delta _{i j} \Delta _{m n}x_{k} \wedge x_{l} - \Delta _{i j} \Delta _{\ln }x_{k} \wedge x_{m} + \Delta _{i j} \Delta _{k n}x_{l} \wedge x_{m} \nonumber  \\
 &  - \Delta _{i k}x_{j} \wedge x_{l} \wedge x_{m} \wedge x_{n} + \Delta _{i k} \Delta _{m n}x_{j} \wedge x_{l} - \Delta _{i k} \Delta _{\ln }x_{j} \wedge x_{m} + \Delta _{i k} \Delta _{j n}x_{l} \wedge x_{m} \nonumber  \\
 &  + \Delta _{i l}x_{j} \wedge x_{k} \wedge x_{m} \wedge x_{n} - \Delta _{i l} \Delta _{m n}x_{j} \wedge x_{k} + \Delta _{i l} \Delta _{k n}x_{j} \wedge x_{m} - \Delta _{i l} \Delta _{j n}x_{k} \wedge x_{m} \nonumber  \\
 &  + \Delta _{j k}x_{i} \wedge x_{l} \wedge x_{m} \wedge x_{n} - \Delta _{j k} \Delta _{m n}x_{i} \wedge x_{l} + \Delta _{j k} \Delta _{\ln }x_{i} \wedge x_{m} - \Delta _{j k} \Delta _{i n}x_{l} \wedge x_{m} \nonumber  \\
 &  - \Delta _{j l}x_{i} \wedge x_{k} \wedge x_{m} \wedge x_{n} + \Delta _{j l} \Delta _{m n}x_{i} \wedge x_{k} - \Delta _{j l} \Delta _{k n}x_{i} \wedge x_{m} + \Delta _{j l} \Delta _{i n}x_{k} \wedge x_{m} \nonumber  \\
 &  + \Delta _{k l}x_{i} \wedge x_{j} \wedge x_{m} \wedge x_{n} - \Delta _{k l} \Delta _{m n}x_{i} \wedge x_{j} + \Delta _{k l} \Delta _{j n}x_{i} \wedge x_{m} - \Delta _{k l} \Delta _{i n}x_{j} \wedge x_{m} \nonumber  \\
 &  - \Delta _{i m}x_{j} \wedge x_{k} \wedge x_{l} \wedge x_{n} - \Delta _{i m} \Delta _{\ln }x_{j} \wedge x_{k} + \Delta _{i m} \Delta _{k n}x_{j} \wedge x_{l} - \Delta _{i m} \Delta _{j n}x_{k} \wedge x_{l} \nonumber  \\
 &  + \Delta _{l m}x_{i} \wedge x_{j} \wedge x_{k} \wedge x_{n} - \Delta _{l m} \Delta _{k n}x_{i} \wedge x_{j} + \Delta _{l m} \Delta _{j n}x_{i} \wedge x_{k} - \Delta _{l m} \Delta _{i n}x_{j} \wedge x_{k} \nonumber  \\
 &  - \Delta _{k m}x_{i} \wedge x_{j} \wedge x_{l} \wedge x_{n} + \Delta _{k m} \Delta _{\ln }x_{i} \wedge x_{j} - \Delta _{k m} \Delta _{j n}x_{i} \wedge x_{l} + \Delta _{k m} \Delta _{i n}x_{j} \wedge x_{l} \nonumber  \\
 &  + \Delta _{j m}x_{i} \wedge x_{k} \wedge x_{l} \wedge x_{n} - \Delta _{j m} \Delta _{\ln }x_{i} \wedge x_{k} + \Delta _{j m} \Delta _{k n}x_{i} \wedge x_{l} - \Delta _{j m} \Delta _{i n}x_{k} \wedge x_{l} \nonumber  \\
 &  + \Delta _{l m} \Delta _{k n}x_{i} \wedge x_{j} - \Delta _{l m} \Delta _{j n}x_{i} \wedge x_{k} + \Delta _{l m} \Delta _{i n}x_{j} \wedge x_{k} \nonumber  \\
 &  - \Delta _{k m} \Delta _{\ln }x_{i} \wedge x_{j} + \Delta _{k m} \Delta _{j n}x_{i} \wedge x_{l} - \Delta _{k m} \Delta _{i n}x_{j} \wedge x_{l} \nonumber  \\
 &  + \Delta _{j m} \Delta _{\ln }x_{i} \wedge x_{k} - \Delta _{j m} \Delta _{k n}x_{i} \wedge x_{l} + \Delta _{j m} \Delta _{i n}x_{k} \wedge x_{l} \nonumber  \\
 &  - \Delta _{i m} \Delta _{\ln }x_{j} \wedge x_{k} + \Delta _{i m} \Delta _{k n}x_{j} \wedge x_{l} - \Delta _{i m} \Delta _{j n}x_{k} \wedge x_{l} \nonumber  \\
 &  +\ensuremath{\operatorname*{Pf}} \left ( \Delta _{j k l m}\right ) x_{i} \wedge x_{n} +\ensuremath{\operatorname*{Pf}} \left ( \Delta _{j k l m}\right )  \Delta _{i n} -\ensuremath{\operatorname*{Pf}} \left ( \Delta _{i k l m}\right ) x_{j} \wedge x_{n} \nonumber  \\
 &  -\ensuremath{\operatorname*{Pf}} \left ( \Delta _{i k l m}\right )  \Delta _{j n} +\ensuremath{\operatorname*{Pf}} \left ( \Delta _{i j l m}\right ) x_{k} \wedge x_{n} +\ensuremath{\operatorname*{Pf}} \left ( \Delta _{i j l m}\right ) \Delta _{k n} \nonumber  \\
 &  -\ensuremath{\operatorname*{Pf}} \left ( \Delta _{i j k m}\right ) x_{l} \wedge x_{n} -\ensuremath{\operatorname*{Pf}} \left ( \Delta _{i j k m}\right )  \Delta _{\ln } +\ensuremath{\operatorname*{Pf}} \left ( \Delta _{i j k l}\right ) x_{m} \wedge x_{n} +\ensuremath{\operatorname*{Pf}} \left ( \Delta _{i j k l}\right ) \Delta _{m n}\end{align}
\begin{align} & x_{i} x_{j} x_{k} x_{l} x_{m} x_{n} =x_{i} \wedge x_{j} \wedge x_{k} \wedge x_{l} \wedge x_{m} \wedge x_{n} \nonumber  \\
 &  + \Delta _{i j}x_{k} \wedge x_{l} \wedge x_{m} \wedge x_{n} - \Delta _{i k}x_{j} \wedge x_{l} \wedge x_{m} \wedge x_{n} + \Delta _{i l}x_{j} \wedge x_{k} \wedge x_{m} \wedge x_{n} - \Delta _{i m}x_{j} \wedge x_{k} \wedge x_{l} \wedge x_{n} \nonumber  \\
 &  + \Delta _{i n}x_{j} \wedge x_{k} \wedge x_{l} \wedge x_{m} \nonumber  \\
 &  + \Delta _{j k}x_{i} \wedge x_{l} \wedge x_{m} \wedge x_{n} - \Delta _{j l}x_{i} \wedge x_{k} \wedge x_{m} \wedge x_{n} + \Delta _{j m}x_{i} \wedge x_{k} \wedge x_{l} \wedge x_{n} - \Delta _{j n}x_{i} \wedge x_{k} \wedge x_{l} \wedge x_{m} \nonumber  \\
 &  + \Delta _{k l}x_{i} \wedge x_{j} \wedge x_{m} \wedge x_{n} - \Delta _{k m}x_{i} \wedge x_{j} \wedge x_{l} \wedge x_{n} + \Delta _{k n}x_{i} \wedge x_{j} \wedge x_{l} \wedge x_{m} \nonumber  \\
 &  + \Delta _{l m}x_{i} \wedge x_{j} \wedge x_{k} \wedge x_{n} - \Delta _{\ln }x_{i} \wedge x_{j} \wedge x_{k} \wedge x_{m} + \Delta _{m n}x_{i} \wedge x_{j} \wedge x_{k} \wedge x_{l} \nonumber  \\
 &  +\left \{ \Delta _{l m} \Delta _{k n} -  \Delta _{l m} \Delta _{j n} -  \Delta _{k m} \Delta _{\ln }\right \} x_{i} \wedge x_{j} +\left \{ \Delta _{j m} \Delta _{\ln } -  \Delta _{l m} \Delta _{j n} +  \Delta _{j l} \Delta _{m n}\right \} x_{i} \wedge x_{k} \nonumber  \\
 &  +\left \{ \Delta _{k m} \Delta _{j n} -  \Delta _{j m} \Delta _{k n} -  \Delta _{j k} \Delta _{m n}\right \} x_{i} \wedge x_{l} +\left \{ \Delta _{j k} \Delta _{\ln } -  \Delta _{j l} \Delta _{k n} +  \Delta _{k l} \Delta _{j n}\right \} x_{i} \wedge x_{m} \nonumber  \\
 &  -\left \{ \Delta _{i m} \Delta _{\ln } +  \Delta _{l m} \Delta _{i n} -  \Delta _{i l} \Delta _{m n}\right \} x_{j} \wedge x_{k} -\left \{ \Delta _{k m} \Delta _{i n} +  \Delta _{i m} \Delta _{k n} +  \Delta _{i k} \Delta _{m n}\right \} x_{j} \wedge x_{l} \nonumber  \\
 &  -\left \{ \Delta _{i k} \Delta _{\ln } +  \Delta _{i l} \Delta _{k n} -  \Delta _{k l} \Delta _{i n}\right \} x_{j} \wedge x_{m} +\left \{ \Delta _{j m} \Delta _{i n} -  \Delta _{i m} \Delta _{j n} +  \Delta _{i j} \Delta _{m n}\right \} x_{k} \wedge x_{l} \nonumber  \\
 &  -\left \{ \Delta _{i j} \Delta _{\ln } +  \Delta _{j l} \Delta _{i n} -  \Delta _{i l} \Delta _{j n}\right \} x_{k} \wedge x_{m} +\left \{ \Delta _{i j} \Delta _{k n} -  \Delta _{j k} \Delta _{i n} +  \Delta _{i k} \Delta _{j n}\right \} x_{l} \wedge x_{m} \nonumber  \\
 &  +\ensuremath{\operatorname*{Pf}} \left (\mathbf{\Delta }_{j k l m}\right ) x_{i} \wedge x_{n} -\ensuremath{\operatorname*{Pf}} \left (\mathbf{\Delta }\right )_{i k l m} x_{j} \wedge x_{n} -\ensuremath{\operatorname*{Pf}} \left (\mathbf{\Delta }\right )_{i j k m} x_{l} \wedge x_{n} \nonumber  \\
 &  +\ensuremath{\operatorname*{Pf}} \left (\mathbf{\Delta }\right )_{i j l m} x_{k} \wedge x_{n} +\ensuremath{\operatorname*{Pf}} \left (\mathbf{\Delta }\right )_{i j k l} x_{m} \wedge x_{n} +\ensuremath{\operatorname*{Pf}} \left (\mathbf{\Delta }\right )_{j k l m}  \Delta _{i n} -\ensuremath{\operatorname*{Pf}} \left (\mathbf{\Delta }\right )_{i j k m} \Delta _{\ln } \nonumber  \\
 &  -\ensuremath{\operatorname*{Pf}} \left (\mathbf{\Delta }\right )_{i k l m}  \Delta _{j n} +\ensuremath{\operatorname*{Pf}} \left (\mathbf{\Delta }\right )_{i j k l}  \Delta _{m n} +\ensuremath{\operatorname*{Pf}} \left (\mathbf{\Delta }\right )_{i j l m} \Delta _{k n}\end{align}and finally
\begin{align} & x_{j_{1}} x_{j} x_{j_{3}} x_{j_{4}} x_{j_{5}} x_{j_{6}} =x_{j_{1}} \wedge x_{j_{2}} \wedge x_{j_{3}} \wedge x_{j_{4}} \wedge x_{j_{5}} \wedge x_{j_{6}} \nonumber  \\
 &  + \Delta _{j_{1} j_{2}}x_{j_{3}} \wedge x_{j_{4}} \wedge x_{j_{5}} \wedge x_{j_{6}} - \Delta _{j_{1} j_{3}}x_{j_{2}} \wedge x_{j_{4}} \wedge x_{j_{5}} \wedge x_{j_{6}} + \Delta _{j_{1} j_{4}}x_{j_{2}} \wedge x_{j_{3}} \wedge x_{j_{5}} \wedge x_{j_{6}} - \Delta _{j_{1} j_{5}}x_{j_{2}} \wedge x_{j_{3}} \wedge x_{j_{4}} \wedge x_{j_{6}} \nonumber  \\
 &  + \Delta _{j_{1} j_{6}}x_{j_{2}} \wedge x_{j_{3}} \wedge x_{j_{4}} \wedge x_{j_{5}} + \Delta _{j_{2} j_{3}}x_{j_{1}} \wedge x_{j_{4}} \wedge x_{j_{5}} \wedge x_{j_{6}} - \Delta _{j_{2} j_{4}}x_{j_{1}} \wedge x_{j_{3}} \wedge x_{j_{5}} \wedge x_{j_{6}} + \Delta _{j_{2} j_{5}}x_{j_{1}} \wedge x_{j_{3}} \wedge x_{j_{4}} \wedge x_{j_{6}} \nonumber  \\
 &  - \Delta _{j_{2} j_{6}}x_{j_{1}} \wedge x_{j_{3}} \wedge x_{j_{4}} \wedge x_{j_{5}} + \Delta _{j_{3} j_{4}}x_{j_{1}} \wedge x_{j_{2}} \wedge x_{j_{5}} \wedge x_{j_{6}} - \Delta _{j_{3} j_{5}}x_{j_{1}} \wedge x_{j_{2}} \wedge x_{j_{4}} \wedge x_{j_{6}} + \Delta _{j_{3} j_{6}}x_{j_{1}} \wedge x_{j_{2}} \wedge x_{j_{4}} \wedge x_{j_{5}} \nonumber  \\
 &  + \Delta _{j_{4} j_{5}}x_{j_{1}} \wedge x_{j_{2}} \wedge x_{j_{3}} \wedge x_{j_{6}} - \Delta _{\ln }x_{j_{1}} \wedge x_{j_{2}} \wedge x_{j_{3}} \wedge x_{j_{5}} + \Delta _{j_{5} j_{6}}x_{j_{1}} \wedge x_{j_{2}} \wedge x_{j_{3}} \wedge x_{j_{4}} \nonumber  \\
 &  +\ensuremath{\operatorname*{Pf}} \left ( \Delta _{j_{4} j_{5} j_{3} j_{6}}\right ) x_{j_{1}} \wedge x_{j_{2}} +\ensuremath{\operatorname*{Pf}} \left ( \Delta _{j_{2} j_{5}\ln }\right ) x_{j_{1}} \wedge x_{j_{3}} +\ensuremath{\operatorname*{Pf}} \left ( \Delta _{j_{3} j_{5} j_{2} j_{6}}\right ) x_{j_{1}} \wedge x_{j_{4}} +\ensuremath{\operatorname*{Pf}} \left ( \Delta _{j_{2} j_{3}\ln }\right ) x_{j_{1}} \wedge x_{j_{5}} -\ensuremath{\operatorname*{Pf}} \left ( \Delta _{j_{1} j_{5}\ln }\right ) x_{j_{2}} \wedge x_{j_{3}} \nonumber  \\
 &  -\ensuremath{\operatorname*{Pf}} \left ( \Delta _{j_{3}\min }\right ) x_{j_{2}} \wedge x_{j_{4}} -\ensuremath{\operatorname*{Pf}} \left ( \Delta _{j_{1} j_{3}\ln }\right ) x_{j_{2}} \wedge x_{j_{5}} +\ensuremath{\operatorname*{Pf}} \left ( \Delta _{j_{2}\min }\right ) x_{j_{3}} \wedge x_{j_{4}} -\ensuremath{\operatorname*{Pf}} \left ( \Delta _{j_{1} j_{2}\ln }\right ) x_{j_{3}} \wedge x_{j_{5}} +\ensuremath{\operatorname*{Pf}} \left ( \Delta _{j_{1} j_{2} j_{3} j_{6}}\right ) x_{j_{4}} \wedge x_{j_{5}} \nonumber  \\
 &  +\ensuremath{\operatorname*{Pf}} \left ( \Delta _{j_{2} j_{3} j_{4} j_{5}}\right ) x_{j_{1}} \wedge x_{j_{6}} -\ensuremath{\operatorname*{Pf}} \left ( \Delta _{j_{1} j_{3} j_{4} j_{5}}\right ) x_{j_{2}} \wedge x_{j_{6}} -\ensuremath{\operatorname*{Pf}} \left ( \Delta _{j_{1} j_{2} j_{3} j_{5}}\right ) x_{j_{4}} \wedge x_{j_{6}} +\ensuremath{\operatorname*{Pf}} \left ( \Delta _{j_{1} j_{2} j_{4} j_{5}}\right ) x_{j_{3}} \wedge x_{j_{6}} +\ensuremath{\operatorname*{Pf}} \left ( \Delta _{j_{1} j_{2} j_{3} j_{4}}\right ) x_{j_{5}} \wedge x_{j_{6}} \nonumber  \\
 &  +\left \{\ensuremath{\operatorname*{Pf}} \left ( \Delta _{j_{1} j_{2} j_{3} j_{4} j_{5} j_{6}}\right ) -\ensuremath{\operatorname*{Pf}} \left ( \Delta _{j_{1} j_{2} j_{4} j_{3} j_{5} j_{6}}\right ) -\ensuremath{\operatorname*{Pf}} \left ( \Delta _{j_{1} j_{3} j_{2} j_{4} j_{5} j_{6}}\right ) +\ensuremath{\operatorname*{Pf}} \left ( \Delta _{j_{1} j_{2} j_{5} j_{3} j_{4} j_{6}}\right ) +\ensuremath{\operatorname*{Pf}} \left ( \Delta _{j_{1} j_{2} j_{4} j_{5} j_{3} j_{6}}\right )\right \} I_{\mathcal{F}}\text{,}\end{align}where
\begin{align} & \ensuremath{\operatorname*{Pf}} \left ( \Delta _{j_{1} j_{2} j_{3} j_{4} j_{5} j_{6}}\right ) =\ensuremath{\operatorname*{Pf}} \left ( \Delta _{j_{1} j_{2} j_{3} j_{4}}\right )  \Delta _{j_{5} j_{6}} -\ensuremath{\operatorname*{Pf}} \left ( \Delta _{j_{1} j_{2} j_{3} j_{5}}\right ) \Delta _{j_{4} j_{6}} \nonumber  \\
 &  +\ensuremath{\operatorname*{Pf}} \left ( \Delta _{j_{1} j_{2} j_{4} j_{5}}\right )  \Delta _{j_{3} j_{6}} -\ensuremath{\operatorname*{Pf}} \left ( \Delta _{j_{1} j_{3} j_{4} j_{5}}\right )  \Delta _{j_{2} j_{6}} +\ensuremath{\operatorname*{Pf}} \left ( \Delta _{j_{2} j_{3} j_{4} j_{5}}\right ) \Delta _{j_{1} j_{6}}\text{.}\end{align}

In summary 

2. The second degree products are given by
\begin{equation*}x_{i} x_{j} =x_{i} \wedge x_{j} + \Delta _{i j}\text{.}
\end{equation*}

3. The third degree products are given by 

\begin{equation}x_{j_{1}} x_{j_{2}} x_{j_{3}} =x_{j_{1}} \wedge x_{j_{2}} \wedge x_{j_{3}} + \Delta _{j_{2} j_{3}}x_{j_{1}} - \Delta _{j_{1} j_{3}}x_{j_{2}} + \Delta _{j_{1} j_{2}}x_{j_{3}}\text{.}
\end{equation}

4. The fourth degree products are given by
\begin{align} & x_{j_{1}} x_{j_{2}} x_{j_{3}} x_{j_{4}} =x_{j_{1}} \wedge x_{j} \wedge x_{j_{3}} \wedge x_{j_{4}} \nonumber  \\
 &  + \Delta _{j_{1} j_{2}}x_{j_{3}} \wedge x_{j_{4}} - \Delta _{j_{1} j_{3}}x_{j_{2}} \wedge x_{j_{4}} + \Delta _{j_{1} j_{4}}x_{j_{2}} \wedge x_{j_{3}} + \Delta _{j_{2} j_{3}}x_{j_{1}} \wedge x_{j_{4}} - \Delta _{j_{2} j_{4}}x_{j_{1}} \wedge x_{j_{3}} + \Delta _{j_{3} j_{4}}x_{j_{1}} \wedge x_{j_{2}} \nonumber  \\
 &  +\ensuremath{\operatorname*{Pf}} \left (\mathbf{\Delta }_{j_{1} j_{2} j_{3} j_{4}}\right ) I_{\mathcal{F}}\text{.}\end{align}

5. The fifth degree products are given by
\begin{align} & x_{j_{1}} x_{j_{2}} x_{j_{3}} x_{j_{4}} x_{j_{5}} =x_{j_{1}} \wedge x_{j_{2}} \wedge x_{j_{3}} \wedge x_{j_{4}} \wedge x_{j_{5}} \nonumber  \\
 &  + \Delta _{j_{1} j_{2}}x_{j_{3}} \wedge x_{j_{4}} \wedge x_{j_{5}} - \Delta _{j_{1} j_{3}}x_{j_{2}} \wedge x_{j_{4}} \wedge x_{j_{5}} + \Delta _{j_{1} j_{4}}x_{j_{2}} \wedge x_{j_{3}} \wedge x_{j_{5}} + \Delta _{j_{2} j_{3}}x_{j_{1}} \wedge x_{j_{4}} \wedge x_{j_{5}} \nonumber  \\
 &  - \Delta _{j_{2} j_{4}}x_{j_{1}} \wedge x_{j_{3}} \wedge x_{j_{5}} + \Delta _{j_{3} j_{4}}x_{j_{1}} \wedge x_{j_{2}} \wedge x_{j_{5}} + \Delta _{j_{4} j_{5}}x_{j_{1}} \wedge x_{j_{2}} \wedge x_{j_{3}} - \Delta _{j_{3} j_{5}}x_{j_{1}} \wedge x_{j_{2}} \wedge x_{j_{4}} \nonumber  \\
 &  + \Delta _{j_{2} j_{5}}x_{j_{1}} \wedge x_{j_{3}} \wedge x_{j_{4}} - \Delta _{j_{1} j_{5}}x_{j_{2}} \wedge x_{j_{3}} \wedge x_{j_{4}} \nonumber  \\
 &  +\ensuremath{\operatorname*{Pf}} \left (\mathbf{\Delta }_{j_{2} j_{3} j_{4} j_{5}}\right ) x_{j_{1}} -\ensuremath{\operatorname*{Pf}} \left (\mathbf{\Delta }_{j_{1} j_{3} j_{4} j_{5}}\right ) x_{j} +\ensuremath{\operatorname*{Pf}} \left (\mathbf{\Delta }_{j_{1} j_{2} j_{4} j_{5}}\right ) x_{j_{3}} -\ensuremath{\operatorname*{Pf}} \left (\mathbf{\Delta }_{j_{1} j_{2} j_{3} j_{5}}\right ) x_{j_{4}} +\ensuremath{\operatorname*{Pf}} \left (\mathbf{\Delta }_{j_{1} j_{2} j_{3} j_{4}}\right ) x_{j_{5}}\text{,}\end{align}which can also be expressed as
\begin{align} &  =x_{j_{1}} \wedge x_{j_{2}} \wedge x_{j_{3}} \wedge x_{j_{4}} \wedge x_{j_{5}} \nonumber  \\
 &  +\sum _{\sigma  \in S_{5}/\left (S_{2} \times S_{3}\right )}\pi  \left (\sigma \right )  \Delta _{j_{\sigma  \left (1\right )} j \sigma  \left (_{2}\right )}x_{j_{\sigma  \left (_{3}\right )}} \wedge x_{j_{\sigma  \left (_{4}\right )}} \wedge x_{j_{\sigma  \left (5\right )}} +\sum _{\sigma  \in S_{5}/\left (S_{4} \times S_{1}\right )}\pi  \left (\sigma \right ) \ensuremath{\operatorname*{Pf}} \left (\mathbf{\Delta }_{j_{\sigma  \left (1\right )} j_{\sigma  \left (_{2}\right )} j_{\sigma  \left (_{3}\right )} j_{\sigma  \left (_{4}\right )}}\right ) x_{j_{\sigma  \left (5\right )}}\text{.}\end{align}

6. The sixth degree products are given by
\begin{align} & x_{j_{1}} x_{j_{2}} x_{j_{3}} x_{j_{4}} x_{j_{5}} x_{j_{6}} =x_{j_{1}} \wedge x_{j_{2}} \wedge x_{j_{3}} \wedge x_{j_{4}} \wedge x_{j_{5}} \wedge x_{j_{6}} \nonumber  \\
 &  + \Delta _{j_{1} j_{2}}x_{j_{3}} \wedge x_{j_{4}} \wedge x_{j_{5}} \wedge x_{j_{6}} - \Delta _{j_{1} j_{3}}x_{j_{2}} \wedge x_{j_{4}} \wedge x_{j_{5}} \wedge x_{j_{6}} + \Delta _{j_{1} j_{4}}x_{j_{2}} \wedge x_{j_{3}} \wedge x_{j_{5}} \wedge x_{j_{6}} \nonumber  \\
 &  - \Delta _{j_{1} j_{5}}x_{j_{2}} \wedge x_{j_{3}} \wedge x_{j_{4}} \wedge x_{j_{6}} + \Delta _{j_{1} j_{6}}x_{j_{2}} \wedge x_{j_{3}} \wedge x_{j_{4}} \wedge x_{j_{5}} + \Delta _{j_{2} j_{3}}x_{j_{1}} \wedge x_{j_{4}} \wedge x_{j_{5}} \wedge x_{j_{6}} \nonumber  \\
 &  - \Delta _{j_{2} j_{4}}x_{j_{1}} \wedge x_{j_{3}} \wedge x_{j_{5}} \wedge x_{j_{6}} + \Delta _{j_{2} j_{5}}x_{j_{1}} \wedge x_{j_{3}} \wedge x_{j_{4}} \wedge x_{j_{6}} - \Delta _{j_{2} j_{6}}x_{j_{1}} \wedge x_{j_{3}} \wedge x_{j_{4}} \wedge x_{j_{5}} \nonumber  \\
 &  + \Delta _{j_{3} j_{4}}x_{j_{1}} \wedge x_{j_{2}} \wedge x_{j_{5}} \wedge x_{j_{6}} - \Delta _{j_{3} j_{5}}x_{j_{1}} \wedge x_{j_{2}} \wedge x_{j_{4}} \wedge x_{j_{6}} + \Delta _{j_{3} j_{6}}x_{j_{1}} \wedge x_{j_{2}} \wedge x_{j_{4}} \wedge x_{j_{5}} \nonumber  \\
 &  + \Delta _{j_{4} j_{5}}x_{j_{1}} \wedge x_{j_{2}} \wedge x_{j_{3}} \wedge x_{j_{6}} - \Delta _{j_{4} j_{6}}x_{j_{1}} \wedge x_{j_{2}} \wedge x_{j_{3}} \wedge x_{j_{5}} + \Delta _{j_{5} j_{6}}x_{j_{1}} \wedge x_{j_{2}} \wedge x_{j_{3}} \wedge x_{j_{4}} \nonumber  \\
 &  +\ensuremath{\operatorname*{Pf}} \left ( \Delta _{j_{4} j_{5} j_{3} j_{6}}\right ) x_{j_{1}} \wedge x_{j_{2}} +\ensuremath{\operatorname*{Pf}} \left ( \Delta _{j_{2} j_{5} j_{4} j_{6}}\right ) x_{j_{1}} \wedge x_{j_{3}} +\ensuremath{\operatorname*{Pf}} \left ( \Delta _{j_{3} j_{5} j_{2} j_{6}}\right ) x_{j_{1}} \wedge x_{j_{4}} +\ensuremath{\operatorname*{Pf}} \left ( \Delta _{j_{2} j_{3} j_{4} j_{6}}\right ) x_{j_{1}} \wedge x_{j_{5}} \nonumber  \\
 &  -\ensuremath{\operatorname*{Pf}} \left ( \Delta _{j_{1} j_{5} j_{4} j_{6}}\right ) x_{j_{2}} \wedge x_{j_{3}} -\ensuremath{\operatorname*{Pf}} \left ( \Delta _{j_{3} j_{5} j_{1} j_{6}}\right ) x_{j_{2}} \wedge x_{j_{4}} -\ensuremath{\operatorname*{Pf}} \left ( \Delta _{j_{1} j_{3} j_{4} j_{6}}\right ) x_{j_{2}} \wedge x_{j_{5}} +\ensuremath{\operatorname*{Pf}} \left ( \Delta _{j_{2} j_{1} j_{6}}\right ) x_{j_{3}} \wedge x_{j_{4}} \nonumber  \\
 &  -\ensuremath{\operatorname*{Pf}} \left ( \Delta _{j_{1} j_{2} j_{4} j_{6}}\right ) x_{j_{3}} \wedge x_{j_{5}} +\ensuremath{\operatorname*{Pf}} \left ( \Delta _{j_{1} j_{2} j_{3} j_{6}}\right ) x_{j_{4}} \wedge x_{j_{5}} +\ensuremath{\operatorname*{Pf}} \left ( \Delta _{j_{1} j_{2} j_{3} j_{4}}\right ) x_{j_{5}} \wedge x_{j_{6}} +\ensuremath{\operatorname*{Pf}} \left ( \Delta _{j_{2} j_{3} j_{4} j_{5}}\right ) x_{j_{1}} \wedge x_{j_{6}} \nonumber  \\
 &  -\ensuremath{\operatorname*{Pf}} \left ( \Delta _{j_{1} j_{3} j_{4} j_{5}}\right ) x_{j_{2}} \wedge x_{j_{6}} -\ensuremath{\operatorname*{Pf}} \left ( \Delta _{j_{1} j_{2} j_{3} j_{5}}\right ) x_{j_{4}} \wedge x_{j_{6}} +\ensuremath{\operatorname*{Pf}} \left ( \Delta _{j_{1} j_{2} j_{4} j_{5}}\right ) x_{j_{3}} \wedge x_{j_{6}} \nonumber  \\
 &  +\left \{\ensuremath{\operatorname*{Pf}} \left ( \Delta _{j_{1} j_{2} j_{3} j_{4} j_{5} j_{6}}\right ) +\ensuremath{\operatorname*{Pf}} \left ( \Delta _{j_{1} j_{2} j_{3} j_{4} j_{5} j_{6}}\right ) -\ensuremath{\operatorname*{Pf}} \left ( \Delta _{j_{1} j_{2} j_{4} j_{3} j_{5} j_{6}}\right ) -\ensuremath{\operatorname*{Pf}} \left ( \Delta _{j_{1} j_{3} j_{2} j_{4} j_{5} j_{6}}\right ) +\ensuremath{\operatorname*{Pf}} \left ( \Delta _{j_{1} j_{2} j_{5} j_{3} j_{4} j_{6}}\right )\right \} I_{\mathcal{F}}\text{,}\end{align}which can also be expressed as
\begin{align} &  =x_{j_{1}} \wedge x_{j_{2}} \wedge x_{j_{3}} \wedge x_{j_{4}} \wedge x_{j_{5}} \wedge x_{j_{6}} \nonumber  \\
 &  +\sum _{\sigma  \in S_{6}/\left (S_{2} \times S_{4}\right )}\pi  \left (\sigma \right )  \Delta _{j_{\sigma  \left (1\right )} j_{\sigma  \left (2\right )}}x_{j_{\sigma  \left (3\right )}} \wedge x_{j_{\sigma  \left (4\right )}} \wedge x_{j_{\sigma  \left (5\right )}} \wedge x_{j_{\sigma  \left (6\right )}} \nonumber  \\
 & \text{\quad \quad \quad \quad \quad \quad \quad \quad } +\sum _{\sigma  \in S_{6}/\left (S_{4} \times S_{2}\right )}\pi  \left (\sigma \right ) \ensuremath{\operatorname*{Pf}} \left (\mathbf{\Delta }_{j_{\sigma  \left (1\right )} j_{\sigma  \left (2\right )} j_{\sigma  \left (3\right )} j_{\sigma  \left (4\right )}}\right ) x_{j_{\sigma  \left (5\right )}} \wedge x_{j_{\sigma  \left (6\right )}} +\ensuremath{\operatorname*{Pf}} \left (\mathbf{\Delta }_{j_{1} j_{2} j_{3} j_{4} j_{5} j_{6}}\right ) I_{\mathcal{F}}\text{.}\end{align}

7. The seventh degree products are given by
\begin{align} & x_{j_{1}} x_{j_{2}} x_{j_{3}} x_{j_{4}} x_{j_{5}} x_{j_{6}} x_{j_{7}} =x_{j_{1}} \wedge x_{j_{2}} \wedge x_{j_{3}} \wedge x_{j_{4}} \wedge x_{j_{5}} \wedge x_{j_{6}} \wedge x_{j_{7}} \nonumber  \\
 &  +\sum _{\sigma  \in S_{7}/\left (S_{2} \times S_{5}\right )}\pi  \left (\sigma \right )  \Delta _{j_{\sigma  \left (1\right )} j \sigma  \left (_{2}\right )}x_{j_{\sigma  \left (_{3}\right )}} \wedge x_{j_{\sigma  \left (_{4}\right )}} \wedge x_{j_{\sigma  \left (5\right )}} \wedge x_{j_{\sigma  \left (_{6}\right )}} \wedge x_{j_{\sigma  \left (7\right )}} \nonumber  \\
 &  +\sum _{\sigma  \in S_{7}/\left (S_{4} \times S_{3}\right )}\pi  \left (\sigma \right ) \ensuremath{\operatorname*{Pf}} \left ( \Delta _{j_{\sigma  \left (1\right )} j_{\sigma  \left (_{2}\right )} j_{\sigma  \left (_{3}\right )} j_{\sigma  \left (_{4}\right )}}\right ) x_{j_{\sigma  \left (_{5}\right )}} \wedge x_{j_{\sigma  \left (_{6}\right )}} \wedge x_{j_{\sigma  \left (7\right )}} \nonumber  \\
 &  +\sum _{\sigma  \in S_{7}/\left (S_{6} \times S_{1}\right )}\pi  \left (\sigma \right ) \ensuremath{\operatorname*{Pf}} \left ( \Delta _{j_{\sigma  \left (1\right )} j_{\sigma  \left (2\right )} j_{\sigma  \left (3\right )} j_{\sigma  \left (4\right )} j_{\sigma  \left (5\right )} j_{\sigma  \left (6\right )}}\right ) x_{j_{\sigma  \left (7\right )}}\end{align}

8. The eigth degree products are given by
\begin{align} & x_{j_{1}} x_{j_{2}} x_{j_{3}} x_{j_{4}} x_{j_{5}} x_{j_{6}} x_{j_{7}} x_{j_{8}} =x_{j_{1}} \wedge x_{j_{2}} \wedge x_{j_{3}} \wedge x_{j_{4}} \wedge x_{j_{5}} \wedge x_{j_{6}} \wedge x_{j_{7}} \wedge x_{j_{8}} \nonumber  \\
 &  +\sum _{\sigma  \in S_{8}/\left (S_{2} \times S_{6}\right )}\pi  \left (\sigma \right )  \Delta _{j_{\sigma  \left (1\right )} j_{\sigma  \left (2\right )}}x_{j_{\sigma  \left (3\right )}} \wedge x_{j_{\sigma  \left (4\right )}} \wedge x_{j_{\sigma  \left (5\right )}} \wedge x_{j_{\sigma  \left (6\right )}} \wedge x_{j_{\sigma  \left (7\right )}} \wedge x_{j_{\sigma  \left (8\right )}} \nonumber  \\
 &  +\sum _{\sigma  \in S_{8}/\left (S_{4} \times S_{4}\right )}\pi  \left (\sigma \right ) \ensuremath{\operatorname*{Pf}} \left ( \Delta _{j_{\sigma  \left (1\right )} j_{\sigma  \left (2\right )} j_{\sigma  \left (3\right )} j_{\sigma  \left (4\right )}}\right ) x_{j_{\sigma  \left (5\right )}} \wedge x_{j_{\sigma  \left (6\right )}} \wedge x_{j_{\sigma  \left (7\right )}} \wedge x_{j_{\sigma  \left (8\right )}} \nonumber  \\
 &  +\sum _{\sigma  \in S_{8}/\left (S_{6} \times S_{2}\right )}\pi  \left (\sigma \right ) \ensuremath{\operatorname*{Pf}} \left ( \Delta _{j_{\sigma  \left (1\right )} j_{\sigma  \left (2\right )} j_{\sigma  \left (3\right )} j_{\sigma  \left (4\right )} j_{\sigma  \left (5\right )} j_{\sigma  \left (6\right )}}\right ) x_{j_{\sigma  \left (7\right )}} \wedge x_{j_{\sigma  \left (8\right )}} +\ensuremath{\operatorname*{Pf}} \left ( \Delta _{j_{1} j_{2} j_{3} j_{4} j_{5} j_{6} j_{7} j_{8}}\right ) I_{\mathcal{F}}\text{.}\end{align}


\begin{theorem}
Grassmann-Clifford Algebra Equivalence (G.P.Wilmount, J. Math. Phys. \textbf{29,}2338-2345 (1988)) 
\end{theorem}

In
general one has the following theorem that relates the exterior and geometric (Clifford) products 

\begin{equation}x_{j_{1}} \cdots  x_{j_{2 q}} ={\displaystyle\prod _{1 \leq k \leq 2 q}}x_{j_{k}} ={\displaystyle\sum _{0 \leq l \leq q}}{\displaystyle\sum _{\begin{array}{c}\sigma  \in C_{2}^{2 l} \left (\mathbb{N}_{1}^{2 l}\right ) \\
l_{1} +l_{2} =l\end{array}}}\pi  \left (\sigma \right ) x_{j_{\sigma  \left (_{1}\right )}} \wedge \cdots  \wedge x_{j_{\sigma  \left (l_{1}\right )}} \ensuremath{\operatorname*{Pf}} \left (\mathbf{\Delta }_{j_{\sigma  \left (l_{1} +1\right )} \cdots  j_{\sigma  \left (2 l_{2}\right )}}\right )\text{,}
\end{equation}where
\begin{equation}l_{1} =0 \Longrightarrow x_{j_{\sigma  \left (_{1}\right )}} \wedge \cdots  \wedge x_{j_{\sigma  \left (l_{1}\right )}} =I_{\mathcal{F}}\text{.}
\end{equation}

\subsection{Quantization}
The above theorem defines the \emph{Symbol} map\emph{, }$\mathfrak{S}\text{,}$
\begin{align}\mathfrak{S} &  :\mathcal{C} \mathcal{L} \left (V ,\Delta \right ) \longrightarrow \bigwedge \left (V\right ) \\
\mathfrak{S} \left (v_{j_{1}} \cdots  v_{j_{2 q}}\right ) &  ={\displaystyle\sum _{0 \leq l \leq q}}{\displaystyle\sum _{\begin{array}{c}\sigma  \in C_{2}^{2 l} \left (\mathbb{N}_{1}^{2 l}\right ) \\
l_{1} +l_{2} =l\end{array}}}\pi  \left (\sigma \right ) v_{j_{\sigma  \left (_{1}\right )}} \wedge \cdots  \wedge v_{j_{\sigma  \left (l_{1}\right )}} \ensuremath{\operatorname*{Pf}} \left (\mathbf{\Delta }_{j_{\sigma  \left (l_{1} +1\right )} \cdots  j_{\sigma  \left (2 l_{2}\right )}}\right )\text{.}\end{align}\emph{\ } 

The inverse of the Symbol map
is the \emph{Quantization} map $\mathfrak{Q}$
\begin{align}\mathfrak{Q} &  :\bigwedge \left (V\right ) \longrightarrow \mathcal{C} \mathcal{L} \left (V ,\Delta \right ) \\
\mathfrak{Q} \left (v_{j_{1}} \wedge \cdots  \wedge v_{j_{k}}\right ) &  =\frac{1}{k !} \sum _{\sigma  \in S_{k}}\pi  \left (\sigma \right ) v_{j_{\sigma  \left (_{1}\right )}} \cdots  v_{j_{\sigma  \left (k\right )}}\text{.}\end{align}

This can also be expressed as
\begin{equation}\mathfrak{Q} \left (v_{j_{1}} \wedge \cdots  \wedge v_{j_{k}}\right ) = -{\displaystyle\sum _{0 \leq l \leq k}}{\displaystyle\sum _{\begin{array}{c}\sigma  \in C_{2}^{2 l} \left (\mathbb{N}_{1}^{2 l}\right ) \\
l_{1} +l_{2} =l\end{array}}}\pi  \left (\sigma \right ) x_{j_{\sigma  \left (_{1}\right )}} \wedge \cdots  \wedge x_{j_{\sigma  \left (l_{1}\right )}} \ensuremath{\operatorname*{Pf}} \left (\mathbf{\Delta }_{j_{\sigma  \left (l_{1} +1\right )} \cdots  j_{\sigma  \left (2 l_{2}\right )}}\right )\text{,}
\end{equation}which gives in low orders
\begin{align}\mathfrak{Q} \left (I\right ) &  =I \\
\mathfrak{Q} \left (v\right ) &  =v \\
\mathfrak{Q} \left (v_{1} \wedge v_{2}\right ) &  =v_{1} v_{2} - \Delta \left (v_{1} ,v_{2}\right ) I \\
\mathfrak{Q} \left (v_{1} \wedge v_{2} \wedge v_{3}\right ) &  =v_{1} v_{2} v_{3} - \Delta \left (v_{2} ,v_{3}\right ) v_{1} + \Delta \left (v_{1} ,v_{3}\right ) v_{2} - \Delta \left (v_{1} ,v_{2}\right ) v_{3}\text{,}\end{align}where
\begin{equation} \Delta \left (v_{1} ,v_{2}\right ) = \langle  \langle v_{1} v_{2} \rangle  \rangle _{0} =\left (\left (\left .v_{1}\right \vert v_{2}\right )\right )\text{.}
\end{equation}

\subsection{Idempotent Operators}


\subsubsection{Second Degree Operators}
First consider the homogeneous second degree operator
\begin{equation}A_{2} =\sum _{1 \leq j_{1} <j_{2} \leq 2 r}A_{2 j_{1} j_{2}} x_{j_{1}} x_{j_{2}}
\end{equation}and its Clifford/Operator square
\begin{equation}\left (A_{2}\right )^{2} =\sum _{\begin{array}{c}1 \leq j_{1} <j_{2} \leq 2 r \\
1 \leq j_{3} <j_{4} \leq 2 r\end{array}}A_{2 j_{1} j_{2}} A_{2 j_{3} j_{4}} x_{j_{1}} x_{j_{2}} x_{j_{3}} x_{j_{4}}\text{.}
\end{equation}Utilizing the relationship
\begin{align} & x_{j_{1}} x_{j_{2}} x_{j_{3}} x_{j_{4}} =x_{j_{1}} \wedge x_{j_{2}} \wedge x_{j_{3}} \wedge x_{j_{4}} \nonumber  \\
 &  +\left (\left .x_{j_{1}}\right \vert x_{j_{2}}\right ) x_{j_{3}} \wedge x_{j_{4}} -\left (\left .x_{j_{1}}\right \vert x_{j_{3}}\right ) x_{j_{2}} \wedge x_{j_{4}} +\left (\left .x_{j_{1}}\right \vert x_{j_{4}}\right ) x_{j_{2}} \wedge x_{j_{3}} \nonumber  \\
 &  +\left (\left .x_{j_{2}}\right \vert x_{j_{3}}\right ) x_{j_{1}} \wedge x_{j_{4}} -\left (\left .x_{j_{2}}\right \vert x_{j_{4}}\right ) x_{j_{1}} \wedge x_{j_{3}} +\left (\left .x_{j_{3}}\right \vert x_{j_{4}}\right ) x_{j_{1}} \wedge x_{j_{2}} \nonumber  \\
 &  +\left (\left .x_{j_{1}}\right \vert x_{j_{2}}\right ) \left (\left .x_{j_{3}}\right \vert x_{j_{4}}\right ) -\left (\left .x_{j_{1}}\right \vert x_{j_{3}}\right ) \left (\left .x_{j_{2}}\right \vert x_{j_{4}}\right ) +\left (\left .x_{j_{1}}\right \vert x_{j_{4}}\right ) \left (\left .x_{j_{2}}\right \vert x_{j_{3}}\right )\text{,}\end{align}one obtains
\begin{align} & \left (A_{2}\right )^{2} =\sum _{\begin{array}{c}1 \leq j_{1} <j_{2} \leq 2 r \\
1 \leq j_{3} <j_{4} \leq 2 r\end{array}}\left \{A_{2 j_{1} j_{2}} A_{2 j_{3} j_{4}} x_{j_{1}} \wedge x_{j_{2}} \wedge x_{j_{3}} \wedge x_{j_{4}}\right . \nonumber  \\
 &  - \Delta _{j_{1} j_{3}}A_{2 j_{1} j_{2}} A_{2 j_{3} j_{4}} x_{j_{2}} \wedge x_{j_{4}} + \Delta _{j_{1} j_{4}}A_{2 j_{1} j_{2}} A_{2 j_{3} j_{4}} x_{j_{2}} \wedge x_{j_{3}} \nonumber  \\
 &  + \Delta _{j_{2} j_{3}}A_{2 j_{1} j_{2}} A_{2 j_{3} j_{4}} x_{j_{1}} \wedge x_{j_{4}} - \Delta _{j_{2} j_{4}}A_{2 j_{1} j_{2}} A_{2 j_{3} j_{4}} x_{j_{1}} \wedge x_{j_{3}} \nonumber  \\
 &  +\left .A_{2 j_{1} j_{2}}  \Delta _{j_{1} j_{4}}A_{2 j_{3} j_{4}}  \Delta _{j_{2} j_{3}}I_{\mathcal{F}} -A_{2 j_{1} j_{2}}  \Delta _{j_{1} j_{3}}A_{2 j_{3} j_{4}}  \Delta _{j_{2} j_{4}}I_{\mathcal{F}}\right \}\text{.}\end{align}The "scalar" term is easily seen to be
\begin{equation} -\sum _{1 \leq j_{1} <j_{2} \leq 2 r}A_{2 j_{1} j_{2}}  \Delta _{j_{1} j_{1}}A_{2 j_{1} j_{2}} \Delta _{j_{2} j_{2}} = -\frac{1}{4} \sum _{1 \leq j_{1} <j_{2} \leq 2 r}\left (A_{2 j_{1} j_{2}}\right )^{2} I_{\mathcal{F}}\text{,}
\end{equation}while one can examine the bivector term by considering $r =2\text{,}$ leading to the terms
\begin{align} & A_{2 1 2} A_{2 1 3} x_{1} x_{2} x_{1} x_{3} +A_{2 1 2} A_{2 1 4} x_{1} x_{2} x_{1} x_{4} +A_{2 1 2} A_{2 2 3} x_{1} x_{2} x_{2} x_{3} \nonumber  \\
 &  +A_{2 1 2} A_{2 2 4} x_{1} x_{2} x_{2} x_{4} +A_{2 1 3} A_{2 1 2} x_{1} x_{3} x_{1} x_{2} +A_{2 1 3} A_{2 1 4} x_{1} x_{3} x_{1} x_{4} \nonumber  \\
 &  +A_{2 1 3} A_{2 2 3} x_{1} x_{3} x_{2} x_{3} +A_{2 1 3} A_{2 3 4} x_{1} x_{3} x_{3} x_{4} +A_{2 1 4} A_{2 1 2} x_{1} x_{4} x_{1} x_{2} \nonumber  \\
 &  +A_{2 1 4} A_{2 1 3} x_{1} x_{4} x_{1} x_{3} +A_{2 1 4} A_{2 1 4} x_{1} x_{4} x_{1} x_{4} +A_{2 1 4} A_{2 2 4} x_{1} x_{4} x_{2} x_{4} \nonumber  \\
 &  +A_{2 1 4} A_{2 3 4} x_{1} x_{4} x_{3} x_{4} +A_{2 2 3} A_{2 1 2} x_{2} x_{3} x_{1} x_{2} +A_{2 2 3} A_{2 1 3} x_{2} x_{3} x_{1} x_{3} \nonumber  \\
 &  +A_{2 2 3} A_{2 2 4} x_{2} x_{3} x_{2} x_{4} +A_{2 2 3} A_{2 3 4} x_{2} x_{3} x_{3} x_{4} +A_{2 2 4} A_{2 1 2} x_{2} x_{4} x_{1} x_{2} \nonumber  \\
 &  +A_{2 2 4} A_{2 1 3} x_{2} x_{4} x_{1} x_{3} +A_{2 2 4} A_{2 1 4} x_{2} x_{4} x_{1} x_{4} +A_{2 2 4} A_{2 2 3} x_{2} x_{4} x_{2} x_{3} \nonumber  \\
 &  +A_{2 2 4} A_{2 2 4} x_{2} x_{4} x_{2} x_{4} +A_{2 2 4} A_{2 3 4} x_{2} x_{4} x_{3} x_{4} +A_{2 3 4} A_{2 1 3} x_{3} x_{4} x_{1} x_{3} \nonumber  \\
 &  +A_{2 3 4} A_{2 1 4} x_{3} x_{4} x_{1} x_{4} +A_{2 3 4} A_{2 2 3} x_{3} x_{4} x_{2} x_{3} +A_{2 3 4} A_{2 2 4} x_{3} x_{4} x_{2} x_{4}\end{align}and see that the term involving e.g. $x_{2} \wedge x_{3}$ is given by
\begin{align} &  - \Delta _{1 1}A_{2 1 2} A_{2 1 3} x_{2} \wedge x_{3} - \Delta _{1 1}A_{2 1 3} A_{2 1 2} x_{3} \wedge x_{2} - \Delta _{4 4}A_{2 2 4} A_{2 3 4} x_{2} \wedge x_{3} \nonumber  \\
 &  - \Delta _{4 4}A_{2 3 4} A_{2 2 4} x_{3} \wedge x_{2}. =.0\text{,}\end{align}and the term involving e.g. $x_{1} \wedge x_{2}$ is given by
\begin{align} &  - \Delta _{3 3}A_{2 1 3} A_{2 2 3} x_{1} \wedge x_{2} - \Delta _{4 4}A_{2 1 4} A_{2 2 4} x_{1} \wedge x_{2} - \Delta _{3 3}A_{2 2 3} A_{2 1 3} x_{2} \wedge x_{1} \nonumber  \\
 &  - \Delta _{4 4}A_{2 2 4} A_{2 1 4} x_{2} \wedge x_{1}. =.0\text{,}\end{align}thus the entire two vector term must be zero. The four vector term in this case is 


\begin{align} & A_{2 1 2} A_{2 3 4} x_{1} x_{2} x_{3} x_{4} +A_{2 1 3} A_{2 2 4} x_{1} x_{3} x_{2} x_{4} +A_{2 1 4} A_{2 2 3} x_{1} x_{4} x_{2} x_{3} \nonumber  \\
 &  +A_{2 2 3} A_{2 1 4} x_{2} x_{3} x_{1} x_{4} +A_{2 2 4} A_{2 1 3} x_{2} x_{4} x_{1} x_{3} +A_{2 3 4} A_{2 1 2} x_{3} x_{4} x_{1} x_{2} \nonumber  \\
 &  =2 \left (A_{2 1 2} A_{2 3 4} -A_{2 1 3} A_{2 2 4} +A_{2 1 4} A_{2 2 3}\right ) x_{1} x_{2} x_{3} x_{4} \nonumber  \\
 &  =2 \ensuremath{\operatorname*{Pf}}\left \{\mathbb{A}_{2_{1 2 3 4}}\right \} x_{1} x_{2} x_{3} x_{4} =2 \ensuremath{\operatorname*{Pf}}\left \{\mathbb{A}_{2 1 2 3 4}\right \} x_{1} \wedge x_{2} \wedge x_{3} \wedge x_{4} =A_{2} \wedge A_{2}\text{.}\end{align}Hence
\begin{equation}\left (A_{2}\right )^{2} =\frac{1}{8} \ensuremath{\operatorname*{Tr}}\left \{\mathbb{A}_{2}^{t} \mathbb{A}_{2}\right \} I_{\mathcal{F}} +A_{2} \wedge A_{2}\text{.}
\end{equation}Thus the square of the operator $A_{0} I_{\mathcal{F}} +A_{2}$ is given by
\begin{equation}\left (A_{0} I_{\mathcal{F}} +A_{2}\right )^{2} =A_{0}^{2} I_{\mathcal{F}} +2 A_{0} A_{2} +A_{2}^{2}\text{.}
\end{equation}The idempotency condition for such operators is
\begin{align} & A_{0}^{2} I_{\mathcal{F}} +2 A_{2} +\frac{1}{8} \ensuremath{\operatorname*{Tr}}\left \{\mathbb{A}_{2}^{t} \mathbb{A}_{2}\right \} I_{\mathcal{F}} +A_{2} \wedge A_{2} -A_{0} I_{\mathcal{F}} -A_{2} \nonumber  \\
 &  =\left (A_{0}^{2} +\frac{1}{8} \ensuremath{\operatorname*{Tr}}\left \{\mathbb{A}_{2}^{t} \mathbb{A}_{2}\right \} -A_{0}\right ) I_{\mathcal{F}} +A_{2} \wedge A_{2} +\left (2 A_{0} -1\right ) A_{2} =0\text{,}\end{align}thus $A$ is idempotent $I f f$
\begin{align}\left (A_{0}^{2} +\frac{1}{8} \ensuremath{\operatorname*{Tr}}\left \{\mathbb{A}_{2}^{t} \mathbb{A}_{2}\right \} -A_{0}\right ) &  =\left (A_{0} \left (A_{0} -1\right ) +\frac{1}{4} \sum _{1 \leq j_{1} <j_{2} \leq 2 r}\left (A_{2 j_{1} j_{2}}\right )^{2}\right ) =0 \\
\left (2 A_{0} -1\right ) A_{2} &  =0 \\
A_{2} \wedge A_{2} &  =0.\end{align}This gives that
\begin{gather}A_{0} =\frac{1}{2}\text{and/or}A_{2} =0 \\
A_{2}\text{is a rank one two-vector or}A_{2} =0\text{,}\end{gather}which imply that
\begin{equation}\left (A_{0}^{2} +\frac{1}{8} \ensuremath{\operatorname*{Tr}}\left \{\mathbb{A}_{2}^{t} \mathbb{A}_{2}\right \} -A_{0}\right ) =\frac{1}{4} +\frac{2}{8} -\frac{1}{2} =0
\end{equation}or
\begin{equation}0 +0 -0 =0.
\end{equation}\emph{Thus for idempotent }second degree operator $A^{ \prime } s$
\begin{equation}A =\frac{1}{2} I_{\mathcal{F}} +a \wedge b\text{.}
\end{equation}

\subsubsection{Fourth Degree Gauge Invariant Operators}
We need to analyse the squares of fourth degree operators that consist of the sum of zero, $A_{0}$, two, $A_{2}\text{,}$ and four, $A_{4}\text{,}$ homogeneous parts
\begin{align}A_{0} &  =\alpha  I_{\mathcal{F}} \\
A_{2} &  =\sum _{1 \leq j_{1} <j_{2} \leq 2 r}A_{2 j_{1} j_{2}} x_{j_{1}} x_{j_{2}} \\
A_{4} &  =\sum _{1 \leq j_{1} <j_{2} <j_{3} <j_{4} \leq 2 r}A_{4 j_{1} j_{2} j_{3} j_{4}} x_{j_{1}} x_{j_{2}} x_{j_{3}} x_{j_{4}}\text{.}\end{align}

First we consider the product of the homogeneous fourth degree operator with the homogeneous second
degree operator
\begin{equation}A_{4} A_{2} =\sum _{\begin{array}{c}1 \leq j_{1} <j_{2} <j_{3} <j_{4} \leq 2 r \\
1 \leq j_{5} <j_{6} \leq 2 r\end{array}}A_{4 j_{1} j_{2} j_{3} j_{4}} A_{2 j_{5} j_{6}} x_{j_{1}} x_{j_{2}} x_{j_{3}} x_{j_{4}} x_{j_{5}} x_{j_{6}}
\end{equation}Utilising the relationship
\begin{align}x_{j_{1}} x_{j_{2}} x_{j_{3}} x_{j_{4}} x_{j_{5}} x_{j_{6}} &  =x_{j_{1}} \wedge x_{j_{2}} \wedge x_{j_{3}} \wedge x_{j_{4}} \wedge x_{j_{5}} \wedge x_{j_{6}} \nonumber  \\
 &  +\sum _{\sigma  \in S_{6}/\left (S_{2} \times S_{4}\right )}\pi  \left (\sigma \right )  \Delta _{j_{\sigma  \left (1\right )} j_{\sigma  \left (2\right )}}x_{j_{\sigma  \left (3\right )}} \wedge x_{j_{\sigma  \left (4\right )}} \wedge x_{j_{\sigma  \left (5\right )}} \wedge x_{j_{\sigma  \left (6\right )}} \nonumber  \\
 & \text{\quad \quad \quad \quad \quad \quad \quad \quad } +\sum _{\sigma  \in S_{6}/\left (S_{4} \times S_{2}\right )}\pi  \left (\sigma \right ) \ensuremath{\operatorname*{Pf}} \left (\mathbf{\Delta }_{j_{\sigma  \left (1\right )} j_{\sigma  \left (2\right )} j_{\sigma  \left (3\right )} j_{\sigma  \left (4\right )}}\right ) x_{j_{\sigma  \left (5\right )}} \wedge x_{j_{\sigma  \left (6\right )}} +\ensuremath{\operatorname*{Pf}} \left (\mathbf{\Delta }_{j_{1} j_{2} j_{3} j_{4} j_{5} j_{6}}\right ) I_{\mathcal{F}}\text{,}\end{align}one obtains
\begin{align}A_{4} A_{2} &  =\sum _{\begin{array}{c}1 \leq j_{1} <j_{2} <j_{3} <j_{4} \leq 2 r \\
1 \leq j_{5} <j_{6} \leq 2 r\end{array}}A_{4 j_{1} j_{2} j_{3} j_{4}} A_{2 j_{5} j_{6}}\left \{x_{j_{1}} \wedge x_{j_{2}} \wedge x_{j_{3}} \wedge x_{j_{4}} \wedge x_{j_{5}} \wedge x_{j_{6}}\right . \nonumber  \\
 &  +\sum _{\sigma  \in S_{6}/\left (S_{4} \times S_{2}\right )}\pi  \left (\sigma \right )  \Delta _{j_{\sigma  \left (1\right )} j_{\sigma  \left (2\right )}}x_{j_{\sigma  \left (3\right )}} \wedge x_{j_{\sigma  \left (4\right )}} \wedge x_{j_{\sigma  \left (5\right )}} \wedge x_{j_{\sigma  \left (6\right )}} \nonumber  \\
 &  +\sum _{\sigma  \in S_{6}/\left (S_{4} \times S_{4}\right )}\pi  \left (\sigma \right ) \ensuremath{\operatorname*{Pf}} \left (\mathbf{\Delta }_{j_{\sigma  \left (1\right )} j_{\sigma  \left (2\right )} j_{\sigma  \left (3\right )} j_{\sigma  \left (4\right )}}\right ) x_{j_{\sigma  \left (5\right )}} \wedge x_{j_{\sigma  \left (6\right )}} \nonumber  \\
 &  +\left .\ensuremath{\operatorname*{Pf}} \left (\mathbf{\Delta }_{j_{1} j_{2} j_{3} j_{4} j_{5} j_{6}}\right ) I_{\mathcal{F}}\right \}\text{.}\end{align}Note that
\begin{equation}\ensuremath{\operatorname*{Pf}} \left (\mathbf{\Delta }_{j_{1} j_{2} j_{3} j_{4}}\right ) = \Delta _{j_{1} j_{2}} \Delta _{j_{3} j_{4}} -  \Delta _{j_{1} j_{3}} \Delta _{j_{2} j_{4}} +  \Delta _{j_{1} j_{4}} \Delta _{j_{2} j_{3}}
\end{equation}and
\begin{align} & \ensuremath{\operatorname*{Pf}} \left (\mathbf{\Delta }\right )_{j_{1} j_{2} j_{3} j_{4} j_{5} j_{6}} =\ensuremath{\operatorname*{Pf}} \left (\mathbf{\Delta }_{j_{1} j_{2} j_{3} j_{4}}\right )  \Delta _{j_{5} j_{6}} -\ensuremath{\operatorname*{Pf}} \left (\mathbf{\Delta }_{j_{1} j_{2} j_{3} j_{5}}\right ) \Delta _{j_{4} j_{6}} \nonumber  \\
 &  +\ensuremath{\operatorname*{Pf}} \left (\mathbf{\Delta }_{j_{1} j_{2} j_{4} j_{5}}\right )  \Delta _{j_{3} j_{6}} -\ensuremath{\operatorname*{Pf}} \left (\mathbf{\Delta }_{j_{1} j_{3} j_{4} j_{5}}\right )  \Delta _{j_{2} j_{6}} +\ensuremath{\operatorname*{Pf}} \left (\mathbf{\Delta }_{j_{2} j_{3} j_{4} j_{5}}\right ) \Delta _{j_{1} j_{6}}\text{.}\end{align}Let $r =3$ then
\begin{align} & \sum _{\begin{array}{c}1 \leq j_{1} <j_{2} <j_{3} <j_{4} \leq 6 \\
1 \leq j_{5} <j_{6} \leq 6\end{array}}A_{4 j_{1} j_{2} j_{3} j_{4}} A_{2 j_{5} j_{6}}\left \{\sum _{\sigma  \in S_{6}/\left (S_{4} \times S_{2}\right )}\pi  \left (\sigma \right )  \Delta _{j_{\sigma  \left (1\right )} j_{\sigma  \left (2\right )}}x_{j_{\sigma  \left (3\right )}} \wedge x_{j_{\sigma  \left (4\right )}} \wedge x_{j_{\sigma  \left (5\right )}} \wedge x_{j_{\sigma  \left (6\right )}}\right . \nonumber  \\
 &  +\sum _{\sigma  \in S_{6}/\left (S_{4} \times S_{2}\right )}\pi  \left (\sigma \right )  \Delta _{j_{\sigma  \left (1\right )} j_{\sigma  \left (2\right )}}x_{j_{\sigma  \left (3\right )}} \wedge x_{j_{\sigma  \left (4\right )}} \wedge x_{j_{\sigma  \left (5\right )}} \wedge x_{j_{\sigma  \left (6\right )}} \nonumber  \\
 &  +\sum _{\sigma  \in S_{6}/\left (S_{4} \times S_{4}\right )}\pi  \left (\sigma \right ) \ensuremath{\operatorname*{Pf}} \left (\mathbf{\Delta }_{j_{\sigma  \left (1\right )} j_{\sigma  \left (2\right )} j_{\sigma  \left (3\right )} j_{\sigma  \left (4\right )}}\right ) x_{j_{\sigma  \left (5\right )}} \wedge x_{j_{\sigma  \left (6\right )}} \nonumber  \\
 &  +\left .\ensuremath{\operatorname*{Pf}} \left (\mathbf{\Delta }_{j_{1} j_{2} j_{3} j_{4} j_{5} j_{6}}\right ) I_{\mathcal{F}}\right \}\text{.}\end{align}Two indices in common gives the second order wedge product term
\begin{align} & 123412 \left (34\right ) ,123512 \left (35\right ) ,123612 \left (36\right ) ,124512 \left (45\right ) ,124612 \left (46\right ) ,125612 \left (56\right ) , \nonumber  \\
 & 123413 \left (24\right ) ,123513 \left (25\right ) ,123613 \left (26\right ) ,134513 \left (45\right ) ,134613 \left (46\right ) ,135613 \left (56\right ) , \nonumber  \\
 & 123414 \left (23\right ) ,124514 \left (25\right ) ,124614 \left (26\right ) ,134514 \left (35\right ) ,134614 \left (36\right ) ,145614 \left (56\right ) , \nonumber  \\
 & 123515 \left (23\right ) ,124515 \left (24\right ) ,125615 \left (26\right ) ,134515 \left (34\right ) ,135615 \left (36\right ) ,145615 \left (46\right ) , \nonumber  \\
 & 123616 \left (23\right ) ,124616 \left (24\right ) ,125616 \left (26\right ) ,134616 \left (34\right ) ,135616 \left (35\right ) ,145616 \left (45\right )\text{,}\end{align}
\begin{align*} & 123423 \left (14\right ) ,123523 \left (15\right ) ,123623 \left (16\right ) ,234523 \left (43\right ) ,234623 \left (46\right ) ,235623 \left (56\right )\text{,} \\
 & 123424 \left (13\right ) ,124524 \left (15\right ) ,124624 \left (16\right ) ,234524 \left (35\right ) ,234624 \left (36\right ) ,245624 \left (56\right )\text{,} \\
 & 123525 \left (13\right ) ,124525 \left (14\right ) ,125625 \left (16\right ) ,234525 \left (34\right ) ,235625 \left (36\right ) ,245625 \left (46\right )\text{,} \\
 & 123626 \left (13\right ) ,124626 \left (14\right ) ,125626 \left (15\right ) ,234626 \left (34\right ) ,235626 \left (35\right ) ,245626 \left (45\right )\text{,}\end{align*}
\begin{align*} & 123434 \left (12\right ) ,134534 \left (15\right ) ,134634 \left (16\right ) ,234534 \left (25\right ) ,234634 \left (26\right ) ,345634 \left (56\right )\text{,} \\
 & 123535 \left (12\right ) ,143535 \left (14\right ) ,135635 \left (16\right ) ,234535 \left (24\right ) ,235635 \left (26\right ) ,345635 \left (46\right )\text{,} \\
 & 123636 \left (12\right ) ,134636 \left (16\right ) ,135636 \left (15\right ) ,234636 \left (24\right ) ,235636 \left (25\right ) ,345636 \left (45\right )\text{,}\end{align*}
\begin{align*} & 124545 \left (12\right ) ,145645 \left (16\right ) ,134545 \left (13\right ) ,234545 \left (23\right ) ,345645 \left (34\right ) ,245645 \left (26\right )\text{,} \\
 & 145646 \left (15\right ) ,124646 \left (12\right ) ,134646 \left (13\right ) ,234646 \left (23\right ) ,245646 \left (25\right ) ,345646 \left (35\right )\text{,}\end{align*}
\begin{equation}125656 \left (12\right ) ,135656 \left (13\right ) ,145656 \left (14\right ) ,235656 \left (23\right ) ,245656 \left (24\right ) ,345656 \left (34\right )\text{.}
\end{equation}Hence the coefficient of $x_{1} \wedge x_{2}$ is
\begin{equation} -\frac{1}{4} \left (A_{1 2 5 6} A_{5 6} +A_{1 2 4 6} A_{4 6} +A_{1 2 3 6} A_{3 6} +A_{1 2 4 5} A_{4 5} +A_{1 2 3 5} A_{3 5} +A_{1 2 3 4} A_{3 4}\right )
\end{equation}

One index in common gives the fourth order wedge product term
\begin{align*} & 134512 \left (3452\right ) ,134612 \left (3462\right ) ,135612 \left (3562\right ) ,145612 \left (4562\right ) ,234512 \left (3451\right ) ,234612 \left (3461\right ) ,235612 \left (3561\right )\text{,} \\
 & 245612 \left (4561\right )\text{,} \\
 & 124513 \left (2453\right ) ,124613 \left (2463\right ) ,135613 \left (3561\right ) ,145613 \left (4561\right ) ,234513 \left (2451\right ) ,234613 \left (2461\right ) ,235613 \left (2561\right )\text{,} \\
 & 345613 \left (4561\right )\text{,} \\
 & 123514 \left (2354\right ) ,123614 \left (2364\right ) ,145614 \left (3561\right ) ,135614 \left (3561\right ) ,234514 \left (2351\right ) ,234614 \left (2361\right ) ,245614 \left (2561\right )\text{,} \\
 & 345614 \left (3561\right )\text{,} \\
 & 123415 \left (2345\right ) ,123615 \left (2365\right ) ,145615 \left (3461\right ) ,134615 \left (3461\right ) ,234515 \left (2341\right ) \ast  , 235615 \left (2361\right ) ,245615 \left (2461\right )\text{,} \\
 & 345615 \left (3461\right )\text{,} \\
 & 123416 \left (2346\right ) ,123516 \left (2356\right ) ,145616 \left (3451\right ) ,134516 \left (3451\right ) ,234616 \left (2341\right ) \ast  , 235616 \left (2351\right ) ,245616 \left (2451\right )\text{,} \\
 & 345616 \left (3451\right )\text{,}\end{align*}
\begin{align*} & 124523 \left (1453\right ) ,125623 \left (1563\right ) ,124623 \left (1463\right ) ,245623 \left (4563\right ) ,134523 \left (1453\right ) ,134623 \left (1462\right ) ,135623 \left (1562\right )\text{,} \\
 & 345623 \left (4562\right )\text{,} \\
 & 123524 \left (1354\right ) ,123624 \left (1364\right ) ,125624 \left (1564\right ) ,235624 \left (3564\right ) ,134524 \left (1453\right ) ,134624 \left (1462\right ) ,145624 \left (1562\right )\text{,} \\
 & 345624 \left (4562\right )\text{,} \\
 & 123425 \left (1345\right ) ,123625 \left (1365\right ) ,124625 \left (1465\right ) ,234625 \left (3465\right ) ,135425 \left (1342\right ) ,135625 \left (1562\right ) ,145625 \left (1462\right )\text{,} \\
 & 354625 \left (3462\right )\text{,} \\
 & 123426 \left (1346\right ) ,123526 \left (1356\right ) ,124526 \left (1456\right ) ,234526 \left (3456\right ) ,136426 \left (1342\right ) ,136526 \left (1352\right ) ,146526 \left (1452\right )\text{,} \\
 & 364526 \left (3452\right )\text{,}\end{align*}
\begin{align*} & 123534 \left (1254\right ) ,123634 \left (1264\right ) ,235634 \left (2564\right ) ,135634 \left (1564\right ) ,124534 \left (1253\right ) ,124634 \left (1263\right ) ,145634 \left (1563\right )\text{,} \\
 & 245634 \left (2563\right )\text{,} \\
 & 123435 \left (1245\right ) ,123635 \left (1265\right ) ,234635 \left (2465\right ) ,134635 \left (1465\right ) ,124535 \left (1243\right ) \ast  , 125635 \left (1263\right ) ,145635 \left (1463\right )\text{,} \\
 & 245635 \left (2463\right )\text{,} \\
 & 123536 \left (1256\right ) ,123436 \left (1246\right ) ,234536 \left (2456\right ) ,134536 \left (1436\right ) ,124636 \left (1243\right ) \ast  , 125636 \left (1253\right ) ,145636 \left (1453\right )\text{,} \\
 & 245636 \left (2463\right )\text{,}\end{align*}
\begin{align*} & 123545 \left (1234\right ) \ast  , 125645 \left (1264\right ) ,235645 \left (2364\right ) ,135645 \left (1364\right ) ,124645 \left (1265\right ) ,134645 \left (1365\right ) ,234645 \left (2365\right )\text{,} \\
 & 123445 \left (1235\right )\text{,} \\
 & 123446 \left (1236\right ) ,124546 \left (1256\right ) ,234546 \left (2356\right ) ,134546 \left (1356\right ) ,123646 \left (1234\right ) \ast  , 125646 \left (1254\right ) ,135646 \left (1354\right )\text{,} \\
 & 235646 \left (2354\right )\text{,}\end{align*}
\begin{align} & 123556 \left (1236\right ) ,124556 \left (1246\right ) ,234556 \left (2346\right ) ,134556 \left (1346\right ) ,124566 \left (1245\right ) ,134656 \left (1345\right ) ,124656 \left (1245\right ) , \nonumber  \\
 & 234656 \left (2345\right )\text{,}\end{align}Hence the coefficient of $x_{1} \wedge x_{2}$ $ \wedge x_{3} \wedge x_{4}$ is
\begin{equation} -\frac{1}{2} \left (A_{1 2 3 6} A_{4 6} +A_{1 2 3 5} A_{4 5} -A_{1 2 4 6} A_{3 6} -A_{1 2 4 5} A_{3 5} -A_{2 3 4 5} A_{1 5} -A_{2 3 4 6} A_{1 6} +A_{1 3 4 5} A_{2 5} +A_{1 3 4 6} A_{2 6}\right )
\end{equation}No indices in common give sixth order wedge product terms
\begin{align} & 345612,245613,235614,234615,234516,145623,135624,134625 , \nonumber  \\
 & 134526,125634,124635,124536,123645,123546,123456.\end{align}Hence the coefficient of $x_{1} \wedge x_{2}$ $ \wedge x_{3} \wedge x_{4}$ $ \wedge x_{5} \wedge x_{6}$ is
\begin{align} & A_{3 4 5 6} A_{1 2} -A_{2 4 5 6} A_{1 3} +A_{2 3 5 6} A_{1 4} -A_{2 3 4 6} A_{1 5} +A_{2 3 4 5} A_{1 6} +A_{1 4 5 6} A_{2 3} -A_{1 3 5 6} A_{2 4} +A_{1 3 4 6} A_{2 5} \nonumber  \\
 &  -A_{1 3 4 5} A_{2 6} +A_{1 2 5 6} A_{3 4} -A_{1 2 4 6} A_{3 5} +A_{1 2 4 5} A_{3 6} +A_{1 2 3 6} A_{4 5} -A_{1 2 3 5} A_{4 6} +A_{1 2 3 4} A_{5 6}\text{,}\end{align}and the sixth order term is
\begin{equation}A_{4} \wedge A_{2} =\left (A_{4} \wedge A_{2}\right )_{1 2 3 4 5 6} x_{1} \wedge x_{2} \wedge x_{3} \wedge x_{4} \wedge x_{5} \wedge x_{6}\text{,}
\end{equation}where $\left (A_{4} \wedge A_{2}\right )_{1 2 3 4 5 6}$ is given by the preceding expression. 

Next we consider the homogeneous fourth degree operator
\begin{equation}A_{4} =\sum _{1 \leq j_{1} <j_{2} <j_{3} <j_{4} \leq 2 r}A_{4 j_{1} j_{2} j_{3} j_{4}} x_{j_{1}} x_{j_{2}} x_{j_{3}} x_{j_{4}}
\end{equation}and its square
\begin{equation}\left (A_{4}\right )^{2} =\sum _{\begin{array}{c}1 \leq j_{1} <j_{2} \leq 2 r \\
1 \leq j_{3} <j_{4} \leq 2 r\end{array}}A_{4 j_{1} j_{2} j_{3} j_{4}} A_{4 j_{5} j_{6} j_{7} j_{8}} x_{j_{1}} x_{j_{2}} x_{j_{3}} x_{j_{4}} x_{j_{5}} x_{j_{6}} x_{j_{7}} x_{j_{8}}\text{.}
\end{equation}Utilizing the relationship
\begin{align} & x_{j_{1}} x_{j_{2}} x_{j_{3}} x_{j_{4}} x_{j_{5}} x_{j_{6}} x_{j_{7}} x_{j_{8}} =x_{j_{1}} \wedge x_{j_{2}} \wedge x_{j_{3}} \wedge x_{j_{4}} \wedge x_{j_{5}} \wedge x_{j_{6}} \wedge x_{j_{7}} \wedge x_{j_{8}} \nonumber  \\
 &  +\sum _{\sigma  \in S_{8}/\left (S_{2} \times S_{6}\right )}\pi  \left (\sigma \right ) \ensuremath{\operatorname*{Pf}} \left (\mathbf{\Delta }_{j_{\sigma  \left (1\right )} j_{\sigma  \left (2\right )}}\right ) x_{j_{\sigma  \left (3\right )}} \wedge x_{j_{\sigma  \left (4\right )}} \wedge x_{j_{\sigma  \left (5\right )}} \wedge x_{j_{\sigma  \left (6\right )}} \wedge x_{j_{\sigma  \left (7\right )}} \wedge x_{j_{\sigma  \left (8\right )}} \nonumber  \\
 &  +\sum _{\sigma  \in S_{8}/\left (S_{4} \times S_{4}\right )}\pi  \left (\sigma \right ) \ensuremath{\operatorname*{Pf}} \left (\mathbf{\Delta }_{j_{\sigma  \left (1\right )} j_{\sigma  \left (2\right )} j_{\sigma  \left (3\right )} j_{\sigma  \left (4\right )}}\right ) x_{j_{\sigma  \left (5\right )}} \wedge x_{j_{\sigma  \left (6\right )}} \wedge x_{j_{\sigma  \left (7\right )}} \wedge x_{j_{\sigma  \left (8\right )}} \nonumber  \\
 &  +\sum _{\sigma  \in S_{8}/\left (S_{6} \times S_{2}\right )}\pi  \left (\sigma \right ) \ensuremath{\operatorname*{Pf}} \left (\mathbf{\Delta }_{j_{\sigma  \left (1\right )} j_{\sigma  \left (2\right )} j_{\sigma  \left (3\right )} j_{\sigma  \left (4\right )} j_{\sigma  \left (5\right )} j_{\sigma  \left (6\right )}}\right ) x_{j_{\sigma  \left (7\right )}} \wedge x_{j_{\sigma  \left (8\right )}} +\ensuremath{\operatorname*{Pf}} \left (\mathbf{\Delta }_{j_{1} j_{2} j_{3} j_{4} j_{5} j_{6} j_{7} j_{8}}\right ) I_{\mathcal{F}}\end{align}one obtains
\begin{align}\left (A_{4}\right )^{2} &  =\sum _{\begin{array}{c}1 \leq j_{1} <j_{2} <j_{3} <j_{4} \leq 2 r \\
1 \leq j_{5} <j_{6}\, <j_{7} <j_{8} \leq 2 r\end{array}}A_{4 j_{1} j_{2} j_{3} j_{4}} A_{4 j_{5} j_{6} j_{7} j_{8}}\left \{x_{j_{1}} \wedge x_{j_{2}} \wedge x_{j_{3}} \wedge x_{j_{4}} \wedge x_{j_{5}} \wedge x_{j_{6}} \wedge x_{j_{7}} \wedge x_{j_{8}}\right . \\
 &  +\sum _{\sigma  \in S_{8}/\left (S_{2} \times S_{6}\right )}\pi  \left (\sigma \right )  \Delta _{j_{\sigma  \left (1\right )} j_{\sigma  \left (2\right )}}x_{j_{\sigma  \left (3\right )}} \wedge x_{j_{\sigma  \left (4\right )}} \wedge x_{j_{\sigma  \left (5\right )}} \wedge x_{j_{\sigma  \left (6\right )}} \wedge x_{j_{\sigma  \left (7\right )}} \wedge x_{j_{\sigma  \left (8\right )}} \\
 &  +\sum _{\sigma  \in S_{8}/\left (S_{4} \times S_{4}\right )}\pi  \left (\sigma \right ) \ensuremath{\operatorname*{Pf}} \left (\mathbf{\Delta }_{j_{\sigma  \left (1\right )} j_{\sigma  \left (2\right )} j_{\sigma  \left (3\right )} j_{\sigma  \left (4\right )}}\right ) x_{j_{\sigma  \left (5\right )}} \wedge x_{j_{\sigma  \left (6\right )}} \wedge x_{j_{\sigma  \left (7\right )}} \wedge x_{j_{\sigma  \left (8\right )}} \\
 &  +\left .\sum _{\sigma  \in S_{8}/\left (S_{6} \times S_{2}\right )}\pi  \left (\sigma \right ) \ensuremath{\operatorname*{Pf}} \left (\mathbf{\Delta }_{j_{\sigma  \left (1\right )} j_{\sigma  \left (2\right )} j_{\sigma  \left (3\right )} j_{\sigma  \left (4\right )} j_{\sigma  \left (5\right )} j_{\sigma  \left (6\right )}}\right ) x_{j_{\sigma  \left (7\right )}} \wedge x_{j_{\sigma  \left (8\right )}} +\ensuremath{\operatorname*{Pf}} \left (\mathbf{\Delta }_{j_{1} j_{2} j_{3} j_{4} j_{5} j_{6} j_{7} j_{8}}\right ) I_{\mathcal{F}}\right \}\text{.}\end{align}One can consider the special case $r =4$ the number of subscripts that have no values in common are
\begin{equation}\binom{8}{4} \binom{8 -4}{4} =\binom{8}{4} =\frac{8 \ast 7 \ast 6 \ast 5}{1 \ast 2 \ast 3 \ast 4} =70
\end{equation}e.g.
\begin{equation}12345678,12354678,12364578,12374568,12384567 ,\cdots \text{,}
\end{equation}which gives the terms
\begin{align} & A_{4 1 2 3 4} A_{4 5 6 7 8} x_{1} \wedge x_{2} \wedge x_{3} \wedge x_{4} \wedge x_{5} \wedge x_{6} \wedge x_{7} \wedge x_{8} \nonumber  \\
 & A_{4 1 2 3 5} A_{4 4 6 7 8} x_{1} \wedge x_{2} \wedge x_{3} \wedge x_{5} \wedge x_{4} \wedge x_{6} \wedge x_{7} \wedge x_{8} \nonumber  \\
 & A_{4 1 2 3 6} A_{4 4 5 7 8} x_{1} \wedge x_{2} \wedge x_{3} \wedge x_{6} \wedge x_{4} \wedge x_{5} \wedge x_{7} \wedge x_{8} \nonumber  \\
 & A_{4 1 2 3 7} A_{4 4 5 6 8} x_{1} \wedge x_{2} \wedge x_{3} \wedge x_{7} \wedge x_{4} \wedge x_{5} \wedge x_{6} \wedge x_{8} \nonumber  \\
 & A_{4 1 2 3 8} A_{4 4 5 6 7} x_{1} \wedge x_{2} \wedge x_{3} \wedge x_{8} \wedge x_{4} \wedge x_{5} \wedge x_{6} \wedge x_{7}\end{align}and the final coefficient is
\begin{equation}\left (A_{4 1 2 3 4} A_{4 5 6 7 8} -A_{4 1 2 3 5} A_{4 4 6 7 8} +A_{4 1 2 3 6} A_{4 4 5 7 8} -A_{4 1 2 3 7} A_{4 4 5 6 8} +A_{4 1 2 3 8} A_{4 4 5 6 7} +\cdots \right ) x_{1} \wedge x_{2} \wedge x_{3} \wedge x_{4} \wedge x_{5} \wedge x_{6} \wedge x_{7} \wedge x_{8}\text{.}
\end{equation}


\begin{enumerate}
\item
\begin{align}A_{1 2 5 6} A_{3 4 7 8} x_{1} x_{2} x_{5} x_{6} x_{3} x_{4} x_{7} x_{8} &  =A_{1 2 5 6} A_{3 4 7 8} x_{1} x_{2} x_{3} x_{4} x_{5} x_{6} x_{7} x_{8} \\
 &  =A_{1 \left (23\right )} A_{\left (14\right ) \left (34\right )} y_{1 2} y_{3 4} y_{5 6} y_{7 8} =A_{1 \left (23\right )} A_{\left (14\right ) \left (34\right )} \mathcal{Y}_{1} \mathcal{Y}_{\left (14\right )} \mathcal{Y}_{\left (23\right )} \mathcal{Y}_{\left (34\right )} \\
A_{1 3 5 6} A_{2 4 7 8} x_{1} x_{3} x_{5} x_{6} x_{2} x_{4} x_{7} x_{8} &  = -A_{1 3 5 6} A_{2 4 7 8} x_{1} x_{2} x_{3} x_{4} x_{5} x_{6} x_{7} x_{8} \\
A_{1 4 5 6} A_{2 3 7 8} x_{1} x_{4} x_{5} x_{6} x_{2} x_{3} x_{7} x_{8} &  =A_{1 4 5 6} A_{2 3 7 8} x_{1} x_{2} x_{3} x_{4} x_{5} x_{6} x_{7} x_{8} \\
A_{2 3 5 6} A_{1 4 7 8} x_{2} x_{3} x_{5} x_{6} x_{1} x_{4} x_{7} x_{8} &  = -A_{2 3 5 6} A_{1 4 7 8} x_{1} x_{2} x_{3} x_{4} x_{5} x_{6} x_{7} x_{8} \\
A_{2 4 5 6} A_{1 3 7 8} x_{2} x_{4} x_{5} x_{6} x_{1} x_{3} x_{7} x_{8} &  = -A_{2 4 5 6} A_{1 3 7 8} x_{1} x_{2} x_{3} x_{4} x_{5} x_{6} x_{7} x_{8} \\
A_{3 4 5 6} A_{1 2 7 8} x_{3} x_{4} x_{5} x_{6} x_{1} x_{2} x_{7} x_{8} &  =A_{3 4 5 6} A_{1 2 7 8} x_{1} x_{2} x_{3} x_{4} x_{5} x_{6} x_{7} x_{8}\end{align}


\begin{align}A_{1 2 5 7} A_{3 4 6 8} x_{1} x_{2} x_{5} x_{7} x_{3} x_{4} x_{6} x_{8} &  = -A_{1 2 5 7} A_{3 4 6 8} x_{1} x_{2} x_{3} x_{5} x_{4} x_{6} x_{7} x_{8} \\
A_{1 3 5 7} A_{2 4 6 8} x_{1} x_{3} x_{5} x_{7} x_{2} x_{4} x_{6} x_{8} &  =A_{1 3 5 7} A_{2 4 6 8} x_{1} x_{2} x_{3} x_{5} x_{4} x_{6} x_{7} x_{8} \\
A_{1 4 5 7} A_{2 3 6 8} x_{1} x_{4} x_{5} x_{7} x_{2} x_{3} x_{6} x_{8} &  = -A_{1 4 5 7} A_{2 3 6 8} x_{1} x_{2} x_{3} x_{5} x_{4} x_{6} x_{7} x_{8} \\
A_{2 3 5 7} A_{1 4 6 8} x_{2} x_{3} x_{5} x_{7} x_{1} x_{4} x_{6} x_{8} &  =A_{2 3 5 7} A_{1 4 6 8} x_{1} x_{2} x_{3} x_{5} x_{4} x_{6} x_{7} x_{8} \\
A_{2 4 5 7} A_{1 3 6 8} x_{2} x_{4} x_{5} x_{7} x_{1} x_{3} x_{6} x_{8} &  =A_{2 4 5 7} A_{1 3 6 8} x_{1} x_{2} x_{3} x_{5} x_{4} x_{6} x_{7} x_{8} \\
A_{3 4 5 7} A_{1 2 6 8} x_{3} x_{4} x_{5} x_{7} x_{1} x_{2} x_{6} x_{8} &  =A_{3 4 5 7} A_{1 2 6 8} x_{1} x_{2} x_{3} x_{5} x_{4} x_{6} x_{7} x_{8}\end{align}


\begin{align}A_{1 2 5 6} A_{3 4 7 8} x_{1} x_{2} x_{5} x_{6} x_{3} x_{4} x_{7} x_{8} &  = \\
A_{1 3 5 6} A_{2 4 7 8} x_{1} x_{3} x_{5} x_{6} x_{2} x_{4} x_{7} x_{8} &  = \\
A_{1 4 5 6} A_{2 3 7 8} x_{1} x_{4} x_{5} x_{6} x_{2} x_{3} x_{7} x_{8} &  = \\
A_{2 3 5 6} A_{1 4 7 8} x_{2} x_{3} x_{5} x_{6} x_{1} x_{4} x_{7} x_{8} &  = \\
A_{2 4 5 6} A_{1 3 7 8} x_{2} x_{4} x_{5} x_{6} x_{1} x_{3} x_{7} x_{8} &  = \\
A_{3 4 5 6} A_{1 2 7 8} x_{3} x_{4} x_{5} x_{6} x_{1} x_{2} x_{7} x_{8} &  =\end{align}


\begin{align}A_{1 2 5 6} A_{3 4 7 8} x_{1} x_{2} x_{5} x_{6} x_{3} x_{4} x_{7} x_{8} &  = \\
A_{1 3 5 6} A_{2 4 7 8} x_{1} x_{3} x_{5} x_{6} x_{2} x_{4} x_{7} x_{8} &  = \\
A_{1 4 5 6} A_{2 3 7 8} x_{1} x_{4} x_{5} x_{6} x_{2} x_{3} x_{7} x_{8} &  = \\
A_{2 3 5 6} A_{1 4 7 8} x_{2} x_{3} x_{5} x_{6} x_{1} x_{4} x_{7} x_{8} &  = \\
A_{2 4 5 6} A_{1 3 7 8} x_{2} x_{4} x_{5} x_{6} x_{1} x_{3} x_{7} x_{8} &  = \\
A_{3 4 5 6} A_{1 2 7 8} x_{3} x_{4} x_{5} x_{6} x_{1} x_{2} x_{7} x_{8} &  =\end{align}


\begin{align}A_{1 2 5 6} A_{3 4 7 8} x_{1} x_{2} x_{5} x_{6} x_{3} x_{4} x_{7} x_{8} &  = \\
A_{1 3 5 6} A_{2 4 7 8} x_{1} x_{3} x_{5} x_{6} x_{2} x_{4} x_{7} x_{8} &  = \\
A_{1 4 5 6} A_{2 3 7 8} x_{1} x_{4} x_{5} x_{6} x_{2} x_{3} x_{7} x_{8} &  = \\
A_{2 3 5 6} A_{1 4 7 8} x_{2} x_{3} x_{5} x_{6} x_{1} x_{4} x_{7} x_{8} &  = \\
A_{2 4 5 6} A_{1 3 7 8} x_{2} x_{4} x_{5} x_{6} x_{1} x_{3} x_{7} x_{8} &  = \\
A_{3 4 5 6} A_{1 2 7 8} x_{3} x_{4} x_{5} x_{6} x_{1} x_{2} x_{7} x_{8} &  =\end{align}


\begin{align}A_{1 2 5 6} A_{3 4 7 8} x_{1} x_{2} x_{5} x_{6} x_{3} x_{4} x_{7} x_{8} &  = \\
A_{1 3 5 6} A_{2 4 7 8} x_{1} x_{3} x_{5} x_{6} x_{2} x_{4} x_{7} x_{8} &  = \\
A_{1 4 5 6} A_{2 3 7 8} x_{1} x_{4} x_{5} x_{6} x_{2} x_{3} x_{7} x_{8} &  = \\
A_{2 3 5 6} A_{1 4 7 8} x_{2} x_{3} x_{5} x_{6} x_{1} x_{4} x_{7} x_{8} &  = \\
A_{2 4 5 6} A_{1 3 7 8} x_{2} x_{4} x_{5} x_{6} x_{1} x_{3} x_{7} x_{8} &  = \\
A_{3 4 5 6} A_{1 2 7 8} x_{3} x_{4} x_{5} x_{6} x_{1} x_{2} x_{7} x_{8} &  =\end{align}

\item $3 -$substitions:
\begin{equation}\binom{4}{3}^{2} =16
\end{equation}
\begin{equation}\begin{array}{c}4678,5478,5648,5674\text{,} \\
3678,5378,5638,5673\text{,} \\
2678,5278,5628,5672\text{,} \\
1678,5178,5618,5671\text{,}\end{array} \begin{array}{c}1235,1236,1237,1238\text{,} \\
1254,1264,1274,1284\text{,} \\
1534,1634,1734,1834\text{,} \\
5234,6234,7234,8234\text{,}\end{array}
\end{equation}
\begin{align}A_{1 2 3 5} A_{4 6 7 8} x_{1} x_{2} x_{3} x_{5} x_{4} x_{6} x_{7} x_{8} &  = -A_{1 2 3 5} A_{4 6 7 8} x_{1} x_{2} x_{3} x_{4} x_{5} x_{6} x_{7} x_{8} \\
A_{1 2 3 6} A_{5 4 7 8} x_{1} x_{2} x_{3} x_{6} x_{5} x_{4} x_{7} x_{8} &  = -A_{1 2 3 6} A_{4 5 7 8} x_{1} x_{2} x_{3} x_{6} x_{5} x_{4} x_{7} x_{8} =A_{1 2 3 6} A_{4 5 7 8} x_{1} x_{2} x_{3} x_{4} x_{5} x_{6} x_{7} x_{8} \\
A_{1 2 3 7} A_{5 6 4 8} x_{1} x_{2} x_{3} x_{7} x_{5} x_{6} x_{4} x_{8} &  = -A_{1 2 3 7} A_{5 6 4 8} x_{1} x_{2} x_{3} x_{4} x_{5} x_{6} x_{7} x_{8} \\
A_{1 2 3 8} A_{5 6 7 4} x_{1} x_{2} x_{3} x_{8} x_{5} x_{6} x_{7} x_{4} &  = -A_{1 2 3 8} A_{5 6 7 4} x_{1} x_{2} x_{3} x_{5} x_{4} x_{6} x_{7} x_{8} =A_{1 2 3 8} A_{5 6 4 7} x_{1} x_{2} x_{3} x_{4} x_{5} x_{6} x_{7} x_{8}\end{align}


\begin{align}A_{1 2 5 4} A_{3 6 7 8} x_{1} x_{2} x_{3} x_{5} x_{4} x_{6} x_{7} x_{8} &  = -A_{1 2 3 5} A_{4 6 7 8} x_{1} x_{2} x_{3} x_{4} x_{5} x_{6} x_{7} x_{8} \\
A_{1 2 6 4} A_{5 3 7 8} x_{1} x_{2} x_{3} x_{6} x_{5} x_{4} x_{7} x_{8} &  = -A_{1 2 3 6} A_{4 5 7 8} x_{1} x_{2} x_{3} x_{6} x_{5} x_{4} x_{7} x_{8} =A_{1 2 3 6} A_{4 5 7 8} x_{1} x_{2} x_{3} x_{4} x_{5} x_{6} x_{7} x_{8} \\
A_{1 2 7 4} A_{5 6 3 8} x_{1} x_{2} x_{3} x_{7} x_{5} x_{6} x_{4} x_{8} &  = -A_{1 2 3 7} A_{5 6 4 8} x_{1} x_{2} x_{3} x_{4} x_{5} x_{6} x_{7} x_{8} \\
A_{1 2 8 4} A_{5 6 7 3} x_{1} x_{2} x_{3} x_{8} x_{5} x_{6} x_{7} x_{4} &  = -A_{1 2 3 8} A_{5 6 7 4} x_{1} x_{2} x_{3} x_{5} x_{4} x_{6} x_{7} x_{8} =A_{1 2 3 8} A_{5 6 4 7} x_{1} x_{2} x_{3} x_{4} x_{5} x_{6} x_{7} x_{8}\end{align}


\begin{align}A_{1 5 3 4} A_{2 6 7 8} x_{1} x_{2} x_{3} x_{5} x_{4} x_{6} x_{7} x_{8} &  = -A_{1 2 3 5} A_{4 6 7 8} x_{1} x_{2} x_{3} x_{4} x_{5} x_{6} x_{7} x_{8} \\
A_{1 6 3 4} A_{5 2 7 8} x_{1} x_{2} x_{3} x_{6} x_{5} x_{4} x_{7} x_{8} &  = -A_{1 2 3 6} A_{4 5 7 8} x_{1} x_{2} x_{3} x_{6} x_{5} x_{4} x_{7} x_{8} =A_{1 2 3 6} A_{4 5 7 8} x_{1} x_{2} x_{3} x_{5} x_{4} x_{6} x_{7} x_{8} \\
A_{1 7 3 4} A_{5 6 2 8} x_{1} x_{2} x_{3} x_{7} x_{5} x_{6} x_{4} x_{8} &  = -A_{1 2 3 7} A_{5 6 4 8} x_{1} x_{2} x_{3} x_{4} x_{5} x_{6} x_{7} x_{8} \\
A_{1 8 3 4} A_{5 6 7 2} x_{1} x_{2} x_{3} x_{8} x_{5} x_{6} x_{7} x_{4} &  = -A_{1 2 3 8} A_{5 6 7 4} x_{1} x_{2} x_{3} x_{5} x_{4} x_{6} x_{7} x_{8} =A_{1 2 3 8} A_{5 6 4 7} x_{1} x_{2} x_{3} x_{5} x_{4} x_{6} x_{7} x_{8}\end{align}


\begin{align}A_{5 2 3 4} A_{1 6 7 8} x_{1} x_{2} x_{3} x_{5} x_{4} x_{6} x_{7} x_{8} &  = -A_{1 2 3 5} A_{4 6 7 8} x_{1} x_{2} x_{3} x_{5} x_{4} x_{6} x_{7} x_{8} \\
A_{6 2 3 4} A_{5 1 7 8} x_{1} x_{2} x_{3} x_{6} x_{5} x_{4} x_{7} x_{8} &  = -A_{1 2 3 6} A_{4 5 7 8} x_{1} x_{2} x_{3} x_{6} x_{5} x_{4} x_{7} x_{8} =A_{1 2 3 6} A_{4 5 7 8} x_{1} x_{2} x_{3} x_{5} x_{4} x_{6} x_{7} x_{8} \\
A_{7 2 3 4} A_{5 6 1 8} x_{1} x_{2} x_{3} x_{7} x_{5} x_{6} x_{4} x_{8} &  = -A_{1 2 3 7} A_{5 6 4 8} x_{1} x_{2} x_{3} x_{5} x_{4} x_{6} x_{7} x_{8} \\
A_{8 2 3 4} A_{5 6 7 1} x_{1} x_{2} x_{3} x_{8} x_{5} x_{6} x_{7} x_{4} &  = -A_{1 2 3 8} A_{5 6 7 4} x_{1} x_{2} x_{3} x_{5} x_{4} x_{6} x_{7} x_{8} =A_{1 2 3 8} A_{5 6 4 7} x_{1} x_{2} x_{3} x_{5} x_{4} x_{6} x_{7} x_{8}\end{align}

\item $4 -$subsitutions:
\begin{equation}\binom{4}{4}^{2} =1
\end{equation}
\begin{equation}56781234
\end{equation}\end{enumerate}


One can consider the special case $r =4$ the number of subscripts that have one values in common are
\begin{equation}\binom{8}{4} \binom{8 -4}{3} =70 \ast 4 =280
\end{equation}e.g.
\begin{align} & 12341567,12351467,12361457,12371456,12451367,12461357,12471356,12561347,12571346,12671356 , \nonumber  \\
 & 13451267,13461257,13471256,13561247,13571246,13671245 , \nonumber  \\
 & 14561237,14571236,14671235 , \nonumber  \\
 & 15671234.\end{align}the coefficients are
\begin{align} & \left (A_{1 2 3 4} A_{1 5 6 7} -A_{1 2 3 5} A_{1 4 6 7} +A_{1 2 3 6} A_{1 4 5 7} -A_{1 2 3 7} A_{1 4 5 6} +A_{1 2 4 5} A_{1 3 6 7} -A_{1 2 4 6} A_{1 3 5 7} +A_{1 2 4 7} A_{1 3 5 6} -A_{1 2 5 6} A_{1 3 4 7}\right . \nonumber  \\
 &  -A_{1 2 5 7} A_{1 3 4 6} +A_{1 2 6 7} A_{1 3 4 5} -A_{1 3 4 5} A_{1 2 6 7} +A_{1 3 4 6} A_{1 2 5 7} -A_{1 3 4 7} A_{1 2 5 6} -A_{1 3 5 6} A_{1 2 4 7} +A_{1 3 5 7} A_{1 2 4 6} -A_{1 3 6 7} A_{1 2 4 5} \nonumber  \\
 & \left . +A_{1 4 5 6} A_{1 2 3 7} -A_{1 4 5 7} A_{1 2 3 6} +A_{1 4 6 7} A_{1 2 3 5} , -A_{1 5 6 7} A_{1 2 3 4}\right )x_{2} \wedge x_{3} \wedge x_{4} \wedge x_{5} \wedge x_{6} \wedge x_{7}\text{.}\end{align}So it seems that this term cancels. 

One can consider the special case $r =4$ the number of subscripts that have two values in common are
\begin{equation}\binom{8}{4} \binom{8 -4}{2} =70 \ast 12 =840
\end{equation}e.g.
\begin{equation}12341256,12351246,12361245,12451236,12461235,12561234.
\end{equation}the coefficients are given by
\begin{align} & A_{1 2 3 4} A_{1 2 5 6} x_{3} \wedge x_{4} \wedge x_{5} \wedge x_{6} \\
 & A_{1 2 3 5} A_{1 2 4 6} x_{3} \wedge x_{5} \wedge x_{4} \wedge x_{6} \\
 & A_{1 2 3 6} A_{1 2 4 5} x_{3} \wedge x_{6} \wedge x_{4} \wedge x_{5} \\
 & A_{1 2 4 5} A_{1 2 3 6} x_{4} \wedge x_{5} \wedge x_{3} \wedge x_{6} \\
 & A_{1 2 4 6} A_{1 2 3 5} x_{4} \wedge x_{6} \wedge x_{3} \wedge x_{5} \\
 & A_{1 2 5 6} A_{1 2 3 4} x_{5} \wedge x_{6} \wedge x_{3} \wedge x_{4}\end{align}which gives on reordering
\begin{align} & A_{1 2 3 4} A_{1 2 5 6} x_{3} \wedge x_{4} \wedge x_{5} \wedge x_{6} -A_{1 2 3 5} A_{1 2 4 6} x_{3} \wedge x_{4} \wedge x_{5} \wedge x_{6} +A_{1 2 3 6} A_{1 2 4 5} x_{3} \wedge x_{4} \wedge x_{5} \wedge x_{6} \\
 &  +A_{1 2 4 5} A_{1 2 3 6} x_{3} \wedge x_{4} \wedge x_{5} \wedge x_{6} -A_{1 2 4 6} A_{1 2 3 5} x_{3} \wedge x_{4} \wedge x_{5} \wedge x_{6} +A_{1 2 5 6} A_{1 2 3 4} x_{3} \wedge x_{4} \wedge x_{5} \wedge x_{6}\end{align}
\begin{equation} =2 \left (A_{1 2 3 4} A_{1 2 5 6} -A_{1 2 3 5} A_{1 2 4 6} +A_{1 2 3 6} A_{1 2 4 5}\right ) x_{3} \wedge x_{4} \wedge x_{5} \wedge x_{6}
\end{equation}which in general is non zero. 

One can consider the special case $r =4$ the subscripts that have three values in common are
\begin{equation}\binom{8}{4} \binom{8 -4}{1} =70 \ast 4 =280
\end{equation}e.g.
\begin{align} & 12341235,12341236,12341237,12341238,12351234,12351236,12351237,12351238\text{,} \\
 & 12361234,12361235,12361237,12361238,12371238\text{,}\end{align}the coefficients are
\begin{equation}A_{1 2 3 4} A_{1 2 3 5} -A_{1 2 3 5} A_{1 2 3 4} +A_{1 2 3 4} A_{1 2 3 6} -A_{1 2 3 6} A_{1 2 3 4} +A_{1 2 3 4} A_{1 2 3 5} -A_{1 2 3 5} A_{1 2 3 4} \cdots \text{.}
\end{equation}So it seems that this term cancels. 

One can consider the special case $r =4$ the number of subscripts that have four values in common are
\begin{equation}\binom{8}{4} \binom{8 -4}{0} =70 \ast 1 =70
\end{equation}
\begin{equation}12341234 +12351235 +\cdots \text{.}
\end{equation}

\subsubsection{Self Adjoint Idempotents i.e. Projections.}
In order that idempotents describe projections and be positive they must be self adjoint. This induces further constriants on the preceeding
conditions. This gives
\begin{align}A_{0} &  =A_{0}^{ \dag } \Rightarrow \alpha  =\overline{\alpha } \Rightarrow \alpha  \in \mathbb{R} \\
A_{2} &  =A_{2}^{ \dag } =\sum _{1 \leq j_{1} <j_{2} \leq 2 r}A_{2 j_{1} j_{2}} x_{j_{1}} x_{j_{2}} = -\sum _{1 \leq j_{1} <j_{2} \leq 2 r}\overline{A}_{2 j_{1} j_{2}} x_{j_{1}} x_{j_{2}} \Rightarrow \mathbb{A}_{2} = -\overline{\mathbb{A}}_{2} \Rightarrow \mathbb{A}_{2} \\
 &  =i \ensuremath{\operatorname*{Im}} \left \{\mathbb{A}_{2}\right \}\, \&\,\mathbb{A}_{2} =\mathbb{A}_{2}^{t} ;\mathbb{A}_{2} \in \mathbb{R}^{2 r} \\
A_{4} &  =A_{4}^{ \dag } =\sum _{1 \leq j_{1} <j_{2} <j_{3} <j_{4} \leq 2 r}A_{4 j_{1} j_{2} j_{3} j_{4}} x_{j_{1}} x_{j_{2}} x_{j_{3}} x_{j_{4}} =\sum _{1 \leq j_{1} <j_{2} <j_{3} <j_{4} \leq 2 r}\overline{A}_{4 j_{1} j_{2} j_{3} j_{4}} x_{j_{4}} x_{j_{3}} x_{j_{2}} x_{j_{1}} \\
 &  =\sum _{1 \leq j_{1} <j_{2} <j_{3} <j_{4} \leq 2 r}\overline{A}_{4 j_{1} j_{2} j_{3} j_{4}} x_{j_{1}} x_{j_{2}} x_{j_{3}} x_{j_{4}} \Rightarrow \mathbb{A}_{4} \in \mathbb{R}^{4 r^{2}}\text{.}\end{align}

\section{Geometric Algebra II}
The operator and Clifford products in $\mathcal{F} =\mathcal{B} \left (\mathcal{H}_{\mathcal{F}}\right )$ are the same and are also known as the \emph{geometric product. }Thus $\mathcal{F}$ can be thought of as an operator algebra, Clifford algebra or Geometric algebra. The product of two grade 1 operators (vectors)
\begin{equation}a =\sum _{1 \leq j \leq 2 r}\alpha _{j} x_{j}\text{and}\,b =\sum _{1 \leq j \leq 2 r}\beta _{j} x_{j}
\end{equation}is given by
\begin{align}a b &  =\sum _{1 \leq j ,k \leq 2 r}\alpha _{j} \beta _{k} x_{j} x_{k} =\sum _{1 \leq j \leq 2 r}\alpha _{j} \beta _{j} x_{j}^{2} +\sum _{1 \leq j \neq k \leq 2 r}\alpha _{j} \beta _{k} x_{j} x_{k} \nonumber  \\
 &  =\frac{1}{2} \sum _{1 \leq j \leq 2 r}\alpha _{j} \beta _{j} +\sum _{1 \leq j <k \leq 2 r}\left (\alpha _{j} \beta _{k} -\alpha _{k} \beta _{j}\right ) x_{j} x_{k} \nonumber  \\
 &  =\frac{1}{2} \left (\left .a\right \vert b\right ) +a \wedge b\end{align}and the product $a b$ expressed in this fashion is called the \emph{geometric product }of the \emph{vectors} $a$ and $b\text{.}$ The products $ \langle \left .a\right \vert b \rangle $ and $a \wedge b$ can be expressed or defined by
\begin{equation*}\left (\left .a\right \vert b\right ) =\frac{1}{2} \left (a b +b a\right ) =\frac{1}{2} \left \{a ,b\right \}\text{,}
\end{equation*}which is $\mathbb{C}$-symmetric and
\begin{equation*}a \wedge b =\frac{1}{2} \left (a b -b a\right ) =\frac{1}{2} \left [a ,b\right ]\text{.}
\end{equation*}

\subsection{Degree of Operators}
$N$-degree operators are defined to be $N^{t h}$ degree polynomials in the generators of the Clifford algebra and are of the form
\begin{equation}A_{N} =\alpha _{0} I_{\mathcal{F}} +\sum _{1 \leq M \leq N}\sum _{1 \leq j_{1} <\cdots  <j_{M} \leq 2 r}\alpha _{M j_{1} \cdots  j_{M}} x_{j_{1}} \cdots  x_{j_{M}}
\end{equation}and grade $N$ operators are homogeneous of degree $N$ i.e.
\begin{equation}B_{N} =\sum _{1 \leq j_{1} <\cdots  <j_{N} \leq 2 r}\alpha _{N j_{1} \cdots  j_{N}} x_{j_{1}} \cdots  x_{j_{N}}\text{.}
\end{equation}

\emph{Blades} are outer $ \equiv $ wedge products of vectors, while \emph{versors} are geometric product of vectors. 

\subsubsection{First Degree Operators}
First degree operators are first degree polynomials in the generators of the Clifford algebra
\begin{equation*}A_{1} =\alpha _{0} I_{\mathcal{F}} +\sum _{1 \leq j \leq 2 r}\alpha _{1 j} x_{j} ,
\end{equation*}these are also called \emph{para vectors}. 

Grade 1 operators are homogeneous of
degree 1 and can be expressed in terms of a non a self adjoint basis as
\begin{align}A_{1} &  =\sum _{1 \leq j \leq 2 r}\alpha _{j} x_{j} =\sum _{1 \leq j \leq r}\left \{\alpha _{j} x_{j} +\alpha _{j +r} x_{j +r}\right \} =\sum _{1 \leq j \leq r}\left \{\alpha _{j} \left (a_{j}^{ \dag } +a_{j}\right ) +\alpha _{j +r} i \left (a_{j}^{ \dag } -a_{j}\right )\right \} \nonumber  \\
 &  =\sum _{1 \leq j \leq r}\left \{\alpha _{j} a_{j}^{ \dag } +\alpha _{j} a_{j} +i \alpha _{j +r} a_{j}^{ \dag } -i \alpha _{j +r} a_{j}\right \} =\sum _{1 \leq j \leq r}\left \{\left (\alpha _{j} +i \alpha _{j +r}\right ) a_{j}^{ \dag } +\left (\alpha _{j} -i \alpha _{j +r}\right ) a_{j}\right \}\text{.}\end{align}If $A_{1}$ is self adjoint then
\begin{equation}A_{1} =\sum _{1 \leq j \leq r}\left \{\left (\alpha _{j} +i \alpha _{j +r}\right ) a_{j}^{ \dag } +\left (\alpha _{j} -i \alpha _{j +r}\right ) a_{j}\right \} =\sum _{1 \leq j \leq r}\left \{z_{j} a_{j}^{ \dag } +\overline{z}_{j} a_{j}\right \}\text{,}
\end{equation}where
\begin{equation}z_{j} =\left (\alpha _{j} +i \alpha _{j +r}\right )
\end{equation}if not
\begin{equation}A_{1} =\sum _{1 \leq j \leq r}\left \{z_{1 j} a_{j}^{ \dag } +z_{2 j} a_{j}\right \}\text{,}
\end{equation}where
\begin{align}z_{1 j} &  =\alpha _{j} +i \alpha _{j +r} \\
z_{2 j} &  =\alpha _{j} -i \alpha _{j +r}\text{.}\end{align}

\subsubsection{Second Degree Operators}
Second degree operators are second degree polynomials in the generators of the Clifford algebra
\begin{equation*}A_{2} =\alpha _{0} I_{\mathcal{F}} +\sum _{1 \leq j \leq 2 r}\alpha _{1 j} x_{j} +\sum _{1 \leq j <k \leq 2 r}\alpha _{2 j k} x_{j} x_{k} .
\end{equation*}Homogeneous second degree operators of degree 2 can be expressed in terms of a non self adjoint basis as
\begin{align}A_{2} &  =\sum _{1 \leq j <k \leq 2 r}\alpha _{j k} x_{j} x_{k} \nonumber  \\
 &  =2 \sum _{1 \leq j <k \leq r}\left (\alpha _{j k} -\alpha _{k j}\right ) \left (a_{j}^{ \dag } +a_{j}\right ) \left (a_{k}^{ \dag } +a_{k}\right ) + \nonumber  \\
 &  +i \sum _{1 \leq j \leq r ,r +1 \leq k \leq 2 r}\left (\alpha _{j k} -\alpha _{k j}\right ) \left (a_{j}^{ \dag } +a_{j}\right ) \left (a_{k}^{ \dag } -a_{k}\right ) -2 \sum _{r +1 \leq j <k \leq 2 r}\left (\alpha _{j k} -\alpha _{k j}\right ) \left (a_{j}^{ \dag } -a_{j}\right ) \left (a_{k}^{ \dag } -a_{k}\right ) \nonumber  \\
 &  =2 \sum _{1 \leq j <k \leq r}\left (\alpha _{j k} -\alpha _{k j}\right ) \left (a_{j}^{ \dag } a_{k}^{ \dag } +a_{j}^{ \dag } a_{k} +a_{j} a_{k}^{ \dag } +a_{j} a_{k}\right ) \nonumber  \\
 &  +i \sum _{1 \leq j \leq r ,r +1 \leq k \leq 2 r}\left (\alpha _{j k} -\alpha _{k j}\right ) \left (a_{j}^{ \dag } a_{k}^{ \dag } -a_{j}^{ \dag } a_{k} +a_{j} a_{k}^{ \dag } -a_{j} a_{k}\right ) \nonumber  \\
 &  -2 \sum _{r +1 \leq j <k \leq 2 r}\left (\alpha _{j k} -\alpha _{k j}\right ) \left (a_{j}^{ \dag } a_{k}^{ \dag } -a_{j}^{ \dag } a_{k} -a_{j} a_{k}^{ \dag } +a_{j} a_{k}\right )\text{.}\end{align}Grade $2$ operators are homogeneous of degree $2$ i.e.
\begin{equation}A_{2} =\sum _{1 \leq j_{1} <j_{2} \leq 2 r}\alpha _{2 j_{1} j_{2}} x_{j_{1}} x_{j_{2}}\text{.}
\end{equation}

\section{Lie Algebra}
The algebra $\mathcal{A}$ becomes a Lie algebra under the commutator product $\left [\text{,}\right ]$ i.e.
\begin{equation}\left [A ,B\right ] =A B -B A\text{.}
\end{equation}Commutator closed subspaces of $\mathcal{A}$ are also Lie algebras, but may not be subalgebras of $\mathcal{A}\text{.}$ 

The \emph{complex} subspace, $\mathcal{A}_{0}^{\bot }\text{,}$ of \emph{traceless} elements of $\mathcal{A}$ is a closed under Lie products as is the \emph{real }subspace , $\mathcal{A}_{s k a}\text{,}$ of \emph{skew adjoint} elements of $\mathcal{A}\text{.}$ 

\section{Jordan Algebras (Euclidean=Formally Real)}
The set of self adjoint operators, $\mathcal{A}_{s a}\text{,}$ is a real subspace of $\mathcal{A}\text{,}$ however it is not a subalgebra. In order to overcome this one can introduce the Jordan product $\left \{\text{,}\right \} \equiv  \ast $ by
\begin{equation}A \ast B =\frac{1}{2} \left \{A ,B\right \}\text{} \forall A ,B \in \mathcal{A}_{s a}\text{.}
\end{equation}Equipped with this product $\mathcal{A}_{s a}$ becomes a \emph{Special} \emph{Jordan} algebra $\left \{\mathcal{A}_{s a} , \ast \right \}$ which is \emph{non associative, but commutative}$ \equiv $\emph{Abelian }i.e. in general
\begin{equation}A \ast \left (B \ast C\right ) \neq \left (A \ast B\right ) \ast C
\end{equation}but
\begin{equation}A \ast B =B \ast A\text{.}
\end{equation}Formally a Jordan algebra satisfies the axioms 


\begin{enumerate}
\item
\begin{equation}A \ast B =B \ast A
\end{equation}

\item
\begin{equation}A^{2} \ast \left (B \ast A\right ) =\left (A^{2} \ast B\right ) \ast A
\end{equation}

\item If
\begin{equation}\sum _{1 \leq j \leq p}A_{j}^{2} =0 \Rightarrow A_{1} =A_{2} =\cdots  =A_{p} =0
\end{equation}the Jordan algebra is \emph{formally real}$ \equiv $\emph{Euclidean}. \end{enumerate}


Any real basis $\left \{I_{\mathcal{A}} ,\left \{\left .x_{j}\right \vert 1 \leq j \leq 2 r\right \}\right \}$ generates $\left \{\mathcal{A}_{s a} , \ast \right \}$ and $I_{\mathcal{A}}$ is the multiplicative unit, as $\mathcal{A}$ is a Clifford/CAR algebra one has the following relationships between the generators
\begin{align}x_{j_{1}} \ast x_{j_{2}} &  =\frac{1}{2} \delta _{j_{1} j_{2}} I_{\mathcal{A}} \\
v_{1} \ast v_{2} &  =\frac{1}{2} \left (\left .v_{1}\right \vert v_{2}\right ) ;v_{1} ,v_{2} \in \mathcal{A}_{1 s a} \\
v_{1} \ast \left (v_{2} \ast v_{3}\right ) &  =v_{1} \ast \left (v_{3} \ast v_{2}\right ) =\left (v_{3} \ast v_{2}\right ) \ast v_{1} =\frac{1}{2} \left (\left .v_{2}\right \vert v_{3}\right ) v_{1} \\
\left (v_{1} \ast v_{2}\right ) \ast v_{3} &  =\left (v_{2} \ast v_{1}\right ) \ast v_{3} =v_{3} \ast \left (v_{2} \ast v_{1}\right ) =\frac{1}{2} \left (\left .v_{1}\right \vert v_{2}\right ) v_{3}\text{.}\end{align}The paravectors form a Euclidean \emph{sub} Jordan algebra. 

$\lceil $
\begin{align}\frac{1}{2} \left (i x_{1} x_{2} i x_{3} x_{4} +i x_{3} x_{4} i x_{1} x_{2}\right ) &  =\text{-}\frac{1}{2} \left (x_{1} x_{2} x_{3} x_{4} +x_{3} x_{4} x_{1} x_{2}\right ) \\
\text{-}\frac{1}{2} \left (x_{1} x_{2} x_{3} x_{4} +x_{3} x_{4} x_{1} x_{2}\right )^{ \dag } &  =\text{-}\frac{1}{2} \left (x_{4} x_{3} x_{2} x_{1} +x_{2} x_{1} x_{4} x_{3}\right ) =\text{-}\frac{1}{2} \left (x_{3} x_{4} x_{1} x_{2} +x_{1} x_{2} x_{3} x_{4}\right )\end{align}
\begin{align}\frac{1}{2} \left (P_{1} P_{2} +P_{2} P_{1}\right )^{2} &  =\frac{1}{4} \left (P_{1} P_{2} +P_{2} P_{1}\right ) \left (P_{1} P_{2} +P_{2} P_{1}\right ) =\frac{1}{4} \left (P_{1} P_{2} P_{1} P_{2} +P_{2} P_{1} P_{1} P_{2} +P_{1} P_{2} P_{2} P_{1} +P_{2} P_{1} P_{2} P_{1}\right ) \\
 &  \neq \frac{1}{2} \left (P_{1} P_{2} +P_{2} P_{1}\right )\end{align}$\rfloor $ 

The vector space $\mathcal{A}_{s a}$ is ordered and generated by the positive cone $\mathcal{A}_{s a +}$ i.e.
\begin{equation}\mathcal{A}_{s a} =\mathcal{A}_{s a +} -\mathcal{A}_{s a +}\text{.}
\end{equation}


\begin{definition}
Jordan Quadratic Product
\begin{equation}U_{x} \left (y\right ) =2 x \ast \left (x \ast y\right ) -x^{2} \ast y =x \ast y \ast x\text{,}
\end{equation}which gives
\begin{equation}U_{U_{x} \left (y\right )} =U_{x} U_{y} U_{x}\text{.}
\end{equation}
\end{definition}


\begin{definition}
Jordan Triple Product
\begin{equation}\left \{x ,y ,z\right \} =2 \left (x \ast \left (y \ast z\right ) +\left (x \ast y\right ) \ast z -\left (x \ast z\right ) \ast y\right ) =x \ast y \ast z +z \ast y \ast x\text{.}
\end{equation}
\end{definition}

\paragraph{Spectral Expansion}
As all elements in a Jordan algebra are self adjoint in the enveloping associative algebra so one always has the spectral theorem
\begin{equation}A =\int _{\ensuremath{\operatorname*{Spec}} \left (A\right ) \subseteq \mathbb{R}}\lambda  d P_{A} \left (\lambda \right )
\end{equation}giving
\begin{equation}A \ast B =\int _{\ensuremath{\operatorname*{Spec}} \left (A \ast B\right ) \subseteq \mathbb{R}}\lambda  d P_{A \ast B} \left (\lambda \right )\text{.}
\end{equation}$\lceil $
\begin{align}\sum _{1 \leq j \leq 2 r}\alpha _{j}\vert x_{j} \rangle  \langle x_{j}\vert \sum _{1 \leq k \leq 2 r}\beta _{k}\vert y_{k} \rangle  \langle y_{k j}\vert  &  =\sum _{1 \leq j .k \leq 2 r}\alpha _{j} \beta _{k}\vert x_{j} \rangle  \langle x_{j}\vert \vert y_{k} \rangle  \langle y_{k}\vert  \\
\sum _{1 \leq k \leq 2 r}\beta _{k}\vert y_{k} \rangle  \langle y_{k j}\vert \sum _{1 \leq j \leq 2 r}\alpha _{j}\vert x_{j} \rangle  \langle x_{j}\vert  &  =\sum _{1 \leq j .k \leq 2 r}\alpha _{j} \beta _{k}\vert y_{k} \rangle  \langle y_{k j}\vert \vert x_{j} \rangle  \langle x_{j}\vert  \\
\sum _{1 \leq j \leq 2 r}\alpha _{j}\vert x_{j} \rangle  \langle x_{j}\vert  \ast \sum _{1 \leq k \leq 2 r} \beta _{k}\vert y_{k} \rangle  \langle y_{k j}\vert  &  =\frac{1}{2} \sum _{1 \leq j .k \leq 2 r}\alpha _{j} \beta _{k} \left \{\vert x_{j} \rangle  \langle x_{j}\vert \vert y_{k} \rangle  \langle y_{k}\vert  +\vert y_{k} \rangle  \langle y_{k}\vert \vert x_{j} \rangle  \langle x_{j}\vert \right \}\end{align}$\rfloor $ 

\section{Change of Basis of $\mathcal{F}$}


\subsection{Resolution of the Identity}
In general
\begin{equation}I_{V} =\sum _{1 \leq j \leq r}\vert v_{j} \rangle  \langle v_{j}^{d}\vert \text{,}
\end{equation}where $\left \{\left .v_{j} ,v_{j}^{d}\right \vert 1 \leq j \leq r\right \}$ are biorthogonal bases for the dual vector spaces $\left \{V ,V^{d}\right \}$. If $V$ is a Hilbert space it is self dual and $V =V^{d}$. If $V$ is a Hibert space one has the SVD identity in $\mathcal{B}_{0} \left (V\right )$
\begin{equation}A =U_{A} \mathbf{\alpha }_{d} W_{A}\text{,}
\end{equation}where $U_{A}\;$and $W_{A}$ are unitary operators
\begin{equation}\mathbf{\alpha }_{d} =\left (A^{ \dag } A\right )^{\frac{1}{2}} =\sum _{1 \leq j \leq r}\alpha _{j}^{\frac{1}{2}} P_{\alpha _{j}}
\end{equation}and
\begin{equation}P_{\alpha _{j}} =P_{\alpha _{j}}^{2} =P_{\alpha _{j}}^{ \dag } ,1 \leq j \leq r\text{.}
\end{equation}If $A$ is invertible
\begin{equation}A^{ -1} =W_{A}^{ \dag } \mathbf{\alpha }_{d}^{ -1} U_{A}^{ \dag }\text{.}
\end{equation}


\begin{proposition}
If $\left \{\left .e_{j}\right \vert 1 \leq j \leq r\right \}$ is a conb for $V$\emph{\ } 


\begin{enumerate}
\item then they form a resolution of the identity
\begin{equation}I_{V} =\sum _{1 \leq j \leq r}\vert e_{j} \rangle  \langle e_{j}\vert 
\end{equation}

\item if $A \in G L \left (V\right )$
\begin{equation}A I_{V} A^{ -1} =\sum _{1 \leq j \leq r}A\vert e_{j} \rangle  \langle e_{j}\vert A^{ -1} =\sum _{1 \leq j \leq r}\vert A e_{j} \rangle  \langle A^{ - \dag } e_{j}\vert  =I_{V}\text{,}
\end{equation}is a resolutions of the identity and $\left \{\left .A e_{j} ,A^{ - \dag } e_{j}\right \vert 1 \leq j \leq r\right \}$ is a \emph{biorthogonal basis} for $V\text{.}$ \end{enumerate}
\end{proposition}

In particular one can consider
\begin{equation}A =\sum _{1 \leq j \leq r}\xi _{j}\vert e_{j} \rangle  \langle e_{j}\vert \text{,}
\end{equation}where
\begin{equation}A \in \mathcal{B}_{1} \left (V\right )
\end{equation}so that
\begin{equation}\sum _{1 \leq j \leq r}\left \vert \xi _{j}\right \vert  <\infty \text{.}
\end{equation}


\begin{remark}
If $r =\infty $ then $A^{ -1} \notin \mathcal{B}_{1} \left (V\right )\text{.}$ 
\end{remark}

\subsection{CAR/Clifford/Hilbert Algebras}
The Hilbert space $\mathcal{F}_{1 H S}$ generates the Hilbert and Exterior algebra $\mathcal{F}_{H S}$ by
\begin{equation}\mathcal{F}_{H S} = \tbigwedge \left (\mathcal{F}_{1 H S}\right )\text{,}
\end{equation}sometimes we will drop the "$H S "" o r$ $"" \ensuremath{\operatorname*{tr}}""$ subscript when it is not crucial. 

A set of resolutions of the identity in $\mathcal{B} \left (\mathcal{F}_{1}\right )$ is given by
\begin{equation}I_{\mathcal{F}_{1}} =\sum _{1 \leq j \leq 2 r}\left \{\vert \widetilde{a}_{j}^{ \dag } \rangle _{\ensuremath{\operatorname*{tr}}}  \langle \widetilde{a}_{j}^{ \dag }\vert  +\vert \widetilde{a}_{j} \rangle _{\ensuremath{\operatorname*{tr}}} \langle \widetilde{a}_{j}\vert \right \} =\sum _{1 \leq j \leq 2 r}K \left (\mathbf{\xi }\right ) \left \{\vert \widetilde{a}_{j}^{ \dag } \rangle _{\ensuremath{\operatorname*{tr}}}  \langle \widetilde{a}_{j}^{ \dag }\vert  +\vert \widetilde{a}_{j} \rangle _{\ensuremath{\operatorname*{tr}}} \langle \widetilde{a}_{j}\vert \right \} K \left (\mathbf{\xi }\right )^{ - \dag }\text{,}
\end{equation}where
\begin{align}\left (\begin{array}{cc}\mathbf{x}_{1} \left (\mathbf{\xi }\right ) & \mathbf{x}_{2} \left (\mathbf{\xi }\right )\end{array}\right ) &  =K \left (\mathbf{\xi }\right ) \left (\begin{array}{cc}\mathbf{a}^{ \dag } & \mathbf{a}\end{array}\right ) =\left (\begin{array}{cc}\mathbf{a}^{ \dag } & \mathbf{a}\end{array}\right ) \mathbb{K} \left (\mathbf{\xi }\right ) =\left (\begin{array}{cc}\mathbf{a}^{ \dag } & \mathbf{a}\end{array}\right ) \left (\begin{array}{cc}\frac{1}{\sqrt{2}} \mathbf{\xi }_{r} & \frac{i}{\sqrt{2}} \overline{\mathbf{\xi }}_{r} \\
\frac{1}{\sqrt{2}} \mathbf{\xi }_{r} &  -\frac{i}{\sqrt{2}} \overline{\mathbf{\xi }}_{r}\end{array}\right ) \\
 &  =\frac{1}{\sqrt{2}} \left (\begin{array}{cc}\mathbf{a}^{ \dag } \mathbf{\xi }_{r} +\mathbf{a} \mathbf{\xi }_{r} & i \left (\mathbf{a}^{ \dag } \overline{\mathbf{\xi }}_{r} -\mathbf{a} \overline{\mathbf{\xi }}_{r}\right )\end{array}\right )\end{align}and
\begin{equation}\mathbf{\xi }_{r} =\left (\begin{array}{ccc}\xi _{1} & \cdots  & 0 \\
\vdots  & \ddots  & \vdots  \\
0 & \cdots  & \xi _{r}\end{array}\right )\text{.}
\end{equation}


\begin{remark}
The transformation $K \left (\mathbf{\xi }\right )\text{,}$where
\begin{equation}\mathcal{R}_{a} \left (K \left (\mathbf{\xi }\right )\right ) =\left (\begin{array}{cc}\frac{1}{\sqrt{2}} \mathbf{\xi }_{r} & \frac{i}{\sqrt{2}} \overline{\mathbf{\xi }}_{r} \\
\frac{1}{\sqrt{2}} \overline{\mathbf{\xi }_{r}} &  -\frac{i}{\sqrt{2}} \overline{\overline{\mathbf{\xi }}_{r}}\end{array}\right )
\end{equation}is a $ \ast $-transformation that belongs to $G L \left (2 r ,\mathbb{R}\right )\text{,}$ but does conserve the CAR so is not cannonical and thus does not belong to $O \left (2 r ,\mathbb{R}\right )\text{.}$ 


\begin{summary}
We summarize some of the properties of the Fock/CAR/Clifford/Hilbert algebra: 


\begin{enumerate}
\item
\begin{equation}\mathcal{F} = \tbigwedge \left (\mathcal{F}_{1}\right )
\end{equation}

\item
\begin{align}\mathcal{F}_{H S} &  =\left \{\left .A\right \vert \ensuremath{\operatorname*{tr}}\left \{A^{ \dag } A\right \} <\infty \right \} \\
\left \Vert A\right \Vert _{\mathcal{F}_{H S}} &  =\left ( \langle \left .A\right \vert A \rangle _{\mathcal{F}_{H S}}\right )^{\frac{1}{2}} =\left (\ensuremath{\operatorname*{tr}}\left \{A^{ \dag } A\right \}\right )^{\frac{1}{2}} \\
 \langle \left .A\right \vert B \rangle _{\mathcal{F}_{H S}} &  =\ensuremath{\operatorname*{tr}}\left \{A^{ \dag } B\right \}\end{align}

\item
\begin{align}\mathcal{F}_{T C} &  =\left \{\left .\overset{}{A}\right \vert \ensuremath{\operatorname*{tr}}\left \{\left (A^{ \dag } A\right )^{\frac{1}{2}}\right \} <\infty \right \} \\
\left \Vert A\right \Vert _{\mathcal{F}_{T C}} &  =\ensuremath{\operatorname*{tr}}\left \{\left (A^{ \dag } A\right )^{\frac{1}{2}}\right \} =\ensuremath{\operatorname*{tr}}\left \{\left \vert A\right \vert \right \}\end{align}

\item
\begin{equation}\mathcal{F}_{H S} =\Phi _{\ensuremath{\operatorname*{tr}}} \left (\mathcal{F}\right ) \equiv \mathcal{H}_{\ensuremath{\operatorname*{tr}}}
\end{equation}
\begin{equation} \langle \left .A\right \vert B \rangle _{\mathcal{F}_{H S}} \equiv  \langle \left .A\right \vert B \rangle _{\ensuremath{\operatorname*{tr}}} =\ensuremath{\operatorname*{tr}} \left \{A^{ \dag } B\right \} = \langle \left .\Phi _{\ensuremath{\operatorname*{tr}}} \left (A\right )\right \vert \Phi _{\ensuremath{\operatorname*{tr}}} \left (B\right ) \rangle _{\mathcal{H}_{\ensuremath{\operatorname*{tr}}}} \equiv  \langle \left .\Phi _{\ensuremath{\operatorname*{tr}}} \left (A\right )\right \vert \Phi _{\ensuremath{\operatorname*{tr}}} \left (B\right ) \rangle _{\ensuremath{\operatorname*{tr}}}
\end{equation}

\item
\begin{equation}\mathcal{F} \supset \mathcal{F}_{H S} \supset \mathcal{F}_{T C}
\end{equation}
\begin{align}\mathcal{F}_{H S} &  =\Phi _{\ensuremath{\operatorname*{tr}}} \left (\mathcal{F}\right ) \\
 \langle \left .A\right \vert B \rangle _{\mathcal{F}_{H S}} &  = \langle \left .\Phi _{\ensuremath{\operatorname*{tr}}} \left (A\right )\right \vert \Phi _{\ensuremath{\operatorname*{tr}}} \left (B\right ) \rangle _{\ensuremath{\operatorname*{tr}}}\end{align}

\item
\begin{align}\dim  \left \{\mathcal{F}\right \} &  =\dim  \left \{\mathcal{F}_{H S}\right \} =\dim  \left \{\Phi _{\ensuremath{\operatorname*{tr}}} \left (\mathcal{F}\right )\right \} =2^{2 r} \\
\dim  \left \{\mathcal{B} \left (\mathcal{F}\right )\right \} &  =\dim  \left \{\mathcal{B} \left (\Phi _{\ensuremath{\operatorname*{tr}}} \left (\mathcal{F}\right )\right )\right \} =2^{4 r} \\
\dim  \left \{\Pi _{\ensuremath{\operatorname*{tr}}} \left (\mathcal{F}\right )\right \} &  =2^{2 r}\end{align}

\item Bounded operators, $\left \{\left .X_{j}\right \vert 1 \leq j \leq n\right \} \in \mathcal{B} \left (\mathcal{F}_{1}\right )\text{,}$ acting in $\mathcal{F}_{1}$ induce bounded operators $X_{1} \overset{ \leftarrow n \rightarrow }{ \boxtimes \cdots  \boxtimes } X_{n} \in \mathcal{B} \left (\mathcal{F}_{n}\right )$ in particular $X \overset{ \leftarrow n \rightarrow }{ \boxtimes \cdots  \boxtimes } X$ when $X_{j} =X ,\,1 \leq j \leq n\text{.}$ 

\item
\begin{equation}\Pi _{\ensuremath{\operatorname*{tr}}} \left (\mathcal{F}\right ) \subset \mathcal{B} \left (\mathcal{F}_{H S}\right ) \equiv \mathcal{B} \left (\Phi _{\ensuremath{\operatorname*{tr}}} \left (\mathcal{F}\right )\right ) \subset \mathcal{B} \left (\mathcal{F}\right )
\end{equation}

\item
\begin{equation}\mathcal{B} \left (\Phi _{\ensuremath{\operatorname*{tr}}} \left (\mathcal{F}\right )\right ) \supseteq \mathcal{B}_{0} \left (\Phi _{\ensuremath{\operatorname*{tr}}} \left (\mathcal{F}\right )\right ) \supseteq \mathcal{B}_{2} \left (\Phi _{\ensuremath{\operatorname*{tr}}} \left (\mathcal{F}\right )\right ) \supseteq \mathcal{B}_{1} \left (\Phi _{\ensuremath{\operatorname*{tr}}} \left (\mathcal{F}\right )\right ) \supseteq \mathcal{B}_{0 0} \left (\Phi _{\ensuremath{\operatorname*{tr}}} \left (\mathcal{F}\right )\right )
\end{equation}

\item $\Pi _{\ensuremath{\operatorname*{tr}}} \left (\mathcal{F}_{H S}\right )$ is a \emph{Hilbert Algebra}. \end{enumerate}
\end{summary}


\end{remark}


\begin{definition}
A Hilbert Algebra is an algebra $\mathfrak{A}$ with involution $\#$ over the field of complex numbers $\mathbb{C}$, equipped with a non-degenerate scalar product $ \langle \left .\text{}\right \vert  \rangle $, for which the following axioms are satisfied: 


\begin{enumerate}
\item $ \langle \left .X\right \vert Y \rangle  = \langle \left .Y^{\#}\right \vert X^{\#} \rangle $ $ \forall X ,Y \in $ $\mathfrak{A}$; 

\item $ \langle \left .X Y\right \vert Z \rangle  = \langle \left .Y\right \vert X^{\#} Z \rangle $ $ \forall X ,Y \in $ $\mathfrak{A}$; 

\item $ \forall X \in $ $\mathfrak{A}$ the mapping $Y \rightarrow X Y$ of $\mathfrak{A}$ into $\mathfrak{A}$ is continuous; 

\item the set of elements of the form $X Y ;$ $X ,Y \in $ $\mathfrak{A}$, is everywhere dense in $\mathfrak{A}$ wrt the norm induced by $ \langle \left .\text{}\right \vert  \rangle $. \end{enumerate}
\end{definition}


\begin{example}
Examples of Hilbert algebras include the algebras $L_{2} (G)$ (with respect to convolution), where $G$ is a compact topological group, and the algebra of Hilbert--Schmidt operators on a given Hilbert space. 
\end{example}


\begin{definition}
a real J\emph{ordan-Hilbert Algebra} (JH-algebra for short). is 


\begin{enumerate}
\item is a real Hilbert space, 

\item real Jordan algebra, so
that $ \exists $ a bilinear product
\begin{equation} \ast \left (X ,Y\right ) \rightarrow X \ast Y
\end{equation}such that 


\begin{enumerate}
\item
\begin{equation}X \ast Y =Y \ast X
\end{equation}

\item
\begin{equation}X^{2} \ast X \ast Y =X \ast X^{2} \ast Y\text{.}
\end{equation}\end{enumerate}


\item
\begin{equation} \langle \left .X \ast Y\right \vert Z \rangle  = \langle \left .Y\right \vert X \ast Z \rangle  \forall X ,Y ,Z \in \mathfrak{A}
\end{equation}\end{enumerate}
\end{definition}


\begin{remark}
One can have HS operators that have infinite trace (i.e. that are not trace class). 
\end{remark}

Some further properties
are 


\begin{enumerate}
\item
\begin{equation}\mathcal{F}_{1} =\mathfrak{f}_{1}  \oplus _{\mathbb{C}}\mathfrak{f}_{2}
\end{equation}

\item
\begin{equation}\mathcal{F} = \tbigwedge \left (\mathfrak{f}_{1}  \oplus _{\mathbb{C}}\mathfrak{f}_{2}\right ) = \tbigwedge \left (\mathfrak{f}_{1}\right ) \otimes  \tbigwedge  \left (\mathfrak{f}_{2}\right )
\end{equation}

\item
\begin{equation}\mathcal{F} = \tbigwedge \left (\mathfrak{f}_{1}  \oplus _{\mathbb{C}}\mathfrak{f}_{1}^{d}\right ) = \tbigwedge \left (\mathfrak{f}_{1}\right ) \otimes  \tbigwedge  \left (\mathfrak{f}_{1}^{d}\right ) \cong \mathcal{B} \left ( \tbigwedge \left (\mathfrak{f}_{1}\right )\right )
\end{equation}

\item
\begin{equation}\left \Vert a_{j}^{ \dag }\right \Vert _{\mathcal{F}_{H S}} =\left \Vert a_{j}\right \Vert _{\mathcal{F}_{H S}} =\frac{1}{\sqrt{2}}
\end{equation}
\begin{equation}\left \Vert x_{j}\right \Vert _{\mathcal{F}_{H S}} =\left \Vert x_{j +r}\right \Vert _{\mathcal{F}_{H S}} =\frac{1}{\sqrt{2}}
\end{equation}

\item
\begin{align}x_{j} &  =\frac{1}{\sqrt{2}} \left (a_{j}^{ \dag } +a_{j}\right ) ;\,1 \leq j \leq r \\
x_{j +r} &  =\frac{i}{\sqrt{2}} \left (a_{j}^{ \dag } -a_{j}\right ) ;\,1 \leq j \leq r \\
x_{j}^{2} &  =\frac{1}{2} \left (a_{j}^{ \dag } +a_{j}\right ) \left (a_{j}^{ \dag } +a_{j}\right ) =\frac{1}{2} \left (a_{j}^{ \dag } a_{j} +a_{j} a_{j}^{ \dag }\right ) =\frac{1}{2} I_{\mathcal{F}} \\
x_{j +r}^{2} &  = -\frac{1}{2} \left (a_{j}^{ \dag } -a_{j}\right ) \left (a_{j}^{ \dag } -a_{j}\right ) =\frac{1}{2} \left (a_{j}^{ \dag } a_{j} +a_{j} a_{j}^{ \dag }\right ) =\frac{1}{2} I_{\mathcal{F}} \\
\left \Vert x_{j}\right \Vert _{\mathcal{F}_{H S}} &  =\left (\ensuremath{\operatorname*{tr}}\left \{x_{j}^{2}\right \}\right )^{\frac{1}{2}} =\left \Vert x_{j +r}\right \Vert _{\mathcal{F}_{H S}} =\left (\ensuremath{\operatorname*{tr}}\left \{x_{j +r}^{2}\right \}\right )^{\frac{1}{2}} =\frac{1}{\sqrt{2}}\text{.}\end{align}

\item
\begin{align}\widetilde{a}_{j}^{ \dag } &  =\sqrt{2} a_{j}^{ \dag } ;1 \leq j \leq r \\
\widetilde{a}_{j} &  =\sqrt{2} a_{j} ;1 \leq j \leq r \\
\widetilde{x}_{j} &  =\sqrt{2} x_{j} ;1 \leq j \leq 2 r\text{.}\end{align}\end{enumerate}


The general linear group $G L \left (\mathcal{F}_{1}\right )$ of the generating subspace $\mathcal{F}_{1 H S}$ of $\mathcal{F}_{H S}$ is given by
\begin{equation}\kappa  \left (\mathbf{\lambda }\right ) =\exp  \left \{\sum _{1 \leq j ,k \leq 2 r}\lambda _{j k} \vert \mathfrak{e}_{j} \rangle _{\ensuremath{\operatorname*{tr}}} \langle \mathfrak{e}_{k}\vert \right \} =\sum _{1 \leq j ,k \leq 2 r}\kappa  \left (\mathbf{\lambda }\right )_{j k} \vert \mathfrak{e}_{j} \rangle _{\ensuremath{\operatorname*{tr}}} \langle \mathfrak{e}_{k}\vert  ;\,\lambda _{j k} ,\kappa  \left (\mathbf{\lambda }\right )_{j k} \in \mathbb{C} ,1 \leq j ,k \leq 2 r
\end{equation}and the subgroup $U \left (\mathcal{F}_{1}\right )$ by
\begin{equation}\varrho  \left (\mathbf{\lambda }\right ) =\exp i \left \{\sum _{1 \leq j ,k \leq 2 r}\lambda _{j k} \vert \mathfrak{e}_{j} \rangle _{\ensuremath{\operatorname*{tr}}} \langle \mathfrak{e}_{k}\vert \right \} =\sum _{1 \leq j ,k \leq 2 r}\varrho  \left (\mathbf{\lambda }\right )_{j k} \vert \mathfrak{e}_{j} \rangle _{\ensuremath{\operatorname*{tr}}} \langle \mathfrak{e}_{k}\vert  ;\,\lambda _{j k} =\overline{\lambda }_{k j} ,\varrho  \left (\mathbf{\lambda }\right )_{j k} \in \mathbb{C} ,1 \leq j ,k \leq 2 r\text{,}
\end{equation}where $\left \{\left .\vert \mathfrak{e}_{j} \rangle \right \vert 1 \leq j \leq 2 r\right \}$ is a conb of $\mathcal{F}_{1}\text{.}$ Both $G L \left (\mathcal{F}_{1}\right )$ and $U \left (\mathcal{F}_{1}\right )$ have reducible representations on $\mathcal{F}_{H S} = \tbigwedge \left (\mathcal{F}_{1 H S}\right )$ and are reduced by the subspaces $ \tbigwedge ^{n}\left (\mathcal{F}_{1 H S}\right )\text{,}$ $1 \leq n \leq 2 r$. 


\begin{remark}
\begin{align}X \wedge Y &  =\left (\sum _{0 \leq n \leq 2 r}\sum _{1 \leq j_{1} <\cdots  <j_{n} \leq 2 r}X_{j_{1} \cdots  j_{n}} x_{j_{1}} \cdots  x_{j_{n}}\right ) \wedge \left (\sum _{0 \leq m \leq 2 r}\sum _{1 \leq k_{1} <\cdots  <k_{m} \leq 2 r}Y_{k_{1} \cdots  k_{m}} x_{k_{1}} \cdots  x_{k_{m}}\right ) \\
 &  =\sum _{0 \leq n ,m \leq 2 r}X_{j_{1} \cdots  j_{n}} Y_{k_{1} \cdots  k_{m}} x_{j_{1}} \cdots  x_{j_{n}} x_{k_{1}} \cdots  x_{k_{m}} =\sum _{0 \leq n ,m \leq 2 r}X_{n} \wedge Y_{m} =\sum _{0 \leq p \leq 2 r}Z_{p}\text{,}\end{align}where
\begin{align}Z_{p} &  =\sum _{n +m =p}X_{n} \wedge Y_{m} ;\,0 \leq n ,m \leq 2 r \\
 &  =\sum _{1 \leq j_{1} <\cdots  <j_{p} \leq 2 r}Z_{j_{1} \cdots  j_{p}} x_{j_{1}} \cdots  x_{j_{p}}\text{.}\end{align}The Clifford algebra definition of exterior product is
\begin{equation}X \wedge Y =\sum _{0 \leq m ,n \leq 2 r} \langle X_{m} Y_{n} \rangle _{m +n}\text{.}
\end{equation}leading to
\begin{equation}Z_{p} =\sum _{n +m =p}X_{n} \wedge Y_{m} =\sum _{n +m =p} \langle X_{m} Y_{n} \rangle _{m +n}\text{.}
\end{equation}
\end{remark}


\begin{example}
\begin{align}X &  =x I_{\mathcal{F}} +X_{1} +X_{2} \\
Y &  =y I_{\mathcal{F}} +Y_{1} +Y_{2} +Y_{3}\end{align}
\begin{align}X \wedge Y &  =\left (x I_{\mathcal{F}} +X_{1} +X_{2}\right ) \wedge \left (y I_{\mathcal{F}} +Y_{1} +Y_{2} +Y_{3}\right ) \\
 &  =x y I_{\mathcal{F}} +x Y_{1} +x Y_{2} +x Y_{3} +y X_{1} +X_{1} \wedge Y_{1} +X_{1} \wedge Y_{2} +X_{1} \wedge Y_{3} \\
 &  +y X_{2} +X_{2} \wedge Y_{1} +X_{2} \wedge Y_{2} +X_{2} \wedge Y_{3} \\
 &  =x y I_{\mathcal{F}} +x Y_{1} +y X_{1} +y X_{2} +x Y_{2} +X_{1} \wedge Y_{1} +x Y_{3} +X_{1} \wedge Y_{2} +X_{2} \wedge Y_{1} \\
 &  +X_{1} \wedge Y_{3} +X_{2} \wedge Y_{2} +X_{2} \wedge Y_{3}\end{align}
\begin{equation}X \wedge Y =\sum _{0 \leq m ,n \leq 2 r} \langle X_{m} Y_{n} \rangle _{m +n}\text{.}
\end{equation}
\begin{align}X \wedge Y &  = \langle X_{0} Y_{0} \rangle _{0} +\fbox{+++} +\fbox{+} +\fbox{
+} \\
 &  +\fbox{+} + \langle X_{2} Y_{3} \rangle _{5}\end{align}
\begin{align}X_{1} Y_{2} &  = \langle X_{1} Y_{2} \rangle _{1} + \langle X_{1} Y_{2} \rangle _{3} \\
X_{1} Y_{3} &  = \langle X_{1} Y_{3} \rangle _{2} + \langle X_{1} Y_{3} \rangle _{4} \\
X_{2} Y_{3} &  = \langle X_{2} Y_{3} \rangle _{1} + \langle X_{2} Y_{3} \rangle _{3} + \langle X_{2} Y_{3} \rangle _{5}\end{align}
\begin{equation}X Y =X \wedge Y + \langle X_{1} Y_{2} \rangle _{1} + \langle X_{1} Y_{3} \rangle _{2} + \langle X_{2} Y_{3} \rangle _{1} + \langle X_{2} Y_{3} \rangle _{3} + \langle X_{2} Y_{1} \rangle _{1}
\end{equation}
\end{example}


\begin{definition}
\begin{equation} \boxtimes ^{n}\kappa  \left (\mathbf{\lambda }\right ) x_{j_{1}} \cdots  x_{j_{n}}\overset{\triangle }{ =}\kappa  \left (\mathbf{\lambda }\right ) \left (x_{j_{1}}\right ) \cdots  \kappa  \left (\mathbf{\lambda }\right ) \left (x_{j_{n}}\right ) ;1 \leq j_{1} <\cdots  <j_{n} \leq r ,0 \leq n \leq 2 r\text{,}
\end{equation}where
\begin{align} \boxtimes ^{1}\kappa  \left (\mathbf{\lambda }\right ) &  =\kappa  \left (\mathbf{\lambda }\right ) \\
 \boxtimes ^{0}\kappa  \left (\mathbf{\lambda }\right ) &  =\tau  \left (\kappa  \left (\mathbf{\lambda }\right )\right ) P_{\mathcal{F}_{0}} ,\,\tau  \left (\kappa  \left (\mathbf{\lambda }\right )\right ) \in \mathbb{C}\text{,}\end{align}which is equivalent to
\begin{equation} \boxtimes ^{n}\kappa  \left (\mathbf{\lambda }\right ) x_{j_{1}} \wedge \cdots  \wedge x_{j_{n}}\overset{\triangle }{ =}\kappa  \left (\mathbf{\lambda }\right ) \left (x_{j_{1}}\right ) \wedge \cdots  \wedge \kappa  \left (\mathbf{\lambda }\right ) \left (x_{j_{n}}\right ) ;1 \leq j_{1} <\cdots  <j_{n} \leq r ,1 \leq n \leq 2 r
\end{equation}as
\begin{equation}x_{j_{1}} x_{j_{2}} =x_{j_{1}} \wedge x_{j_{2}} \neq 0
\end{equation}when
\begin{equation}j_{1} \neq j_{2}
\end{equation}and
\begin{equation}x_{j_{1}} \wedge x_{j_{1}} =0 \neq x_{j_{1}} x_{j_{1}} =x_{j_{1}}^{2}\text{.}
\end{equation}
\end{definition}


\begin{definition}
The exterior algebra map $ \boxtimes \kappa  \left (\mathbf{\lambda }\right )$ is defined as
\begin{equation} \boxtimes \kappa  \left (\mathbf{\lambda }\right ) =\sum _{0 \leq n \leq 2 r} \boxtimes ^{n}\kappa  \left (\mathbf{\lambda }\right ) =\tbigoplus ??? \boxtimes ^{n}\kappa  \left (\mathbf{\lambda }\right )\text{.}
\end{equation}
\end{definition}


\begin{lemma}
\begin{align} \boxtimes ^{n}\kappa  \left (\mathbf{\lambda }\right ) \left (x_{j_{1}} \wedge \cdots  \wedge x_{j_{n}}\right ) &  = \boxtimes ^{p_{1}}\kappa  \left (\mathbf{\lambda }\right ) \left (x_{j_{1}} \wedge \cdots  \wedge x_{j_{p_{1}}}\right ) \wedge \cdots  \wedge  \boxtimes ^{p_{m}}\kappa  \left (\mathbf{\lambda }\right ) \left (x_{j_{m -11 +1}} \wedge \cdots  \wedge x_{j_{p_{m}}}\right ) \\
1 &  \leq p_{1} ,\cdots  ,p_{m} \leq n ;\, \ni p_{1} +\cdots  +p_{m} =n\end{align}


\begin{proof}
Let
\begin{equation}m =n :p_{1} =p_{2} =\cdots  =p_{m} =1
\end{equation}
\begin{equation} \boxtimes ^{n}\kappa  \left (\mathbf{\lambda }\right ) \left (x_{j_{1}} \wedge \cdots  \wedge x_{j_{n}}\right ) =\left ( \boxtimes ^{1}\kappa  \left (\mathbf{\lambda }\right ) \left (x_{j_{1}}\right )\right ) \wedge \cdots  \wedge \left ( \boxtimes ^{1}\kappa  \left (\mathbf{\lambda }\right ) \left (x_{j_{p_{n}}}\right )\right )
\end{equation}
\begin{equation}\vdots 
\end{equation}
\begin{equation}m =2 :p_{1} +p_{2} =n
\end{equation}
\begin{align} \boxtimes ^{n}\kappa  \left (\mathbf{\lambda }\right ) \left (x_{j_{1}} \wedge \cdots  \wedge x_{j_{n}}\right ) &  = \boxtimes ^{p_{1}}\kappa  \left (\mathbf{\lambda }\right ) \left (x_{j_{1}} \wedge \cdots  \wedge x_{j_{p_{1}}}\right ) \wedge  \boxtimes ^{n -p_{1}}\kappa  \left (\mathbf{\lambda }\right ) \left (x_{j_{p_{1} +1}} \wedge \cdots  \wedge x_{j_{n}}\right ) \\
 &  =\left ( \boxtimes ^{1}\kappa  \left (\mathbf{\lambda }\right ) \left (x_{j_{1}}\right )\right ) \wedge \cdots  \wedge \left ( \boxtimes ^{1}\kappa  \left (\mathbf{\lambda }\right ) \left (x_{j_{p_{1}}}\right )\right ) \wedge \left ( \boxtimes ^{1}\kappa  \left (\mathbf{\lambda }\right ) \left (x_{j_{p_{1} +1}}\right )\right ) \wedge \cdots  \wedge \left ( \boxtimes ^{1}\kappa  \left (\mathbf{\lambda }\right ) \left (x_{j_{n}}\right )\right )\end{align}
\begin{equation}m =1 :p_{1} =n
\end{equation}
\begin{equation} \boxtimes ^{n}\kappa  \left (\mathbf{\lambda }\right ) \left (x_{j_{1}} \wedge \cdots  \wedge x_{j_{n}}\right ) = \boxtimes ^{p_{1}}\kappa  \left (\mathbf{\lambda }\right ) \left (x_{j_{1}} \wedge \cdots  \wedge x_{j_{p_{1}}}\right ) =\left ( \boxtimes ^{1}\kappa  \left (\mathbf{\lambda }\right ) \left (x_{j_{1}}\right )\right ) \wedge \cdots  \wedge \left ( \boxtimes ^{1}\kappa  \left (\mathbf{\lambda }\right ) \left (x_{j_{n}}\right )\right )
\end{equation}
\end{proof}


\end{lemma}

The linear maps $\kappa  \left (\mathbf{\lambda }\right )$ induce algebra maps $ \boxtimes \kappa  \left (\mathbf{\lambda }\right )$
\begin{align} \boxtimes \kappa  \left (\mathbf{\lambda }\right ) &  :\mathcal{F} \rightarrow \mathcal{F} \\
 \boxtimes \kappa  \left (\mathbf{\lambda }\right ) &  :\mathcal{F}_{n} \rightarrow \mathcal{F}_{n}\end{align}in $\mathcal{F}$ i.e.
\begin{equation} \boxtimes \kappa  \left (\mathbf{\lambda }\right ) \left (X \wedge Y\right ) =\left (\kappa  \left (\mathbf{\lambda }\right ) \boxtimes \kappa  \left (\mathbf{\lambda }\right )\right ) \left (X \wedge Y\right ) =\left ( \boxtimes \kappa  \left (\mathbf{\lambda }\right ) \left (X\right )\right ) \wedge \left ( \boxtimes \kappa  \left (\mathbf{\lambda }\right ) \left (Y\right )\right )\text{.}
\end{equation}


\begin{lemma}
The linear maps $\kappa  \left (\mathbf{\lambda }\right )$ induce algebra maps $ \boxtimes \kappa  \left (\mathbf{\lambda }\right )$
\begin{align} \boxtimes \kappa  \left (\mathbf{\lambda }\right ) &  :\mathcal{F} \rightarrow \mathcal{F} \\
 \boxtimes \kappa  \left (\mathbf{\lambda }\right ) &  :\mathcal{F}_{n} \rightarrow \mathcal{F}_{n}\end{align}in $\mathcal{F}$ i.e.
\begin{equation} \boxtimes \kappa  \left (\mathbf{\lambda }\right ) \left (X \wedge Y\right ) = \boxtimes \kappa  \left (\mathbf{\lambda }\right ) \left (X\right )  \wedge  \boxtimes \kappa  \left (\mathbf{\lambda }\right ) \left (Y\right )\text{.}
\end{equation}


\begin{proof}
\begin{equation} \boxtimes \kappa  \left (\mathbf{\lambda }\right ) \left (X_{n} \wedge Y_{m}\right ) =\left (\sum _{0 \leq p \leq 2 r} \boxtimes ^{p}\kappa  \left (\mathbf{\lambda }\right )\right ) \left (X_{n} \wedge Y_{m}\right ) = \boxtimes ^{n +m}\kappa  \left (\mathbf{\lambda }\right ) \left (X_{n} \wedge Y_{m}\right ) = \boxtimes \kappa  \left (\mathbf{\lambda }\right ) \left (X_{n}\right )  \wedge  \boxtimes \kappa  \left (\mathbf{\lambda }\right ) \left (Y_{m}\right )
\end{equation}
\begin{equation} \boxtimes \kappa  \left (\mathbf{\lambda }\right ) \left (X_{n} \wedge Y_{m}\right ) = \boxtimes \kappa  \left (\mathbf{\lambda }\right ) \left (X_{n}\right )  \wedge  \boxtimes \kappa  \left (\mathbf{\lambda }\right ) \left (Y_{m}\right )
\end{equation}We now study whether the maps $ \boxtimes \kappa  \left (\mathbf{\lambda }\right )$ are algebraic wrt the geometric product "$ \cdot ""$ i.e whether
\begin{equation}\kappa  \left (\mathbf{\lambda }\right ) \left (X Y\right ) =\kappa  \left (\mathbf{\lambda }\right ) \left (X\right ) \kappa  \left (\mathbf{\lambda }\right ) \left (Y\right )\text{.}
\end{equation}One can first check the simplest case
\begin{equation}X = \langle X \rangle _{1} ,Y = \langle Y \rangle _{1}\text{,}
\end{equation}which gives
\begin{equation}X_{1} Y_{1} = \langle X_{1} Y_{1} \rangle _{0} +X_{1} \wedge Y_{1} = \langle  \langle X_{1} Y_{1} \rangle  \rangle _{0} I_{\mathcal{F}} +X_{1} \wedge Y_{1} =\left (\left .X_{1}\right \vert Y_{1}\right ) I_{\mathcal{F}} +X_{1} \wedge Y_{1} =\ensuremath{\operatorname*{tr}}\left \{X_{1} Y_{1}\right \} I_{\mathcal{F}} +X_{1} \wedge Y_{1}
\end{equation}and
\begin{align} \boxtimes \kappa  \left (\mathbf{\lambda }\right ) \left (X_{1} Y_{1}\right ) &  = \boxtimes ^{0}\kappa  \left (\mathbf{\lambda }\right )  \langle X_{1} Y_{1} \rangle _{0} + \boxtimes ^{2}\kappa  \left (\mathbf{\lambda }\right ) \left (X_{1} \wedge Y_{1}\right ) = \boxtimes ^{0}\kappa  \left (\mathbf{\lambda }\right ) \left (\left .X_{1}\right \vert Y_{1}\right ) I_{\mathcal{F}} + \boxtimes ^{2}\kappa  \left (\mathbf{\lambda }\right ) \left (X_{1} \wedge Y_{1}\right ) \\
 \boxtimes \kappa  \left (\mathbf{\lambda }\right ) \left (X_{1}\right ) \boxtimes \kappa  \left (\mathbf{\lambda }\right ) \left (Y_{1}\right ) &  = \boxtimes ^{1}\kappa  \left (\mathbf{\lambda }\right ) \left (X_{1}\right )  \boxtimes ^{1}\kappa  \left (\mathbf{\lambda }\right ) \left (Y_{1}\right ) = \langle  \boxtimes ^{1}\kappa  \left (\mathbf{\lambda }\right ) \left (X_{1}\right )  \boxtimes ^{1}\kappa  \left (\mathbf{\lambda }\right ) \left (Y_{1}\right ) \rangle _{0} + \boxtimes ^{1}\kappa  \left (\mathbf{\lambda }\right ) \left (X_{1}\right ) \wedge  \boxtimes ^{1}\kappa  \left (\mathbf{\lambda }\right ) \left (Y_{1}\right ) \\
 &  =\left (\left . \boxtimes ^{1}\kappa  \left (\mathbf{\lambda }\right ) \left (X_{1}\right )\right \vert  \boxtimes ^{1}\kappa  \left (\mathbf{\lambda }\right ) \left (Y_{1}\right )\right ) I_{\mathcal{F}} + \boxtimes ^{1}\kappa  \left (\mathbf{\lambda }\right ) \left (X_{1}\right ) \wedge  \boxtimes ^{1}\kappa  \left (\mathbf{\lambda }\right ) \left (Y_{1}\right ) \\
 &  =\left (\left . \boxtimes ^{1}\kappa  \left (\mathbf{\lambda }\right ) \left (X_{1}\right )\right \vert  \boxtimes ^{1}\kappa  \left (\mathbf{\lambda }\right ) \left (Y_{1}\right )\right ) I_{\mathcal{F}} + \boxtimes ^{2}\kappa  \left (\mathbf{\lambda }\right ) \left (X_{1} \wedge Y_{1}\right )\text{.}\end{align}Hence one can observe that
\begin{align} \boxtimes \kappa  \left (\mathbf{\lambda }\right ) \left (X_{1} Y_{1}\right ) - \boxtimes \kappa  \left (\mathbf{\lambda }\right ) \left (X_{1}\right ) \boxtimes \kappa  \left (\mathbf{\lambda }\right ) \left (Y_{1}\right ) &  =\left \{\overset{}{\left ( \boxtimes ^{0}\kappa  \left (\mathbf{\lambda }\right ) \left (\left .X_{1}\right \vert Y_{1}\right )\right ) -\left (\left . \boxtimes ^{1}\kappa  \left (\mathbf{\lambda }\right ) \left (X_{1}\right )\right \vert  \boxtimes ^{1}\kappa  \left (\mathbf{\lambda }\right ) \left (Y_{1}\right )\right )}\right \} I_{\mathcal{F}} \\
 &  =\left \{\overset{}{\left ( \boxtimes ^{0}\kappa  \left (\mathbf{\lambda }\right ) \left (\left .X_{1}\right \vert Y_{1}\right )\right ) -\left (\left .\kappa  \left (\mathbf{\lambda }\right ) \left (X_{1}\right )\right \vert \kappa  \left (\mathbf{\lambda }\right ) \left (Y_{1}\right )\right )}\right \} I_{\mathcal{F}} \\
 &  =\left \{\overset{}{\left ( \boxtimes ^{0}\kappa  \left (\mathbf{\lambda }\right ) \left (\left .X_{1}\right \vert Y_{1}\right )\right ) -\left (\left .X_{1}\right \vert \kappa  \left (\mathbf{\lambda }\right )^{t} \kappa  \left (\mathbf{\lambda }\right ) Y_{1}\right )}\right \} I_{\mathcal{F}}\end{align}and
\begin{equation}\left ( \boxtimes ^{0}\kappa  \left (\mathbf{\lambda }\right ) \left (\left .X_{1}\right \vert Y_{1}\right )\right ) -\left (\left . \boxtimes ^{1}\kappa  \left (\mathbf{\lambda }\right ) \left (X_{1}\right )\right \vert  \boxtimes ^{1}\kappa  \left (\mathbf{\lambda }\right ) \left (Y_{1}\right )\right ) =\kappa  \left (\mathbf{\lambda }\right )_{0} \sum _{1 \leq j \leq 2 r}X_{1 j} Y_{1 j} -\sum _{1 \leq j \leq 2 r}\left (\kappa  \left (\mathbf{\lambda }\right ) \left (X_{1}\right )\right )_{j} \left (\kappa  \left (\mathbf{\lambda }\right ) \left (Y_{1}\right )\right )_{j}\text{.}
\end{equation}
\begin{equation}\kappa  \left (\mathbf{\lambda }\right )^{t} \kappa  \left (\mathbf{\lambda }\right ) =I_{\mathcal{F}}
\end{equation}so $\kappa  \left (\mathbf{\lambda }\right )$ is an algebraic map if it belongs to $O \left (2 r ,\mathbb{C}\right )$ and $\kappa  \left (\mathbf{\lambda }\right )_{0} =1.$
\begin{align}\kappa  \left (\mathbf{\lambda }\right ) &  \notin \mathcal{F} \\
\kappa  \left (\mathbf{\lambda }\right ) &  =e^{\mathbf{\lambda }} =I_{\mathcal{F}} +\mathbf{\lambda } +\mathbf{\cdots } \\
\ensuremath{\operatorname*{tr}}\left \{e^{\mathbf{\lambda }}\right \} &  =I_{\mathcal{F}} +\frac{1}{2} \ensuremath{\operatorname*{tr}}\left \{\mathbf{\lambda }^{2}\right \} +\frac{1}{4 !} \ensuremath{\operatorname*{tr}}\left \{\mathbf{\lambda }^{4}\right \} +\cdots  \\
\det \left \{e^{\mathbf{\lambda }}\right \} &  =e^{\ensuremath{\operatorname*{tr}}\left \{\mathbf{\lambda }\right \}} =1\text{if}\mathbf{\lambda }^{t} = -\mathbf{\lambda }\end{align}
\end{proof}


\end{lemma}


\begin{remark}
$\kappa  \left (\mathbf{\lambda }\right )$ is a $c^{ \ast }$-algebraic map only if $\kappa  \left (\mathbf{\lambda }\right ) \in O \left (2 r ,\mathbb{R}\right )$ 
\end{remark}


\begin{theorem}
\underline {\emph{Skolem-Noether}}: every \underline {algebra} automorphism of the matrix algebra
$M \left (n ,\mathbb{F}\right )$ is an inner automorphism. 
\end{theorem}

In particular we consider the basis $\left \{\left .\vert \widetilde{a}_{j}^{ \dag } \rangle  ,\vert \widetilde{a}_{j} \rangle \right \vert 1 \leq j \leq 2 r\right \}$ of $\mathcal{F}_{1}$, then
\begin{align}\sum _{1 \leq j ,k \leq 2 r}\kappa  \left (\mathbf{\lambda }\right )_{j k} \vert \mathfrak{e}_{j} \rangle _{\ensuremath{\operatorname*{tr}}} \langle \mathfrak{e}_{k}\vert  &  =\sum _{1 \leq j ,k \leq r}\kappa  \left (\mathbf{\lambda }\right )_{j k} \vert \widetilde{a}_{j}^{ \dag } \rangle _{\ensuremath{\operatorname*{tr}}}  \langle \widetilde{a}_{k}^{ \dag }\vert  +\sum _{1 \leq j ,k \leq r}\kappa  \left (\mathbf{\lambda }\right )_{j k +r} \vert \widetilde{a}_{j}^{ \dag } \rangle _{\ensuremath{\operatorname*{tr}}} \langle \widetilde{a}_{k}\vert  \\
 &  +\sum _{1 \leq j ,k \leq r j +r k}\kappa  \left (\mathbf{\lambda }\right )_{j +r k} \vert \widetilde{a}_{j} \rangle _{\ensuremath{\operatorname*{tr}}}  \langle \widetilde{a}_{k}^{ \dag }\vert  +\sum _{1 \leq j ,k \leq r}\kappa  \left (\mathbf{\lambda }\right )_{j +r k +r} \vert \widetilde{a}_{j} \rangle _{\ensuremath{\operatorname*{tr}}} \langle \widetilde{a}_{k}\vert \end{align}giving
\begin{align}\sum _{1 \leq j ,k \leq 2 r}\kappa  \left (\mathbf{\lambda }\right )_{j k} \vert \mathfrak{e}_{j} \rangle _{\ensuremath{\operatorname*{tr}}} \langle \mathfrak{e}_{k}\vert  &  =2 \sum _{1 \leq j ,k \leq r}\kappa  \left (\mathbf{\lambda }\right )_{j k} \vert a_{j}^{ \dag } \rangle _{\ensuremath{\operatorname*{tr}}}  \langle a_{k}^{ \dag }\vert  +2 \sum _{1 \leq j ,k \leq r}\kappa  \left (\mathbf{\lambda }\right )_{j k +r} \vert a_{j}^{ \dag } \rangle _{\ensuremath{\operatorname*{tr}}} \langle a_{k}\vert  \\
 &  +2 \sum _{1 \leq j ,k \leq r j +r k}\kappa  \left (\mathbf{\lambda }\right )_{j +r k} \vert a_{j} \rangle _{\ensuremath{\operatorname*{tr}}}  \langle a_{k}^{ \dag }\vert  +2 \sum _{1 \leq j ,k \leq r}\kappa  \left (\mathbf{\lambda }\right )_{j +r k +r} \vert a_{j} \rangle _{\ensuremath{\operatorname*{tr}}} \langle a_{k}\vert \text{.}\end{align}

The $\kappa  \left (\mathbf{\lambda }\right )$ transformation is defined by
\begin{equation}\left (\begin{array}{cc}\mathbf{x}_{1} & \mathbf{x}_{2}\end{array}\right ) =\kappa  \left (\mathbf{\lambda }\right ) \left (\begin{array}{cc}\mathbf{a}^{ \dag } & \mathbf{a}\end{array}\right ) =\left (\begin{array}{cc}\mathbf{a}^{ \dag } & \mathbf{a}\end{array}\right ) \mathbb{K} \left (\mathbf{\lambda }\right ) =\left (\begin{array}{cc}\mathbf{a}^{ \dag } & \mathbf{a}\end{array}\right ) \mathcal{R}_{a} \left (\kappa  \left (\mathbf{\lambda }\right )\right )
\end{equation}the $K$ transformation
\begin{equation}\mathbb{K} =\left (\begin{array}{cc}\frac{1}{\sqrt{2}} \mathbb{I}_{r} & \frac{i}{\sqrt{2}} \mathbb{I}_{r} \\
\frac{1}{\sqrt{2}} \mathbb{I}_{r} &  -\frac{i}{\sqrt{2}} \mathbb{I}_{r}\end{array}\right )
\end{equation}is a special case. 


\begin{remark}
\begin{equation}\left (\begin{array}{cc}\widetilde{\mathbf{x}}_{1} & \widetilde{\mathbf{x}}_{2}\end{array}\right ) =\kappa  \left (\mathbf{\lambda }\right ) \left (\begin{array}{cc}\widetilde{\mathbf{a}}^{ \dag } & \widetilde{\mathbf{a}}\end{array}\right ) =\left (\begin{array}{cc}\widetilde{\mathbf{a}}^{ \dag } & \widetilde{\mathbf{a}}\end{array}\right ) \mathbb{K} \left (\mathbf{\lambda }\right ) =\left (\begin{array}{cc}\widetilde{\mathbf{a}}^{ \dag } & \widetilde{\mathbf{a}}\end{array}\right ) \mathcal{R}_{\widetilde{a}} \left (\kappa  \left (\mathbf{\lambda }\right )\right )
\end{equation}\emph{AND}
\begin{equation}\mathcal{R}_{a} \left (\kappa  \left (\mathbf{\lambda }\right )\right ) =\mathcal{R}_{\widetilde{a}} \left (\kappa  \left (\mathbf{\lambda }\right )\right ) =\mathbb{K} \left (\mathbf{\lambda }\right )\text{.}
\end{equation}
\end{remark}


\begin{lemma}
The transformation $\kappa  \left (\mathbf{\lambda }\right )$ produces a self adjoint basis from $\left (\begin{array}{cc}\widetilde{\mathbf{a}}^{ \dag } & \widetilde{\mathbf{a}}\end{array}\right )$ $\ensuremath{\operatorname*{Iff}}$
\begin{equation}\mathcal{R}_{a} \left (\kappa  \left (\mathbf{\lambda }\right )\right ) =\left (\begin{array}{cc}\kappa _{1 1} \left (\mathbf{\lambda }\right ) & \overline{\kappa _{2 2}} \left (\mathbf{\lambda }\right ) \\
\overline{\kappa _{1 1}} \left (\mathbf{\lambda }\right ) & \kappa _{2 2} \left (\mathbf{\lambda }\right )\end{array}\right )\text{.}
\end{equation}
\end{lemma}

In order that
\begin{align}\mathbf{x}_{1} &  =\mathbf{x}_{1}^{ \dag } \\
\mathbf{x}_{2} &  =\mathbf{x}_{2}^{ \dag }\end{align}from the transformation
\begin{equation}\left (\begin{array}{cc}\mathbf{a}^{ \dag } & \mathbf{a}\end{array}\right ) \left (\begin{array}{cc}\kappa _{1 1} \left (\mathbf{\lambda }\right ) & \kappa _{1 2} \left (\mathbf{\lambda }\right ) \\
\kappa _{2 1} \left (\mathbf{\lambda }\right ) & \kappa _{2 2} \left (\mathbf{\lambda }\right )\end{array}\right ) =\left (\begin{array}{cc}\mathbf{x}_{1} & \mathbf{x}_{2}\end{array}\right )
\end{equation}one must have that
\begin{align}\mathbf{a}^{ \dag } \kappa _{1 1} \left (\mathbf{\lambda }\right ) +\mathbf{a} \kappa _{2 1} \left (\mathbf{\lambda }\right ) &  =\mathbf{a} \overline{\kappa _{1 1}} \left (\mathbf{\lambda }\right ) +\mathbf{a}^{ \dag } \overline{\kappa _{2 1}} \left (\mathbf{\lambda }\right ) \Rightarrow \kappa _{1 1} \left (\mathbf{\lambda }\right ) =\overline{\kappa _{2 1}} \left (\mathbf{\lambda }\right ) ;\kappa _{2 1} \left (\mathbf{\lambda }\right ) =\overline{\kappa _{1 1}} \left (\mathbf{\lambda }\right ) \\
\mathbf{a}^{ \dag } \kappa _{1 2} \left (\mathbf{\lambda }\right ) +\mathbf{a} \kappa _{2 2} \left (\mathbf{\lambda }\right ) &  =\mathbf{a} \overline{\kappa _{1 2}} \left (\mathbf{\lambda }\right ) +\mathbf{a}^{ \dag } \overline{\kappa _{2 2}} \left (\mathbf{\lambda }\right ) \Rightarrow \kappa _{1 2} \left (\mathbf{\lambda }\right ) =\overline{\kappa _{2 2}} \left (\mathbf{\lambda }\right ) ;\kappa _{2 2} \left (\mathbf{\lambda }\right ) =\overline{\kappa _{1 2}} \left (\mathbf{\lambda }\right )\text{.}\end{align}giving
\begin{equation}\mathbb{K} \left (\mathbf{\lambda }\right ) =\left (\begin{array}{cc}\kappa _{1 1} \left (\mathbf{\lambda }\right ) & \kappa _{1 2} \left (\mathbf{\lambda }\right ) \\
\kappa _{2 1} \left (\mathbf{\lambda }\right ) & \kappa _{2 2} \left (\mathbf{\lambda }\right )\end{array}\right ) =\left (\begin{array}{cc}\kappa _{1 1} \left (\mathbf{\lambda }\right ) & \overline{\kappa _{2 2}} \left (\mathbf{\lambda }\right ) \\
\overline{\kappa _{1 1}} \left (\mathbf{\lambda }\right ) & \kappa _{2 2} \left (\mathbf{\lambda }\right )\end{array}\right )\text{.}
\end{equation}

$\kappa  \left (\mathbf{\lambda }\right )$ is a $ \ast $-map if 


\begin{lemma}
$\kappa  \left (\mathbf{\lambda }\right )$ is a $ \ast $-map i.e.
\begin{equation}\kappa  \left (\mathbf{\lambda }\right ) \left (a_{j}^{ \dag }\right ) =\left (\kappa  \left (\mathbf{\lambda }\right ) \left (a_{j}\right )\right )^{ \dag }
\end{equation}$\ensuremath{\operatorname*{Iff}}$
\begin{equation}\mathcal{R}_{a} \left (\kappa  \left (\mathbf{\lambda }\right )\right ) =\left (\begin{array}{cc}\kappa _{1 1} \left (\mathbf{\lambda }\right ) & \kappa _{1 2} \left (\mathbf{\lambda }\right ) \\
\overline{\kappa _{1 2}} \left (\mathbf{\lambda }\right ) & \overline{\kappa _{1 1}} \left (\mathbf{\lambda }\right )\end{array}\right )
\end{equation}


\begin{proof}
\begin{equation}\sum _{1 \leq j \leq r}\left \{a_{j}^{ \dag } \overline{\kappa  \left (\mathbf{\lambda }\right )}_{j +r l} +a_{j} \overline{\kappa  \left (\mathbf{\lambda }\right )}_{j l}\right \} =\sum _{1 \leq j \leq r}\left \{a_{j}^{ \dag } \kappa  \left (\mathbf{\lambda }\right )_{j l +r} +a_{j} \kappa  \left (\mathbf{\lambda }\right )_{j +r l +r}\right \}
\end{equation}i.e.
\begin{equation}\mathcal{R}_{a} \left (\kappa  \left (\mathbf{\lambda }\right )\right ) =\left (\begin{array}{cc}\kappa _{1 1} \left (\mathbf{\lambda }\right ) & \kappa _{1 2} \left (\mathbf{\lambda }\right ) \\
\overline{\kappa _{1 2}} \left (\mathbf{\lambda }\right ) & \overline{\kappa _{1 1}} \left (\mathbf{\lambda }\right )\end{array}\right )\text{.}
\end{equation}
\end{proof}


\end{lemma}


\begin{proposition}
$\kappa  \left (\mathbf{\lambda }\right )$ cannot be invertible i.e. $ \notin G L \left (\mathcal{F}_{1}\right )$, a $ \ast $-map and produce a self adjoint basis 


\begin{proof}
If $\kappa  \left (\mathbf{\lambda }\right )$ is a $ \ast $-map and produces a self adjoint basis then
\begin{equation}\mathcal{R}_{a} \left (\kappa  \left (\mathbf{\lambda }\right )\right ) =\left (\begin{array}{cc}\kappa _{1 1} \left (\mathbf{\lambda }\right ) & \overline{\kappa _{2 2}} \left (\mathbf{\lambda }\right ) \\
\overline{\kappa _{1 1}} \left (\mathbf{\lambda }\right ) & \kappa _{2 2} \left (\mathbf{\lambda }\right )\end{array}\right ) =\left (\begin{array}{cc}\kappa _{1 1} \left (\mathbf{\lambda }\right ) & \kappa _{1 2} \left (\mathbf{\lambda }\right ) \\
\overline{\kappa _{1 2}} \left (\mathbf{\lambda }\right ) & \overline{\kappa _{1 1}} \left (\mathbf{\lambda }\right )\end{array}\right )\text{,}
\end{equation}which shows that
\begin{equation}\mathcal{R}_{a} \left (\kappa  \left (\mathbf{\lambda }\right )\right ) =\left (\begin{array}{cc}\kappa _{1 1} \left (\mathbf{\lambda }\right ) & \kappa _{1 1} \left (\mathbf{\lambda }\right ) \\
\overline{\kappa _{1 1}} \left (\mathbf{\lambda }\right ) & \overline{\kappa _{1 1}} \left (\mathbf{\lambda }\right )\end{array}\right )\text{.}
\end{equation}However $\mathcal{R}_{a} \left (\kappa  \left (\mathbf{\lambda }\right )\right )$ is clearly not injective as
\begin{equation}\left (\begin{array}{cc}\mathbf{a}^{ \dag } & \mathbf{a}\end{array}\right ) \left (\begin{array}{cc}\kappa _{1 1} \left (\mathbf{\lambda }\right ) & \kappa _{1 1} \left (\mathbf{\lambda }\right ) \\
\overline{\kappa _{1 1}} \left (\mathbf{\lambda }\right ) & \overline{\kappa _{1 1}} \left (\mathbf{\lambda }\right )\end{array}\right ) =\left (\begin{array}{cc}\mathbf{x}_{1} & \mathbf{x}_{1}\end{array}\right )
\end{equation}thus
\begin{equation}\dim  \left \{\ker  \left \{\mathcal{R}_{a} \left (\kappa  \left (\mathbf{\lambda }\right )\right )\right \}\right \} =r\text{.}
\end{equation}
\end{proof}


\end{proposition}

The action of $\kappa  \left (\mathbf{\lambda }\right )$ on basis elements is given by
\begin{equation}\kappa  \left (\mathbf{\lambda }\right ) \left (a_{l}^{ \dag }\right ) =\sum _{1 \leq j \leq r}\left \{a_{j}^{ \dag } \kappa  \left (\mathbf{\lambda }\right )_{j l} +a_{j} \kappa  \left (\mathbf{\lambda }\right )_{j +r l}\right \}
\end{equation}and
\begin{equation}\kappa  \left (\mathbf{\lambda }\right ) \left (a_{l}\right ) =\sum _{1 \leq j \leq r}\left \{a_{j}^{ \dag } \kappa  \left (\mathbf{\lambda }\right )_{j l +r} +a_{j} \kappa  \left (\mathbf{\lambda }\right )_{j +r l +r}\right \}\text{,}
\end{equation}which can also be expressed in more detail as 


\begin{enumerate}
\item
\begin{equation}\sum _{1 \leq j ,\kappa  \leq r}\kappa  \left (\mathbf{\lambda }\right )_{j \kappa } \vert a_{j}^{ \dag } \rangle _{\ensuremath{\operatorname*{tr}}} \langle a_{\kappa }^{ \dag }\vert a_{l}^{ \dag } =\sum _{1 \leq j ,\kappa  \leq r}\kappa  \left (\mathbf{\lambda }\right )_{j \kappa }\vert a_{j}^{ \dag } \rangle \ensuremath{\operatorname*{tr}} \left \{a_{\kappa } a_{l}^{ \dag }\right \} =\frac{1}{2} \sum _{1 \leq j \leq r}\vert a_{j}^{ \dag } \rangle  \kappa  \left (\mathbf{\lambda }\right )_{j l}
\end{equation}

\item
\begin{equation}\sum _{1 \leq j ,\kappa  \leq r}\kappa  \left (\mathbf{\lambda }\right )_{j \kappa  +r} \vert a_{j}^{ \dag } \rangle _{\ensuremath{\operatorname*{tr}}} \langle a_{\kappa }\vert a_{l}^{ \dag } =\sum _{1 \leq j ,\kappa  \leq r}\kappa  \left (\mathbf{\lambda }\right )_{j \kappa  +r}\vert a_{j}^{ \dag } \rangle \ensuremath{\operatorname*{tr}} \left \{a_{\kappa }^{ \dag } a_{l}^{ \dag }\right \} =0
\end{equation}

\item
\begin{equation}\sum _{1 \leq j ,\kappa  \leq r}\kappa  \left (\mathbf{\lambda }\right )_{j +r \kappa } \vert a_{j} \rangle _{\ensuremath{\operatorname*{tr}}} \langle a_{\kappa }^{ \dag }\vert a_{l}^{ \dag } =\sum _{1 \leq j ,\kappa  \leq r}\kappa  \left (\mathbf{\lambda }\right )_{j +r \kappa }\vert a_{j} \rangle \ensuremath{\operatorname*{tr}} \left \{a_{\kappa } a_{l}^{ \dag }\right \} =\frac{1}{2} \sum _{1 \leq j \leq r}\vert a_{j} \rangle  \kappa  \left (\mathbf{\lambda }\right )_{j +r l}
\end{equation}

\item
\begin{equation}\sum _{1 \leq j ,\kappa  \leq r}\kappa  \left (\mathbf{\lambda }\right )_{j +r \kappa  +r} \vert a_{j} \rangle _{\ensuremath{\operatorname*{tr}}} \langle a_{\kappa }\vert a_{l}^{ \dag } =\sum _{1 \leq j ,\kappa  \leq r}\kappa  \left (\mathbf{\lambda }\right )_{j +r \kappa  +r}\vert a_{j} \rangle \ensuremath{\operatorname*{tr}} \left \{a_{\kappa }^{ \dag } a_{l}^{ \dag }\right \} =0.
\end{equation}

\item One can use these relationships to again study the conditions
for a $ \ast $-map \end{enumerate}


\begin{equation}\left (\kappa  \left (\mathbf{\lambda }\right ) \left (a_{l}^{ \dag }\right )\right )^{ \dag } =\sum _{1 \leq j \leq r}\left \{a_{j} \overline{\kappa  \left (\mathbf{\lambda }\right )}_{j l} +a_{j}^{ \dag } \overline{\kappa  \left (\mathbf{\lambda }\right )}_{j +r l}\right \} =\sum _{1 \leq j \leq r}\left \{a_{j}^{ \dag } \overline{\kappa  \left (\mathbf{\lambda }\right )}_{j +r l} +a_{j} \overline{\kappa  \left (\mathbf{\lambda }\right )}_{j l}\right \}\text{.}
\end{equation}


\begin{proposition}
Is $K$ an inner automorphism even though it cannot be generated by the Lie algebra $\mathcal{A}_{2}\text{?}$ It is given by
\begin{equation}K \left (\mathbf{\lambda }\right ) =\exp  \left \{i \sum _{1 \leq j ,k \leq 2 r}\lambda _{j k} \vert \mathfrak{e}_{j} \rangle _{\ensuremath{\operatorname*{tr}}} \langle \mathfrak{e}_{k}\vert \right \} =\sum _{1 \leq j ,k \leq 2 r}K \left (\mathbf{\lambda }\right )_{j k} \vert \mathfrak{e}_{j} \rangle _{\ensuremath{\operatorname*{tr}}} \langle \mathfrak{e}_{k}\vert 
\end{equation}for a given $\mathbf{\lambda }\text{.}$ 
\end{proposition}


\begin{align}\lambda _{1} &  =\sum _{1 \leq l ,m \leq r}\lambda _{1 l m} K \left (\mathbf{\lambda }\right ) \left (a_{l}^{ \dag }\right ) K \left (\mathbf{\lambda }\right ) \left (a_{m}\right ) =\sum _{1 \leq l ,m \leq r}\lambda _{1 l m} \sum _{1 \leq j \leq r}\left \{a_{j}^{ \dag } K \left (\mathbf{\lambda }\right )_{j l} +a_{j} K \left (\mathbf{\lambda }\right )_{j +r l}\right \} \sum _{1 \leq k \leq r}\left \{a_{k}^{ \dag } K \left (\mathbf{\lambda }\right )_{k m +r} +a_{k} K \left (\mathbf{\lambda }\right )_{k +r m +r}\right \} \\
 &  =\sum _{1 \leq j ,k ,l ,m \leq r}\lambda _{1 l m} \left \{a_{j}^{ \dag } K \left (\mathbf{\lambda }\right )_{j l} +a_{j} K \left (\mathbf{\lambda }\right )_{j +r l}\right \} \left \{a_{k}^{ \dag } K \left (\mathbf{\lambda }\right )_{k m +r} +a_{k} K \left (\mathbf{\lambda }\right )_{k +r m +r}\right \} \\
 &  =\sum _{1 \leq j ,k ,l ,m \leq r}\lambda _{1 l m} \left \{a_{j}^{ \dag } a_{k}^{ \dag } K \left (\mathbf{\lambda }\right )_{j l} K \left (\mathbf{\lambda }\right )_{k m +r} +a_{j}^{ \dag } a_{k} K \left (\mathbf{\lambda }\right )_{j l} K \left (\mathbf{\lambda }\right )_{k +r m +r} +a_{j} a_{k}^{ \dag } K \left (\mathbf{\lambda }\right )_{j +r l} K \left (\mathbf{\lambda }\right )_{k m +r} +a_{j} a_{k} K \left (\mathbf{\lambda }\right )_{j +r l} K \left (\mathbf{\lambda }\right )_{k +r m +r}\right \} \\
 &  =\sum _{1 \leq l ,m \leq r}\sum _{1 \leq j <k \leq r\text{,}}\lambda _{1 l m} \left \{a_{j}^{ \dag } a_{k}^{ \dag } \left (K \left (\mathbf{\lambda }\right )_{j l} K \left (\mathbf{\lambda }\right )_{k m +r} -K \left (\mathbf{\lambda }\right )_{k l} K \left (\mathbf{\lambda }\right )_{j m +r}\right ) +a_{j} a_{k} \left (K \left (\mathbf{\lambda }\right )_{j +r l} K \left (\mathbf{\lambda }\right )_{k +r m +r} -K \left (\mathbf{\lambda }\right )_{k +r l} K \left (\mathbf{\lambda }\right )_{j +r m +r}\right )\right \} \\
 &  +\sum _{1 \leq l ,m \leq r}\sum _{1 \leq j ,k \leq r\text{,}}\lambda _{1 l m} a_{j}^{ \dag } a_{k} \left (K \left (\mathbf{\lambda }\right )_{j l} K \left (\mathbf{\lambda }\right )_{k +r m +r} -K \left (\mathbf{\lambda }\right )_{k +r l} K \left (\mathbf{\lambda }\right )_{j m +r}\right ) \\
 &  +\sum _{1 \leq l ,m \leq r}\sum _{1 \leq j \leq r}\lambda _{1 l m} \left \{a_{j}^{ \dag } a_{j} \left (K \left (\mathbf{\lambda }\right )_{j l} K \left (\mathbf{\lambda }\right )_{j +r m +r} -K \left (\mathbf{\lambda }\right )_{j +r l} K \left (\mathbf{\lambda }\right )_{j m +r}\right )\right \} +\sum _{1 \leq j ,l ,m \leq r}\lambda _{1 l m} K \left (\mathbf{\lambda }\right )_{j +r l} K \left (\mathbf{\lambda }\right )_{j m +r} I_{r}\end{align}
\begin{equation}K \in \mathcal{R}_{\mathcal{F}} \left (U \left (2 r\right )\right )
\end{equation}
\begin{equation} \langle \left .K \left (A\right )\right \vert K \left (B\right ) \rangle _{\mathcal{F}_{H S}} = \langle \left .A\right \vert B \rangle _{\mathcal{F}_{H S}}
\end{equation}
\begin{equation}\ensuremath{\operatorname*{tr}}\left \{K \left (A\right )\right \} = \langle \left .K \left (I_{\mathcal{F}}\right )\right \vert K \left (A\right ) \rangle _{\mathcal{F}_{H S}} =\ensuremath{\operatorname*{tr}}\left \{A\right \}
\end{equation}
\begin{align} & \frac{1}{2} \sum _{1 \leq l ,m \leq r}\sum _{1 \leq j <k \leq r\text{,}}\lambda _{1 l m} \left \{a_{j}^{ \dag } a_{k}^{ \dag } \left (\delta _{j l} i \delta _{k m} -i \delta _{k l} \delta _{j m}\right ) +a_{j} a_{k} \left (\delta _{j l} \left ( -i\right ) \delta _{k m} -\delta _{k l} \left ( -i\right ) \delta _{j m}\right )\right \} \\
 &  +\frac{1}{2} \sum _{1 \leq l ,m \leq r}\sum _{1 \leq j ,k \leq r\text{,}}\lambda _{1 l m} a_{j}^{ \dag } a_{k} \left (K \left (\mathbf{\lambda }\right )_{j l} K \left (\mathbf{\lambda }\right )_{k +r m +r} -K \left (\mathbf{\lambda }\right )_{k +r l} K \left (\mathbf{\lambda }\right )_{j m +r}\right ) \\
 &  +\frac{1}{2} \sum _{1 \leq l ,m \leq r}\sum _{1 \leq j \leq r}\lambda _{1 l m} \left \{a_{j}^{ \dag } a_{j} \left (K \left (\mathbf{\lambda }\right )_{j l} K \left (\mathbf{\lambda }\right )_{j +r m +r} -K \left (\mathbf{\lambda }\right )_{j +r l} K \left (\mathbf{\lambda }\right )_{j m +r}\right )\right \} +\sum _{1 \leq j ,l ,m \leq r}\lambda _{1 l m} \left ( -i\right ) \delta _{j l} \delta _{j m} I_{r}\end{align}
\begin{equation}\ensuremath{\operatorname*{Tr}}_{\mathcal{F}} \left \{A\right \} =\sum _{1 \leq j \leq 2^{2 r}} \langle \left .\mathfrak{E}_{j}\right \vert A \mathfrak{E}_{j} \rangle _{\mathcal{F}_{H S}} =\sum _{1 \leq j \leq 2^{2 r}}\ensuremath{\operatorname*{tr}}\left \{\mathfrak{E}_{j}^{ \dag } A \mathfrak{E}_{j}\right \} =2^{2 r} \ensuremath{\operatorname*{tr}}\left \{A\right \}
\end{equation}


\begin{lemma}
$\ensuremath{\operatorname*{Tr}}_{\mathcal{F}} \left \{A\right \}$ is basis independent. 
\end{lemma}

Consider the dual pair
\begin{equation} \langle \mathcal{B} \left (\Phi _{\ensuremath{\operatorname*{tr}}} \left (\mathcal{F}\right )\right ) ,\mathcal{B}_{1} \left (\Phi _{\ensuremath{\operatorname*{tr}}} \left (\mathcal{F}\right )\right ) \rangle  \equiv  \langle \mathcal{F} ,\mathcal{F}_{T C} \rangle 
\end{equation}$\mathcal{B}_{1} \left (\Phi _{\ensuremath{\operatorname*{tr}}} \left (\mathcal{F}\right )\right )$ is a two sided ideal in $\mathcal{B} \left (\Phi _{\ensuremath{\operatorname*{tr}}} \left (\mathcal{F}\right )\right )$ which is in general not closed wrt $\left \Vert \text{}\right \Vert _{\mathcal{B}_{1} \left (\Phi _{\ensuremath{\operatorname*{tr}}} \left (\mathcal{F}\right )\right )}\text{.}$ If $A \in \mathcal{B}_{1} \left (\Phi _{\ensuremath{\operatorname*{tr}}} \left (\mathcal{F}\right )\right )$ then $A^{ -1} \in \mathcal{B} \left (\Phi _{\ensuremath{\operatorname*{tr}}} \left (\mathcal{F}\right )\right )$ if it exists$\text{,}$ \emph{BUT in general }$A^{ -1} \notin \mathcal{B}_{1} \left (\Phi _{\ensuremath{\operatorname*{tr}}} \left (\mathcal{F}\right )\right )\text{.}$\emph{\ }
\begin{align}K \left (\lambda _{1}\right ) &  =\sum _{1 \leq l ,m \leq r}\lambda _{1 l m} K \left (a_{l}^{ \dag }\right ) K \left (a_{m}\right ) =\frac{i}{2} \sum _{1 \leq l ,m \leq r}\lambda _{1 l m} \left (a_{l}^{ \dag } +a_{l}\right ) \left (a_{m}^{ \dag } -a_{m}\right ) =\frac{i}{2} \sum _{1 \leq l ,m \leq r}\lambda _{1 l m} \left (a_{l}^{ \dag } a_{m}^{ \dag } -a_{l}^{ \dag } a_{m} +a_{l} a_{m}^{ \dag } -a_{l} a_{m}\right ) \\
 &  =i \sum _{1 \leq l <m \leq r}\left (\lambda _{1 l m} -\lambda _{1 m l}\right ) \left (a_{l}^{ \dag } a_{m}^{ \dag } -a_{l} a_{m} -a_{l}^{ \dag } a_{m}\right ) +\frac{i}{2} \sum _{1 \leq l \leq r}\lambda _{1 l l} \left (a_{l} a_{l}^{ \dag } -a_{l}^{ \dag } a_{l}\right ) \\
 &  =i \sum _{1 \leq l <m \leq r}\left (\lambda _{1 l m} -\lambda _{1 m l}\right ) \left (a_{l}^{ \dag } a_{m}^{ \dag } -a_{l} a_{m} -a_{l}^{ \dag } a_{m}\right ) +\frac{i}{2} \sum _{1 \leq l \leq r}\lambda _{1 l l} \left (I_{r} -2 a_{l}^{ \dag } a_{l}\right ) \\
 &  =i \sum _{1 \leq l <m \leq r}\left (\lambda _{1 l m} -\lambda _{1 m l}\right ) \left (a_{l}^{ \dag } a_{m}^{ \dag } -a_{l} a_{m} -a_{l}^{ \dag } a_{m}\right ) +\frac{i}{2} \sum _{1 \leq l \leq r}\lambda _{1 l l} I_{r} -i \sum _{1 \leq l \leq r}\lambda _{1 l m} \left (a_{l}^{ \dag } a_{l}\right )\end{align}
\begin{equation}\left (\begin{array}{cc}\frac{1}{\sqrt{2}} \mathbf{\xi } & \frac{i}{\sqrt{2}} \mathbf{\xi } \\
\frac{1}{\sqrt{2}} \overline{\mathbf{\xi }} &  -\frac{i}{\sqrt{2}} \overline{\mathbf{\xi }}\end{array}\right )^{ \dag } \left (\begin{array}{cc}\frac{1}{\sqrt{2}} \mathbf{\xi } & \frac{i}{\sqrt{2}} \mathbf{\xi } \\
\frac{1}{\sqrt{2}} \overline{\mathbf{\xi }} &  -\frac{i}{\sqrt{2}} \overline{\mathbf{\xi }}\end{array}\right ) =\left (\begin{array}{cc}\frac{1}{\sqrt{2}} \overline{\mathbf{\xi }} & \frac{1}{\sqrt{2}} \mathbf{\xi } \\
 -\frac{i}{\sqrt{2}} \overline{\mathbf{\xi }} & \frac{i}{\sqrt{2}} \mathbf{\xi }\end{array}\right ) \left (\begin{array}{cc}\frac{1}{\sqrt{2}} \mathbf{\xi } & \frac{i}{\sqrt{2}} \mathbf{\xi } \\
\frac{1}{\sqrt{2}} \overline{\mathbf{\xi }} &  -\frac{i}{\sqrt{2}} \overline{\mathbf{\xi }}\end{array}\right ) =\left (\begin{array}{cc}\left \vert \mathbf{\xi }\right \vert ^{2} & \mathbf{0} \\
0 & \left \vert \mathbf{\xi }\right \vert ^{2}\end{array}\right )
\end{equation}
\begin{align}K_{\mathbf{\xi }} \left (\lambda _{1}\right ) &  =\sum _{1 \leq l ,m \leq r}\lambda _{1 l m} K_{\mathbf{\xi }} \left (a_{l}^{ \dag }\right ) K_{\mathbf{\xi }} \left (a_{m}\right ) =\frac{i}{2} \sum _{1 \leq l ,m \leq r}\lambda _{1 l m} \left (\xi _{l} a_{l}^{ \dag } +\overline{\xi _{l}} a_{l}\right ) \left (\xi _{m} a_{m}^{ \dag } -\overline{\xi _{m}} a_{m}\right ) \\
 &  =\frac{i}{2} \sum _{1 \leq l ,m \leq r}\lambda _{1 l m} \left (\xi _{l} \xi _{m} a_{l}^{ \dag } a_{m}^{ \dag } -\xi _{l} \overline{\xi _{m}} a_{l}^{ \dag } a_{m} +\overline{\xi _{l}} \xi _{m} a_{l} a_{m}^{ \dag } -\overline{\xi _{l}} \overline{\xi _{m}} a_{l} a_{m}\right ) \\
 &  =i \sum _{1 \leq l <m \leq r}\left (\lambda _{1 l m} -\lambda _{1 m l}\right ) \left (\xi _{l} \xi _{m} a_{l}^{ \dag } a_{m}^{ \dag } -\overline{\xi _{l}} \overline{\xi _{m}} a_{l} a_{m} -\xi _{l} \overline{\xi _{m}} a_{l}^{ \dag } a_{m}\right ) +\frac{i}{2} \sum _{1 \leq l \leq r}\lambda _{1 l l} \left \vert \xi _{l}\right \vert ^{2} \left (a_{l} a_{l}^{ \dag } -a_{l}^{ \dag } a_{l}\right ) \\
 &  =i \sum _{1 \leq l <m \leq r}\left (\lambda _{1 l m} -\lambda _{1 m l}\right ) \left (\xi _{l} \xi _{m} a_{l}^{ \dag } a_{m}^{ \dag } -\overline{\xi _{l}} \overline{\xi _{m}} a_{l} a_{m} -\xi _{l} \overline{\xi _{m}} a_{l}^{ \dag } a_{m}\right ) +\frac{i}{2} \sum _{1 \leq l \leq r}\lambda _{1 l l} \left \vert \xi _{l}\right \vert ^{2} \left (I_{r} -2 a_{l}^{ \dag } a_{l}\right ) \\
 &  =i \sum _{1 \leq l <m \leq r}\left (\lambda _{1 l m} -\lambda _{1 m l}\right ) \left (\xi _{l} \xi _{m} a_{l}^{ \dag } a_{m}^{ \dag } -\overline{\xi _{l}} \overline{\xi _{m}} a_{l} a_{m} -\xi _{l} \overline{\xi _{m}} a_{l}^{ \dag } a_{m}\right ) +\frac{i}{2} \sum _{1 \leq l \leq r}\lambda _{1 l l} \left \vert \xi _{l}\right \vert ^{2} I_{r} -i \sum _{1 \leq l \leq r}\lambda _{1 l m} \left \vert \xi _{l}\right \vert ^{2} \left (a_{l}^{ \dag } a_{l}\right )\end{align}If $\mathbf{\xi } \in \mathcal{B}_{1} \left (\mathbb{C}^{r}\right )$ then $K_{\mathbf{\xi }}$ is bounded and well defined on elements of $\mathcal{B} \left (\mathbb{C}^{r}\right )$ and $\mathcal{B} \left (\mathcal{H}_{a}^{1}\right )$. \#\#\#\#\#
\begin{align}\left (\xi _{l} a_{l}^{ \dag } +\overline{\xi _{l}} a_{l}\right ) \left (\xi _{l} a_{l}^{ \dag } +\overline{\xi _{l}} a_{l}\right ) &  =\left (\xi _{l} a_{l}^{ \dag } \xi _{l} a_{l}^{ \dag } +\overline{\xi _{l}} \xi _{l} a_{l}^{ \dag } a_{l} +\overline{\xi _{l}} a_{l} \xi _{l} a_{l}^{ \dag } +\overline{\xi _{l}} \overline{\xi _{l}} a_{l} a_{l}\right ) \\
 &  =\left \vert \xi _{l}\right \vert ^{2} \left (a_{l}^{ \dag } a_{l} +a_{l} a_{l}^{ \dag }\right ) =\left \vert \xi _{l}\right \vert ^{2} I_{r} =x_{\xi  l}^{2}\end{align}
\begin{equation}\left (\mathbf{x}_{1} \left (\mathbf{\xi }\right ) ,\mathbf{x}_{2} \left (\mathbf{\xi }\right )\right ) =\left (\mathbf{a}^{ \dag } ,\mathbf{a}\right )\left (\begin{array}{cc}\mathbf{\xi }_{r} & \mathbb{O}_{r} \\
\mathbb{O}_{r} & \overline{\mathbf{\xi }_{r}}\end{array}\right ) \left (\begin{array}{cc}\frac{1}{\sqrt{2}} \mathbb{I}_{r} & \frac{i}{\sqrt{2}} \mathbb{I}_{r} \\
\frac{1}{\sqrt{2}} \mathbb{I}_{r} &  -\frac{i}{\sqrt{2}} \mathbb{I}_{r}\end{array}\right )
\end{equation}
\begin{equation}\left (\mathbf{a}^{ \dag } ,\mathbf{a}\right ) \equiv \left (\vert \mathbf{a}^{ \dag } \rangle  ,\vert \mathbf{a} \rangle \right )
\end{equation}
\begin{align}v &  =\left (\mathbf{a}^{ \dag } ,\mathbf{a}\right )\left (\begin{array}{cc}\mathbf{\xi }_{r} & \mathbb{O}_{r} \\
\mathbb{O}_{r} & \overline{\mathbf{\xi }_{r}}\end{array}\right ) \left (\begin{array}{cc}\frac{1}{\sqrt{2}} \mathbb{I}_{r} & \frac{i}{\sqrt{2}} \mathbb{I}_{r} \\
\frac{1}{\sqrt{2}} \mathbb{I}_{r} &  -\frac{i}{\sqrt{2}} \mathbb{I}_{r}\end{array}\right ) \left (\begin{array}{cc}\frac{1}{\sqrt{2}} \mathbb{I}_{r} & \frac{i}{\sqrt{2}} \mathbb{I}_{r} \\
\frac{1}{\sqrt{2}} \mathbb{I}_{r} &  -\frac{i}{\sqrt{2}} \mathbb{I}_{r}\end{array}\right )^{ \dag } \left (\begin{array}{cc}\mathbf{\xi }_{r} & \mathbb{O}_{r} \\
\mathbb{O}_{r} & \overline{\mathbf{\xi }_{r}}\end{array}\right )^{ -1} \left (\begin{array}{c}\mathbf{v}_{a^{ \dag }} \\
\mathbf{v}_{a}\end{array}\right ) \\
 &  =\left (\mathbf{x}_{1} \left (\mathbf{\xi }\right ) ,\mathbf{x}_{2} \left (\mathbf{\xi }\right )\right )\left (\begin{array}{cc}\frac{1}{\sqrt{2}} \mathbb{I}_{r} & \frac{i}{\sqrt{2}} \mathbb{I}_{r} \\
\frac{1}{\sqrt{2}} \mathbb{I}_{r} &  -\frac{i}{\sqrt{2}} \mathbb{I}_{r}\end{array}\right )^{ \dag } \left (\begin{array}{cc}\mathbf{\xi }_{r} & \mathbb{O}_{r} \\
\mathbb{O}_{r} & \overline{\mathbf{\xi }_{r}}\end{array}\right )^{ -1} \left (\begin{array}{c}\mathbf{v}_{a^{ \dag }} \\
\mathbf{v}_{a}\end{array}\right ) \\
 &  =\left (\mathbf{x}_{1} \left (\mathbf{\xi }\right ) ,\mathbf{x}_{2} \left (\mathbf{\xi }\right )\right )\left (\begin{array}{c}\mathbf{v}_{v_{x_{1}}} \\
\mathbf{v}_{v_{x_{2}}}\end{array}\right )\text{,}\end{align}where
\begin{equation}\left (\begin{array}{c}\mathbf{v}_{v_{x_{1}}} \\
\mathbf{v}_{v_{x_{2}}}\end{array}\right ) =\left (\begin{array}{cc}\frac{1}{\sqrt{2}} \mathbb{I}_{r} & \frac{i}{\sqrt{2}} \mathbb{I}_{r} \\
\frac{1}{\sqrt{2}} \mathbb{I}_{r} &  -\frac{i}{\sqrt{2}} \mathbb{I}_{r}\end{array}\right )^{ \dag } \left (\begin{array}{cc}\mathbf{\xi }_{r} & \mathbb{O}_{r} \\
\mathbb{O}_{r} & \overline{\mathbf{\xi }_{r}}\end{array}\right )^{ -1} \left (\begin{array}{c}\mathbf{v}_{a^{ \dag }} \\
\mathbf{v}_{a}\end{array}\right )\text{.}
\end{equation}
\begin{align}A &  =\vert \overset{\overset{}{}}{\left (\mathbf{a}^{ \dag } ,\mathbf{a}\right )} \rangle  \langle \left (\begin{array}{c}\mathbf{a}^{ \dag } \\
\mathbf{a}\end{array}\right )\vert A\vert \overset{\overset{}{}}{\left (\mathbf{a}^{ \dag } ,\mathbf{a}\right )} \rangle  \langle \left (\begin{array}{c}\mathbf{a}^{ \dag } \\
\mathbf{a}\end{array}\right )\vert  \equiv \left (\vert \mathbf{a}^{ \dag } \rangle  ,\vert \mathbf{a} \rangle \right )\left (\begin{array}{c} \langle \mathbf{a}^{ \dag }\vert  \\
 \langle \mathbf{a}\vert \end{array}\right ) A \left (\vert \mathbf{a}^{ \dag } \rangle  ,\vert \mathbf{a} \rangle \right ) \left (\begin{array}{c} \langle \mathbf{a}^{ \dag }\vert  \\
 \langle \mathbf{a}\vert \end{array}\right ) \\
 &  =\left (\vert \mathbf{a}^{ \dag } \rangle  ,\vert \mathbf{a} \rangle \right )\left (\begin{array}{cc} \langle \left .\mathbf{a}^{ \dag }\right \vert A \mathbf{a}^{ \dag } \rangle  &  \langle \left .\mathbf{a}^{ \dag }\right \vert A \mathbf{a} \rangle  \\
 \langle \left .\mathbf{a}\right \vert A \mathbf{a}^{ \dag } \rangle  &  \langle \left .\mathbf{a}\right \vert A \mathbf{a} \rangle \end{array}\right ) \left (\begin{array}{c} \langle \mathbf{a}^{ \dag }\vert  \\
 \langle \mathbf{a}\vert \end{array}\right ) \\
 &  =\left (\vert \mathbf{a}^{ \dag } \rangle  ,\vert \mathbf{a} \rangle \right )\left (\begin{array}{cc}\mathbb{A}_{1 1} & \mathbb{A}_{1 2} \\
\mathbb{A}_{2 1} & \mathbb{A}_{2 2}\end{array}\right ) \left (\begin{array}{c} \langle \mathbf{a}^{ \dag }\vert  \\
 \langle \mathbf{a}\vert \end{array}\right ) =\vert \mathbf{a}^{ \dag } \rangle \mathbb{A}_{1 1}  \langle \mathbf{a}^{ \dag }\vert  +\vert \mathbf{a}^{ \dag } \rangle  \mathbb{A}_{1 2}  \langle \mathbf{a}\vert  +\vert \mathbf{a} \rangle  \mathbb{A}_{2 1}  \langle \mathbf{a}^{ \dag }\vert  +\vert \mathbf{a} \rangle  \mathbb{A}_{2 2} \langle \mathbf{a}\vert  \\
 &  =\sum _{1 \leq j ,k \leq r}A_{1 1 j k} a_{j}^{ \dag } a_{k}^{ \dag } +\sum _{1 \leq j ,k \leq r}A_{1 2 j k} a_{j}^{ \dag } a_{k} +\sum _{1 \leq j ,k \leq r}A_{2 1 j k} a_{j} a_{k}^{ \dag } +\sum _{1 \leq j ,k \leq r}A_{2 2 j k} a_{j} a_{k} \\
\lambda _{1} &  =\sum _{1 \leq l ,m \leq r}\lambda _{1 l m} a_{l}^{ \dag } a_{m}\end{align}

\section{One Particle Operators}
\textsc{(these do not have to be particle conserving operators so could be better described as one Bogoluybon Operators)} 

One can first consider one particle operators of 

\subsubsection{Diagonal form}
Diagonal one particle operators expressed in terms of the generators of the Clifford algebra are
\begin{align}\lambda _{1} &  =\sum _{1 \leq j \leq r}\beta _{1 j} b_{j}^{ \dag } b_{j} =\frac{1}{2} \sum _{1 \leq j ,k \leq r}\beta _{1 j} \left (y_{j} -i y_{j +r}\right ) \left (y_{j} +i y_{j +r}\right ) \\
 &  =\frac{1}{2} \sum _{1 \leq j \leq r}\beta _{1 j} \left \{y_{j} y_{j} +y_{j +r} y_{j +r} +i y_{j} y_{j +r} -i y_{j +r} y_{j}\right \} \\
 &  =\frac{1}{2} \left (\left .\ensuremath{\operatorname*{Tr}}\right \vert _{1} \lambda _{1}\right ) I_{\mathcal{F}} +\frac{i}{2} \sum _{1 \leq j \leq r}\beta _{1 j} \left \{\left (y_{j} y_{j +r} -y_{j +r} y_{j}\right )\right \} \\
 &  =\frac{1}{2} \left (\left .\ensuremath{\operatorname*{Tr}}\right \vert _{1} \lambda _{1}\right ) I_{\mathcal{F}} +i \sum _{1 \leq j \leq r}\beta _{1 j} y_{j} y_{j +r}\text{,}\end{align}$\lceil $checking the inverse
\begin{align} & \frac{1}{2} \left (\left .\ensuremath{\operatorname*{Tr}}\right \vert _{1} \lambda _{1}\right ) I_{\mathcal{F}} +i \sum _{1 \leq j \leq r}\beta _{1 j} y_{j} y_{j +r} = \\
 & \frac{1}{2} \left (\left .\ensuremath{\operatorname*{Tr}}\right \vert _{1} \lambda _{1}\right ) I_{\mathcal{F}} +i \sum _{1 \leq j \leq r}\beta _{1 j} \frac{1}{\sqrt{2}} \left (b_{j}^{ \dag } +b_{j}\right ) \frac{i}{\sqrt{2}} \left (b_{j}^{ \dag } -b_{j}\right ) = \\
 & \frac{1}{2} \left (\left .\ensuremath{\operatorname*{Tr}}\right \vert _{1} \lambda _{1}\right ) I_{\mathcal{F}} -\frac{1}{2} \sum _{1 \leq j \leq r}\beta _{1 j} \left ( -b_{j}^{ \dag } b_{j} +b_{j} b_{j}^{ \dag }\right ) = \\
 & \frac{1}{2} \left (\left .\ensuremath{\operatorname*{Tr}}\right \vert _{1} \lambda _{1}\right ) I_{\mathcal{F}} -\frac{1}{2} \sum _{1 \leq j \leq r}\beta _{1 j} \left (I_{\mathcal{F}} -2 b_{j}^{ \dag } b_{j}\right ) = \\
 & \frac{1}{2} \left (\left .\ensuremath{\operatorname*{Tr}}\right \vert _{1} \lambda _{1}\right ) I_{\mathcal{F}} -\frac{1}{2} \left (\left .\ensuremath{\operatorname*{Tr}}\right \vert _{1} \lambda _{1}\right ) I_{\mathcal{F}} +\sum _{1 \leq j \leq r}\beta _{1 j} b_{j}^{ \dag } b_{j} =\sum _{1 \leq j \leq r}\beta _{1 j} b_{j}^{ \dag } b_{j}\end{align}$\rfloor $ 

which is quite general as any \emph{normal} one particle operator can be diagonalized. If the
operator is self adjoint $\left \{\left .\beta _{1 j} \in \mathbb{R}\right \vert 1 \leq j \leq r\right \}\text{.}$ One can observe from the preceding that a one particle operator has both grade zero and one components and is a first degree polynomial.


We can find the square of the diagonal form of a normal one particle operator when $r =2$ as
\begin{align} & \left (\frac{1}{2} \left (\left .\ensuremath{\operatorname*{Tr}}\right \vert _{1} \lambda _{1}\right ) I_{\mathcal{F}} +i \left (\beta _{1 1} y_{1} y_{3} +\beta _{1 2} y_{2} y_{4}\right )\right )^{2} \nonumber  \\
 &  =\frac{1}{4} \left (\left .\ensuremath{\operatorname*{Tr}}\right \vert _{1} \lambda _{1}\right )^{2} I_{\mathcal{F}} +i \left (\left .\ensuremath{\operatorname*{Tr}}\right \vert _{1} \lambda _{1}\right ) \left (\beta _{1 1} y_{1} y_{3} +\beta _{1 2} y_{2} y_{4}\right ) -\left (\beta _{1 1} y_{1} y_{3} +\beta _{1 2} y_{2} y_{4}\right ) \left (\beta _{1 1} y_{1} y_{3} +\beta _{1 2} y_{2} y_{4}\right ) \nonumber  \\
 &  =\frac{1}{4} \left (\left .\ensuremath{\operatorname*{Tr}}\right \vert _{1} \lambda _{1}\right )^{2} I_{\mathcal{F}} +i \left (\left .\ensuremath{\operatorname*{Tr}}\right \vert _{1} \lambda _{1}\right ) \left (\beta _{1 1} y_{1} y_{3} +\beta _{1 2} y_{2} y_{4}\right ) +\frac{1}{4} \beta _{1 1}^{2} I_{\mathcal{F}} +\frac{1}{4} \beta _{1 2}^{2} I_{\mathcal{F}} \nonumber  \\
 & \text{\quad \quad \quad \quad \quad \quad \quad \quad \quad \quad \quad \quad \quad \quad \quad \quad \quad \quad } +\beta _{1 1} \beta _{1 2} y_{1} y_{3} y_{2} y_{4} +\beta _{1 1} \beta _{1 2} y_{2} y_{4} y_{1} y_{3} \nonumber  \\
 &  =\frac{1}{4} \left (\left .\ensuremath{\operatorname*{Tr}}\right \vert _{1} \lambda _{1}\right )^{2} I_{\mathcal{F}} +i \left (\left .\ensuremath{\operatorname*{Tr}}\right \vert _{1} \lambda _{1}\right ) \left (\beta _{1 1} y_{1} y_{3} +\beta _{1 2} y_{2} y_{4}\right ) +\frac{1}{4} \beta _{1 1}^{2} I_{\mathcal{F}} +\frac{1}{4} \beta _{1 2}^{2} I_{\mathcal{F}} -2 \beta _{1 1} \beta _{1 2} y_{1} y_{2} y_{3} y_{4} \nonumber  \\
 &  =\frac{1}{4} \left \{\left (\left .\ensuremath{\operatorname*{Tr}}\right \vert _{1} \lambda _{1}\right )^{2} +\beta _{1 1}^{2} +\beta _{1 2}^{2}\right \} I_{\mathcal{F}} +i \left (\left .\ensuremath{\operatorname*{Tr}}\right \vert _{1} \lambda _{1}\right ) \left (\beta _{1 1} y_{1} y_{3} +\beta _{1 2} y_{2} y_{4}\right ) -2 \beta _{1 1} \beta _{1 2} y_{1} y_{2} y_{3} y_{4} \nonumber  \\
 &  =\frac{1}{4} \left \{\left (\left .\ensuremath{\operatorname*{Tr}}\right \vert _{1} \lambda _{1}\right )^{2} +\left .\ensuremath{\operatorname*{Tr}}\right \vert _{1} \lambda _{1}\right \} I_{\mathcal{F}} +i \left (\left .\ensuremath{\operatorname*{Tr}}\right \vert _{1} \lambda _{1}\right ) \left (\beta _{1 1} y_{1} y_{3} +\beta _{1 2} y_{2} y_{4}\right ) -2 \beta _{1 1} \beta _{1 2} y_{1} y_{2} y_{3} y_{4}\end{align}Thus the idempotency condition gives
\begin{align} & \frac{1}{2} \left (\left .\ensuremath{\operatorname*{Tr}}\right \vert _{1} \lambda _{1} \left (\frac{1}{2} \left .\ensuremath{\operatorname*{Tr}}\right \vert _{1} \lambda _{1} -1\right ) +\frac{1}{2} \left (\beta _{1 1}^{2} +\beta _{1 2}^{2}\right )\right ) I_{\mathcal{F}} +i \left (\left .\ensuremath{\operatorname*{Tr}}\right \vert _{1} \lambda _{1} -1\right ) \left (\beta _{1 1} y_{1} y_{3} +\beta _{1 2} y_{2} y_{4}\right ) \nonumber  \\
 &  -2 \beta _{1 1} \beta _{1 2} y_{1} y_{2} y_{3} y_{4} =0 \Longrightarrow  \nonumber  \\
 & \frac{1}{2} \left (\left .\ensuremath{\operatorname*{Tr}}\right \vert _{1} \lambda _{1} \left (\frac{1}{2} \left .\ensuremath{\operatorname*{Tr}}\right \vert _{1} \lambda _{1} -1\right ) +\frac{1}{2} \left .\ensuremath{\operatorname*{Tr}}\right \vert _{1} \lambda _{1}\right ) I_{\mathcal{F}} +i \left (\left .\ensuremath{\operatorname*{Tr}}\right \vert _{1} \lambda _{1} -1\right ) \left (\beta _{1 1} y_{1} y_{3} +\beta _{1 2} y_{2} y_{4}\right ) \nonumber  \\
 &  -2 \beta _{1 1} \beta _{1 2} y_{1} y_{2} y_{3} y_{4} =0.\end{align}One can express this more transparently as
\begin{align}\left (\left .\ensuremath{\operatorname*{Tr}}\right \vert _{1} \lambda _{1}\right ) I_{\mathcal{F}} &  =2 \left (\left .\ensuremath{\operatorname*{Tr}}\right \vert _{1} \lambda _{1}\right ) \left (\left (\left .\ensuremath{\operatorname*{Tr}}\right \vert _{1} \lambda _{1}\right ) -\frac{1}{2}\right ) I_{\mathcal{F}} \\
i \left (\left .\ensuremath{\operatorname*{Tr}}\right \vert _{1} \lambda _{1}\right ) \left (\beta _{1 1} y_{1} y_{3} +\beta _{1 2} y_{2} y_{4}\right ) &  =i \left (\beta _{1 1} y_{1} y_{3} +\beta _{1 2} y_{2} y_{4}\right ) \\
2 \beta _{1 1} \beta _{1 2} y_{1} y_{2} y_{3} y_{4} &  =0 \\
\left .\ensuremath{\operatorname*{Tr}}\right \vert _{1} \lambda _{1} &  =\beta _{1 1} +\beta _{1 2}\end{align}Noting
\begin{align}y_{1} y_{2} y_{3} y_{4} &  = -\frac{1}{4} \left (a_{1}^{ \dag } +a_{1}\right ) \left (a_{2}^{ \dag } +a_{2}\right ) \left (a_{3}^{ \dag } -a_{3}\right ) \left (a_{4}^{ \dag } -a_{4}\right ) \nonumber  \\
 &  = -\frac{1}{4} \left (a_{1}^{ \dag } a_{2}^{ \dag } +a_{1}^{ \dag } a_{2} +a_{1} a_{2}^{ \dag } +a_{1} a_{2}\right ) \left (a_{3}^{ \dag } a_{4}^{ \dag } -a_{3}^{ \dag } a_{4} -a_{3} a_{4}^{ \dag } +a_{3} a_{4}\right ) \neq 0\end{align}one can deduce that either or both $\beta _{1 1} ,\beta _{1 2} =0$ thus
\begin{align}\left (\left .\ensuremath{\operatorname*{Tr}}\right \vert _{1} \lambda _{1}\right ) \left (\beta _{1 1} y_{1} y_{3}\right ) &  =\beta _{1 1} y_{1} y_{3} \nonumber  \\
 &  \Rightarrow \left .\ensuremath{\operatorname*{Tr}}\right \vert _{1} \lambda _{1} =\beta _{1 1} =1.\end{align}For any $r$ one obtains
\begin{align} & \left (\frac{1}{2} \left (\left .\ensuremath{\operatorname*{Tr}}\right \vert _{1} \lambda _{1}\right ) I_{\mathcal{F}} +i \sum _{1 \leq j \leq r}\beta _{1 j} y_{j} y_{j +r}\right )^{2} \nonumber  \\
 &  =\frac{1}{4} \left (\left .\ensuremath{\operatorname*{Tr}}\right \vert _{1} \lambda _{1}\right )^{2} I_{\mathcal{F}} +i \left (\left .\ensuremath{\operatorname*{Tr}}\right \vert _{1} \lambda _{1}\right ) \sum _{1 \leq j \leq r}\beta _{1 j} y_{j} y_{j +r} -\sum _{1 \leq j \leq r}\beta _{1 j} y_{j} y_{j +r} \sum _{1 \leq j \leq r}\beta _{1 j} y_{j} y_{j +r} \nonumber  \\
 &  =\frac{1}{4} \left (\left .\ensuremath{\operatorname*{Tr}}\right \vert _{1} \lambda _{1}\right )^{2} I_{\mathcal{F}} +i \left (\left .\ensuremath{\operatorname*{Tr}}\right \vert _{1} \lambda _{1}\right ) \sum _{1 \leq j \leq r}\beta _{1 j} y_{j} y_{j +r} +\frac{1}{4} \sum _{1 \leq j \leq r}\beta _{1 j}^{2} I_{\mathcal{F}} +2 \sum _{1 \leq j_{1} <j_{2} \leq r}\beta _{1 j_{1}} \beta _{1 j_{2}} y_{j_{1}} y_{j_{1} +r} y_{j_{2}} y_{j_{2} +r} \nonumber  \\
 &  =\frac{1}{4} \left (\left .\ensuremath{\operatorname*{Tr}}\right \vert _{1} \lambda _{1}\right )^{2} I_{\mathcal{F}} +i \left (\left .\ensuremath{\operatorname*{Tr}}\right \vert _{1} \lambda _{1}\right ) \sum _{1 \leq j \leq r}\beta _{1 j} y_{j} y_{j +r} +\frac{1}{4} \sum _{1 \leq j \leq r}\beta _{1 j}^{2} I_{\mathcal{F}} -2 \sum _{1 \leq j_{1} <j_{2} \leq r}\beta _{1 j_{1}} \beta _{1 j_{2}} y_{j_{1}} y_{j_{2}} y_{j_{1} +r} y_{j_{2} +r} \nonumber  \\
 &  =\frac{1}{4} \left \{\left (\left .\ensuremath{\operatorname*{Tr}}\right \vert _{1} \lambda _{1}\right )^{2} +\sum _{1 \leq j \leq r}\beta _{1 j}^{2} I_{\mathcal{F}}\right \} I_{\mathcal{F}} +i \left (\left .\ensuremath{\operatorname*{Tr}}\right \vert _{1} \lambda _{1}\right ) \sum _{1 \leq j \leq r}\beta _{1 j} y_{j} y_{j +r} -2 \sum _{1 \leq j_{1} <j_{2} \leq r}\beta _{1 j_{1}} \beta _{1 j_{2}} y_{j_{1}} y_{j_{2}} y_{j_{1} +r} y_{j_{2} +r} \nonumber  \\
 &  =\frac{1}{4} \left \{\left (\left .\ensuremath{\operatorname*{Tr}}\right \vert _{1} \lambda _{1}\right )^{2} +\left (\left .\ensuremath{\operatorname*{Tr}}\right \vert _{1} \lambda _{1}\right ) I_{\mathcal{F}}\right \} I_{\mathcal{F}} +i \left (\left .\ensuremath{\operatorname*{Tr}}\right \vert _{1} \lambda _{1}\right ) \sum _{1 \leq j \leq r}\beta _{1 j} y_{j} y_{j +r} -2 \sum _{1 \leq j_{1} <j_{2} \leq r}\beta _{1 j_{1}} \beta _{1 j_{2}} y_{j_{1}} y_{j_{2}} y_{j_{1} +r} y_{j_{2} +r}\end{align}The idempotency condition
\begin{equation}\left (\frac{1}{2} \left (\left .\ensuremath{\operatorname*{Tr}}\right \vert _{1} \lambda _{1}\right ) I_{\mathcal{F}} +i \sum _{1 \leq j \leq r}\beta _{1 j} y_{j} y_{j +r}\right )^{2} -\left (\frac{1}{2} \left (\left .\ensuremath{\operatorname*{Tr}}\right \vert _{1} \lambda _{1}\right ) I_{\mathcal{F}} +i \sum _{1 \leq j \leq r}\beta _{1 j} y_{j} y_{j +r}\right ) =\,0
\end{equation}gives
\begin{align} & \frac{1}{4} \left \{\left (\left .\ensuremath{\operatorname*{Tr}}\right \vert _{1} \lambda _{1}\right ) \left (\left (\left .\ensuremath{\operatorname*{Tr}}\right \vert _{1} \lambda _{1}\right ) I_{\mathcal{F}} -2\right ) +\left (\left .\ensuremath{\operatorname*{Tr}}\right \vert _{1} \lambda _{1}\right )\right \} I_{\mathcal{F}} +i \left (\left .\ensuremath{\operatorname*{Tr}}\right \vert _{1} \lambda _{1} -1\right ) \sum _{1 \leq j \leq r}\beta _{1 j} y_{j} y_{j +r} \nonumber  \\
 & . -.2 \sum _{1 \leq j_{1} <j_{2} \leq r}\beta _{1 j_{1}} \beta _{1 j_{2}} y_{j_{1}} y_{j_{2}} y_{j_{1} +r} y_{j_{2} +r} =0.\end{align}One can apply the same reasoning as above to deduce that $\left .\ensuremath{\operatorname*{Tr}}\right \vert _{1} \lambda _{1} =\beta _{1 1} =1$ and all other $\beta _{1 j}$ 's are zero 

\subsubsection{General form}
Any one particle operator $\lambda _{1}$ can be expressed as
\begin{align}\lambda _{1} &  =\sum _{1 \leq j ,k \leq r}\lambda _{1 j k} a_{j}^{ \dag } a_{k} =\frac{1}{2} \sum _{1 \leq j ,k \leq r}\lambda _{1 j k} \left (x_{j} -i x_{j +r}\right ) \left (x_{k} +i x_{k +r}\right ) \nonumber  \\
 &  =\frac{1}{2} \sum _{1 \leq j ,k \leq r}\lambda _{1 j k} \left \{x_{j} x_{k} +x_{j +r} x_{k +r} +i x_{j} x_{k +r} -i x_{j +r} x_{k}\right \}\text{,}\end{align}which can be further developed as
\begin{align}\lambda _{1} &  =\frac{1}{2} \left (\left .\ensuremath{\operatorname*{Tr}}\right \vert _{1} \lambda _{1}\right ) I_{\mathcal{F}} +\frac{1}{2} \sum _{1 \leq j <k \leq r}\left (\lambda _{1 j k} -\lambda _{1 k j}\right ) \left \{x_{j} x_{k} +x_{j +r} x_{k +r} +i \left (x_{j} x_{k +r} -x_{j +r} x_{k}\right )\right \} \nonumber  \\
 &  =\frac{1}{2} \left (\left .\ensuremath{\operatorname*{Tr}}\right \vert _{1} \lambda _{1}\right ) I_{\mathcal{F}} +\sum _{1 \leq j <k \leq 2 r}\eta _{j k} x_{j} x_{k}\text{,}\end{align}where the coefficients $\eta _{j k}$ are antisymmetric and given by
\begin{align}\eta _{j k} &  =\frac{1}{2} \left (\lambda _{1 j k} -\lambda _{1 k j}\right ) ;\text{\quad }1 \leq j <k \leq r ;r +1 \leq j <k \leq 2 r \\
\eta _{j k} &  = -\frac{i}{2} \left (\lambda _{1 j k} -\lambda _{1 k j}\right ) ;\text{\quad }1 \leq j \leq r ;r +1 \leq k \leq 2 r \\
\eta _{j j} &  =0 ;\text{\quad }1 \leq j \leq 2 r\text{.}\end{align}

One can compute the square of $\lambda _{1}$ when $r =2$ i.e.
\begin{align*}\left (\lambda _{1}\right )^{2} &  =\left (\frac{1}{2} \left (\left .\ensuremath{\operatorname*{Tr}}\right \vert _{1} \lambda _{1}\right ) I_{\mathcal{F}} +\eta _{1 2} x_{1} x_{2} +\eta _{1 3} x_{1} x_{3} +\eta _{1 4} x_{1} x_{4} +\eta _{2 3} x_{2} x_{3} +\eta _{2 4} x_{2} x_{4} +\eta _{3 4} x_{3} x_{4}\right )^{2} \\
 &  =\frac{1}{4} \left (\left .\ensuremath{\operatorname*{Tr}}\right \vert _{1} \lambda _{1}\right )^{2} I_{\mathcal{F}} +\left (\left .\ensuremath{\operatorname*{Tr}}\right \vert _{1} \lambda _{1}\right ) \left (\eta _{1 2} x_{1} x_{2} +\eta _{1 3} x_{1} x_{3} +\eta _{1 4} x_{1} x_{4} +\eta _{2 3} x_{2} x_{3} +\eta _{2 4} x_{2} x_{4} +\eta _{3 4} x_{3} x_{4}\right ) \\
 &  -\frac{1}{4} \left (\eta _{1 2}^{2} +\eta _{1 3}^{2} +\eta _{1 4}^{2} +\eta _{2 3}^{2} +\eta _{2 4}^{2} +\eta _{3 4}^{2}\right ) I_{\mathcal{F}} \\
 &  +\eta _{1 2} \eta _{1 3} x_{1} x_{2} x_{1} x_{3} +\eta _{1 2} \eta _{1 4} x_{1} x_{2} x_{1} x_{4} +\eta _{1 2} \eta _{2 3} x_{1} x_{2} x_{2} x_{3} +\eta _{1 2} \eta _{2 4} x_{1} x_{2} x_{2} x_{4} \\
 &  +\eta _{1 3} \eta _{1 2} x_{1} x_{3} x_{1} x_{2} +\eta _{1 3} \eta _{1 4} x_{1} x_{3} x_{1} x_{4} +\eta _{1 3} x_{1} x_{3} \eta _{2 3} x_{2} x_{3} +\eta _{1 3} \eta _{3 4} x_{1} x_{3} x_{3} x_{4} \\
 &  +\eta _{1 4} \eta _{1 2} x_{1} x_{4} x_{1} x_{2} +\eta _{1 4} \eta _{1 3} x_{1} x_{4} x_{1} x_{3} +\eta _{1 4} x_{1} x_{4} \eta _{2 4} x_{2} x_{4} +\eta _{1 4} \eta _{3 4} x_{1} x_{4} x_{3} x_{4} \\
 &  +\eta _{2 3} \eta _{1 2} x_{2} x_{3} x_{1} x_{2} +\eta _{2 3} \eta _{1 3} x_{2} x_{3} x_{1} x_{3} +\eta _{2 3} \eta _{2 4} x_{2} x_{3} x_{2} x_{4} +\eta _{2 3} \eta _{3 4} x_{2} x_{3} x_{3} x_{4} \\
 &  +\eta _{2 4} \eta _{1 2} x_{2} x_{4} x_{1} x_{2} +\eta _{2 4} \eta _{1 4} x_{2} x_{4} x_{1} x_{4} +\eta _{2 4} \eta _{2 3} x_{2} x_{4} x_{2} x_{3} +\eta _{2 4} \eta _{3 4} x_{2} x_{4} x_{3} x_{4} \\
 &  +\eta _{3 4} \eta _{1 3} x_{3} x_{4} x_{1} x_{3} +\eta _{3 4} \eta _{1 4} x_{3} x_{4} x_{1} x_{4} +\eta _{3 4} \eta _{2 3} x_{3} x_{4} x_{2} x_{3} +\eta _{3 4} \eta _{2 4} x_{3} x_{4} x_{2} x_{4} \\
 &  +\eta _{1 2} \eta _{3 4} x_{1} x_{2} x_{3} x_{4} +\eta _{1 3} \eta _{2 4} x_{1} x_{3} x_{2} x_{4} +\eta _{1 4} \eta _{2 3} x_{1} x_{4} x_{2} x_{3} +\eta _{2 3} \eta _{1 4} x_{2} x_{3} x_{1} x_{4} \\
 &  +\eta _{2 4} x \eta _{1 3 2} x_{4} x_{1} x_{3} +\eta _{3 4} \eta _{1 2} x_{3} x_{4} x_{1} x_{2}\end{align*}
\begin{align*} &  =\left (\left (\left .\ensuremath{\operatorname*{Tr}}\right \vert _{1} \lambda _{1}\right )^{2} -\frac{1}{4} \left (\eta _{1 2}^{2} +\eta _{1 3}^{2} +\eta _{1 4}^{2} +\eta _{2 3}^{2} +\eta _{2 4}^{2} +\eta _{3 4}^{2}\right )\right ) I_{\mathcal{F}} \\
 &  +\left (\left (\ensuremath{\operatorname*{Tr}} \lambda _{1}\right ) \eta _{2 3} +\eta _{3 4} \eta _{2 4} -\eta _{1 2} \eta _{1 3} +\eta _{1 3} \eta _{1 2} -\eta _{2 4} \eta _{3 4}\right ) x_{2} x_{3} \\
 &  +\left (\left (\ensuremath{\operatorname*{Tr}} \lambda _{1}\right ) \eta _{2 4} +\eta _{2 3} \eta _{3 4} -\eta _{1 2} \eta _{1 4} +\eta _{1 4} \eta _{1 2} -\eta _{3 4} \eta _{2 3}\right ) x_{2} x_{4} \\
 &  +\left (\left (\ensuremath{\operatorname*{Tr}} \lambda _{1}\right ) \eta _{1 3} +\eta _{1 2} \eta _{2 3} -\eta _{1 4} \eta _{3 4} -\eta _{2 3} \eta _{1 2} +\eta _{3 4} \eta _{1 4 3}\right ) x_{1} x_{3} \\
 &  +\left (\left (\ensuremath{\operatorname*{Tr}} \lambda _{1}\right ) \eta _{1 4} +\eta _{1 2} \eta _{2 4} -\eta _{2 4} \eta _{1 2} +\eta _{1 3} \eta _{3 4} -\eta _{3 4} \eta _{1 3}\right ) x_{1} x_{4} \\
 &  +\left (\left (\ensuremath{\operatorname*{Tr}} \lambda _{1}\right ) \eta _{3 4} +\eta _{1 4} \eta _{1 3} -\eta _{1 3} \eta _{1 4} -\eta _{2 3} \eta _{2 4} +\eta _{2 4} \eta _{2 3}\right ) x_{3} x_{4} \\
 &  +\left (\left (\ensuremath{\operatorname*{Tr}} \lambda _{1}\right ) \eta _{1 2} +\eta _{2 3} \eta _{1 3} -\eta _{1 3} \eta _{2 3} -\eta _{1 4} \eta _{2 4} +\eta _{2 4} \eta _{1 4}\right ) x_{1} x_{2} \\
 &  +2 \left (\eta _{1 2} \eta _{3 4} -\eta _{1 3} \eta _{2 4} +\eta _{1 4} \eta _{2 3}\right ) x_{1} x_{2} x_{3} x_{4}\end{align*}
\begin{align*} &  =\left (\left (\ensuremath{\operatorname*{Tr}}\lambda _{1}\right )^{2} -\frac{1}{4} \left (\eta _{1 2}^{2} +\eta _{1 3}^{2} +\eta _{1 4}^{2} +\eta _{2 3}^{2} +\eta _{2 4}^{2} +\eta _{3 4}^{2}\right )\right ) I_{\mathcal{F}} \\
 &  +\left (\ensuremath{\operatorname*{Tr}}\lambda _{1}\right ) \left (\eta _{2 3} x_{2} x_{3} +\eta _{2 4} x_{2} x_{4} +\eta _{1 3} x_{1} x_{3} +\eta _{1 4} x_{1} x_{4} +\eta _{3 4} x_{3} x_{4} +\eta _{1 2} x_{1} x_{2}\right ) \\
 &  +2 \left (\eta _{1 2} \eta _{3 4} -\eta _{1 3} \eta _{2 4} +\eta _{1 4} \eta _{2 3}\right ) x_{1} x_{2} x_{3} x_{4}\end{align*}
\begin{align} &  =\left (\ensuremath{\operatorname*{Tr}}\lambda _{1} \left (\ensuremath{\operatorname*{Tr}}\lambda _{1} -1\right ) -\left (\eta _{1 2}^{2} +\eta _{1 3}^{2} +\eta _{1 4}^{2} +\eta _{2 3}^{2} +\eta _{2 4}^{2} +\eta _{3 4}^{2}\right )\right ) I_{\mathcal{F}} \nonumber  \\
 &  +\left (\left .\ensuremath{\operatorname*{Tr}}\right \vert _{1} \lambda _{1} -1\right ) \left (\eta _{2 3} x_{2} x_{3} +\eta _{2 4} x_{2} x_{4} +\eta _{1 3} x_{1} x_{3} +\eta _{1 4} x_{1} x_{4} +\eta _{3 4} x_{3} x_{4} +\eta _{1 2} x_{1} x_{2}\right ) \nonumber  \\
 & . +.2 \left (\eta _{1 2} \eta _{3 4} -\eta _{1 3} \eta _{2 4} +\eta _{1 4} \eta _{2 3}\right ) x_{1} x_{2} x_{3} x_{4} =0.\end{align}recall
\begin{align}\eta _{j k} &  =\frac{1}{2} \left (\lambda _{1 j k} -\lambda _{1 k j}\right ) ;\text{\quad }1 \leq j <k \leq r ;r +1 \leq j <k \leq 2 r \\
\eta _{j k} &  = -\frac{i}{2} \left (\lambda _{1 j k} -\lambda _{1 k j}\right ) ;\text{\quad }1 \leq j \leq r ;r +1 \leq k \leq 2 r\text{.}\end{align}This gives in general for any $r$
\begin{align}\left (\lambda _{1}\right )^{2} = & \left (\left (\ensuremath{\operatorname*{Tr}}\lambda _{1}\right ) \left (\ensuremath{\operatorname*{Tr}}\lambda _{1} -1\right ) -\sum _{1 \leq j <k \leq r}\eta _{j k}^{2}\right ) I_{\mathcal{F}} +\left (\left .\ensuremath{\operatorname*{Tr}}\right \vert _{1} \lambda _{1} -1\right ) \sum _{1 \leq j <k \leq r}\eta _{j k} x_{j} x_{k} \nonumber  \\
 & . +.2 \sum _{1 \leq j_{1} <j_{2} <j_{3} <j_{4} \leq r}\ensuremath{\operatorname*{Pf}} \left (\eta _{j_{1} j_{2} j_{3} j_{4}}\right ) x_{j_{1}} x_{j_{2}} x_{j_{3}} x_{j_{4}} =0.\end{align}

\subsubsection{One Particle and One Hole Idempotents}
The creators and annihilators
\begin{align}a_{j}^{ \dag } &  =\frac{1}{\sqrt{2}} \left (x_{j} -i x_{j +r}\right ) \\
a_{j} &  =\frac{1}{\sqrt{2}} \left (x_{j} +i x_{j +r}\right )\end{align}produce the one particle operators
\begin{align}a_{j}^{ \dag } a_{k} &  =\frac{1}{2} \left (x_{j} -i x_{j +r}\right ) \left (x_{k} +i x_{k +r}\right ) =\frac{1}{2} \left (x_{j} x_{k} +i x_{j} x_{k +r} -i x_{j +r} x_{k} +x_{j +r} x_{k +r}\right ) \\
 &  =\frac{1}{2} \left \{\left (x_{j} x_{k} +x_{j +r} x_{k +r}\right ) +i \left (x_{j} x_{k +r} -x_{j +r} x_{k}\right )\right \}\text{.}\end{align}One particle idempotents are given by
\begin{align}n_{j} &  =a_{j}^{ \dag } a_{j} =\frac{1}{2} \left (x_{j} -i x_{j +r}\right ) \left (x_{j} +i x_{j +r}\right ) =\frac{1}{2} \left (x_{j}^{2} +i x_{j} x_{j +r} -i x_{j +r} x_{j} +x_{j +r}^{2}\right ) \\
 &  =\frac{1}{2} I_{\mathcal{F}} +i x_{j} x_{j +r} \\
\left (n_{j}\right )^{2} &  =\left (\frac{1}{2} I_{\mathcal{F}} +i x_{j} x_{j +r}\right )^{2} =\left (\frac{1}{2} I_{\mathcal{F}} +i x_{j} x_{j +r}\right ) \left (\frac{1}{2} I_{\mathcal{F}} +i x_{j} x_{j +r}\right ) \\
 &  =\frac{1}{4} I_{\mathcal{F}} +\frac{1}{2} i x_{j} x_{j +r} +\frac{1}{2} i x_{j} x_{j +r} -x_{j} x_{j +r} x_{j} x_{j +r} \\
 &  =\frac{1}{4} I_{\mathcal{F}} +i x_{j} x_{j +r} +\frac{1}{4} I_{\mathcal{F}} =\frac{1}{2} I_{\mathcal{F}} +i x_{j} x_{j +r} =n_{j}\text{.}\end{align}and one hole idempotents are given by
\begin{align}h_{j} &  =a_{j} a_{j}^{ \dag } =\frac{1}{2} \left (x_{j} +i x_{j +r}\right ) \left (x_{j} -i x_{j +r}\right ) =\frac{1}{2} \left (x_{j}^{2} -i x_{j} x_{j +r} +i x_{j +r} x_{j} +x_{j +r}^{2}\right ) \\
 &  =\frac{1}{2} I_{\mathcal{F}} -i x_{j} x_{j +r} =\left (\frac{1}{2} I_{\mathcal{F}} +i x_{j} x_{j +r}\right )^{t} =n_{j}^{t}\text{.}\end{align}\emph{Note that }$h_{j}$ is the \emph{transpose} of $n_{j}\text{,}$ i.e.
\begin{equation}h_{j} =n_{j}^{t}\text{.}
\end{equation}

The projectors $\left \{n_{j} ,h_{j}\right \}$ are \emph{infinite dimensional} when $r =\infty $. 

One can form other projectors of similiar form e.g. consider
\begin{equation}p =\frac{1}{2} I_{\mathcal{F}} +i x y\text{,}
\end{equation}where
\begin{equation}x^{2} =y^{2} =\frac{1}{2} I_{\mathcal{F}}\text{,}
\end{equation}then
\begin{equation}p^{2} =\left (\frac{1}{2} I_{\mathcal{F}} +i x y\right ) \left (\frac{1}{2} I_{\mathcal{F}} +i x y\right ) =\frac{1}{4} I_{\mathcal{F}} +\frac{1}{2} i x y +\frac{1}{2} i x y -x y x y =\frac{1}{2} I_{\mathcal{F}} +i x y =p\text{.}
\end{equation}Noting that
\begin{align}x_{j} &  =\frac{1}{\sqrt{2}} \left (a_{j}^{ \dag } +a_{j}\right ) =x_{j}^{ \dag } \\
x_{j +r} &  =\frac{i}{\sqrt{2}} \left (a_{j}^{ \dag } -a_{j}\right ) =x_{j +r}^{ \dag }\end{align}one can see that
\begin{equation}p_{j k} =\left (\frac{1}{2} I_{\mathcal{F}} +i x_{j} x_{k}\right ) =\frac{1}{2} \left (I_{\mathcal{F}} +i \left (a_{j}^{ \dag } a_{k}^{ \dag } +a_{j} a_{k}^{ \dag } +a_{j}^{ \dag } a_{k} +a_{j} a_{k}\right )\right )
\end{equation}are self adjoint idempotents. 

One can also note that one can form \emph{para vector}
idempotents
\begin{equation}p_{x} =\frac{1}{\sqrt{2}} \left (I_{\mathcal{F}} +\sqrt{2} x\right )\text{,}
\end{equation}where
\begin{equation}x^{2} =\frac{1}{2} I_{\mathcal{F}}\text{.}
\end{equation}

\subsubsection{Clifford product versus (vector) exterior product}
\begin{align}n_{1}^{2} &  =\left (\frac{1}{2} I_{\mathcal{F}} +i x_{1} x_{2}\right ) \left (\frac{1}{2} I_{\mathcal{F}} +i x_{1} x_{2}\right ) =\frac{1}{4} I_{\mathcal{F}} +\frac{i}{2} \left (x_{1} x_{2} +x_{1} x_{2}\right ) -x_{1} x_{2} x_{1} x_{2} =\frac{1}{2} I_{\mathcal{F}} +i x_{1} x_{2} =n_{1} \\
n_{1} \wedge n_{1} = & \left (\frac{1}{2} I_{\mathcal{F}} +i x_{1} x_{2}\right ) \wedge \left (\frac{1}{2} I_{\mathcal{F}} +i x_{1} x_{2}\right ) =\frac{1}{4} I_{\mathcal{F}} \wedge I_{\mathcal{F}} +\frac{1}{2} i I_{\mathcal{F}} \wedge x_{1} \wedge x_{2} +\frac{1}{2} i x_{1} \wedge x_{2} \wedge I_{\mathcal{F}} \\
 &  =\frac{1}{4} I_{\mathcal{F}} +i x_{1} \wedge x_{2} =n_{1} -\frac{1}{4} I_{\mathcal{F}}\text{.}\end{align}$\lceil $
\begin{align}x_{j} &  =\frac{1}{\sqrt{2}} \left (a_{j}^{ \dag } +a_{j}\right ) \\
x_{j +r} &  =\frac{i}{\sqrt{2}} \left (a_{j}^{ \dag } -a_{j}\right )\end{align}
\begin{align}\frac{1}{\sqrt{2}} \left (x_{j} +i x_{j +r}\right ) &  =a_{j} \\
\frac{1}{\sqrt{2}} \left (x_{j} -i x_{j +r}\right ) &  =a_{j}^{ \dag }\end{align}$\rfloor $ 

\section{Two Particle Operators}
\textsc{(these do not have to be particle conserving operators so could be better described as two Bogoliubon Operators). These are not homogeneous
polynomials when expressed in terms of }$\left \{x_{j}\right \}$ 

One can first consider two particle operators of 

\subsubsection{Decomposable Diagonal form}
These type of particle operators can be expanded as
\begin{align}\lambda _{2} &  =\frac{1}{2} \sum _{1 \leq j_{1} ,j_{2} \leq r}\lambda _{2 j_{1} j_{2} j_{1} j_{2}} a_{j_{1}}^{ \dag } a_{j_{2}}^{ \dag } a_{j_{2}} a_{j_{1}} \nonumber  \\
 &  =\sum _{1 \leq j_{1} <j_{2} \leq r}\lambda _{2 j_{1} j_{2} j_{1} j_{2}} n_{j_{1}} n_{j_{2}} =\sum _{1 \leq j_{1} <j_{2} \leq r}\lambda _{2 j_{1} j_{2} j_{1} j_{2}} \left (\frac{1}{2} I_{\mathcal{F}} +i x_{j_{1}} x_{j_{1} +r}\right ) \left (\frac{1}{2} I_{\mathcal{F}} +i x_{j_{2}} x_{j_{2} +r}\right ) \\
 &  =\frac{1}{8} \sum _{1 \leq j_{1} ,j_{2} \leq r}\lambda _{2 j_{1} j_{2} j_{1} j_{2}} \left (x_{j_{1}} -i x_{j_{1} +r}\right ) \left (x_{j_{2}} -i x_{j_{2} +r}\right ) \left (x_{j_{2}} +i x_{j_{2} +r}\right ) \left (x_{j_{1}} +i x_{j_{1} +r}\right ) \nonumber  \\
 &  =\frac{1}{8} \sum _{1 \leq j_{1} ,j_{2} \leq r}\lambda _{2 j_{1} j_{2} j_{1} j_{2}} \left (x_{j_{1}} x_{j_{2}} -i x_{j_{1} +r} x_{j_{2}} -i x_{j_{1}} x_{j_{2} +r} -x_{j_{1} +r} x_{j_{2} +r}\right ) \times  \nonumber  \\
 & \left (x_{j_{2}} x_{j_{1}} +i x_{j_{2} +r} x_{j_{1}} +i x_{j_{2}} x_{j_{1} +r} -x_{j_{2} +r} x_{j_{1} +r}\right ) \nonumber  \\
 &  =\frac{1}{8} \sum _{1 \leq j_{1} ,j_{2} \leq r}\lambda _{2 j_{1} j_{2} j_{1} j_{2}}\left \{\text{}\right . \nonumber  \\
 &  +x_{j_{1}} x_{j_{2}} x_{j_{2}} x_{j_{1}} -i x_{j_{1} +r} x_{j_{2}} x_{j_{2}} x_{j_{1}} -i x_{j_{1}} x_{j_{2} +r} x_{j_{2}} x_{j_{1}} -x_{j_{1} +r} x_{j_{2} +r} x_{j_{2}} x_{j_{1}} \nonumber  \\
 &  +i x_{j_{1}} x_{j_{2}} x_{j_{2} +r} x_{j_{1}} +x_{j_{1} +r} x_{j_{2}} x_{j_{2} +r} x_{j_{1}} +x_{j_{1}} x_{j_{2} +r} x_{j_{2} +r} x_{j_{1}} -i x_{j_{1} +r} x_{j_{2} +r} x_{j_{2} +r} x_{j_{1}} \nonumber  \\
 &  +i x_{j_{1}} x_{j_{2}} x_{j_{2}} x_{j_{1} +r} +x_{j_{1} +r} x_{j_{2}} x_{j_{2}} x_{j_{1} +r} +x_{j_{1}} x_{j_{2} +r} x_{j_{2}} x_{j_{1} +r} -i x_{j_{1} +r} x_{j_{2} +r} x_{j_{2}} x_{j_{1} +r} \nonumber  \\
 &  -\left .x_{j_{1}} x_{j_{2}} x_{j_{2} +r} x_{j_{1} +r} +i x_{j_{1} +r} x_{j_{2}} x_{j_{2} +r} x_{j_{1} +r} +i x_{j_{1}} x_{j_{2} +r} x_{j_{2} +r} x_{j_{1} +r} +x_{j_{1} +r} x_{j_{2} +r} x_{j_{2} +r} x_{j_{1} +r}\right \}\end{align}The zeroth degree terms are
\begin{align} & \frac{1}{8} \sum _{1 \leq j_{1} ,j_{2} \leq r}\lambda _{2 j_{1} j_{2} j_{1} j_{2}} \left \{x_{j_{1}} x_{j_{2}} x_{j_{2}} x_{j_{1}} +x_{j_{1}} x_{j_{2} +r} x_{j_{2} +r} x_{j_{1}} +x_{j_{1} +r} x_{j_{2}} x_{j_{2}} x_{j_{1} +r} +x_{j_{1} +r} x_{j_{2} +r} x_{j_{2} +r} x_{j_{1} +r}\right \} \nonumber  \\
 &  =\frac{1}{4} \sum _{1 \leq j_{1} <j_{2} \leq r}\lambda _{2 j_{1} j_{2} j_{1} j_{2}} I_{\mathcal{F}} =\frac{1}{4} \ensuremath{\operatorname*{Tr}}\left \{\lambda _{2}\right \} I_{\mathcal{F}}\end{align}The second degree terms are
\begin{align*} &  =\frac{1}{8} \sum _{1 \leq j_{1} ,j_{2} \leq r}\lambda _{2 j_{1} j_{2} j_{1} j_{2}}\left \{ -i x_{j_{1} +r} x_{j_{2}} x_{j_{2}} x_{j_{1}} -i x_{j_{1}} x_{j_{2} +r} x_{j_{2}} x_{j_{1}}\right . \\
 &  +i x_{j_{1}} x_{j_{2}} x_{j_{2} +r} x_{j_{1}} -i x_{j_{1} +r} x_{j_{2} +r} x_{j_{2} +r} x_{j_{1}} +i x_{j_{1}} x_{j_{2}} x_{j_{2}} x_{j_{1} +r} \\
 &  -\left .i x_{j_{1} +r} x_{j_{2}} x_{j_{2} +r} x_{j_{1} +r} -i x_{j_{1} +r} x_{j_{2} +r} x_{j_{2}} x_{j_{1} +r} +i x_{j_{1}} x_{j_{2} +r} x_{j_{2} +r} x_{j_{1} +r}\right \}\end{align*}
\begin{align} &  =\frac{1}{16} \sum _{1 \leq j_{1} ,j_{2} \leq r}\lambda _{2 j_{1} j_{2} j_{1} j_{2}}\left \{ -i x_{j_{1} +r} x_{j_{1}} -i x_{j_{2} +r} x_{j_{2}}\right . \nonumber  \\
 &  +i x_{j_{2}} x_{j_{2} +r} -i x_{j_{1} +r} x_{j_{1}} +i x_{j_{1}} x_{j_{1} +r} -i x_{j_{2} +r} x_{j_{2}} -\left .i x_{j_{2} +r} x_{j_{2}} +i x_{j_{1}} x_{j_{1} +r}\right \} \nonumber  \\
 &  =\frac{i}{2} \sum _{1 \leq j_{1} ,j_{2} \leq r}\lambda _{2 j_{1} j_{2} j_{1} j_{2}} x_{j_{1}} x_{j_{1} +r} =\frac{i}{2} \sum _{1 \leq j_{1} \leq r}L_{2}^{1} \left (\lambda _{2}\right )_{j_{1} j_{1}} x_{j_{1}} x_{j_{1} +r}\text{,}\end{align}where the contraction operator is defined by
\begin{equation}L_{2}^{1} \left (\lambda _{2}\right )_{j_{1} j_{2}} =\sum _{1 \leq j_{3} \leq r}\lambda _{2 j_{1} j_{3} j_{2} j_{3}}\text{.}
\end{equation}The fourth degree terms are
\begin{align} & \frac{1}{8} \sum _{1 \leq j_{1} ,j_{2} \leq r}\lambda _{2 j_{1} j_{2} j_{1} j_{2}} \left \{x_{j_{1}} x_{j_{2}} x_{j_{1} +r} x_{j_{2} +r} +x_{j_{1}} x_{j_{2}} x_{j_{1} +r} x_{j_{2} +r} +x_{j_{1}} x_{j_{2}} x_{j_{1} +r} x_{j_{2} +r} +x_{j_{1}} x_{j_{2}} x_{j_{1} +r} x_{j_{2} +r}\right \} \nonumber  \\
 &  =\frac{1}{2} \sum _{1 \leq j_{1} ,j_{2} \leq r}\lambda _{2 j_{1} j_{2} j_{1} j_{2}} x_{j_{1}} x_{j_{2}} x_{j_{1} +r} x_{j_{2} +r} =\sum _{1 \leq j_{1} <j_{2} \leq r}\lambda _{2 j_{1} j_{2} j_{1} j_{2}} x_{j_{1}} x_{j_{2}} x_{j_{1} +r} x_{j_{2} +r}\text{.}\end{align}

Hence
\begin{align}\lambda _{2} &  =\frac{1}{4} \left .\ensuremath{\operatorname*{Tr}}\right \vert _{2} \lambda _{2} I_{\mathcal{F}} +\frac{i}{2} \sum _{1 \leq j_{1} \leq r}L_{2}^{1} \left (\lambda _{2}\right )_{j_{1} j_{1}} x_{j_{1}} x_{j_{1} +r} +\sum _{1 \leq j_{1} <j_{2} \leq r}\lambda _{2 j_{1} j_{2} j_{1} j_{2}} x_{j_{1}} x_{j_{2}} x_{j_{1} +r} x_{j_{2} +r} \nonumber  \\
 &  =\frac{1}{4} L_{2}^{0} \left (\lambda _{2}\right ) I_{\mathcal{F}} +\frac{i}{2} \sum _{1 \leq j_{1} \leq r}L_{2}^{1} \left (\lambda _{2}\right )_{j_{1} j_{1}} x_{j_{1}} x_{j_{1} +r} +\sum _{1 \leq j_{1} <j_{2} \leq r}L_{2}^{2} \left (\lambda _{2}\right )_{j_{1} j_{2} j_{1} j_{2}} x_{j_{1}} x_{j_{2}} x_{j_{1} +r} x_{j_{2} +r}\text{.}\end{align}

\subsubsection{General form}
General two particle operators can be expanded as
\begin{align}\lambda _{2} &  =\frac{1}{4} \sum _{1 \leq j_{1} ,j_{2} ,j_{3} ,j_{4} \leq r}\lambda _{2 j_{1} j_{2} j_{3} j_{4}} a_{j_{1}}^{ \dag } a_{j_{2}}^{ \dag } a_{j_{4}} a_{j_{3}} \nonumber  \\
 &  =\frac{1}{16} \sum _{1 \leq j_{1} ,j_{2} ,j_{3} ,j_{4} \leq r}\lambda _{2 j_{1} j_{2} j_{3} j_{4}} \left (x_{j_{1}} -i x_{j_{1} +r}\right ) \left (x_{j_{2}} -i x_{j_{2} +r}\right ) \left (x_{j_{4}} +i x_{j_{4} +r}\right ) \left (x_{j_{3}} +i x_{j_{3} +r}\right ) \nonumber  \\
 &  =\frac{1}{16} \sum _{1 \leq j_{1} ,j_{2} ,j_{3} ,j_{4} \leq r}\lambda _{2 j_{1} j_{2} j_{3} j_{4}} \left (x_{j_{1}} x_{j_{2}} -i x_{j_{1} +r} x_{j_{2}} -i x_{j_{1}} x_{j_{2} +r} -x_{j_{1} +r} x_{j_{2} +r}\right ) \nonumber  \\
 &  \times \left (x_{j_{4}} x_{j_{3}} +i x_{j_{4} +r} x_{j_{3}} +i x_{j_{4}} x_{j_{3} +r} -x_{j_{4} +r} x_{j_{3} +r}\right ) \\
 &  =\frac{1}{16} \sum _{1 \leq j_{1} ,j_{2} ,j_{3} ,j_{4} \leq r}\lambda _{2 j_{1} j_{2} j_{3} j_{4}}\left \{x_{j_{1}} x_{j_{2}} x_{j_{4}} x_{j_{3}} -i x_{j_{1} +r} x_{j_{2}} x_{j_{4}} x_{j_{3}} -i x_{j_{1}} x_{j_{2} +r} x_{j_{4}} x_{j_{3}} -x_{j_{1} +r} x_{j_{2} +r} x_{j_{4}} x_{j_{3}}\right . \nonumber  \\
 &  +i x_{j_{1}} x_{j_{2}} x_{j_{4} +r} x_{j_{3}} +x_{j_{1} +r} x_{j_{2}} x_{j_{4} +r} x_{j_{3}} +x_{j_{1}} x_{j_{2} +r} x_{j_{4} +r} x_{j_{3}} -i x_{j_{1} +r} x_{j_{2} +r} x_{j_{4} +r} x_{j_{3}} \nonumber  \\
 &  +i x_{j_{1}} x_{j_{2}} x_{j_{4}} x_{j_{3} +r} +x_{j_{1} +r} x_{j_{2}} x_{j_{4}} x_{j_{3} +r} +x_{j_{1}} x_{j_{2} +r} x_{j_{4}} x_{j_{3} +r} -i x_{j_{1} +r} x_{j_{2} +r} x_{j_{4}} x_{j_{3} +r} \nonumber  \\
 & \left . -x_{j_{1}} x_{j_{2}} x_{j_{4} +r} x_{j_{3} +r} +i x_{j_{1} +r} x_{j_{2}} x_{j_{4} +r} x_{j_{3} +r} +i x_{j_{1}} x_{j_{2} +r} x_{j_{4} +r} x_{j_{3} +r} +x_{j_{1} +r} x_{j_{2} +r} x_{j_{4} +r} x_{j_{3} +r}\right \}\text{.}\end{align}(Note that $\lambda _{2}$ is not homogeneous in $\left \{x_{j}\right \})\text{.}$ 

Let $r =4$ then
\begin{align*}\lambda _{2} &  =\sum _{1 \leq j_{1} <j_{2} \leq 4 ;1 \leq  ,j_{3} <j_{4} \leq 4}\lambda _{2 j_{1} j_{2} j_{3} j_{4}} a_{j_{1}}^{ \dag } a_{j_{2}}^{ \dag } a_{j_{4}} a_{j_{3}} \\
 &  =\frac{1}{4} \sum _{1 \leq j_{1} <j_{2} \leq 4 ;1 \leq  ,j_{3} <j_{4} \leq 4}\lambda _{2 j_{1} j_{2} j_{3} j_{4}} \left (x_{j_{1}} -i x_{j_{1} +r}\right ) \left (x_{j_{2}} -i x_{j_{2} +r}\right ) \left (x_{j_{4}} +i x_{j_{4} +r}\right ) \left (x_{j_{3}} +i x_{j_{3} +r}\right )\end{align*}
\begin{align} &  =\frac{1}{4} \sum _{1 \leq j_{1} <j_{2} \leq 4 ;1 \leq  ,j_{3} <j_{4} \leq 4}\lambda _{2 j_{1} j_{2} j_{3} j_{4}}\left \{x_{j_{1}} x_{j_{2}} x_{j_{4}} x_{j_{3}} -i x_{j_{1} +r} x_{j_{2}} x_{j_{4}} x_{j_{3}} -i x_{j_{1}} x_{j_{2} +r} x_{j_{4}} x_{j_{3}} -x_{j_{1} +r} x_{j_{2} +r} x_{j_{4}} x_{j_{3}}\right . \nonumber  \\
 &  +i x_{j_{1}} x_{j_{2}} x_{j_{4} +r} x_{j_{3}} +x_{j_{1} +r} x_{j_{2}} x_{j_{4} +r} x_{j_{3}} +x_{j_{1}} x_{j_{2} +r} x_{j_{4} +r} x_{j_{3}} -i x_{j_{1} +r} x_{j_{2} +r} x_{j_{4} +r} x_{j_{3}} \nonumber  \\
 &  +i x_{j_{1}} x_{j_{2}} x_{j_{4}} x_{j_{3} +r} +x_{j_{1} +r} x_{j_{2}} x_{j_{4}} x_{j_{3} +r} +x_{j_{1}} x_{j_{2} +r} x_{j_{4}} x_{j_{3} +r} -i x_{j_{1} +r} x_{j_{2} +r} x_{j_{4}} x_{j_{3} +r} \nonumber  \\
 & \left . -x_{j_{1}} x_{j_{2}} x_{j_{4} +r} x_{j_{3} +r} +i x_{j_{1} +r} x_{j_{2}} x_{j_{4} +r} x_{j_{3} +r} +i x_{j_{1}} x_{j_{2} +r} x_{j_{4} +r} x_{j_{3} +r} +x_{j_{1} +r} x_{j_{2} +r} x_{j_{4} +r} x_{j_{3} +r}\right \}\text{.}\end{align}

\bigskip $\lceil $Zero differences
\begin{align} &  =\frac{1}{4} \lambda _{2 1 2 1 2}\left \{x_{1} x_{2} x_{2} x_{1} -i x_{5} x_{2} x_{2} x_{1} -i x_{1} x_{6} x_{2} x_{1} -x_{5} x_{6} x_{2} x_{1}\right . \nonumber  \\
 &  +i x_{1} x_{2} x_{6} x_{1} +x_{5} x_{2} x_{6} x_{1} +x_{1} x_{6} x_{6} x_{1} -i x_{5} x_{6} x_{6} x_{1} \nonumber  \\
 &  +i x_{1} x_{2} x_{2} x_{5} +x_{5} x_{2} x_{2} x_{5} +x_{1} x_{6} x_{2} x_{5} -i x_{5} x_{6} x_{2} x_{5} \nonumber  \\
 & \left . -x_{1} x_{2} x_{6} x_{5} +i x_{5} x_{2} x_{6} x_{5} +i x_{1} x_{6} x_{6} x_{5} +x_{5} x_{6} x_{6} x_{5}\right \}\text{.}\end{align}


\begin{align} &  =\frac{1}{4} \lambda _{2 1 2 1 2}\left \{\frac{1}{4} I_{\mathcal{F}} +i \frac{1}{2} x_{1} x_{5} +i \frac{1}{2} x_{2} x_{6} +x_{5} x_{2} x_{6} x_{1}\right . \nonumber  \\
 &  +i \frac{1}{2} x_{2} x_{6} +x_{5} x_{2} x_{6} x_{1} +\frac{1}{4} I_{\mathcal{F}} +i \frac{1}{2} x_{1} x_{5} \nonumber  \\
 &  +i \frac{1}{2} x_{1} x_{5} +\frac{1}{4} I_{\mathcal{F}} +x_{5} x_{2} x_{6} x_{1} +i \frac{1}{2} x_{2} x_{6} \nonumber  \\
 &  +\left .x_{5} x_{2} x_{1} x_{6} +i \frac{1}{2} x_{2} x_{6} +i \frac{1}{2} x_{1} x_{5} +\frac{1}{4} I_{\mathcal{F}}\right \}\text{.}\end{align}
\begin{equation} =\lambda _{2 1 2 1 2} \left \{\frac{1}{4} I_{\mathcal{F}} +\frac{i}{2} x_{1} x_{5} +\frac{i}{2} i x_{2} x_{6} +x_{1} x_{2} x_{5} x_{6}\right \}
\end{equation}

\bigskip One difference 


\begin{align} &  =\frac{1}{4} \lambda _{2 1 2 1 3}\left \{x_{1} x_{2} x_{3} x_{1} -i x_{5} x_{2} x_{3} x_{1} -i x_{1} x_{6} x_{3} x_{1} -x_{5} x_{6} x_{3} x_{1}\right . \nonumber  \\
 &  +i x_{1} x_{2} x_{7} x_{1} +x_{5} x_{2} x_{7} x_{1} +x_{1} x_{6} x_{7} x_{1} -i x_{5} x_{6} x_{7} x_{1} \nonumber  \\
 &  +i x_{1} x_{2} x_{3} x_{5} +x_{5} x_{2} x_{3} x_{5} +x_{1} x_{6} x_{3} x_{5} -i x_{5} x_{6} x_{3} x_{5} \nonumber  \\
 & \left . -x_{1} x_{2} x_{7} x_{5} +i x_{5} x_{2} x_{7} x_{5} +i x_{1} x_{6} x_{7} x_{5} +x_{5} x_{6} x_{7} x_{5}\right \}\text{.}\end{align}


\begin{align} &  =\frac{1}{4} \lambda _{2 1 2 1 3}\left \{\frac{1}{2} x_{2} x_{3} -i x_{5} x_{2} x_{3} x_{1} -i \frac{1}{2} x_{6} x_{3} -x_{5} x_{6} x_{3} x_{1}\right . \nonumber  \\
 &  +i \frac{1}{2} x_{2} x_{7} +x_{5} x_{2} x_{7} x_{1} +\frac{1}{2} x_{6} x_{7} -i x_{5} x_{6} x_{7} x_{1} \nonumber  \\
 &  +i x_{1} x_{2} x_{3} x_{5} +\frac{1}{2} x_{2} x_{3} +x_{1} x_{6} x_{3} x_{5} -i \frac{1}{2} x_{6} x_{3} \nonumber  \\
 & \left . -x_{1} x_{2} x_{7} x_{5} +i \frac{1}{2} x_{2} x_{7} +i x_{1} x_{6} x_{7} x_{5} +\frac{1}{2} x_{6} x_{7}\right \}\text{.}\end{align}


\begin{align*} &  =\frac{1}{4} \lambda _{2 1 2 1 3}\left \{x_{2} x_{3} +x_{6} x_{7} +2 x_{1} x_{3} x_{5} x_{6} +2 x_{1} x_{2} x_{5} x_{7} +i x_{3} x_{6}\right . \\
 &  +i x_{2} x_{7} -i x_{1} x_{2} x_{3} x_{5} -i x_{1} x_{5} x_{6} x_{7} +i x_{1} x_{2} x_{3} x_{5} +\left .i x_{1} x_{5} x_{6} x_{7}\right \} .\end{align*}

\bigskip Two differences 


\begin{align} &  =\frac{1}{4} \lambda _{2 1 2 4 3}\left \{x_{1} x_{2} x_{3} x_{4} -i x_{5} x_{2} x_{3} x_{4} -i x_{1} x_{6} x_{3} x_{4} -x_{5} x_{6} x_{3} x_{4}\right . \nonumber  \\
 &  +i x_{1} x_{2} x_{8} x_{3} +x_{5} x_{2} x_{8} x_{3} +x_{1} x_{6} x_{8} x_{3} -i x_{5} x_{6} x_{8} x_{3} \nonumber  \\
 &  +i x_{1} x_{2} x_{4} x_{8} +x_{5} x_{2} x_{4} x_{8} +x_{1} x_{6} x_{4} x_{8} -i x_{5} x_{6} x_{4} x_{8} \nonumber  \\
 & \left . -x_{1} x_{2} x_{7} x_{8} +i x_{5} x_{2} x_{7} x_{8} +i x_{1} x_{6} x_{7} x_{8} +x_{5} x_{6} x_{7} x_{8}\right \}\text{.}\end{align}


\begin{align} &  =\frac{1}{4} \lambda _{2 1 2 4 3}\left \{x_{1} x_{2} x_{3} x_{4} -x_{3} x_{4} x_{5} x_{6} -x_{2} x_{3} x_{5} x_{8} +x_{1} x_{3} x_{6} x_{8}\right . \nonumber  \\
 &  +x_{2} x_{4} x_{5} x_{8} -x_{1} x_{4} x_{6} x_{8} -x_{1} x_{2} x_{7} x_{8} +x_{5} x_{6} x_{7} x_{8} \nonumber  \\
 &  +i x_{2} x_{3} x_{4} x_{5} -i x_{1} x_{3} x_{4} x_{6} -i x_{1} x_{2} x_{3} x_{8} +i x_{3} x_{5} x_{6} x_{8} \\
 &  +i x_{1} x_{2} x_{4} x_{8} -i x_{4} x_{5} x_{6} x_{8}\left . -i x_{2} x_{5} x_{7} x_{8} +i x_{1} x_{6} x_{7} x_{8}\right \} . \nonumber \end{align}$\rfloor $ 

The zeroth degree terms are
\begin{align} & \frac{1}{16} \sum _{1 \leq j_{1} ,j_{2} \leq r}\lambda _{2 j_{1} j_{2} j_{1} j_{2}} \left \{x_{j_{1}} x_{j_{2}} x_{j_{2}} x_{j_{1}} +x_{j_{1}} x_{j_{2} +r} x_{j_{2} +r} x_{j_{1}} +x_{j_{1} +r} x_{j_{2}} x_{j_{2}} x_{j_{1} +r} +x_{j_{1} +r} x_{j_{2} +r} x_{j_{2} +r} x_{j_{1} +r}\right \} \nonumber  \\
 &  =\frac{1}{4} \sum _{1 \leq j_{1} <j_{2} \leq r}\lambda _{2 j_{1} j_{2} j_{1} j_{2}} I_{\mathcal{F}} =\frac{1}{4} \left .\ensuremath{\operatorname*{Tr}}\right ._{2} \left \{\lambda _{2}\right \} I_{\mathcal{F}}\text{.}\end{align}

The second degree terms are after setting $j_{3} =j_{1} ,j_{3} =j_{2} ,j_{4} =j_{1}\text{,}$ and $j_{4} =j_{2}\text{,}$ then $j_{4} =j_{3}$ and interchanging $j_{1}$ and $j_{2}$ (in the following sums we can consider $\left (j_{2} ,j_{2} +r\right )$ but not $\left (j_{2} ,j_{2}\right ))$
\begin{align*} & \frac{1}{16} \sum _{1 \leq j_{1} ,j_{2} \neq j_{3} \leq r}\lambda _{2 j_{1} j_{2} j_{1} j_{3}}\left \{x_{j_{1}} x_{j_{2}} x_{j_{3}} x_{j_{1}} +x_{j_{1}} x_{j_{2}} x_{j_{3}} x_{j_{1}} -i x_{j_{1}} x_{j_{2} +r} x_{j_{3}} x_{j_{1}}\right . \\
 &  +i x_{j_{1}} x_{j_{2}} x_{j_{3} +r} x_{j_{1}} +x_{j_{1}} x_{j_{2} +r} x_{j_{3} +r} x_{j_{1}} -i x_{j_{1} +r} x_{j_{2} +r} x_{j_{3}} x_{j_{1} +r} \\
 + & \left .i x_{j_{1} +r} x_{j_{2}} x_{j_{3} +r} x_{j_{1} +r} +x_{j_{1} +r} x_{j_{2} +r} x_{j_{3} +r} x_{j_{1} +r}\right \}\end{align*}
\begin{align*} - & \frac{1}{16} \sum _{1 \leq j_{1} ,j_{2} \neq j_{3} \leq r}\lambda _{2 j_{1} j_{2} j_{1} j_{3}}\left \{x_{j_{1}} x_{j_{2}} x_{j_{1}} x_{j_{3}} -i x_{j_{1}} x_{j_{2} +r} x_{j_{1}} x_{j_{3}}\right . \\
 &  +x_{j_{1} +r} x_{j_{2}} x_{j_{1} +r} x_{j_{3}} -i x_{j_{1} +r} x_{j_{2} +r} x_{j_{1} +r} x_{j_{3}} \\
 &  +i x_{j_{1}} x_{j_{2}} x_{j_{1}} x_{j_{3} +r} +x_{j_{1}} x_{j_{2} +r} x_{j_{1}} x_{j_{3} +r} \\
 &  +\left .i x_{j_{1} +r} x_{j_{2}} x_{j_{1} +r} x_{j_{3} +r} +x_{j_{1} +r} x_{j_{2} +r} x_{j_{1} +r} x_{j_{3} +r}\right \}\end{align*}
\begin{align*} - & \frac{1}{16} \sum _{1 \leq j_{1} ,j_{2} \neq j_{3} \leq r}\lambda _{2 j_{1} j_{2} j_{1} j_{3}}\left \{x_{j_{2}} x_{j_{1}} x_{j_{3}} x_{j_{2}} -i x_{j_{2} +r} x_{j_{1}} x_{j_{3}} x_{j_{1}}\right . \\
 &  +i x_{j_{1}} x_{j_{1}} x_{j_{3} +r} x_{j_{1}} +x_{j_{2} +r} x_{j_{1}} x_{j_{3} +r} x_{j_{1}} \\
 &  +x_{j_{2}} x_{j_{1} +r} x_{j_{3}} x_{j_{1} +r} -i x_{j_{2} +r} x_{j_{1} +r} x_{j_{3}} x_{j_{1} +r} \\
 &  +\left .i x_{j_{2}} x_{j_{1} +r} x_{j_{3} +r} x_{j_{1} +r} +x_{j_{2} +r} x_{j_{1} +r} x_{j_{3} +r} x_{j_{1} +r}\right \}\end{align*}
\begin{align*} + & \frac{1}{16} \sum _{1 \leq j_{1} ,j_{2} \neq j_{3} \leq r}\lambda _{2 j_{1} j_{2} j_{1} j_{3}}\left \{x_{j_{2}} x_{j_{1}} x_{j_{1}} x_{j_{3}} -i x_{j_{2} +r} x_{j_{1}} x_{j_{1}} x_{j_{3}}\right . \\
 &  +x_{j_{2}} x_{j_{1} +r} x_{j_{1} +r} x_{j_{3}} -i x_{j_{2} +r} x_{j_{1} +r} x_{j_{1} +r} x_{j_{3}} \\
 &  +i x_{j_{2}} x_{j_{1}} x_{j_{1}} x_{j_{3} +r} +x_{j_{2} +r} x_{j_{1}} x_{j_{1}} x_{j_{3} +r} \\
 &  +\left .i x_{j_{2}} x_{j_{1} +r} x_{j_{1} +r} x_{j_{3} +r} +x_{j_{2} +r} x_{j_{1} +r} x_{j_{1} +r} x_{j_{3} +r}\right \}\end{align*}
\begin{align*} = & \frac{1}{32} \sum _{1 \leq j_{1} ,j_{2} \neq j_{3} \leq r}\lambda _{2 j_{1} j_{2} j_{1} j_{3}}\left \{x_{j_{2}} x_{j_{3}} -i x_{j_{2} +r} x_{j_{3}}\right . \\
 &  +i x_{j_{2}} x_{j_{3} +r} +x_{j_{2} +r} x_{j_{3} +r} -i x_{j_{2} +r} x_{j_{3}} +\left .i x_{j_{2}} x_{j_{3} +r} +x_{j_{2} +r} x_{j_{3} +r}\right \}\end{align*}
\begin{align*} &  +\frac{1}{32} \sum _{1 \leq j_{1} ,j_{2} \neq j_{3} \leq r}\lambda _{2 j_{1} j_{2} j_{1} j_{3}}\left \{x_{j_{2}} x_{j_{3}} -i x_{j_{2} +r} x_{j_{3}}\right . \\
 &  +x_{j_{2}} x_{j_{3}} -i x_{j_{2} +r} x_{j_{3}} +i x_{j_{2}} x_{j_{3} +r} +x_{j_{2} +r} x_{j_{3} +r} +\left .i x_{j_{2}} x_{j_{3} +r} +x_{j_{2} +r} x_{j_{3} +r}\right \}\end{align*}
\begin{align*} + & \frac{1}{32} \sum _{1 \leq j_{1} ,j_{2} \neq j_{3} \leq r}\left \{x_{j_{2}} x_{j_{3}} -i x_{j_{2} +r} x_{j_{3}}\right . \\
 &  +i x_{j_{1}} x_{j_{3} +r} +x_{j_{2} +r} x_{j_{3} +r} +x_{j_{2}} x_{j_{3}} -i x_{j_{2} +r} x_{j_{3}} +\left .i x_{j_{2}} x_{j_{3} +r} +x_{j_{2} +r} x_{j_{3} +r}\right \}\end{align*}
\begin{align*} + & \frac{1}{32} \sum _{1 \leq j_{1} ,j_{2} \neq j_{3} \leq r}\left \{x_{j_{2}} x_{j_{3}} -i x_{j_{2} +r} x_{j_{3}}\right . \\
 &  +x_{j_{2}} x_{j_{3}} -i x_{j_{2} +r} x_{j_{3}} +i x_{j_{2}} x_{j_{3} +r} +x_{j_{2} +r} x_{j_{3} +r} +\left .i x_{j_{2}} x_{j_{3} +r} +x_{j_{2} +r} x_{j_{3} +r}\right \}\end{align*}
\begin{equation*} =\frac{1}{4} \sum _{1 \leq j_{2} \neq j_{3} \leq r}L_{2}^{1} \left (\lambda _{2}\right )_{j_{2} j_{3}} \left \{x_{j_{2}} x_{j_{3}} +x_{j_{2} +r} x_{j_{3} +r} +i x_{j_{3}} x_{j_{2} +r} +i x_{j_{2}} x_{j_{3} +r}\right \}
\end{equation*}as the products are antisymmetric this gives that the second degree terms are
\begin{equation}\frac{1}{2} \sum _{1 \leq j_{2} <j_{3} \leq r}\left (L_{2}^{1} \left (\lambda _{2}\right )_{j_{2} j_{3}} -L_{2}^{1} \left (\lambda _{2}\right )_{j_{3} j_{2}}\right ) \left \{x_{j_{2}} x_{j_{3}} +x_{j_{2} +r} x_{j_{3} +r} +i \left (x_{j_{3}} x_{j_{2} +r} +x_{j_{2}} x_{j_{3} +r}\right )\right \}\text{.}
\end{equation}$\lceil $
\begin{equation}\frac{1}{2} \left (L_{2}^{1} \left (\lambda _{2}\right )_{1 2} -L_{2}^{1} \left (\lambda _{2}\right )_{2 1}\right ) \left \{x_{1} x_{2} +x_{3} x_{4} +i \left (x_{2} x_{3} +x_{2} x_{4}\right )\right \}
\end{equation}$\rfloor $ 

The fourth degree terms are
\begin{align*} & \frac{1}{16} \sum _{1 \leq j_{1} \neq j_{2} \neq j_{3} \neq j_{4} \leq r}\lambda _{2 j_{1} j_{2} j_{3} j_{4}}\left \{x_{j_{1}} x_{j_{2}} x_{j_{4}} x_{j_{3}} -x_{j_{1} +r} x_{j_{2} +r} x_{j_{4}} x_{j_{3}} +x_{j_{1} +r} x_{j_{2}} x_{j_{4} +r} x_{j_{3}} +x_{j_{1}} x_{j_{2} +r} x_{j_{4} +r} x_{j_{3}}\right . \\
 &  -x_{j_{1}} x_{j_{2}} x_{j_{4} +r} x_{j_{3} +r} +x_{j_{1} +r} x_{j_{2}} x_{j_{4}} x_{j_{3} +r} +x_{j_{1}} x_{j_{2} +r} x_{j_{4}} x_{j_{3} +r} +x_{j_{1} +r} x_{j_{2} +r} x_{j_{4} +r} x_{j_{3} +r} \\
 &  +i x_{j_{1}} x_{j_{2}} x_{j_{4} +r} x_{j_{3}} -i x_{j_{1} +r} x_{j_{2}} x_{j_{4}} x_{j_{3}} -i x_{j_{1}} x_{j_{2} +r} x_{j_{4}} x_{j_{3}} \\
 &  +i x_{j_{1}} x_{j_{2}} x_{j_{4}} x_{j_{3} +r} -i x_{j_{1} +r} x_{j_{2} +r} x_{j_{4}} x_{j_{3} +r} +\left .i x_{j_{1}} x_{j_{2} +r} x_{j_{4} +r} x_{j_{3} +r} -i x_{j_{1} +r} x_{j_{2} +r} x_{j_{4} +r} x_{j_{3}} +i x_{j_{1} +r} x_{j_{2}} x_{j_{4} +r} x_{j_{3} +r}\right \}\text{.}\end{align*}

\subsubsection{General Even Fourth Degree Operators}
Two particle operators are \emph{even }fourth degree operators, which are fourth degree polynomials in the generators of the
Clifford algebra
\begin{equation}A_{2} =\alpha _{0} I_{\mathcal{F}} +\sum _{1 \leq j_{1} <j_{2} \leq 2 r}\alpha _{2 j_{1} j_{2}} x_{j_{1}} x_{j_{2}} +\sum _{1 \leq j_{1} <j_{2} <j_{3} <j_{4} \leq 2 r}\alpha _{4 j_{1} j_{2} j_{3} j_{4}} x_{j_{1}} x_{j_{2}} x_{j_{3}} x_{j_{4}}\text{.}
\end{equation}The square of such operators is
\begin{align}\left (A_{2}\right )^{2} &  =\left (A_{2 0} +A_{2 2} +A_{2 4}\right )^{2} =\left (A_{2 0} +A_{2 2} +A_{2 4}\right ) \left (A_{2 0} +A_{2 2} +A_{2 4}\right ) \nonumber  \\
 &  =A_{2 0} A_{2 0} +A_{2 0} A_{2 2} +A_{2 0} A_{2 4} +A_{2 2} A_{2 0} +A_{2 2} A_{2 2} +A_{2 2} A_{2 4} +A_{2 4} A_{2 0} +A_{2 4} A_{2 2} +A_{2 4} A_{2 4} \nonumber  \\
 &  =\alpha _{0}^{2} I_{\mathcal{F}} +2 \alpha _{0} A_{2 2} +2 \alpha _{0} A_{2 4} +\left (A_{2 2}\right )^{2} +A_{2 2} A_{2 4} +A_{2 4} A_{2 2} +\left (A_{2 4}\right )^{2}\text{,}\end{align}where
\begin{align}\left (A_{2 0}\right )^{2} &  =\alpha _{0}^{2} I_{\mathcal{F}} \\
A_{2 0} A_{2 2} &  =\alpha _{0} A_{2 2} \\
A_{2 0} A_{2 4} &  =\alpha _{0} A_{2 4}\end{align}and 

\begin{equation}A_{2 2} A_{2 2} =\sum _{\begin{array}{c}1 \leq j_{1} <j_{2} \leq 2 r \\
1 \leq j_{3} <j_{4} \leq 2 r\end{array}}\alpha _{2 j_{1} j_{2}} \alpha _{2 j_{3} j_{4}} x_{j_{1}} x_{j_{2}} x_{j_{3}} x_{j_{4}}\text{.}
\end{equation}$\lceil $In the case $r =2$ one has
\begin{align} & \alpha _{2 1 2} \alpha _{2 1 2} x_{1} x_{2} x_{1} x_{2} +\alpha _{2 1 2} \alpha _{2 1 3} x_{1} x_{2} x_{1} x_{3} +\alpha _{2 1 2} \alpha _{2 1 4} x_{1} x_{2} x_{1} x_{4} \nonumber  \\
 &  +\alpha _{2 1 2} \alpha _{2 2 3} x_{1} x_{2} x_{2} x_{3} +\alpha _{2 1 2} \alpha _{2 2 4} x_{1} x_{2} x_{2} x_{4} +\alpha _{2 1 2} \alpha _{2 3 4} x_{1} x_{2} x_{3} x_{4} \nonumber  \\
 &  +\alpha _{2 1 3} \alpha _{2 1 2} x_{1} x_{3} x_{1} x_{2} +\alpha _{2 1 3} \alpha _{2 1 3} x_{1} x_{3} x_{1} x_{3} +\alpha _{2 1 3} \alpha _{2 1 4} x_{1} x_{3} x_{1} x_{4} \nonumber  \\
 &  +\alpha _{2 1 3} \alpha _{2 2 3} x_{1} x_{3} x_{2} x_{3} +\alpha _{2 1 3} \alpha _{2 2 4} x_{1} x_{3} x_{2} x_{4} +\alpha _{2 1 3} \alpha _{2 3 4} x_{1} x_{3} x_{3} x_{4} \nonumber  \\
 &  +\alpha _{2 1 4} \alpha _{2 1 2} x_{1} x_{4} x_{1} x_{2} +\alpha _{2 1 4} \alpha _{2 1 3} x_{1} x_{4} x_{1} x_{3} +\alpha _{2 1 4} \alpha _{2 1 4} x_{1} x_{4} x_{1} x_{4} \nonumber  \\
 &  +\alpha _{2 1 4} \alpha _{2 2 3} x_{1} x_{4} x_{2} x_{3} +\alpha _{2 1 4} \alpha _{2 2 4} x_{1} x_{4} x_{2} x_{4} +\alpha _{2 1 4} \alpha _{2 3 4} x_{1} x_{4} x_{3} x_{4} \nonumber  \\
 &  +\alpha _{2 2 3} \alpha _{2 1 2} x_{2} x_{3} x_{1} x_{2} +\alpha _{2 2 3} \alpha _{2 1 3} x_{2} x_{3} x_{1} x_{3} +\alpha _{2 2 3} \alpha _{2 1 4} x_{2} x_{3} x_{1} x_{4} \nonumber  \\
 &  +\alpha _{2 2 3} \alpha _{2 2 3} x_{2} x_{3} x_{2} x_{3} +\alpha _{2 2 3} \alpha _{2 2 4} x_{2} x_{3} x_{2} x_{4} +\alpha _{2 2 3} \alpha _{2 3 4} x_{2} x_{3} x_{3} x_{4} \nonumber  \\
 &  +\alpha _{2 2 4} \alpha _{2 1 2} x_{2} x_{4} x_{1} x_{2} +\alpha _{2 2 4} \alpha _{2 1 3} x_{2} x_{4} x_{1} x_{3} +\alpha _{2 2 4} \alpha _{2 1 4} x_{2} x_{4} x_{1} x_{4} \nonumber  \\
 &  +\alpha _{2 2 4} \alpha _{2 2 3} x_{2} x_{4} x_{2} x_{3} +\alpha _{2 2 4} \alpha _{2 2 4} x_{2} x_{4} x_{2} x_{4} +\alpha _{2 2 4} \alpha _{2 3 4} x_{2} x_{4} x_{3} x_{4} \nonumber  \\
 &  +\alpha _{2 3 4} \alpha _{2 1 2} x_{3} x_{4} x_{1} x_{2} +\alpha _{2 3 4} \alpha _{2 1 3} x_{3} x_{4} x_{1} x_{3} +\alpha _{2 3 4} \alpha _{2 1 4} x_{3} x_{4} x_{1} x_{4} \nonumber  \\
 &  +\alpha _{2 3 4} \alpha _{2 2 3} x_{3} x_{4} x_{2} x_{3} +\alpha _{2 3 4} \alpha _{2 2 4} x_{3} x_{4} x_{2} x_{4} +\alpha _{2 3 4} \alpha _{2 3 4} x_{3} x_{4} x_{3} x_{4}\end{align}
\begin{align} = &  -\frac{1}{4} \alpha _{2 1 2} \alpha _{2 1 2} I_{\mathcal{F}} -\frac{1}{2} \alpha _{2 1 2} \alpha _{2 1 3} x_{2} x_{3} -\frac{1}{2} \alpha _{2 1 2} \alpha _{2 1 4} x_{2} x_{4} +\frac{1}{2} \alpha _{2 1 2} \alpha _{2 2 3} x_{1} x_{3} +\frac{1}{2} \alpha _{2 1 2} \alpha _{2 2 4} x_{1} x_{4} +\alpha _{2 1 2} \alpha _{2 3 4} x_{1} x_{2} x_{3} x_{4} \nonumber  \\
 &  -\frac{1}{2} \alpha _{2 1 3} \alpha _{2 1 2} x_{3} x_{2} -\alpha _{2 1 3} \alpha _{2 1 3} \frac{1}{4} I_{\mathcal{F}} -\frac{1}{2} \alpha _{2 1 3} \alpha _{2 1 4} x_{3} x_{4} -\frac{1}{2} \alpha _{2 1 3} \alpha _{2 2 3} x_{1} x_{2} -\alpha _{2 1 3} \alpha _{2 2 4} x_{1} x_{2} x_{3} x_{4} +\frac{1}{2} \alpha _{2 1 3} \alpha _{2 3 4} x_{1} x_{4} \nonumber  \\
 &  -\frac{1}{2} \alpha _{2 1 4} \alpha _{2 1 2} x_{4} x_{2} -\frac{1}{2} \alpha _{2 1 4} \alpha _{2 1 3} x_{4} x_{3} -\alpha _{2 1 4} \alpha _{2 1 4} \frac{1}{4} I_{\mathcal{F}} +\alpha _{2 1 4} \alpha _{2 2 3} x_{1} x_{2} x_{3} x_{4} -\frac{1}{2} \alpha _{2 1 4} \alpha _{2 2 4} x_{1} x_{2} -\frac{1}{2} \alpha _{2 1 4} \alpha _{2 3 4} x_{1} x_{3} \nonumber  \\
 &  +\frac{1}{2} \alpha _{2 2 3} \alpha _{2 1 2} x_{3} x_{1} -\frac{1}{2} \alpha _{2 2 3} \alpha _{2 1 3} x_{2} x_{1} +\alpha _{2 2 3} \alpha _{2 1 4} x_{1} x_{2} x_{3} x_{4} -\alpha _{2 2 3} \alpha _{2 2 3} \frac{1}{4} I_{\mathcal{F}} -\frac{1}{2} \alpha _{2 2 3} \alpha _{2 2 4} x_{3} x_{4} +\frac{1}{2} \alpha _{2 2 3} \alpha _{2 3 4} x_{2} x_{4} \nonumber  \\
 &  +\frac{1}{2} \alpha _{2 2 4} \alpha _{2 1 2} x_{4} x_{1} -\alpha _{2 2 4} \alpha _{2 1 3} x_{1} x_{2} x_{3} x_{4} -\frac{1}{2} \alpha _{2 2 4} \alpha _{2 1 4} x_{2} x_{1} -\frac{1}{2} \alpha _{2 2 4} \alpha _{2 2 3} x_{4} x_{3} -\alpha _{2 2 4} \alpha _{2 2 4} \frac{1}{4} I_{\mathcal{F}} -\alpha _{2 2 4} \alpha _{2 3 4} \frac{1}{2} x_{2} x_{3} \nonumber  \\
 &  +\alpha _{2 3 4} \alpha _{2 1 2} x_{1} x_{2} x_{3} x_{4} +\frac{1}{2} \alpha _{2 3 4} \alpha _{2 1 3} x_{4} x_{1} -\frac{1}{2} \alpha _{2 3 4} \alpha _{2 1 4} x_{3} x_{1} +\frac{1}{2} \alpha _{2 3 4} \alpha _{2 2 3} x_{4} x_{2} -\frac{1}{2} \alpha _{2 3 4} \alpha _{2 2 4} x_{3} x_{2} -\alpha _{2 3 4} \alpha _{2 3 4} \frac{1}{4} I_{\mathcal{F}}\end{align}
\begin{align} &  -\frac{1}{4} I_{\mathcal{F}} \left \{\alpha _{2 1 2} \alpha _{2 1 2} +\alpha _{2 1 3} \alpha _{2 1 3} +\alpha _{2 1 4} \alpha _{2 1 4} +\alpha _{2 2 3} \alpha _{2 2 3} +\alpha _{2 2 4} \alpha _{2 2 4} +\alpha _{2 3 4} \alpha _{2 3 4}\right \} \nonumber  \\
 &  +\frac{1}{2} \left \{ -\alpha _{2 1 2} \alpha _{2 1 3} +\alpha _{2 1 3} \alpha _{2 1 2} -\alpha _{2 2 4} \alpha _{2 3 4} +\alpha _{2 3 4} \alpha _{2 2 4}\right \} x_{2} x_{3} \nonumber  \\
 &  +\frac{1}{2} \left \{ -\alpha _{2 1 2} \alpha _{2 1 4} +\alpha _{2 1 2} \alpha _{2 1 4} +\alpha _{2 2 3} \alpha _{2 3 4} -\alpha _{2 3 4} \alpha _{2 2 3}\right \} x_{2} x_{4} \nonumber  \\
 &  +\frac{1}{2} \left \{\alpha _{2 1 2} \alpha _{2 2 3} -\alpha _{2 1 4} \alpha _{2 3 4} -\alpha _{2 2 3} \alpha _{2 1 2} +\alpha _{2 3 4} \alpha _{2 1 4}\right \} x_{1} x_{3} \nonumber  \\
 &  +\frac{1}{2} \left \{\alpha _{2 1 2} \alpha _{2 2 4} -\alpha _{2 2 4} \alpha _{2 1 2} -\alpha _{2 3 4} \alpha _{2 1 3} +\alpha _{2 1 3} \alpha _{2 3 4}\right \} x_{1} x_{4} \nonumber  \\
 &  +\frac{1}{2} \left \{ -\alpha _{2 1 3} \alpha _{2 1 4} +\alpha _{2 1 4} \alpha _{2 1 3} -\alpha _{2 2 3} \alpha _{2 2 4} +\alpha _{2 2 4} \alpha _{2 2 3}\right \} x_{3} x_{4} \nonumber  \\
 &  +\frac{1}{2} \left \{ -\alpha _{2 1 3} \alpha _{2 2 3} -\alpha _{2 1 4} \alpha _{2 2 4} +\alpha _{2 2 3} \alpha _{2 1 3} +\alpha _{2 2 4} \alpha _{2 1 4}\right \} x_{1} x_{2} \nonumber  \\
 &  +2 \left \{\alpha _{2 1 2} \alpha _{2 3 4} -\alpha _{2 1 3} \alpha _{2 2 4} +\alpha _{2 1 4} \alpha _{2 2 3}\right \} x_{1} x_{2} x_{3} x_{4}\end{align}
\begin{align} &  = -\frac{1}{4} \left (\sum _{1 \leq j_{1} <j_{2} \leq 2 r}\alpha _{2 j_{1} j_{2}}^{2}\right ) I_{\mathcal{F}} +2 \ensuremath{\operatorname*{Pf}}\left \{\mathbf{\alpha }_{2}\right \} x_{1} x_{2} x_{3} x_{4} \nonumber  \\
 &  =\frac{1}{8} \left .\ensuremath{\operatorname*{Tr}}\right \vert _{1} \left \{\mathbf{\alpha }_{2}^{t} \mathbf{\alpha }_{2}\right \} I_{\mathcal{F}} +2 \ensuremath{\operatorname*{Pf}}\left \{\mathbf{\alpha }_{2}\right \} x_{1} x_{2} x_{3} x_{4}\text{.}\end{align}One should observe
\begin{equation}2 \ensuremath{\operatorname*{Pf}}\left \{\mathbf{\alpha }_{2}\right \} x_{1} x_{2} x_{3} x_{4} =\alpha _{2} \wedge \alpha _{2}\text{,}
\end{equation}however
\begin{equation}\alpha _{2} \wedge \alpha _{2} \neq \mathbf{\alpha }_{2} \boxtimes \mathbf{\alpha }_{2} =\mathcal{C}_{2} \left (\mathbf{\alpha }_{2}\right )\text{.}
\end{equation}$\rfloor $
\begin{align}A_{2 2} A_{2 2} &  =\frac{1}{8} \left .\ensuremath{\operatorname*{Tr}}\right \vert _{1} \left \{\mathbf{\alpha }_{2}^{t} \mathbf{\alpha }_{2}\right \} I_{\mathcal{F}} +2 \sum _{\begin{array}{c}1 \leq j_{1} <j_{2} \leq 2 r \\
1 \leq j_{3} <j_{4} \leq 2 r\end{array}}\ensuremath{\operatorname*{Pf}}\left \{\mathbf{\alpha }_{2}\right \}_{j_{1} j_{2} j_{3} j_{4}} x_{j_{1}} x_{j_{2}} x_{j_{3}} x_{j_{4}} \\
 &  =\frac{1}{8} \left .\ensuremath{\operatorname*{Tr}}\right \vert _{1} \left \{\mathbf{\alpha }_{2}^{t} \mathbf{\alpha }_{2}\right \} I_{\mathcal{F}} +\alpha _{2} \wedge \alpha _{2}\text{.}\end{align}

\bigskip \begin{equation}A_{2 2} A_{2 4} =A_{2 4} A_{2 2} =\sum _{\begin{array}{c}1 \leq j_{1} <j_{2} \leq 2 r \\
1 \leq j_{3} <j_{4} <j_{5} <j_{6} \leq 2 r\end{array}}\alpha _{2 j_{1} j_{2}} \alpha _{4 j_{5} j_{6} j_{7} j_{8}} x_{j_{1}} x_{j_{2}} x_{j_{3}} x_{j_{4}} x_{j_{5}} x_{j_{6}}
\end{equation}
\begin{align}121234 &  : -\frac{1}{4} \alpha _{2 1 2} \alpha _{4 1 2 3 4} x_{3} x_{4} \\
121235 &  : -\frac{1}{4} \alpha _{2 1 2} \alpha _{4 1 2 3 5} x_{3} x_{5} \\
121236 &  : -\frac{1}{4} \alpha _{2 1 2} \alpha _{4 1 2 3 6} x_{3} x_{6} \\
121245 &  : -\frac{1}{4} \alpha _{2 1 2} \alpha _{4 1 2 3 4} x_{4} x_{5} \\
121246 &  : -\frac{1}{4} \alpha _{2 1 2} \alpha _{4 1 2 4 6} x_{4} x_{6} \\
121256 &  : -\frac{1}{4} \alpha _{2 1 2} \alpha _{4 1 2 5 6} x_{5} x_{6} ;131356 : -\frac{1}{4} \alpha _{2 1 3} \alpha _{4 1 3 5 6} x_{5} x_{6} ;141456 : -\frac{1}{4} \alpha _{2 1 4} \alpha _{4 1 4 5 6} x_{5} x_{6} ;232356 : -\frac{1}{4} \alpha _{2 2 3} \alpha _{4 2 3 5 6} x_{5} x_{6} ; \\
122456 &  : -\frac{1}{4} \alpha _{2 2 4} \alpha _{4 2 4 5 6} x_{5} x_{6} ;343456 : -\frac{1}{4} \alpha _{2 3 4} \alpha _{4 3 4 5 6} x_{5} x_{6}\end{align}


\begin{align} & 121345 : -\frac{1}{2} \alpha _{2 1 2} \alpha _{4 1 3 4 5} x_{2} x_{3} x_{4} x_{5} \\
 & 121346 : \\
 & 121356 : \\
 & 121456 : \\
 & 122345 : \\
 & 122346 : \\
 & 122356 : \\
 & 122456 : \\
 & 123456 :\end{align}
\begin{equation}A_{2 2} A_{2 4} =A_{2 4} A_{2 2} =\left (\alpha _{2 1 2} \alpha _{4 3 4 5 6} -\alpha _{2 1 3} \alpha _{4 2 4 5 6} +\alpha _{2 1 4} \alpha _{4 3 2 5 6}\right ) x_{1} x_{2} x_{3} x_{4} x_{5} x_{6}
\end{equation}

\bigskip \begin{equation}A_{2 4} A_{2 4} =\sum _{\begin{array}{c}1 \leq j_{1} <j_{2} <j_{3} <j_{4} \leq 2 r \\
1 \leq j_{5} <j_{6} <j_{7} <j_{8} \leq 2 r\end{array}}\alpha _{4 j_{1} j_{2} j_{3} j_{4}} \alpha _{4 j_{5} j_{6} j_{7} j_{8}} x_{j_{1}} x_{j_{2}} x_{j_{3}} x_{j_{4}} x_{j_{5}} x_{j_{6}} x_{j_{7}} x_{j_{8}}
\end{equation}
\begin{equation}A_{2 4} A_{2 4} =\alpha _{4} \wedge \alpha _{4} +\frac{1}{2} L_{3}^{4} \left (\alpha _{4} \wedge \alpha _{4}\right ) -\frac{1}{4} L_{2}^{4} \left (\alpha _{4} \wedge \alpha _{4}\right ) +\frac{1}{8} L_{1}^{4} \left (\alpha _{4} \wedge \alpha _{4}\right ) -\frac{1}{16} L_{0}^{4} \left (\alpha _{4} \wedge \alpha _{4}\right )I_{\mathcal{F}}\text{.}
\end{equation}

\section{Grassmann Powers}
Let $A_{p}$ and $B_{q}$ be homogeneous of degrees $p$ and $q$ respectively then one has the identies and definitions
\begin{align}A_{p} \wedge B_{q} &  = \langle A_{p} B_{q} \rangle _{p +q} \\
A_{p} \wedge B_{q} &  =\left ( -1\right )^{p q} B_{q} \wedge A_{p}\text{,}\end{align}$\lceil $
\begin{align}I \wedge I &  = \langle I^{2} \rangle _{0} =I \\
\alpha  I \wedge \beta  I &  = \langle \alpha  \beta  I^{2} \rangle _{0} =\alpha  \beta  I\end{align}
\begin{equation}I \wedge B =B \wedge I\text{\quad } \forall B \in \mathcal{F}\text{,}
\end{equation}$\rfloor $ 

where
\begin{equation}A_{p} = \langle A \rangle _{p} ,B_{q} = \langle B \rangle _{q}\text{.}
\end{equation}This gives
\begin{equation}A_{p} \wedge A_{p} =\left ( -1\right )^{p^{2}} A_{p} \wedge A_{p}
\end{equation}thus if $p^{2}$ is odd then
\begin{equation}A_{p} \wedge A_{p} =0\text{,}
\end{equation}i.e.
\begin{equation}A_{p} \wedge A_{p} =0 ;p =1 ,3 ,5 ,\cdots 
\end{equation}


\begin{enumerate}
\item Let $p =2$ then
\begin{align}A_{2} \wedge A_{2} &  =\sum _{1 \leq j_{1} <j_{2} <j_{3} <j_{4} \leq 2 r}2 \left (A_{j_{1} j_{2}} A_{j_{3} j_{4}} -A_{j_{1} j_{3}} A_{j_{2} j_{4}} +A_{j_{1} j_{4}} A_{j_{2} j_{3}}\right ) x_{j_{1}} x_{j_{2}} x_{j_{3}} x_{j_{4}} \\
 &  =\sum _{1 \leq j_{1} <j_{2} <j_{3} <j_{4} \leq 2 r}2 \ensuremath{\operatorname*{Pf}} \left (\mathbb{A}_{2} \left (j_{1} j_{2} j_{3} j_{4}\right )\right ) x_{j_{1}} x_{j_{2}} x_{j_{3}} x_{j_{4}}\text{,}\end{align}where $\mathbb{A}_{2} \left (j_{1} j_{2} j_{3} j_{4}\right )$ is the antisymmetric matrix
\begin{equation}\left (\begin{array}{cccc}0 & A_{j_{1} j_{2}} & A_{j_{1} j_{3}} & A_{j_{1} j_{4}} \\
 -A_{j_{1} j_{2}} & 0 & A_{j_{2} j_{3}} & A_{j_{2} j_{4}} \\
 -A_{j_{1} j_{3}} &  -A_{j_{2} j_{3}} & 0 & A_{j_{3} j_{4}} \\
 -A_{j_{1} j_{4}} &  -A_{j_{2} j_{4}} &  -A_{j_{3} j_{4}} & 0\end{array}\right )
\end{equation}and
\begin{equation}\ensuremath{\operatorname*{Pf}} \left (\mathbb{A}_{2} \left (j_{1} j_{2} j_{3} j_{4}\right )\right ) =A_{j_{1} j_{2}} A_{j_{3} j_{4}} -A_{j_{1} j_{3}} A_{j_{2} j_{4}} +A_{j_{1} j_{4}} A_{j_{2} j_{3}}\text{.}
\end{equation}

\item Let $p =4$ then
\begin{align}A_{4} \wedge A_{4} &  =\sum _{1 \leq j_{1} <j_{2} <j_{3} <j_{4} <j_{5} <j_{6} <j_{7} <j_{8} \leq 2 r}2 \left (A_{j_{1} j_{2} j_{3} j_{4}} A_{j_{5} j_{6} j_{7} j_{8}} +\cdots \right ) x_{j_{1}} x_{j_{2}} x_{j_{3}} x_{j_{4}} x_{j_{5}} x_{j_{6}} x_{j_{7}} x_{j_{8}} \\
 &  =\sum _{1 \leq j_{1} <j_{2} <j_{3} <j_{4} <j_{5} <j_{6} <j_{7} <j_{8} \leq 2 r}2 \ensuremath{\operatorname*{Pf}} \left (\mathbb{A}_{4} \otimes \mathbb{A}_{4} \left (j_{1} j_{2} j_{3} j_{4} j_{5} j_{6} j_{7} j_{8}\right )\right ) x_{j_{1}} x_{j_{2}} x_{j_{3}} x_{j_{4}} x_{j_{5}} x_{j_{6}} x_{j_{7}} x_{j_{8}}\text{.}\end{align}$\lceil $
\begin{equation}\ensuremath{\operatorname*{Pf}} \left (a_{1} ,\cdots  ,a_{2 n}\right ) =\sum _{j =2}^{2 n}\ensuremath{\operatorname*{Pf}} \left (a_{1} ,a_{j}\right ) \left ( -1\right )^{j} \ensuremath{\operatorname*{Pf}} \left (a_{2} ,\cdots  ,\hat{a}_{j} ,\cdots  ,a_{2 n}\right )
\end{equation}$\rfloor $ 

One can consider the special case $r =4$ the subscripts that have no values in common are
\begin{equation}\binom{8}{4} \binom{8 -4}{4} =\binom{8}{4} =\frac{8 \ast 7 \ast 6 \ast 5}{1 \ast 2 \ast 3 \ast 4} =70
\end{equation}e.g.
\begin{equation}12345678,12354678,12364578,12374568,12384567\text{,}
\end{equation}which gives the terms
\begin{align} & A_{1 2 3 4} A_{5 6 7 8} x_{1} \wedge x_{2} \wedge x_{3} \wedge x_{4} \wedge x_{5} \wedge x_{6} \wedge x_{7} \wedge x_{8} \nonumber  \\
 & A_{1 2 3 5} A_{4 6 7 8} x_{1} \wedge x_{2} \wedge x_{3} \wedge x_{5} \wedge x_{4} \wedge x_{6} \wedge x_{7} \wedge x_{8} \nonumber  \\
 & A_{1 2 3 6} A_{4 5 7 8} x_{1} \wedge x_{2} \wedge x_{3} \wedge x_{6} \wedge x_{4} \wedge x_{5} \wedge x_{7} \wedge x_{8} \nonumber  \\
 & A_{1 2 3 7} A_{4 5 6 8} x_{1} \wedge x_{2} \wedge x_{3} \wedge x_{7} \wedge x_{4} \wedge x_{5} \wedge x_{6} \wedge x_{8} \nonumber  \\
 & A_{1 2 3 8} A_{4 5 6 7} x_{1} \wedge x_{2} \wedge x_{3} \wedge x_{8} \wedge x_{4} \wedge x_{5} \wedge x_{6} \wedge x_{7}\end{align}and the final coefficient is
\begin{equation}\left (A_{1 2 3 4} A_{5 6 7 8} -A_{1 2 3 5} A_{4 6 7 8} +A_{1 2 3 6} A_{4 5 7 8} -A_{1 2 3 7} A_{4 5 6 8} +A_{1 2 3 8} A_{4 5 6 7}\right ) x_{1} \wedge x_{2} \wedge x_{3} \wedge x_{4} \wedge x_{5} \wedge x_{6} \wedge x_{7} \wedge x_{8}
\end{equation}
\begin{align} & A_{1 2 3 4} A_{5 6 7 8} x_{1} \wedge x_{2} \wedge x_{3} \wedge x_{4} \wedge x_{5} \wedge x_{6} \wedge x_{7} \wedge x_{8} \nonumber  \\
 & A_{1 2 5 4} A_{3 6 7 8} x_{1} \wedge x_{2} \wedge x_{3} \wedge x_{4} \wedge x_{5} \wedge x_{6} \wedge x_{7} \wedge x_{8} \nonumber  \\
 & A_{1 2 6 4} A_{5 3 7 8} x_{1} \wedge x_{2} \wedge x_{3} \wedge x_{4} \wedge x_{5} \wedge x_{6} \wedge x_{7} \wedge x_{8} \nonumber  \\
 & A_{1 2 7 4} A_{5 6 3 8} x_{1} \wedge x_{2} \wedge x_{3} \wedge x_{4} \wedge x_{5} \wedge x_{6} \wedge x_{7} \wedge x_{8} \nonumber  \\
 & A_{1 2 8 4} A_{5 6 7 3} x_{1} \wedge x_{2} \wedge x_{3} \wedge x_{4} \wedge x_{5} \wedge x_{6} \wedge x_{7} \wedge x_{8}\end{align}
\begin{align}\ensuremath{\operatorname*{Pf}} \left (a_{1} ,\cdots  ,a_{4}\right ) &  =\sum _{j =2}^{4}\ensuremath{\operatorname*{Pf}} \left (a_{1} ,a_{j}\right ) \left ( -1\right )^{j} \ensuremath{\operatorname*{Pf}} \left (a_{2} ,\cdots  ,\hat{a}_{j} ,\cdots  ,a_{4}\right ) \\
 &  =\ensuremath{\operatorname*{Pf}} \left (a_{1} ,a_{2}\right ) \left ( -1\right )^{2} \ensuremath{\operatorname*{Pf}} \left (a_{3} ,a_{4}\right ) +\ensuremath{\operatorname*{Pf}} \left (a_{1} ,a_{3}\right ) \left ( -1\right )^{3} \ensuremath{\operatorname*{Pf}} \left (a_{2} ,a_{4}\right ) \\
 &  +\ensuremath{\operatorname*{Pf}} \left (a_{1} ,a_{4}\right ) \left ( -1\right )^{4} \ensuremath{\operatorname*{Pf}} \left (a_{2} ,a_{3}\right ) \\
 &  =\frac{1}{4} \left (a_{1} a_{2} -a_{2} a_{1}\right ) \left (a_{3} a_{4} -a_{4} a_{3}\right ) -\frac{1}{4} \left (a_{1} a_{3} -a_{3} a_{1}\right ) \left (a_{2} a_{4} -a_{4} a_{2}\right ) \\
 &  +\frac{1}{4} \left (a_{1} a_{4} -a_{4} a_{1}\right ) \left (a_{2} a_{3} -a_{3} a_{2}\right )\end{align}\end{enumerate}


\section{Exterior Product Between General Elements}
\begin{equation}A \wedge B =\sum _{\begin{array}{c}0 \leq m ,n \leq 2 r \\
0 \leq m +n \leq 2 r\end{array}} \langle A_{m} B_{n} \rangle _{m +n}
\end{equation}
\begin{gather}\left (x_{j_{1}} \cdots  x_{j_{m}}\right ) \wedge \left (x_{k_{1}} \cdots  x_{k_{m}}\right ) =\left (x_{j_{1}} \wedge \cdots  \wedge x_{j_{m}}\right ) \wedge \left (x_{k_{1}} \wedge \cdots  \wedge x_{k_{m}}\right ) = \langle \left (x_{j_{1}} \cdots  x_{j_{m}}\right ) \left (x_{k_{1}} \cdots  x_{k_{m}}\right ) \rangle _{m +n} \\
1 \leq j_{1} <\cdots  <j_{m} \leq 2 r ;1 \leq k_{1} \cdots  <k_{m} \leq 2 r\end{gather}

\section{Geometric Product Between General Elements}


\subparagraph{\begin{equation}A B =\sum _{0 \leq m ,n \leq 2 r}A_{m} B_{n}
\end{equation}\begin{equation}A_{m} B_{n} =\sum _{j =0}^{\min \left (m ,n\right )} \langle A_{m} B_{n} \rangle _{\left \vert m -n\right \vert  +2 j}\text{.}
\end{equation}In particular for squares}
\begin{align}A^{2} &  =\sum _{0 \leq m ,n \leq 2 r}A_{m} A_{n} \\
A_{m} A_{n} &  =\sum _{j =0}^{\min \left (m ,n\right )} \langle A_{m} A_{n} \rangle _{\left \vert m -n\right \vert  +2 j}\text{.}\end{align}

Thus for fourth order operators
\begin{align}A^{2} &  =A_{0}^{2} +A_{1}^{2} +A_{2}^{2} +A_{3}^{2} +A_{4}^{2} \nonumber  \\
 &  +A_{0} A_{1} +A_{0} A_{2} +A_{0} A_{3} +A_{0} A_{4} +A_{1} A_{0} +A_{2} A_{0} +A_{3} A_{0} +A_{4} A_{0} \nonumber  \\
 &  +A_{1} A_{2} +A_{1} A_{3} +A_{1} A_{4} +A_{2} A_{1} +A_{3} A_{1} +A_{4} A_{1} \nonumber  \\
 &  +A_{2} A_{3} +A_{2} A_{4} +A_{3} A_{2} +A_{4} A_{2} \nonumber  \\
 &  +A_{3} A_{4} +A_{4} A_{3}\text{,}\end{align}where 

1.
\begin{subequations}
\begin{align}A_{4}^{2} &  = \langle A_{4}^{2} \rangle _{0} + \langle A_{4}^{2} \rangle _{2} + \langle A_{4}^{2} \rangle _{4} + \langle A_{4}^{2} \rangle _{6} + \langle A_{4}^{2} \rangle _{8} \\
A_{3}^{2} &  = \langle A_{3}^{2} \rangle _{0} + \langle A_{3}^{2} \rangle _{2} + \langle A_{3}^{2} \rangle _{4} + \langle A_{3}^{2} \rangle _{6} \\
A_{2}^{2} &  = \langle A_{2}^{2} \rangle _{0} + \langle A_{2}^{2} \rangle _{2} + \langle A_{2}^{2} \rangle _{4} \\
A_{1}^{2} &  = \langle A_{1}^{2} \rangle _{0} + \langle A_{1}^{2} \rangle _{2} \\
A_{0}^{2} &  = \langle A_{0}^{2} \rangle _{0}\text{,}\end{align}thus
\end{subequations}
\begin{align*}\left (\left .A\right \vert A\right ) &  = \langle A_{4}^{2} \rangle _{0} + \langle A_{3}^{2} \rangle _{0} + \langle A_{2}^{2} \rangle _{0} + \langle A_{1}^{2} \rangle _{0} + \langle A_{0}^{2} \rangle _{0} \\
 &  =\left (\left .A_{0}\right \vert A_{0}\right ) +\left (\left .A_{1}\right \vert A_{1}\right ) +\left (\left .A_{2}\right \vert A_{2}\right ) +\left (\left .A_{3}\right \vert A_{3}\right ) +\left (\left .A_{4}\right \vert A_{4}\right ) \\
 &  =\left \{\alpha _{0}^{2} +\frac{1}{2} \sum _{1 \leq j_{1} \leq 2 r}\left (\alpha _{1 j_{1}}\right )^{2} -\frac{1}{4} \sum _{1 \leq j_{1} <j_{2} \leq 2 r}\left (\alpha _{2 j_{1} j_{2}}\right )^{2} -\frac{1}{8} \sum _{1 \leq j_{1} <j_{2} <j_{3} \leq 2 r}\left (\alpha _{3 j_{1} j_{2} j_{3}}\right )^{2} +\frac{1}{16} \sum _{1 \leq j_{1} <j_{2} <j_{3} <j_{4} \leq 2 r}\left (\alpha _{4 j_{1} j_{2} j_{3} j_{4}}\right )^{2}\right \} I_{\mathcal{F}}\end{align*}One can also express the inner products of the even number components as
\begin{align}\left (\left .A_{0}\right \vert A_{0}\right ) &  =\left .\ensuremath{\operatorname*{Tr}}\right \vert _{0} \left \{\mathbb{A}_{0}^{t} \mathbb{A}_{0}\right \} \ensuremath{\operatorname*{Tr}} I_{\mathcal{F}} \nonumber  \\
\left (\left .A_{2}\right \vert A_{2}\right ) &  =\frac{1}{4} \left .\ensuremath{\operatorname*{Tr}}\right \vert _{1} \left \{\mathbb{A}_{2}^{t} \mathbb{A}_{2}\right \} I_{\mathcal{F}} \nonumber  \\
\left (\left .A_{4}\right \vert A_{4}\right ) &  =\frac{1}{16} \left .\ensuremath{\operatorname*{Tr}}\right \vert _{2} \left \{\mathbb{A}_{4}^{t} \mathbb{A}_{4}\right \} I_{\mathcal{F}}\text{,}\end{align}where we have used the ordered form of the basis. 

2.
\begin{align}A_{3} A_{4} &  = \langle A_{3} A_{4} \rangle _{1} + \langle A_{3} A_{4} \rangle _{3} + \langle A_{3} A_{4} \rangle _{5} + \langle A_{3} A_{4} \rangle _{7} \\
A_{4} A_{3} &  = \langle A_{4} A_{3} \rangle _{1} + \langle A_{4} A_{3} \rangle _{3} + \langle A_{4} A_{3} \rangle _{5} + \langle A_{4} A_{3} \rangle _{7}\text{,}\end{align}

3.
\begin{align}A_{2} A_{3} &  = \langle A_{2} A_{3} \rangle _{1} + \langle A_{2} A_{3} \rangle _{3} + \langle A_{2} A_{3} \rangle _{5} \\
A_{3} A_{2} &  = \langle A_{3} A_{2} \rangle _{1} + \langle A_{3} A_{2} \rangle _{3} + \langle A_{3} A_{2} \rangle _{5}\text{,}\end{align}

4.
\begin{align}A_{2} A_{4} &  = \langle A_{2} A_{4} \rangle _{2} + \langle A_{2} A_{4} \rangle _{4} + \langle A_{2} A_{4} \rangle _{6} \\
A_{4} A_{2} &  = \langle A_{4} A_{2} \rangle _{2} + \langle A_{4} A_{2} \rangle _{4} + \langle A_{4} A_{2} \rangle _{6}\text{.}\end{align}

5.
\begin{align}A_{1} A_{2} &  = \langle A_{1} A_{2} \rangle _{1} + \langle A_{1} A_{2} \rangle _{3} \\
A_{2} A_{1} &  = \langle A_{2} A_{1} \rangle _{1} + \langle A_{2} A_{1} \rangle _{3}\end{align}

6.
\begin{align}A_{1} A_{3} &  = \langle A_{1} A_{3} \rangle _{2} + \langle A_{1} A_{3} \rangle _{4} \\
A_{3} A_{1} &  = \langle A_{3} A_{1} \rangle _{2} + \langle A_{3} A_{1} \rangle _{4}\end{align}

7.
\begin{align}A_{1} A_{4} &  = \langle A_{1} A_{4} \rangle _{3} + \langle A_{1} A_{4} \rangle _{5} \\
A_{4} A_{1} &  = \langle A_{4} A_{1} \rangle _{3} + \langle A_{4} A_{1} \rangle _{5}\text{.}\end{align}In our case
\begin{equation}A =A_{0} +A_{2} +A_{4}\text{,}
\end{equation}hence
\begin{align}A^{2} &  = \langle A_{4}^{2} \rangle _{0} + \langle A_{4}^{2} \rangle _{2} + \langle A_{4}^{2} \rangle _{4} + \langle A_{4}^{2} \rangle _{6} + \langle A_{4}^{2} \rangle _{8} + \langle A_{2}^{2} \rangle _{0} + \langle A_{2}^{2} \rangle _{2} + \langle A_{2}^{2} \rangle _{4} \nonumber  \\
 &  + \langle A_{2} A_{4} \rangle _{2} + \langle A_{2} A_{4} \rangle _{4} + \langle A_{2} A_{4} \rangle _{6} + \langle A_{4} A_{2} \rangle _{2} + \langle A_{4} A_{2} \rangle _{4} + \langle A_{4} A_{2} \rangle _{6} +2 A_{0} A_{2} +2 A_{0} A_{4}\text{,}\end{align}which can be rearranged as
\begin{align}A^{2} &  = \langle A_{4}^{2} \rangle _{0} + \langle A_{2}^{2} \rangle _{0} +A_{0}^{2} \nonumber  \\
 &  + \langle A_{4}^{2} \rangle _{2} + \langle A_{2}^{2} \rangle _{2} + \langle A_{2} A_{4} \rangle _{2} + \langle A_{4} A_{2} \rangle _{2} +2 A_{0} A_{2} \nonumber  \\
 &  +2 A_{0} A_{4} + \langle A_{4}^{2} \rangle _{4} + \langle A_{2} A_{4} \rangle _{4} + \langle A_{4} A_{2} \rangle _{4} + \langle A_{2}^{2} \rangle _{4} \nonumber  \\
 &  + \langle A_{4}^{2} \rangle _{6} + \langle A_{2} A_{4} \rangle _{6} + \langle A_{4} A_{2} \rangle _{6} \nonumber  \\
 &  + \langle A_{4}^{2} \rangle _{8}\text{,}\end{align}As
\begin{equation} \langle A_{r} B_{s} \rangle _{r +s -2 j} =\left ( -1\right )^{r s -j}  \langle B_{s} A_{r} \rangle _{r +s -2 j}
\end{equation}this gives
\begin{align}A^{2} &  =A_{0}^{2} + \langle A_{4}^{2} \rangle _{0} + \langle A_{2}^{2} \rangle _{0} \nonumber  \\
 &  +2  \langle A_{2} A_{4} \rangle _{2} +2 A_{0} A_{2} \nonumber  \\
 &  + \langle A_{4}^{2} \rangle _{4} + \langle A_{2}^{2} \rangle _{4} +2 A_{0} A_{4} \nonumber  \\
 &  +2  \langle A_{2} A_{4} \rangle _{6} \nonumber  \\
 &  + \langle A_{4}^{2} \rangle _{8}\text{,}\end{align}as
\label{contraccond}
\begin{subequations}
\begin{align} \langle A_{4} A_{4} \rangle _{8 -6} &  =\left ( -1\right )^{16 -3}  \langle A_{4} A_{4} \rangle _{8 -6} \Rightarrow  \langle A_{4}^{2} \rangle _{2} =0 \label{4-2-2} \\
 \langle A_{4} A_{4} \rangle _{8 -4} &  =\left ( -1\right )^{16 -2}  \langle A_{4} A_{4} \rangle _{8 -4} = \langle A_{4} A_{4} \rangle _{4} \label{4-4-4} \\
 \langle A_{4} A_{4} \rangle _{8 -2} &  =\left ( -1\right )^{16 -1}  \langle A_{4} A_{4} \rangle _{8 -2} \Rightarrow  \langle A_{4}^{2} \rangle _{6} =0 \label{4-4-6} \\
 \langle A_{4} A_{4} \rangle _{8 -0} &  =\left ( -1\right )^{16 -0}  \langle A_{4} A_{4} \rangle _{8} = \langle A_{4}^{2} \rangle _{8} \label{4-4-8} \\
 \langle A_{4} A_{2} \rangle _{6 -4} &  =\left ( -1\right )^{8 -2}  \langle A_{2} A_{4} \rangle _{6 -4} = \langle A_{2} A_{4} \rangle _{2} \label{4-2--2} \\
 \langle A_{2} A_{4} \rangle _{6 -2} &  =\left ( -1\right )^{8 -1}  \langle A_{4} A_{2} \rangle _{6 -2} \Rightarrow  \langle A_{2} A_{4} \rangle _{4} = - \langle A_{4} A_{2} \rangle _{4} \label{2-4-4} \\
 \langle A_{4} A_{2} \rangle _{6 -0} &  =\left ( -1\right )^{8 -0}  \langle A_{2} A_{4} \rangle _{6 -0} = \langle A_{2} A_{4} \rangle _{6} \label{4-2-6} \\
 \langle A_{2} A_{2} \rangle _{4 -2} &  =\left ( -1\right )^{4 -1}  \langle A_{2} A_{2} \rangle _{4 -2} \Rightarrow  \langle A_{2}^{2} \rangle _{2} =0 \label{2-2-2} \\
 \langle A_{2} A_{2} \rangle _{4 -0} &  =\left ( -1\right )^{4 -0}  \langle A_{2} A_{2} \rangle _{4 -0} = \langle A_{2}^{2} \rangle _{4} . \label{2-2-0}\end{align}It is easy to show that
\end{subequations}
\begin{subequations}
\begin{align} \langle A_{4}^{2} \rangle _{0} &  =\left (\left .A_{4}\right \vert A_{4}\right ) I_{\mathcal{F}} \\
 \langle A_{2}^{2} \rangle _{0} &  =\frac{1}{8} \left .\ensuremath{\operatorname*{Tr}}\right \vert _{1} \left \{\mathbb{A}_{2}^{t} \mathbb{A}_{2}\right \} I_{\mathcal{F}} =\left (\left .A_{2}\right \vert A_{2}\right ) I_{\mathcal{F}}\end{align}and
\end{subequations}
\begin{subequations}
\begin{align} \langle A_{4}^{2} \rangle _{8} &  =A_{4} \wedge A_{4} \\
 \langle A_{2}^{2} \rangle _{4} &  =A_{2} \wedge A_{2} \\
 \langle A_{2} A_{4} \rangle _{6} &  =A_{2} \wedge A_{4}\text{.}\end{align}One can check this by considering representative tensor elements for eq.$\left (\ref{4-2-2}\right )$
\end{subequations}
\begin{equation} \langle A_{4} A_{2} \rangle _{\left (2\right ) 12} = -\left (A_{1 2 5 6} A_{5 6} +A_{1 2 4 6} A_{4 6} +A_{1 2 3 6} A_{3 6} +A_{1 2 4 5} A_{4 5} +A_{1 2 3 5} A_{3 5} +A_{1 2 3 4} A_{3 4}\right ) \neq 0
\end{equation}
\begin{equation} \langle A_{2} A_{4} \rangle _{\left (2\right ) 12} = -\left (A_{5 6} A_{1 2 5 6} +A_{4 6} A_{1 2 4 6} +A_{3 6} A_{1 2 3 6} +A_{4 5} A_{1 2 4 5} +A_{3 5} A_{1 2 3 5} +A_{3 4} A_{1 2 3 4}\right ) \neq 0
\end{equation}for eq.$\left (\ref{4-4-6}\right )$
\begin{align} \langle A_{4}^{2} \rangle _{\left (6\right ) 234567} = & A_{1 2 3 4} A_{1 5 6 7} -A_{1 2 3 5} A_{1 4 6 7} +A_{1 2 3 6} A_{1 4 5 7} -A_{1 2 3 7} A_{1 4 5 6} +A_{1 2 4 5} A_{1 3 6 7} -A_{1 2 4 6} A_{1 3 5 7} +A_{1 2 4 7} A_{1 3 5 6} -A_{1 2 5 6} A_{1 3 4 7} \nonumber  \\
 &  -A_{1 2 5 7} A_{1 3 4 6} +A_{1 2 6 7} A_{1 3 4 5} -A_{1 3 4 5} A_{1 2 6 7} +A_{1 3 4 6} A_{1 2 5 7} -A_{1 3 4 7} A_{1 2 5 6} -A_{1 3 5 6} A_{1 2 4 7} +A_{1 3 5 7} A_{1 2 4 6} -A_{1 3 6 7} A_{1 2 4 5} \nonumber  \\
 &  +A_{1 4 5 6} A_{1 2 3 7} -A_{1 4 5 7} A_{1 2 3 6} +A_{1 4 6 7} A_{1 2 3 5} , -A_{1 5 6 7} A_{1 2 3 4} =0\end{align}for eq. $\left (\ref{4-4-4}\right )$
\begin{align} \langle A_{4}^{2} \rangle _{\left (4\right ) 3456} &  = -A_{1 2 3 4} A_{1 2 5 6} +A_{1 2 3 5} A_{1 2 4 6} -A_{1 2 3 6} A_{1 2 4 5} -A_{1 2 4 5} A_{1 2 3 6} +A_{1 2 4 6} A_{1 2 3 5} -A_{1 2 5 6} A_{1 2 3 4} \\
 &  =2 \left (A_{1 2 3 5} A_{1 2 4 6} -A_{1 2 3 4} A_{1 2 5 6} -A_{1 2 3 6} A_{1 2 4 5}\right )\end{align}for eq. $\left (\ref{4-4-2}\right )$
\begin{equation} \langle A_{4}^{2} \rangle _{\left (2\right ) 45} =A_{1 2 3 4} A_{1 2 3 5} -A_{1 2 3 5} A_{1 2 3 4} +A_{1 2 3 4} A_{1 2 3 6} -A_{1 2 3 6} A_{1 2 3 4} +A_{1 2 3 4} A_{1 2 3 5} -A_{1 2 3 5} A_{1 2 3 4} \cdots  =0.
\end{equation}

The \emph{reflection or involution }condition $A^{2} =I_{\mathcal{F}}$ gives
\begin{equation}A^{2} =A_{0}^{2} + \langle A_{4}^{2} \rangle _{0} + \langle A_{2}^{2} \rangle _{0} +2  \langle A_{2} A_{4} \rangle _{2} +2 A_{0} A_{2} + \langle A_{4}^{2} \rangle _{4} + \langle A_{2}^{2} \rangle _{4} +2 A_{0} A_{4} +2  \langle A_{2} A_{4} \rangle _{6} + \langle A_{4}^{2} \rangle _{8} =I_{\mathcal{F}}\text{,}
\end{equation}leading to
\begin{align}A_{0}^{2} + \langle A_{4}^{2} \rangle _{0} + \langle A_{2}^{2} \rangle _{0} &  =I_{\mathcal{F}} \\
 \langle A_{2} A_{4} \rangle _{2} +A_{0} A_{2} &  =0 \\
 \langle A_{4}^{2} \rangle _{4} + \langle A_{2}^{2} \rangle _{4} +2 A_{0} A_{4} &  =0 \\
 \langle A_{2} A_{4} \rangle _{6} &  =A_{2} \wedge A_{4} =0 \\
 \langle A_{4}^{2} \rangle _{8} &  =A_{4} \wedge A_{4} =0\end{align}

If $A$ is a reflection then
\begin{align*}\mathcal{A}_{ +} &  =\frac{1}{\sqrt{2}} \left (I_{\mathcal{F}} +A\right ) \\
\mathcal{A}_{ -} &  =\frac{1}{\sqrt{2}} \left (I_{\mathcal{F}} -A\right )\end{align*}are \emph{orthogonal idempotents}, as
\begin{equation*}\mathcal{A}_{ +}^{2} =\left (\frac{1}{\sqrt{2}} \left (I_{\mathcal{F}} +A\right )\right )^{2} =\frac{1}{2} \left (I_{\mathcal{F}} +2 A +I_{\mathcal{F}}\right ) =\mathcal{A}_{ +}\text{,}
\end{equation*}
\begin{equation*}\mathcal{A}_{ -}^{2} =\left (\frac{1}{\sqrt{2}} \left (I_{\mathcal{F}} -A\right )\right )^{2} =\frac{1}{2} \left (I_{\mathcal{F}} -2 A +I_{\mathcal{F}}\right ) =\mathcal{A}_{ -}
\end{equation*}and
\begin{equation*}\mathcal{A}_{ +} \mathcal{A}_{ -} =\mathcal{A}_{ -} \mathcal{A}_{ +} =0 =\frac{1}{\sqrt{2}} \left (I_{\mathcal{F}} +A\right ) \frac{1}{\sqrt{2}} \left (I_{\mathcal{F}} -A\right ) =\frac{1}{2} \left (I_{\mathcal{F}} +A\right ) \left (I_{\mathcal{F}} -A\right )\text{.}
\end{equation*}

The \emph{idempotency} condition $A^{2} =A$ gives
\begin{align}A^{2} -A &  =\left \{A_{0} \left (A_{0} -1\right ) + \langle A_{4}^{2} \rangle _{0} + \langle A_{2}^{2} \rangle _{0}\right \} +\left \{2  \langle A_{2} A_{4} \rangle _{2} +\left (2 A_{0} -1\right ) A_{2}\right \} +\left \{ \langle A_{4}^{2} \rangle _{4} + \langle A_{2}^{2} \rangle _{4} +\left (2 A_{0} -1\right ) A_{4}\right \} \nonumber  \\
 &  +2  \langle A_{2} A_{4} \rangle _{6} + \langle A_{4}^{2} \rangle _{8} =0\text{,}\end{align}which gives
\label{idem cond}
\begin{subequations}
\begin{align}A_{0} \left (A_{0} -I_{\mathcal{F}}\right ) + \langle A_{4}^{2} \rangle _{0} + \langle A_{2}^{2} \rangle _{0} &  =0 \\
2  \langle A_{2} A_{4} \rangle _{2} +\left (2 A_{0} -I_{\mathcal{F}}\right ) A_{2} &  =0 \\
 \langle A_{4}^{2} \rangle _{4} + \langle A_{2}^{2} \rangle _{4} +\left (2 A_{0} -I_{\mathcal{F}}\right ) A_{4} &  =0 \\
 \langle A_{2} A_{4} \rangle _{6} &  =0 \\
 \langle A_{4}^{2} \rangle _{8} &  =0\end{align}and thus
\end{subequations}
\begin{align}A_{0} \left (A_{0} -1\right ) +\left (\left .A_{4}\right \vert A_{4}\right ) +\left (\left .A_{2}\right \vert A_{2}\right ) &  =0 \\
2  \langle A_{2} A_{4} \rangle _{2} +\left (2 A_{0} -1\right ) A_{2} &  =0 \\
 \langle A_{4}^{2} \rangle _{4} +A_{2} \wedge A_{2} +\left (2 A_{0} -1\right ) A_{4} &  =0 \\
A_{2} \wedge A_{4} &  =0 \\
A_{4} \wedge A_{4} &  =0\end{align}and in more detail
\begin{align}A_{0} \left (A_{0} -1\right ) +\sum _{1 \leq j_{1} <j_{2} <j_{3} <j_{4} \leq 2 r}A_{4 j_{1} j_{2} j_{3} j_{4}}^{2} +\sum _{1 \leq j_{1} <j_{2} \leq 2 r}A_{2 j_{1} j_{2}}^{2} &  =0 \\
2  \langle A_{2} A_{4} \rangle _{2} +\left (2 \alpha  -1\right ) A_{2} &  =0 \\
 \langle A_{4}^{2} \rangle _{4} +A_{2} \wedge A_{2} +\left (2 \alpha  -1\right ) A_{4} &  =0 \\
A_{2} \wedge A_{4} &  =0 \\
A_{4} \wedge A_{4} &  =0.\end{align}

\subsection{Canonical form of 4-forms}
(check what happens if $A_{4}$ are self adjoint) 

The homogeneous operator $A_{4}$ can be expanded as
\begin{align}A_{4} &  =\frac{1}{4 !} \sum _{\begin{array}{c}1 \leq j_{1} ,j_{2} ,j_{3} ,j_{4} \leq 2 r\end{array}}A_{j_{1} j_{2} j_{3} j_{4}} x_{j_{1}} x_{j_{2}} x_{j_{3}} x_{j_{4}} =\frac{1}{4 !} \sum _{\begin{array}{c}1 \leq j_{1} ,j_{2} ,j_{3} ,j_{4} \leq 2 r\end{array}}A_{j_{1} j_{2} j_{3} j_{4}} x_{j_{1}} \wedge x_{j_{2}} \wedge x_{j_{3}} \wedge x_{j_{4}} \\
 &  =\sum _{\begin{array}{c}1 \leq j_{1} <j_{2} <j_{3} <j_{4} \leq 2 r\end{array}}A_{j_{1} j_{2} j_{3} j_{4}} x_{j_{1}} x_{j_{2}} x_{j_{3}} x_{j_{4}} =\frac{1}{4} \sum _{\begin{array}{c}1 \leq j_{1} <j_{2} \leq 2 r \\
1 <j_{3} <j_{4} \leq 2 r\end{array}}A_{j_{1} j_{2} j_{3} j_{4}} x_{j_{1}} x_{j_{2}} x_{j_{3}} x_{j_{4}}\text{,}\end{align}where the coefficients $\left \{A_{j_{1} j_{2} j_{3} j_{4}}\right \}$ are antisymmetric in the subscripts and
\begin{equation*}A_{j_{1} j_{2} j_{3} j_{4}} =0\text{if}j_{1} =j_{3}\text{or}j_{4} ,\,j_{2} =j_{3}\text{or}j_{4} ,\,j_{1} =j_{2}\text{or}\,j_{3} =j_{4} .\text{\quad \quad \quad }
\end{equation*}$\lceil \;$consider the example for $r =2$ arranged as $12 \times 12$ matrix
\begin{equation*}\begin{array}{ccccccccccccc}\, & 12 & 13 & 14 & 21 & 23 & 24 & 31 & 32 & 34 & 41 & 42 & 43 \\
12 & 0 & 0 & 0 & 0 & 0 & 0 & 0 & 0 & A_{1 2 3 4} & 0 & 0 & A_{1 2 4 3} \\
13 & 0 & 0 & 0 & 0 & 0 & A_{1 3 2 4} & 0 & 0 & 0 & 0 &  \ast  & 0 \\
14 & 0 & 0 & 0 & 0 & A_{1 4 2 3} & 0 & 0 & A_{1 4 3 2} & 0 & 0 & 0 & 0 \\
21 & 0 & 0 & 0 & 0 & 0 & 0 & 0 & 0 &  \ast  & 0 & 0 &  \ast  \\
23 & 0 & 0 &  \ast  & 0 & 0 & 0 & 0 & 0 & 0 &  \ast  & 0 & 0 \\
24 & 0 &  \ast  & 0 & 0 & 0 & 0 &  \ast  & 0 & 0 & 0 & 0 & 0 \\
31 & 0 & 0 & 0 & 0 & 0 &  \ast  & 0 & 0 & 0 & 0 &  \ast  & 0 \\
32 & 0 & 0 &  \ast  & 0 & 0 & 0 & 0 & 0 & 0 &  \ast  & 0 & 0 \\
34 &  \ast  & 0 & 0 &  \ast  & 0 & 0 & 0 & 0 & 0 & 0 & 0 & 0 \\
41 & 0 & 0 & 0 & 0 &  \ast  & 0 & 0 &  \ast  & 0 & 0 & 0 & 0 \\
42 & 0 &  \ast  & 0 & 0 & 0 & 0 &  \ast  & 0 & 0 & 0 & 0 & 0 \\
43 &  \ast  & 0 & 0 &  \ast  & 0 & 0 & 0 & 0 & 0 & 0 & 0 & 0\end{array}
\end{equation*}or as a $6 \times 6$ \emph{symmetric} matrix
\begin{equation*}\left (\begin{array}{ccccccc}\, & 12 & 13 & 14 & 23 & 24 & 34 \\
12 & 0 & 0 & 0 & 0 & 0 & A_{1 2 3 4} \\
13 & 0 & 0 & 0 & 0 & A_{1 3 2 4} & 0 \\
14 & 0 & 0 & 0 & A_{1 4 2 3} & 0 & 0 \\
23 & 0 & 0 & A_{2 3 1 4} & 0 & 0 & 0 \\
24 & 0 & A_{2 4 1 3} & 0 & 0 & 0 & 0 \\
34 & A_{3 4 1 2} & 0 & 0 & 0 & 0 & 0\end{array}\right )
\end{equation*}$\rfloor $ 

As
\begin{equation*}x_{j_{1}} x_{j_{2}} x_{j_{3}} x_{j_{4}} =x_{j_{3}} x_{j_{4}} x_{j_{1}} x_{j_{2}}
\end{equation*}one has
\begin{equation*}A_{j_{1} j_{2} j_{3} j_{4}} =A_{j_{3} j_{4} j_{1} j_{2}}
\end{equation*}and thus the matrix $\mathbb{A} =\left \{\left .A_{J_{1 2} J_{3 4}}\right \vert 1 \leq J_{1 2} ,J_{3 4} \leq r \left (2 r -1\right )\right \}$ is symmetric i.e.
\begin{equation*}A_{J_{1 2} J_{3 4}} =A_{J_{3 4} J_{1 2}} ;1 \leq J_{1 2} ,J_{3 4} \leq r \left (2 r -1\right )
\end{equation*}and the operator $A_{4}$ can be expanded as
\begin{equation*}A_{4} =\frac{1}{4} \sum _{1 \leq J_{1 2} ,J_{3 4} \leq r \left (2 r -1\right )}A_{4 J_{1 2} J_{3 4}} Y_{J_{1 2}} Y_{J_{3 4}} =\frac{1}{4} \sum _{1 \leq J_{1 2} ,J_{3 4} \leq r \left (2 r -1\right )}A_{4 J_{1 2} J_{3 4}} Y_{J_{1 2}} \wedge Y_{J_{3 4}}\text{,}
\end{equation*}where
\begin{align*}Y_{J_{1 2}} &  =x_{j_{1}} x_{j_{2}} \\
Y_{J_{3 4}} &  =x_{j_{3}} x_{j_{4}}\text{.}\end{align*}If the matrix $\mathbb{A} =\left \{\left .A_{J_{1 2} J_{3 4}}\right \vert 1 \leq J_{1 2} ,J_{3 4} \leq r \left (2 r -1\right )\right \}$ is real symmetric, (which it must be if $A_{4}$ is self adjoint), one can diagonalize it with real orthogonal transformations i.e. $\mathbb{O} \in O (r \left (2 r -1\right ) ,\mathbb{R})\text{,}$ to obtain
\begin{equation*}\mathbb{O}^{t} \mathbb{A} \mathbb{O} =\left (\begin{array}{ccc}\alpha _{1} & \cdots  & 0 \\
\vdots  & \ddots  & \vdots  \\
0 & \cdots  & \alpha _{r \left (2 r -1\right )}\end{array}\right ) \Longleftrightarrow \mathbb{A} =\mathbb{O} \left (\begin{array}{ccc}\alpha _{1} & \cdots  & 0 \\
\vdots  & \ddots  & \vdots  \\
0 & \cdots  & \alpha _{r \left (2 r -1\right )}\end{array}\right ) \mathbb{O}^{t}
\end{equation*}giving
\begin{equation*}A_{J_{1 2} J_{3 4}} =\sum _{1 \leq K \leq r \left (2 r -1\right )}O_{J_{1 2} K} \alpha _{K} O_{J_{3 4} K}
\end{equation*}and thus
\begin{equation*}A_{4} =\frac{1}{4} \sum _{1 \leq J_{1 2} ,J_{3 4} \leq r \left (2 r -1\right )}A_{4 J_{1 2} J_{3 4}} Y_{J_{1 2}} Y_{J_{3 4}} =\frac{1}{4} \sum _{1 \leq J_{1 2} ,J_{3 4} \leq r \left (2 r -1\right )}\sum _{1 \leq K \leq r \left (2 r -1\right )}O_{J_{1 2} K} \alpha _{K} O_{J_{3 4} K} Y_{J_{1 2}} Y_{J_{3 4}}\text{.}
\end{equation*}One can define the generalized geminal operator
\begin{equation*}\mathcal{O}_{K} =\sum _{1 \leq J_{1 2} \leq r \left (2 r -1\right )}O_{J_{1 2} K} Y_{J_{1 2}} =\sum _{1 \leq j_{1} \neq j_{2} \leq 2 r}O_{k j_{1} j_{2}} x_{j_{1}} x_{j_{2}} ;\text{\quad }O_{k j_{1} j_{2}} = -O_{k j_{2} j_{1}}\text{,}
\end{equation*}which has the canonical form
\begin{equation*}\mathcal{O}_{K} =\sum _{1 \leq l \leq r}c_{K l} y_{K l} y_{K l +r} =\sum _{1 \leq l \leq r}c_{K l} y_{K l} \wedge y_{K l +r} ,\text{\quad }c_{K l} \in \mathbb{R}^{ +}\text{.}
\end{equation*}Note that $\left \{y_{K l}\right \}$ \emph{are} \emph{self adjoint }operators belonging to $\mathcal{F}$ that are linear combinations of single creation and annihilation operators$\text{.}$ The operator $A_{4}$ can thus be expressed as
\begin{align*}A_{4} &  =\frac{1}{4} \sum _{1 \leq J_{1 2} ,J_{3 4} \leq r \left (2 r -1\right )}A_{4 J_{1 2} J_{3 4}} Y_{J_{1 2}} Y_{J_{3 4}} =\frac{1}{4} \sum _{1 \leq K \leq r \left (2 r -1\right )}\alpha _{K} \sum _{1 \leq J_{1 2} ,J_{3 4} \leq r \left (2 r -1\right )}O_{J_{12 r \left (2 r +1\right ) 2} K} O_{J_{3 4} K} Y_{J_{1 2}} Y_{J_{3 4}} \\
 &  =\frac{1}{4} \sum _{1 \leq K \leq r \left (2 r -1\right )}\alpha _{K} \sum _{1 \leq J_{1 2} \leq r \left (2 r -1\right )}O_{J_{1 2} K} Y_{J_{1 2}} \sum _{1 \leq J_{3 4} \leq r \left (2 r -1\right )}O_{J_{3 4} K} Y_{J_{3 4}} =\frac{1}{4} \sum _{1 \leq K \leq r \left (2 r -1\right )}\alpha _{K} \mathcal{O}_{K} \wedge \mathcal{O}_{K} \\
 &  =\frac{1}{4} \sum _{1 \leq K \leq 2 r \left (2 r -1\right )}\alpha _{K} \mathcal{O}_{K}^{2 \wedge }\text{.}\end{align*}Using the canonical form for a second degree antisymmetric polynomial this gives the canonical form for a fourth degree
antisymmetric polynomial as
\begin{align*}A_{4} &  =\frac{1}{4} \sum _{1 \leq K \leq r \left (2 r -1\right )}\alpha _{K} \sum _{1 \leq l_{1} ,l_{2} \leq r}c_{K l_{1}} c_{K l_{2}} y_{K l_{1}} y_{K l_{1} +r} y_{K l_{2}} y_{K l_{2} +r} =\frac{1}{4} \sum _{1 \leq K \leq r \left (2 r -1\right )}\alpha _{K} \sum _{1 \leq l_{1} \neq l_{2} \leq r}c_{K l_{1}} c_{K l_{2}} y_{K l_{1}} y_{K l_{1} +r} y_{K l_{2}} y_{K l_{2} +r} \\
 &  =\frac{1}{2} \sum _{1 \leq K \leq r \left (2 r -1\right )}\alpha _{K} \sum _{1 \leq l_{1} <l_{2} \leq r}c_{K l_{1}} c_{K l_{2}} y_{K l_{1}} y_{K l_{1} +r} y_{K l_{2}} y_{K l_{2} +r} =\sum _{1 \leq K \leq r \left (2 r -1\right )}\alpha _{K} \mathcal{O}_{K} \wedge \mathcal{O}_{K}\text{.}\end{align*}One should note that $\left \{\alpha _{K} \in \mathbb{R}\right \}$ and $\left \{c_{K l} \geq 0\right \}$ and $\left \{\mathcal{O}_{K}\right \}$ are self adjoint. However the coefficients $\alpha _{K}$ could be negative \emph{unless the matrix} $\mathbb{A}$ is positive, which it is if it ia a projector. 

$\lceil $ If
\begin{equation*}\mathcal{O}_{K} \wedge \mathcal{O}_{K} =\mathcal{O}_{K}^{2 \wedge } =0
\end{equation*}then
\begin{equation*}\text{rank}\left \{\mathcal{O}_{K}\right \} =1
\end{equation*}as
\begin{align*}\mathcal{O}_{K}^{2 \wedge } &  =2 \sum _{1 \leq l <m \leq r}c_{K l} c_{K m} y_{K l} y_{K l +r} y_{K m} y_{K m +r} \\
 &  =2 \sum _{1 \leq l <m \leq r}c_{K l} c_{K m} y_{K l} y_{K m} y_{K l +r} y_{K m +r}\end{align*}and
\begin{equation*}\left (\left .y_{K p}\right \vert y_{K q}\right ) =\delta _{p q} I_{\mathcal{F}}
\end{equation*}so that
\begin{equation*}y_{K l} y_{K l +r} y_{K m} y_{K m +r} \neq 0\text{if}l \neq m\text{.}
\end{equation*}$\rfloor $ 

The Grassmann square of $A_{4}$ is
\begin{align*}A_{4}^{2 \wedge } &  =A_{4} \wedge A_{4} =\genfrac{(}{)}{}{}{1}{4}^{2} \sum _{1 \leq K \leq r \left (2 r -1\right )}\alpha _{K} \mathcal{O}_{K} \wedge \mathcal{O}_{K} \wedge \sum _{1 \leq L \leq r \left (2 r -1\right )}\alpha _{L} \mathcal{O}_{L} \wedge \mathcal{O}_{L} \\
 &  =\genfrac{(}{)}{}{}{1}{4}^{2} \sum _{\begin{array}{c}1 \leq K \leq r \left (2 r -1\right ) \\
1 \leq L \leq r \left (2 r -1\right )\end{array}}\alpha _{K} \alpha _{L} \mathcal{O}_{K} \wedge \mathcal{O}_{L} \wedge \mathcal{O}_{K} \wedge \mathcal{O}_{L} \\
 &  =\genfrac{(}{)}{}{}{1}{4}^{2} \sum _{\begin{array}{c}1 \leq K \leq r \left (2 r -1\right ) \\
1 \leq L \leq r \left (2 r -1\right )\end{array}}\alpha _{K} \alpha _{L} \mathcal{O}_{K L} \wedge \mathcal{O}_{K L} \\
 &  =\genfrac{(}{)}{}{}{1}{4}^{2} \sum _{1 \leq K ,L \leq r \left (2 r -1\right )}\alpha _{K} \alpha _{L} \sum _{1 \leq m ,n ,p ,q \leq r}c_{K m} c_{K n} c_{L p} c_{L q} \left (y_{K m} y_{K m +r} y_{K n} y_{K n +r} y_{L p} y_{L p +r} y_{L q} y_{L q +r}\right ) \\
 &  =\genfrac{(}{)}{}{}{1}{4}^{2} \sum _{1 \leq K ,L \leq r \left (2 r -1\right )}\alpha _{K} \alpha _{L} \sum _{1 \leq m ,n ,p ,q \leq r}c_{K m} c_{K n} c_{L p} c_{L q} \left (y_{K m} y_{K n} y_{L p} y_{L q} y_{K m +r} y_{K n +r} y_{L p +r} y_{L q +r}\right )\text{.}\end{align*}

If $A$ is a reflection then
\begin{align*}\mathcal{A}_{ +} &  =\frac{1}{\sqrt{2}} \left (I_{\mathcal{F}} +A\right ) \\
\mathcal{A}_{ -} &  =\frac{1}{\sqrt{2}} \left (I_{\mathcal{F}} -A\right )\end{align*}are idempotents, as
\begin{equation*}\mathcal{A}_{ +}^{2} =\left (\frac{1}{\sqrt{2}} \left (I_{\mathcal{F}} +A\right )\right )^{2} =\frac{1}{2} \left (I_{\mathcal{F}} +2 A +I_{\mathcal{F}}\right ) =\mathcal{A}_{ +}
\end{equation*}and
\begin{equation*}\mathcal{A}_{ -}^{2} =\left (\frac{1}{\sqrt{2}} \left (I_{\mathcal{F}} -A\right )\right )^{2} =\frac{1}{2} \left (I_{\mathcal{F}} -2 A +I_{\mathcal{F}}\right ) =\mathcal{A}_{ -}\text{.}
\end{equation*}Also
\begin{equation*}\mathcal{A}_{ +} \mathcal{A}_{ -} =\mathcal{A}_{ -} \mathcal{A}_{ +} =0 =\frac{1}{\sqrt{2}} \left (I_{\mathcal{F}} +A\right ) \frac{1}{\sqrt{2}} \left (I_{\mathcal{F}} -A\right ) =\frac{1}{2} \left (I_{\mathcal{F}} +A\right ) \left (I_{\mathcal{F}} -A\right )\text{.}
\end{equation*}

\begin{equation} \langle A_{4}^{2} \rangle _{4} =\frac{1}{4} \left (A_{1 2 3 4} A_{1 2 5 6} -A_{1 2 3 5} A_{1 2 4 6} +A_{1 2 3 6} A_{1 2 4 5}\right ) x_{3} \wedge x_{4} \wedge x_{5} \wedge x_{6}
\end{equation}
\begin{equation} \langle A_{2} A_{4} \rangle _{\left (2\right ) 12} = -\frac{1}{4} \left (A_{1 2 5 6} A_{5 6} +A_{1 2 4 6} A_{4 6} +A_{1 2 3 6} A_{3 6} +A_{1 2 4 5} A_{4 5} +A_{1 2 3 5} A_{3 5} +A_{1 2 3 4} A_{3 4}\right ) x_{1} \wedge x_{2}
\end{equation}

\section{Zeon Algebra}
\qquad The algebra formed from the even Fock algebra
\begin{equation}\mathcal{F}_{ +} ={\displaystyle\bigwedge _{0 \leq p \leq 2 r}}\mathcal{F}_{2 p}
\end{equation}is a \emph{Zeon algebra} as
\begin{equation} \langle A \rangle _{2 p} \wedge  \langle B \rangle _{2 q} =( -1)^{2 \left (p +q\right )}  \langle B \rangle _{2 q} \wedge  \langle A \rangle _{2 p} = \langle B \rangle _{2 q} \wedge  \langle A \rangle _{2 p}
\end{equation}and is generated by
\begin{equation}\left \{\left .e_{j k}\right \vert e_{j k} =e_{j} \wedge e_{k} =e_{j} e_{k} ;1 \leq j <k \leq 2 r\right \}\text{and}I_{\mathcal{F}}\text{,}
\end{equation}where
\begin{equation}\left (e_{j k}\right )^{2} =\left (e_{j k}\right )^{2 \wedge } =0
\end{equation}i.e. are nilpotent. \emph{This condition makes it different from a standerd symmetric tensor algebra. }The
first and second order subspaces are generated by the nilpotent elements
\begin{align}\mathcal{F}_{2} &  =\ensuremath{\operatorname*{linspan}}\left \{\left .e_{j k}\right \vert 1 \leq j <k \leq 2 r\right \} \\
\mathcal{F}_{4} &  =\ensuremath{\operatorname*{linspan}}\left \{\left .e_{j k l m}\right \vert e_{j k l m} =e_{j k} e_{l m} ;1 \leq j <k <l <m \leq 2 r\right \}\text{.}\end{align}A general element of $\mathcal{F}_{2}$ is given by
\begin{equation}A_{2} ={\displaystyle\sum _{1 \leq j <k \leq 2 r}}A_{2 j k} e_{j k}
\end{equation}and of $\mathcal{F}_{4}$ by
\begin{equation}A_{4} ={\displaystyle\sum _{1 \leq j <k <l <m \leq 2 r}}A_{4 j k l m} e_{j k l m} ={\displaystyle\sum _{1 \leq j <k <l <m \leq 2 r}}A_{4 j k l m} e_{j k} e_{l m} =\frac{1}{2} {\displaystyle\sum _{\begin{array}{c}1 \leq j <k \leq 2 r \\
1 \leq l <m \leq 2 r\end{array}}}\tilde{A}_{4 j k l m} e_{j k} e_{l m} =\frac{1}{2} {\displaystyle\sum _{1 \leq J ,K \leq r \left (2 r -1\right )}}\tilde{A}_{4 J K} e_{J} e_{K}\text{.}
\end{equation}The basis $\left \{\left .e_{J} e_{K}\right \vert 1 \leq J <K \leq r \left (2 r -1\right )\right \}$ is overcomplete (i.e. linearly dependent) as e.g.
\begin{equation}e_{1 2 3 4} =e_{1 2} e_{3 4} = -e_{1 3} e_{2 4} =e_{1 4} e_{2 3} \equiv e_{\mathbf{1}} e_{\mathbf{14}} = -e_{\mathbf{2}} e_{\mathbf{9}} =e_{\mathbf{3}} e_{\mathbf{8}}
\end{equation}giving
\begin{equation}A_{4 1 2 3 4} =\tilde{A}_{4 1 2 3 4} -\tilde{A}_{4 1 3 2 4} +\tilde{A}_{4 1 4 2 3}\text{.}
\end{equation}

$\lceil $For example $r =3$ then the bases is indexed by 

\begin{equation}\left \{12 -1 ,13 -2 ,14 -3 ,15 -4 ,16 -5 ,23 -6 ,24 -7 ,25 -8 ,26 -9 ,34 -10 ,35 -11 ,36 -12 ,45 -13 ,46 -14 ,56 -15\right \}
\end{equation}
\begin{align} & \left \{12/34 ,12/35 ,12/36 ,12/45 ,12/46 ,12/56 ,13/45 ,13/46 ,13/56 ,14/56 ,23/45 ,23/46 ,23/56 ,24/56 ,34/56\right \} \\
 &  \equiv \left \{1/10 ,1/11 ,1/12 ,1/13 ,1/14 ,1/15 ,2/13 ,2/14 ,2/15 ,3/15 ,4/13 ,4/14 ,4/15 ,5/15 ,6/15\right \}\end{align}
\begin{align}A_{4} &  ={\displaystyle\sum _{1 \leq j <k <l <m \leq 6}}A_{4 j k l m} e_{j k l m} ={\displaystyle\sum _{1 \leq j <k <l <m \leq 6}}A_{4 j k l m} e_{j k} e_{l m} \\
 &  =A_{4 1 2 3 4} e_{1 2} e_{3 4} +A_{4 1 2 3 5} e_{1 2} e_{3 5} +A_{4 1 2 3 6} e_{1 2} e_{3 6} +A_{4 1 2 4 5} e_{1 2} e_{4 5} +A_{4 1 2 4 6} e_{1 2} e_{4 6} \\
 &  +A_{4 1 2 5 6} e_{1 2} e_{5 6} +A_{4 1 3 4 5} e_{1 3} e_{4 5} +A_{4 1 3 4 6} e_{1 3} e_{4 6} +A_{4 1 3 5 6} e_{1 3} e_{5 6} +A_{4 1 4 5 6} e_{1 4} e_{5 6} \\
 &  +A_{4 2 3 4 5} e_{2 3} e_{4 5} +A_{4 1 2 4 6} e_{2 3} e_{4 6} +A_{4 2 3 5 6} e_{2 3} e_{5 6} +A_{4 2 4 5 6} e_{2 4} e_{5 6} +A_{4 3 4 5 6} e_{3 4} e_{5 6}\text{.}\end{align}$\rfloor $ 

The matrix of coefficients $\widetilde{\mathbb{A}}_{4}$ is symmetric and real if $A_{4}$ is self adjoint (as $\left \{e_{j k}\right \}$ are self adjoint), it also has zero diagonal, thus $\ensuremath{\operatorname*{Tr}}\widetilde{\mathbb{A}}_{4} =0$ and can be diagonalized to give
\begin{equation}\mathbb{G}^{t} \widetilde{\mathbb{A}}_{4} \mathbb{G} =\mathbb{A}_{4 D}\text{,}
\end{equation}where
\begin{equation}\mathbb{G}^{t} \mathbb{G} =\mathbb{G} \mathbb{G}^{t} =\mathbb{I}_{r \left (2 r -1\right )}
\end{equation}and
\begin{equation}\widetilde{\mathbb{A}}_{4 D} =\left (\begin{array}{ccc}\alpha _{1} & \mathbb{O} & \mathbb{O} \\
\mathbb{O} & \ddots  & \mathbb{O} \\
\mathbb{O} & \mathbb{O} & \alpha _{r \left (2 r -1\right )}\end{array}\right )\text{.}
\end{equation}Thus
\begin{equation}\widetilde{\mathbb{A}}_{4} =\mathbb{G} \widetilde{\mathbb{A}}_{4 D} \mathbb{G}^{t}
\end{equation}and
\begin{equation}\widetilde{A}_{4 j k l m} ={\displaystyle\sum _{\begin{array}{c}1 \leq p <q \leq 2 r \\
1 \leq r <s \leq 2 r\end{array}}}G_{j k p q} \widetilde{A}_{4 D p q r s} \delta _{\left (p q\right ) \left (r s\right )} G_{l m r s} ={\displaystyle\sum _{1 \leq p <q \leq 2 r}}G_{j k p q} \widetilde{A}_{4 D p q p q} G_{l m p q}
\end{equation}hence
\begin{align}A_{4} &  ={\displaystyle\sum _{1 \leq j <k <l <m \leq 2 r}}\widetilde{A}_{4 j k l m} e_{j k} e_{l m} =\frac{1}{2} {\displaystyle\sum _{\begin{array}{c}1 \leq j <k \leq 2 r \\
1 \leq l <m \leq 2 r\end{array}}}{\displaystyle\sum _{1 \leq p <q \leq 2 r}}G_{j k p q} \widetilde{A}_{4 D p q p q} G_{l m p q} e_{j k} e_{l m} \\
 &  =\frac{1}{2} {\displaystyle\sum _{1 \leq p <q \leq 2 r}}{\displaystyle\sum _{\begin{array}{c}1 \leq j <k \leq 2 r \\
1 \leq l <m \leq 2 r\end{array}}}G_{j k p q} \widetilde{A}_{4 D p q p q} G_{l m p q} e_{j k} e_{l m} \\
 &  ={\displaystyle\sum _{1 \leq p <q \leq 2 r}}\widetilde{A}_{4 D p q p q} E_{p q} E_{p q}\text{,}\end{align}where
\begin{equation}E_{p q} =\frac{1}{\sqrt{2}} {\displaystyle\sum _{1 \leq l <m \leq 2 r}}G_{l m p q} e_{l m}\text{.}
\end{equation}This can be expressed in canonical form as
\begin{equation}E_{p q} ={\displaystyle\sum _{1 \leq l \leq r}}c_{p q l} \varphi _{p q l l +r} ;c_{p q l} \in \mathbb{R}^{ +}\text{,}
\end{equation}where
\begin{equation}\varphi _{p q l l +r} =\varphi _{p q l} \varphi _{p q l +r} =\varphi _{p q l} \wedge \varphi _{p q l +r}\text{.}
\end{equation}If $A_{4 D p q} <0$ then we can replace the self adjoint operator $E_{p q}$ with the anti selfadjoint operator $\widetilde{E_{p q}} =i E_{p q}\text{,}$ while if it is positive $\widetilde{E_{p q}} =E_{p q}$ to obtain
\begin{equation}A_{4} ={\displaystyle\sum _{1 \leq p <q \leq 2 r}}\left \vert \widetilde{A}_{4 D p q}\right \vert  \left (\widetilde{E_{p q}}\right )^{2 \wedge }\text{.}
\end{equation}

\subsection{Nilpotency}
The condition
\begin{equation}\left (A_{4}\right )^{2 \wedge } ={\displaystyle\sum _{\begin{array}{c}1 \leq p <q \leq 2 r \\
1 \leq r <s \leq 2 r\end{array}}}\left \vert A_{4 D p q}\right \vert  \left \vert A_{4 D r s}\right \vert  \widetilde{E_{p q}} \wedge \widetilde{E_{p q}} \wedge \widetilde{E_{r s}} \wedge \widetilde{E_{r s}} =0
\end{equation}gives
\begin{equation}\left (A_{4}\right )^{2 \wedge } ={\displaystyle\sum _{\begin{array}{c}1 \leq p <q \leq 2 r \\
1 \leq r <s \leq 2 r\end{array}}}\left \vert \widetilde{A}_{4 D p q}\right \vert  \left \vert \widetilde{A}_{4 D r s}\right \vert  \widetilde{E_{p q}} \wedge \widetilde{E_{r s}} \wedge \widetilde{E_{p q}} \wedge \widetilde{E_{r s}} ={\displaystyle\sum _{\begin{array}{c}1 \leq p <q \leq 2 r \\
1 \leq r <s \leq 2 r\end{array}}}\left \vert \widetilde{A}_{4 D p q}\right \vert  \left \vert \widetilde{A}_{4 D r s}\right \vert  \left (\widetilde{E_{p q}} \wedge \widetilde{E_{r s}}\right )^{2 \wedge } =0.
\end{equation}

A basis for $\mathcal{F}_{4}$ is
\begin{equation}\left \{\left .e_{j k l m}\right \vert 1 \leq j <k <l <m \leq 2 r\right \}\text{,}
\end{equation}where
\begin{equation}e_{j k l m} =e_{j k} e_{l m} =e_{j} e_{k} e_{l} e_{m}\text{.}
\end{equation}Products of different basis elements of $\mathcal{F}_{2}$ can produce the same basis element of $\mathcal{F}_{4}$ e.g.
\begin{equation}e_{1 2 3 4} =e_{1 2} e_{3 4} = -e_{1 3} e_{2 4} =e_{1 4} e_{2 3}
\end{equation}and if any indices are in common the product is zero e.g.
\begin{equation}e_{1 2} e_{1 3} =0.
\end{equation}\emph{Note} that these properties are $\mathbf{n} \mathbf{o} \mathbf{t}$ true for the operators $\left \{\widetilde{E_{p q}}\right \}$ even though they form a basis for $\mathcal{F}_{2}\text{.}$ 

This analysis can be extended to $\mathcal{F}_{8}$ e.g.
\begin{equation}e_{1 2 3 4 5 6 7 8} =e_{1 2 3 4} e_{5 6 7 8} = -e_{1 2 3 5} e_{4 6 7 8}
\end{equation}etc.. 

$\lceil $The number of such equalities is $\binom{8}{4} =\frac{8 \times 7 \times 6 \times 5}{1 \times 2 \times 3 \times 4} =70\rfloor \text{.}$ 


\begin{align}A_{4} \wedge A_{4} &  ={\displaystyle\sum _{1 \leq j <k <l <m \leq 2 r}}A_{4 j k l m} e_{j k l m} \wedge {\displaystyle\sum _{1 \leq j <k <l <m \leq 2 r}}A_{4 j k l m} e_{j k l m} ={\displaystyle\sum _{\begin{array}{c}1 \leq j <k <l <m \leq 2 r \\
1 \leq n <o <p <q \leq 2 r\end{array}}}A_{4 j k l m} A_{4 n o p q} e_{j k l m} \wedge e_{n o p q} \\
 &  ={\displaystyle\sum _{1 \leq j <k <l <m <n <o <p <q \leq 2 r}}f \left (A_{4 j k l m} ,A_{4 n o p q}\right ) e_{j k l m n o p q} ={\displaystyle\sum _{1 \leq j <k <l <m <n <o <p <q \leq 2 r}}f_{j k l m n o p q} \left (A_{4}\right ) e_{j k l m n o p q}\end{align}$\rfloor $ 

The nilpotency condition is then
\begin{equation}\left (f \left (A_{4 j k l m} ,A_{4 n o p q}\right ) =f_{j k l m n o p q} \left (A_{4}\right ) =f_{j k l m n o p q} \left (\mathbb{A}_{4}\right ) =0\right )\text{or}e_{j k l m n o p q} =0\, \forall j k l m n o p q\text{.}
\end{equation}
\begin{equation}A_{4} \wedge A_{4} ={\displaystyle\sum _{\begin{array}{c}1 \leq J \leq \binom{2 r}{4} \\
1 \leq K \leq \binom{2 r}{4}\end{array}}}A_{4 J} A_{4 K} e_{J} \wedge e_{K} =\frac{1}{2} {\displaystyle\sum _{1 \leq J <K \leq \binom{2 r}{4}}}A_{4 J} A_{4 K} e_{J} \wedge e_{K}
\end{equation}
\begin{align}e_{J} \wedge e_{K} &  =e_{K} \wedge e_{J} \\
e_{J} \wedge e_{K} &  =0 \Leftrightarrow J \cap K \neq \phi  \\
e_{J} \wedge e_{K} &  \neq 0 \Leftrightarrow J \cap K =\phi \end{align}
\begin{equation}e_{J_{1}} \wedge e_{K_{1}} =\left ( -1\right )^{\pi  \left (\sigma \right )} e_{J_{2}} \wedge e_{K_{2}} \Leftrightarrow J_{1} \cup K_{1} =J_{2} \cup K_{2}
\end{equation}where
\begin{equation}\sigma  \left (\left \{J_{1} K_{1}\right \}\right ) =\left \{J_{2} K_{2}\right \}
\end{equation}and $\pi $ is the parity of $\sigma \text{.}$ 

If
\begin{equation}A_{4} ={\displaystyle\sum _{1 \leq j <k <l <m \leq 6}}A_{4 j k l m} e_{j k l m}
\end{equation}then
\begin{equation}A_{4} \wedge A_{4} =0\text{,}
\end{equation}while if 


\begin{align}A_{4} &  ={\displaystyle\sum _{1 \leq j <k <l <m \leq 8}}A_{4 j k l m} e_{j k l m} \\
A_{4} \wedge A_{4} &  \neq 0\end{align}unless
\begin{equation}f_{1 2 3 4 5 6 7 8} \left (\mathbb{A}_{4}\right ) =0.
\end{equation}Zero replacements
\begin{equation} +1234/5678\text{,}
\end{equation}one replacement
\begin{align} &  -1235/4678 , +1236/4578 , -1237/4568 , +1238/4567\text{,} \\
 &  +1245/3678 , -1246/3578 , +1247/3568 , -1248/3567\text{,} \\
 &  -1345/2678 , +1346/2578 , -1347/2568 , +1348/2567\text{,} \\
 &  +2345/1678 , -2346/1578 , +2347/1568 , -2348/1567\text{,}\end{align}two replacements (are self transpose)
\begin{align} & 3478/1256,3468/1257,3467/1258,3458/1267,3457/1268,3456/1278\text{,} \\
 & 2478/1356,2468/1357,2467/1358,2458/1367,2457/1368,2456/1378\text{,} \\
 & 2378/1456,2368/1457,2367/1458,2358/1467,2357/1468,2356/1478\text{,} \\
 & 1478/2356,1468/2357,1467/2358,1458/2367,1457/2368,1456/2378\text{,} \\
 & 1378/2456,1368/2457,1367/2458,1358/2467,1357/2468,1356/2478\text{,} \\
 & 1278/3456,1268/3457,1267/3458,1258/3467,1257/3468,1256/3478\text{,}\end{align}three replacements are the transpose of the one replacements and four replacements the transpose of the zero replacements.


The linearized subscripts give
\begin{align} & f_{1 2 3 4 5 6 7 8} \left (\mathbb{A}_{4}\right ) =f_{1 ,70} \left (\mathbb{A}_{4}\right ) =A_{1} A_{7 0} \\
 &  -A_{2} A_{6 9} +A_{3} A_{6 8} -A_{4} A_{6 7} +A_{5} A_{6 6} \\
 &  +A_{6} A_{6 5} -A_{7} A_{6 4} +A_{8} A_{6 3} -A_{9} A_{6 2} \\
 &  -A_{1 6} A_{5 5} +A_{1 7} A_{5 4} -A_{1 8} A_{5 3} +A_{1 9} A_{5 2} \\
 &  +A_{3 6} A_{3 5} -A_{3 7} A_{3 4} +A_{3 8} A_{3 3} -A_{3 9} A_{3 2} \\
 &  +A_{6 1} A_{1 0} -A_{6 0} A_{1 1} +A_{5 9} A_{1 2} -A_{5 8} A_{1 3} +A_{5 7} A_{1 4} -A_{5 6} A_{1 5} \\
 &  +A_{5 1} A_{2 0} -A_{5 0} A_{2 1} +A_{4 9} A_{2 2} -A_{4 8} A_{2 3} +A_{4 7} A_{2 4} -A_{4 6} A_{2 5} \\
 &  +A_{4 5} A_{2 6} -A_{4 4} A_{2 7} +A_{4 3} A_{2 8} -A_{4 2} A_{2 9} +A_{4 1} A_{3 0} -A_{4 0} A_{3 1}\end{align}and the quadratic subscripts give
\begin{align} & f_{1 2 3 4 5 6 7 8} \left (\mathbb{A}_{4}\right ) =f_{\left (1 ,14\right ) ,\left (23 ,28\right )} \left (\mathbb{A}_{4}\right ) =A_{1 ,14} A_{23 ,28} \\
 &  -A_{1 ,15} A_{20 ,28} +A_{1 ,16} A_{19 ,28} -A_{1 ,17} A_{19 ,27} +A_{1 ,18} A_{19 ,26} \\
 &  +A_{1 ,19} A_{16 ,28} -A_{1 ,20} A_{15 ,28} +A_{1 ,21} A_{15 ,27} -A_{1 ,22} A_{15 ,26} \\
 &  -A_{2 ,19} A_{11 ,28} +A_{2 ,19} A_{10 ,28} -A_{2 ,19} A_{10 ,27} +A_{2 ,19} A_{10 ,26} \\
 &  +A_{8 ,19} A_{5 ,28} -A_{8 ,20} A_{4 ,28} +A_{8 ,21} A_{4 ,27} -A_{8 ,22} A_{4 ,26}\end{align}

\subsection{Compound Indices}
The ordered pairs $\left (j ,k\right )$ can be expressed as a single index $J \left (j ,k\right )\text{,}$where the transformation is
\begin{equation}J \left (j ,k\right ) =2 r \left (j -1\right ) -j \left (j +1\right )/2 +k ;1 \leq j <k \leq 2 r\text{.}
\end{equation}One can check this formula for $r =3$ by: 

\begin{equation}\begin{array}{ccc}j & k & J \left (j ,k\right ) \\
1 & 2 & 6 \left (1 -1\right ) -1 \left (1 +1\right )/2 +2 =1 \\
1 & 3 & 6 \left (1 -1\right ) -1 \left (1 +1\right )/2 +3 =2 \\
1 & 4 & 6 \left (1 -1\right ) -1 \left (1 +1\right )/2 +4 =3 \\
1 & 5 & 6 \left (1 -1\right ) -1 \left (1 +1\right )/2 +5 =4 \\
1 & 6 & 6 \left (1 -1\right ) -1 \left (1 +1\right )/2 +6 =5 \\
2 & 3 & 6 \left (2 -1\right ) -2 \left (2 +1\right )/2 +3 =6 \\
2 & 4 & 6 \left (2 -1\right ) -2 \left (2 +1\right )/2 +4 =7 \\
2 & 5 & 6 \left (2 -1\right ) -2 \left (2 +1\right )/2 +5 =8 \\
2 & 6 & 6 \left (2 -1\right ) -2 \left (2 +1\right )/2 +6 =9 \\
3 & 4 & 6 \left (3 -1\right ) -3 \left (3 +1\right )/2 +4 =10 \\
3 & 5 & 6 \left (3 -1\right ) -3 \left (3 +1\right )/2 +5 =11 \\
3 & 6 & 6 \left (3 -1\right ) -3 \left (3 +1\right )/2 +6 =12 \\
4 & 5 & 6 \left (4 -1\right ) -4 \left (4 +1\right )/2 +5 =13 \\
4 & 6 & 6 \left (4 -1\right ) -4 \left (4 +1\right )/2 +6 =14 \\
5 & 6 & 6 \left (5 -1\right ) -5 \left (5 +1\right )/2 +6 =15\end{array}
\end{equation}

\bigskip One can classify this term into $0 -$substition, $1 -$substition, $2 -$substitions, $3 -$substitions, and $4 -$substitions: 


\begin{enumerate}
\item $0 -$substition :
\begin{equation}12345678
\end{equation}

\item $1 -$substition :
\begin{equation}\binom{4}{1}^{2} =16
\end{equation}
\begin{equation}\begin{array}{cc}1235,1236,1237,1238\text{,} & 4678,5478,5648,5674\text{,} \\
1254,1264,1274,1284\text{,} & 3678,5378,5638,5673\text{,} \\
1534,1634,1734,1834\text{,} & 2678,5278,5628,5672\text{,} \\
5234,6234,7234,8234\text{,} & 1678,5178,5618,5671\text{,}\end{array}
\end{equation}
\begin{align}A_{1 2 3 5} A_{4 6 7 8} x_{1} x_{2} x_{3} x_{5} x_{4} x_{6} x_{7} x_{8} &  = -A_{1 2 3 5} A_{4 6 7 8} x_{1} x_{2} x_{3} x_{4} x_{5} x_{6} x_{7} x_{8} \\
A_{1 2 3 6} A_{5 4 7 8} x_{1} x_{2} x_{3} x_{6} x_{5} x_{4} x_{7} x_{8} &  = -A_{1 2 3 6} A_{4 5 7 8} x_{1} x_{2} x_{3} x_{6} x_{5} x_{4} x_{7} x_{8} =A_{1 2 3 6} A_{4 5 7 8} x_{1} x_{2} x_{3} x_{4} x_{5} x_{6} x_{7} x_{8} \\
A_{1 2 3 7} A_{5 6 4 8} x_{1} x_{2} x_{3} x_{7} x_{5} x_{6} x_{4} x_{8} &  = -A_{1 2 3 7} A_{5 6 4 8} x_{1} x_{2} x_{3} x_{4} x_{5} x_{6} x_{7} x_{8} \\
A_{1 2 3 8} A_{5 6 7 4} x_{1} x_{2} x_{3} x_{8} x_{5} x_{6} x_{7} x_{4} &  = -A_{1 2 3 8} A_{5 6 7 4} x_{1} x_{2} x_{3} x_{5} x_{4} x_{6} x_{7} x_{8} =A_{1 2 3 8} A_{5 6 4 7} x_{1} x_{2} x_{3} x_{4} x_{5} x_{6} x_{7} x_{8}\end{align}


\begin{align}A_{1 2 5 4} A_{3 6 7 8} x_{1} x_{2} x_{3} x_{5} x_{4} x_{6} x_{7} x_{8} &  = -A_{1 2 3 5} A_{4 6 7 8} x_{1} x_{2} x_{3} x_{4} x_{5} x_{6} x_{7} x_{8} \\
A_{1 2 6 4} A_{5 3 7 8} x_{1} x_{2} x_{3} x_{6} x_{5} x_{4} x_{7} x_{8} &  = -A_{1 2 3 6} A_{4 5 7 8} x_{1} x_{2} x_{3} x_{6} x_{5} x_{4} x_{7} x_{8} =A_{1 2 3 6} A_{4 5 7 8} x_{1} x_{2} x_{3} x_{4} x_{5} x_{6} x_{7} x_{8} \\
A_{1 2 7 4} A_{5 6 3 8} x_{1} x_{2} x_{3} x_{7} x_{5} x_{6} x_{4} x_{8} &  = -A_{1 2 3 7} A_{5 6 4 8} x_{1} x_{2} x_{3} x_{4} x_{5} x_{6} x_{7} x_{8} \\
A_{1 2 8 4} A_{5 6 7 3} x_{1} x_{2} x_{3} x_{8} x_{5} x_{6} x_{7} x_{4} &  = -A_{1 2 3 8} A_{5 6 7 4} x_{1} x_{2} x_{3} x_{5} x_{4} x_{6} x_{7} x_{8} =A_{1 2 3 8} A_{5 6 4 7} x_{1} x_{2} x_{3} x_{4} x_{5} x_{6} x_{7} x_{8}\end{align}


\begin{align}A_{1 5 3 4} A_{2 6 7 8} x_{1} x_{2} x_{3} x_{5} x_{4} x_{6} x_{7} x_{8} &  = -A_{1 2 3 5} A_{4 6 7 8} x_{1} x_{2} x_{3} x_{4} x_{5} x_{6} x_{7} x_{8} \\
A_{1 6 3 4} A_{5 2 7 8} x_{1} x_{2} x_{3} x_{6} x_{5} x_{4} x_{7} x_{8} &  = -A_{1 2 3 6} A_{4 5 7 8} x_{1} x_{2} x_{3} x_{6} x_{5} x_{4} x_{7} x_{8} =A_{1 2 3 6} A_{4 5 7 8} x_{1} x_{2} x_{3} x_{4} x_{5} x_{6} x_{7} x_{8} \\
A_{1 7 3 4} A_{5 6 2 8} x_{1} x_{2} x_{3} x_{7} x_{5} x_{6} x_{4} x_{8} &  = -A_{1 2 3 7} A_{5 6 4 8} x_{1} x_{2} x_{3} x_{4} x_{5} x_{6} x_{7} x_{8} \\
A_{1 8 3 4} A_{5 6 7 2} x_{1} x_{2} x_{3} x_{8} x_{5} x_{6} x_{7} x_{4} &  = -A_{1 2 3 8} A_{5 6 7 4} x_{1} x_{2} x_{3} x_{5} x_{4} x_{6} x_{7} x_{8} =A_{1 2 3 8} A_{5 6 4 7} x_{1} x_{2} x_{3} x_{4} x_{5} x_{6} x_{7} x_{8}\end{align}


\begin{align}A_{5 2 3 4} A_{1 6 7 8} x_{1} x_{2} x_{3} x_{5} x_{4} x_{6} x_{7} x_{8} &  = -A_{1 2 3 5} A_{4 6 7 8} x_{1} x_{2} x_{3} x_{4} x_{5} x_{6} x_{7} x_{8} \\
A_{6 2 3 4} A_{5 1 7 8} x_{1} x_{2} x_{3} x_{6} x_{5} x_{4} x_{7} x_{8} &  = -A_{1 2 3 6} A_{4 5 7 8} x_{1} x_{2} x_{3} x_{6} x_{5} x_{4} x_{7} x_{8} =A_{1 2 3 6} A_{4 5 7 8} x_{1} x_{2} x_{3} x_{4} x_{5} x_{6} x_{7} x_{8} \\
A_{7 2 3 4} A_{5 6 1 8} x_{1} x_{2} x_{3} x_{7} x_{5} x_{6} x_{4} x_{8} &  = -A_{1 2 3 7} A_{5 6 4 8} x_{1} x_{2} x_{3} x_{4} x_{5} x_{6} x_{7} x_{8} \\
A_{8 2 3 4} A_{5 6 7 1} x_{1} x_{2} x_{3} x_{8} x_{5} x_{6} x_{7} x_{4} &  = -A_{1 2 3 8} A_{5 6 7 4} x_{1} x_{2} x_{3} x_{5} x_{4} x_{6} x_{7} x_{8} =A_{1 2 3 8} A_{5 6 4 7} x_{1} x_{2} x_{3} x_{4} x_{5} x_{6} x_{7} x_{8}\end{align}

\item $2 -$substitions:
\begin{equation}\binom{4}{2}^{2} =36
\end{equation}
\begin{equation}\begin{array}{cc}1256,1356,1456,2356,2456,3456\text{,} & 3478,2478,2378,1478,1378,1278 \\
1257,1357,1457,2357,2457,3457\text{,} & 3468,2468,2368,1468,1368,1268 \\
1258,1358,1458,2358,2458,3458\text{,} & 3467,2467,2367,1467,1367,1267 \\
1267,1367,1467,2367,2467,3467\text{,} & 3458,2458,2358,1458,1358,1258 \\
1268,1368,1468,2368,2468,3468\text{,} & 3457,2457,2357,1457,1357,1257 \\
1278,1378,1478,2378,2478,3478\text{,} & 3456,2456,2356,1456,1356,1256\end{array}
\end{equation}
\begin{equation}J \left (j ,k\right ) =2 r \left (j -1\right ) -j \left (j +1\right )/2 +k ;1 \leq j <k \leq 2 r\text{.}
\end{equation}
\begin{equation}J \left (5 ,6\right ) =8 \left (5 -1\right ) -5 \left (5 +1\right )/2 +6 =32 -15 +6 =23
\end{equation}
\begin{equation}J \left (3 ,4\right ) =8 \left (3 -1\right ) -3 \left (3 +1\right )/2 +4 =16 -6 +4 =14
\end{equation}
\begin{equation}J \left (7 ,8\right ) =8 \left (7 -1\right ) -7 \left (7 +1\right )/2 +8 =54 -28 +8 =34
\end{equation}\end{enumerate}


Ordered 4 to 1 


\begin{align} & \mathbf{1}1234 \left (\mathbf{1} ,\mathbf{14}\right ) ,\mathbf{2}1235 \left (\mathbf{1} ,\mathbf{15}\right ) ,\mathbf{3}1236 \left (\mathbf{1} ,\mathbf{16}\right ) ,\mathbf{4}1237 \left (\mathbf{1} ,\mathbf{17}\right ) ,\mathbf{5}1238 \left (\mathbf{1} ,\mathbf{18}\right )\text{,} \\
 & \mathbf{6}1245 \left (\mathbf{1} ,\mathbf{19}\right ) ,\mathbf{7}1246 \left (\mathbf{1} ,\mathbf{20}\right ) ,\mathbf{8}1247 \left (\mathbf{1} ,\mathbf{21}\right ) ,\mathbf{9}1248 \left (\mathbf{1} ,\mathbf{22}\right )\text{,} \\
 & \mathbf{10}1256 \left (\mathbf{1} ,\mathbf{23}\right ) ,\mathbf{11}1257 \left (\mathbf{1} ,\mathbf{24}\right ) ,\mathbf{12}1258 \left (\mathbf{1} ,\mathbf{25}\right )\text{,} \\
 & \mathbf{13}1267 \left (\mathbf{1} ,\mathbf{26}\right ) ,\mathbf{14}1268 \left (\mathbf{1} ,\mathbf{27}\right )\text{,} \\
 & \mathbf{15}1278 \left (\mathbf{1} ,\mathbf{28}\right )\text{,}\end{align}
\begin{align} & \mathbf{16}1345 \left (\mathbf{2} ,\mathbf{19}\right ) ,\mathbf{17}1346 \left (\mathbf{2} ,\mathbf{20}\right ) ,\mathbf{18}1347 \left (\mathbf{2} ,\mathbf{21}\right ) ,\mathbf{19}1348 \left (\mathbf{2} ,\mathbf{22}\right )\text{,} \\
 & \mathbf{20}1356 \left (\mathbf{2} ,\mathbf{23}\right ) ,\mathbf{21}1357 \left (\mathbf{2} ,\mathbf{24}\right ) ,\mathbf{22}1358 \left (\mathbf{2} ,\mathbf{25}\right )\text{,} \\
 & \mathbf{23}1367 \left (\mathbf{2} ,\mathbf{26}\right ) ,\mathbf{24}1368 \left (\mathbf{2} ,\mathbf{27}\right )\text{,} \\
 & \mathbf{25}1378 \left (\mathbf{2} ,\mathbf{28}\right )\text{,}\end{align}
\begin{align} & \mathbf{26}1456 \left (\mathbf{3} ,\mathbf{23}\right ) ,\mathbf{27}1457 \left (\mathbf{3} ,\mathbf{24}\right ) ,\mathbf{28}1458 \left (\mathbf{3} ,\mathbf{25}\right )\text{,} \\
 & \mathbf{29}1467 \left (\mathbf{3} ,\mathbf{26}\right ) ,\mathbf{30}1468 \left (\mathbf{3} ,\mathbf{27}\right )\text{,} \\
 & \mathbf{31}1478 \left (\mathbf{3} ,\mathbf{28}\right )\text{,}\end{align}
\begin{align} & \mathbf{32}1567 \left (\mathbf{4} ,\mathbf{26}\right ) ,\mathbf{33}1568 \left (\mathbf{4} ,\mathbf{27}\right )\text{,} \\
 & \mathbf{34}1578 \left (\mathbf{4} ,\mathbf{28}\right )\text{,}\end{align}
\begin{equation}\mathbf{35}1678 \left (\mathbf{5} ,\mathbf{28}\right )\text{,}
\end{equation}
\begin{align} & \mathbf{36}2345 \left (\mathbf{8} ,\mathbf{19}\right ) ,\mathbf{37}2346 \left (\mathbf{8} ,\mathbf{20}\right ) ,\mathbf{38}2347 \left (\mathbf{8} ,\mathbf{21}\right ) ,\mathbf{39}2348 \left (\mathbf{8} ,\mathbf{22}\right )\text{,} \\
 & \mathbf{40}2356 \left (\mathbf{8} ,\mathbf{23}\right ) ,\mathbf{41}2357 \left (\mathbf{8} ,\mathbf{24}\right ) ,\mathbf{42}2358 \left (\mathbf{8} ,\mathbf{25}\right )\text{,} \\
 & \mathbf{43}2367 \left (\mathbf{8} ,\mathbf{26}\right ) ,\mathbf{44}2368 \left (\mathbf{8} ,\mathbf{27}\right )\text{,} \\
 & \mathbf{45}2378 \left (\mathbf{8} ,\mathbf{28}\right )\text{,}\end{align}
\begin{align} & \mathbf{46}2456 \left (\mathbf{9} ,\mathbf{23}\right ) ,\mathbf{47}2457 \left (\mathbf{9} ,\mathbf{24}\right ) ,\mathbf{48}2458 \left (\mathbf{9} ,\mathbf{25}\right )\text{,} \\
 & \mathbf{49}2467 \left (\mathbf{9} ,\mathbf{26}\right ) ,\mathbf{50}2468 \left (\mathbf{9} ,\mathbf{27}\right )\text{,} \\
 & \mathbf{51}2478 \left (\mathbf{9} ,\mathbf{28}\right )\text{,}\end{align}
\begin{align} & \mathbf{52}2567 \left (\mathbf{10} ,\mathbf{26}\right ) ,\mathbf{53}2568 \left (\mathbf{10} ,\mathbf{27}\right )\text{,} \\
 & \mathbf{54}2578 \left (\mathbf{10} ,\mathbf{28}\right )\text{,}\end{align}
\begin{equation}\mathbf{55}2678 \left (\mathbf{11} ,\mathbf{28}\right )\text{,}
\end{equation}
\begin{align} & \mathbf{56}3456 \left (\mathbf{14} ,\mathbf{23}\right ) ,\mathbf{57}3457 \left (\mathbf{14} ,\mathbf{24}\right ) ,\mathbf{58}3458 \left (\mathbf{14} ,\mathbf{25}\right )\text{,} \\
 & \mathbf{59}3467 \left (\mathbf{14} ,\mathbf{26}\right ) ,\mathbf{60}3468 \left (\mathbf{14} ,\mathbf{27}\right )\text{,} \\
 & \mathbf{61}3478 \left (\mathbf{14} ,\mathbf{28}\right )\text{,}\end{align}
\begin{align} & \mathbf{62}3567 \left (\mathbf{15} ,\mathbf{26}\right ) ,\mathbf{63}3568 \left (\mathbf{15} ,\mathbf{27}\right )\text{,} \\
 & \mathbf{64}3578 \left (\mathbf{15} ,\mathbf{28}\right )\text{,}\end{align}
\begin{equation}\mathbf{65}3678 \left (\mathbf{15} ,\mathbf{28}\right )\text{,}
\end{equation}
\begin{align} & \mathbf{66}4567 \left (\mathbf{19} ,\mathbf{26}\right ) ,\mathbf{67}4568 \left (\mathbf{19} ,\mathbf{27}\right )\text{,} \\
 & \mathbf{68}4578 \left (\mathbf{19} ,\mathbf{28}\right )\text{,}\end{align}
\begin{equation}\mathbf{69}4678 \left (\mathbf{20} ,\mathbf{28}\right )\text{,}
\end{equation}
\begin{equation}\mathbf{70}5678 \left (\mathbf{21} ,\mathbf{28}\right )\text{.}
\end{equation}
\begin{equation}\binom{8}{4} =\frac{8 \times 7 \times 6 \times 5}{1 \times 2 \times 3 \times 4} =70
\end{equation}Ordered 2 to 1
\begin{align} & \mathbf{1}12 ,\mathbf{2}13 ,\mathbf{3}14 ,\mathbf{4}15 ,\mathbf{5}16 ,\mathbf{6}17 ,\mathbf{7}18\text{,} \\
 & \mathbf{8}23 ,\mathbf{9}24 ,\mathbf{10}25 ,\mathbf{11}26 ,\mathbf{12}27 ,\mathbf{13}28\text{,} \\
 & \mathbf{14}34 ,\mathbf{15}35 ,\mathbf{16}36 ,\mathbf{17}37 ,\mathbf{18}38\text{,} \\
 & \mathbf{19}45 ,\mathbf{20}46 ,\mathbf{21}47 ,\mathbf{22}48\text{,} \\
 & \mathbf{23}56 ,\mathbf{24}57 ,\mathbf{25}58\text{,} \\
 & \mathbf{26}67 ,\mathbf{27}68\text{,} \\
 & \mathbf{28}78.\end{align}
\begin{equation}\binom{8}{2} =\frac{8 \times 7}{1 \times 2} =28
\end{equation}

\subsection{Theorem}
\begin{theorem}
$A_{4} \wedge A_{4} =0$ Iff $A_{4} ={\displaystyle\sum _{1 \leq j <k <l <m \leq 7}}A_{4 j k l m} e_{j k l m}$ 


\begin{proof}
The if part is easy as $\left \{i j k l\right \} \cap \left \{m n o p\right \} \neq \phi $ for this range of subscripts thus the products $e_{j k l m} \wedge e_{m n o p} =0$ $ \forall \left \{i ,j ,k ,l ,m ,n ,o ,p\right \} \in \left \{1 ,2 ,3 ,4 ,5 ,6 ,7\right \}\text{.}$ 

The only if part: Let $A_{4} \wedge A_{4} =0$ and $A_{4} ={\displaystyle\sum _{1 \leq j <k <l <m \leq 8}}A_{4 j k l m} e_{j k l m}$ then $f_{1 2 3 4 5 6 7 8} \left (\mathbb{A}_{4}\right ) =0.$ 

This can easily be shown to be possible, let $f_{1 2 3 4 5 6 7 8} \left (\mathbb{A}_{4}\right ) =A_{1} A_{7 0} -A_{2} A_{6 9} =0 \Leftrightarrow A_{1} A_{7 0} =A_{2} A_{6 9} \Leftrightarrow \frac{A_{1}}{A_{2}} =\frac{A_{6 9}}{A_{7 0}} =\alpha \text{,}$where 


\begin{align}A_{4} &  =A_{1} e_{1 2 3 4} +A_{2} e_{1 2 3 5} +A_{6 9} e_{4 6 7 8} +A_{7 0} e_{5 6 7 8} \\
 &  =\alpha  A_{2} e_{1 2 3 4} +A_{2} e_{1 2 3 5} +\alpha  A_{6 9} e_{4 6 7 8} +A_{6 9} e_{5 6 7 8} \\
 &  =A_{2} \left (\alpha  e_{1 2 3 4} +e_{1 2 3 5}\right ) +A_{6 9} \left (\alpha  e_{4 6 7 8} +e_{5 6 7 8}\right )\end{align}
\begin{align} & \left (\overset{}{} A_{2} \left (\alpha  e_{1 2 3 4} +e_{1 2 3 5}\right ) +A_{6 9} \left (\alpha  e_{4 6 7 8} +e_{5 6 7 8}\right )\right ) \wedge \left (\overset{}{} A_{2} \left (\alpha  e_{1 2 3 4} +e_{1 2 3 5}\right ) +A_{6 9} \left (\alpha  e_{4 6 7 8} +e_{5 6 7 8}\right )\right ) \\
 &  =\alpha  A_{2} e_{1 2 3 4} \wedge A_{6 9} e_{5 6 7 8} +A_{2} e_{1 2 3 5} \wedge A_{6 9} \alpha  e_{4 6 7 8} +A_{6 9} \alpha  e_{4 6 7 8} \wedge A_{2} e_{1 2 3 5} +A_{6 9} e_{5 6 7 8} \wedge \alpha  A_{2} e_{1 2 3 4} \\
 &  =\alpha  A_{2} A_{6 9} e_{1 2 3 4} \wedge e_{5 6 7 8} +\alpha  A_{2} A_{6 9} e_{1 2 3 5} \wedge e_{4 6 7 8} +\alpha  A_{6 9} A_{2} e_{4 6 7 8} \wedge e_{1 2 3 5} +\alpha  A_{2} A_{6 9} e_{5 6 7 8} \wedge e_{1 2 3 4} \\
 &  =2 \alpha  A_{2} A_{6 9} \left (e_{1 2 3 4} \wedge e_{5 6 7 8} +e_{1 2 3 5} \wedge e_{4 6 7 8}\right ) =2 \alpha  A_{2} A_{6 9} \left (e_{1 2 3 4 5 6 7 8} -e_{1 2 3 4 5 6 7 8}\right ) =0\end{align}
\end{proof}


\end{theorem}


\begin{remark}
Unlike the tensor algebra the Clifford algebra is not Z-graded, since two vectors can multiply to a scalar. Nevertheless it is Z2-graded,
and this Z2-grading is important. We can define an algebra automorphism {\small\fbox{3BB}}  on Cl(Q) by taking {\small\fbox{3BB}} (v) = -v for v {\small\fbox{2208}}  V and extending this to
be an algebra automorphism. 
\end{remark}


\end{document}